% =====================================================================
%  main.tex — La gramática: de las nociones teóricas al lenguaje de todos los días
%  Compilar:  latexmk            (usa .latexmkrc: LuaLaTeX + biber)
%         o:  latexmk -xelatex main.tex
% =====================================================================
\documentclass[
  fontsize=10.5pt,
  twoside,
  open=right,
  chapterprefix=true,
  numbers=noendperiod,
  bibliography=totoc,
  toc=flat,
]{scrbook}

% =====================================================================
%  config/paquetes.tex — Paquetes del manual
%  Motor: XeLaTeX o LuaLaTeX  |  Bibliografía: biblatex + biber
%  Solo carga paquetes y fija opciones de paquete; el diseño y los
%  entornos propios están en config/estilo.tex.
% =====================================================================

% --- Detección de motor ----------------------------------------------
\RequirePackage{iftex}
\RequireTUTeX  % aborta con pdfLaTeX: el proyecto requiere XeLaTeX o LuaLaTeX

% --- Idioma, fuentes y tipografía -------------------------------------
\usepackage{fontspec}
\usepackage[spanish,
            es-tabla,          % «Tabla» en lugar de «Cuadro»
            es-noshorthands,   % evita conflictos de " < > con TikZ/forest
            es-nolists,        % no altera las listas (gb4e, enumitem)
            es-noquoting]{babel}
\usepackage{csquotes}     % \enquote{...}: comillas “ ” y ‘ ’ como en los originales
\DeclareQuoteStyle{manual}
  {\textquotedblleft}{\textquotedblright}
  {\textquoteleft}{\textquoteright}
\setquotestyle{manual}

\ifLuaTeX
  \usepackage[protrusion=true,expansion=true]{microtype}
\else
  \usepackage[protrusion=true]{microtype}
\fi

% --- Página, encabezados y color --------------------------------------
\usepackage{geometry}
\usepackage[automark,headsepline=0.4pt]{scrlayer-scrpage}
\usepackage[dvipsnames,svgnames]{xcolor}
\usepackage{graphicx}

% --- Tablas -----------------------------------------------------------
\usepackage{array}
\usepackage{booktabs}
\usepackage{tabularx}
\usepackage{multirow}
\usepackage{longtable}
\usepackage{makecell}

% --- Listas y recuadros -----------------------------------------------
\usepackage{enumitem}
\usepackage{multicol}     % glosario general a dos columnas
\usepackage[most]{tcolorbox}
\usepackage{etoolbox}
\usepackage{xparse}

% --- Lingüística ------------------------------------------------------
\usepackage{gb4e}        % ejemplos numerados y glosas (\gll ... \glt)
\noautomath              % gb4e no debe activar _ y ^ fuera de modo matemático
\usepackage{leipzig}     % abreviaturas de glosa (versalitas): \Acc, \Dat, \Pl...
\usepackage{forest}      % árboles sintácticos
\useforestlibrary{linguistics}
\usepackage{tikz}

% --- Bibliografía -----------------------------------------------------
\usepackage[backend=biber,
            style=apa,
            sortcites=true,
            maxcitenames=2,
            uniquename=false,
            uniquelist=false]{biblatex}
\DeclareLanguageMapping{spanish}{spanish-apa}

% --- Hipervínculos (al final) -----------------------------------------
\usepackage[hidelinks,
            pdfusetitle,
            bookmarksnumbered,
            bookmarksopen]{hyperref}
\usepackage{bookmark}

% =====================================================================
%  config/estilo.tex — Diseño y entornos propios del manual
%  Requiere config/paquetes.tex cargado antes.
% =====================================================================
\makeatletter

% ---------------------------------------------------------------------
% 1. Página
% ---------------------------------------------------------------------
\geometry{
  paperwidth  = 170mm,
  paperheight = 240mm,
  inner       = 20mm,
  outer       = 22mm,
  top         = 24mm,
  bottom      = 26mm,
  headsep     = 7mm,
  footskip    = 11mm,
}
\setlength{\parindent}{1.2em}
\linespread{1.06}
\frenchspacing
\setlength{\emergencystretch}{1.5em}
\raggedbottom   % sin estirar el espacio vertical entre ejemplos y recuadros

% ---------------------------------------------------------------------
% 2. Fuentes (Libertinus: incluida en TeX Live, con versalitas reales)
% ---------------------------------------------------------------------
\setmainfont{LibertinusSerif}[
  Extension      = .otf,
  UprightFont    = *-Regular,
  ItalicFont     = *-Italic,
  BoldFont       = *-Semibold,
  BoldItalicFont = *-SemiboldItalic,
  Ligatures      = TeX,
  Numbers        = Lining]
\setsansfont{LibertinusSans}[
  Extension      = .otf,
  UprightFont    = *-Regular,
  ItalicFont     = *-Italic,
  BoldFont       = *-Bold,
  Ligatures      = TeX,
  Scale          = MatchLowercase]
\newfontfamily\versalitasespaciadas{LibertinusSans}[
  Extension = .otf, UprightFont = *-Regular, BoldFont = *-Bold,
  LetterSpace = 12, Scale = MatchLowercase]

% ---------------------------------------------------------------------
% 3. Colores
% ---------------------------------------------------------------------
\definecolor{acento}{HTML}{1B5E6B}       % petróleo: títulos, ejercicios
\definecolor{acentoclaro}{HTML}{EAF2F3}  % fondo de recuadros
\definecolor{acento2}{HTML}{A4502A}      % terracota: ejemplos resueltos
\definecolor{gristexto}{HTML}{555555}

% ---------------------------------------------------------------------
% 4. Títulos (KOMA-Script)
% ---------------------------------------------------------------------
\setkomafont{disposition}{\sffamily\bfseries\color{acento}}
\setkomafont{chapterprefix}{\normalsize\versalitasespaciadas\bfseries\color{acento}}
\setkomafont{chapter}{\sffamily\bfseries\LARGE\color{black}}
\setkomafont{section}{\sffamily\bfseries\Large}
\setkomafont{subsection}{\sffamily\bfseries\large}
\setkomafont{subsubsection}{\sffamily\bfseries\normalsize}

\renewcommand*{\chapterformat}{\MakeUppercase{\chaptername}~\thechapter}
\renewcommand*{\chapterlinesformat}[3]{%
  \Ifstr{#1}{chapter}{%
    \raggedright
    \Ifstr{#2}{}{}{{\usekomafont{chapterprefix}#2}\par\nobreak\vspace{.5\baselineskip}}%
    #3\par
  }{\@hangfrom{#2}{#3}}}
\RedeclareSectionCommand[beforeskip=10mm, afterskip=4mm]{chapter}
\RedeclareSectionCommand[beforeskip=-2.2ex plus -.5ex minus -.2ex,
                         afterskip=1.1ex plus .2ex]{section}
\RedeclareSectionCommand[beforeskip=-1.8ex plus -.4ex minus -.2ex,
                         afterskip=.8ex plus .1ex]{subsection}
\setcounter{secnumdepth}{\subsubsectionnumdepth}
% Secciones numeradas con el número de capítulo: 2.1, 2.3.1
\setcounter{tocdepth}{\sectiontocdepth}

% \capitulo[título corto]{Título}{Autor/es}
%   Abre un capítulo numerado con firma de autor (en el cuerpo y en el índice).
\NewDocumentCommand{\capitulo}{o m m}{%
  \IfNoValueTF{#1}{\chapter{#2}}{\chapter[#1]{#2}}%
  \label{cap:\capclave}%
  \xappto\capsincluidos{\capclave,}%
  \gdef\autorcapitulo{#3}%
  {\raggedright\normalfont\large\itshape #3\par}%
  \vspace{2mm}\noindent{\color{acento}\rule{\linewidth}{.6pt}}\par
  \vspace{7mm}%
  \addtocontents{toc}{\protect\autorentoc{#3}}%
  \@afterindentfalse\@afterheading}
\newcommand*{\autorcapitulo}{}
\newcommand*{\autorentoc}[1]{%
  \par\nopagebreak
  {\leftskip=1.5em\relax\rightskip=\@tocrmarg\relax
   \small\itshape\color{gristexto}#1\par}%
  \addvspace{2pt}}

% Secciones finales de capítulo (no numeradas, fuera del índice,
% con marcador PDF): Lecturas, Solucionario, Glosario, Referencias.
\newcounter{seccionfinal}
\newcommand*{\seccionfinal}[1]{%
  \stepcounter{seccionfinal}%
  \addsec*{#1}%
  \pdfbookmark[1]{#1}{seccionfinal.\theseccionfinal}%
  \markright{#1}}

% ---------------------------------------------------------------------
% 5. Encabezados y pies
% ---------------------------------------------------------------------
\setkomafont{pageheadfoot}{\small\sffamily\color{gristexto}}
\setkomafont{pagenumber}{\small\sffamily\bfseries\color{acento}}
\renewcommand*{\chaptermarkformat}{\thechapter.\enskip}
\renewcommand*{\sectionmarkformat}{\thesection\enskip}
\automark[section]{chapter}
\clearpairofpagestyles
\ohead{\pagemark}
\ihead{\headmark}
\cfoot[\pagemark]{}   % páginas «plain» (apertura de capítulo)
\KOMAoptions{headsepline=false}

% ---------------------------------------------------------------------
% 6. Notas al pie, epígrafes de tablas y figuras
% ---------------------------------------------------------------------
\deffootnote[1.2em]{1.2em}{1em}{\textsuperscript{\thefootnotemark}\,}
\setfootnoterule[.4pt]{3cm}
\addtokomafont{footnote}{\linespread{1}\selectfont}

\renewcommand*{\thetable}{\arabic{table}}     % «Tabla 1», reinicia por capítulo
\renewcommand*{\thefigure}{\arabic{figure}}
\renewcommand*{\captionformat}{.\enskip}
\setkomafont{caption}{\small}
\setkomafont{captionlabel}{\small\sffamily\bfseries\color{acento}}
\setcapindent{0pt}
\KOMAoptions{captions=tablesignature}         % epígrafe debajo de la tabla

% Tablas: columnas centradas y cuerpo en \small
\newcolumntype{C}{>{\centering\arraybackslash}X}
\newcolumntype{L}{>{\raggedright\arraybackslash}X}
\newcolumntype{Y}[1]{>{\centering\arraybackslash}p{#1}}
\AtBeginEnvironment{tabularx}{\small\renewcommand{\arraystretch}{1.2}}
\AtBeginEnvironment{tabular}{\small\renewcommand{\arraystretch}{1.2}}
% Valores de los cuadros de ejercicios
\newcommand*{\mas}{+}
\newcommand*{\menos}{\textendash}
\newcommand*{\masmenos}{+/\textendash}
% Encabezado de columna en cursiva (como en los originales)
\newcommand*{\enc}[1]{\textit{#1}}

% ---------------------------------------------------------------------
% 7. Ejemplos (gb4e)
%    Texto:      numeración continua por capítulo, entre paréntesis: (1), (2a)
%    Ejercicios: numeración propia, reinicia en cada ejercicio, sin
%                paréntesis: 1., 2. a.
% ---------------------------------------------------------------------
\counterwithin*{exx}{chapter}
\newcounter{ejx}
\newcommand*{\exnumcounter}{exx}
\newcommand*{\exlabelformat}[1]{(#1)}

\patchcmd{\@exe}{(\thexnumi)}{\exlabelformat{\thexnumi}}{}%
  {\PackageError{estilo}{No se pudo adaptar gb4e (\string\@exe)}{}}
\patchcmd{\@exe}{\setcounter{xnumi}{\value{exx}}}%
                {\setcounter{xnumi}{\value{\exnumcounter}}}{}%
  {\PackageError{estilo}{No se pudo adaptar gb4e (contador)}{}}
\patchcmd{\endexe}{\setcounter{exx}{\value{xnumi}}}%
                  {\setcounter{\exnumcounter}{\value{xnumi}}}{}%
  {\PackageError{estilo}{No se pudo adaptar gb4e (\string\endexe)}{}}

% Dentro de ejercicios, cada bloque exe reinicia la numeración
\pretocmd{\exe}{\ifdefstring{\exnumcounter}{ejx}{\setcounter{ejx}{0}}{}}{}%
  {\PackageError{estilo}{No se pudo adaptar gb4e (\string\exe)}{}}

% Modo «ejercicio»: contador ejx, etiqueta sin paréntesis
\newcommand*{\modoejercicio}{%
  \setcounter{ejx}{0}%
  \renewcommand*{\exnumcounter}{ejx}%
  \renewcommand*{\exlabelformat}[1]{##1.}%
  \exewidth{99.}}

% Remisiones: \exref{ex:ambig} → (7);  \exrefs{ex:a}{ex:b} → (7)-(9)
\newcommand*{\exref}[1]{(\ref{#1})}
\newcommand*{\exrefs}[2]{(\ref{#1})-(\ref{#2})}

% Juicio en línea: \juicio{*}{haiga} → *haiga
\newcommand*{\juicio}[2]{#1#2}

% Abreviaturas de glosa propias (leipzig): \Dom
\newleipzig{dom}{dom}{marcado diferencial de objeto}

% Formas mencionadas y significados
\newcommand*{\forma}[1]{\textit{#1}}               % lengua objeto
\newcommand*{\sig}[1]{\textquoteleft#1\textquoteright} % ‘significado’
\newcommand*{\termino}[1]{\textit{#1}}             % término técnico introducido
\newcommand*{\concepto}[1]{\textsc{\MakeLowercase{#1}}} % VERSALITA de las citas

% Bloques de cita (APA, ≥ 40 palabras)
\newenvironment{cita}
  {\begin{quote}\small\linespread{1}\selectfont}
  {\end{quote}}

% ---------------------------------------------------------------------
% 8. Árboles (forest)
% ---------------------------------------------------------------------
\forestset{
  gramatica/.style={
    for tree={
      parent anchor=south, child anchor=north,
      s sep=6mm, l sep=4mm, inner sep=1pt,
      font=\small, align=center,
    },
    baseline,   % alinea la raíz con la etiqueta del ejemplo
  },
  triangulo/.style={roof},
}
\NewDocumentEnvironment{arbol}{+b}
  {\par\medskip\begin{center}\begin{forest}gramatica #1\end{forest}\end{center}\medskip}{}

% ---------------------------------------------------------------------
% 9. Ejercicios
%    \begin{ejercicio}[subtítulo opcional] ... \end{ejercicio}
%    \begin{resuelto}[Ejemplo resuelto] ... \end{resuelto}
% ---------------------------------------------------------------------
\newcounter{ejercicio}[chapter]
\renewcommand*{\theejercicio}{\arabic{ejercicio}}

\newtcolorbox{cajaejercicio}[1]{%
  enhanced, breakable, blanker,
  borderline west={2.2pt}{0pt}{acento},
  left=4mm, right=0mm, top=0mm, bottom=1mm,
  before skip=6mm plus 1mm, after skip=6mm plus 1mm,
  fontupper=\small,
  title={#1}, coltitle=acento,
  fonttitle=\sffamily\bfseries,
  attach title to upper={\par\smallskip},
  parbox=false,
}
\NewDocumentEnvironment{ejercicio}{o}
  {\par\needspace{8\baselineskip}\refstepcounter{ejercicio}%
   \begin{cajaejercicio}{Ejercicio~\theejercicio%
       \IfValueT{#1}{\enskip{\mdseries\color{gristexto}·\enskip #1}}}%
   \modoejercicio\setlength{\parindent}{0pt}\setlength{\parskip}{.5\baselineskip}}
  {\end{cajaejercicio}}

\NewDocumentEnvironment{resuelto}{O{Ejemplo resuelto}}
  {\par\smallskip\small\noindent
   {\sffamily\bfseries\color{acento2}#1.}\enskip\ignorespaces}
  {\par}

% Listas del texto: espaciado compacto
\setlist[itemize]{itemsep=.25\baselineskip, topsep=.4\baselineskip,
                  parsep=0pt, leftmargin=1.6em}
\setlist[itemize,2]{label=--, topsep=.2\baselineskip}
\setlist[enumerate]{itemsep=.25\baselineskip, topsep=.4\baselineskip, parsep=0pt}

% Preguntas o consignas internas del ejercicio: a), b)...
\newlist{consignas}{enumerate}{1}
\setlist[consignas]{label=\alph*), leftmargin=1.8em, topsep=.3\baselineskip,
                    itemsep=.2\baselineskip}

% ---------------------------------------------------------------------
% 10. Recuadros destacados
%    \begin{recuadro}[Título] ... \end{recuadro}  (con título)
%    \begin{destacado} ... \end{destacado}         (sin título)
% ---------------------------------------------------------------------
\newtcolorbox{recuadro}[1][]{%
  enhanced, breakable,
  colback=acentoclaro, colframe=acento,
  boxrule=0pt, toprule=1.4pt, bottomrule=0pt, leftrule=0pt, rightrule=0pt,
  arc=0pt, outer arc=0pt,
  left=5mm, right=5mm, top=3mm, bottom=3mm,
  before skip=5mm, after skip=5mm,
  fontupper=\small,
  title={#1}, coltitle=acento, fonttitle=\sffamily\bfseries,
  attach title to upper={\par\smallskip},
}
\newtcolorbox{destacado}{%
  enhanced, breakable, colback=acentoclaro, frame hidden,
  arc=0pt, left=5mm, right=5mm, top=2.5mm, bottom=2.5mm,
  before skip=4mm, after skip=4mm,
}

% ---------------------------------------------------------------------
% 11. Secciones finales del capítulo
% ---------------------------------------------------------------------
% Clave del capítulo para filtrar la bibliografía: cap01, cap02...
\newcommand*{\capclave}{\two@digits{\value{chapter}}}

\defbibheading{seccionfinal}[]{\seccionfinal{#1}}

% Bibliografía general (al final del libro): entradas de los capítulos
% incluidos (palabras clave capNN y lectNN) más todo lo citado con \cite.
% Las entradas marcadas «nobib» (duplicados para lecturas) se omiten.
\newcommand*{\capsincluidos}{}
\DeclareBibliographyCategory{citados}
\AtEveryCitekey{\addtocategory{citados}{\thefield{entrykey}}}
\newcommand*{\bib@chequearcap}[1]{%
  \ifkeyword{cap#1}{\def\bib@ok{1}}{}%
  \ifkeyword{lect#1}{\def\bib@ok{1}}{}}
\defbibcheck{libro}{%
  \ifkeyword{nobib}
    {\skipentry}
    {\def\bib@ok{0}%
     \ifcategory{citados}{\def\bib@ok{1}}{}%
     \forcsvlist{\bib@chequearcap}{\capsincluidos}%
     \ifnum\bib@ok=0 \skipentry\fi}}
\newcommand*{\bibliografiageneral}{%
  \printbibliography[check=libro, heading=bibintoc,
                     title={Referencias bibliográficas}]}

\newcommand*{\lecturasrecomendadas}{%
  \begingroup
  \edef\x{\noexpand\printbibliography[keyword=lect\capclave,
          heading=seccionfinal, title={Lecturas recomendadas}]}%
  \x\endgroup}

% Tipos de editor en español: editortype = {coordinator|director|compiler}
%   → (Coord./Coords.), (Dir./Dirs.), (Comp./Comps.)
\csdef{tipoeditor@coordinator}{Coord\adddot}
\csdef{tipoeditor@coordinators}{Coords\adddot}
\csdef{tipoeditor@director}{Dir\adddot}
\csdef{tipoeditor@directors}{Dirs\adddot}
\csdef{tipoeditor@compiler}{Comp\adddot}
\csdef{tipoeditor@compilers}{Comps\adddot}
\renewbibmacro*{apaeditorstrg}[1]{%
  \iffieldundef{#1type}
    {\ifthenelse{\value{#1}>1\OR\ifandothers{#1}}
       {\bibcpstring{editors}}
       {\bibcpstring{editor}}}
    {\edef\tipoed@{\thefield{#1type}}%
     \ifthenelse{\value{#1}>1\OR\ifandothers{#1}}
       {\ifcsdef{tipoeditor@\tipoed@ s}
          {\csuse{tipoeditor@\tipoed@ s}}
          {\bibcpstring{type\thefield{#1type}s}}}
       {\ifcsdef{tipoeditor@\tipoed@}
          {\csuse{tipoeditor@\tipoed@}}
          {\bibcpstring{type\thefield{#1type}}}}}}

% APA en español: «y» en lugar de «&»; «Obra original publicada en»;
% partícula pospuesta («Saussure, F. de»)
\ExecuteBibliographyOptions{useprefix=false}
\DeclareDelimFormat[bib,biblist]{finalnamedelim}{%
  \ifthenelse{\value{listcount}>\maxprtauth}
    {}
    {\addspace y\space}}
\DeclareDelimFormat[parencite,fullcite,fullcitebib]{finalnamedelim}{\addspace y\space}
\renewbibmacro*{origyear}{%
  \ifboolexpr{%
    not test {\iffieldundef{origyear}}
    and not test {\iffieldundef{labelyear}}
    and not test {\iffieldsequal{labelyear}{origyear}}}
    {\printtext[parens]{Obra original publicada en\addnbspace\printorigdate}%
     \renewcommand*{\finentrypunct}{\relax}}
    {}}

% Dos autores sin coma antes de «y» (Bosque, I. y Demonte, V.); con tres
% o más, coma entre todos salvo antes de «y».
\renewcommand*{\apablx@ifrevnameappcomma}{%
  \iftoggle{apablx@wantcomma}
    {\ifboolexpr{test {\ifnumless{\value{listcount}}{\value{liststop}-1}}
                 or test \ifmorenames}}
    {\@secondoftwo}}

% Capítulos de libro: (Vol. 2, pp. 23-45), como en APA 7
\renewbibmacro*{addinfo1}{%
  \printtexte[parens]{%
    \usebibmacro{editor+trans+narr}%
    \setunit*{\addsemicolon\addspace}%
    \printfield{edition}%
    \setunit*{\addcomma\addspace}%
    \ifthenelse{\ifentrytype{report}\OR\ifentrytype{unpublished}\OR
                \ifentrytype{dataset}\OR\ifentrytype{software}\OR
                \ifentrytype{audio}\OR\ifentrytype{patent}\OR
                \ifentrytype{image}}{\usebibmacro{addinfo3}}{}%
    \setunit*{\addcomma\addspace}%
    \ifthenelse{\ifentrytype{book}\OR\ifentrytype{inbook}}{\printfield{volumes}}{}%
    \setunit*{\addcomma\addspace}%
    \ifthenelse{\ifentrytype{book}\OR\ifentrytype{inbook}}{\printfield{volume}}{}%
    \setunit*{\addcomma\addspace}%
    \ifentrytype{inbook}{\printfield{pages}}{}}}

% Bibliografía: sangría francesa y cuerpo menor
\setlength{\bibhang}{1.5em}
\renewcommand*{\bibfont}{\small\raggedright}

% Solucionario
\newcommand*{\solucionario}{\seccionfinal{Solucionario}}
\NewDocumentEnvironment{solucion}{m}
  {\par\addvspace{\medskipamount}%
   \needspace{4\baselineskip}%
   \noindent{\sffamily\bfseries\color{acento}Ejercicio~#1}\par\nobreak
   \small\modoejercicio\setlength{\parindent}{0pt}\setlength{\parskip}{.4\baselineskip}}
  {\par}

% Glosario (por capítulo y general)
\newlist{glosario}{description}{4}
\setlist[glosario]{font=\normalfont\bfseries\itshape, leftmargin=1.5em,
                   labelindent=0pt, itemsep=.35\baselineskip,
                   style=sameline, before=\small}
\setlist[glosario,2]{leftmargin=1.2em, labelindent=0pt, topsep=.2\baselineskip}
\setlist[glosario,3]{leftmargin=1.2em, labelindent=0pt, topsep=.2\baselineskip}
\setlist[glosario,4]{leftmargin=1.2em, labelindent=0pt, topsep=.2\baselineskip}
% Término con subentradas: \item[Distribución]\subterminos  (fuerza salto de línea)
\newcommand*{\subterminos}{\mbox{}}
\newcommand*{\glosariocapitulo}{\seccionfinal{Glosario}}

% Glosario general (finales/glosario.tex): dos columnas
%   \begin{glosariogeneral} \item[Término:] Definición. \end{glosariogeneral}
\setlength{\columnsep}{7mm}
\setlength{\columnseprule}{0pt}
\NewDocumentEnvironment{glosariogeneral}{}
  {\begin{multicols}{2}\raggedcolumns\begin{glosario}[before=\small\raggedright]}
  {\end{glosario}\end{multicols}}

% Siglario
\newlist{siglario}{description}{1}
\setlist[siglario]{font=\normalfont\sffamily\bfseries, labelwidth=4.5em,
                   leftmargin=5.5em, labelsep=1em, align=left,
                   itemsep=.15\baselineskip, style=standard}

% ---------------------------------------------------------------------
% 12. Índice general
% ---------------------------------------------------------------------
\setkomafont{chapterentry}{\sffamily\bfseries}
\RedeclareSectionCommand[tocindent=0pt, tocnumwidth=1.5em]{chapter}
\RedeclareSectionCommand[tocindent=1.5em, tocnumwidth=2.3em]{section}


% ---------------------------------------------------------------------
% 13. Utilidades de edición y de ejercicios
% ---------------------------------------------------------------------
% Marca visible de algo a revisar en el original (buscar «\pendiente»)
\newcommand*{\pendiente}[1]{{\color{red!70!black}\sffamily\small[#1]}}

% Respuesta de un ítem resuelto dentro de un ejercicio
\newcommand*{\respuesta}[1]{{\color{acento2}\itshape #1}}

% Líneas para completar (número de renglones, por defecto 2)
\NewDocumentCommand{\completar}{O{2}}{%
  \enspace\dotfill\mbox{}%
  \foreach \i in {2,...,#1}{\newline\mbox{}\dotfill\mbox{}}}

% Palabra subrayada en consignas
\newcommand*{\subrayado}[1]{\uline{#1}}

% Texto de lectura (cuentos, fragmentos para analizar)
\newtcolorbox{textolectura}[1][]{%
  enhanced, breakable,
  colback=white, colframe=acento!60, boxrule=.5pt, arc=0pt,
  left=6mm, right=6mm, top=4mm, bottom=4mm,
  before skip=5mm, after skip=5mm,
  parbox=false,
  before upper={{\noindent\sffamily\bfseries\small\color{acento}#1\par}%
                \medskip\setlength{\parindent}{1.2em}},
}

% Figura cuya imagen todavía no está en figuras/: deja un recuadro con
% el epígrafe definitivo. Uso: \figurapendiente{Epígrafe}{fig:etiqueta}
\NewDocumentCommand{\figurapendiente}{m m O{45mm}}{%
  \begin{figure}[htbp]
    \centering
    \fcolorbox{gristexto}{white}{\parbox[c][#3][c]{.8\linewidth}{\centering
      \sffamily\small\color{gristexto}Imagen pendiente\\[2pt]\footnotesize #1}}
    \caption{#1}\label{#2}
  \end{figure}}

% ---------------------------------------------------------------------
% 14. Enlaces
% ---------------------------------------------------------------------
\hypersetup{
  linkcolor = acento,     % ejemplos, secciones, notas, índice
  citecolor = acento2,    % citas bibliográficas
  urlcolor  = acento,     % direcciones web
}

% Remisión a otro capítulo: \capref{4} → «capítulo 4», con enlace si ese
% capítulo ya está incluido en el libro; texto plano si todavía no.
\newcommand*{\capref}[1]{%
  \@ifundefined{r@cap:\two@digits{#1}}%
    {capítulo~#1}%
    {\hyperref[cap:\two@digits{#1}]{capítulo~#1}}}
% Variante abreviada: \Capref{4} → «cap. 4»
\newcommand*{\Capref}[1]{%
  \@ifundefined{r@cap:\two@digits{#1}}%
    {cap.~#1}%
    {\hyperref[cap:\two@digits{#1}]{cap.~#1}}}

% Etiquetas de corchetes y nudos con barra:
%   [\etq{SN} el\etq{D/Esp} libro\etq{N/n}]   N\prima (N')
\newcommand*{\etq}[1]{\textsubscript{#1}}
\newcommand*{\prima}{\ensuremath{'}}

% Transcripción fonética [ ] y fonológica / /
\newcommand*{\fon}[1]{[#1]}
\newcommand*{\fonol}[1]{/#1/}

% Imagen o audio que todavía no está en figuras/ (no flotante, para
% ejercicios): \imagenpendiente[alto]{Descripción}
\NewDocumentCommand{\imagenpendiente}{O{22mm} m}{%
  \par\smallskip\noindent\fcolorbox{gristexto}{white}{\parbox[c][#1][c]{\dimexpr\linewidth-2\fboxsep-2\fboxrule}{\centering
    \sffamily\footnotesize\color{gristexto}#2}}\par\smallskip}

% Texto ajeno que se incorpora desde el original (letras de canciones,
% textos literarios completos): deja un recuadro con la indicación.
\newcommand*{\textoaincorporar}[1]{%
  \imagenpendiente[18mm]{Incorporar aquí #1 desde el original}}

\makeatother


\addbibresource{bib/referencias.bib}

% --- Datos del libro --------------------------------------------------
\newcommand*{\titulolibro}{La gramática}
\newcommand*{\subtitulolibro}{De las nociones teóricas al lenguaje de todos los días}
\newcommand*{\responsables}{Mabel Giammatteo y Laura Ferrari\\(coordinadoras)}
\newcommand*{\editorial}{\textellipsis}
\newcommand*{\aniolibro}{2026}

\hypersetup{pdftitle={\titulolibro. \subtitulolibro},
            pdfauthor={Mabel Giammatteo y Laura Ferrari (coords.)}, pdflang={es}}

% --- Capítulos a compilar ---------------------------------------------
% Para trabajar en un solo capítulo, descomentar y listar solo ese:
% \includeonly{capitulos/cap01}

\begin{document}

% Todas las entradas del .bib quedan disponibles. La bibliografía general
% (al final) imprime las de los capítulos incluidos y todo lo citado; las
% lecturas recomendadas de cada capítulo se filtran por lectNN.
\nocite{*}

\frontmatter
% preliminares/portada.tex
\begin{titlepage}
  \raggedright
  \vspace*{30mm}
  {\color{acento}\rule{\linewidth}{1.2pt}}\par
  \vspace{8mm}
  {\sffamily\bfseries\fontsize{26}{30}\selectfont \titulolibro\par}
  \vspace{4mm}
  {\sffamily\Large\color{gristexto}\subtitulolibro\par}
  \vspace{8mm}
  {\color{acento}\rule{\linewidth}{.6pt}}\par
  \vfill
  {\large\itshape\responsables\par}
  \vspace{20mm}
  {\sffamily\small\editorial\enskip·\enskip\aniolibro\par}
\end{titlepage}
\cleardoublepage

\tableofcontents
% preliminares/prologo.tex
% Fuente: Prólogo.docx
\addchap{Prólogo}
\addtocontents{toc}{\protect\autorentoc{Mabel Giammatteo y Laura Ferrari}}

La gramática, surgida en la Antigüedad clásica, venerada y vilipendiada a lo
largo de su desarrollo en Occidente, es actualmente considerada «la estructura
del lenguaje». Según \textcite{saussure1945}, este sistema es el encargado de
aparear determinadas porciones de la masa amorfa de nuestros pensamientos con
determinadas porciones del sonido para conformar signos lingüísticos que hagan
comunicables nuestros pensamientos, deseos y sentimientos. Por su complejidad,
esta estructura ha sido comparada con el código genético, ya que, así como este,
mediante la combinación de una mínima cantidad de elementos, puede generar la
inmensa variedad de lo viviente, la gramática nos permite «combinar un número
relativamente pequeño de unidades y obtener un número potencialmente infinito de
resultados» \parencite[18]{bosque2009}. Dado, entonces, que tanto para precisar
nuestras ideas como para interactuar con otras personas nos valemos del
lenguaje, podemos afirmar que nuestro desarrollo cognitivo y nuestra vida social
están bajo el dominio de la gramática.

Por un lado, la gramática se nos revela como un cuerpo de doctrina atesorado por
siglos de estudios sobre el lenguaje y las distintas lenguas y, por el otro,
constituye un conocimiento vivo, según iremos viendo, internalizado en la mente
de cada hablante \parencite{chomsky1957} e indispensable para la interacción
social. Este manual surge en la intersección entre el saber teórico y la
necesidad de hacer consciente el conocimiento sobre cómo funcionan las
estructuras que hacen posible el lenguaje. Es fruto de un proyecto compartido
nacido en el seno de la cátedra B de Gramática de la Facultad de Filosofía y
Letras de la Universidad de Buenos Aires. Su objetivo ha sido ofrecer las
nociones y herramientas básicas para analizar la lengua, pero, sobre todo,
propiciar un acercamiento reflexivo, que permita comprender cuáles son sus
unidades y cómo se combinan para construir significados relevantes.

La obra, que recoge la experiencia docente y el trabajo de investigación del
equipo, surgió inicialmente como una simple guía de ejercicios para los y las
estudiantes de nuestros cursos de primer año de la Facultad y fue poco a poco
creciendo con la incorporación de las nociones teóricas que fundamentaban las
distintas prácticas solicitadas. Asimismo, el núcleo inicial de temas abordados
también se fue ampliando hasta convertirse en lo que hoy presentamos: un manual
completo teórico-práctico, que refleja las teorías actuales sobre el sistema y
las aplica al estudio de la lengua de todos los días.

Se trata de un texto colectivo, con un espíritu común, pero cuyos capítulos, que
pueden ser aprovechados en forma independiente, conservan la impronta que les
han asignado quienes los han producido. La idea central ha sido proveer un
material actualizado, a la vez que accesible, que pueda servir de guía o hilo
conductor para cualquier curso de gramática del español en las etapas iniciales
del nivel superior, ya sea en la universidad o en profesorados de Lengua y
Literatura, pero que también puede ser de utilidad en cursos superiores de grado
y posgrado y en capacitaciones docentes. Asimismo, también puede servir de base
de comparación para quienes estudian alguna lengua extranjera, y como texto de
referencia para profesionales que trabajan con el lenguaje, como periodistas,
especialistas en traducción, corrección de textos o edición, y para quienes estén
interesados en profundizar sus conocimientos sobre los distintos aspectos de la
gramática del español.

El libro se inicia con un capítulo introductorio que presenta diferentes teorías
lingüísticas y su incidencia en la enseñanza de la gramática, seguido de otro
dedicado a distinguir entre nociones básicas, como las de gramaticalidad,
corrección y aceptabilidad. Los capítulos siguientes conforman un trayecto que
aborda las principales cuestiones de la gramática del español, desde los sonidos
y su caracterización, pasando por la estructura de la palabra y sus procesos de
formación, tanto los tradicionales como otros más informales y novedosos, hasta
la actitud comunicativa del hablante y su reflejo en la modalidad de las
oraciones. La parte central de la obra considera la estructura oracional desde
una perspectiva que integra tres niveles: el nivel interior o del sintagma
verbal, relacionado con la predicación; el nivel central, vinculado con la
flexión, donde se establece la concordancia entre el sujeto y el predicado; y el
nivel externo o periferia oracional, en el que se manifiesta la fuerza ilocutiva
del hablante. Se tratan también cuestiones relacionadas con aspectos
problemáticos de la gramática del español, como las construcciones pronominales,
las perífrasis verbales o los predicativos, y los nuevos planteos para el
tratamiento de las estructuras interoracionales, con un capítulo dedicado a las
oraciones coordinadas y otro a las subordinadas. La obra incluye un capítulo
final dedicado a la variación, en el que se discuten las nociones de lengua y
dialecto y se presentan las principales características del español del Río de
la Plata.

Las nociones teóricas se abordan desde un enfoque léxico-sintáctico que no solo
tiene en cuenta las estructuras de la lengua, sino también los significados que
resultan de sus combinaciones y las interpretaciones que producen en los textos
en los que se insertan. Así, lo más destacable de la obra es que no se limita a
la descripción de las formas lingüísticas ni a su explicación en relación con
principios gramaticales generales, sino que las estructuras se conectan con el
uso en diferentes tipos de textos, tanto literarios como cotidianos (noticias
periodísticas, afiches callejeros, memes de internet, canciones y proverbios,
entre otros géneros que la obra aborda). Cada capítulo contiene numerosos
ejemplos que dan cuenta de las estructuras estudiadas y proporciona abundantes y
variadas actividades que permiten poner en práctica los contenidos. La obra
también incorpora un solucionario con las claves de resolución de los
ejercicios, y se completa con un siglario y un glosario que hacen más accesibles
su empleo y su consulta.

Hasta aquí, nuestro trabajo se ha concentrado en el diseño de la estructura del
libro, en la supervisión de la escritura de los capítulos y en la revisión
final. A partir de este momento, con el texto ya concluido, solo nos queda
confiar en que sea de utilidad para quienes deciden aventurarse por los caminos
del estudio de la gramática y también en que despierte el interés para seguir
leyendo. Para ello, los distintos capítulos proporcionan una lista de textos de
referencia que permitirán ampliar y profundizar los conocimientos.

\bigskip
\begin{flushright}
  \itshape Las coordinadoras\\
  \upshape Buenos Aires, 9 de octubre de 2026
\end{flushright}


\mainmatter
% Un \include por capítulo, en orden: la numeración del capítulo
% determina qué referencias imprime (palabra clave capNN del .bib).
% =====================================================================
%  capitulos/cap01.tex — Capítulo 1. La gramática y su enseñanza
%  Autora: Mabel Giammatteo
%  Fuente: Capitulo_1_La_gramatica_y_su_ensenanza_normalizado.docx
% =====================================================================
\capitulo{La gramática y su enseñanza}{Mabel Giammatteo}

\section{Introducción}

La lengua es una entidad que ni siquiera necesita a sus hablantes para existir,
como prueban los numerosos idiomas muertos que han sobrevivido a sus pueblos.
Los individuos, en cambio, no podemos hacer prácticamente nada ---o nada
relevante para la vida en comunidad--- sin utilizarla. Pero la lengua no es un
instrumento externo con el que debamos familiarizarnos o que tengamos que
aprender a manejar, como si fuera un celular, un electrodoméstico o un
automóvil, sino que, como sostiene \textcite[18]{bosque2018}, se trata de «un
sistema complejo que los hablantes tenemos interiorizado y que constituye una
parte esencial de nuestra naturaleza». Conforma una estructura que, junto con
el vocabulario, se encuentra almacenada en nuestra mente y a diario la
empleamos en nuestros intercambios. Con solo abrir la boca\ldots, podríamos
decir.

Comenzamos este capítulo introductorio con un recorrido evaluativo sobre
distintos planteos que a lo largo de la historia se han hecho respecto de la
gramática y su utilidad para la enseñanza. Revisamos así conceptos centrales
como los de lengua y habla, competencia y actuación, sistema y uso, entre
muchos otros. Luego, introducimos la propuesta léxico-gramatical que se irá
desarrollando a lo largo de todo este libro y que tiene como objetivo avanzar
desde las nociones teóricas a la reflexión sobre el lenguaje de todos los días.
Para ello, además de describir y explicar las estructuras lingüísticas, nos
interesa conectarlas con el uso. Esperamos así poder colaborar en el
reconocimiento de los valores expresivos y comunicativos de las configuraciones
morfosintácticas y léxico-semánticas que se irán estudiando en este libro.

\section{Debates e interrogantes sobre la gramática y su enseñanza}

Si la lengua forma parte del conocimiento internalizado de los hablantes, que
diariamente la utilizamos para comunicarnos a través de distintos medios y
redes sociales, ¿para qué se necesita enseñarla en la escuela? ¿Para qué
insistir con paradigmas verbales y complicados análisis de oraciones?
Detengámonos en esto. Como señaló \textcite{saussure1945}, el lenguaje tiene
un doble costado. Por un lado, es lo que cada uno de nosotros «dice» en cada
una de las distintas circunstancias en que lo empleamos, lo que denominó
\termino{habla} y caracterizó como momentánea, individual y variable. Pero
también, por otro lado, el lenguaje tiene un costado social, la
\termino{lengua} propiamente dicha, que constituye un factor de herencia,
cultura y tradición. La lengua es, sin duda, un componente de nuestra
identidad, que nos atraviesa, interpela y define. Como facultad humana propia
de nuestra especie \parencite{chomsky1989}, el lenguaje es universal, pero la
lengua es propia de cada cultura, de cada pueblo y también de cada hablante que
la hace suya. Desde esta perspectiva, la educación lingüística en la escuela
cumple primariamente una función social: enseñar la lengua estándar a
poblaciones muchas veces heterogéneas y que se expresan mediante diferentes
variedades de uso. Por supuesto, esto no quiere decir que estas formas
particulares del habla de cada grupo o región, también conocidas como
lectos,\footnote{\pendiente{Falta el texto de esta nota en el original.}}
deban ser abandonadas o menospreciadas, ya que su dominio enriquece y
caracteriza al que lo posee. Pero la incorporación de una variedad estándar
facilita el acercamiento a los textos que se necesita comprender y producir en
el ámbito escolar. Asimismo, la enseñanza también debe hacer consciente ese
conocimiento tácito y subliminal, que naturalmente el hablante posee sobre su
lengua. Por tanto, la escuela tiene que ofrecer lo siguiente:

\begin{itemize}
  \item un \termino{conocimiento del uso}, que respete la variación, pero que
        también incorpore los aspectos formales necesarios para acceder a los
        textos con los que se debe interactuar en los estudios;
  \item un \termino{conocimiento reflexivo de la estructura y posibilidades},
        que pueda no solo colaborar con la mejora de las habilidades
        cognitivo-lingüísticas ---argumentar, fundamentar, comparar, explicar,
        entre otras---, sino también incentivar el desarrollo de la capacidad
        de abstracción y el juicio crítico, que se requieren tanto para los
        estudios como para la vida cotidiana.
\end{itemize}

\begin{recuadro}[Síntesis]
Toda actividad reflexiva sobre la lengua redunda en un mayor conocimiento de
las posibilidades expresivas que ofrece, lo cual, hacia el exterior, se
proyecta en mejores intercambios comunicativos y, hacia el interior, incide en
la capacidad de conceptualización y de establecer relaciones entre fenómenos
del mundo circundante.
\end{recuadro}

En esta perspectiva, la enseñanza de la lengua materna en la escuela debe
apelar a un conocimiento metacognitivo\footnote{El término
\textit{metacognición} \parencite{flavell1979} refiere a la capacidad de
reflexionar sobre los propios procesos cognitivos, es decir, tener control
consciente sobre las operaciones mentales que realizamos.} que forme hablantes
conscientes, que no solo reconozcan las estructuras de la lengua, sino que
también puedan identificar los valores comunicativos que transmiten y avanzar
desde la interpretación de los significados convencionales, que son los que
tenemos internalizados, a otros más novedosos o complejos, que exigen mayor
esfuerzo interpretativo, como los sentidos metafóricos o los de mayor nivel de
abstracción.

Para ilustrar la problemática que estamos tratando, podemos partir de una
historieta publicada hace ya un tiempo en el diario \textit{Clarín} de Buenos
Aires, en la que tiene lugar el siguiente diálogo entre una madre y su hija:

\figurapendiente{Donatela, por Raznovich (\textit{Clarín}, 2/11/12)}{fig:donatela}

Más allá de la ironía presente en la historieta, \forma{comunicarnos} aquí
refiere a uno de los usos más básicos del lenguaje: el simple \sig{entrar en
contacto} con el otro, al que Jakobson (1984)\pendiente{Jakobson 1984 no está
en la bibliografía} denominó \termino{función fática},\footnote{Las funciones
del lenguaje, según Jakobson (1984), son seis: referencial, centrada en la
transmisión de información; emotiva, que manifiesta la expresividad del emisor;
conativa, que apela al receptor del mensaje; metalingüística, que se emplea
cuando usamos el código lingüístico para hacer referencia al lenguaje mismo;
poética, que pone el acento en la conformación misma del mensaje; y fática,
centrada en el funcionamiento del canal a través del cual se produce el
intercambio.} la cual es muy útil en el inicio de los estudios de una lengua
extranjera, cuando básicamente lo que queremos es «conectarnos» con sus
hablantes mediante saludos e intercambios corteses. Sin embargo, estudiar la
lengua materna implica una actitud muy diferente, ya que el objetivo no es
aprender a hablarla, sino que lo que se busca es adquirir la práctica de
reflexionar sobre los matices significativos que sus diferentes construcciones
permiten transmitir e interpretar.

Volviendo a la inclusión de la gramática en la enseñanza, cuestión que, como
todos sabemos, viene siendo debatida desde hace ya largo tiempo, las distintas
posturas han oscilado entre la negación absoluta de su utilidad, por un lado, y
encendidas defensas de su contribución a la formación lingüística, por el otro.
Por eso, antes de avanzar, debemos especificar a qué nos referimos con
\forma{gramática}, ya que, según veremos a continuación, a lo largo de la
historia, no siempre se entendió el término de la misma manera.

\subsection{De la gramática tradicional al estructuralismo}

En el panorama de los estudios de la Antigüedad clásica y el Medioevo, la
gramática, junto con la dialéctica y la retórica, conformaba el
Trivium,\footnote{El Trivium y el Quadrivium ---formado por aritmética,
geometría, astronomía y música--- constituían el primer y segundo ciclo de las
llamadas artes liberales, surgidas en el siglo \textsc{vi} y representativas
del saber de la época.} destinado a comprender «los secretos de las palabras».
Tal como transmite \textcite[108]{barthes1993}, las gramáticas latinas de
Donato (hacia 350) y Prisciano (fines del s.~\textsc{v} y comienzos del
s.~\textsc{vi}) «representaron la ley absoluta, salvo cuando no coinciden con
la Vulgata:\footnote{Es la «versión latina de la Sagrada Escritura, [\ldots]
declarada auténtica por la Iglesia católica» (\textit{DLE}, en línea).} la
gramática no podía ser entonces más que normativa,\footnote{La gramática
normativa o prescriptiva es la que se ocupa de los usos correctos o incorrectos
según la norma aceptada socialmente en un momento determinado. Estos conceptos
se amplían en el capítulo~2.} puesto que se pensaba que las \enquote*{reglas}
de la locución habían sido inventadas por los gramáticos [\ldots]. Hasta el
siglo \textsc{xii}, la \textit{Grammatica} abarca la gramática y la poesía,
trata a la vez de la \enquote*{precisión} y de la \enquote*{imaginación}; de las
letras, las sílabas, la oración, el período,\footnote{En gramática,
\textit{período} hace referencia a la «estructura oracional, generalmente
bimembre, formada por una oración principal y una subordinada» (\textit{DLE},
en línea), sobre todo las construcciones condicionales y concesivas.} las
figuras\footnote{Las figuras retóricas son empleos no convencionales de las
palabras con fines expresivos.} y la métrica».\footnote{Estudio de la medida o
estructura de los versos y de sus distintas combinaciones.}

Las teorías sobre el lenguaje basadas en los estudios de los gramáticos griegos
y latinos constituyen lo que se conoce como gramática tradicional, cuya
influencia en la enseñanza se mantuvo vigente hasta la primera mitad del
siglo~\textsc{xx}. Su enfoque estaba centrado en la palabra y en sus
«accidentes» ---variaciones en la terminación para manifestar las categorías
morfológicas (género, número, caso, persona, tiempo, modo, aspecto,
voz)\footnote{El concepto de categoría morfológica se trata en el capítulo~4.
Para ampliar, v. \textcite{giammatteo2018}.}---. La sintaxis ocupaba un lugar
bastante marginal, de modo que, en las gramáticas y manuales escolares, las
nociones sintácticas solían colocarse al final y se orientaban sobre todo a
favorecer la buena formación de las oraciones y a evitar los llamados «vicios»
o errores de construcción.\footnote{Entre los considerados vicios del lenguaje
se encuentran el barbarismo (uso de extranjerismos), el solecismo (alteración
del orden sintáctico), el pleonasmo (agregado de palabras innecesarias), la
anfibología (uso de palabras con doble sentido), la cacofonía (combinación
desagradable de sonidos en un sintagma), entre otros.}

\begin{destacado}
Para la gramática tradicional, el interés fundamental radicaba en el «buen uso
de la lengua», por lo que las reglas y preceptos eran de tipo normativo y no
había una explicitación del funcionamiento de la lengua.
\end{destacado}

En el siglo \textsc{xix} la lingüística se encontraba influida por el afán
cientificista del positivismo,\footnote{Teoría filosófica formulada por
Auguste Comte, para la cual la única fuente del conocimiento es la experiencia,
por lo que el conocimiento debe ser observable, verificable y objetivo.} que
buscaba para las humanidades el rigor metodológico de las ciencias naturales,
por lo que el estudio del lenguaje comenzó a ser «cuidadosa y objetivamente
investigado y luego explicado en términos de hipótesis inductivas»
\parencite[22]{lyons1985}. Asimismo, la atracción que ejercieron los planteos
evolucionistas de Darwin hizo que la observación y el análisis de los datos se
enfocaran al estudio histórico y comparativo de las lenguas. Las
investigaciones de la época se basaban fundamentalmente en el análisis de las
distintas lenguas de la familia indoeuropea,\footnote{La mayoría de las lenguas
de Europa (el griego, el latín y sus derivadas las lenguas romances, las
lenguas germánicas, las eslavas, las célticas, etc.) y algunas de Asia (el
sánscrito y otras lenguas indoiranias, el persa, etc.) pertenecen a esta
familia.} que, por provenir de un tronco común y poseer textos escritos,
constituían un excelente campo de aplicación de la metodología comparativa. Los
estudios se hacían contrastando las unidades lingüísticas entre sí,
considerando su origen común y las características de su desarrollo histórico
en cada lengua, pero sin ir más allá, es decir que no se contemplaban sus
relaciones respecto del funcionamiento global de la lengua.

La idea de que en la lengua todos los elementos están interrelacionados, es
decir, que constituyen un \termino{sistema}, surge recién a principios del
siglo~\textsc{xx}, a partir de los planteos de \textcite{saussure1945},
considerado el padre de la lingüística moderna y el iniciador de la corriente
conocida como estructuralismo.\footnote{No obstante, Saussure no usó la palabra
\textit{estructura}, sino \textit{sistema}.} Los planteos de Saussure
impulsaron un cambio de paradigma en los estudios sobre el lenguaje, que
pasaron a ocuparse del conocimiento

\begin{itemize}
  \item no solo de las lenguas de prestigio, sino de \textit{todas las lenguas};
  \item no solo de las formas correctas (lingüística prescriptiva), sino de
        \textit{todas las formas} (lingüística descriptiva);
  \item no solo de la lengua escrita, sino primordialmente de \textit{la lengua
        hablada};
  \item no solo del acto individual, sino de \textit{la naturaleza social del
        lenguaje};
  \item no solo del desarrollo histórico de las lenguas (\textit{diacronía}),
        sino de \textit{su estado actual} (\textit{sincronía});
  \item no de la lengua como una nomenclatura, constituida por elementos
        aislados, sino como \textit{un sistema de signos formados por un
        significado} (contenido) \textit{y un significante} (soporte material
        de la expresión).\footnote{El soporte material debe ser perceptible,
        como el sonido en la lengua hablada, los gestos en la lengua de señas,
        o los signos gráficos en la escritura.}
\end{itemize}

Para Saussure, la tarea de la lingüística debía centrarse en los siguientes
puntos:

\begin{itemize}
  \item describir todas las lenguas y familias de lenguas;
  \item formular leyes generales;
  \item delimitarse y definirse a sí misma, es decir, constituirse como
        disciplina autónoma e independiente.
\end{itemize}

\begin{destacado}
Frente al carácter normativo de la gramática tradicional, el estructuralismo, a
partir del concepto central de \termino{sistema}, aportó un modelo consistente
para la descripción de las lenguas.
\end{destacado}

Dentro del estructuralismo, además de conceptos fundamentales para la
disciplina, como \textit{signo, sistema, lengua/habla, diacronía/sincronía,
significado/significante}, también se desarrollaron nociones clave, como las de
\termino{fonema} y \termino{morfema},\footnote{La noción de fonema se trata en
el capítulo~3 y la de morfema, en el capítulo~4.} que se incorporaron al
análisis junto a las unidades tradicionales, como la palabra y la oración.
Asimismo, para distinguir las «partes de la oración» o clases de palabras, se
buscó superar las distinciones basadas en conceptos semántico-nocionales y
establecer criterios más nítidos, asentados en pautas formales. Se puede
ejemplificar esta diferencia con las dos definiciones del verbo siguientes:

\begin{exe}
  \ex\label{ex:01-defverbo}
    \begin{xlist}
      \ex\label{ex:01-defverbo-a} palabra que manifiesta \sig{acción}.
      \ex\label{ex:01-defverbo-b} palabra que concuerda con el sustantivo núcleo del sujeto en
          número y persona.
    \end{xlist}
\end{exe}

La definición \exref{ex:01-defverbo-a}, que descansa en un criterio
semántico-nocional como el concepto de \sig{acción}, presenta la dificultad de
que no es aplicable a todos los verbos, ya que además de los que efectivamente
manifiestan una \sig{acción}\footnote{\label{fn:01-accion}Los verbos de
\sig{acción} responden a la pregunta \textit{¿Qué hizo X?} Así, frente a
\textit{¿Qué hizo Juan?} podemos responder: \textit{Juan corrió, comió} o
\textit{estudió}, pero no \textit{Juan creció} o \textit{nació}, que son verbos
de proceso, ni tampoco \textit{Juan tuvo fiebre} o \textit{era alto}, que son
verbos de estado (v. capítulo~6).} ---\forma{correr, comer, estudiar}---,
existen otros que señalan \sig{procesos} ---\forma{crecer, nacer} o
\forma{florecer}---, y otros que refieren \sig{estados} ---\forma{tener, ser,
durar}---. En cambio, \exref{ex:01-defverbo-b}, basada en un criterio formal,
como la concordancia en número y persona entre el verbo y su sujeto, resulta
aplicable en todos los casos.

En relación con la sintaxis, que fue adquiriendo cada vez más espacio frente a
la posición a la que la relegaba la gramática tradicional, el concepto de
\termino{constituyente inmediato} (CI) (v. capítulo~5), propuesto por el
estructuralista norteamericano Rulon \textcite{wells1947}, se convirtió en la
base del análisis sintáctico (ingl. \textit{parsing}), tal como lo conocemos
hoy en día.

Aunque el estructuralismo aportó un modelo completo para el estudio de la
lengua y un fundamento para el análisis sintáctico, con el correr del tiempo su
enfoque se consideró reduccionista porque, a partir de la distinción
saussureana entre \termino{lengua} y \termino{habla}, se dejó de lado el uso y
se puso el énfasis en el sistema. Este enfoque también recibió críticas por
establecer un modelo estático, que se centraba en el aspecto sincrónico y
dejaba de lado el cambio y la variación, aspectos esenciales de la dinámica de
las lenguas. Finalmente, también se lo cuestionó por ser fundamentalmente
taxonómico, es decir, enfocarse en la clasificación de las unidades para su
reconocimiento y etiquetado, dejando de lado su funcionamiento real en los
textos.

\subsection{La llegada del enfoque comunicativo}

A fin de subsanar los problemas que acabamos de mencionar, a mediados de los
años 80 del siglo pasado, se introdujo el denominado enfoque comunicativo,
surgido en la enseñanza de lenguas segundas y extranjeras. El objetivo de su
incorporación en el estudio de la lengua materna era subsanar la simplificación
y el alejamiento de los textos que el modelo estructural ofrecía. Se
incorporaron, entonces, propuestas provenientes de la gramática textual, del
análisis del discurso y de las teorías de la enunciación y de los actos de
habla, a las que se completaba con planteos tomados de los enfoques
predominantes en la enseñanza de lenguas extranjeras.

Asimismo, buscando remediar las limitaciones de la gramática oracional, se
introdujeron conceptos de la propuesta funcional\footnote{\label{fn:01-funcionalismo}El
funcionalismo parte de lo afirmado por los lingüistas estructuralistas de la
escuela de Praga (1929): «la lengua es un sistema de medios de expresión
apropiados para un fin» (\textit{TCLP}, I; \cite{aavv1929}). Los
funcionalistas conciben el lenguaje principalmente como un sistema de
comunicación, que se crea en la interacción social.} de
\textcite{halliday1976}. Se incorporaron en las aulas y manuales nociones
referidas al análisis textual, como coherencia, cohesión, progresión temática,
registro, entre otras. Sin embargo, como plantea \textcite[104]{gaiser2011},
«lo que el discurso de la didáctica de la lengua olvidó, es que, como dice el
propio Halliday, para aplicar a un texto la gramática funcional, se debe tener
una sólida base en gramática oracional».

En línea con lo señalado para Argentina, \textcite[18]{munguia2015} hace la
siguiente evaluación respecto de la aplicación del enfoque comunicativo en
México:

\begin{cita}
esta nueva reforma trató de enseñar «haciendo», es decir, los alumnos
aprenderían a escribir «escribiendo», a leer «leyendo» textos cuyas temáticas
fueran cercanas a sus intereses y motivaciones. Entonces, se prescindió de toda
teorización, por ello se eliminaron casi todos los temas relacionados con
gramática [\ldots]

Pero ya transcurrieron más de 20 años y los resultados siguen siendo muy
desalentadores [\ldots]. Las diferentes evaluaciones aplicadas muestran que el
enfoque comunicativo está lejos de haber logrado sus fines.
\end{cita}

En el ámbito local, lo que se observa en la actualidad es que, aunque la
propuesta comunicativa no cuajó en un modelo eficaz para la enseñanza de la
lengua materna, lo que logró fue arrinconar los contenidos gramaticales hasta
llevarlos a su casi completa desaparición de las aulas. En cuanto a la
sintaxis, o bien se ignora su tratamiento o bien se aplica el modelo más
conocido y que reproduce la mayoría de los manuales, que es el de cuño
estructural. Asimismo, la práctica tiende a ser realizada de un modo
predominantemente clasificatorio, sin establecer una genuina reflexión entre
los constituyentes que se reconocen en el análisis y los valores semánticos y
comunicativos que transmiten. Así, como expresa \textcite{crystal2007},

\begin{cita}
[n]o es sorprendente que la mayoría de la gente que aprendió análisis
sintáctico en el colegio termine sin poder ver la utilidad de este ejercicio, y
solo pueda recordar la gramática como una disciplina muerta e irrelevante, la
anatomía forense del lenguaje.
\end{cita}

\noindent Sin embargo, el mismo \textcite{crystal2007} añade que

\begin{cita}
[l]a realidad dice lo contrario. Las técnicas del análisis gramatical pueden
ser usadas para demostrar el enorme poder creativo del lenguaje ---cómo a
partir de un conjunto finito de patrones gramaticales, aún un niño pequeño
puede producir un conjunto infinito de oraciones--- (traducción nuestra).
\end{cita}

\subsection{Los aportes de la lingüística generativa}\label{sec:01-generativa}

La propuesta de la corriente conocida como gramática
generativa,\footnote{La denominación de la teoría se vincula con la concepción
de gramática que propone \textcite{chomsky1974a}: «la gramática de una lengua
será, por consiguiente, un dispositivo que genera todas las secuencias
gramaticales de la lengua y ninguna de las agramaticales» (v. capítulo~2).}
basada en los planteos de \textcite{chomsky1970,chomsky1974a}, en un sentido
representa una continuidad con el estructuralismo en lo referido a la
preocupación por el rigor metodológico, que en el caso de Chomsky lo lleva a
buscar una descripción del lenguaje cercana a las formulaciones lógicas y
matemáticas. Pero, en otro sentido, constituye una fuerte ruptura con todas las
teorías anteriores. A partir de los trabajos de Chomsky desde los años 60 del
siglo pasado, se produjo una revolución en el campo de la lingüística y las
ciencias cognitivas:\footnote{La lingüística integra el espacio de las ciencias
cognitivas junto a otras disciplinas interesadas en la exploración de la mente
(psicología, antropología, filosofía, inteligencia artificial, neurología,
etc.).}

\begin{cita}
Frente a los posicionamientos teóricos propios del estructuralismo europeo y
americano que concebían el lenguaje como institución social [\ldots], surge una
visión del lenguaje como facultad, que forma parte de la
mente/cerebro\footnote{Esta dupla hace referencia a la parte psíquica (mente) y
a su soporte físico (cerebro).} y debe ser aprehendido con técnicas y métodos
similares a los empleados en las ciencias naturales\ldots{}
\parencite[54]{mairal2012}.
\end{cita}

El lenguaje deja de considerarse un producto exclusivamente cultural y pasa a
estudiarse como una \termino{facultad}, en sentido biológico, una capacidad
propia y exclusiva de nuestra especie, con la que estamos dotados
naturalmente. Así, para \textcite[14]{chomsky1992},

\begin{cita}
[e]n el habla normal, uno no repite mecánicamente lo que ha oído, sino que
produce formas lingüísticas nuevas ---a menudo nuevas en la experiencia de uno
o incluso en la historia de la lengua--- y no hay límites para dicha
innovación.
\end{cita}

Pero ¿cómo es que esto se produce? La respuesta se vincula con lo que Chomsky
denominó \termino{el problema de Descartes} o \termino{de la infinitud
discreta}, que se refiere a la productividad del lenguaje, ya que, como
explican \textcite[17]{eguren2004}: «Los hablantes son capaces de producir un
número ilimitado de oraciones nuevas porque conocen un número finito de
unidades lingüísticas y saben, además, cómo combinarlas». Esta perspectiva
lleva a Chomsky a proponer que el lenguaje humano es recursivo\footnote{Un
proceso recursivo es aquel que se puede volver a aplicar al resultado de su
aplicación previa, en forma potencialmente indefinida. En gramática, la
subordinación es un proceso recursivo porque dentro de una oración se pueden
incrustar otras. Así en una oración simple como [\textsubscript{O} Tengo una
casa], podemos incluir otra, como en [\textsubscript{O1} Tengo una casa
[\textsubscript{O2} que está alejada del centro de la ciudad]], y conformar
una oración compuesta (v. capítulo~13).} y a plantear que la recursividad es
su característica más prominente, ya que la combinación de elementos finitos
puede dar lugar a un número infinito de oraciones.

El otro dilema que Chomsky ha planteado es el que se conoce como \termino{el
problema de Platón} o \termino{de la pobreza de los estímulos}, y tiene que ver
con otra importante cuestión abordada por esta teoría: ¿cómo se adquiere el
lenguaje? A diferencia de lo que sostenían los psicólogos
conductistas,\footnote{La psicología conductista, que tuvo gran influencia en
el estructuralismo norteamericano, considera que las conductas se producen
como respuesta a los estímulos del ambiente.} para quienes el lenguaje es una
conducta aprendida mediante ensayo y error, Chomsky sostiene que el lenguaje se
adquiere. ¿Pero en qué consiste la diferencia? Dentro de la mente hay
actividades cognitivas diferentes, con características y funciones
particulares. \textcite{chomsky1993} distingue entre «tareas naturales», que se
adquieren sin esfuerzo ---respirar, caminar, hablar--- y otras más complejas,
que tienen que ser aprendidas con empeño ---tocar el piano, jugar al ajedrez,
escribir, o construir teorías científicas---.

Sin embargo, aunque es natural, el desarrollo del lenguaje requiere del
estímulo que proviene de escucharlo en el medio ambiente, de la misma manera
que una semilla necesita la tierra y el agua para germinar en una planta:

\begin{cita}
[C]ualquier organismo necesita un entorno rico y estimulante para que sus
capacidades naturales emerjan. Volviendo a la imagen de que enseñar era como
permitir crecer bien a una flor, si no se riega la planta no va a crecer hasta
dar flor. No es que aprenda del agua a ser flor ---si fuera un árbol usaría la
misma agua para crecer hasta convertirse en árbol---. Pienso que sucede más o
menos lo mismo en el desarrollo humano, sin excluir el desarrollo del lenguaje
y del pensamiento \parencite[160]{chomsky1992}.
\end{cita}

Pero ¿qué implica esta nueva perspectiva que propone la gramática generativa
para el estudio del lenguaje? Si el lenguaje es una facultad alojada en
nuestra mente, se puede decir que cada hablante nativo tiene internalizada la
gramática de su lengua, en tanto puede identificar, formar e interpretar las
palabras y oraciones que se producen en ella. A este conocimiento tácito o no
consciente que el hablante posee acerca de su lengua, \textcite{chomsky1970}
lo denomina \termino{competencia} y lo diferencia de la \termino{actuación} o
uso real de la lengua en situaciones concretas.

Dado que el lenguaje es un conocimiento, \textcite{chomsky1989} plantea un
programa de investigación en torno a tres preguntas básicas:

\begin{itemize}
  \item ¿Qué es lo que constituye el conocimiento del lenguaje?
  \item ¿Cómo se adquiere ese conocimiento?
  \item ¿Cómo se utiliza?
\end{itemize}

\begin{ejercicio}\label{ej:01-facultad}
Teniendo en cuenta lo que se ha planteado en este capítulo y después de leer
muy atentamente el texto siguiente, complete lo que se pide más abajo, según
se muestra en el ejemplo resuelto~1.

\begin{exe}
  \ex El lenguaje es una facultad porque \respuesta{es una capacidad propia de
      nuestra especie}.\hfill{\footnotesize\sffamily\color{acento2}ejemplo resuelto}
  \ex Se trata de una facultad universal porque\completar
  \ex Es privativa de nuestra especie porque\completar
  \ex El lenguaje sirve para la comunicación, pero también cumple otros fines,
      como ser\completar
  \ex La forma oral del lenguaje es natural para el hombre porque\completar
  \ex Desde esta perspectiva, la escritura puede considerarse\completar
  \ex Sin el lenguaje la vida social no sería posible porque\completar
  \ex El lenguaje es parte de la dotación biológica de nuestra especie
      porque\completar
  \ex Para el desarrollo del lenguaje es necesario interactuar con el medio
      porque\completar
  \ex Somos capaces de producir combinaciones de unidades que constituyen
      nuevas oraciones que jamás hemos oído porque\completar
\end{exe}
\end{ejercicio}

\subsubsection{De la facultad del lenguaje a la gramática}

Para Chomsky la facultad del lenguaje consiste en \termino{principios} muy
generales que están presentes en todas las lenguas y constituyen la
\termino{gramática universal} (GU), de naturaleza mental o cognitiva. Estos
principios son comunes a todos los individuos de la especie y hacen posible que
los hablantes adquieran su lengua, es decir, que desarrollen la competencia
gramatical en el idioma hablado en su entorno.

Ahora bien, si los principios están en todas las lenguas, ¿por qué estas
difieren tanto? La respuesta de la teoría es que en cada lengua los principios
se realizan mediante \termino{parámetros}, que son los que establecen las
diferencias que se reconocen entre las lenguas.

Para explicar las diferencias, Chomsky recurre a la comparación con un
conmutador, cuyas complejas redes representan los distintos principios de la
facultad del lenguaje. Para que el sistema funcione, la red posee distintos
interruptores que deben fijarse en una de las dos posiciones disponibles para
cada uno. Cada posición adoptada por un interruptor representa el parámetro que
la lengua en cuestión adopta para cada principio y el tener que optar entre dos
posibilidades (y no entre infinitas) facilita a los niños la adquisición de la
propia lengua.

\figurapendiente{Una metáfora técnica}{fig:01-metafora}

Ahora veamos un sencillo ejemplo que permitirá ilustrar mejor la diferencia
entre principios y parámetros. Si observamos las siguientes oraciones, vemos
que mientras el español \exref{ex:01-sujnulo-a} y el italiano
\exref{ex:01-sujnulo-b} permiten omitir el sujeto,\footnote{Estas nociones
sobre la omisión del sujeto se amplían en el capítulo~6.} en otras lenguas,
como el inglés \exref{ex:01-sujnulo-c} y el francés \exref{ex:01-sujnulo-d},
una oración sin sujeto explícito está mal formada, es decir que resulta
agramatical.\footnote{El asterisco señala que se trata de una oración mal
formada o agramatical. El concepto de agramaticalidad se desarrolla en el
capítulo~2.}

\begin{exe}
  \ex\label{ex:01-sujnulo}
    \begin{xlist}
      \ex\label{ex:01-sujnulo-a} Comen perdices.
      \ex\label{ex:01-sujnulo-b} Mangiano le pernici.
      \ex\label{ex:01-sujnulo-c} \juicio{*}{Eat partridges} / They eat partridges.
      \ex\label{ex:01-sujnulo-d} \juicio{*}{Mangent des perdrix} / Ils mangent des perdrix.
    \end{xlist}
\end{exe}

¿Cómo podemos explicar estas diferencias? La variación que se observa en
\exref{ex:01-sujnulo} radica en la diferente forma en que se adopta el
principio que establece que todas las oraciones poseen un sujeto del cual se
predica. Como todo principio, este también tiene dos parámetros posibles de
realización (dos posiciones para el interruptor, siguiendo la metáfora
introducida arriba).

\begin{itemize}
  \item \textit{Opción 1}. El sujeto está presente obligatoriamente en todas
        las oraciones como sucede, por ejemplo, en inglés y en francés.
  \item \textit{Opción 2}. El sujeto puede omitirse, como, por ejemplo, en
        español y en italiano. Esta opción se conoce como \termino{parámetro
        del sujeto nulo} y las lenguas que permiten la omisión se denominan
        \textit{pro-drop} (con caída del pronombre sujeto).
\end{itemize}

\begin{recuadro}[Para recordar]
Los \termino{principios} son propiedades que tienen aplicación en todas las
lenguas, es decir que carecen de excepciones y constituyen la facultad misma
del lenguaje. Los \termino{parámetros} son las formas de realización de los
principios en las diferentes lenguas. La gramática de cada lengua no es más que
el conjunto de valores paramétricos que esa lengua adopta.
\end{recuadro}

\section{La estructura de la gramática}

Desde esta perspectiva la \termino{gramática} de una lengua es el mecanismo
cognitivo específico que permite formar adecuadamente las expresiones
lingüísticas encargadas de transmitir significados a través de sonidos. Y, dado
que son las estructuras del \termino{componente sintáctico}, al que
\textcite{chomsky1999} denomina el \termino{sistema computacional central}, las
encargadas de establecer la articulación entre ambos dominios ---el físico del
sonido y el conceptual del significado---, se explica la centralidad otorgada a
este componente en el esquema de gramática que plantea. La gramática también
posee \termino{dos componentes de interfaz}, \termino{forma fonética} y
\termino{forma lógica}, encargados, respectivamente, de determinar la
interpretación fonológica y semántica de las cadenas generadas por la sintaxis,
para lo cual conectan las representaciones abstractas producidas, con el
sistema articulatorio-perceptivo (que emite y percibe los sonidos) y con el
sistema conceptual-intencional (que interpreta y comunica los
significados).\footnote{El modelo presentado en \textcite{hauser2002}
distingue entre: a) la Facultad del Lenguaje en sentido amplio (FLA), que
incluye los sistemas «periféricos» de procesamiento (articulatorio-perceptivo)
y de conceptualización (conceptual-intencional); y b) la Facultad del Lenguaje
en sentido estricto (FLE), cuyo componente central es el dispositivo
computacional, dotado de la que Chomsky ha señalado como la principal
innovación del lenguaje: la recursividad.} El esquema de la gramática se
completa con el componente léxico ---también conocido como \termino{diccionario}
o \termino{lexicón mental}---, que almacena y proporciona las unidades léxicas
que necesita el componente sintáctico para alimentar las estructuras
lingüísticas que produce. La figura~\ref{fig:01-modelo} muestra el esquema
global de la gramática, según el modelo presentado en \textcite{chomsky1999}:

\begin{figure}[htbp]
  \centering
  \begin{tikzpicture}[
      font=\sffamily\small,
      caja/.style={draw=acento, thick, fill=acentoclaro, align=center,
                   minimum width=34mm, minimum height=9mm, inner sep=3pt},
      externo/.style={draw=gristexto, dashed, align=center,
                      minimum width=34mm, minimum height=9mm, inner sep=3pt,
                      text=gristexto},
      flecha/.style={-{Stealth[length=2.2mm]}, thick, acento},
      enlace/.style={{Stealth[length=2mm]}-{Stealth[length=2mm]}, gristexto}]
    \node[caja] (lex) at (0,0) {Léxico\\{\footnotesize (diccionario o lexicón mental)}};
    \node[caja] (sint) at (0,-1.9) {Sistema computacional\\{\footnotesize (componente sintáctico)}};
    \node[caja] (ff) at (-2.6,-3.9) {Forma fonética};
    \node[caja] (fl) at (2.6,-3.9) {Forma lógica};
    \node[externo] (ap) at (-2.6,-5.6) {Sistema\\articulatorio-perceptivo};
    \node[externo] (ci) at (2.6,-5.6) {Sistema\\conceptual-intencional};
    \draw[flecha] (lex) -- (sint);
    \draw[flecha] (sint) -- (ff);
    \draw[flecha] (sint) -- (fl);
    \draw[enlace] (ff) -- (ap);
    \draw[enlace] (fl) -- (ci);
  \end{tikzpicture}
  \caption{El modelo de la gramática}\label{fig:01-modelo}
\end{figure}

Si bien desde la perspectiva de la propuesta de Chomsky resulta perfectamente
entendible la centralidad otorgada a la sintaxis, en numerosas ocasiones esta
postura ha sido criticada por lingüistas que ponen el acento en los aspectos
semánticos\footnote{La corriente conocida como Semántica Generativa, que se
desarrolló entre 1968 y 1975, con representantes como G.~Lakoff, J.~Ross,
J.~McCawley y Ch.~Fillmore, entre otros lingüistas formados dentro del tronco
generativo, cuestionaba la centralidad de la sintaxis.} o en el valor
comunicativo de las combinaciones. Los cuestionamientos recibidos han ido
llevando a la teoría a ir otorgando cada vez más relevancia a las cuestiones
léxico-semánticas.\footnote{En \textit{Observaciones sobre la nominalización}
\parencite{chomsky1974b}, como respuesta a los cuestionamientos de los
semánticos generativos, Chomsky propuso la hipótesis lexicalista, que otorga
mayor prominencia a la palabra.} En este sentido ha sido fundamental el
\termino{principio de proyección} propuesto por \textcite{chomsky1981}, el cual
plantea que las estructuras sintácticas se proyectan a partir de las unidades
léxicas. Pero ¿qué implicaciones tiene este principio en la combinatoria de las
palabras?

\subsection{La proyección sintagmática de las palabras}\label{sec:01-proyeccion}

Para \textcite[21]{bosque2009} «podemos al abordar el análisis sintáctico
\enquote*{empezar por las oraciones}, como se hacía en la tradición, o bien
\enquote*{empezar por las palabras}, según se recomienda ahora». Los
mencionados autores se pronuncian por esta segunda opción fundados en que las
que se van combinando entre sí son las palabras, primero «en grupos pequeños y
luego en secuencias cada vez más complejas». Esos «grupos» en que se combinan
las palabras son las denominadas unidades intermedias de la gramática, o sea
los \termino{sintagmas}, que son la máxima proyección o expansión de un núcleo.
Y es el núcleo, según la clase de palabra a la que pertenece ---nombre (N),
verbo (V), adjetivo (A), adverbio (Adv), preposición (P)\footnote{\label{fn:01-lexicas}Nos
limitamos aquí a la proyección de las clases de palabras léxicas, que son las
que poseen contenido descriptivo inherente: \textit{árbol} (N), \textit{viejo}
(A), \textit{rápidamente} (Adv), \textit{correr} (V), \textit{para} (P). Las
palabras léxicas se diferencian de las palabras funcionales, como los
determinativos (\textit{el, este, un, mi}, etc.) y las conjunciones
(\textit{y, pero, o, que}, etc.), cuyo significado se relaciona con la función
gramatical que cumplen (v. \cite{giammatteo2009}).}---, el que le da su clase o
categoría al sintagma (SN, SV, SA, SAdv, SP). A su vez, los grupos así
conformados se seguirán recombinando para ir conformando secuencias cada vez
mayores, como \exref{ex:01-roma-a}, representada también en el diagrama
arbóreo \exref{ex:01-roma-b}:

\begin{exe}
  \ex\label{ex:01-roma}
    \begin{xlist}
      \ex\label{ex:01-roma-a} Francisco vivía en Roma.
      \ex\label{ex:01-roma-b}
        \begin{forest} gramatica
          [O
            [SN [Francisco, roof]]
            [SV
              [V [vivía]]
              [SP [P [en]] [SN [Roma, roof]]]]]
        \end{forest}
    \end{xlist}
\end{exe}

Dado que en la oración los elementos se van combinando desde el interior, es
decir, de derecha a izquierda, la lectura de \exref{ex:01-roma-b} procede de
abajo hacia arriba: en primer lugar, la P \forma{en} se une con su complemento
\forma{Roma}, formando un SP ---\forma{en Roma}---; luego ese SP se combina con
la forma verbal \forma{vive} y conforma el SV \forma{vive en Roma}, que, a su
vez, se combina con el SN \forma{Francisco} formando una oración. Asimismo,
dado que \forma{vive} concuerda con el SN \forma{Francisco} en persona y número
(ambos están en 3.ª persona singular), este último hará las veces de sujeto de
la oración resultante.\pendiente{El ejemplo dice \textit{vivía}; el texto,
\textit{vive}}

Ahora bien, ¿cualquier combinación de elementos léxicos es válida? ¿Qué pasa si
combinamos el SP \forma{en Roma} con el verbo \forma{nieva} y, al SV resultante
\forma{nieva en Roma}, lo unimos con el SN \forma{Francisco}? Si bien los SS
combinados siguen siendo de la misma clase de palabras, la oración
\juicio{*}{\forma{Francisco nieva en Roma}} no estaría bien formada porque
\forma{nevar} es un V impersonal que no admite un sujeto léxico como
\forma{Francisco}.

\begin{ejercicio}\label{ej:01-combinaciones}
Considere las siguientes combinaciones e intente explicar por qué resultan
anómalas o bajo qué condiciones podrían ser aceptables, según se muestra en el
ejemplo resuelto~1.

\begin{exe}
  \ex Francisco maúlla en Roma.
  \ex La lluvia vive en Roma.
  \ex Francisco tiene en Roma.
  \ex Francisco vive en triste.
  \ex Francisco reza la piedra.
  \ex La piedra reza en Francisco.
\end{exe}

\begin{resuelto}
1. El V \forma{maullar} exige un sujeto con los rasgos semánticos
\sig{animado} y \sig{felino}. La O sería aceptable si \forma{Francisco}
remitiera a un individuo con esas características.
\end{resuelto}
\end{ejercicio}

Los casos del ejercicio~\ref{ej:01-combinaciones} muestran que las unidades
léxicas tienen exigencias combinatorias tanto sintácticas como semánticas. En
cuanto a los requisitos sintácticos, vemos que la P requiere un SN como
complemento y que un V transitivo como \forma{tener} exige un complemento
objeto directo (COD) (v. capítulo~7). Respecto de las exigencias semánticas:
\forma{maullar} pide un sujeto animado \forma{gato} (o algún otro felino);
\forma{rezar} exige un COD del tipo \forma{oración} o \forma{plegaria}; la P
\forma{en} requiere no solo un SN como complemento, sino que este debe referir
a un lugar ---\forma{en Roma, en la colina}--- y también, con un V como
\forma{rezar}, puede indicar una lengua determinada ---\forma{Francisco reza en
latín}---. Cuestiones de este tipo, entre otras muchas en las que no podemos
extendernos aquí, pero que se irán desarrollando a lo largo de este libro,
hacen que la gramática no pueda quedarse exclusivamente en el estudio de las
formas, sino que también deba atender a los significados.

\begin{destacado}
En palabras de \textcite[13]{bosque2009}, «el análisis de la combinatoria
gramatical no puede \enquote*{poner entre paréntesis} el significado, sino más
bien todo lo contrario: es tarea esencial de la sintaxis explicar lo que
significan las combinaciones de palabras».
\end{destacado}

\subsection{Las palabras en contexto}

Ahora bien, ¿es lo que acabamos de plantear suficiente para dar cuenta
acabadamente de lo que manifiestan las expresiones lingüísticas? Veamos un
nuevo ejemplo en contexto:

\begin{exe}
  \ex\label{ex:01-francisco} Francisco vivía en Roma, en la casa de Santa
      Marta, un edificio adyacente a la Basílica de San Pedro. Durante un
      tiempo, en la misma Ciudad del Vaticano, convivió con su predecesor, el
      papa emérito, Benedicto XVI, que había fallecido en 2022.
\end{exe}

Aunque el sujeto y el verbo de la primera oración de \exref{ex:01-francisco}
siguen siendo los mismos de \exref{ex:01-roma}, nuestro conocimiento de mundo
---la enciclopedia que cada hablante tiene incorporada en su mente según sus
experiencias de vida--- y las pistas contextuales que aportan los elementos
léxicos ---\forma{la casa de Santa Marta, la Basílica de San Pedro, la Ciudad
del Vaticano, el papa emérito Benedicto XVI}--- nos llevan a asociar el nombre
\forma{Francisco} con el del papa fallecido en 2025. Si no fuéramos capaces de
establecer esta vinculación, nuestra interpretación de \exref{ex:01-francisco}
quedaría incompleta. Así, si bien dijimos que la perspectiva comunicativa que
excluye a la gramática no resultó eficaz para la enseñanza, esto no significa
que debamos desentendernos completamente del funcionamiento de las estructuras
en contexto. Como sostiene \textcite[s/p]{brucart2018},

\begin{cita}
[n]o integrar la gramática oracional en una perspectiva comunicativa equivale a
diseñar una ciudad con la trama urbana y los servicios preparados, pero sin
edificios ni principios constructivos.
\end{cita}

\subsection{Gramática y comprensión}

Si queremos integrar la gramática en la comprensión textual, es necesario
revisar las distintas concepciones acerca de la lectura a lo largo de la
historia. Según plantea \textcite{pozo2014}, en 5000 años de desarrollo, la
lectura ha ido cambiando, por lo que reconoce tres etapas:

\begin{itemize}
  \item La de \termino{lectura reproductiva o literal}, propia de los primeros
        tiempos, cuando aprender a leer era repetir y la función del lector no
        era encontrar el sentido de lo que leía, sino retener el texto
        almacenado en su mente.
  \item La de la \termino{lectura escolástica}, hacia el fin de la Edad Media,
        cuando la función de la lectura ya no era repetir lo que decía el
        texto, sino reproducir una interpretación autorizada de él.
  \item La de la \termino{lectura moderna}, que surge a partir del
        Renacimiento con la invención de la imprenta, cuando aparece el lector
        crítico o analítico. Se trata de un lector que va a interpretar,
        analizar, criticar, dar sentido y relacionarse epistémicamente con el
        texto.
\end{itemize}

A las tres etapas señaladas por Pozo, hoy día ya es necesario incorporar una
cuarta: la de \termino{la lectura hipertextual}, caracterizada por la
multimodalidad (que conecta con otros medios), la interactividad (que vincula
con otros usuarios), la flexibilidad (que permite otros recorridos) y la
intertextualidad (que relaciona con otros textos).

\figurapendiente{La lectura hipertextual}{fig:01-hipertextual}

En la actualidad, la comprensión lectora se considera un proceso formado por
múltiples procedimientos cognitivos que posibilitan al lector modificar su
pensamiento respecto de la realidad mediante la información que descubre en el
texto. Por tanto, lo que interesa en la enseñanza es formar lectores críticos
capaces de interpelar al texto, de interactuar con él y ser capaces de
relacionarlo con otros textos, para lo cual la simple habilidad de decodificar
grafemas\footnote{Los grafemas son las unidades mínimas de la escritura de una
lengua (\textit{DLE}, en línea).} no alcanza. Se requiere una alfabetización
integral que sume otros saberes «lectores» ---los aspectos enunciativos, los
géneros discursivos, los factores histórico-culturales de producción y
recepción---, pero sin excluir un conocimiento gramatical explícito, que ayude
a reconocer las «pistas» que ofrecen los textos para su cabal interpretación.

Pero ¿cómo se puede vincular las estructuras ---motor de la lengua--- con los
significados, que les dan sentido y, a su vez estos, con los valores
comunicativos que los hablantes les adjudican en sus distintos usos? Volvamos a
recurrir a una historieta. La viñeta de abajo muestra tres personajes: el
hablante, un niño aborigen; su interlocutora, a la que se dirige con el
apelativo familiar \forma{ma}; y un tercer personaje que lleva un traje de la
época de la conquista, a quien podríamos identificar con Cristóbal Colón. Pero
¿qué dice el texto que acompaña la imagen?

\figurapendiente{Fuente:
\url{https://m.facebook.com/866981573634920/photos/a.937101296622947/1274222342910839/?type=3}}{fig:01-colon}

Aunque breve, podemos decir que el texto de la viñeta constituye una
argumentación completa. Sin pretender un análisis exhaustivo, señalaremos
algunos de los principales recursos gramaticales que emplea el texto para
lograr su efectividad.

\begin{itemize}
  \item En la primera oración, el niño se vale del discurso referido
        indirecto\footnote{A diferencia del discurso referido directo, en el
        que el hablante reproduce literalmente lo dicho por otro ---p.~ej.
        \textit{Juan dijo: Ahora no puedo ir}---, en el indirecto el hablante
        ajusta lo dicho a su situación comunicativa ---\textit{Juan dijo que en
        ese momento no podía venir}---.} como estrategia\footnote{Se trata de un
        recurso evidencial mediante el cual el hablante expresa la fuente de la
        que deriva la información. El tema de modalidad y evidencialidad se
        desarrolla en el capítulo~14.} que le permite reproducir lo dicho por el
        conquistador ---\forma{hay un señor que dice que nos descubrió}--- y, al
        mismo tiempo, tomar distancia y desligar su responsabilidad respecto de
        un argumento con el que no concuerda.
  \item En el inicio de la segunda oración, hay una palabra
        funcional,\footnote{V. nota~\ref{fn:01-lexicas}.} la conjunción
        adversativa \forma{pero}, que conecta con la oración anterior e
        introduce el contraargumento mediante el cual el niño refuta lo dicho
        por el conquistador dejando en evidencia que no estaban dadas las
        condiciones para «ser descubiertos»,\footnote{La acepción 3 de
        \textit{descubrir} en el \textit{DLE}, en línea, dice: «Hallar lo que
        estaba ignorado o escondido, principalmente tierras o mares
        desconocidos».} ya que «no estaban escondidos».
  \item En la tercera oración, el adverbio deíctico \forma{acá} seguido del SP
        \forma{en casa}, que lo desarrolla,\footnote{Constituye una aposición
        explicativa (v. capítulo~5).} remite al contexto de situación y le
        permite al niño reafirmarse no solo en el lugar de hablante, sino en el
        de «dueño de casa». En el contexto de la historieta, nuestro
        conocimiento de mundo nos lleva a inferir que tanto \forma{acá} como
        \forma{en casa} refieren a América, el territorio que el «señor» dice
        haber descubierto. Asimismo, \forma{en casa}, junto con el apelativo
        \forma{ma} de la primera oración, remiten a un contexto familiar que
        evidencia el carácter de intruso del conquistador y desbarata sus
        dichos y pretensiones.
  \item La oración final contiene la conclusión que, en línea con el trato
        respetuoso que el niño le dispensa al conquistador al referirse a él
        como \forma{señor}, también se manifiesta en un tono morigerado, aunque
        no exento de sarcasmo: \forma{me parece que él está perdido}. La ironía
        que sobrevuela todo el texto remata en una doble posibilidad
        interpretativa para \forma{estar perdido}:
        \begin{itemize}
          \item \forma{está perdido}\textsubscript{1} porque «ha errado el
                camino», ya que, como se sabe, Colón se dirigía a las Indias y
                no a América,
          \item \forma{está perdido}\textsubscript{2} porque «está confundido»
                y no es cierto lo que afirma.
        \end{itemize}
\end{itemize}

En suma, como venimos señalando, los marcos teóricos de corte comunicativo que
hacen hincapié en el uso, pero se desentienden del estudio formal de los
fenómenos no son suficientes para la enseñanza, pero tampoco lo son los
enfoques formales cuando no tienen en cuenta que la lengua es un dispositivo
mental que los hablantes empleamos con múltiples intenciones para innumerables
fines. En este sentido, la gramática ha de servir para que los conocimientos
sobre los recursos lingüísticos estudiados puedan proyectarse en los textos y
colaborar en su comprensión.

\section{La gramática como descubrimiento}

Según venimos viendo, la lengua constituye un mecanismo de gran complejidad
diseñado para transmitir, mediante el uso de sonidos distintivos (fonemas o
bien grafemas, si se trata de la escritura), información (es decir,
significados) sobre el mundo y sobre nosotros mismos, para lo cual se vale de
una estructura con diferentes componentes interconectados.

El planteo central de la propuesta que se desarrolla en este libro, y que aquí
estamos solo presentando, se asienta en una base formal que permite dar cuenta
de las restricciones y correlaciones que conforman la combinatoria lingüística.

\begin{destacado}
Desde esta perspectiva, estudiar gramática no es otra cosa que hacer explícito
ese saber operativo (\textit{know how}, saber hacer o conocimiento tácito) que
el hablante común posee y que le permite comprender y producir las infinitas
oraciones de su lengua.
\end{destacado}

Sin embargo, como sostienen \textcite[69]{bosque2016}, por lo general, no se
enseña «a observar, no tanto lo anómalo o lo inusitado, sino más bien lo
sorprendente cuando se halla dentro de lo cotidiano y es relevante para
comprender las generalizaciones». Por eso, en esta propuesta el énfasis está
puesto en llamar la atención sobre el lenguaje y sus particularidades, a fin de
comenzar a reflexionar sobre él.

El punto de partida será la propia competencia, lo que como hablante nativo
cada uno «sabe» acerca de su lengua.

El objetivo es hacer consciente (explícito) ese conocimiento que tenemos
internalizado para que pueda ser utilizado eficazmente en distintas
situaciones, y así

\begin{itemize}[label=$\rightarrow$]
  \item poder comprender mejor lo que nos dicen, y
  \item ser mejor comprendidos en nuestros intercambios comunicativos.
\end{itemize}

Este proceso que va del «saber hacer» al conocer consciente comprende dos
momentos:

\begin{itemize}
  \item \termino{el «extrañamiento»}, que implica una toma de distancia de la
        inmediatez que todos tenemos respecto de la propia lengua, para poder
        considerarla como un objeto de estudio, cuyo abordaje implica
        categorías y criterios de análisis específicos;
  \item \termino{el «descubrimiento»}, que lleva a reconocer no solo qué
        formas existen en la lengua y qué restricciones y correlaciones de
        formación presentan, sino también cuáles resultan posibles y cuáles
        imposibles y por qué. Asimismo, estas mismas constricciones permitirán
        entender los aparentes «desvíos» en que incurren los hablantes (v. los
        ejemplos de \exref{ex:01-mente-nombres}).
\end{itemize}

El enfoque que se propone, al que denominamos léxico-gramatical porque se
asienta en estos dos componentes esenciales de la estructura lingüística, tiene
en cuenta que son los elementos léxicos, con una determinada configuración
fonológica y morfosemántica, los que, al operar en relación sintagmática, van
conformando unidades significativas mayores cada vez más amplias y abarcativas
\parencite{giammatteo2017}. Teniendo en cuenta las unidades resultantes, hemos
organizado la propuesta en torno a cuatro ejes principales:

\begin{figure}[htbp]
  \centering
  \begin{tikzpicture}[
      font=\sffamily\small,
      eje/.style={draw=acento, thick, fill=acentoclaro, align=center,
                  text width=25mm, minimum height=10mm, inner sep=3pt},
      unid/.style={align=center, text width=25mm, font=\sffamily\footnotesize},
      flecha/.style={-{Stealth[length=2.2mm]}, thick, acento}]
    \node[eje] (e1) at (0,0)    {\textbf{1.\textsuperscript{er} eje}\\Nivel subléxico};
    \node[eje] (e2) at (3.2,0)  {\textbf{2.\textsuperscript{do} eje}\\Nivel léxico};
    \node[eje] (e3) at (6.4,0)  {\textbf{3.\textsuperscript{er} eje}\\Nivel sintagmático};
    \node[eje] (e4) at (9.6,0)  {\textbf{4.\textsuperscript{to} eje}\\Nivel textual};
    \node[unid, below=2mm of e1, anchor=north] {Fonemas\\Morfemas};
    \node[unid, below=2mm of e2, anchor=north] {Palabras};
    \node[unid, below=2mm of e3, anchor=north] {Sintagmas\\Oraciones};
    \node[unid, below=2mm of e4, anchor=north] {Párrafos\\Secuencias\\Textos};
    \draw[flecha] (e1) -- (e2);
    \draw[flecha] (e2) -- (e3);
    \draw[flecha] (e3) -- (e4);
  \end{tikzpicture}
  \caption{Ejes de la propuesta}\label{fig:01-ejes}
\end{figure}

En lo que sigue, se presentan, a modo de ejemplos, algunos casos que se podrían
considerar a fin de producir el «extrañamiento» respecto del uso cotidiano y
propiciar, a través de la reflexión, el «descubrimiento» de sus posibilidades e
imposibilidades combinatorias, lo cual debe conducir a reconocer el
funcionamiento de las estructuras y de sus significados, y a identificar sus
valores comunicativos y expresivos.

\subsection{Del interior de la palabra al reconocimiento de sus valores en contexto}

\subsubsection{Restricciones y correlaciones combinatorias en el nivel subléxico}\label{sec:01-sublexico}

En el nivel subléxico referido al sonido se encuentran los \termino{fonemas},
que son unidades fónicas sin significado en sí mismas, pero cuyo contraste
permite distinguir significados como, por ejemplo, entre
/ˈpaso/\footnote{En la transcripción fonológica, los fonemas convencionalmente
van entre barras inclinadas y el acento se coloca delante de la sílaba (v.
capítulo~3).} y /ˈpeso/ (la diferencia es la presencia de
/-a-/\footnote{Los guiones indican que una forma no es libre, sino que tiene
algún contexto antes /-a/ o después /a-/, o bien en ambas posiciones /-a-/.}
frente a /-e-/) o entre /ˈpeso/ y /ˈbeso/ (la diferencia se da entre /p-/ y
/b-/). En cuanto a la combinatoria de estas unidades, podemos decir que la
estructura fonológica del español admite algunas combinaciones y rechaza otras;
así, por ejemplo, es posible formar cadenas como \forma{\textbf{tr}aza} o
\forma{\textbf{tr}ama}, en las que interviene el grupo consonántico
\forma{tr-}, pero nunca podríamos encontrar \juicio{*}{\forma{\textbf{sr}aza}}
o \juicio{*}{\forma{\textbf{sr}ama}}, ya que la combinación \forma{sr-} no
constituye un patrón fonológico de nuestra lengua. A veces, para responder a los
esquemas propios de la lengua, se necesitan ajustes fonológicos: por ejemplo,
algunas palabras provenientes del latín o de otras lenguas que admiten grupos
consonánticos iniciales que en español no se aceptan, incorporan una /e-/
delante, para asimilarse a la norma de pronunciación del
español:\footnote{El signo > se lee: «se convierte en».}

\begin{exe}
  \ex\label{ex:01-eprotetica}
    \begin{tabular}[t]{@{}l l c l l@{}}
      \forma{schola}  & (lat.)  & > & \forma{e-scuela} (esp.) & (ingl. \forma{school}) \\
      \forma{stilus}  & (lat.)  & > & \forma{e-stilo} (esp.)  & (ingl. \forma{style}) \\
      \forma{stress}  & (ingl.) & > & \forma{e-strés} (esp.)  & \\
      \forma{scanner} & (ingl.) & > & \forma{e-scáner} (esp.) & \\
    \end{tabular}
\end{exe}

Al combinarse dentro de la palabra, los fonemas van formando unidades
significativas\footnote{Según se verá en el capítulo~4, el significado de los
morfemas puede ser léxico (de contenido) o puramente gramatical (funcional). En
ese capítulo también se desarrolla la noción de morfema y los conceptos
relacionados con la estructura interna de la palabra.} denominadas
\termino{morfemas}; así, a una cadena como /ˈutil/ (= \sig{que produce
beneficio}\footnote{Convencionalmente los significados se colocan entre comillas
simples.}) se le puede añadir un prefijo de significado negativo como /in-/
para formar /inˈutil/ (\sig{que no produce beneficio}) y luego, si añadimos el
sufijo \forma{-mente}, el resultado será /inˈutilmente/ (\sig{de manera inútil
o que no produce beneficio}). Pero también existen restricciones para la
combinación de morfemas: con el mismo mecanismo de \forma{inútilmente} también
podemos formar \forma{alegremente, rápidamente, infantilmente} (\sig{de manera
alegre, rápida o infantil}), pero no resultan posibles combinaciones como
\juicio{*}{\forma{vasomente}}, \juicio{*}{\forma{libromente}} o
\juicio{*}{\forma{mesamente}}, ya que el sufijo \forma{-mente} se puede unir a
adjetivos, como \forma{alegre, rápido} o \forma{infantil}, pero no a
sustantivos, como \forma{vaso, libro} o \forma{mesa}.

Estos conocimientos sobre la conformación de las palabras los tenemos
internalizados y, por lo general, no reflexionamos sobre ellos, los usamos
automáticamente y no dudamos en emplearlos cuando formamos nuevas palabras;
así, a partir del adjetivo coloquial \forma{trucho}, no titubeamos en formar el
derivado \forma{truchamente} (\sig{en forma trucha}, es decir, \sig{falsa o
ilegal}).

De este modo, según vamos viendo, el hilo conductor de la reflexión sobre la
lengua son los propios juicios acerca de la buena o mala formación de las
distintas construcciones, es decir, sobre su gramaticalidad o agramaticalidad.
Pero muchas veces estos conceptos también deberán ser complementados con la
noción de aceptabilidad que, en ciertos contextos, puede permitir extender los
límites de formación (v. capítulo~2), tal como muestran los ejemplos
siguientes, en los que, contrariamente a lo que hemos dicho más arriba, el
sufijo \forma{-mente} aparece unido a sustantivos:

\begin{exe}
  \ex\label{ex:01-mente-nombres}
    \begin{xlist}
      \ex Llegará, \forma{tanga\textbf{mente}}, mi muerte enamorada (H.~Ferrer,
          \textit{Balada para mi muerte}).
      \ex Zarpar en el barco sin mirar atrás, subirse a la patera que aguantará
          sólo mientras pueda,\ldots{} Las mayorías soportan
          \forma{golondrina\textbf{mente}} un presente atroz y aplastador. O
          terminan mansamente confinados como esclavos en talleres textiles
          \ldots{} (C.~Rafael, \textit{La tos de la
          esclavitud}).\pendiente{Falta el dominio del enlace en el original}
      \ex Muestra de esta insensibilidad es cuando por ejemplo estos bancos
          neoliberales provocan una y otra vez las olas especulativas en los
          tipos de cambio\ldots{} lo que les permite a través del Mercado de
          Divisas mover sus reservas y utilidades hacia donde
          \forma{golondrina\textbf{mente}} les conviene (A.~Tirado Pazos, 2016,
          \textit{La paradoja de México: Del neoliberalismo al liberalismo
          social},
          \url{https://www.ruizhealytimes.com/economia-y-negocios/la-paradoja-de-mexico-del-neoliberalismo-al-liberalismo-social}).
    \end{xlist}
\end{exe}

¿Qué explicación podemos ofrecer para estas «aparentes» anomalías? A diferencia
de \juicio{*}{\forma{vaso\textbf{mente}}}, \juicio{*}{\forma{libro\textbf{mente}}}
o \juicio{*}{\forma{mesa\textbf{mente}}}, en los que es imposible asignar
significado al conjunto (\juicio{*}{\sig{de manera vaso, libro o mesa}}), los
ejemplos de \exref{ex:01-mente-nombres}, en cambio, son aceptables para los
hablantes porque resultan interpretables en un contexto determinado. Así,
\forma{tangamente} significa \sig{de manera tanguera} y
\forma{golondrinamente} se entiende en el sentido de \sig{como se comportan las
golondrinas migratorias, en forma estacional}.

\begin{ejercicio}\label{ej:01-diccionariamente}
Lea atentamente el siguiente fragmento y responda si es o no aceptable la
palabra subrayada en el texto. Si la respuesta es afirmativa, explique cuál es
su significado.

\begin{cita}
Miller también fue pionero del estudio de cómo los chicos aprenden palabras,
no \subrayado{diccionariamente} sino en contexto: nadie puede estudiar desde el
diccionario (D.~Golombek, «Las palabras (y las cosas). Y los viajes por el
léxico», \textit{La Nación}, 22/01/2017).
\end{cita}
\end{ejercicio}

\subsubsection{Comparación y contraste: los pares mínimos}

Según estamos viendo, entonces, se parte de una focalización que lleva a
descubrir a través de ejemplos, las estructuras que subyacen a la lengua, lo
que luego podrá hacerse extensivo a otros casos similares para llegar, si cabe,
a establecer generalizaciones. El proceso cognitivo en que se basa «el
descubrimiento» es la comparación, que permite reconocer, a través del cotejo
de los que se conocen como \termino{pares mínimos}
\parencite{bosque2016,brucart2018}, qué comparten y en qué se diferencian las
construcciones examinadas.\footnote{En \textcite{bosque2016} se plantea que los
pares mínimos pueden ser dos secuencias sencillas y gramaticales que se
diferencian en un solo factor o dos secuencias, una de las cuales es
agramatical.} Aunque son muy utilizados en fonología, ya que sirven para
distinguir sonidos contrastivos como los fonemas
(v.~§\ref{sec:01-sublexico}), los pares mínimos no suelen emplearse para
establecer comparaciones en otros niveles de la estructura lingüística. Para
mostrar su utilidad, a modo de ejemplo, podemos comparar entre sí los
siguientes pares de palabras cuyas estructuras ---un prefijo seguido de una
base\footnote{El concepto de base es relacional: es de donde se parte para
hacer una derivación; por ejemplo, en \textit{inútil} la base es \textit{útil},
pero en \textit{inútilmente}, la base es \textit{inútil} (v. capítulo~4:
§6.1).}--- pueden parecer idénticas, aunque, como veremos, no lo son:

\begin{exe}
  \ex\label{ex:01-despistar} \forma{descentrar} / \forma{despistar}
  \ex\label{ex:01-autopista} \forma{autorretrato} / \forma{autopista}
\end{exe}

En \exref{ex:01-despistar}, ambas palabras están formadas por un sustantivo
respectivamente, \forma{centro} y \forma{pista}, a los que se les adjunta el
mismo prefijo \forma{des-}. Sin embargo, solo en la primera es posible derivar
un verbo del sustantivo \forma{centro} > \forma{centrar}, al que se le puede
añadir el prefijo \forma{des-}, que implica la realización de la \sig{acción
opuesta} \forma{centrar} ≠ \forma{descentrar}, y formar así una palabra
derivada. En \forma{despistar}, en cambio, no podemos partir de un inexistente
verbo \juicio{*}{\forma{pistar}}, lo que hace que \forma{despistar} provenga
directamente del nombre \forma{pista}, al que se le adjuntan, simultáneamente,
un prefijo y un sufijo, según el proceso de formación de palabras que se conoce
como parasíntesis (v. capítulo~4: §7.3).

\begin{exe}
  \ex\label{ex:01-parasintesis}
    \begin{tabular}[t]{@{}l@{\quad}l@{}}
      \forma{centro} > \forma{centr\textbf{ar}} > \forma{\textbf{des}centrar} & (palabra derivada) \\
      \forma{pista} > \forma{\textbf{des}pist\textbf{ar}}                     & (palabra parasintética) \\
    \end{tabular}
\end{exe}

En cuanto a los significados, ambas palabras constituyen el tipo de opuestos
que \textcite[18]{cruse2001} denomina \termino{reversos}: «pares tales como:
\textit{entrar: salir, subir: bajar, avanzar: retroceder, ascender: descender,
aparecer: desaparecer, unir: desunir, vestir: desvestir}. Todos ellos denotan
cambios de estado de algún tipo que son, en principio, reversibles». Desde esta
perspectiva, \forma{descentrar} es el proceso inverso de \forma{centrar}, o sea
«sacar del centro» y, de modo semejante, aunque \juicio{*}{\forma{pistar}} no
exista, \forma{despistar}, en una primera acepción se entiende como su opuesto:
«hacer perder la pista» (\textit{DLE}, en línea), pero, además,
\forma{despistar} adopta un valor derivado metafóricamente: «distraer,
desorientar o confundir» (\textit{DLE}, en línea), ya que quien se sale de
pista queda, por lo general, desorientado o confundido.

En cuanto a las palabras de \exref{ex:01-auto-estructura},
\forma{autorretrato}, formada por la unión del prefijo \forma{auto-} y el
sustantivo \forma{retrato}, se corresponde con el significado esperado de
\sig{retrato de uno mismo}. Pero si aplicamos el mismo procedimiento en
\forma{autopista}, nos lleva a un significado inexistente \juicio{\#}{\sig{pista
de uno mismo}}. En este segundo caso, como en otros muchos similares en español
(\forma{bicisenda, homofóbico},\footnote{Se trata también de palabras
compuestas con un acortamiento en el primer constituyente:
\textit{bici(cleta)senda} (= «senda para bicicletas»);
\textit{homo(sexual)fóbico} (= «el que padece fobia a los homosexuales»).}
etc.), \forma{auto-} representa un acortamiento de una palabra más extensa, el
compuesto \forma{automóvil}; de modo que la interpretación correcta es
\sig{pista de/para automóvil(es)}.

\begin{exe}
  \ex\label{ex:01-auto-estructura}
    \begin{xlist}
      \ex Palabra prefijada: \forma{retrato} > \forma{\textbf{auto}rretrato}
          (= «retrato de uno mismo»)
      \ex Palabra compuesta con acortamiento del primer constituyente:
          \forma{automóvil} + \forma{pista} > \forma{auto(móvil)pista}
          (= «pista para automóvil(es)»)
    \end{xlist}
\end{exe}

Ahora bien, ¿se pueden conectar estos planteos, focalizados y centrados en el
nivel de la palabra, con el uso real de la lengua y la interpretación en
contexto? Los siguientes ejemplos, tomados de la prensa, pueden servir para
mostrar cómo el hablante puede hacer uso eficaz de su conocimiento de la lengua
para generar y reconocer nuevas formas.

\begin{exe}
  \ex\label{ex:01-descomer}
    \begin{xlist}
      \ex\label{ex:01-descomer-a} «Contratar y despedir debería ser natural como
          \forma{comer y descomer}». El debate por una presunta estrategia de
          «flexibilización laboral» por parte del Gobierno se deslizó en las
          últimas horas en medio de polémicas declaraciones de funcionarios del
          área de Trabajo. El secretario de Empleo y ex CEO del grupo Techint,
          Miguel Ángel Ponte, opinó que incorporar o despedir personal debería
          ser para las empresas natural como «comer y descomer»
          (\url{https://www.perfil.com/noticias/politica/ponte-contratar-y-despedir-deberia-ser-natural-como-comer-y-descomer.phtml},
          09/01/17).
      \ex\label{ex:01-descomer-b} \textsc{Flexibilización laboral, pasantías y
          reducción de aportes patronales}. La teoría «comer y descomer». El ex
          Techint y actual secretario de Empleo, Miguel Ponte, ilustró con que
          la movilidad en el mercado laboral tiene que ser similar al proceso de
          ingesta y evacuación de comida
          (\url{https://www.pagina12.com.ar/tags/empleo},
          21/01/17).\pendiente{Fecha confusa en el original}
    \end{xlist}
\end{exe}

El ejemplo \exref{ex:01-descomer-a} incluye una palabra, \forma{descomer}, que
en su momento llamó la atención de la ciudadanía en relación con el planteo
respecto de la flexibilización laboral.\footnote{La flexibilización laboral o
desregulación del mercado de trabajo es un modelo que elimina regulaciones
para contratar y despedir empleados por parte de las empresas y organizaciones
privadas.} Este desconcierto se refleja en \exref{ex:01-descomer-b}, donde se
trata de descifrar el significado de la palabra a partir de su relación con su
opuesto reverso \parencite[18]{cruse2001} \forma{comer}. Si \forma{comer} es
\sig{ingerir alimentos}, el proceso inverso, como bien se aclara en
\exref{ex:01-descomer-b}, es su \forma{evacuación} o «expulsión de excrementos
del interior del cuerpo». Ahora bien, ¿si ya existe en la lengua una palabra
como \forma{evacuación} para dar cuenta de este proceso, por qué se elige
«crear» otra? Indudablemente en una cuestión sensible para la sociedad, como la
flexibilización laboral, el hablante busca atenuar la expresión evitando un
término tabú y creando otro, que en este contexto cuenta como un eufemismo.

\begin{recuadro}[Para tener en cuenta]
Aun elementos menores, del nivel subléxico, como los afijos, pueden hacer
aportes significativos a la comprensión del texto, por lo que, además de
considerar sus significados descontextualizados, es importante integrarlos al
sentido global del texto, el cual es producto de las relaciones y oposiciones
que cada constituyente entabla con el resto en un contexto determinado
\parencite{giammatteo2015}.
\end{recuadro}

\subsection{Del nivel léxico a las relaciones sintagmáticas}

\subsubsection{Las conexiones sintácticas entre las palabras}

Para el hablante común, la palabra presenta una «visibilidad» de la que carecen
otras unidades lingüísticas, como podrían ser el morfema o la oración, lo que le
confiere una mayor apariencia de realidad psicológica y prominencia cognitiva
en la conciencia \parencite{julien2007}. Además, como sostiene
\textcite[49]{jakobson1967}, los ítems léxicos son entidades de doble vía
porque «están ligados necesariamente con el código por una relación interna y
con el mensaje por una relación externa». Esta doble relación implica dos
posibles orientaciones:

\begin{itemize}
  \item \textit{hacia el interior del sistema lingüístico}, las palabras
        constituyen una interfaz \parencite{jackendoff1990} que conecta todos
        los componentes de la estructura lingüística ---fonológico,
        morfosintáctico, semántico---; y
  \item \textit{hacia el exterior}, las palabras se proyectan en los múltiples
        mensajes con los que nos comunicamos, relacionando así el sistema
        lingüístico con el conocimiento de mundo.
\end{itemize}

Dado que son los elementos léxicos los que se proyectan sintagmáticamente para
formar unidades mayores (v.~§\ref{sec:01-proyeccion}), también en este nivel
podemos valernos de comparaciones para identificar las reglas que rigen la
formación de sintagmas y oraciones. Así, \exref{ex:01-concordancia} y
\exref{ex:01-modo} muestran la importancia fundamental de las categorías
gramaticales que, aunque se manifiestan morfológicamente en el interior de la
palabra, su posición periférica en ella indica que su función va más allá de la
palabra, ya que consiste en establecer conexiones entre los constituyentes
sintagmáticos (v. capítulo~4: §9):

\begin{exe}
  \ex\label{ex:01-concordancia}
    \begin{xlist}
      \ex L-o-s extenuad-o-s viajer-o-s recorri-ero-n el parque.
      \ex[*]{El extenuad-o viajer-o-s recorriero-n el parque.}
      \ex[*]{L-o-s extenuad-o-s viajer-o-s recorri-ó el parque.}
    \end{xlist}
  \ex\label{ex:01-modo}
    \begin{xlist}
      \ex\label{ex:01-modo-a} Creo que vienen / \juicio{*}{vengan}.
      \ex\label{ex:01-modo-b} No creo que vengan / \juicio{*}{vienen}.
    \end{xlist}
\end{exe}

En \exref{ex:01-concordancia}, los rasgos \sig{masculino} y \sig{plural} del
sustantivo \forma{viajer-o-s} se repiten en todos los elementos que integran el
SN, como el artículo \forma{l-o-s} y el adjetivo \forma{extenuad-o-s};
asimismo, el rasgo \sig{plural} también se retoma en el verbo
\forma{recorriero-n}, como señal de su vinculación con el núcleo del sujeto
\forma{viajero-s}.\pendiente{El original dice «‘masculino’ y ‘singular’»} En
\exref{ex:01-modo} el modo verbal es una categoría exigida por la correlación
entre el verbo de la oración principal ---\forma{creo}--- y el de la
subordinada: cuando el verbo principal es afirmativo requiere modo indicativo
\forma{vienen} \exref{ex:01-modo-a} y rechaza el subjuntivo; y viceversa
respecto de la forma negativa del verbo, que exige subjuntivo \forma{vengan}
\exref{ex:01-modo-b}. Cualquier alteración de estas correlaciones convierte a
las estructuras resultantes en agramaticales.

\subsubsection{La interfaz sintaxis-semántica}

En relación con las oraciones, es necesario considerar que conforman
estructuras típicas, a las que denominamos \termino{esquemas oracionales}
transitivos, intransitivos, reflexivos, pasivos, impersonales, etc. Estas
configuraciones se encargan de transmitir, mediante los recursos
léxico-gramaticales de la lengua, las representaciones mentales de los
distintos sucesos y situaciones que experimentamos, deseamos o simplemente
imaginamos. Se generan en torno a los verbos, que son las palabras que refieren
a los sucesos que se describen en la oración. En consecuencia, son las
exigencias de los verbos las que determinan la configuración oracional
resultante (aunque, según veremos, un mismo verbo puede adoptar distintas
configuraciones).

Teniendo lo anterior como punto de partida, podemos plantear que toda oración
se organiza en torno a tres estructuras interrelacionadas:

\begin{itemize}
  \item \termino{la estructura eventiva}, que tiene en cuenta el tipo de verbo
        que constituye el núcleo de la oración: de estado, de acción, de
        proceso,\footnote{V. nota~\ref{fn:01-accion}.} etc.;
  \item \termino{la estructura argumental}, constituida por las funciones
        semánticas, conocidas como papeles temáticos, que asumen los argumentos
        que el verbo exige en función de su significado: agente, tema/paciente,
        experimentante, etc.;
  \item \termino{la estructura sintáctica}, que da cuenta de las funciones que
        cumplen los distintos constituyentes presentes en la
        oración.\footnote{La estructura sintáctica y su relación con la
        estructura argumental se trata en el capítulo~7 y la estructura
        eventiva, en el capítulo~8.}
\end{itemize}

Veamos solo algunos ejemplos sencillos para ilustrar estos planteos:

\begin{exe}
  \ex\label{ex:01-argumentos}
    \begin{xlist}
      \ex\label{ex:01-argumentos-a} Tres hombres robaron una tienda de productos electrodomésticos.
      \ex\label{ex:01-argumentos-b} Un intendente del conurbano entregó viviendas a los habitantes del municipio.
      \ex\label{ex:01-argumentos-c} Los cimientos cedieron y la casa se derrumbó.
    \end{xlist}
\end{exe}

Tanto en \exref{ex:01-argumentos-a} como en \exref{ex:01-argumentos-b} se trata
de verbos de \sig{acción} ---\forma{robar, entregar}---, pero mientras el
primero exige solo dos participantes, un agente que realiza la acción
---\forma{tres hombres}--- y un tema afectado por el suceso ---\forma{una tienda
de productos electrodomésticos}---, el segundo verbo, además del agente
---\forma{un intendente del conurbano}--- y del tema ---(las)
\forma{viviendas}---, requiere también un tercer participante en el evento, la
meta o término de la acción ---\forma{los habitantes del municipio}---. Pero
¿qué ocurre en \exref{ex:01-argumentos-c}? Los sucesos que se describen no se
presentan como resultado de la acción de ningún agente, sino que se trata de
procesos ---\forma{cedieron} y \forma{se derrumbó}--- que afectan a un solo
argumento cada uno respectivamente ---\forma{los cimientos} y \forma{la
casa}---, que cumplen la función de temas afectados por los eventos. Si bien no
podemos extendernos mayormente en estas consideraciones, que serán desarrolladas
en los capítulos correspondientes de este libro, los ejemplos presentados
muestran que los verbos difieren en sus exigencias argumentales y, lo que
resulta más interesante, es que estas diferencias no son aleatorias, sino que,
en gran medida, derivan de sus diferencias semánticas.

Asimismo, desde la perspectiva sintáctica, estos significados que conforman la
estructura argumental del verbo se reflejan en la organización de la estructura
oracional, estableciendo configuraciones canónicas o por defecto para los
distintos tipos de verbos. Así, los verbos de \sig{acción}, como \forma{robar}
en \exref{ex:01-argumentos-a}, típicamente involucran dos argumentos, agente y
tema, que funcionan, respectivamente, como sujeto y COD. \forma{Entregar}, en
cambio, forma parte del subgrupo de verbos de \sig{acción} que implican una
\sig{transferencia}, integrado también por verbos como \forma{dar, prestar,
regalar} o \forma{ayudar}, entre otros. Los verbos de \sig{transferencia} son
ditransitivos (o doblemente transitivos) y exigen, como vimos en
\exref{ex:01-argumentos-b}, tres argumentos: el agente y el tema funcionan
sintácticamente igual que con los transitivos, es decir, como sujeto y COD
respectivamente; pero el tercer argumento, la meta, cumple la función
sintáctica de COI. En el caso de los verbos de \exref{ex:01-argumentos-c}, que
manifiestan procesos que se presentan como no agentivos, el único argumento
exigido es el tema afectado, que en estos casos cumple la función sintáctica de
sujeto del verbo, como lo muestra en \exref{ex:01-argumentos-c} la
concordancia: plural en \forma{cedieron} porque el sujeto es \forma{los
cimientos}, y singular en \forma{se derrumbó}, cuyo sujeto es \forma{la casa}.

En función de los planteos que estamos presentando de modo ilustrativo y no
exhaustivo, ya que la mayoría se irá ampliando a lo largo de los capítulos que
siguen, resulta también importante resaltar que las distintas estructuras
oracionales poseen, además de la forma canónica, otras configuraciones
alternativas que le permiten al hablante mostrar distintas «visiones» o
perspectivas respecto de los sucesos que presenta. Dentro de estas
posibilidades de alternancia la más reconocida es la que se da entre la voz
activa y la pasiva, que en español puede manifestarse utilizando tanto la
pasiva perifrástica \exref{ex:01-pasivas-a}\footnote{La voz pasiva y demás
perífrasis verbales del español se desarrollan en el capítulo~11; y la pasiva
con \textit{se} y otras construcciones pronominales, en el capítulo~10.} como
la llamada pasiva con \forma{se} \exref{ex:01-pasivas-b}. Ambas pasivas
permiten promover al argumento tema/paciente\footnote{Ambos, tema y paciente,
representan al mismo argumento, pero se emplea la denominación de
\textit{tema} cuando el sustantivo que lo manifiesta posee el rasgo \sig{no
animado}; y la de \textit{paciente}, cuando es \sig{animado}.} a la posición
más prominente de la oración: la de sujeto, y mover al agente a una posición
marginal \exref{ex:01-pasivas-a} o incluso eliminarlo
\exref{ex:01-pasivas-b}.

\begin{exe}
  \ex\label{ex:01-pasivas}
    \begin{xlist}
      \ex\label{ex:01-pasivas-a} Una tienda de productos electrónicos fue robada (por tres hombres).
      \ex\label{ex:01-pasivas-b} Se entregaron viviendas a los habitantes del municipio
          \juicio{*}{por el intendente}.
    \end{xlist}
\end{exe}

Además de la alternancia activa/pasiva, la lengua presenta para los verbos
distintas formas de construcción correlativas de diferentes maneras de enfocar
los sucesos, tal como se muestra en los ejemplos de \exref{ex:01-perspectiva}:

\begin{exe}
  \ex\label{ex:01-perspectiva}
    \begin{xlist}
      \ex\label{ex:01-perspectiva-a} Un fuerte temporal de lluvia y viento azotó varios distritos de
          la provincia de Buenos Aires tras la intensa ola de calor que sufrió
          gran parte del país
          (\url{https://www.lanacion.com.ar/sociedad/temporal-en-chascomus-la-desesperacion-de-la-gente-que-fue-sorprendida-por-la-tormenta-nid},
          16/01/22).
      \ex\label{ex:01-perspectiva-b} Se derrumbó una casa en Rosario y buscan a un hombre de 52 años
          perdido entre los escombros
          (\url{https://www.perfil.com/noticias/actualidad/se-derrumbo-una-casa-en-rosario-y-buscan-a-un-hombre-de-52-anos-perdido-entre-los-escombros.phtml},
          26/07/22).
      \ex\label{ex:01-perspectiva-c} Las casas fabricadas con este material se construyen
          rápidamente y se necesita personal especializado para colocar
          tornillos y soldar las uniones
          (\url{https://mejorconsalud.as.com/casas-prefabricadas-debes-saber-comprar-una/}).
    \end{xlist}
\end{exe}

Si comparamos \exref{ex:01-perspectiva-a} con \exref{ex:01-argumentos-c}, en la
que los dos verbos refieren procesos que se presentan como no agentivos, el
primero de ellos, \exref{ex:01-perspectiva-a}, introduce una variante de este
papel temático, la causa ---\forma{un fuerte temporal de lluvia y viento}---,
que, en este caso, representa al fenómeno natural que provoca el hecho y que,
además, puede cumplir la función sintáctica de sujeto de la oración. En
\exref{ex:01-perspectiva-b}, el verbo transitivo \forma{derrumbar} se
intransitiviza mediante la pronominalización que lo convierte en
\forma{derrumbarse}, lo que habilita una lectura según la cual el suceso se
presenta como no provocado por ningún agente o causa, de modo que la oración
puede interpretarse \forma{La casa se derrumbó sola}. En
\exref{ex:01-perspectiva-c}, aunque \forma{construir} sea también un verbo de
\sig{acción} transitivo que supone un agente, se diferencia de
\exref{ex:01-perspectiva-b} porque hay en la oración ciertas marcas gramaticales
específicas, como el uso del presente y la presencia de un adverbio con valor
modal, que hacen que el agente se interprete como indeterminado y la lectura sea
de tipo genérico: \forma{Cualquiera puede construir rápidamente / en un abrir y
cerrar de ojos estas casas}.

Todas estas diferentes posibilidades de construcción, que configuran los
distintos esquemas oracionales estables de la lengua, los cuales pueden ser
reconocidos a partir de las marcas léxicas y gramaticales con que se
construyen, no pueden ser ignoradas cuando estudiamos una lengua. No se trata
solo de describir estructuras y de realizar minuciosos análisis de sus
componentes, sino de reflexionar acerca de la amplia gama de matices de
significación que estas estructuras permiten manifestar; se trata, en suma, de
que el conocimiento de estos esquemas sintácticos no se reduzca a un tema a
memorizar para la evaluación de la materia, sino de que se integre al bagaje de
estrategias del que podamos disponer como hablantes competentes de la lengua,
ya que este tipo de estructuras son las que a diario encontramos en los textos
con los que debemos interactuar.

\section{El sistema en uso: las pistas gramaticales en el texto}

En relación con el último eje planteado no vamos a hacer aquí un desarrollo
extenso, pero, al menos, nos interesa aportar un ejemplo, que nos permitirá
comprobar si el enfoque que estamos planteando, al considerar conjuntamente el
funcionamiento de los ítems léxicos y las estructuras morfosintácticas, resulta
adecuado no solo para aumentar el poder descriptivo y explicativo del análisis,
sino también para orientar la comprensión lectora y «convertir el saber
conceptual sobre el sistema de la lengua, en uso en géneros discursivos
concretos» \parencite{fontich2006}.

Para ilustrar estos conceptos, el texto elegido es el microrrelato de Jorge
Luis Borges \textit{El cautivo}, sobre el que no proponemos realizar un
análisis exhaustivo, sino simplemente poner de relieve cómo los recursos
léxico-gramaticales de la lengua intervienen en su construcción y, por tanto,
para su comprensión resulta útil prestar atención a su funcionamiento. Leamos
primero el texto:

\begin{textolectura}[El cautivo, de J. L. Borges]
En Junín o en Tapalqué refieren la historia. Un chico desapareció después de un
malón; se dijo que lo habían robado los indios. Sus padres lo buscaron
inútilmente; al cabo de los años, un soldado que venía de tierra adentro les
habló de un indio de ojos celestes que bien podía ser su hijo. Dieron al fin con
él (la crónica ha perdido las circunstancias y no quiero inventar lo que no sé)
y creyeron reconocerlo.

El hombre, trabajado por el desierto y por la vida bárbara, ya no sabía oír las
palabras de la lengua natal, pero se dejó conducir, indiferente y dócil, hasta
la casa. Ahí se detuvo, tal vez porque los otros se detuvieron. Miró la puerta,
como sin entenderla. De pronto bajó la cabeza, gritó, atravesó corriendo el
zaguán y los dos largos patios y se metió en la cocina. Sin vacilar, hundió el
brazo en la ennegrecida campana y sacó el cuchillito de mango de asta que había
escondido ahí, cuando chico. Los ojos le brillaron de alegría y los padres
lloraron porque habían encontrado al hijo.

Acaso a este recuerdo siguieron otros, pero el indio no podía vivir entre
paredes y un día fue a buscar su desierto. Yo querría saber qué sintió en aquel
instante de vértigo en que el pasado y el presente se confundieron; yo querría
saber si el hijo perdido renació y murió en aquel éxtasis o si alcanzó a
reconocer, siquiera como una criatura o un perro, los padres y la casa.
\end{textolectura}

Empecemos por considerar algunas cuestiones relativas al encuadre general del
texto. \textit{El cautivo} se presenta como un relato construido por un
narrador no omnisciente que, como es típico en Borges, simula no conocer todo
lo que concierne a la historia y la mente de sus personajes. Esta perspectiva
asumida por el narrador da ocasión a que el relato tenga dos finales sucesivos.
En el segundo párrafo, siguiendo lo que describe a partir de las señales
externas del comportamiento de los personajes, la historia del cautivo tiene un
primer final feliz ---\forma{Los ojos le brillaron de alegría y los padres
lloraron porque habían encontrado al hijo}---. Pero en el párrafo siguiente, el
narrador decide continuar la historia y propone un segundo final, que se
contrapone al anterior, lo que se explicita mediante el uso de la conjunción
coordinante adversativa inicial, encargada de señalar el cambio
---\forma{\textbf{pero} el indio no podía vivir entre paredes y un día fue a
buscar su desierto}---.

En cuanto a su organización global, el cuento presenta: 1) un \textbf{encuadre
de referencia}, que introduce lo que se cuenta en la región acerca de la
historia del cautivo; 2) el \textbf{núcleo narrativo}, que relata la historia
de la pérdida, recuperación y fuga final del protagonista; y 3) la
\textbf{evaluación}, que contiene los comentarios subjetivos del narrador
acerca de la historia. Ahora bien, un primer aspecto a destacar es que las
distinciones anteriores se correlacionan con cambios en la categoría gramatical
de persona. En el encuadre de referencia, el efecto de indeterminación se
transmite mediante formas propias de la impersonalidad, como el uso de la
3.ª~persona plural con sujeto indeterminado ---\forma{refieren (la
historia)}---, la 3.ª~singular pasiva (pasiva con \forma{se})
---\forma{se dijo}--- y la inclusión de un sujeto no personal ---\forma{la
crónica (ha perdido)}---. En el núcleo narrativo, en el que se relatan los
hechos fundamentales de la historia, se emplean formas de 3.ª~persona singular
o plural, pero referidas a sujetos personales y determinados, correspondientes
a los personajes de la historia ---el cautivo, sus padres, un soldado---. En la
evaluación final, la perspectiva del narrador se introduce mediante la
1.ª~persona, resaltada por la presencia repetida del pronombre \forma{yo}. La
tabla~\ref{tab:01-cautivo} resume lo expuesto:

\begin{table}[htbp]
  \centering
  \footnotesize
  \begin{tabularx}{\linewidth}{@{}>{\raggedright\arraybackslash}p{.24\linewidth}
                                  >{\raggedright\arraybackslash}X
                                  >{\raggedright\arraybackslash}p{.24\linewidth}@{}}
    \toprule
    \textbf{Encuadre de referencia} & \textbf{Núcleo narrativo} & \textbf{Evaluación} \\
    \textsc{lo que cuentan}          & \textsc{la historia del cautivo} & \textsc{comentarios del narrador} \\
    \textit{Uso de formas impersonales} & \textit{Uso de formas de 3.ª pers. sg. y pl. referidas a sujetos personales} & \textit{Uso de la 1.ª pers.} \\
    \midrule
    \textbf{Refieren} la historia (3.ª impersonal) \newline \textbf{Se dijo} (pasiva con \textit{se})
      & \textit{Desaparición}: Un chico desapareció. Lo habían robado los indios. Sus padres lo buscaron inútilmente.
      & \\
    \addlinespace
    \textbf{La crónica} ha perdido las circunstancias (sujeto no personal)
      & \textit{Encuentro/Regreso}: Un soldado les habló de un indio\ldots{} (Los padres) Dieron al fin con él. El hombre se dejó conducir.
      & No quiero inventar lo que no \textbf{sé} \\
    \addlinespace
      & \textit{Reencuentro}: Se detuvo, miró la puerta\ldots{} Gritó, bajó la cabeza, atravesó corriendo, se metió en la cocina. Sin vacilar, hundió el brazo, sacó el cuchillito. Los ojos le brillaron, los padres lloraron.
      & \\
    \addlinespace
      & \textit{Desencuentro/Huida}: El hombre no podía vivir entre paredes. Un día el indio se fue a buscar su desierto.
      & \textbf{Yo} querría saber qué sintió en aquel instante de vértigo\ldots; \textbf{yo} querría saber\ldots \\
    \bottomrule
  \end{tabularx}
  \caption[Correlación de las personas gramaticales]{Correlación de las
  personas gramaticales y las personas gramaticales\protect\pendiente{Epígrafe
  repetido en el original}}\label{tab:01-cautivo}
\end{table}

Uno de los aspectos más destacados del relato es que se presenta atravesado por
la incertidumbre. Ahora bien, ¿cómo apoyan los elementos léxico-gramaticales
esta sensación que transmite el texto? El relato se inicia con una coordinación
disyuntiva, que no tiene el típico valor de \sig{elección} de estas
construcciones, como el que se observa, por ejemplo, en \forma{¿Querés un
helado o una bebida fría? ¿Vamos al cine o al teatro?} La disyunción inicial
\forma{en Junín o en Tapalqué}, que sugiere que los hechos sucedieron «en algún
lugar de La Pampa», genera el efecto de presentar como poco seguro el lugar en
que transcurren los hechos, a la vez que remite intertextualmente al comienzo
del \textit{Quijote} («En un lugar de la Mancha, de cuyo nombre no quiero
acordarme»). Los elementos léxicos también colaboran en señalar falta de
certeza: así, por ejemplo, encontramos la combinación \forma{bien podía ser},
en la que el adverbio equivale a \forma{tal vez}; el uso del verbo
\forma{creer} ---\forma{creyeron reconocerlo}---, que le quita seguridad al
reconocimiento; adverbios de duda como \forma{tal vez} y \forma{acaso}; la
construcción aproximativa \forma{como sin entenderla}. En la parte
correspondiente a la evaluación, también abundan las marcas de incertidumbre:
el uso reiterado del condicional seguido del verbo \forma{saber}, que enfatiza
el desconocimiento por parte del narrador ---\forma{yo querría saber}---; la
interrogativa indirecta que manifiesta lo que ignora pero desea saber
---\forma{qué sintió en aquel instante de vértigo}---; y finalmente las dos
subordinadas condicionales coordinadas ---\forma{si el hijo perdido renació y
murió en aquel éxtasis o si alcanzó a reconocer, siquiera como una criatura o
un perro, los padres y la casa}---, que señalan enfáticamente y con un dejo de
ironía el deseo del narrador no omnisciente de conocer los resortes íntimos de
la conducta de su personaje. La segunda condicional plantea una disyuntiva
tácita: \forma{si alcanzó a reconocer o no}, y una expectativa de mínima:
\forma{siquiera}. En contraste con el clima de incertidumbre que enmarca la
historia, es de destacar que dentro del núcleo narrativo el narrador incluye,
contrastivamente, una afirmación enfática: \forma{sin vacilar}, que es la que
impulsa al lector a aceptar el primer final feliz, que luego será desplazado
por el relato de la huida del indio al desierto.

Por último, también se destaca el papel que en este microrrelato cumplen los
SSNN con que se va sucesivamente nombrando al protagonista, los cuales, poco a
poco, van revelando su verdadera identidad. \forma{El cautivo} es la
denominación inicial del título, que luego no aparece dentro del texto, pero
que alude a la dualidad interna que atraviesa al protagonista: cautivo de
\sig{dos civilizaciones} o \sig{modos de vida}, tal como los presenta el
narrador. Dentro del texto, el protagonista es presentado inicialmente
precedido del determinativo\footnote{Los determinativos son palabras
funcionales que tienen un significado fundamentalmente gramatical, que se
precisa en su relación con otra(s) palabra(s). Los determinativos débiles
---\textit{un, alguno, muchos, varios, otros}, etc.--- sirven para introducir
entidades nuevas en el discurso, mientras que los determinativos fuertes
---\textit{los, mis, estos, todos, ambos}, etc.--- permiten hacer mención a
entidades conocidas previamente \parencite{giammatteo2009}.} indefinido, como
\forma{un chico}, según lo esperado en la primera mención de una entidad nueva
en el discurso. Luego, cuando el soldado cree descubrirlo, el texto vuelve a
utilizar otro sustantivo precedido por un determinativo indefinido ---\forma{un
indio de ojos azules}---, y la reiteración del indefinido hace que el sintagma
funcione como una segunda presentación del protagonista, ahora según «su
nueva» naturaleza adquirida «en el desierto y la vida bárbara». En la tercera
mención, el protagonista ya aparece como corresponde a una entidad que ya fue
mencionada en el discurso, por lo que va precedida del determinante definido
---\forma{el hombre}---. Este sintagma contrasta doblemente por el carácter
definido y por el sustantivo empleado, con la mención inicial ---\forma{un
chico}---, señalando, además, el paso del tiempo y el crecimiento del niño. En
el momento del reconocimiento, el narrador elige nombrarlo como \forma{el
hijo}, para ponerlo en correlación con la anterior referencia a sus padres
---\forma{y los padres lloraron porque habían encontrado al hijo}---. Se
manifiesta así que la serie denominativa \forma{el cautivo / un chico / un
indio / el hombre} culmina en \forma{el hijo} que los padres reencuentran.
Pero, como ya hemos dicho, la historia continúa y, en esta segunda parte, el
narrador emplea otro sintagma definido ---\forma{el indio}---, con el que
revela la verdadera identidad del protagonista, aquello en lo que se ha
convertido a través del cautiverio. Por último, en su planteo dubitativo final,
el narrador se refiere nuevamente al cautivo como \forma{el hijo}, pero le
añade el calificativo \forma{perdido}, que contrasta con la mención en que
aparecía como \forma{el hijo que los padres habían recuperado}. Así, el
narrador culmina la caracterización del personaje y completa la serie
denominativa: \forma{el cautivo / un chico / un indio de ojos azules / el
hombre / el hijo} (recuperado) / \forma{el indio / el hijo perdido}.

El análisis del texto no se ha agotado con el enfoque presentado; no obstante,
la perspectiva léxico-gramatical desarrollada, que parte del conocimiento del
sistema, pero que también se interesa por el uso, no solo hace posible
sistematizar el funcionamiento de los diferentes recursos lingüísticos en el
nivel oracional, sino que también posibilita considerar su proyección textual.

\begin{recuadro}[Reflexión final]
El saber gramatical cumple un papel esencial: hace explícito el conocimiento
tácito del hablante, permitiéndole distanciarse del uso cotidiano para comenzar
a reflexionar sobre la naturaleza del lenguaje. Asimismo, constituye un
conocimiento fundamental para la comprensión de los sutiles matices
involucrados en los múltiples mensajes que damos y recibimos cotidianamente. De
este modo, el reconocimiento de los valores expresivos que conjugan los
elementos léxicos y gramaticales, sin duda, contribuye a captar la intención
del productor y a una interpretación más plena del texto. Este conocimiento
reflexivo es de gran importancia porque permite entender cómo funciona la
lengua y, en consecuencia, operar más eficazmente con ella en el uso cotidiano.
En esta perspectiva, el aprendizaje de la lengua materna puede dejar de ser un
procedimiento estéril y repetitivo para convertirse en un espacio para el
desarrollo de las habilidades cognitivo-lingüísticas necesarias para acompañar
la lectura y estimular la producción de los textos requeridos en el proceso de
enseñanza-aprendizaje.
\end{recuadro}

% --- Secciones finales -------------------------------------------------
\lecturasrecomendadas

\solucionario

\begin{solucion}{\ref{ej:01-facultad}}
\begin{enumerate}[label=\arabic*., start=2, leftmargin=2em, itemsep=.2\baselineskip]
  \item Se trata de una facultad universal porque todos los seres humanos la
        poseen.
  \item Es privativa de nuestra especie porque, al menos hasta donde sabemos,
        ninguna otra especie ha desarrollado un lenguaje con las
        características del humano.
  \item El lenguaje sirve para la comunicación, pero también cumple otros
        fines, como desarrollar el pensamiento abstracto, organizar la actividad
        mental y planificar nuestra actuación futura.
  \item La forma oral del lenguaje es natural para el hombre porque es la que
        todos los humanos sin distinción poseen.
  \item Desde esta perspectiva, la escritura puede considerarse una técnica
        desarrollada para conservar las palabras, la cual debe ser aprendida.
  \item Sin el lenguaje la vida social no sería posible porque es básicamente
        el medio por el que planificamos y organizamos todas nuestras
        actividades comunitarias.
  \item El lenguaje es parte de la dotación biológica de nuestra especie porque
        es parte de nuestra herencia genética y, por lo tanto, estamos
        naturalmente capacitados para desarrollarlo.
  \item Para el desarrollo del lenguaje es necesario interactuar con el medio
        porque necesitamos del estímulo que proviene de escucharlo en el medio
        ambiente (de la misma manera que una semilla necesita la tierra y el
        agua para germinar en una planta).
  \item Somos capaces de producir combinaciones de unidades que constituyen
        nuevas oraciones que jamás hemos oído porque conocemos las unidades y
        sabemos cómo combinarlas (o bien: porque tenemos internalizada la
        gramática de nuestra propia lengua).
\end{enumerate}
\end{solucion}

\begin{solucion}{\ref{ej:01-combinaciones}}
\begin{enumerate}[label=\arabic*., start=2, leftmargin=2em, itemsep=.2\baselineskip]
  \item El V \forma{vivir} exige un sujeto \sig{animado} y no es compatible con
        un sujeto como \forma{la lluvia}.
  \item El V \forma{tener}, además de sujeto, exige un COD que dé cuenta de lo
        que el sujeto posee. Si no se explicita, la O queda incompleta.
  \item La P exige como complemento un SN, no un SA como \forma{triste}.
  \item El V \forma{rezar} exige un COD compatible con su significado
        ---\forma{una plegaria, un padrenuestro}, etc.---.
  \item La P \forma{en} exige un complemento referido a lugar ---\forma{en
        Roma, en la iglesia}--- o bien con el V \forma{rezar} (que es un V de
        \sig{habla}) podría referirse a un idioma ---\forma{en latín, en
        chino}---.
\end{enumerate}
\end{solucion}

\begin{solucion}{\ref{ej:01-diccionariamente}}
La palabra \forma{diccionariamente} que usa el científico Diego Golombek está
formada por el N \forma{diccionario} y el sufijo \forma{-mente}. Resulta
aceptable porque tiene interpretación en el contexto del artículo: se refiere a
\sig{la manera en que aprendemos las palabras en los diccionarios, es decir, en
forma descontextualizada}.
\end{solucion}

\referenciascapitulo

% =====================================================================
%  capitulos/cap02.tex — Capítulo 2. Las reglas de la gramática
%  Autor: Augusto M. Trombetta
%  Fuente: Capitulo_02_normalizado.docx
% =====================================================================
\capitulo{Las reglas de la gramática}{Augusto M. Trombetta}

\section{Introducción}

Las reglas de la gramática no parecen ser todas iguales: unas definen las
características constitutivas propias de los objetos de una lengua, otras
establecen normas de corrección y etiqueta propias del uso social de esa
lengua. La diferencia entre ambos tipos de reglas es correlativa de la
distinción moderna entre gramáticas científicas (descriptivas y explicativas de
los fenómenos lingüísticos observables) y gramáticas tradicionales (normativas
y prescriptivas de los usos lingüísticos correctos). En un caso, las reglas se
descubren, se postulan como conjeturas vinculadas con el funcionamiento interno
de la lengua (por ejemplo, \forma{la oración consta de un sujeto y un
predicado}); en el otro, las reglas se fijan, se establecen como normas
externas relativas a las convenciones sociales de uso de la lengua en un lugar
y época determinados (por ejemplo, \forma{no debe haber coma entre el sujeto y
el predicado}). Este capítulo gira en torno a las reglas de la gramática. En
cuanto a su organización, en primer lugar, se esquematizan las diferentes
concepciones de la gramática y se contrastan los criterios evaluativos de
corrección y gramaticalidad presentes en estas; en segundo lugar, se
especifican las nociones de gramaticalidad y aceptabilidad incorporadas en las
gramáticas científicas, incluyendo las dimensiones normativas como parte de las
condiciones de aceptabilidad; finalmente, se vincula la noción de
gramaticalidad con las variedades lectales que presentan las lenguas.

\section{Concepciones de la gramática}

Como se dijo en el \capref{1}, tradicionalmente, la gramática ha sido
considerada una disciplina sujeta a reglas que rigen el buen uso, es decir, una
disciplina \termino{normativa} o \termino{prescriptiva}. Más modernamente, la
gramática, ya sea considerada en un sentido estático como \textit{la
descripción de un cierto estado de lengua}, ya sea definida en un sentido
dinámico como \textit{el dispositivo mental interno que hace posible articular
las estructuras propias de una lengua}, se entiende como constituida por un
sistema de relaciones, reglas y operaciones, que se deben descubrir para
comprender la conducta lingüística: en este sentido, se trata de una disciplina
\termino{descriptiva} o \termino{científica}.

La valoración de una combinación de la lengua (sílaba, palabra, sintagma,
oración) como \termino{correcta} o \termino{incorrecta} solo adquiere sentido en
el marco de una gramática normativa o prescriptiva, con reglas que controlan
externamente las emisiones lingüísticas. En contraste, una gramática
descriptiva o científica, con reglas que explican internamente el
funcionamiento de la lengua, valora las estructuras de la lengua como
\termino{gramaticales} o \termino{agramaticales}. Esquematizando, se pueden
definir dos maneras complementarias de concebir la gramática y de tipificar las
reglas que forman parte de esta, asociadas a su vez a dos maneras de evaluar
las formas propias del lenguaje humano:

\begin{table}[htbp]
  \centering
  \begin{tabularx}{\linewidth}{@{}>{\itshape}l C C@{}}
    \toprule
    Concepciones de la gramática & Normativa o prescriptiva & Descriptiva o científica \\
    \midrule
    Tipos de reglas          & Regulativas (externas)  & Constitutivas (internas) \\
    Criterios de evaluación  & Corrección/\hspace{0pt}Incorrección & Gramaticalidad/\hspace{0pt}Agramaticalidad \\
    \bottomrule
  \end{tabularx}
  \caption{Concepciones de la gramática}\label{tab:02-concepciones}
\end{table}

\begin{ejercicio}\label{ej:02-definiciones}
Considere las siguientes definiciones 1-5, representativas de las dos
concepciones señaladas sobre la gramática y, posteriormente, complete con los
valores «\mas» o «\menos» (\sig{sí} o \sig{no}) el cuadro presentado más abajo,
como se muestra en el ejemplo resuelto~1.

\begin{exe}
  \ex La \textit{gramática de la actual lengua castellana} no es otra cosa que
      \textit{el conjunto ordenado de las reglas de lenguaje que vemos
      observadas en los escritos o conversación de las personas doctas que
      hablan el castellano o español} \parencite[1]{salva1852}.
  \ex La \textit{gramática} de una lengua es el arte de hablarla correctamente,
      esto es, conforme al buen uso, que es el de la gente educada. Se prefiere
      este uso porque es el más uniforme en las varias provincias y pueblos que
      hablan una misma lengua, y por lo tanto el que hace que más fácil y
      generalmente se entienda lo que se dice; al paso que las palabras y frases
      propias de la gente ignorante varían mucho de unos pueblos y provincias a
      otros, y no son fácilmente entendidas fuera de aquel estrecho recinto en
      que las usa el vulgo \parencite[§§\,1-2]{bello2002}.
  \ex La lingüística estática o descripción de un estado de lengua se puede
      llamar \textit{gramática} en el sentido muy preciso, y por lo demás usual
      [\ldots] donde se trata de un objeto complejo y sistemático que pone en
      juego valores coexistentes. La gramática estudia la lengua como sistema
      de medios de expresión; quien dice gramatical dice sincrónico y
      significativo \parencite[223]{saussure1945}.
  \ex {[La lengua]} Es un tesoro depositado por la práctica del habla en los
      sujetos que pertenecen a una misma comunidad, un sistema gramatical
      virtualmente existente en cada cerebro, o, más exactamente, en los
      cerebros de un conjunto de individuos, pues la lengua no está completa en
      ninguno, no existe perfectamente más que en la masa
      \parencite[57]{saussure1945}.
  \ex La gramática de L será, pues, un ingenio que genere todas las secuencias
      gramaticales de L y ninguna de las agramaticales
      \parencite[27]{chomsky1974a}.
\end{exe}

\begin{tabularx}{\linewidth}{@{}L *{5}{Y{.8cm}}@{}}
  \toprule
  De acuerdo con las definiciones dadas, la gramática de una lengua\ldots & 1 & 2 & 3 & 4 & 5 \\
  \midrule
  \ldots\ remite al uso lingüístico de ciertos grupos sociales.               & \mas   &&&& \\
  \ldots\ representa la descripción de un sistema de valores lingüísticos.    & \menos &&&& \\
  \ldots\ constituye un mecanismo capaz de producir estructuras lingüísticas. & \menos &&&& \\
  \ldots\ implica una regulación modélica, exterior al individuo que la usa.  & \mas   &&&& \\
  \ldots\ supone un sistema o mecanismo internalizado, interior al individuo. & \menos &&&& \\
  \ldots\ conforma una disciplina normativa o prescriptiva.                   & \mas   &&&& \\
  \ldots\ conforma una disciplina descriptiva o científica.                   & \menos &&&& \\
  \bottomrule
\end{tabularx}

\textit{Nota}: La columna correspondiente a 1 muestra el ejemplo resuelto.
\end{ejercicio}

\begin{ejercicio}\label{ej:02-searle}
Considere las siguientes citas de \textcite[55]{searle1964},\footnote{En su
ensayo de 1969, \textit{Actos de habla}, Searle retoma su distinción entre
reglas regulativas y constitutivas, dando a entender que las reglas
gramaticales son regulativas: «Los diferentes lenguajes humanos, en la medida
en que son intertraducibles, pueden considerarse como plasmaciones
convencionales diferentes de las mismas reglas subyacentes. El hecho de que en
francés pueda hacerse una promesa diciendo \enquote*{Je promets} y que en
castellano pueda hacerse diciendo \enquote*{Yo prometo}, es un asunto de
convención. Pero el hecho de que una emisión de un dispositivo de prometer
cuente como (bajo condiciones apropiadas) la asunción de una obligación, es un
asunto de reglas y no un asunto de convenciones del francés o del castellano»
\parencite[48-49]{searle1994}. Es decir, las lenguas naturales aparecen como
convenciones que plasman o materializan las reglas subyacentes que gobiernan
los actos de habla. En el mismo sentido, considerando el dominio estricto de
las estructuras de una lengua, se mantiene la distinción introducida por Searle
entre constitutivo y regulativo: una gramática subyacente, formada por reglas
constitutivas que son la condición de posibilidad de las estructuras
lingüísticas, y una normativa gramatical que regula convencionalmente el uso de
tales estructuras posibles.} donde se establece la diferencia entre
\termino{reglas regulativas} y \termino{reglas constitutivas}. Posteriormente
complete con los valores «\mas» o «\menos» (\sig{sí} o \sig{no}) el cuadro
presentado más abajo, como se muestra en el ejemplo resuelto~(a).

\begin{exe}
  \ex Algunas reglas regulan formas de comportamiento previamente existentes.
      Por ejemplo, las reglas de etiqueta en la mesa regulan el evento de
      comer, pero el evento de comer existe con independencia de estas reglas.
      Existen reglas, por otro lado, que no se limitan a regular, sino que
      crean o definen nuevas formas de comportamiento: las reglas del ajedrez,
      por ejemplo, no se limitan a regular una actividad previamente existente
      llamada jugar al ajedrez; tales reglas, por así decirlo, crean la
      posibilidad o definen esa actividad.
  \ex La distinción que estoy tratando de hacer fue adelantada por la
      distinción de Kant entre principios regulativos y constitutivos, así que
      adoptemos su terminología y describamos nuestra distinción como una
      distinción entre reglas regulativas y constitutivas. Las reglas
      regulativas regulan actividades cuya existencia es independiente de tales
      reglas; las reglas constitutivas constituyen (y también regulan) formas de
      actividad cuya existencia depende lógicamente de estas reglas.
\end{exe}

\begin{tabularx}{\linewidth}{@{}L Y{1.7cm} Y{1.7cm}@{}}
  \toprule
  ¿Qué tipo de regla encierran los siguientes enunciados? & Regulativa & Constitutiva \\
  \midrule
  a) En una oración, todo predicado tiene un sujeto. & \menos & \mas \\
  b) No corresponde decir «pienso de que mañana empiezo la Facultad». & & \\
  c) No se dice «haiga», se dice «haya». & & \\
  d) «Si seguís hablando, la vas a pasar mal» se entiende como amenaza. & & \\
  e) Toda estructura lingüística contiene un núcleo. & & \\
  f) Usar el léxico lunfardo queda prohibido en la radiodifusión argentina. & & \\
  \midrule
  ¿Qué tipo de regla ofrecen las gramáticas normativas o prescriptivas? & & \\
  ¿Qué tipo de regla presentan las gramáticas descriptivas o científicas? & & \\
  ¿Qué tipo de regla es responsable de la regulación interna de una lengua? & & \\
  ¿Qué tipo de regla es responsable de la regulación externa de una lengua? & & \\
  \bottomrule
\end{tabularx}

\textit{Nota}: La fila correspondiente a (a) muestra el ejemplo resuelto.
\end{ejercicio}

\section{Corrección y gramaticalidad}

Según lo presentado hasta acá, las gramáticas normativas o prescriptivas
normalmente se encuentran asociadas a un conjunto de reglas que hacen a los
aspectos regulativos de la lengua, o sea, a las condiciones del buen uso de las
formas lingüísticas. Tales reglas regulativas, en la medida en que representan
regulaciones conscientes de actividades o conductas cuya existencia no depende
de ellas, constituyen factores externos a la estructura de la lengua, que
frecuentemente se justifican de acuerdo con criterios históricos y «que son
producto de las relaciones sociales\ldots» \parencite[29]{bosque2009}. Aunque el
tradicional concepto de \termino{corrección}, aspecto central de las gramáticas
normativas o prescriptivas, fue objeto de reconsideración a lo largo del
siglo~\textsc{xx}, conviene observar que la existencia de «gramáticas
normativas [\ldots] es esperable, e incluso necesaria por razones educativas, y
en general, sociales» \parencite[29]{bosque2009}.

Podemos ilustrar las reglas regulativas con las que figuran en \textit{El buen
uso del español}, una obra normativa elaborada por la Real Academia Española y
la Asociación de Academias de la Lengua \parencite{rae2013}, en la que se
pueden encontrar las siguientes recomendaciones:

\begin{exe}
  \ex\label{ex:02-coma} No se escribe coma entre el sujeto y el predicado,
      independientemente del orden que adopten, ni aunque el sujeto sea
      complejo o una subordinada sustantiva (p.~66).
  \ex\label{ex:02-loismo} Recibe el nombre de loísmo la utilización de las
      formas masculinas de acusativo \forma{lo/los} en lugar de las de dativo
      (\forma{le/les}) [\ldots] Aunque aparece ocasionalmente en el español
      clásico, nunca entró de lleno en la lengua literaria. Hoy se registra
      solamente en algunas zonas de Castilla. El loísmo está fuertemente
      desprestigiado, por lo que se recomienda evitarlo en todos los niveles de
      la lengua (p.~261).
  \ex\label{ex:02-dequeismo} Se llama dequeísmo el uso incorrecto de la
      secuencia \forma{de que} cuando la preposición \forma{de} no está
      justificada [\ldots] El dequeísmo se ha extendido de forma desigual en
      los países hispanohablantes, más en la lengua oral que en la escrita y
      algo más en el español americano que en el europeo, aunque se documenta
      ampliamente en ambos. A pesar de esta difusión, el dequeísmo constituye
      una incorrección gramatical que se recomienda evitar (p.~382).
\end{exe}

Asociado a las gramáticas descriptivas o científicas, el concepto de
\termino{gramaticalidad}, de incorporación más reciente en el campo de la
lingüística, apunta a los aspectos constitutivos de la estructura del lenguaje,
a los factores internos que caracterizan su naturaleza o propiedades. Siguiendo
a \textcite{bosque2009}, las reglas constitutivas, por representar pautas «que
constituyen los fundamentos de algún sistema» (p.~28), ofrecen los «principios
que articulan el sistema lingüístico» (p.~29). Ajeno a la dicotomía normativa
entre formas \termino{correctas} y formas \termino{incorrectas}, el
descubrimiento de tales principios fundamenta tanto la aceptación de las
secuencias bien formadas (\termino{gramaticales}) como la exclusión de las
secuencias mal formadas (\termino{agramaticales}).

Las siguientes observaciones o descripciones extraídas de la
\textit{Introducción a la sintaxis del español}, una obra orientada a la
enseñanza \parencite{camacho2018}, permiten ejemplificar las reglas
constitutivas:

\begin{exe}
  \ex\label{ex:02-haiga} ¿Es lo mismo «gramatical» que «correcto»? No, el
      primero es un concepto cognitivo, el segundo es un concepto social. Por
      ejemplo, en español latinoamericano se oye frecuentemente la forma verbal
      \forma{haiga} (\forma{ojalá que haiga comida en la fiesta}), y algunas
      personas dicen que solo la primera forma es «correcta». Esto refleja una
      jerarquía social: los hablantes que dicen \forma{haya} tienen más
      prestigio social que los que dicen \forma{haiga}. Pero esto no tiene nada
      que ver con la gramaticalidad: la gramática de los hablantes produce
      \forma{haiga} y lo reconoce como posible, por lo tanto, es gramatical
      (p.~17).
  \ex\label{ex:02-faxear} Más arriba hablábamos de los
      morfemas\footnote{El concepto de morfema y los demás términos de la
      morfología se desarrollan en el \capref{4}.} derivativos como morfemas
      que permiten relacionar palabras con una raíz común, pero con
      significados parcialmente distintos, como \forma{mente} y \forma{mental}.
      Este mecanismo también nos sirve para formar palabras nuevas. Por
      ejemplo, a partir de la palabra \forma{fax}, se puede introducir el verbo
      \forma{faxear}, y a partir de \forma{chat}, \forma{chatear} (p.~75).
  \ex\label{ex:02-recursividad} Una de las propiedades más interesantes del
      lenguaje humano es la recursividad,\footnote{Esta noción se presentó en
      el \capref{1}.} que nos permite expresar un número potencialmente
      infinito de pensamientos usando recursos limitados (p.~311).
\end{exe}

\begin{ejercicio}\label{ej:02-correccion}
Considere las siguientes secuencias lingüísticas y evalúe su corrección y
gramaticalidad completando con los valores «\mas» o «\menos» (\sig{sí} o
\sig{no}) el cuadro presentado más abajo, como se muestra en el ejemplo
resuelto~1. Luego responda las preguntas sobre 4a y 4b.

\begin{tabularx}{\linewidth}{@{}L Y{1.5cm} Y{1.5cm}@{}}
  \toprule
  ¿Qué evaluación merecen las siguientes secuencias lingüísticas? & Correcta & Gramatical \\
  \midrule
  1. a. Pierden derechos y carecen por completo de toda protección. & \mas & \mas \\
  \phantom{1.} b. Pierden y carecen por completo. & \menos & \menos \\
  \phantom{1.} c. Pierden derechos y carecen de toda protección. & \mas & \mas \\
  \addlinespace
  2. a. Las diputadas están saliendo y están entrando a otra reunión. & & \\
  \phantom{2.} b. Las diputadas están saliendo de una reunión y están entrando a otra. & & \\
  \addlinespace
  3. a. Trajimos tres remeras iguales para que no haya diferencias de color. & & \\
  \phantom{3.} b. Trajimos tres remeras iguales para que no haiga diferencias de color. & & \\
  \phantom{3.} c. Trajimos tres remeras iguales para que no hay diferencias de color. & & \\
  \addlinespace
  4. a. Tu mujer trabaja todo el día, parece \textit{blecansain}. & & \\
  \phantom{4.} b. Tu \textit{jermu} trabaja todo el día, parece incansable. & & \\
  \bottomrule
\end{tabularx}

\textit{Nota}: El bloque correspondiente a 1 muestra el ejemplo resuelto.

\begin{resuelto}[Aclaraciones sobre el ejemplo resuelto]
1. a. [Pierden [derechos]] y [carecen [por completo] [de toda protección]] /
b. [Pierden Ø] y [carecen [por completo] Ø]. Si bien \forma{pierden} puede
usarse sin COD (\forma{derechos}), si se omite el CR de \forma{carecer}
(\forma{de toda protección}), la oración resulta agramatical. / c. [Pierden
[derechos]] y [carecen Ø [de toda protección]]. La omisión del circunstancial
\forma{por completo}, al no ser un elemento exigido por el verbo, no impide que
la oración siga siendo gramatical. En el contraste de a-b-c se advierte que la
buena formación (gramaticalidad) de las oraciones no depende de la mera
acumulación de palabras: ciertos elementos pueden ser requeridos por la
semántica del verbo (argumentos), frente a otros que resultan facultativos
(adjuntos). Tomado en un sentido amplio, se puede decir que la incorrección de
b depende de su falta de gramaticalidad («es incorrecta porque es
agramatical»), pero no se puede hacer la afirmación opuesta.
\end{resuelto}

\begin{resuelto}[Preguntas sobre 4a y 4b]
¿Es lo mismo \forma{incansable} que \forma{blecansain}? [Esta pregunta la podría
contestar cualquier hablante nativo del español.] ¿Es lo mismo \forma{mujer}
que \forma{jermu}? [Esta pregunta parece apropiada para que la conteste
cualquier hablante nativo del español de la Argentina.]
\end{resuelto}
\end{ejercicio}

\begin{ejercicio}\label{ej:02-registros}
Considere la siguiente cita de \textcite[30]{bosque2009} y conteste las dos
preguntas que se formulan más abajo.

\begin{cita}
Los diferentes modos de comportarse lingüísticamente en una determinada
comunidad se denominan \textsc{registros} o \textsc{estilos}. Los registros se
parecen a los atuendos, los modales, los gestos y otras manifestaciones del
comportamiento social en que son relativos a las situaciones y resultan
apropiados o inapropiados en función de ellas [\ldots] Al igual que sucede con
los atuendos, una parte del conocimiento del idioma afecta a las
\textsc{condiciones de uso} que se asocian con muchas palabras y con algunas
construcciones sintácticas. En muchos casos, los juicios que cabe hacer sobre
las construcciones han de tener en cuenta la variante del sistema gramatical en
la que resultan apropiados.
\end{cita}

\begin{consignas}
  \item ¿Hay algún contraste entre las respuestas esperadas para los puntos 3a
        y 3b del ejercicio~\ref{ej:02-correccion}? [En general, ¿hay algún
        contraste entre las secuencias gramaticales y las secuencias
        correctas?]
  \item Al señalar las condiciones de uso, los registros y los estilos de la
        lengua, ¿los autores están aludiendo al concepto de corrección o al de
        gramaticalidad?
\end{consignas}
\end{ejercicio}

\section{Gramaticalidad y aceptabilidad}\label{sec:02-aceptabilidad}

Los conceptos de \termino{gramaticalidad} y \termino{aceptabilidad} forman
parte del desarrollo de la gramática generativa\footnote{Ver \capref{1}, §\ref{sec:01-generativa} y
§\ref{sec:01-proyeccion}.} y derivan de las obras de Chomsky. Inicialmente, ambas nociones
aparecían muy relacionadas, ya que la idea de \textit{aceptable} venía a
explicar la de \textit{gramaticalidad}, según se puede leer en
\textit{Estructuras sintácticas} \parencite{chomsky1974a}:

\begin{cita}
El propósito fundamental del análisis lingüístico de una lengua L es el de
separar secuencias gramaticales que son oraciones de L, de las secuencias
agramaticales que no son oraciones de L, y estudiar la estructura de las
secuencias gramaticales [\ldots] Una manera de probar (\enquote*{test}) si una
gramática propuesta para L es adecuada, es determinar si las oraciones que
genera son o no son realmente gramaticales, es decir, aceptables para un
hablante nativo (p.~27).
\end{cita}

Son dos conceptos diferentes, como el propio Chomsky se encarga de destacar en
\textit{Aspectos de la teoría de la sintaxis} \parencite{chomsky1970}, donde se
estableció una clara diferencia al relacionarlos con las nociones de
\termino{competencia} (conocimiento lingüístico implícito del hablante de una
lengua) y \termino{actuación} (uso concreto de la lengua por parte de un
hablante). En todo caso, la gramaticalidad es condición necesaria de la
aceptabilidad, pero no es el único factor que interviene en ella:

\begin{cita}
La distinción relacionada entre locuciones \enquote*{aceptables} y
\enquote*{gramaticales} aparece también en SS-LSLT [\textit{Syntactic
structures} y \textit{The logical structure of linguistic theory}], aunque la
terminología es introducida, claramente sólo años después. Podemos tomar
\enquote*{aceptabilidad} como un concepto de la teoría de la actuación, y
\enquote*{gramaticalidad} como una noción relacionada de la teoría de la
competencia. Las oraciones generadas por la gramática directamente son las
oraciones \enquote*{gramaticales} (o \enquote*{gramaticales en el más alto
grado}) de la lengua. Entre ellas hay oraciones que son verdaderas, falsas,
divertidas, ininteligibles, triviales, oraciones que no tienen sentido, etc.
Puede ser muy apropiado, en algunas ocasiones, construir y usar locuciones que
se desvían de las reglas gramaticales en un respecto u otro, y puede muy bien
haber principios generales para interpretar tales oraciones [\ldots] Podemos
decir que la gramática \enquote*{genera derivativamente} las oraciones que se
desvían de las reglas de modos específicos, y que estas oraciones tienen los
grados más bajos de gramaticalidad. Las oraciones son aceptables (o quizá
aceptables bajo circunstancias concretas) si son convenientes, apropiadas,
adecuadas al propósito del momento, etc. La gramática de la competencia
contribuye a determinar la aceptabilidad, pero este concepto implica muchos
otros factores \parencite[6-7]{chomsky1974a}.\pendiente{¿Chomsky 1970
[1965]? El texto atribuye la cita a \textit{Aspectos}}
\end{cita}

Según \textcite[32]{bosque2009}, la \termino{gramaticalidad}, como expresión de
los puros factores gramaticales que «corresponden propiamente al sistema
lingüístico», aparece como una extensión de la teoría de la
\termino{competencia}, el conocimiento tácito (no explícito) del sistema de
reglas constitutivas de su lengua. En tanto, la \termino{aceptabilidad}, que es
la expresión de factores tanto gramaticales como extragramaticales, se presenta
como una extensión de la teoría de la \termino{actuación}, que conjuga factores
de naturaleza diversa (no únicamente gramatical), entre los que se incluyen la
memoria, la consciencia, la rapidez de procesamiento, el cansancio, el estado
de ánimo, el sueño, el nerviosismo, la ingesta de alcohol, etcétera.

\begin{table}[htbp]
  \centering
  \begin{tabularx}{\linewidth}{@{}L L@{}}
    \toprule
    \textbf{Aceptabilidad} & \textbf{Gramaticalidad} \\
    \midrule
    Podemos tomar \enquote*{aceptabilidad} como un concepto de la teoría de la
    \textit{actuación} [\ldots]
      & [\ldots] y \enquote*{gramaticalidad} como una noción relacionada de la
        teoría de la \textit{competencia}. \\
    \addlinespace
    Las oraciones son aceptables [\ldots] si son convenientes, apropiadas,
    adecuadas al propósito del momento, etc. \textit{La gramática de la
    competencia contribuye a determinar la aceptabilidad, pero este concepto
    implica muchos otros factores}.
      & Las oraciones \textit{generadas por la gramática directamente} son las
        oraciones \enquote*{gramaticales} [\ldots] de la lengua. Entre ellas hay
        oraciones que son verdaderas, falsas, divertidas, ininteligibles,
        triviales, oraciones que no tienen sentido, etc. \\
    \bottomrule
  \end{tabularx}
  \caption{Aceptabilidad vs. gramaticalidad}\label{tab:02-aceptabilidad}
\end{table}

\begin{ejercicio}\label{ej:02-aceptabilidad}
Considere las siguientes secuencias lingüísticas y evalúe su gramaticalidad y
aceptabilidad completando el siguiente cuadro con los valores «\mas» o
«\menos» (\sig{sí} o \sig{no}), como se muestra en el ejemplo resuelto~1.

\begin{tabularx}{\linewidth}{@{}L Y{1.5cm} Y{1.5cm}@{}}
  \toprule
  ¿Qué evaluación merecen las siguientes secuencias lingüísticas? & Gramatical & Aceptable \\
  \midrule
  1. a. Estoy seguro de que algunas personas no saben que ciertos periodistas
  afirman que determinados funcionarios no entienden que la razón de su
  nombramiento no viene dada por el currículum que declaran. & \mas & \menos \\
  \phantom{1.} b. Estoy seguro de que algunas personas no lo saben. Ciertos
  periodistas afirman que determinados funcionarios no entienden una cosa
  simple: la razón de su nombramiento no viene dada por el currículum que
  declaran. & \mas & \mas \\
  \addlinespace
  2. La madre del primo de la novia del hermano de mi suegra está contenta. & & \\
  \addlinespace
  3. a. «Toda fuerza armada o reunión de personas que se atribuya los derechos
  del pueblo y peticione a nombre de éste, comete delito de sedición»
  (Constitución Nacional, art.~22). & & \\
  \phantom{3.} b. Toda fuerza armada o reunión de chabones que se atribuya los
  derechos del pueblo y peticione a nombre de éste, comete delito de sedición. & & \\
  \bottomrule
\end{tabularx}

\textit{Nota}: El bloque correspondiente a 1 muestra el ejemplo resuelto.

\begin{resuelto}[Aclaraciones sobre el ejemplo resuelto]
Tanto 1a como 1b son gramaticales. Frente a 1b, la aceptabilidad de 1a se ve
comprometida fuertemente por la cantidad de oraciones subordinadas, incrustadas
unas dentro de las otras, lo que dificulta el procesamiento cognitivo de la
oración final: [Estoy seguro de [que algunas personas no saben [que ciertos
periodistas afirman [que determinados funcionarios no entienden [que la razón
de su nombramiento no viene dada por el currículum [que declaran]]]]]].
\end{resuelto}
\end{ejercicio}

\begin{ejercicio}\label{ej:02-ingrid}
Considere los efectos del alcohol (¿acaso el fernet?) como uno de los factores
no gramaticales que intervienen en la aceptabilidad de las emisiones
lingüísticas. En un video que circuló bastante durante 2015 («El hit del
momento: ¡Dale Ingrid dale!»,
\url{https://www.youtube.com/watch?v=u1AAwMVNxs8}) se puede ver a un joven que
entona una canción cuya letra va cambiando.

\begin{consignas}
  \item ¿Las oraciones de la canción son gramaticales? [Tenga presente la
        estructura de las oraciones: todas las repeticiones que entona.]
  \item ¿Las oraciones de la canción son aceptables? [Tenga presente la
        organización léxica de las oraciones: las primeras dos repeticiones
        frente a las siguientes.]
\end{consignas}
\end{ejercicio}

\section{Gramaticalidad, aceptabilidad y corrección}

El criterio de corrección, o sea la valoración entre \termino{formas correctas}
y \termino{formas incorrectas}, propio de las gramáticas tradicionales de tipo
normativo o prescriptivo vinculadas de manera típica a la enseñanza, se puede
entender como uno de los componentes incluidos en la noción de aceptabilidad (la
dicotomía entre \termino{estructuras aceptables} y \termino{estructuras no
aceptables}) que, según se señaló, es un concepto vinculado con la teoría de la
actuación dentro de un enfoque formal de la gramática. En este sentido,
volviendo sobre la cita de Chomsky (§\ref{sec:02-aceptabilidad}), se puede advertir que la
gramaticalidad (la \textit{buena formación} oracional, es decir, la forma que
adquieren las estructuras lingüísticas según las pautas constitutivas internas
de la gramática) es la condición necesaria para que se materialicen las
estructuras oracionales, con independencia de factores tales como su valor de
verdad y sentido, su elegancia y corrección, su farragosidad o su elocuencia,
etcétera. Queda planteada, en consecuencia, una distinción fundamental entre la
noción de gramaticalidad, asociada con la buena formación de las construcciones
lingüísticas, y la de aceptabilidad, dentro de la que se incluye el criterio de
corrección, y vinculada a los factores sociales e individuales que intervienen
en la actuación lingüística.

\begin{table}[htbp]
  \centering
  \begin{tabularx}{\linewidth}{@{}L L@{}}
    \toprule
    \textbf{Competencia} & \textbf{Actuación} \\
    \midrule
    Reglas y principios constitutivos de la gramática de una lengua natural
      & Factores regulativos y limitantes (adecuación, estilo, corrección,
        memoria, etc.) \\
    \addlinespace
    Gramaticalidad (condición de buena formación de las construcciones de una
    lengua)
      & Aceptabilidad (depende de la situación social, del procesamiento
        mental, etc.) \\
    \bottomrule
  \end{tabularx}
  \caption{Competencia vs. actuación}\label{tab:02-competencia}
\end{table}

Finalmente, la teoría de la competencia (el conocimiento interno del conjunto de
principios y reglas de la gramática de una lengua natural) y la teoría de la
actuación (la puesta en uso por parte del individuo de su competencia
lingüística), con sus conceptos asociados de gramaticalidad y aceptabilidad, no
deja de remitir a la distinción saussureana entre la lengua (el sistema de
valores opositivos internalizado por los miembros de una comunidad lingüística)
y el habla (el uso del sistema de la lengua por parte de un individuo). La
diferencia entre ambas perspectivas pasa por el carácter dinámico que tiene la
gramática interna en Chomsky, que es considerada un dispositivo generativo,
frente al carácter estático y descriptivo que tiene la lengua en Saussure, que
es planteada como un conjunto de relaciones y oposiciones entre los elementos
que la integran.

En lo que sigue se va a contrastar la consistencia formal del concepto de
gramaticalidad con distintas dimensiones que pueden incidir en la aceptabilidad
de una secuencia bien formada. Para ello se van a revisar algunas cuestiones
interpretativas, como la claridad, la verosimilitud o la significatividad, como
así también nociones vinculadas con el uso social del lenguaje, como la
elegancia.

\subsection{Ambigüedad}\label{sec:02-ambiguedad}

El concepto de ambigüedad refiere a secuencias bien formadas que generan más de
una interpretación, como, por ejemplo, \forma{Mi libro}, que puede hacer
referencia a \sig{el libro que poseo} o a \sig{el que escribí}. En las
gramáticas normativas tradicionales se prescribía evitar las expresiones
ambiguas cuando son anfibológicas, es decir, cuando se trata de formas con
múltiples interpretaciones que no resultan claras. Según
\textcite[35]{bosque2009}, que abordan el concepto desde una perspectiva formal,
una secuencia posee ambigüedad gramatical o estructural «cuando puede
corresponderle más de una estructura sintáctica, y por tanto más de una
interpretación». Desde esta perspectiva, el concepto de ambigüedad es
constitutivo de la gramática y apunta a que algunas oraciones por lo menos
admiten «dos análisis sintácticos en el sistema gramatical del español, lo que
da lugar a otros tantos significados». Por ejemplo, una oración como
\exref{ex:02-secarropa} admite dos lecturas, según la manera en que se
interpreten los vínculos entre los elementos que la constituyen:\footnote{La
noción de constituyente se introduce en el \capref{5}.}

\begin{exe}
  \ex\label{ex:02-secarropa} Por favor, cuando termine de usar el secarropa,
      limpie el filtro de pelusas.
\end{exe}

En \exref{ex:02-secarropa}, se puede considerar que [\forma{el filtro}] y
[\forma{de pelusas}] son dos modificadores independientes del verbo
\exref{ex:02-lectura1-a} y que, por lo tanto, sus posiciones son
intercambiables, o bien que [\forma{de pelusas}] modifica a [\forma{el filtro}]
\exref{ex:02-lectura2-a}, por lo que constituyen una unidad. Cada una de estas
lecturas implica un significado diferente de la oración
\exref{ex:02-lectura1-b} y \exref{ex:02-lectura2-b}:

\begin{exe}
  \ex\label{ex:02-lectura1}
    \begin{xlist}
      \ex\label{ex:02-lectura1-a} Por favor, cuando termine de usar el
          secarropa, limpie [el filtro] [de pelusas].
      \ex\label{ex:02-lectura1-b} \textit{Por favor, cuando termine de usar el
          secarropa, limpie [de pelusas] [el filtro].}
    \end{xlist}
  \ex\label{ex:02-lectura2}
    \begin{xlist}
      \ex\label{ex:02-lectura2-a} Por favor, cuando termine de usar el
          secarropa, limpie [el filtro de pelusas].
      \ex\label{ex:02-lectura2-b} \textit{Por favor, cuando termine de usar el
          secarropa, limpie [el filtro [de pelusas]].}
    \end{xlist}
\end{exe}

La ambigüedad gramatical, también conocida como homonimia estructural, puede
darse tanto en el nivel sintáctico como morfológico.\footnote{La ambigüedad u
homonimia estructural en el nivel sintáctico se trata en el \capref{5}, §\ref{sec:05-homonimia} y,
en el nivel morfológico, se estudia en el \capref{4}, §\ref{sec:04-homonimia}.} Pero además de
gramatical, la ambigüedad puede ser léxica, es decir, que puede estar
relacionada con el doble significado que tienen muchas palabras, como se
muestra en el siguiente meme en el que el humor surge a partir de las dos
interpretaciones que la palabra \forma{privado} admite en el mismo contexto:

\begin{exe}
  \ex\label{ex:02-privado} El barrio era privado: privado de carne, de leche, de
      pan\ldots, privado de todo (Diego Maradona, entrevista con J.~Guinzburg,
      \url{https://www.facebook.com/watch/?v=133231901641222}).
    \begin{xlist}
      \ex \forma{privado}\textsubscript{1} adj. que no es propiedad pública o
          estatal, sino que pertenece a particulares.
      \ex \forma{privado}\textsubscript{2} participio pasado del verbo
          \forma{privar}: despojado de (\textit{DLE}, en línea).
    \end{xlist}
\end{exe}

\begin{ejercicio}\label{ej:02-ambiguedad}
Considere las siguientes oraciones y conteste las preguntas que se formulan
abajo, como se muestra en el ejemplo resuelto~1.

\begin{exe}
  \ex Me trajeron esta escultura de Francia.
  \ex Antes de salir a navegar, acomodaron correctamente las velas.
  \ex Yo no recibiría a tu hermano en mi oficina.
\end{exe}

\begin{consignas}
  \item Las tres oraciones anteriores pueden tener más de una interpretación.
        Realice paráfrasis que pongan de relieve los diferentes significados y
        explique de qué tipo de ambigüedad se trata en cada caso.
  \item ¿Observa alguna mala formación de estructura gramatical? ¿Se podría
        resolver de alguna manera el problema de ambigüedad observado en (a)?
\end{consignas}

\begin{resuelto}
1. (a) La oración está bien formada y es gramatical, aunque resulte ambigua, ya
que, fuera de contexto, admite múltiples significados. En este caso, la
ambigüedad se relaciona con factores sintácticos: a. [Me trajeron [esta
escultura [de porcelana [de Francia]]]] / a'. Me trajeron esta escultura
(hecha) de porcelana de Francia. / b. [Me trajeron [esta escultura [de
porcelana]] [de Francia]] / b'. Me trajeron desde Francia esta escultura de
porcelana. (b) La ambigüedad se podría resolver reestructurando el orden de los
constituyentes, como en b'.\pendiente{El ítem 1 no dice «de porcelana»}
\end{resuelto}
\end{ejercicio}

\subsection{Elegancia en la expresión: cortesía/descortesía, redundancia}

La regla de elegancia «no constituye un principio gramatical», ya que al
infringirla no se obtiene una secuencia agramatical, sino «una construcción poco
apropiada o poco elegante» \parencite[36]{bosque2009}. En este sentido, el mayor
o menor grado de elegancia presente en una construcción «se distingue de su
gramaticalidad» y se puede decir que una regla de elegancia constituye «más
bien una norma de cortesía», en tanto se trata de una forma de expresarse que
manifiesta \textit{atención, respeto o afecto} hacia otro (\textit{DLE}, en
línea). Lo contrario de la cortesía, la descortesía, se encuentra cuando alguien
se manifiesta mediante expresiones descalificantes, como las que son frecuentes
en redes sociales:

\begin{exe}
  \ex\label{ex:02-carroza} Ahí lo tenés al pelotudo (meme de amplio uso en redes
      sociales, creado a partir de una escena del filme \textit{Esperando la
      carroza}, de A.~Doria, 1985).
\end{exe}

Asimismo, también se producen expresiones poco elegantes cuando se incurre en
una redundancia o repetición innecesaria de palabras o construcciones. Sin
embargo, la normativa sigue una doble vara:

\begin{cita}
Si nos fijamos ahora en la aproximación retórica y gramatical a la redundancia,
observaremos una extraña dicotomía: las caracterizaciones de la redundancia que
suelen hacerse desde la filología dan a entender que se trata de una opción
legítima si procede de la pluma de los buenos escritores, a la vez que
ilegítima en los demás casos. La definición de \textit{pleonasmo} en el
\textit{DLE} ilustra esta dicotomía con claridad: «Empleo en la oración de uno o
más vocablos innecesarios para que tenga sentido completo, pero con los cuales
se añade expresividad a lo dicho, como en \textit{en fuga irrevocable huye la
hora}» \parencite[30]{bosque2020}.
\end{cita}

Según lo planteado, entonces, las repeticiones deben evitarse cuando reiteran la
información ya expresada en la misma oración \exref{ex:02-nuevamente}, pero son
aceptadas cuando refuerzan una idea o señalan diferencias entre cosas
semejantes \exref{ex:02-puertas}:

\begin{exe}
  \ex\label{ex:02-nuevamente} Así, si el acuífero cumple con las condiciones
      adecuadas, el agua puede \forma{volverse nuevamente} pura («¿Qué hacer
      ante la contaminación de aguas subterráneas?», Sondagua,
      \url{https://www.sondagua.cl/blog/que-hacer-ante-la-contaminacion-de-aguas-subterraneas/}).
  \ex\label{ex:02-puertas} Pasamos \forma{puertas y puertas y cuartos}, hasta que
      llegamos a una gran pieza oscura (A.~Uslar Pietri, «Maichak, el hombre del
      río», p.~170).
\end{exe}

En el caso de \exref{ex:02-matamos}, una estructura condicional tautológica
(\textit{NGLE}, §\,47.5e) pone de relieve la escasa importancia otorgada a lo que
le puede suceder (como si se dijera: \forma{si nos matamos, no pasa nada}) al
mismo tiempo que manifiesta que se trata de una situación irremediable.

\begin{exe}
  \ex\label{ex:02-matamos} Si nos matamos, nos matamos.
  \ex\label{ex:02-nopasa} \textit{No importa, no pasa nada}.
\end{exe}

En los ejemplos \exrefs{ex:02-carroza}{ex:02-nopasa} no se advierten problemas
respecto de la constitución interna de las oraciones: ninguna de las
consideraciones previas afecta la buena formación oracional.

\begin{ejercicio}\label{ej:02-elegancia}
Considere las siguientes oraciones y conteste las dos preguntas que se formulan
más abajo, como se muestra en el ejemplo resuelto~1.

\begin{exe}
  \ex No lo voy a repetir de nuevo.
  \ex Yo y mi hermano fuimos al cine.
  \ex Y si se comen unas patatas, se comen, y si no, no se comen (C.~Martín
      Gaite, \textit{El balneario}).
\end{exe}

\begin{consignas}
  \item ¿En qué se parecen y en qué contrastan estas tres oraciones? [La idea es
        considerar los elementos presentes o su disposición en la estructura.]
  \item ¿Qué problema plantea la estructura gramatical? ¿Se podría resolver de
        alguna manera lo observado en (a)?
\end{consignas}

\begin{resuelto}
1. (a) La oración está bien formada, aunque pueda señalarse como redundante (el
verbo ya encierra el valor de \sig{nuevamente}). (b) La oración no presenta
problemas gramaticales. La falta de elegancia se podría resolver ofreciendo
cambios léxicos o estructurales (\forma{No lo voy a repetir, No lo voy a volver
a decir, No lo voy a decir de nuevo}).
\end{resuelto}
\end{ejercicio}

\subsection{Gramaticalidad y significado referencial}

La buena formación de las expresiones de las lenguas naturales no depende de su
mayor o menor grado de relación con el referente, es decir, con aquella entidad
de la realidad extralingüística a la que remite:

\begin{cita}
Las secuencias que construimos combinando palabras pueden ser gramaticales y a
la vez relativamente incomprensibles, si su significación no se ajusta a lo que
independientemente sabemos sobre el mundo que nos rodea. Lo gramatical no es lo
\textsc{cierto}, ni siquiera lo \textsc{verosímil} [\ldots] Las secuencias que
se ajustan a los principios fundamentales de la gramática no tienen que denotar
situaciones ciertas o verosímiles \parencite[36]{bosque2009}.
\end{cita}

El contraste entre la gramaticalidad y la inconsistencia semántica fue
subrayado por Chomsky, al advertir que una frase como \forma{Las verdes ideas
incoloras duermen furiosamente} constituye, al mismo tiempo, una oración bien
formada desde el punto de vista sintáctico y un sinsentido desde la perspectiva
del significado. Sin embargo, lejos de representar una anomalía, esta clase de
contraste remite al \termino{desplazamiento}, una de las propiedades
constitutivas del lenguaje humano \parencite{hockett1971}, según la cual las
expresiones lingüísticas permiten hacer referencia tanto a lo que se encuentra
en el entorno inmediato (\forma{Hay un olor desagradable en la oficina}) como a
realidades distantes (\forma{Napoleón murió en Santa Elena en 1821}) o
inventadas (\forma{Superman nació en el planeta Krypton}). Así, el amplio
espectro de la ficcionalización hace interpretables oraciones que construyen no
solo sentidos paradójicos sino también mundos irreales:

\begin{cita}
Cuidamos las mancuspias hasta bastante tarde, ahora con el calor del verano se
llenan de caprichos y versatilidades, las más atrasadas reclaman alimentación
especial y les llevamos avena malteada en grandes fuentes de loza; las mayores
están mudando el pelaje del lomo, de manera que es preciso ponerlas aparte,
atarles una manta de abrigo y cuidar que no se junten de noche con las
mancuspias que duermen en jaulas y reciben alimento cada ocho horas
(J.~Cortázar, «Cefalea», p.~43).
\end{cita}

La contradicción (más en general, la oposición de sentidos, como en \forma{Se
produjo un silencio ensordecedor}), la autoconsistencia (es decir, la identidad
analítica de sentidos, como en \forma{Los triángulos tienen tres ángulos}), la
tautología (es decir, la unicidad en la asignación del valor de verdad, como en
\forma{Un peso es un peso}) y otros conceptos semejantes no entran en conflicto
con la noción de gramaticalidad, más bien constituyen «anomalías» desde el punto
de vista interpretativo. Por eso, la inserción de tales oraciones en contextos
situacionales muy específicos o en contextos lingüísticos ampliatorios puede
mejorar su significatividad (inferencia de implícitos, contrastes semánticos).
Así, la contradicción de \exref{ex:02-solteros-a} entre el significado de
\forma{soltero} (que no está casado) y \forma{casados} (que han contraído
matrimonio) se resuelve mediante el añadido de \forma{ahora}
\exref{ex:02-solteros-b}, haciendo referencia a dos momentos en el tiempo
(\forma{los que estaban solteros, ahora ya se han casado}). En
\exref{ex:02-catulo}, en el fragmento inicial del poema 85 de Catulo, la
contraposición entre los antónimos verbales \forma{odio} y \forma{amo} no
constituye en absoluto un sinsentido, sino que, por el contrario, refuerza lo
contradictorio del sentimiento amoroso.

\begin{exe}
  \ex\label{ex:02-solteros}
    \begin{xlist}
      \ex\label{ex:02-solteros-a} Todos los solteros de mi pueblo están casados.
      \ex\label{ex:02-solteros-b} Todos los solteros de mi pueblo están casados
          ahora.
    \end{xlist}
  \ex\label{ex:02-catulo} Odio y amo. ¿Por qué hago esto?, quizá preguntas. / No
      lo sé, pero siento que es así y me torturo.
\end{exe}

\begin{ejercicio}\label{ej:02-referencia}
Considere las siguientes oraciones y conteste las dos preguntas que se formulan
más abajo.

\begin{exe}
  \ex Una selenita nos atacó en la otra cuadra.
  \ex Asimismo agradecemos a Ireneo Fernando Cruz el habernos iniciado, durante
      un viaje a San Juan, en el conocimiento de las mancuspias (J.~Cortázar,
      «Cefalea», p.~43).
\end{exe}

\begin{consignas}
  \item ¿En qué se parecen y en qué contrastan estas dos oraciones?
  \item ¿Qué problema plantea la estructura gramatical? ¿Se podría resolver de
        alguna manera lo observado en (a)?
\end{consignas}
\end{ejercicio}

\begin{ejercicio}\label{ej:02-contradiccion}
Considere las siguientes oraciones y conteste las dos preguntas que se formulan
abajo, como se muestra en el ejemplo resuelto~1.

\begin{exe}
  \ex Los huérfanos de esta ciudad tienen unos padres maravillosos.
  \ex El ciego de la otra cuadra fue el primero que vio el accidente.
  \ex {[Un cliente llamando al mozo]} Este café caliente está frío.
  \ex El nuevo entrenador es simpático y antipático.
\end{exe}

\begin{consignas}
  \item ¿En qué se parecen y en qué contrastan estas oraciones? [Considere la
        manera en que se dan las relaciones de significados entre las palabras
        de las oraciones.]
  \item ¿Qué problema plantea la estructura gramatical? ¿Se podría resolver de
        alguna manera lo observado en (a)?
\end{consignas}

\begin{resuelto}
1. (a) La oración presenta estructuras bien formadas. Sin embargo, se establece
una contradicción entre lo que expresa la oración y el significado de una de las
palabras que contiene: \forma{huérfano} (persona menor de edad a quien se le han
muerto el padre y la madre o uno de los dos) es incompatible con el hecho de
\forma{tener unos padres maravillosos}. (b) No hay problemas de estructura
gramatical, aunque la contradicción es irresoluble porque entran en conflicto
los significados de las palabras: \forma{padres} y \forma{huérfanos}.
\end{resuelto}
\end{ejercicio}

\section{Gramaticalidad y variedades lingüísticas}

Finalmente, conviene observar que la gramaticalidad de las expresiones
lingüísticas, su buena formación, está referida a las distintas variedades de
una lengua, no se trata de una propiedad constitutiva abstracta. Si la gramática
y el léxico de una lengua tienen naturaleza mental (como se acepta desde
Saussure) e individual (como ha planteado Chomsky), es esperable que existan
diferencias explicables, tal vez, por el cambio en la disposición de las reglas
internas de la gramática (la modificación interna que se da en el dispositivo
gramatical). Bajo esta perspectiva, una lengua natural (el español, por
ejemplo) se debe entender como un conjunto de variedades lingüísticas o
lectos\footnote{El concepto de lecto se presentó en el \capref{1}.}
(cronolectos, dialectos, sociolectos e idiolectos, por usar denominaciones
clásicas), que comparten un altísimo número de principios y reglas
constitutivas internas y un número no menos importante de unidades léxicas
básicas. La gramaticalidad de las expresiones de una lengua, ante todo, tiene
que ver con sus distintas variedades: así, por ejemplo, ciertas variedades del
español aceptan la presencia del determinante delante del nombre propio
---\forma{la María, el Pedro}---, mientras que otras lo rechazan.

En términos generales, se pueden considerar cuatro tipos básicos de variaciones
lingüísticas: diacrónica, diatópica, diastrática y diafásica. Los conceptos
presentes en esta terminología, el primero proveniente de \textcite{saussure1945},
los tres restantes tomados de \textcite{coseriu1969}, remiten a las dimensiones
temporal, espacial, social y situacional en que se desarrollan los hechos del
lenguaje.

\subsection{Variación diacrónica}

Hace referencia a la dimensión temporal que origina las diferencias en el habla
entre las distintas generaciones, que se conocen como \termino{cronolectos}. Las
lenguas cambian a lo largo del tiempo, como muestran las dos versiones
siguientes del \textit{Auto de los Reyes Magos} (vv.~1-4 y 29-30, anónimo del
siglo~\textsc{xii}):

\begin{exe}
  \ex\label{ex:02-reyes}
    \begin{xlist}
      \ex\label{ex:02-reyes-a} ¡Dios criador, cuál maravila, / no sé cuál es
          aquesta strela! / Agora primas la e veída; / poco tiempo ha que es
          nacida. [\ldots] ¿En todo, en todo es nacido? / Non se si algo e veido
          (versión original).
      \ex\label{ex:02-reyes-b} ¡Dios Creador! ¡Qué maravilla! / ¿Qué estrella
          será esa que brilla? / Hasta ahora, no la he advertido; / hace bien
          poco que ha nacido. [\ldots] ¿Será verdad que ya ha nacido? / Dudo de
          lo que he advertido (versión de F.~Lázaro Carreter,
          1958).\pendiente{El original dice «L. Carreter»}
    \end{xlist}
\end{exe}

Además de las evidentes diferencias léxicas entre las dos versiones, ambas se
distinguen porque en \exref{ex:02-reyes-a}, que ejemplifica el castellano
medieval, los verbos con forma pasiva y significado activo, conocidos como
verbos deponentes en las gramáticas latinas (verbos inacusativos), tomaban como
auxiliar de los tiempos compuestos el verbo \forma{ser} (\forma{es nacida, es
nacido}), mientras que los demás verbos tomaban el auxiliar \forma{haber}
(\forma{e veída, e veido}). En el castellano moderno \exref{ex:02-reyes-b} esta
diferencia de auxiliares se perdió en favor del auxiliar \forma{haber}
(\forma{ha nacido, he advertido}).

\subsection{Variación diatópica}

Se trata de la variación que se produce en la dimensión geográfica entre
distintas zonas, regiones o lugares, y produce la diferencia entre
\termino{dialectos}, como las que se observan entre los ejemplos de
\exref{ex:02-queres} y \exref{ex:02-quetu}:

\begin{exe}
  \ex\label{ex:02-queres} ¿Qué querés vos para el 2021?
      (\url{https://www.instagram.com/reel/CFiNsjKltdb/}, del usuario comercial
      argentino @dalepaprika, de la red social Instagram).
  \ex\label{ex:02-quetu} ¿Qué tú quieres, mamita? ¿Qué tú quieres? («Qué tú
      quieres», del álbum \textit{Visualízate} de Gente de Zona, grupo cubano de
      salsa y reguetón).
\end{exe}

Mientras en la mayoría de las variedades del español el sujeto se pospone al
verbo en las oraciones interrogativas \exref{ex:02-queres}, en el español de
Cuba, el sujeto de las interrogativas mantiene la posición preverbal
\exref{ex:02-quetu}. Los ejemplos también muestran la diferencia entre el uso
del pronombre de segunda persona típico del español de Buenos Aires,
\forma{vos}, y la forma \forma{tú}, característica del español peninsular y de
la mayoría de las regiones hispanohablantes, como Cuba. Cada una de ellas
implica también diferencias en las formas verbales correspondientes:
\forma{vos querés / tú quieres}.

\subsection{Variación diastrática}

Constituye la dimensión referida a la estratificación o diferenciación social
referida a que las lenguas presentan variaciones relacionadas con los grupos o
estratos sociales, los \termino{sociolectos}, según la ocupación, profesión,
nivel de instrucción, entre otras variables sociales, tal como muestran las dos
interpretaciones siguientes del tango \textit{Cambalache} de Enrique Santos
Discépolo (1934):

\begin{exe}
  \ex\label{ex:02-sosa} ¡Que allá en el Horno / se vamo' a encontrar\ldots!
      (interpretación de J.~Sosa).
  \ex\label{ex:02-gardel} ¡Que allá en el Horno / nos vamo' a encontrar\ldots!
      (interpretación de C.~Gardel).
\end{exe}

En \exref{ex:02-sosa} y \exref{ex:02-gardel} se advierten diferencias en el
empleo del pronombre clítico\footnote{Los clíticos son pronombres personales
átonos que se apoyan para su pronunciación en la palabra que los precede o
sigue. Sus características y funciones se tratan en el \capref{10}.} de primera
persona plural. La forma \forma{se}, que aparece en \exref{ex:02-sosa}, estuvo
asociada durante décadas con los usos de los sectores populares urbanos,
mientras que \forma{nos}, la forma que aparece en \exref{ex:02-gardel}, siempre
se ha considerado como propia de la población escolarizada.

\subsection{Variación diafásica}\label{sec:02-diafasica}

Refiere a la dimensión vinculada con la situación social en que se usa la
lengua. Las formas lingüísticas se van adaptando a los distintos contextos
situacionales en que se produce el acto de habla, a las características del
entorno y del interlocutor, al tipo de discurso empleado y al grado de
formalidad de la comunicación. Tales aspectos constituyen el
\termino{registro}. Así, en \exref{ex:02-solicitud} y \exref{ex:02-guita} se
muestran dos situaciones diferentes en las que la misma persona, un alumno de la
carrera de Ingeniería Eléctrica de la Universidad Nacional del Sur, podría
formular un pedido:

\begin{exe}
  \ex\label{ex:02-solicitud} Por la presente, y en mi carácter de alumno de la
      carrera de Ingeniería Eléctrica, plan 2010, solicito una excepción que me
      permita cursar la asignatura que se detalla a continuación, sin cumplir
      con la totalidad de requisitos exigidos por el plan de estudios
      correspondiente (Solicitud de excepción para cursar o rendir materias,
      según el modelo de la UNS).
  \ex\label{ex:02-guita} Che, tengo que pedirte un favor muy grande. Este mes se
      me está haciendo muy largo y me estoy quedando corto de guita. ¿No me
      aguantás unos mangos hasta que cobre? En una semana te devuelvo todo.
\end{exe}

\begin{ejercicio}\label{ej:02-torito}
Lea atentamente el siguiente fragmento adaptado de \textit{Torito} de
J.~Cortázar y luego resuelva lo pedido más abajo, como se muestra en el ejemplo
resuelto.

\begin{cita}
Qué le vas a hacer, \subrayado{ñato}, cuando estás abajo todos te
\subrayado{fajan}. Todos, che, hasta el más maula. Te sacuden contra las sogas,
te encajan la \subrayado{biaba}. Vos te crees que yo me desespero, lo que pasa
es que no doy más aquí tumbado todo el día. Las noches de invierno son más
largas que esperanza'e pobre. Estoy todo el tiempo mirando el techo. Ahí tenés
otra cosa que no sé hacer, mirar p'arriba. Todos dijeron que me hubiera
convenido, que hice la gran \subrayado{macana} de levantarme a los dos
segundos, \subrayado{cabrero} como la gran flauta. Tienen razón, si me quedo
hasta las ocho no me agarra tan mal el rubio.
\end{cita}

\begin{consignas}
  \item Busque en el \textit{Diccionario de la lengua española} (\textit{DLE})
        de la Real Academia el significado de las palabras subrayadas y
        transcriba la acepción correspondiente a su uso en el texto.
  \item Una de las palabras subrayadas no está en este diccionario. ¿Por qué
        cree que sucede y en qué diccionario podría encontrarla?
  \item ¿Qué tipo de variación ejemplifican estas palabras?
\end{consignas}

\begin{resuelto}
\forma{ñato}: Am. chato (‖ de nariz poco prominente).
\end{resuelto}
\end{ejercicio}

\begin{ejercicio}\label{ej:02-variacion}
Considere las siguientes oraciones y resuelva lo pedido más abajo, como se
muestra en el ejemplo resuelto~1.

\begin{exe}
  \ex Hay que estudiar, \subrayado{chaval}, que la vida está muy mala
      (I.~Aldecoa).
  \ex ¿Es \subrayado{guapo}? No me he fijado en él.
  \ex Palermo, en estos días, era un \subrayado{rancherío} de muy mala muerte
      (M.~Caparrós, \textit{Página 12}, 1993).
  \ex Esto que \subrayado{agora} emprenderéis ---agregó el
      \subrayado{lectoral}--- será en servicio de la santa Iglesia de Cristo
      (E.~Larreta, \textit{La gloria de Don Ramiro}).
  \ex Ante el alza de los alquileres, algunos sectores del gobierno sugieren
      como alternativa el sistema de \subrayado{\textit{leasing}}.
  \ex Es una de las regiones más calurosas y \subrayado{paupérrimas} de todo el
      territorio nepalés.
\end{exe}

\begin{consignas}
  \item Busque en el diccionario el significado de las palabras subrayadas y
        transcriba la acepción correspondiente a su uso en las oraciones.
  \item Indique qué tipo de variación ejemplifica cada una de las palabras
        subrayadas.
\end{consignas}

\begin{resuelto}
1. (a) \forma{Chaval}: m. y f. coloq. Esp. Niño o joven. (b) Variación
dialectal.
\end{resuelto}
\end{ejercicio}

\begin{recuadro}[Reflexión final]
Aunque pueda resultar obvio, debería quedar claro que las tipologías anteriores
referidas a las diferentes variedades lingüísticas coexisten en cada individuo
que articula una oración en su variedad de lengua, una variedad que es a la vez
temporal (hablantes de diferentes generaciones), geográfica (hablantes de
diferentes regiones) y social (hablante con diferentes grados de instrucción o
con distintas ocupaciones), y que se produce en un determinado contexto
situacional (diferentes grados de formalidad y cuidado en el habla). Sin
embargo, muchas veces se pone el eje explicativo del lenguaje en la variedad por
la variedad misma, como si una sola dimensión determinara la conducta
lingüística de los hablantes y oyentes. Alternativamente, son los conocimientos
tácitos e interiores al individuo, su competencia lingüística (gramatical) y su
competencia comunicativa (situacional), los elementos que hacen posible una
conducta lingüística concreta por parte del hablante/oyente. En este sentido,
la competencia del individuo, protagonista central de la lengua, puede parecer
una explicación general y abstracta, pero resulta más esclarecedora que la mera
clasificación de variedades y lectos, que constituyen generalizaciones
exteriores y ajenas a los individuos.
\end{recuadro}

% --- Secciones finales -------------------------------------------------
\lecturasrecomendadas

\solucionario

\begin{solucion}{\ref{ej:02-definiciones}}
\begin{tabularx}{\linewidth}{@{}L *{4}{Y{.8cm}}@{}}
  \toprule
  De acuerdo con las definiciones dadas, la gramática de una lengua\ldots & 2 & 3 & 4 & 5 \\
  \midrule
  \ldots\ representa la descripción de un sistema de valores lingüísticos.    & \textbf{\menos} & \textbf{\mas} & \textbf{\mas} & \textbf{\menos} \\
  \ldots\ constituye un mecanismo capaz de producir estructuras lingüísticas. & \textbf{\menos} & \textbf{\menos} & \textbf{\menos} & \textbf{\mas} \\
  \ldots\ implica una regulación, exterior al individuo que la usa.           & \textbf{\mas} & \textbf{\menos} & \textbf{\menos} & \textbf{\menos} \\
  \ldots\ supone un sistema o mecanismo interior al individuo.                & \textbf{\menos} & \textbf{\mas} & \textbf{\mas} & \textbf{\mas} \\
  \ldots\ conforma una disciplina normativa o prescriptiva.                   & \textbf{\mas} & \textbf{\menos} & \textbf{\menos} & \textbf{\menos} \\
  \ldots\ conforma una disciplina descriptiva o científica.                   & \textbf{\menos} & \textbf{\mas} & \textbf{\mas} & \textbf{\mas} \\
  \bottomrule
\end{tabularx}
\pendiente{Falta la fila «remite al uso lingüístico de ciertos grupos sociales»}
\end{solucion}

\begin{solucion}{\ref{ej:02-searle}}
\begin{tabularx}{\linewidth}{@{}L Y{1.7cm} Y{1.7cm}@{}}
  \toprule
  ¿Qué tipo de regla encierran los siguientes enunciados? & Regulativa & Constitutiva \\
  \midrule
  b) No corresponde decir «pienso de que mañana empiezo la Facultad». & \textbf{\mas} & \textbf{\menos} \\
  c) No se dice «haiga», se dice «haya». & \textbf{\mas} & \textbf{\menos} \\
  d) «Si seguís hablando, la vas a pasar mal» se entiende como amenaza. & \textbf{\menos} & \textbf{\mas} \\
  e) Toda estructura lingüística contiene un núcleo. & \textbf{\menos} & \textbf{\mas} \\
  f) Usar el léxico lunfardo queda prohibido en la radiodifusión argentina. & \textbf{\mas} & \textbf{\menos} \\
  \midrule
  ¿Qué tipo de regla ofrecen las gramáticas normativas o prescriptivas? & \textbf{\mas} & \textbf{\menos} \\
  ¿Qué tipo de regla presentan las gramáticas descriptivas o científicas? & \textbf{\menos} & \textbf{\mas} \\
  ¿Qué tipo de regla es responsable de la regulación interna de una lengua? & \textbf{\menos} & \textbf{\mas} \\
  ¿Qué tipo de regla es responsable de la regulación externa de una lengua? & \textbf{\mas} & \textbf{\menos} \\
  \bottomrule
\end{tabularx}
\end{solucion}

\begin{solucion}{\ref{ej:02-correccion}}
\begin{tabularx}{\linewidth}{@{}L Y{1.5cm} Y{1.5cm}@{}}
  \toprule
  ¿Qué evaluación merecen las siguientes secuencias lingüísticas? & Correcta & Gramatical \\
  \midrule
  2. a. Las diputadas están saliendo y están entrando a otra reunión. & \textbf{\menos} & \textbf{\mas} \\
  \phantom{2.} b. Las diputadas están saliendo de una reunión y están entrando a otra. & \textbf{\mas} & \textbf{\mas} \\
  \addlinespace
  3. a. Trajimos tres remeras iguales para que no haya diferencias de color. & \textbf{\mas} & \textbf{\mas} \\
  \phantom{3.} b. Trajimos tres remeras iguales para que no haiga diferencias de color. & \textbf{\menos} & \textbf{\mas} \\
  \phantom{3.} c. Trajimos tres remeras iguales para que no hay diferencias de color. & \textbf{\menos} & \textbf{\menos} \\
  \addlinespace
  4. a. Tu mujer trabaja todo el día, parece \textit{blecansain}. & \textbf{\menos} & \textbf{\menos} \\
  \phantom{4.} b. Tu \textit{jermu} trabaja todo el día, parece incansable. & \textbf{\masmenos} & \textbf{\mas} \\
  \bottomrule
\end{tabularx}

2. a. [[Las diputadas] [están saliendo Ø] y [están entrando [a otra
reunión]].] La falta de la preposición locativa direccional \forma{de} no impide
la gramaticalidad de la oración, ya que su omisión se considera una cuestión
estilística. b. [[Las diputadas] [están saliendo [de una reunión]] y [están
entrando [a otra]]]. En el contraste de 2a-2b se observa que pueden existir
elementos no manifiestos (presupuestos) dentro de una estructura oracional bien
formada. Los verbos de movimiento como \forma{salir} y \forma{entrar} pueden
tener implícito el argumento locativo direccional que requieren porque se pueden
inferir del significado del mismo verbo, ya que \forma{salir} es \sig{egresar
del lugar donde se está} (presuposición de origen) y \forma{entrar} es
\sig{ingresar al lugar donde no se está} (presuposición de destino). A
diferencia del caso planteado en el ejemplo resuelto 1, en el punto b, en el que
la omisión de los argumentos del verbo produce una estructura agramatical, la
omisión de 2b no resulta en una estructura mal formada.

3. a. [Trajimos tres remeras iguales [para que no [haya] diferencias de
color].] b. [Trajimos tres remeras iguales [para que no [haiga] diferencias de
color].] c. [Trajimos tres remeras iguales [para que no [hay] diferencias de
color].] En el contraste de 3a-3b-3c se muestra el carácter histórico de la
normativa, en el sentido de que a lo largo del tiempo puede ir cambiando: en
3b, una forma existente en otras épocas como \forma{haiga}, hoy día resulta
incorrecta; en cambio, en 3a se evidencia la exigencia del modo subjuntivo en
construcciones finales, lo que torna agramatical la secuencia de 3c.

4. a. [Tu mujer trabaja todo el día, parece \juicio{*}{\forma{ble-cansa-in}}.]
b. [Tu \forma{jer-mu} trabaja todo el día, parece incansable.] En las formas de
4a-4b se extiende el concepto de gramaticalidad, inicialmente propio de la
sintaxis, a la morfología. En 4a, la estructura interna de la palabra
\forma{in-cansa-ble} aparece modificada, con las posiciones del prefijo y el
sufijo invertidas. En 4b la modificación afecta a la estructura silábica:
\forma{mu-jer} se invierte y el resultado es una combinación apropiada
\forma{jer-mu}, frente a otras que serían imposibles para el español, como
\juicio{*}{\forma{jre-um}}, \juicio{*}{\forma{um-jre}}.

\textit{Preguntas sobre 4a y 4b}. (4a) La forma
\juicio{*}{[\forma{ble-cansa-in}]}, además de que carece de un significado
predecible, viola las reglas de formación de palabras (los prefijos deben ir
delante de la raíz y los sufijos, detrás). Esto es el opuesto exacto de lo que
ocurre con la forma \forma{in-cansa-ble}, no solo calculable en cuanto a su
significado, sino también porque responde a las reglas de formación de palabras
en español. (4b) En contraste con lo anterior, ambas formas \forma{mu-jer} y
\forma{jer-mu} tienen viabilidad fonológica, pero difieren en el registro
(v.~§\ref{sec:02-diafasica}): la primera corresponde al español general y la segunda es una
palabra propia del lunfardo (léxico popular argentino). En \forma{jermu} se
aplica el procedimiento típico del lunfardo conocido como \termino{vesre} o
inversión silábica, como en \forma{feca} (\forma{café}), \forma{gotán}
(\forma{tango}), \forma{gomía} (\forma{amigo}). Por otra parte, la forma
\forma{jermu} tiene una interpretación más específica o restringida, ya que no
refiere a cualquier \forma{mujer}, sino a la «esposa o pareja de alguien».
\end{solucion}

\begin{solucion}{\ref{ej:02-registros}}
\begin{consignas}
  \item En 4a, la forma \juicio{*}{[\forma{ble-cansa-in}]} resulta anómala para
        cualquier hablante del castellano, frente a la bien formada
        [\forma{in-cansa-ble}]. En 4b, la forma [\forma{jer.mu}], dado que
        corresponde al léxico lunfardo, tiene aceptación restringida al registro
        coloquial del Río de la Plata.\pendiente{La consigna pregunta por 3a y 3b}
  \item Condiciones de uso, registros y estilos de lengua remiten al ámbito de
        la corrección, ya que lo agramatical, dado que contraviene las reglas
        internas o constitutivas del sistema, es lo que «en condiciones
        normales» el hablante no produce. Así, la gramaticalidad de
        [\forma{jer.mu}], cuyo empleo está fehacientemente atestiguado, se
        encuentra salvada frente a la agramaticalidad de
        \juicio{*}{[\forma{ble-in-cansa}]}, que no sería una forma que los
        hablantes emitirían naturalmente.
\end{consignas}
\end{solucion}

\begin{solucion}{\ref{ej:02-aceptabilidad}}
\begin{tabularx}{\linewidth}{@{}L Y{1.5cm} Y{1.5cm}@{}}
  \toprule
  ¿Qué evaluación merecen las siguientes secuencias lingüísticas? & Gramatical & Aceptable \\
  \midrule
  2. La madre del primo de la novia del hermano de mi suegra está contenta. & \textbf{\mas} & \textbf{\masmenos} \\
  \addlinespace
  3. a. «Toda fuerza armada o reunión de personas que se atribuya los derechos
  del pueblo y peticione a nombre de éste, comete delito de sedición» (CN,
  art.~22). & \textbf{\mas} & \textbf{\mas} \\
  \phantom{3.} b. «Toda fuerza armada o reunión de chabones que se atribuya los
  derechos del pueblo y peticione a nombre de éste, comete delito de sedición»
  (¿CN, art.~22?). & \textbf{\mas} & \textbf{\menos} \\
  \bottomrule
\end{tabularx}

2. La estructura de 2 (como en el ejemplo resuelto 1) no presenta problemas de
gramaticalidad, pero la aceptabilidad resulta al menos dudosa: en cada instancia
se introducen sustantivos que remiten a nuevos referentes (\forma{madre, primo,
novia, hermano, suegra}), cuya acumulación dificulta el procesamiento cognitivo
de la oración. [[La madre [del primo [de la novia [del hermano [de mi
suegra]]]]] está contenta.]

3. La selección léxica de los textos formales, como el de la Constitución
Nacional, explica las diferencias en la aceptabilidad de 3a frente a 3b. En
ambos casos la buena formación oracional está asegurada, solo que el lexema
\forma{chabones} en 3b la torna inaceptable por cuestiones de estilo;
\forma{chabones} es una palabra demasiado coloquial para el registro que
corresponde a un texto formal como la Constitución. En contraste, la forma
\forma{personas} en 3a es la esperable dentro de un texto de naturaleza
jurídica.
\end{solucion}

\begin{solucion}{\ref{ej:02-ingrid}}
\begin{consignas}
  \item La estructura de las oraciones, en todas las repeticiones que se entonan
        en la canción, resulta gramatical; no se advierten formas incompletas,
        como podría ser la falta de algún argumento verbal. Lo que se advierte
        son problemas en la construcción del discurso, lo que se conoce como
        anacoluto, y también arranques de frases que no se completan y que se
        retoman desde el inicio con una nueva oración.
  \item La aceptabilidad de las oraciones formadas resulta dispar debido a la
        selección léxica algo errática que se produce en las últimas dos
        repeticiones de la canción. Con todo, se observa que la elección de
        palabras, aunque presenta variaciones, se produce dentro del mismo campo
        semántico de la forma original; así, según se puede escuchar, aparece
        \forma{ventilador} en lugar de \forma{aire (acondicionado)} y
        \forma{cerrar}, en vez de \forma{abrir (la botella)}.
\end{consignas}
\end{solucion}

\begin{solucion}{\ref{ej:02-ambiguedad}}
(a) Las tres oraciones del ejercicio están bien formadas, son gramaticales,
aunque ambiguas, ya que, fuera de contexto, admiten múltiples significados. La
ambigüedad se relaciona con factores sintácticos en el caso resuelto 1 y en 3,
y con factores léxicos en 2.

2. a. \textit{Antes de salir a navegar, acomodaron correctamente las velas (que
utilizaban para alumbrar)}. b. \textit{Antes de salir, acomodaron correctamente
las velas (que servían para impulsar el barco)}.

3. a. [Yo [no recibiría] a tu hermano en mi oficina.] a'. \textit{Yo no
recibiría a tu hermano en mi oficina (es más, lo ignoraría)}. b. [[Yo] no
recibiría a tu hermano en mi oficina.] b'. \textit{Yo no recibiría a tu hermano
en mi oficina (pero sí mi secretario)}. c. [Yo no recibiría [a tu hermano] en mi
oficina.] c'. \textit{Yo no recibiría a tu hermano en mi oficina (pero sí a su
enviado)}. d. [Yo no recibiría a tu hermano [en mi oficina].] d'. \textit{Yo no
recibiría a tu hermano en mi oficina (pero sí en mi casa)}.

(b) Las estructuras no presentan problemas gramaticales, y la ambigüedad
interpretativa (multiplicidad de significados) deriva de la combinatoria
sintagmática o léxica. En 2, se podrían establecer alternativas léxicas o
especificaciones para desambiguar los lexemas: \forma{Antes de salir a navegar,
acomodaron correctamente las velas para alumbrar} frente a \forma{Antes de salir
a navegar, acomodaron correctamente el velamen del barco}. En 3 se podría dar
una marcación especial al constituyente focalizado,\footnote{El foco es el
segmento sintáctico del enunciado al que se otorga mayor relieve informativo
(\textit{DLE}, en línea). Este concepto se introduce en el \capref{6}.} tal
como ocurre en la oralidad, en la que se enfatiza con la entonación el
constituyente sobre el que se centra la interpretación, frente a la entonación
neutra que caracteriza a la negación del núcleo predicativo: \forma{Yo no
recibiría a tu hermano en mi oficina, \textup{YO} no recibiría a tu hermano en
mi oficina, Yo no recibiría \textup{A TU HERMANO} en mi oficina, Yo no
recibiría a tu hermano \textup{EN MI OFICINA}}.
\end{solucion}

\begin{solucion}{\ref{ej:02-elegancia}}
(a) Las oraciones 2 y 3 están bien formadas. 2 puede señalarse como descortés
(sobre este tipo de oraciones pesa aquello de «el burro adelante, para que no se
espante»). La elegancia de una expresión lingüística, en todo caso, no tiene que
ver con la gramática y la estructura de la lengua, sino con una valoración de
naturaleza social. En el caso de 3, lo que se quiere indicar con la repetición
es la importancia de tomar una decisión, sea esta la que fuera (\textit{NGLE},
§\,48.5e). (b) Las estructuras no presentan problemas gramaticales. La falta de
elegancia se podría resolver en 2 restableciendo el orden de los constituyentes
(\forma{Mi hermano y yo fuimos al cine, Fui al cine con mi hermano}). La
tautología de 3, por su parte, se resolvería con cambios en la estructura
(\forma{Y si se comen unas patatas, todos felices, y si no, también})
sacrificando el valor de expresividad que añade la repetición.
\end{solucion}

\begin{solucion}{\ref{ej:02-referencia}}
(a) Las oraciones 1 y 2 están bien formadas, aunque las referencias de
\forma{selenita} (habitante imaginario de la luna) y de \forma{mancuspia}
(palabra inventada por Cortázar) no remitan a entidades reconocidas en el mundo
real: la buena formación oracional no depende de su mayor o menor veracidad, es
decir, de que su contenido sea cierto o por lo menos verosímil («se non è vero,
è ben trovato»). Ni las ficciones de la vida cotidiana ni la ficción literaria
(la de Cortázar, en el ejemplo) serían posibles si la viabilidad de las
oraciones dependiera de su veracidad. (b) Las estructuras están bien formadas,
carecen de problemas gramaticales. La escasa verosimilitud de 1 y la falta de
certidumbre respecto de a qué remite \forma{mancuspias} en 2 se podrían
resolver introduciendo otros lexemas que, aunque cambien o precisen el
significado, permitirían corroborar que ambas oraciones están bien formadas. Así
en 1 se podría reemplazar \forma{selenita} por \forma{delincuente}: \forma{Un
delincuente nos atacó en la otra cuadra}; y en 2, en vez de \forma{mancuspias},
se podría colocar \forma{araucarias}: \forma{Asimismo agradecemos a Ireneo
Fernando Cruz el habernos iniciado, durante un viaje a San Juan, en el
conocimiento de las araucarias}.
\end{solucion}

\begin{solucion}{\ref{ej:02-contradiccion}}
(a) Las oraciones 2-4 presentan estructuras bien formadas. Sin embargo, en 2 (al
igual que en el ejemplo resuelto 1) se establecen contradicciones entre lo que
expresa la oración y el significado de una de las palabras que contienen. En 2,
la incompatibilidad se produce entre \forma{ciego} (privado de la vista) y el
hecho de \forma{ver el accidente}. En 3 se produce una contraposición entre
antónimos, en este caso, \forma{caliente} y \forma{frío}. En 4 la
contradicción se da entre una palabra y su antónimo. (b) No hay problemas de
estructura gramatical en los ejemplos. En 2 la contradicción es irresoluble
porque entran en conflicto los significados de las palabras: \forma{ciego} y
\forma{vio}. Los casos de 3 y 4 podrían ser salvables. En 3 el contexto
situacional permitiría interpretar la contradicción como aparente asignando un
sentido irónico a las palabras del cliente, que se estaría quejando porque el
café caliente que pidió, el mozo se lo trajo frío. En 4 la contradicción
también se puede resolver interpretando que el nuevo entrenador en ocasiones es
simpático y en otras no, señalando que se trata de una característica variable
de la persona de la que se habla, que puede depender de las circunstancias o del
momento.
\end{solucion}

\begin{solucion}{\ref{ej:02-torito}}
(a) \forma{Fajar}: tr. Can. y Am. Pegar a alguien, golpearlo. \forma{Biaba}: no
se encuentra en el \textit{DLE}. \forma{Macana}: f. coloq. Arg., Bol., Par.,
Perú y Ur. Mentira, desatino. \forma{Cabrero}: adj. coloq. Arg. Que es, o suele
estar, malhumorado, agresivo o irritado. (b) La palabra que no está en el
\textit{DLE} es \forma{biaba}. Se puede encontrar en un diccionario del lunfardo
como el \textit{DEL} \parencite{conde2004} y significa \sig{golpe}. (c)
\forma{Chato} y \forma{fajar} son americanismos, \forma{macana} se usa en varios
países de Sudamérica y \forma{cabrero} es un argentinismo; por lo tanto,
ejemplifican la variedad dialectal. El lunfardismo \forma{biaba} ejemplifica la
variedad diastrática.
\end{solucion}

\begin{solucion}{\ref{ej:02-variacion}}
(a) 2. \forma{Guapo}: adj. Bien parecido. 3. \forma{Rancherío}: m. Arg., Chile
y Ur. Conjunto desordenado y miserable de viviendas precarias en las zonas
rurales. 4. \forma{Agora}: adv. dem. desus. ahora (‖ en este momento). U. c.
rur. \forma{Lectoral}: canónigo lectoral. m. Prebendado de oficio. Es el
teólogo del cabildo, y deberá ser licenciado o doctor en teología. 5.
\forma{Leasing}: Econ. Arrendamiento con opción de compra del objeto arrendado.
6. \forma{Paupérrimo}: adj. sup. cult. de \forma{pobre}. (b) 2 y 3: variación
dialectal. 4: variación diacrónica (\forma{agora}). 5 y 6: variación
diastrática.
\end{solucion}


% =====================================================================
%  capitulos/cap03.tex — Capítulo 3. El sonido y la estructura fonológica
%  Autora: Mercedes Güemes
%  Fuente: Capitulo_03_normalizado.docx
%  Transcripciones: \fon{...} = [ ], \fonol{...} = / /
% =====================================================================
\capitulo{El sonido y la estructura fonológica}{Mercedes Güemes}

\section{Introducción}

En este capítulo se abordan las principales características de los sonidos y
su valor en los sistemas lingüísticos. En las primeras secciones se revisa el
sistema ortográfico del español, considerando el principio fonémico en que se
basa y sus irregularidades. Luego, se presenta el sistema de notación del
Alfabeto Fonético Internacional (AFI), como herramienta para la descripción
articulatoria de los sonidos vocálicos y consonánticos. Posteriormente, el
capítulo aborda nociones centrales de la teoría fonológica, como la distinción
entre fono y fonema, la distribución contrastiva, complementaria y la variación
libre. Asimismo, se describen dos de los principales procesos fonológicos del
español: la asimilación y la neutralización. También se introduce el concepto
de sílaba y se presenta la configuración de su estructura interna. Por último,
se explican los rasgos suprasegmentales más relevantes del español: el acento y
la entonación.

\section{Dos disciplinas complementarias: fonética y fonología}

Las lenguas están formadas por un número limitado de unidades fónicas sin
significado, cuyo contraste permite distinguir un número ilimitado de
significados. Estas unidades, conocidas como \termino{fonemas}, deben
combinarse entre sí para poder construir unidades significativas mayores, como
por ejemplo los morfemas,\footnote{Este concepto se presenta y desarrolla en el
\capref{4}.} las palabras, los enunciados y los textos. Como manifiestan
\textcite[12]{casado2011}, «el conjunto potencialmente infinito de expresiones
que se pueden producir en una lengua está formado por un inventario limitado de
unidades sonoras discretas». A pesar de que en la actualidad buena parte de las
interacciones lingüísticas se basan en la escritura, el lenguaje humano es
fundamentalmente un fenómeno de naturaleza oral.\footnote{La mayor parte de las
lenguas humanas son orales. No obstante, también existen las lenguas de señas;
su gramática y léxico se manifiestan mediante la producción de gestos que se
perciben visualmente.} Las dos disciplinas que se encargan del estudio del
sonido humano y su función comunicativa son la \termino{fonética} y la
\termino{fonología}.

La \termino{fonética} se ocupa de estudiar el aspecto material de los sonidos
del habla, es decir, los que efectivamente pronunciamos, y se interesa
principalmente en tres aspectos que definen las tres ramas que la conforman,
según se muestra en la tabla~\ref{tab:03-ramas}.

\begin{table}[htbp]
  \centering
  \begin{tabularx}{\linewidth}{@{}L L L@{}}
    \toprule
    \textbf{Fonética acústica} & \textbf{Fonética articulatoria} & \textbf{Fonética auditiva} \\
    \midrule
    Se enfoca en la estructura de los sonidos y grupos de sonidos. Se
    concentra principalmente en las características físicas de la onda sonora.
    & Analiza cómo se producen y se articulan los sonidos en el aparato
    fonatorio humano e investiga la fisiología y la anatomía humana vinculada
    con la producción del sonido. También clasifica los sonidos de acuerdo con
    los parámetros que los generan.
    & Estudia la manera en la que el oído recibe y percibe los sonidos. Su
    objetivo es entender el funcionamiento psicoacústico del sistema auditivo y
    determinar el alcance de la audición en general. \\
    \bottomrule
  \end{tabularx}
  \caption{Ramas de la fonética}\label{tab:03-ramas}
\end{table}

La \termino{fonología}, por otra parte, es la subdisciplina de la lingüística
encargada de identificar y describir el comportamiento sistemático de los
sonidos de las lenguas. En consecuencia, para la fonología no es del todo
relevante la materialidad del sonido, sino, según iremos viendo a lo largo de
este capítulo, la manera en que los rasgos sonoros se configuran en un sistema
fonológico particular para contrastar significados. La fonología se encarga de
estudiar el conjunto abstracto de representaciones sonoras compartido por una
comunidad de hablantes y las variaciones recurrentes de los sonidos de habla
que se dan dentro de cada variedad lingüística particular.

\section{Segmentación y representación grafémica}\label{sec:03-grafemas}

Dado que las lenguas son entidades principalmente orales, el ser humano cuenta
con la escritura como forma de representar y almacenar las palabras que se
emiten. Para poder hacer un estudio sobre los sonidos de las lenguas, es
importante comprender las formas en las que la escritura alfabética transcribe
los sonidos lingüísticos. Actualmente, una gran mayoría de las lenguas del
mundo cuenta con un sistema de escritura que se basa en el principio fonémico;
esto significa que la ortografía intenta representar los sonidos del habla
asignando para cada sonido una letra. Sin embargo, los abecedarios no cumplen
con exactitud este principio; por ejemplo, algunas letras (\termino{grafemas})
no representan un único sonido: hay sonidos que se transcriben mediante dos o
más grafemas, y hay grafemas que no se pronuncian. Dicho de otra manera, la
correspondencia fonema-grafema no es perfecta. Por ejemplo, la palabra
\forma{higuera} presenta un grafema que no tiene manifestación sonora, como la
\forma{h}, y la dupla \forma{gu} se pronuncia como un solo sonido, equivalente
al primero de la palabra \forma{gala}.

Como puede observarse, la escritura alfabética presenta particularidades que la
alejan, en mayor o menor medida, de la realidad oral. La distancia que se
produce entre la oralidad y la escritura varía según las lenguas. En comparación
con otras lenguas, el español cuenta con un sistema ortográfico mayormente
transparente, dado que las irregularidades ortográficas no son tantas y que
muchas pueden resumirse en un conjunto de reglas predecibles. Sin embargo,
lenguas como el inglés o el francés\footnote{La ortografía del inglés y el
francés no es transparente debido a su evolución histórica, las influencias de
otros idiomas y la falta de reformas ortográficas que mantuvieran la ortografía
al día con los cambios que se fueron produciendo en la pronunciación.} cuentan
con sistemas ortográficos altamente opacos: no es suficiente conocer cómo se
pronuncia una palabra para saber cómo se escribe y tampoco es muy claro cómo
debe leerse una secuencia de letras. Por este motivo, para poder estudiar los
fenómenos vinculados a la oralidad, es muy importante poder tomar distancia de
la representación alfabética del sonido. Esta no es tarea fácil ya que, en el
proceso de alfabetización se producen cambios estructurales en el cerebro que
generan que nuestras actividades perceptuales de escucha estén vinculadas a la
escritura \parencite[253]{dehaene2014}. Para lograr entrar en los estudios de
los sonidos del habla, este capítulo se propone como paso inicial identificar
los límites de la escritura alfabética del español para, posteriormente,
introducir un sistema de representación fonética más adecuado que permita
reconocer, estudiar y contrastar los sonidos lingüísticos a partir de
principios más generales.

\begin{ejercicio}[\textit{El viento}, de Manu Chao]\label{ej:03-viento}
A continuación se presenta la canción de Manu Chao \textit{El viento}. Luego, a
partir de la letra y del audio, proponemos realizar una serie de ejercicios. El
audio puede escucharse en \url{https://www.youtube.com/watch?v=OK-j4WpClMs}. En
cada caso, se muestra como guía la resolución del primer ítem.

% La letra de la canción se incorpora a mano desde el original (texto protegido).
\textoaincorporar{la letra de \textit{El viento} (Manu Chao)}

\begin{consignas}
  \item Indique cuántos segmentos del habla oral se producen en cada una de las
        siguientes palabras:
        \begin{exe}
          \ex \forma{Ruta}: \respuesta{4}
          \ex \forma{Por}:
          \ex \forma{Va}:
          \ex \forma{Manu Chao}:
          \ex \forma{Carretera}:
          \ex \forma{Hombre}:
        \end{exe}
        \begin{resuelto}
        1. 4. Todos los grafemas de la palabra representan un sonido.
        \end{resuelto}
  \item Escuche la canción y observe la pronunciación de los grafemas marcados
        en negrita. ¿Cómo se puede describir la relación de las letras con los
        sonidos que representan?
        \begin{exe}
          \ex \forma{\textbf{B}abylon} -- \forma{\textbf{v}a}
          \ex \forma{\textbf{R}uta} -- \forma{ca\textbf{rr}etera} -- \forma{volve\textbf{r}á}
          \ex \forma{\textbf{S}in} -- \forma{ra\textbf{z}ón}
          \ex \forma{\textbf{H}ombre} -- \forma{\textbf{h}ambre}
        \end{exe}
        \begin{resuelto}
        1. La pronunciación de los grafemas \forma{b} y \forma{v} no presenta
        diferencias en español (con excepción de algunos hablantes que los
        distinguen a causa del contacto con otras lenguas en las que sí hay
        diferencia, como ocurre en francés o en catalán).
        \end{resuelto}
  \item El grafema inicial de las palabras \forma{cielo} y \forma{celos}
        comparte características de pronunciación con los siguientes pares,
        ¿qué ocurre en cada caso? Marque la respuesta correcta.

        \begin{tabular}{@{}l Y{2.6cm} Y{2.6cm}@{}}
          \toprule
          & Comparten grafema & Comparten fonema \\
          \midrule
          \forma{\textbf{s}in} -- \forma{ra\textbf{z}ón} & & \\
          \forma{\textbf{c}arretera} -- \forma{\textbf{c}uando} & & \\
          \bottomrule
        \end{tabular}
  \item A partir de lo visto en los puntos a) y b), ¿qué conclusiones son
        verdaderas sobre la relación grafema-fonema en español? Justifique cada
        respuesta y aporte ejemplos ilustrativos de cada caso.

        \begin{tabularx}{\linewidth}{@{}l L Y{.5cm} Y{.5cm} L@{}}
          \toprule
          & & V & F & Ejemplo \\
          \midrule
          1. & Dos o más grafemas pueden representar un único sonido. & X & &
               Los grafemas \forma{s}, \forma{c} y \forma{z} se pronuncian igual
               en comunidades seseantes. \\
          2. & Un único grafema puede representar varios sonidos. & & & \\
          3. & Cada grafema del español representa un único sonido. & & & \\
          4. & Dos grafemas contiguos pueden corresponderse con un único sonido. & & & \\
          5. & Todos los grafemas representan un sonido. & & & \\
          6. & Según el contexto, algunos grafemas cambian su pronunciación. & & & \\
          \bottomrule
        \end{tabularx}
  \item De las afirmaciones verdaderas encontradas en el punto anterior,
        identifique cuál ocurre con los grafemas señalados en estos grupos de
        palabras:
        \begin{exe}
          \ex \forma{\textbf{g}ente -- \textbf{j}efe -- \textbf{j}irafa -- \textbf{g}irar}
          \ex \forma{\textbf{g}uante -- \textbf{g}ala -- \textbf{g}erminar -- \textbf{g}ol -- \textbf{g}itano}
          \ex \forma{ta\textbf{x}i -- e\textbf{x}amen -- lé\textbf{x}ico -- ne\textbf{x}o}
          \ex \forma{\textbf{ll}amar -- \textbf{y}a -- \textbf{y}erba -- \textbf{ll}eno}
          \ex \forma{\textbf{g}anar -- \textbf{gu}eto -- \textbf{g}ol -- \textbf{gu}iño -- \textbf{g}ula}
          \ex \forma{\textbf{c}osa -- \textbf{c}álido -- \textbf{qu}echua -- \textbf{qu}ímica -- \textbf{c}ubo}
        \end{exe}
        \begin{resuelto}
        1. Dos o más grafemas pueden representar un único sonido. En este caso,
        \forma{g} (delante de vocales anteriores \fon{e} \fon{i}) y \forma{j}
        representan el mismo fonema \fonol{x}.
        \end{resuelto}
\end{consignas}
\end{ejercicio}

\begin{recuadro}[Seseo y yeísmo]\label{rec:03-seseo}
En algunas regiones hispanohablantes las palabras \forma{caza} y \forma{casa}
se distinguen fonológicamente: \fonol{ˈkaθa} y \fonol{ˈkasa} respectivamente.
Sin embargo, para los hablantes de otras regiones, a las que se denomina
seseantes, existe una única pronunciación para el fonema \fonol{s}, ya que el
fonema \fonol{θ} no se registra, tal como ocurre por ejemplo en el español de
Argentina. Se denomina \termino{seseo}, entonces, a la no diferenciación entre
los sonidos \fonol{s} y \fonol{θ} de \forma{caza} y \forma{casa}. A nivel
ortográfico, las comunidades seseantes representan el fonema \fonol{s} con los
grafemas \forma{c}, \forma{z} o \forma{s}, como ocurre en el primer segmento de
\forma{sapo} \fonol{ˈsapo}, \forma{zapato} \fonol{saˈpato} y \forma{cima}
\fonol{ˈsima}. Por otra parte, muchas comunidades hispanohablantes no presentan
el contraste entre las palabras \forma{calló} y \forma{cayó}. Esto se conoce
como \termino{yeísmo}. La fusión de ambos fonemas, en general, se produce con el
sonido palatal \fon{ʝ}. Sin embargo, en el Río de la Plata, se articula con los
sonidos postalveolares \fon{ʒ} y \fon{ʃ}, fenómeno que se conoce como
\termino{yeísmo rehilado} (v. \capref{15}).
\end{recuadro}

\begin{ejercicio}\label{ej:03-ortografia}
\begin{consignas}
  \item Las siguientes imágenes presentan errores que se vinculan con las
        irregularidades del sistema ortográfico del español. Encuentre los
        errores y explique, a partir de la relación grafema-fonema, por qué se
        han producido. Se muestra la resolución del primer ejemplo como guía.

        \imagenpendiente{Imágenes del ejercicio (carteles con errores
        ortográficos)}

        \begin{exe}
          \ex La ortografía correcta para la palabra \forma{targetas} es
              \respuesta{tarjetas}.
          \ex La ortografía correcta para la palabra \forma{benta} es\ldots
          \ex La ortografía correcta para la palabra \forma{ielo} es\ldots
          \ex La ortografía correcta para la palabra \forma{fabor} es\ldots
          \ex La ortografía correcta para la palabra \forma{cilencio} es\ldots
          \ex La ortografía correcta para la palabra \forma{descancen} es\ldots
          \ex La ortografía correcta para la palabra \forma{tullo} es\ldots
          \ex La ortografía correcta para la palabra \forma{poyo} es\ldots
        \end{exe}
        \begin{resuelto}
        1. La ortografía correcta es \forma{tarjetas}. La confusión se produce
        porque los grafemas \forma{j} y \forma{g} delante de las vocales
        \fon{e}/\fon{i} tienen una pronunciación idéntica.
        \end{resuelto}
  \item Busque ejemplos que evidencien confusiones ortográficas recurrentes
        como el caso que se produce entre \forma{haya} (árbol; verbo
        \forma{haber}) y \forma{halla} (verbo \forma{hallar}).
\end{consignas}
\end{ejercicio}

\section{Producción del sonido}

Para comprender las características básicas de los sonidos, primero es
fundamental visualizar cómo se producen y cuáles son los factores que permiten
lograr la configuración particular de cada uno de ellos. A nivel evolutivo, la
especie humana no ha desarrollado un sistema o aparato específico para producir
sonidos lingüísticos, sino que ha utilizado algunos órganos del aparato
respiratorio (pulmones, laringe, faringe, nariz) y del aparato masticador (boca,
dientes, alvéolos) para la función del habla. Así conformado, el aparato
fonatorio humano es un sistema que permite generar energía sonora a partir de
diferentes mecanismos. La respiración pulmonar es uno de esos mecanismos: el
aire procede de los pulmones y se dirige hacia el exterior atravesando la
cavidad oral o nasal para, de esta forma, emitir sonidos.

La creación de sonido es posible gracias a una serie de presiones y movimientos
que va experimentando el aire en su salida. En el trayecto del pulmón a las
cavidades de salida, boca y nariz (ver figura~\ref{fig:03-aparato}), el aire
pasa por la laringe hacia la faringe. En la laringe, la corriente de aire se
encuentra con los pliegues o cuerdas vocales que forman una abertura llamada
\termino{glotis}. El grado de abertura de la glotis varía según el sonido que se
produzca. Cuando respiramos sin emitir sonido o durante la producción de algunas
consonantes, la glotis está abierta (ver figura~\ref{fig:03-glotis}A). Sin
embargo, para que se configuren los distintos sonidos, las cuerdas deben tensarse
o vibrar de diversas maneras; esto genera que el espacio glotal varíe o se cierre
cuando hablamos (ver figura~\ref{fig:03-glotis}B). A este mecanismo se lo
denomina \termino{fonación}.

\figurapendiente{Aparato fonatorio\protect\pendiente{Falta la figura en el original}}{fig:03-aparato}

\figurapendiente{Salida del aire por la boca y por la nariz\protect\pendiente{Falta la figura en el original}}{fig:03-velo}

\figurapendiente{Glotis y cuerdas vocales. A la izquierda (A), glotis abierta;
a la derecha (B), glotis cerrada}{fig:03-glotis}

Una vez que el aire pasa por la glotis, se encuentra con una serie de cavidades,
comúnmente denominadas supraglóticas: la faringe, la cavidad oral y la cavidad
nasal. Su función principal es generar resonancia para que los sonidos se
modulen; por ello se los denomina órganos resonadores. Los sonidos
faríngeos\footnote{Los sonidos consonánticos faríngeos se producen cuando el
punto de articulación involucra la faringe. Pueden ser oclusivos o fricativos.
Si bien su presencia es escasa en español, son principalmente abundantes en
lenguas semíticas.} se producen por el golpe de aire contra la pared posterior
de la faringe. Los sonidos nasales, como \fon{m} y \fon{n}, se producen cuando el
velo del paladar permanece bajo y permite el egreso del aire a través de la nariz
(ver figura~\ref{fig:03-velo}). Los sonidos orales, por ejemplo \fon{p} y
\fon{s}, se producen cuando el aire sale por la boca; esto es posible cuando el
velo del paladar asciende y bloquea el paso del aire hacia la nariz.

La mayor cantidad de sonidos del habla se produce gracias a la complejidad que
presenta la cavidad bucal humana, integrada por órganos móviles como los labios,
la lengua (con sus partes ápice, dorso y raíz), el paladar blando, el velo del
paladar y la úvula, denominada popularmente \forma{campanilla}. También la boca
contiene puntos fijos de articulación: los dientes, las encías (alvéolos) y el
paladar duro. La diversidad de sonidos que encontramos en las lenguas va a estar
determinada por la acción de las cuerdas vocales en la glotis (si vibran o no),
el punto de contacto de los órganos móviles y fijos de la boca, y cómo se regula
la salida de aire durante todo el proceso. En la próxima sección, se explicarán
los rasgos que permiten clasificar los sonidos de las lenguas humanas.

\begin{ejercicio}\label{ej:03-aparato}
\begin{consignas}
  \item ¿Cuál es el orden correcto en el que sale el aire desde los pulmones
        hacia la boca?
        \begin{exe}
          \ex Pulmones, tráquea, faringe, laringe, cavidad bucal.
          \ex Pulmones, tráquea, cavidad bucal, laringe, faringe.
          \ex Pulmones, tráquea, laringe, faringe, cavidad bucal.
        \end{exe}
        \begin{resuelto}[Punto (a) resuelto]
        El orden correcto es 3.
        \end{resuelto}
  \item Si el velo del paladar está bajo, el sonido que se produce es:
        \begin{exe}
          \ex Oral.
          \ex Nasal.
          \ex Faríngeo.
        \end{exe}
  \item ¿Qué es la \termino{glotis}?
        \begin{exe}
          \ex Un órgano encargado de hacer vibrar las cuerdas vocales.
          \ex Un músculo que produce la fonación.
          \ex El espacio que se forma entre las cuerdas vocales.
        \end{exe}
  \item Distinga las siguientes partes de la boca en órganos fijos y móviles:

        lengua -- alvéolos -- paladar blando -- velo del paladar -- dientes --
        labios -- paladar duro -- úvula
\end{consignas}
\end{ejercicio}

\section{Transcripción fonética y descripción articulatoria de los sonidos}

Como pudo observarse en las secciones anteriores, los sistemas de escritura no
representan con exactitud la oralidad. Para graficar la amplia diversidad que
presentan los sonidos del habla y evitar confusiones, se dispone del Alfabeto
Fonético Internacional (AFI), que identifica cada sonido lingüístico humano ---o
\termino{fono}--- en virtud de sus características articulatorias. El AFI, creado
en 1888 y actualizado cada cierta cantidad de años, consiste en un alfabeto
compuesto de símbolos gráficos y diacríticos\footnote{Los diacríticos son signos
gráficos que modifican el valor de un signo de representación fonética.}
aprobados por la Asociación Fonética Internacional para representar los sonidos
del habla de todas las lenguas. Para no confundir su representación, los fonos o
sonidos lingüísticos se escriben entre corchetes; por ejemplo, cuando leemos
\fon{a} o \fon{p} se está hablando de la pronunciación oral y no de los
grafemas.

En una primera aproximación, los sonidos se pueden distinguir en dos grandes
clases: vocales y consonantes. Las \termino{vocales} son aquellos segmentos que
se producen sin obstrucción en el paso del aire, mientras que las
\termino{consonantes} ocurren cuando se produce una obstrucción total o parcial
durante la salida de la corriente de aire.\footnote{Si bien existen sonidos
vocálicos y consonánticos prototípicos, no todos los fonos se ajustan a esta
distinción. No hay consenso para determinar qué parámetro establece la oposición
clara entre una vocal y una consonante.}

\begin{recuadro}[Para ampliar]
El Alfabeto Fonético Internacional, también conocido como IPA (\textit{International
Phonetic Alphabet}) por sus siglas en inglés, es uno de los alfabetos más
utilizados. Sin embargo, no es el único. Para la transcripción del español,
muchos textos académicos han utilizado el alfabeto diseñado por la
\textit{Revista de Filología Española} (RFE). Asimismo, para entornos digitales
se emplea el alfabeto SAMPA (\textit{Speech Assessment Methods Phonetic
Alphabet}), que es compatible con cualquier teclado.\footnote{\textcite{gurlekian2001}
han hecho una adaptación de SAMPA al español.} La necesidad de contar con
alfabetos de este tipo no solo se circunscribe a las áreas de fonética y
fonología. Numerosas disciplinas requieren de instrumentos adecuados para
transcribir los fenómenos orales, entre ellas se encuentran la dialectología, la
enseñanza de lenguas, la etnolingüística y la lingüística computacional.
\end{recuadro}

\subsection{Vocales}

Dado que no hay interrupción en el paso del aire, los segmentos vocálicos se
articulan en relación con tres parámetros:

\begin{itemize}
  \item grado de abertura de la cavidad bucal,
  \item posición de la lengua respecto del paladar,
  \item redondeamiento de los labios.
\end{itemize}

El grado de abertura bucal permite distinguir entre vocales abiertas, medias y
cerradas (y otros grados intermedios). En el trapecio vocálico que acompaña al
AFI (ver figura~\ref{fig:03-vocales}), se puede observar que la máxima abertura
la tienen las vocales de la base \fon{a, ɶ, ɑ, ɒ}, mientras que en las filas
superiores se encuentran las más cerradas \fon{i, y, ɨ, ʉ, ɯ, u}, que se
articulan cuando la lengua se aproxima al paladar.

\figurapendiente{Vocales (AFI). En el esquema, las vocales que se presentan a la
derecha de otra son redondeadas; así por ejemplo \fon{œ} es redondeada respecto
de \fon{ɛ}, y \fon{o} respecto de \fon{ɤ}. Enlace al AFI:
\url{http://www.internationalphoneticalphabet.org}}{fig:03-vocales}[55mm]

Si miramos el trapecio de izquierda a derecha, las columnas muestran la
distinción entre vocales anteriores, centrales y posteriores. Esta clasificación
está determinada por la posición de la lengua. En la producción de las vocales
anteriores, el dorso de la lengua se eleva hacia el paladar (está adelantado),
como ocurre en \fon{i}, mientras que en las posteriores, la lengua se eleva
sobre el velo del paladar (se posiciona hacia atrás), como cuando pronunciamos
\fon{u}. Por último, el redondeamiento de los labios o labialización distingue
entre los pares izquierda y derecha de cada punto del trapecio. La vocal \fon{i}
se articula sin redondeamiento de los labios, a la vez que \fon{y} es producida
con redondeamiento labial. La figura~\ref{fig:03-resonancia} ilustra los
parámetros articulatorios que distinguen las vocales \fon{a, e, i, o, u}.

\figurapendiente{Articulación de vocales. En las imágenes, tomadas a partir de
resonancias magnéticas funcionales, se puede ver la posición de la lengua, la
abertura de la boca y el redondeamiento de los labios en la producción de las
vocales \fon{a, e, i, o, u}. Video completo:
\url{https://www.youtube.com/watch?v=jaIquq_4560} (Matt Edwards,
EdwardsVoice.com)}{fig:03-resonancia}

\begin{ejercicio}\label{ej:03-vocales}
Las siguientes palabras tienen marcado un grafema que representa un segmento
vocálico. Seleccione la descripción articulatoria que lo describe, como se
muestra en el ejemplo resuelto~1.

\begin{tabular}{@{}l l @{\qquad} l@{}}
  1. \forma{v\textbf{a}}      & \fon{a} & a. cerrado anterior no redondeado \\
  2. \forma{raz\textbf{ó}n}   & \fon{o} & b. cerrado posterior redondeado \\
  3. \forma{r\textbf{u}ta}    & \fon{u} & c. abierto central no redondeado \\
  4. \forma{s\textbf{i}n}     & \fon{i} & d. medio cerrado posterior redondeado \\
  5. \forma{\textbf{e}l}      & \fon{e} & e. medio cerrado anterior no redondeado \\
\end{tabular}

\begin{resuelto}
1. \forma{va} \fon{a}: c. abierto central no redondeado.
\end{resuelto}
\end{ejercicio}

\subsection{Consonantes}

Las \termino{consonantes} se producen cuando los articuladores crean alguna
obstrucción durante el paso del aire. Los sonidos consonánticos se clasifican a
partir de cuatro parámetros articulatorios:

\begin{itemize}
  \item modo de articulación,
  \item punto o zona de articulación,
  \item acción de las cuerdas vocales,
  \item acción del velo del paladar.
\end{itemize}

El \termino{modo de articulación} es la manera en la cual sale el aire durante
la producción de la consonante. Esto depende del grado de acercamiento de los
órganos articuladores y del tipo de salida de la corriente de aire. Los órganos
pueden obstruirse completamente, lo que da lugar a sonidos oclusivos, como
\fon{p} en \forma{\textbf{p}az} o \fon{g} en el sonido inicial de
\forma{\textbf{g}ato}. También pueden estrecharse entre sí o aproximarse; así se
generan los sonidos fricativos, como por ejemplo \fon{s} en
\forma{\textbf{s}apo}, y también los aproximantes, como \fon{j} en
\forma{rub\textbf{i}o}. Los sonidos africados incluyen dos fases de articulación
contiguas: una oclusiva y otra fricativa. La única africada del español es la
que se encuentra en \forma{o\textbf{ch}o}. Esta consonante tiene una fase
oclusiva \fon{t} y otra fricativa \fon{ʃ}; por lo tanto, se representa
\fon{t̠͡ʃ}.\footnote{Las consonantes africadas, dado que se producen como
articulaciones dobles, se representan con los símbolos correspondientes unidos
por un diacrítico superior \fon{t̠͡ʃ}.} Si se produce un contacto rápido de los
órganos, se generan los sonidos vibrantes, que pueden ser múltiples (con muchos
contactos de la lengua contra el paladar, como \fon{r} en \forma{pa\textbf{rr}a})
o simples (un solo contacto, como \fon{ɾ} de \forma{pa\textbf{r}a}). Por último,
en los sonidos laterales la lengua genera una constricción central que ocasiona
que el aire fluya por los costados de la cavidad bucal, como ocurre en el sonido
inicial de \forma{\textbf{l}interna} \fon{l}. Los modos de articulación aparecen
representados en las filas del cuadro del AFI
(figura~\ref{fig:03-consonantes}).

El \termino{punto o zona de articulación} sirve para clasificar los sonidos
consonánticos teniendo en cuenta el lugar o la zona en la cual los articuladores
que intervienen en su producción se conectan o aproximan. El cuadro del AFI
(figura~\ref{fig:03-consonantes}) presenta horizontalmente, de izquierda a
derecha, las posibilidades de articulación del sonido desde la parte más
externa de la boca, los labios, hasta el último punto de contacto, que es la
glotis. De esta manera, los sonidos se clasifican en bilabiales (\fon{b} en
inicio de palabra como \forma{\textbf{b}arco}), labiodentales (\fon{f} en
\forma{je\textbf{f}e}), dentales (\fon{θ} en \forma{ca\textbf{z}a} para las
variedades peninsulares no seseantes), alveolares (\fon{t} en
\forma{\textbf{t}alco}), postalveolares (\fon{ʃ} en \forma{llegar} y \forma{yo}
para la variedad rioplatense), palatales (\fon{ɲ} en \forma{a\textbf{ñ}o}) y
velares (\fon{k} en \forma{\textbf{c}abeza}). A pesar de que las articulaciones
de consonantes posteriores (de la zona uvular, faríngea y glotal) no son
frecuentes en español, son comunes en otras lenguas. Los puntos de articulación
aparecen representados en las columnas del cuadro del AFI.

\figurapendiente{Consonantes pulmonares (AFI, 2020). Las casillas en gris
representan articulaciones consideradas imposibles. Las casillas en blanco son
posibles, pero no fueron documentadas hasta el momento}{fig:03-consonantes}[60mm]

\begin{recuadro}[Para ampliar]
Si bien en las lenguas de tipo indoeuropeo (v. \Capref{1}, nota~\ref{fn:01-indoeuropea}),
casi la totalidad de los sonidos son pulmonares, otras lenguas poseen sonidos
con otros mecanismos de iniciación de la corriente de aire. Existen sonidos no
pulmonares que se generan en la glotis. Se dividen en \termino{eyectivos},
cuando el aire sale de la glotis hacia la boca, e \termino{implosivos}, cuando
el aire ingresa. Asimismo, los llamados \textit{clicks}, que presentan algunas
lenguas originarias africanas, se producen por mecanismos de corriente de aire
que se inician en el velo del paladar. En este caso, el aire ingresa y un cierre
bucal genera el \textit{click}. La figura~\ref{fig:03-nopulmonares} incluye la
transcripción de sonidos no pulmonares del AFI.
\end{recuadro}

\figurapendiente{Consonantes no pulmonares (AFI)\protect\pendiente{El original
dice «AFI, 2026»}}{fig:03-nopulmonares}

La \termino{acción de las cuerdas vocales} es un parámetro que distingue los
sonidos en sordos o sonoros. Si se observa el AFI, se puede apreciar que ciertos
puntos y modos de articulación presentan dos alternativas. Los símbolos de la
izquierda representan los sonidos sordos, que son aquellos que se producen sin
actividad de las cuerdas vocales. Los sonidos de la derecha son sonoros, es decir
que se realizan con vibración de las cuerdas vocales. Por ejemplo, el par de
consonantes \fon{b} y \fon{p} presentan el mismo punto y modo de articulación,
pero se distinguen por la acción de las cuerdas: \fon{b} es sonora y \fon{p} es
sorda. Este rasgo permite distinguir palabras como \forma{bala} y \forma{pala}.

Por último, la \termino{acción del velo del paladar} hace referencia a la salida
principal del aire. Si el velo del paladar desciende durante la articulación, la
salida del aire se produce por la nariz, lo que da lugar a sonidos nasales, como
\fon{n} en \forma{\textbf{n}e\textbf{n}e} y \fon{m} en
\forma{\textbf{m}a\textbf{m}á}. Por otra parte, el resto de los sonidos, que se
conocen como orales, se producen cuando el velo asciende y la salida del aire
ocurre por la boca.

En resumen, cada consonante se puede describir a partir de los cuatro
parámetros mencionados. Por ejemplo, el fono \fon{p} es oclusivo (modo de
articulación), bilabial (zona de articulación), sordo (acción de las cuerdas
vocales) y oral (acción del velo del paladar). Se denomina \termino{descripción
articulatoria} a la enumeración de rasgos que configuran un sonido.

\begin{ejercicio}\label{ej:03-transcripcion}
Una cada palabra con su correspondiente transcripción fonética. (Se muestra el
primer ejemplo resuelto).

\begin{tabular}{@{}l @{\qquad} l@{}}
  A. ruta       & 1. \fon{ˈombɾe} \\
  B. por        & 2. \fon{ˈbjento} \\
  C. suerte     & 3. \fon{ba} \\
  D. va         & 4. \fon{ˈsweɾte} \\
  E. Manu Chao  & 5. \fon{manuˈtʃao} \\
  F. viento     & 6. \fon{ˈruta} \\
  G. el         & 7. \fon{poɾ} \\
  H. hombre     & 8. \fon{el} \\
  I. carretera  & 9. \fon{kareˈteɾa} \\
\end{tabular}

\begin{resuelto}
A. ruta -- 6. \fon{ˈruta}.
\end{resuelto}
\end{ejercicio}

\begin{ejercicio}\label{ej:03-articulacion}
\begin{consignas}
  \item Encuentre en la canción de Manu Chao al menos dos palabras que comiencen
        con sonidos cuyo modo de articulación sea oclusivo o fricativo y complete
        el cuadro. (Se muestra el primer ejemplo resuelto).

        \begin{tabular}{@{}Y{4cm} Y{4cm}@{}}
          \toprule
          \textbf{Oclusivos} & \textbf{Fricativos} \\
          \midrule
          \textit{Ej.}: \forma{carretera} & \textit{Ej.}: \forma{se} \\
          & \\
          \bottomrule
        \end{tabular}
  \item Distribuya las siguientes palabras en la tabla de abajo según el punto
        de articulación de la consonante que las inicia.

        \fon{ˈsweɾte} \fon{ˈmanu} \fon{ˈkwando} \fon{ˈbjento} \fon{tʃao}
        \fon{fɾonˈteɾa} \fon{raˈson} \fon{poɾ}

        \begin{tabular}{@{}*{5}{Y{1.9cm}}@{}}
          \toprule
          \textbf{bilabial} & \textbf{labiodental} & \textbf{alveolar} & \textbf{palatal} & \textbf{velar} \\
          \midrule
          \fon{ˈmanu} & & & & \\
          \bottomrule
        \end{tabular}
  \item En los ejercicios a) y b) encontramos palabras cuyo segmento inicial se
        caracteriza por ser bilabial y oclusivo, ¿qué parámetro articulatorio
        permite distinguir estos sonidos?
  \item Encuentre palabras en la canción de Manu Chao que contengan segmentos
        con las siguientes características articulatorias:
        \begin{exe}
          \ex Alveolar oclusiva nasal sonora: \respuesta{\fon{n} en \fon{ˈbjene}}
          \ex Bilabial oclusiva oral sonora:
          \ex Alveolar vibrante múltiple oral sonora:
          \ex Labiodental fricativa oral sorda:
          \ex Alveolar lateral oral sonora:
          \ex Alveolar oclusiva oral sorda:
        \end{exe}
  \item Realice la descripción articulatoria de los siguientes fonos:
        \begin{exe}
          \ex \fon{β}: \respuesta{bilabial fricativa (aproximante) oral
              sonora}\footnote{V. §\ref{sec:03-alofonos} para la aclaración
              sobre la caracterización de este sonido y la diferencia entre
              fricativas y aproximantes.}
          \ex \fon{ɛ}
          \ex \fon{h}
          \ex \fon{x}
          \ex \fon{ɡ}
          \ex \fon{ʒ}
          \ex \fon{ɔ}
        \end{exe}
  \item Transcriba ortográficamente las siguientes palabras:

        \fon{βeˈɾano} \fon{ˈsizmo} \fon{akoˈplaβa} \fon{ˈoŋɡo} \fon{ˈleʧe}
        \fon{ˈsjero} \fon{arkeˈoloɣa} \fon{ˈweso} \fon{ˈpeɾa} \fon{ˈkaɲa}
        \fon{ˈtaksi}

        \begin{resuelto}
        1. \fon{βeˈɾano}: \forma{verano}.
        \end{resuelto}
  \item Transcriba fonéticamente las siguientes palabras:

        \forma{manzana, caño, llave, dura, neurona, pienso, palo, hora, cereza,
        puerro, cama, ley, terapeuta, jota}

        \begin{resuelto}
        1. \forma{manzana}: \fon{manˈsana}.
        \end{resuelto}
\end{consignas}
\end{ejercicio}

\subsection{Los diptongos}

Al observar la transcripción de las palabras \forma{viento} y \forma{suerte}
(\fon{ˈbjento} \fon{ˈsweɾte}), se puede advertir que la primera vocal del
diptongo se representa de manera diferente de la vocal cuando está sola en la
sílaba (este concepto se desarrolla más adelante en §\ref{sec:03-silaba}). Si
dentro de una sílaba se incluye un diptongo, el núcleo silábico es complejo:
una de las vocales de la sílaba es plena, mientras que las que la preceden o
anteceden tienden a ser aproximantes o transicionales, por lo que su
articulación se acerca, sin llegar a serlo, a la de una consonante. Son
\termino{semiconsonantes} los sonidos \fon{j} y \fon{w} que están en primera
posición del diptongo, como en las palabras \forma{pie} y \forma{suela}, o en
posición final, como en \forma{peine} y \forma{deuda}.

\subsection{La representación del acento}

Se puede observar que en la transcripción de las palabras aparece el diacrítico
\fon{ˈ} que antecede a un segmento o sílaba de la palabra. Se trata de la
transcripción del \termino{acento léxico}. El acento es el realce que se realiza
sobre una sílaba de la palabra en contraste con otras. Debe distinguirse el
\termino{acento prosódico}, que es propio de la oralidad, del acento
ortográfico o \termino{tilde}. En el nivel prosódico, todas las palabras tienen
acento;\footnote{En §\ref{sec:03-suprasegmentales} se explican algunas
excepciones a esta afirmación. A su vez, se ampliarán los detalles sobre este
fenómeno y sus implicancias.} en cambio, a nivel ortográfico, solo en ciertas
palabras y siguiendo reglas específicas, se destacan algunas sílabas acentuadas
con un signo (´). Como para la transcripción fonética lo que cuenta es la
pronunciación oral, el acento prosódico se marca de forma obligatoria siempre e
independientemente de la ortografía que rige en cada lengua.

\section{Los sonidos distintivos o fonemas}

Las lenguas humanas utilizan y combinan un número relativamente limitado de
sonidos para la expresión de todas sus unidades: palabras, enunciados y textos.
Algunos de estos sonidos son distintivos y se denominan \termino{fonemas}. Cada
lengua dispone de una cantidad restringida de unidades distintivas o fonemas
que, si bien carecen de significado, permiten distinguir los significados de
unidades mayores. Por ejemplo, en español las palabras \forma{sol} y \forma{sal}
se diferencian entre sí por los segmentos vocálicos \fonol{o} y \fonol{a},
respectivamente. Esto evidencia que para el español estos sonidos son fonemas,
ya que su contraste permite distinguir dos palabras de la lengua. Lo mismo
ocurre con los pares \forma{palo-polo}, \forma{dato-pato} y \forma{lado-lodo}.

Estos pares de palabras conocidos como \termino{pares mínimos} posibilitan
identificar si dos sonidos constituyen o no fonemas de una lengua mediante la
prueba de conmutación en un mismo contexto. Se trata de una prueba diagnóstica
para reconocer el inventario de unidades opositivas de una lengua (fonemas), en
la medida en que la aparición de uno u otro fonema contrasta significados. Las
palabras \forma{toro} y \forma{coro} constituyen un par mínimo,\footnote{V.
\Capref{1}, §\ref{sec:01-paresminimos}.} ya que se diferencian únicamente por la
consonante inicial. La aparición del fonema \fonol{t} o \fonol{k} en posición
inicial conlleva un cambio de significado producido por la aparición del fono
\fon{t} o \fon{k} en el mismo contexto \fon{\_oro}. Con el descubrimiento de este
par mínimo y ante la posibilidad de contrastar significados para las secuencias
\fon{ˈtoɾo} y \fon{ˈkoɾo}, es posible concluir que los sonidos \fon{t} y \fon{k}
son fonemas distintos, ya que dentro del sistema fonológico del español los
fonos \fon{t} y \fon{k} tienen una distribución contrastiva, es decir, que su
presencia en el mismo contexto opone significados diferentes.

Es importante destacar que los fonemas son unidades abstractas: no se trata de
sonidos en sí, sino de entidades mentales que forman parte de la competencia
lingüística de los hablantes. Se forman a partir de una serie de
\termino{rasgos distintivos} que permiten que los reconozcamos e identifiquemos
frente a otros fonemas. La identidad de un fonema se constituye a partir de las
posibilidades opositivas que genera en el sistema frente a otras unidades. Las
pruebas diagnósticas para el descubrimiento del inventario fonológico de una
lengua tienen sentido únicamente dentro del sistema de cada lengua particular, ya
que cada una tiene su propio conjunto de fonemas. La cantidad de fonemas y su
combinatoria dependen del sistema lingüístico que se esté estudiando.

\begin{recuadro}[Para ampliar]
Los fonemas constituyen un número limitado de elementos que están representados
por rasgos que los diferencian unos de otros en la realidad lingüística mental.
Por ejemplo, el fonema \fonol{p} se puede caracterizar como \sig{consonántico},
\sig{oral}, \sig{oclusivo}, \sig{bilabial}, \sig{sordo}. Los cuatro primeros
rasgos los comparte con \fonol{b}, que es también \sig{consonántico},
\sig{oral}, \sig{oclusivo} y \sig{bilabial}, pero del que se diferencia porque
\fonol{b} es \sig{sonoro}; en cambio, respecto de \fonol{m}, los rasgos que
comparte son solo tres: \sig{consonántico}, \sig{oclusivo},
\sig{bilabial}, ya que \fonol{m} es \sig{nasal} y \sig{sonoro}. A pesar de que
cada lengua difiera en la cantidad de fonemas que posee, todas las lenguas
incluyen un grupo pequeño de fonemas que tienen una amplia posibilidad
combinatoria entre sí. En español, la cantidad de fonemas varía según la
variedad lingüística que se estudie. Se cuentan unos 24 fonemas: diecinueve son
consonantes y cinco son vocales. En algunas variedades ese número cambia por el
seseo y el yeísmo, a los que ya nos referimos en §\ref{sec:03-grafemas}.
\end{recuadro}

\begin{ejercicio}\label{ej:03-paresminimos}
\begin{consignas}
  \item Observe el siguiente corpus e identifique un par mínimo.
  \item Proporcione pares mínimos para las otras palabras del corpus.
\end{consignas}

\fon{mas} \fon{ba} \fon{ˈmanu} \fon{ˈbjento} \fon{la} \fon{lo} \fon{sin}
\fon{ˈombɾe} \fon{ˈsweɾte} \fon{poɾ} \fon{ˈambɾe} \fon{se}

\begin{resuelto}[Punto (a) resuelto]
Par mínimo presente en el corpus: \fon{ˈombɾe} y \fon{ˈambɾe}.
\end{resuelto}
\end{ejercicio}

\begin{ejercicio}\label{ej:03-ue}
\begin{consignas}
  \item Observe el siguiente corpus del español. ¿Se puede identificar algún par
        mínimo?
  \item ¿Qué conclusión se puede sacar de la distribución de los fonos \fon{u}
        y \fon{e} en el sistema fonológico del español?
\end{consignas}

\fon{pus} \fon{seɾ} \fon{ˈluʧe} \fon{ˈpeɣo} \fon{saˈβeɾ} \fon{ˈleʧe}
\fon{ˈuniko} \fon{suˈβiɾ} \fon{suɾ} \fon{enˈtjero} \fon{pes}
\end{ejercicio}

\begin{ejercicio}\label{ej:03-aleman}
\begin{consignas}
  \item Observe el siguiente corpus del alemán. ¿Se puede identificar algún par
        mínimo?
  \item ¿Qué conclusión se puede sacar de la distribución de los fonos \fon{uː}
        y \fon{ʊ} en el sistema fonológico del alemán?
\end{consignas}

\begin{tabular}{@{}lll@{}}
  \fon{ˈp͡fuːl} \sig{estanque}    & \fon{ˈʃpuːkən} \sig{embrujado} & \fon{ˈbaʊm} \sig{árbol} \\
  \fon{ˈhuːn} \sig{pollito}       & \fon{ˈʊntən} \sig{abajo}       & \fon{ˈbʊʃ} \sig{arbusto} \\
  \fon{ˈbʊnt} \sig{colorido}      & \fon{ˈbuːx} \sig{libro}        & \fon{ˈhaʊs} \sig{casa} \\
  \fon{ˈʃpʊkən} \sig{escupir}     & & \\
\end{tabular}
\end{ejercicio}

\section{Los alófonos o variantes de un fonema}\label{sec:03-alofonos}

El sistema fonológico de cada lengua en particular cuenta con un grupo reducido
de fonemas cuya combinación genera unidades mayores significativas; así, por
ejemplo, la combinación de los fonemas \fonol{a} \fonol{l} \fonol{a} produce la
palabra \forma{ala}, mientras que los fonemas \fonol{a} \fonol{s} \fonol{a}
forman la palabra \forma{asa}. Sin embargo, dado que, como ya dijimos, el fonema
es una unidad abstracta, la producción real del sonido lingüístico no ocurre
siempre de la misma manera. En otras palabras, la pronunciación de un fonema
puede variar según el contexto en el que aparezca.

Los sonidos de una lengua nunca se pronuncian exactamente igual, incluso si se
trata de un mismo hablante. Existen factores propios del habla que influyen en
la producción de los sonidos lingüísticos: el contexto, la velocidad, la
intención comunicativa, la variación dialectal o sociolectal, entre otros.

Cada lengua puede presentar variaciones sistemáticas en la pronunciación de un
mismo fonema. Este fenómeno es conocido como \termino{alofonía}. Como puede
observarse en el corpus del ejercicio~\ref{ej:03-bd}, en la palabra
\forma{babero} \fon{baˈβeɾo} podemos reconocer dos segmentos distintos \fon{b} y
\fon{β}, que son alófonos del fonema \fonol{b}, porque su mutación no implica un
cambio de significado. Ambos segmentos son bilabiales orales sonoros, pero
difieren en el modo de articulación, ya que \fon{b} es oclusivo y \fon{β}
representa un sonido fricativo o aproximante.\footnote{Si bien se utilizan los
símbolos del AFI que corresponden al modo de articulación fricativo para las
consonantes intervocálicas \fon{β ð ɣ}, en español estos alófonos se realizan
como aproximantes, ya que se producen con un grado un tanto mayor de apertura
que los sonidos fricativos. A pesar de que algunas transcripciones utilizan un
diacrítico para marcar la diferencia entre una consonante fricativa \fon{ð} y
una aproximante \fon{ð̞}, en este capítulo se simplificará la notación y se
tomarán los símbolos \fon{β ð ɣ} como representaciones de sonidos
aproximantes.} Entonces, la distinción entre el rasgo oclusivo o aproximante en
estos dos fonos en particular no es contrastiva para el español e incluso esta
variación, aunque sea sistemática, resulta imperceptible para los hablantes no
entrenados.

\begin{ejercicio}\label{ej:03-bd}
\begin{consignas}
  \item Lea en voz alta las siguientes palabras y observe la pronunciación de
        las consonantes iniciales y medias marcadas en negrita. ¿Son similares o
        es posible percibir una diferencia en su articulación?

        \begin{tabular}{@{}l@{\qquad}l@{}}
          \forma{\textbf{B}a\textbf{b}ylon} & \forma{\textbf{D}eu\textbf{d}a} \\
          \forma{\textbf{B}e\textbf{b}ida}  & \forma{\textbf{D}e\textbf{d}ica\textbf{d}o} \\
          \forma{\textbf{B}a\textbf{b}ero}  & \forma{\textbf{D}e\textbf{d}al} \\
        \end{tabular}
  \item A continuación se presenta la transcripción fonética de las palabras del
        ejercicio anterior, ¿qué se puede observar?

        \fon{ˈbaβilon} \fon{baˈβeɾo} \fon{deˈðal} \fon{ˈdeṷða} \fon{deðiˈkaðo}
        \fon{beˈβiða}
  \item Realice la descripción articulatoria de \fon{b}, \fon{β}, \fon{d} y
        \fon{ð}. (Se muestra el primer ejemplo resuelto).
        \begin{resuelto}[Ejemplos resueltos]
        \fon{b} bilabial oclusiva oral sonora y \fon{β} bilabial aproximante
        (fricativa) oral sonora.
        \end{resuelto}
  \item ¿Cuál es el contexto que explica la aparición de \fon{b} \fon{β} y
        \fon{d} \fon{ð}?
\end{consignas}
\end{ejercicio}

\section{Distribución contrastiva y distribución complementaria}\label{sec:03-distribucion}

El concepto de contexto es clave para comprender cómo se distribuyen los
sonidos de una lengua. En fonética y fonología, cuando se habla de contexto, se
hace referencia a los segmentos (sonidos) o propiedades suprasegmentales (como
el acento o el tono, ver §\ref{sec:03-suprasegmentales}) que rodean al elemento
que se está analizando. Puede tratarse de una vocal o de una consonante que
antecede o sucede a un segmento. Por ejemplo, en la palabra \forma{sabe}, el
fonema \fonol{b} está en contexto intervocálico, ya que aparece entre las
vocales \fon{a} y \fon{e}, lo que hace que, como vimos, su articulación sea
aproximante \fon{β}. Asimismo, el contexto puede referirse a una unidad
suprasegmental, como el acento. Por ejemplo, en el par de palabras
\forma{paso-pasó}, la vocal \fonol{o} en el primer caso se presenta en una
sílaba no acentuada, mientras que en \forma{pasó} la vocal está en una sílaba
tónica o acentuada.

Según ya hemos dicho, los fonemas no son equivalentes a un sonido «real». Se
trata de unidades abstractas que pueden manifestarse concretamente mediante uno,
dos o más alófonos. En algunos casos, la presencia de uno u otro alófono depende
del contexto en el que aparezcan. Por ejemplo, en español \fon{d} y \fon{ð} son
alófonos de un mismo fonema: el alófono \fon{ð} aparece en contexto
intervocálico (entre vocales) y el alófono \fon{d} después de pausa o después de
consonante nasal, como ocurre en \fon{deðiˈkaðo} y \fon{kondoˈminjo}. Esto
significa que en el sistema fonológico del español estos fonos están en
\termino{distribución complementaria}. Juntos forman una unidad abstracta: el
fonema \fonol{d}. La diferencia material que hay entre estos sonidos no genera
oposiciones de significado para el español; ambos alófonos conforman una única
entidad mental. Este fenómeno es recurrente en español y, además de con
\fonol{b} y \fonol{d}, lo podemos observar también en el fonema \fonol{g}, cuyos
alófonos son \fon{g}-\fon{ɣ}, como en \fon{ˈgato} y \fon{aɣoˈtaðo}. Los ejemplos
muestran que, cuando el contexto fonético cambia, el alófono correspondiente se
realiza en función de las vocales adyacentes. Esto se debe a la coarticulación
de las vocales con las consonantes que les siguen: los órganos no llegan a
cerrarse del todo cuando tienen que producir una consonante entre dos vocales;
por lo tanto, no pueden realizarse como oclusivos.

Sin embargo, cada lengua presenta su propia configuración fonológica. En
inglés, por ejemplo, los mismos sonidos \fon{d} y \fon{ð} anteriores no son
alófonos, sino fonemas distintos porque distinguen significados. El par mínimo
que forman las palabras \forma{day} \fon{deɪ} \sig{día} y \forma{they}
\fon{ðeɪ} \sig{ellos} muestra que en el mismo contexto la producción de un fono
o de otro genera un cambio de significado. En el sistema fonológico del inglés
\fonol{d} y \fonol{ð} son dos fonemas distintos; por lo tanto, los sonidos o
fonos \fon{d} y \fon{ð} se encuentran en \termino{distribución contrastiva},
como ocurre con los fonos \fon{e} y \fon{u} en español. En la
figura~\ref{fig:03-espectrograma} se muestra la imagen de un espectrograma y un
oscilograma en el que aparece una variación alofónica.

\figurapendiente{Alófonos del fonema \fonol{b}. Espectrograma y oscilograma de
la palabra \textit{bebé}. Se puede observar que la producción de la primera
consonante no es similar a la segunda, que está en posición intervocálica. Se
trata de una consonante oclusiva y otra aproximante,
respectivamente}{fig:03-espectrograma}

\begin{ejercicio}\label{ej:03-sz-espanol}
\begin{consignas}
  \item Observe el siguiente corpus del español. Defina las características
        articulatorias de \fon{s} y \fon{z}.
  \item ¿El corpus presenta algún par mínimo?
  \item ¿Qué distribución tienen los fonos \fon{s} y \fon{z} en español?
        Justifique y explicite los contextos en los que aparecen.
\end{consignas}

\begin{tabular}{@{}llll@{}}
  \fon{dezˈðen} & \fon{sin}         & \fon{sin ˈθeɾo} & \fon{oˈsaða} \\
  \fon{raˈson}  & \fon{konsuˈmizmo} & \fon{aβizˈmal}  & \fon{fizɣoneˈar} \\
\end{tabular}

\begin{resuelto}[Punto (a) resuelto]
El sonido \fon{s} es alveolar, fricativo, oral, sordo y \fon{z} es alveolar,
fricativo, sonoro.
\end{resuelto}
\end{ejercicio}

\begin{ejercicio}\label{ej:03-sz-frances}
\begin{consignas}
  \item Observe el siguiente corpus del francés. ¿Presenta algún par mínimo?
  \item ¿Qué distribución tienen los fonos \fon{s} y \fon{z} en francés?
        Justifique y explicite los contextos en los que aparecen.
\end{consignas}

\begin{tabular}{@{}llll@{}}
  \fon{pwasɔ̃} \sig{pez}     & \fon{dəsu} \sig{debajo}   & \fon{dezɛʀ} \sig{desierto} & \fon{sak} \sig{bolsa} \\
  \fon{gɑzɔ̃} \sig{césped}   & \fon{fos} \sig{fosa}      & \fon{pwazɔ̃} \sig{veneno}   & \fon{zon} \sig{zona} \\
  \fon{desɛʀ} \sig{postre}  & \fon{pusjɛʀ} \sig{polvo}  & \fon{vizɔ̃} \sig{visón}     & \\
\end{tabular}

\begin{resuelto}[Punto (a) resuelto]
Pares mínimos del corpus: \fon{pwasɔ̃} \sig{pez} / \fon{pwazɔ̃} \sig{veneno};
\fon{dezɛʀ} \sig{desierto} / \fon{desɛʀ} \sig{postre}.
\end{resuelto}
\end{ejercicio}

\begin{ejercicio}\label{ej:03-quechua}
\begin{consignas}
  \item Observe los siguientes corpus del quechua y del español. Caracterice las
        vocales \fon{i}, \fon{e}, \fon{o} y \fon{u} desde el punto de vista
        articulatorio.
  \item Indique el tipo de distribución y los contextos de aparición de las
        vocales \fon{i} y \fon{e} para el quechua.
  \item Indique el tipo de distribución y los contextos de aparición de las
        vocales \fon{i} y \fon{e} para el español.
  \item Indique el tipo de distribución y los contextos de aparición de las
        vocales \fon{o} y \fon{u} para el quechua.
  \item Indique el tipo de distribución y los contextos de aparición de las
        vocales \fon{o} y \fon{u} para el español.
\end{consignas}

\textit{Quechua}

\begin{tabular}{@{}llll@{}}
  \fon{weqe} \sig{lágrima}   & \fon{suti} \sig{nombre}      & \fon{qespi} \sig{vidrio}    & \fon{niy} \sig{decir} \\
  \fon{apiy} \sig{agarrar}   & \fon{pisiŋga} \sig{maíz}     & \fon{qožqe} \sig{dinero}    & \fon{sapi} \sig{raíz} \\
  \fon{puriy} \sig{andar}    & \fon{t͡ʃeqɑ} \sig{cierto}     & \fon{urpila} \sig{paloma}   & \fon{šiši} \sig{hormiga} \\
  \fon{t͡ʃoqɑ} \sig{arrojar}  & \fon{punčaw} \sig{día}       & \fon{rumi} \sig{piedra}     & \fon{ušpa} \sig{cenizas} \\
  \fon{žoqe} \sig{izquierda} & \fon{aku} \sig{harina}       & \fon{ñuñu} \sig{seno}       & \fon{ažqo} \sig{perro} \\
\end{tabular}

\medskip
\textit{Español}

\begin{tabular}{@{}llll@{}}
  \fon{ˈmisa}       & \fon{seˈkreto}     & \fon{ˈmesa}      & \fon{termiˈno} \\
  \fon{ˈpesa}       & \fon{miˈlimetro}   & \fon{ˈpisa}      & \fon{ˈkeso} \\
  \fon{nesesiˈðad}  & \fon{senˈtiðo}     & \fon{puˈðjente}  & \fon{ˈmuðo} \\
  \fon{xiˈrafa}     & \fon{konoˈsjendo}  & \fon{ˈoxos}      & \fon{ˈkulpa} \\
  \fon{guˈsano}     & \fon{luˈɣar}       & \fon{ˈpuso}      & \fon{kurjosiˈðad} \\
  \fon{ˈposo}       & \fon{lusjeɾˈnaɣa}  & \fon{multiˈtud}  & \fon{ˈmoðo} \\
\end{tabular}

\begin{resuelto}[Ejemplo resuelto del punto (a)]
La vocal \fon{i} es anterior, cerrada y no redondeada.
\end{resuelto}
\end{ejercicio}

\begin{ejercicio}\label{ej:03-ruso}
\begin{consignas}
  \item Observe el siguiente corpus del ruso. ¿Se puede identificar algún par
        mínimo dentro del corpus?
  \item ¿Qué parámetro articulatorio diferencia \fon{ɨ} de \fon{i}?
  \item ¿Qué distribución tienen estos fonos en ruso?
  \item ¿Cuál es el contexto de aparición de cada sonido?
\end{consignas}

\begin{tabular}{@{}llll@{}}
  \fon{dˠɨm} \sig{humo}         & \fon{pastuxʲi} \sig{pastores} & \fon{tenʲisʲist} \sig{tenistas} & \fon{mʲixka} \sig{suave} \\
  \fon{pˠɨtkə} \sig{tortura}    & \fon{stolˠɨ} \sig{mesas}      & \fon{zakonˠɨ} \sig{leyes}       & \fon{vʲitɐmʲin} \sig{vitamina} \\
  \fon{pˠɨˈtatsə} \sig{tratar}  & \fon{ʌlfʌvʲit} \sig{alfabeto} & \fon{zubˠɨ} \sig{dientes}       & \fon{kʲijev} \sig{Kiev} \\
  \fon{vʲit} \sig{monedas}      & \fon{nosˠɨ} \sig{narices}     & & \\
\end{tabular}
\end{ejercicio}

\section{Variación libre}

La distribución que presentan dos alófonos de un mismo fonema es complementaria
si su producción puede predecirse de acuerdo con la posición silábica, el
contacto con otros sonidos o el contexto de aparición, entre otros factores. En
estos casos, la aparición de uno u otro es excluyente y no son intercambiables:
el contexto determina si ocurre uno u otro alófono.

No obstante, existen variaciones alofónicas que no dependen del contexto
fonético, sino de las características de los hablantes. Un mismo hablante puede
pronunciar de una manera un sonido cuando su estilo es más cuidado o formal y de
otra cuando está en una situación más informal. Por ejemplo, en el español de la
Argentina, la realización de \fonol{s} en final de palabra puede alternar entre
tres posibles realizaciones: alveolar, aspirada o directamente elidida. El
sintagma \forma{los amigos} puede pronunciarse entonces \fon{losaˈmiɣos},
\fon{lohaˈmiɣoh}, \fon{loaˈmiɣo}, respectivamente.

También un mismo fonema puede pronunciarse de distinta manera según la región
geográfica a la que pertenece un hablante. Por ejemplo, la realización de la
consonante rótica \fonol{r} (la que se pronuncia en \forma{co\textbf{rr}o}) en
Argentina es muy sensible a la variación geográfica. En muchas regiones del
país, principalmente del noroeste, la articulación de esta consonante es de tipo
asibilada \fon{ř},\footnote{Un sonido asibilado se aproxima a un silbido. No
existe un símbolo exacto para la descripción de este sonido en el AFI. En la
bibliografía, se lo encuentra como \fon{ř}.} mientras que en la región de
Buenos Aires, parte de Santa Fe y Patagonia se articula como vibrante múltiple
\fon{r}. De igual forma, hay zonas en las que conviven ambas pronunciaciones y
se distribuyen según las condiciones sociológicas de los hablantes, como se
muestra en el ejercicio~\ref{ej:03-video}.

La coexistencia de estas alternancias en la pronunciación de ciertos fonemas se
puede explicar a partir de factores de índole sociolectal (variaciones sociales
de los hablantes dentro de una misma comunidad, como el sexo, el nivel de
escolarización, la edad, entre otros), de registro (contextos formales vs.
informales) y dialectales (que dependen de la región
geográfica).\footnote{Estas diferencias de lectos se introdujeron en el
\capref{2}.} A este tipo de variación alofónica se la conoce como
\termino{variación libre} o \termino{estilística}.

\begin{ejercicio}\label{ej:03-bsas}
\begin{consignas}
  \item Observe el siguiente corpus del español de Buenos Aires. ¿Qué rasgos
        fonéticos caracterizan a los sonidos \fon{ʃ} y \fon{ʒ}? Realice su
        descripción articulatoria.
  \item ¿Hay pares mínimos en el corpus? ¿En qué contexto aparecen estos
        segmentos?
  \item ¿Cuál es la distribución de estos sonidos en el español de Buenos Aires?
  \item Escuche los audios e identifique a qué hablante corresponde cada corpus.
\end{consignas}

\imagenpendiente[12mm]{Audios A y B del ejercicio}

\begin{tabular}{@{}l Y{2.4cm} Y{2.4cm}@{}}
  \toprule
  Transcripción ortográfica & \multicolumn{2}{c}{Español de Buenos Aires} \\
  \cmidrule(l){2-3}
  & Hablante 1 & Hablante 2 \\
  \midrule
  Río        & \fon{ˈrio}     & \fon{ˈrio} \\
  Desmayo    & \fon{desˈmaʃo} & \fon{desˈmaʒo} \\
  Carretera  & \fon{kareˈteɾa} & \fon{kareˈteɾa} \\
  Llave      & \fon{ˈʃaβe}    & \fon{ˈʒaβe} \\
  Boya       & \fon{ˈboʃa}    & \fon{ˈboʒa} \\
  Callo      & \fon{ˈkaʃo}    & \fon{ˈkaʒo} \\
  Pollo      & \fon{ˈpoʃo}    & \fon{ˈpoʒo} \\
  Parrilla   & \fon{paˈriʃa}  & \fon{paˈriʒa} \\
  \bottomrule
\end{tabular}

\begin{resuelto}[Punto (a) resuelto]
El fono \fon{ʃ} es postalveolar, fricativo, oral, sordo y el fono \fon{ʒ} es
postalveolar, fricativo, oral, sonoro.
\end{resuelto}
\end{ejercicio}

\begin{ejercicio}\label{ej:03-catamarca}
\begin{consignas}
  \item Observe el siguiente corpus del español de
        Catamarca.\footnote{Agradezco a la Prof. Alejandra Seco por su
        colaboración para elaborar este corpus.} ¿Cuál es la distribución de los
        sonidos \fon{ř} y \fon{r} en el español de Catamarca?
  \item ¿Qué motiva la aparición de los fonos en cada muestra?
  \item Escuche los audios e identifique a qué hablante corresponde cada corpus.
\end{consignas}

\imagenpendiente[12mm]{Audios A y B del ejercicio}

\begin{tabular}{@{}l Y{2.4cm} Y{2.4cm}@{}}
  \toprule
  Transcripción ortográfica & \multicolumn{2}{c}{Español de Catamarca} \\
  \cmidrule(l){2-3}
  & Hablante 1 & Hablante 2 \\
  \midrule
  Río        & \fon{ˈřio}      & \fon{ˈrio} \\
  Desmayo    & \fon{desˈmajo}  & \fon{desˈmajo} \\
  Carretera  & \fon{kařeˈteɾa} & \fon{kareˈteɾa} \\
  Llave      & \fon{ˈjaβe}     & \fon{ˈjaβe} \\
  Pollo      & \fon{ˈpojo}     & \fon{ˈpojo} \\
  Parrilla   & \fon{paˈřija}   & \fon{paˈrija} \\
  Borra      & \fon{ˈbořa}     & \fon{ˈbora} \\
  Carro      & \fon{ˈkařo}     & \fon{ˈkaro} \\
  \bottomrule
\end{tabular}
\end{ejercicio}

\begin{ejercicio}\label{ej:03-video}
\begin{consignas}
  \item Lea la siguiente transcripción fonética. ¿Qué ocurre con los fonos
        \fon{r} y \fon{ř}?
  \item ¿Qué puede motivar la aparición de cada sonido?
  \item Realice la transcripción ortográfica del fragmento.
  \item Escuche el video para analizar el fenómeno.\pendiente{Falta el enlace
        al video}
\end{consignas}

\fon{ˈes to es un ˈři o ˈes te es el ka ˈnal de la ma ˈrja no mo ˈre no ˈtje ne
kas ˈka ða a ˈo ra ˈmi ren lo ke ez la ma ˈrja no mo ˈre no las ˈka jes ke ˈβje
nen de a ˈri βa ðe la ma ˈrja no mo ˈre no la ˈpra ðo eh im po ˈsi βle tran si
ˈtar la ˈmi ren lo ke eh ˈes ta ez la ˈro xas ˈmi ren lo ke ez la ˈro xas}

\begin{resuelto}[Punto (a) resuelto]
Los alófonos \fon{ř} y \fon{r} son producidos por el mismo hablante en
distintas partes de su enunciado.
\end{resuelto}
\end{ejercicio}

\begin{ejercicio}\label{ej:03-yendo}
Observe el error ortográfico presente en la siguiente imagen. ¿Con qué
característica fonológica dialectal se lo puede vincular?

\imagenpendiente{Imagen del ejercicio (error ortográfico con \textit{yendo})}
\end{ejercicio}

\section{Procesos fonológicos}\label{sec:03-procesos}

Cuando hablamos, los sonidos se producen dentro de secuencias cada vez más
complejas. De esta manera, la cadena de habla está formada por sonidos que
presentan características combinatorias cuando se articulan conjuntamente, es
decir, cuando se coarticulan. No es lo mismo estudiar cada sonido de forma
aislada que cuando se produce en un contexto particular. Existen sonidos que al
combinarse con otros se modifican como resultado de esa coarticulación.

Se denomina \termino{proceso fonológico} a toda modificación sistemática que se
produce en un fonema cuando ocurre en determinado contexto. La fonética
combinatoria se ocupa de estudiar estos fenómenos. Si bien existen numerosos
procesos fonológicos, en este apartado nos ocuparemos de dos de los más
frecuentes en español: la neutralización y la asimilación.

Hualde (2013, p.~93)\pendiente{Hualde 2013 no está en la bibliografía (¿2014?)}
define \termino{neutralización} como «situaciones en que dos sonidos que en
general se encuentran en oposición fonémica dejan de contrastar en un contexto
específico». Para ampliar este concepto, vamos a partir de lo que se explicó
anteriormente. Para poder establecer que dos sonidos son contrastivos en una
lengua, se debe encontrar un par mínimo que lo verifique. Si la conmutación de
dos segmentos en la misma posición del mismo contexto resulta en un cambio de
significado, se trata de dos sonidos con distribución contrastiva. Este sería el
caso de \fon{b} y \fon{p}; gracias a la existencia de pares mínimos como
\fon{ˈbata} y \fon{ˈpata} o \fon{ˈpeso} y \fon{ˈbeso} se puede afirmar que se
trata de dos fonemas distintos del español: \fonol{b} y \fonol{p}.

Sin embargo, dos sonidos que contrastan en un contexto fonético pueden no
hacerlo en otro. Así, los mismos dos fonos anteriores, según veremos, en otro
contexto no contrastan. Por ejemplo, en la palabra \forma{óptica} la realización
del segundo segmento puede variar sin alterar el significado: \fon{ˈobtika}
\fon{ˈoptika}. Se observa que en la posición de final de sílaba, conocida como
coda silábica, estos dos fonos pierden su valor opositivo. La representación de
este proceso de neutralización se realiza a través de un
\termino{archifonema}, que consiste en una representación en letra mayúscula
del fonema, y que indica el rasgo en común que tienen los segmentos que se
neutralizan; en el ejemplo, la representación fonológica sería
\fonol{ˈoBtika}.

La \termino{asimilación} es un fenómeno articulatorio que ocurre cuando un
sonido incorpora rasgos de otro sonido del contexto cercano. El ejemplo más
claro de esto es lo que ocurre con las nasales en español en contexto
preconsonántico. Si observamos los ejemplos de \exref{ex:03-nasales}, podemos
concluir que la nasal toma el punto de articulación de la consonante que le
sigue. En el caso de \exref{ex:03-nasales-a}, la nasal se produce como bilabial,
dada la adyacencia con \fon{p}; en \exref{ex:03-nasales-b}, se articula como
labiodental por el contacto con \fon{f}. Los casos de \exref{ex:03-nasales-c} y
\exref{ex:03-nasales-d} ilustran el mismo fenómeno: la nasal se articula como
dental/alveolar frente a \fon{d} y como velar antes de \fon{g}.

\begin{exe}
  \ex\label{ex:03-nasales}
    \begin{xlist}
      \ex\label{ex:03-nasales-a} \forma{un poco}  \fon{umˈpoko}
      \ex\label{ex:03-nasales-b} \forma{un faro}  \fon{uɱˈfaɾo}
      \ex\label{ex:03-nasales-c} \forma{un dólar} \fon{unˈdolaɾ}
      \ex\label{ex:03-nasales-d} \forma{un gato}  \fon{uŋˈgato}
    \end{xlist}
\end{exe}

Este tipo de asimilación es la más frecuente en las lenguas; se trata de la
asimilación regresiva, ya que los rasgos de un segmento son asimilados en el
segmento anterior. Cuando la asimilación es progresiva, la direccionalidad
cambia y es el primer segmento el que afecta al siguiente. Según
\textcite{ullman2006}, el fonema nasal también se ve involucrado en procesos de
asimilación progresiva. Este tipo de proceso se da en palabras que incluyen la
secuencia de fonemas \fonol{d} + \fonol{n} como en \forma{llamadnos}, o a través
de fronteras léxicas como en \forma{red nacional}. Según el autor, en muchas
comunidades hispanohablantes el fonema nasal se realiza como apicodental
\fon{n̪} debido a la influencia de \fon{d}, como sucede en
\fon{redn̪asjoˈnal}.\pendiente{En (1) el original transcribe [umpo'ko],
[uɱ'faro], [undo'lar]; se corrigieron el acento y las vibrantes}

\begin{ejercicio}\label{ej:03-griego}
\begin{consignas}
  \item Observe los siguientes corpus del español y del griego. ¿Qué rasgo
        fonético permite distinguir las consonantes \fon{ɾ} y \fon{r}?
  \item ¿Qué distribución presentan estos sonidos en español? ¿Hay un par
        mínimo?
  \item ¿Ocurre esto en todos los contextos?
  \item ¿Cómo se distribuyen estos sonidos en griego? Justifique a través de los
        contextos de aparición.
  \item Formule una conclusión sobre la distribución de estos fonos en ambas
        lenguas.
  \item Represente fonológicamente las palabras \forma{acusarte} y
        \forma{dador}.
\end{consignas}

\textit{Español}

\begin{tabular}{@{}llll@{}}
  \fon{ˈpero}  & \fon{poˈɾoto} & \fon{akuˈsarte} & \fon{ˈpaɾa} \\
  \fon{ˈroxo}  & \fon{daˈðoɾ}  & \fon{ˈpeɾo}     & \fon{akuˈsaɾte} \\
  \fon{ˈpweðo} & \fon{daˈðor}  & \fon{ˈraβja}    & \fon{ˈpwero} \\
  \fon{ˈpara}  & & & \\
\end{tabular}

\medskip
\textit{Griego}

\begin{tabular}{@{}llll@{}}
  \fon{ˈartos} \sig{pan}           & \fon{kaˈɾoto} \sig{zanahoria} & \fon{ˈɾafi} \sig{armario} & \fon{arˈluba} \sig{tontería} \\
  \fon{ˈorkos} \sig{juramento}     & \fon{ˈzaɾi} \sig{dado}        & \fon{ˈturta} \sig{tarta}  & \fon{aˈɾiða} \sig{pie} \\
  \fon{maˈɾuli} \sig{repollo}      & \fon{ˈkerveɾos} \sig{cerbero} & & \\
\end{tabular}

\begin{resuelto}[Punto (a) resuelto]
Las consonantes \fon{ɾ} y \fon{r} son alveolares, orales y sonoras y se
diferencian por su modo de articulación: \fon{ɾ} es vibrante simple y \fon{r}
es vibrante múltiple.
\end{resuelto}
\end{ejercicio}

\begin{ejercicio}\label{ej:03-arabe}
\begin{consignas}
  \item Observe el siguiente corpus de árabe estándar. En cada cuadro se
        encuentra la transcripción fonética y su traducción. Las palabras están
        separadas con un guion. A partir de las diferencias entre la columna A
        y la columna B, explique qué ocurre con el artículo definido
        \fonol{ʔal}\footnote{El sonido \fon{ʔ} es oclusivo y de articulación
        glotal. El español no cuenta con este sonido en su inventario, pero es
        frecuente en la lengua árabe.} en esta lengua.
  \item Identifique el rasgo que determina la aparición de este fenómeno.
\end{consignas}

\begin{tabular}{@{}ll@{\qquad}ll@{}}
  \toprule
  \multicolumn{2}{@{}l}{\textbf{A}} & \multicolumn{2}{l@{}}{\textbf{B}} \\
  \textit{Transcripción} & \textit{Traducción} & \textit{Transcripción} & \textit{Traducción} \\
  \midrule
  ʔal-qamr           & la luna         & ʔaʃ-ʃams           & el sol \\
  ʔal-walad          & el chico        & ʔat-tiːn           & los higos \\
  ʔal-bint           & la chica        & ʔas-safar          & el viaje \\
  ʔal-ism            & el nombre       & ʔaz-zaffa          & la boda \\
  ʔal-kitaːbul-kabiːr & el gran libro  & ʔan-nahruṭ-ṭawiːl  & el largo río \\
  \bottomrule
\end{tabular}
\end{ejercicio}

\section{La sílaba}\label{sec:03-silaba}

La agrupación de fonemas dentro de una palabra se denomina sílaba. La
\termino{sílaba} es una unidad suprasegmental con estructura jerárquica, dado
que se constituye a partir de un núcleo. El núcleo silábico, que en español es
siempre una vocal, puede presentar elementos que lo preceden y reciben la
denominación de \termino{ataque}. A su vez, puede incluir fonemas posteriores
al núcleo, que se conocen como \termino{coda}. El núcleo y la coda conforman la
\termino{rima}. La estructura completa de la sílaba se visualiza en la
figura~\ref{fig:03-silaba}.

\begin{figure}[htbp]
  \centering
  \begin{forest}
    gramatica, for tree={s sep=4mm}
    [Palabra
      [σ
        [Ataque [Ø]]
        [Rima [Núcleo [a]] [Coda [l]]]]
      [σ
        [Ataque [m]]
        [Rima [Núcleo [i]]]]
      [σ
        [Ataque [d]]
        [Rima [Núcleo [o]] [Coda [n]]]]]
  \end{forest}
  \caption{Estructura interna de la sílaba. Núcleo, ataque y coda de las sílabas
  de \textit{almidón}}\label{fig:03-silaba}
\end{figure}

La palabra \forma{almidón} está formada por tres sílabas: \forma{al -- mi --
dón}. Cada una de esas sílabas muestra una estructura diferente en cuanto a sus
márgenes. En la primera, el núcleo \fonol{a} se acompaña por la coda \fonol{l}.
En la segunda, la sílaba \forma{mi} se forma a partir del ataque \fonol{m}
seguido del núcleo \fonol{i}. La última sílaba está constituida por el núcleo
\fonol{o}, el ataque \fonol{d} y la coda \fonol{n}. Como se puede ver en la
figura~\ref{fig:03-silaba}, las sílabas presentan estructuras diversas. En
español, la sílaba mínima es la que se forma a partir de una única vocal
nuclear, como las primeras sílabas de las palabras \forma{a.mi.go},
\forma{e.le.gir}, \forma{i.le.gal}, \forma{o.la}, \forma{u.so}. El núcleo
vocálico, que constituye la máxima sonoridad de la sílaba, puede estar
acompañado de una semiconsonante que posee menor sonoridad y que se configura,
en términos de Hualde (2013), como un «satélite» de la vocal. Este tipo de
núcleo se observa en el caso de los diptongos de las palabras \forma{ai.re},
\forma{au.la}, \forma{pei.ne}, \forma{pau.ta}.

Asimismo, el ataque o la coda también pueden ser complejos, cuando presentan más
de una consonante, como la primera sílaba de \forma{fra.ter.no},
\forma{trans.por.te} o \forma{cons.tru.ir}. Los grupos consonánticos de ataque
silábico en español son aquellos formados por \fonol{p}, \fonol{b}, \fonol{g},
\fonol{k}, \fonol{f}, \fonol{t} + \fonol{l} o \fonol{r}, por ejemplo
\forma{bro.ma, frá.gil, cro.a.ta}, y el grupo \fonol{dr}, como en
\forma{dra.gón}. Las codas complejas contienen como elemento final siempre
\fonol{s}. Sin embargo, es habitual que las consonantes en coda silábica sufran
procesos de elisión o neutralización.

\begin{ejercicio}\label{ej:03-silabas}
\begin{consignas}
  \item Divida las siguientes palabras en sílabas.

        \forma{barco, implante, halcón, suerte, paseo, protector, antiguo,
        particular, transmutable}
  \item Describa estructuralmente todas las sílabas de las palabras en el
        siguiente cuadro:

        {\footnotesize
        \begin{tabular}{@{}*{9}{c}@{}}
          \toprule
          \multicolumn{3}{c}{Sílaba 3} & \multicolumn{3}{c}{Sílaba 2} & \multicolumn{3}{c}{Sílaba 1} \\
          \cmidrule(r){1-3}\cmidrule(lr){4-6}\cmidrule(l){7-9}
          Ataque & Núcleo & Coda & Ataque & Núcleo & Coda & Ataque & Núcleo & Coda \\
          \midrule
          & & & B & A & R & C & O & \\
          \bottomrule
        \end{tabular}}
  \item Busque palabras que contengan sílabas con las estructuras pedidas abajo
        (V: vocal; C: consonante):
        \begin{exe}
          \ex V: \forma{te.\textbf{a}.tro}
          \ex CV:
          \ex VC:
          \ex CVC:
          \ex CCV:
          \ex CCVC:
          \ex CCVCC:
        \end{exe}
\end{consignas}
\end{ejercicio}

\section{Los rasgos suprasegmentales}\label{sec:03-suprasegmentales}

Además de los segmentos vocálicos y consonánticos que conforman las sílabas, las
palabras y los grupos de palabras presentan \termino{rasgos suprasegmentales}
que se superponen a los segmentos. El acento y la entonación son los rasgos más
importantes para el español. Ambos tienen valor lingüístico tanto a nivel de la
palabra como del sintagma, respectivamente.

\subsection{El acento}

El acento es un rasgo suprasegmental que ocurre dentro del dominio de la sílaba.
Consiste en resaltar una sílaba sobre las otras de su entorno. Esto se realiza
en español a través de distintos recursos acústicos: aumento del tono, la
duración y la energía de pronunciación de la sílaba, entre otras propiedades.
La mayoría de las palabras presentan al menos un acento léxico. En español, la
posición de la sílaba acentuada puede variar.

\begin{itemize}
  \item Las palabras agudas u oxítonas son aquellas que están acentuadas en la
        última sílaba, por ejemplo, \forma{cantar} y \forma{pisó}.
  \item Las palabras graves o paroxítonas son aquellas cuyo acento está en la
        anteúltima sílaba, como \forma{difícil} y \forma{piso}.
  \item Las palabras esdrújulas o proparoxítonas son aquellas con acento en la
        antepenúltima sílaba, como \forma{fantástico} y \forma{cálida}.
\end{itemize}

Si bien cada palabra presenta el acento en una posición, un mismo grupo de
segmentos puede presentar variación en el acento. Este es el caso de las
palabras \fon{ˈteɾmino}, \fon{teɾˈmino} y \fon{teɾmiˈno}. Se puede observar que
la variación de la posición de la sílaba tónica conlleva un cambio de
significado léxico. De esta manera, el acento en español se caracteriza por su
capacidad de distinguir palabras idénticas y, por este motivo, es un rasgo
prosódico que tiene valor fonológico. En resumen, el español es una lengua cuyo
acento es \termino{libre y contrastivo}.

No todas las lenguas presentan esta propiedad. Existen lenguas de \termino{acento
fijo}, en las cuales el acento cae siempre sobre la misma sílaba, como el
francés y el polaco, entre otras.

Si bien de forma aislada todas las palabras presentan un acento primario,
también existen excepciones a la obligatoriedad del acento en español. La mayor
parte de las preposiciones y las palabras funcionales, como las conjunciones o
determinantes, forman un grupo prosódico con la palabra que les sigue; por lo
tanto, no presentan un acento léxico propio.

Asimismo, los adverbios terminados en \forma{-mente} y algunas palabras
compuestas pueden presentar dos acentos. Por ejemplo, la palabra
\forma{fácilmente} posee una sílaba acentuada en la raíz y otra sílaba acentuada
en el sufijo \fon{ˈfasilˈmente}.

\begin{ejercicio}\label{ej:03-acento}
\begin{consignas}
  \item Observe los siguientes corpus del inglés británico y del turco. ¿Qué
        tipo de acento presenta el inglés? ¿Cómo se justifica a partir del
        corpus?
  \item ¿Qué consecuencias trae el cambio en la sílaba tónica en estos pares?
  \item ¿Qué diferencia presenta con el turco? ¿Por qué?
\end{consignas}

\textit{Inglés británico}

\begin{tabular}{@{}lll@{}}
  \fon{ˈʤɜːmən} \sig{alemán}        & \fon{ˈpɾəʊgɾɛs} \sig{el progreso} & \fon{ɪnˈkɾiːs} \sig{aumentar} \\
  \fon{ɪnˈsʌlt} \sig{insultar}      & \fon{ˈtɛlɪfəʊn} \sig{teléfono}    & \fon{ˈɪnkɾiːs} \sig{el aumento} \\
  \fon{ˈmɪnɪstəɾ} \sig{ministro}    & \fon{ˈɪnsʌlt} \sig{el insulto}    & \fon{pɾəʊˈgɾɛs} \sig{progresar} \\
  \fon{ˈfɑːðə} \sig{padre}          & & \\
\end{tabular}

\medskip
\textit{Turco}

\begin{tabular}{@{}lll@{}}
  \fon{baˈkan} \sig{ministro}  & \fon{deˈniz} \sig{mar}          & \fon{teleˈfon} \sig{teléfono} \\
  \fon{aˈçi} \sig{ángulo}      & \fon{baˈba} \sig{padre}         & \fon{meˈraːk} \sig{curiosidad} \\
  \fon{seˈpet} \sig{canasta}   & & \\
\end{tabular}

\begin{resuelto}[Punto (a) resuelto]
El inglés presenta acento libre y contrastivo, como lo muestra el par
\forma{insultar} \fon{ɪnˈsʌlt} y \forma{el insulto} \fon{ˈɪnsʌlt}.
\end{resuelto}
\end{ejercicio}

\subsection{La entonación}

El habla humana transmite información lingüística a través de recursos
acústicos que traspasan el dominio de la palabra. Este fenómeno se conoce como
\termino{entonación} y se produce a partir de variaciones en el tono, la
duración o la velocidad, entre otros.

La entonación posee varias funciones en el lenguaje. Por un lado, permite
distinguir información discursiva al indicar el tipo de enunciado que se
realiza. En este sentido, se vincula con la \termino{modalidad},\footnote{La
modalidad se trata en el \capref{14} de este volumen.} dado que es un recurso
que utiliza el hablante para expresar sus intenciones comunicativas. Por
ejemplo, la secuencia \forma{vas a llegar temprano} se puede interpretar como un
enunciado declarativo \exref{ex:03-modalidad-a}, interrogativo
\exref{ex:03-modalidad-b} o exclamativo \exref{ex:03-modalidad-c}, según el tipo
de entonación que lo acompañe.

\begin{exe}
  \ex\label{ex:03-modalidad}
    \begin{xlist}
      \ex\label{ex:03-modalidad-a} Vas a llegar temprano.
      \ex\label{ex:03-modalidad-b} ¿Vas a llegar temprano?
      \ex\label{ex:03-modalidad-c} ¡Vas a llegar temprano!
    \end{xlist}
\end{exe}

Por otra parte, la entonación permite agrupar palabras en unidades mayores y así
delimitar los enunciados. A través de la presencia o ausencia de pausas o de
cambios tonales, se pueden generar distintos significados para la misma
secuencia lineal de palabras. Así, la entonación, gracias a diversos recursos
acústicos, permite definir o desambiguar las frases que constituyen un
enunciado, fenómeno que se conoce como \termino{fraseo}. La oración \forma{Lupe
dice que cocines} se puede interpretar de dos maneras. Gracias a la presencia de
una pausa, representada con una coma en \exref{ex:03-lupe-a}, el nombre propio
es un vocativo, dado que está usado por el sujeto hablante en función
apelativa, para dirigirse a un interlocutor; mientras que en
\exref{ex:03-lupe-b} el nombre propio funciona como el sujeto del verbo
\forma{decir}.

\begin{exe}
  \ex\label{ex:03-lupe}
    \begin{xlist}
      \ex\label{ex:03-lupe-a} Lupe, dice que cocines. (Alguien se dirige a Lupe
          para pedirle que cocine).
      \ex\label{ex:03-lupe-b} Lupe dice que cocines. (Lupe se dirige a otra
          persona para pedirle que cocine).
    \end{xlist}
\end{exe}

La entonación, además, es el vehículo para codificar información de tipo
pragmático (cortesía, mandato, etc.), sociolectal (dialecto, edad, etc.) o
afectivo (estados de ánimo, humor, etc.).

\begin{recuadro}[El tono]
El tono es un rasgo suprasegmental cuyo correlato acústico es la modulación de
la frecuencia fundamental en la que vibran las cuerdas vocales (F0). La F0 se
calcula sobre la base de la frecuencia de abertura y cierre de las cuerdas
vocales. El tono puede ser alto, medio o bajo, según el rango de la F0. En
español, las variaciones del tono poseen una función pragmática o afectiva; sin
embargo, en otras lenguas llamadas \termino{lenguas tonales}, el tono tiene una
función distintiva a nivel léxico. Por ejemplo, en chino, la sílaba \fonol{ma}
posee diversos significados según su tono: «Con tono alto sostenido significa
\sig{madre}; con tono ascendente empezando desde una altura media,
\sig{cáñamo}; con un contorno bajo ascendente-descendente, \sig{caballo}, y con
tono descendente desde un punto de partida alto, \sig{reñir}» (Hualde, 2013,
p.~260).
\end{recuadro}

\begin{ejercicio}\label{ej:03-puntuacion}
\begin{consignas}
  \item Realice la transcripción ortográfica del siguiente fragmento.
  \item Identifique distintos tipos de enunciados y coloque los signos de
        puntuación correspondientes.
\end{consignas}

\fon{ˈo la ˈko mo es ˈtas te ʃa ˈma βa ˈpa ra kon sul ˈtar te pu ˈðis te kon se
ɣi ˈr el ˈli βro ðe ˈleŋ gwa por fa ˈβor a βi ˈsa me no te ol ˈβi ðes}
\end{ejercicio}

\begin{ejercicio}\label{ej:03-ambiguedad}
Encuentre la ambigüedad que se genera en los siguientes pares de enunciados,
según el tipo de entonación con que se pronuncien. Explique cada caso y
relacione la entonación con el uso de la puntuación.

\begin{exe}
  \ex Cuando llegó, Emilia me entregó el paquete. / Cuando llegó Emilia, me
      entregó el paquete.
  \ex No los aceptaron rápidamente. / No, los aceptaron rápidamente.
  \ex Los clientes que tuvieron problemas deben reclamar en la oficina. / Los
      clientes, que tuvieron problemas, deben reclamar en la oficina.
  \ex Fernando te molestaba por un turno. / Fernando, te molestaba por un
      turno.
\end{exe}

\begin{resuelto}
1. En el primer caso, Emilia es el sujeto del verbo principal
\forma{entregó}, mientras que en el segundo caso es el sujeto del verbo de la
subordinada \forma{llegó}. Si bien es posible interpretar que la persona que
llega y la que entrega el paquete pueden diferir, en el primer caso es más
lógico pensar que las dos acciones son realizadas por la misma persona. En el
primer caso, la frase entonativa termina en \forma{llegó}, mientras que en el
segundo termina en \forma{Emilia}.
\end{resuelto}
\end{ejercicio}

\begin{ejercicio}\label{ej:03-humor}
En estas imágenes se produce un efecto humorístico debido a la ambigüedad en la
interpretación de los rasgos suprasegmentales. Describa qué interpretaciones se
desprenden y cómo se producen los efectos de ambigüedad.

\imagenpendiente{Imágenes A y B del ejercicio}
\end{ejercicio}

\begin{recuadro}[Reflexión final]
En este capítulo se abordaron los contenidos básicos para la comprensión del
aspecto sonoro del lenguaje. En primer lugar, se presentaron las reglas y
excepciones que rigen el sistema de escritura del español; en segundo lugar, se
describieron brevemente los métodos y enfoques para el estudio de los sonidos
del habla, es decir, de la fonética. Finalmente, se expusieron los principios
que explican cómo la organización y modulación del sonido fundamentan la
construcción del significado lingüístico, objeto de estudio de la fonología.
Las herramientas conceptuales y analíticas desarrolladas en este capítulo
constituyen una base indispensable para diversas áreas de aplicación, tales como
la lingüística computacional, los estudios forenses, la lingüística clínica y
la fonoaudiología, entre otras. Con frecuencia no se considera suficientemente
la relevancia de este campo disciplinar. Por ello, es importante destacar que la
fonética y la fonología son saberes transversales que se vinculan no solo con
los demás niveles de la gramática, sino también con múltiples dimensiones del
estudio de la actividad humana.
\end{recuadro}

% --- Secciones finales -------------------------------------------------
\lecturasrecomendadas

\solucionario

\begin{solucion}{\ref{ej:03-viento}}
a. 2 y 3. \forma{por} y \forma{va}: todos los grafemas de la palabra representan
un sonido. 4. \forma{Manu Chao}: hay 7 segmentos; la secuencia \forma{c+h}, en
palabras como \forma{Chao} o \forma{choclo}, representa un solo sonido. 5.
\forma{carretera} \fon{kareˈteɾa}: hay 8 segmentos; la consonante vibrante
múltiple \fonol{r} se escribe en español mediante la grafía doble \forma{rr}
entre vocales. 6. \forma{hombre} \fon{ˈombɾe}: hay 5 segmentos; la letra
\forma{h} en español no representa ningún fonema.

b. 2. \forma{ruta -- carretera -- volverá}. El sistema de representación
ortográfica de las vibrantes puede distinguirse de acuerdo con el contexto de
aparición: la vibrante múltiple \fon{r} se grafica con \forma{rr} entre vocales,
como en \forma{perro}, y con \forma{r} al inicio de palabra o después de las
consonantes \fonol{n}, \fonol{l} o \fonol{s}, como en \forma{rojo},
\forma{Enrique}, \forma{alrededor} y \forma{desratizar}. Por otro lado, la letra
\forma{r} entre vocales corresponde a la vibrante simple \fon{ɾ}, como en
\forma{pero}. 3. \forma{sin -- razón}. En el español de América, los grafemas
\forma{s}, \forma{c} y \forma{z} se pronuncian de la misma manera (seseo). 4.
\forma{hombre -- hambre}. La letra inicial de ambas palabras no representa
ningún sonido.

c. Las palabras \forma{cielo} y \forma{celos} comparten el grafema inicial con
\forma{carretera} y \forma{cuando}. La pronunciación del grafema \forma{c}
delante de las vocales \fon{i}, \fon{e} es similar a la de \forma{sin} y
\forma{razón} en el español americano y canario (comunidades seseantes); en
cambio, delante de \fon{a}, \fon{o} y \fon{u}, \forma{c} se pronuncia como
\fon{k}, como en \forma{casa}, \forma{comida} y \forma{cuenco}.

d. 2. Verdadero; por ejemplo, el grafema \forma{g} se pronuncia \fon{g} en
\forma{gato} y \fon{x} en \forma{gema}.\pendiente{El original dice [ç]} 3.
Falso; por ejemplo, en algunos casos, el grafema \forma{x} representa dos
sonidos, la unión de \fon{k} + \fon{s}, como en \forma{taxi}. 4. Verdadero; por
ejemplo, \forma{rr} en \forma{perro} o \forma{ch} en \forma{chancho}. 5. Falso;
por ejemplo, la letra \forma{h} no representa ningún sonido. 6. Verdadero; por
ejemplo, el grafema \forma{r} puede representar un sonido vibrante múltiple al
inicio de palabra, como en \forma{rata}, y un sonido vibrante simple entre
vocales, como en \forma{para}.

e. 2. \forma{\textbf{g}uante -- \textbf{g}ala -- \textbf{g}erminar --
\textbf{g}ol -- \textbf{g}itano}. Un único grafema puede representar varios
sonidos, como el grafema \forma{g}, que representa el sonido \fon{g} en posición
inicial y delante de \fon{a}, \fon{o} y \fon{u}, pero delante de \fon{e},
\fon{i} se pronuncia como \fon{x}.\pendiente{El original dice [ç]} 3.
\forma{ta\textbf{x}i -- e\textbf{x}amen -- lé\textbf{x}ico -- ne\textbf{x}o}.
Algunos grafemas pueden representar más de un sonido; por ejemplo, el grafema
\forma{x} representa dos sonidos, la unión de \fon{k} + \fon{s}. 4.
\forma{\textbf{ll}amar -- \textbf{y}a -- \textbf{y}erba -- \textbf{ll}eno}. Dos
o más grafemas pueden representar un único sonido; así, en regiones yeístas (ver
recuadro «Seseo y yeísmo»), estos dos grafemas corresponden a la misma
pronunciación, aunque la pronunciación de estas consonantes puede variar en
otras regiones. 5. \forma{\textbf{g}anar -- \textbf{gu}eto -- \textbf{g}ol --
\textbf{gu}iño -- \textbf{g}ula}. Dos grafemas contiguos pueden corresponder a
un único sonido; por ejemplo, la dupla \forma{gu} representa la pronunciación
\fon{g} delante de las vocales anteriores \fon{e}/\fon{i}. 6.
\forma{\textbf{c}osa -- \textbf{c}álido -- \textbf{qu}echua --
\textbf{qu}ímica -- \textbf{c}ubo}. Distintos grafemas representan un solo
sonido; por ejemplo, la dupla \forma{qu} representa la pronunciación \fon{k}
delante de las vocales anteriores \fon{e}/\fon{i}.
\end{solucion}

\begin{solucion}{\ref{ej:03-ortografia}}
a. 1. La ortografía correcta es \forma{tarjetas}. La confusión se produce porque
los grafemas \forma{j} y \forma{g} delante de las vocales \fon{e}/\fon{i} tienen
una pronunciación idéntica. 2 y 4. La ortografía correcta es \forma{venta} y
\forma{favor}. La confusión entre \forma{b} y \forma{v} es habitual ya que ambos
grafemas representan el mismo fonema en todas las variedades del español. 3. La
ortografía correcta es \forma{hielo}. La omisión del grafema \forma{h} se
produce por su falta de manifestación acústica. 5. La ortografía correcta es
\forma{silencio}. 6. La ortografía correcta es \forma{descansen}. Para ciertos
hablantes es habitual confundir los grafemas \forma{s} y \forma{c} delante de
\fonol{e}, \fonol{i}, dado que se pronuncian de igual manera en contextos
seseantes. 7. La ortografía correcta es \forma{tuyo}. 8. La ortografía correcta
es \forma{pollo}. Para los hablantes yeístas es habitual confundir los grafemas
\forma{ll} y \forma{y}, dado que se pronuncian de igual manera. En el caso de
\forma{tuyo} (pronombre posesivo de 2.ª persona singular) y \forma{tullo}
(1.ª persona del verbo \forma{tullir}) se produce un caso de homofonía de dos
palabras que en la escritura diferencian significado.

b. Al igual que ocurre con \forma{tuyo/tullo} y \forma{haya/halla}, otros pares
como \forma{hecho} (verbo \forma{hacer}) / \forma{echo} (verbo \forma{echar}),
\forma{valla} (cerca o vallado) / \forma{vaya} (verbo \forma{ir}), \forma{ay}
(interjección) / \forma{hay} (verbo \forma{haber}) presentan el mismo fenómeno de
homofonía (al menos en la variante rioplatense). La homofonía es un fenómeno en
el cual dos palabras se pronuncian de manera idéntica, pero significan distintas
cosas. En este caso, dado que las palabras se escriben de diferente manera a
pesar de que suenan igual, aunque son homófonas, no son homógrafas.
\end{solucion}

\begin{solucion}{\ref{ej:03-aparato}}
b. Si el velo del paladar está bajo, el sonido que se produce es nasal
(respuesta 2). c. La glotis es el espacio que se forma entre las cuerdas vocales
(respuesta 3). d. Son órganos móviles: lengua, labios, paladar blando, velo del
paladar y úvula. Los alvéolos, el paladar duro y los dientes son articuladores
fijos.
\end{solucion}

\begin{solucion}{\ref{ej:03-vocales}}
\forma{el} \fon{e}: medio cerrado anterior no redondeado. \forma{razón}
\fon{o}: medio cerrado posterior redondeado. \forma{ruta} \fon{u}: cerrado
posterior redondeado. \forma{sin} \fon{i}: cerrado anterior no redondeado.
\end{solucion}

\begin{solucion}{\ref{ej:03-transcripcion}}
B. por -- 7. \fon{poɾ}; C. suerte -- 4. \fon{ˈsweɾte}; D. va -- 3. \fon{ba};
E. Manu Chao -- 5. \fon{manuˈtʃao}; F. viento -- 2. \fon{ˈbjento}; G. el -- 8.
\fon{el}; H. hombre -- 1. \fon{ˈombɾe}; I. carretera -- 9. \fon{kareˈteɾa}.
\end{solucion}

\begin{solucion}{\ref{ej:03-articulacion}}
a.

\begin{tabular}{@{}>{\raggedright\arraybackslash}p{5cm} >{\raggedright\arraybackslash}p{4cm}@{}}
  \toprule
  \textbf{Oclusivos} & \textbf{Fricativos} \\
  \midrule
  \forma{viene; viento; va; carretera; cuando; Babylon; volverá; por} &
  \forma{se; frontera; sin; suerte} \\
  \bottomrule
\end{tabular}

b.

\begin{tabular}{@{}*{5}{Y{1.9cm}}@{}}
  \toprule
  \textbf{bilabial} & \textbf{labiodental} & \textbf{alveolar} & \textbf{palatal} & \textbf{velar} \\
  \midrule
  \fon{ˈmanu} \fon{ˈbjento} \fon{poɾ} & \fon{fɾonˈteɾa} & \fon{raˈson} \fon{ˈsweɾte} & \fon{tʃao} & \fon{ˈkwando} \\
  \bottomrule
\end{tabular}

c. Las palabras \forma{viento} \fon{ˈbjento}, \forma{por} \fon{poɾ} y
\forma{Manu} \fon{ˈmanu} están introducidas por consonantes oclusivas
bilabiales. Las consonantes \fon{p} y \fon{b} contrastan entre sí por la acción
de las cuerdas vocales: \fon{p} es sorda y \fon{b} es sonora. El sonido inicial
de \fon{ˈmanu} es bilabial oclusivo y sonoro. Se diferencia de \fon{b} por la
acción del velo del paladar: \fon{m} es nasal y \fon{b} es oral.

d. Bilabial oclusiva oral sonora: \fon{b} en \fon{ˈbjento}. Alveolar vibrante
múltiple oral sonora: \fon{r} en \fon{ˈruta}. Labiodental fricativa oral sorda:
\fon{f} en \fon{fɾonˈteɾa}. Alveolar lateral oral sonora: \fon{l} en \fon{la}.
Alveolar oclusiva oral sorda: \fon{t} en \fon{kareˈteɾa}.

e. 2. \fon{ɛ}: semiabierta anterior no redondeada.\pendiente{El original dice
«redondeada»} 3. \fon{h}: glotal fricativa oral sorda. 4. \fon{x}: velar
fricativa oral sorda. 5. \fon{ɡ}: velar oclusiva oral sonora. 6. \fon{ʒ}:
postalveolar fricativa oral sonora. 7. \fon{ɔ}: semiabierta posterior
redondeada.

f. 2. \fon{ˈsizmo}: \forma{sismo}. 3. \fon{akoˈplaβa}: \forma{acoplaba}. 4.
\fon{ˈoŋɡo}: \forma{hongo}. 5. \fon{ˈleʧe}: \forma{leche}. 6. \fon{ˈsjero}:
\forma{cierro}. 7. \fon{teˈxaðo}: \forma{tejado}.\pendiente{No figura en la
consigna} 8. \fon{arkeˈoloɣa}: \forma{arqueóloga}. 9. \fon{ˈweso}:
\forma{hueso}. 10. \fon{ˈpeɾa}: \forma{pera}. 11. \fon{ˈkaɲa}: \forma{caña}. 12.
\fon{ˈtaksi}: \forma{taxi}.

g. 2. \forma{caño} \fon{ˈkaɲo}. 3. \forma{llave} \fon{ˈʃaβe}. 4. \forma{dura}
\fon{ˈduɾa}. 5. \forma{neurona} \fon{newˈɾona}. 6. \forma{pienso} \fon{ˈpjenso}.
7. \forma{palo} \fon{ˈpalo}. 8. \forma{hora} \fon{ˈoɾa}. 9. \forma{cereza}
\fon{seˈɾesa}. 10. \forma{puerro} \fon{ˈpwero}. 11. \forma{cama} \fon{ˈkama}.
12. \forma{ley} \fon{lei̯}. 13. \forma{terapeuta} \fon{teɾaˈpewta}. 14.
\forma{jota} \fon{ˈxota}.
\end{solucion}

\begin{solucion}{\ref{ej:03-paresminimos}}
a. Par mínimo presente en el corpus: \fon{ˈombɾe} y \fon{ˈambɾe}. b. Otros
ejemplos de pares mínimos: \fon{mas} \fon{mes}; \fon{ba} \fon{bi}; \fon{ˈmanu}
\fon{ˈmano}; \fon{ˈbjento} \fon{ˈsjento}; \fon{sin} \fon{fin}; \fon{ˈsweɾte}
\fon{ˈmweɾte}; \fon{poɾ} \fon{pon}; \fon{se} \fon{me}.
\end{solucion}

\begin{solucion}{\ref{ej:03-ue}}
a. Se identifican tres pares mínimos en el corpus: \fon{pus} \forma{pus} y
\fon{pes} \forma{pez}, \fon{suɾ} \forma{sur} y \fon{seɾ} \forma{ser},
\fon{ˈluʧe} \forma{luche} y \fon{ˈleʧe} \forma{leche}. b. Dado que la
sustitución de los fonos \fon{e} y \fon{u} en el mismo contexto produce cambio
de significado, como se evidencia en los pares mínimos encontrados, se puede
afirmar que estos fonos son fonemas distintos en español y mantienen una
distribución contrastiva.
\end{solucion}

\begin{solucion}{\ref{ej:03-aleman}}
a. Se puede observar el par mínimo \fon{ˈʃpuːkən} \sig{embrujado} y
\fon{ˈʃpʊkən} \sig{escupir} en el corpus del alemán. b. A partir de la
identificación de este par mínimo se puede concluir que los fonos \fon{uː} y
\fon{ʊ} son fonemas distintos en el sistema fonológico del alemán. Se observa
que su intersustitución genera un cambio de significado; por lo tanto, tienen
una distribución contrastiva en esta lengua.
\end{solucion}

\begin{solucion}{\ref{ej:03-bd}}
a. Si se observa con detenimiento la pronunciación de las consonantes marcadas
en negrita en posición inicial o en medio de la palabra, se puede apreciar que
el contacto de los labios no es idéntico en ambas articulaciones. b. En la
transcripción del corpus se aprecia que después de silencio (es decir, al inicio
de la palabra) los fonemas \fonol{b} y \fonol{d} se producen con cierre completo,
o sea, como oclusivos \fon{b} \fon{d}, mientras que entre vocales ese cierre no
llega a producirse completo y el aire se escapa por entre los articuladores,
produciendo una articulación aproximante \fon{β} \fon{ð}. c. Descripción
articulatoria: \fon{d} alveolar oclusiva oral sonora y \fon{ð} dental
aproximante (fricativa) oral sonora. d. Regla: después de silencio, ocurre el
alófono oclusivo, mientras que en contexto intervocálico la consonante se
realiza como aproximante.
\end{solucion}

\begin{solucion}{\ref{ej:03-sz-espanol}}
b. El corpus no presenta ningún par mínimo. c. En español, \fon{s} y \fon{z}
son alófonos de un mismo fonema y tienen una distribución complementaria, ya que
no comparten contexto. El alófono \fon{s} se realiza después de silencio y entre
vocales, como en \fon{raˈson} y \fon{sin}, y el alófono \fon{z} ocurre delante
de consonante sonora, como en \fon{dezˈðen} y \fon{konsuˈmizmo}.
\end{solucion}

\begin{solucion}{\ref{ej:03-sz-frances}}
b. En francés, los fonos \fon{s} y \fon{z} son fonemas distintos. Ambos ocurren
en los mismos contextos y permiten distinguir significados, es decir que tienen
distribución contrastiva, como se evidencia en los pares mínimos encontrados.
\end{solucion}

\begin{solucion}{\ref{ej:03-quechua}}
a. La vocal \fon{e} es anterior, medio cerrada y no redondeada. La vocal \fon{o}
es posterior, medio cerrada y redondeada. La vocal \fon{u} es posterior, cerrada
y redondeada. b. Los fonos \fon{i} y \fon{e} son alófonos del mismo fonema en
quechua. No comparten contexto y no contrastan significado. El contexto de
aparición de \fon{e} es en contacto con el fonema \fonol{q}, mientras que el
alófono \fon{i} aparece en el resto de los contextos. c. En español, los fonos
\fon{i} y \fon{e} son fonemas distintos, como lo evidencia el par mínimo
\fon{ˈmisa} y \fon{ˈmesa}. Tienen distribución contrastiva y su
intersustitución produce cambio de significado. d. Con los fonos \fon{o} y
\fon{u} del quechua sucede lo mismo que se describió en (b) respecto de \fon{i}
y \fon{e}. e. En español, con los fonos \fon{o} y \fon{u} sucede lo mismo que se
describió en el punto (c) respecto de \fon{i} y \fon{e}, como lo evidencia el
par mínimo \fon{ˈposo} y \fon{ˈpuso}.
\end{solucion}

\begin{solucion}{\ref{ej:03-ruso}}
a. El corpus no presenta ningún par mínimo. b. Ambas vocales son cerradas, pero
se diferencian en la posición que adopta la lengua al producirlas: \fon{i} es
anterior, mientras que \fon{ɨ} es central. c. Estas vocales tienen distribución
complementaria. d. La vocal \fon{i} aparece después de consonantes
palatalizadas (ʲ) y la vocal \fon{ɨ} se produce después de consonantes
velarizadas (ˠ). Palatalización y velarización son procesos fonológicos de
asimilación (ver §\ref{sec:03-procesos}) en los cuales se produce un cambio en
el punto de articulación.\pendiente{El original remite a §8.1.2} Una consonante
palatalizada adelanta hacia el paladar su punto de articulación, mientras que
una consonante velarizada lo desplaza hacia el velo del paladar.
\end{solucion}

\begin{solucion}{\ref{ej:03-bsas}}
b. No hay pares mínimos con contraste de significado. Los dos fonos aparecen en
los mismos contextos, tanto después de silencio, \fon{ˈʃaβe} / \fon{ˈʒaβe}, como
en contexto intervocálico, \fon{ˈpoʃo} / \fon{ˈpoʒo}. c. Dado que no presentan
contraste ni oposición significativa, estos dos fonos son alófonos del mismo
fonema en variación libre; lo que motiva la aparición de uno u otro es el estilo
de habla o las características sociolectales o dialectales de los hablantes que
los producen. d. El hablante 1 corresponde al audio B, ya que produce el sonido
sordo; mientras que el audio A es el hablante 2, que pronuncia la consonante
sonora.
\end{solucion}

\begin{solucion}{\ref{ej:03-catamarca}}
a. No hay pares mínimos con contraste de significado. Las dos consonantes
aparecen en los mismos contextos, tanto después de silencio, \fon{ˈřio} /
\fon{ˈrio}, como en contexto intervocálico, \fon{ˈkařo} / \fon{ˈkaro}. b. La
respuesta es la misma que la del punto (c) del ejercicio~\ref{ej:03-bsas}. c. El
hablante 1 corresponde al audio A, ya que produce una consonante asibilada;
mientras que el audio B es el hablante 2, que pronuncia la vibrante múltiple.
\end{solucion}

\begin{solucion}{\ref{ej:03-video}}
a y b. Es posible que al inicio del video la pronunciación sea más informal o
relajada, mientras que al final se produzca una pronunciación más controlada o
formal. La distribución de ambos alófonos puede estar motivada por el estilo o
el registro. c. Transcripción ortográfica: «Esto es un río. Este es el canal de
la Mariano Moreno, tiene cascada ahora. Miren lo que es la Mariano Moreno, las
calles que vienen de arriba de la Mariano Moreno. La Prado es imposible
transitarla. Miren lo que es\ldots{} esta es la\ldots{} Rojas, miren lo que es
la Rojas». En el video se puede escuchar la alternancia de los alófonos que se
produce durante la grabación, posiblemente debido a la variación del registro.
\end{solucion}

\begin{solucion}{\ref{ej:03-yendo}}
Ciertos errores de ortografía suelen revelar características dialectológicas de
los autores. En este caso, se trata de un hablante que pronuncia la palatal de
inicio de palabra \forma{yendo} como la semiconsonante \fon{j}, lo cual es un
rasgo característico de la pronunciación de la región noroeste de Argentina.
\end{solucion}

\begin{solucion}{\ref{ej:03-griego}}
b. Estos dos sonidos son dos fonemas diferentes y presentan una distribución
contrastiva en español, como evidencia el par mínimo \fon{ˈpaɾa} y \fon{ˈpara}.
c. Sin embargo, hay pares que no muestran contraste, como el caso de
\fon{daˈðoɾ} / \fon{daˈðor} y \fon{akuˈsarte} / \fon{akuˈsaɾte}. Se puede
concluir que en posición final o antes de consonante, las vibrantes se
neutralizan en español. d. En griego, los sonidos \fon{ɾ} y \fon{r} son
alófonos del mismo fonema \fonol{r} y están en distribución complementaria.
Según la descripción fonológica clásica (Kouti, 2010)\pendiente{Kouti 2010 no
está en la bibliografía}, en griego moderno la vibrante se realiza como múltiple
en coda silábica y como simple en otros contextos. e. En español, las
consonantes \fon{ɾ} y \fon{r} son fonemas diferentes, pero en ciertos contextos
no se observa contraste fonémico (hay neutralización). En griego moderno, los
sonidos son alófonos del mismo fonema con distribución complementaria, es
decir, no comparten contexto. f. La representación fonológica de
\forma{acusarte} y de \forma{dador} incluye un archifonema vibrante:
\fonol{akuˈsaRte}, \fonol{daˈdoR}.
\end{solucion}

\begin{solucion}{\ref{ej:03-arabe}}
a. En la columna A, el artículo definido \fonol{ʔal} del árabe estándar
conserva sus rasgos con todas las consonantes iniciales de la palabra que
antecede, mientras que en la columna B la última consonante del artículo
\fonol{l} sufre un proceso de asimilación regresiva. Esto significa que se
asimila a la consonante que le sigue. En este caso, se trata de una asimilación
total, ya que toma todos sus rasgos. b. La asimilación se produce frente a
consonantes coronales, que son aquellas que tienen el punto de articulación en
la zona dentoalveolar, lo que engloba a las consonantes dentales, alveolares y
postalveolares. Con el resto de las consonantes no ocurre este fenómeno, como se
observa en la columna A.
\end{solucion}

\begin{solucion}{\ref{ej:03-silabas}}
a. bar.co -- im.plan.te -- hal.cón -- suer.te -- pa.se.o -- pro.tec.tor --
an.ti.guo -- par.ti.cu.lar -- trans.mu.ta.ble

b.\pendiente{Revisar el orden de las sílabas y algunas celdas del cuadro original}

{\footnotesize
\begin{tabular}{@{}*{9}{c}@{}}
  \toprule
  \multicolumn{3}{c}{Sílaba 3} & \multicolumn{3}{c}{Sílaba 2} & \multicolumn{3}{c}{Sílaba 1} \\
  \cmidrule(r){1-3}\cmidrule(lr){4-6}\cmidrule(l){7-9}
  Ataque & Núcleo & Coda & Ataque & Núcleo & Coda & Ataque & Núcleo & Coda \\
  \midrule
     &    &    & B  & A  & R & C & O  &   \\
     & I  & M  & PL & A  & N & T & E  &   \\
     &    &    &    & HA & L & C & O  & N \\
     &    &    & S  & UE & R & T & E  &   \\
  P  & A  &    & S  & E  &   &   & O  &   \\
  PR & O  &    & T  & E  & C & T & O  & R \\
     & A  & N  & T  & I  &   & G & UO &   \\
  \bottomrule
\end{tabular}

\medskip
\begin{tabular}{@{}*{12}{c}@{}}
  \toprule
  \multicolumn{3}{c}{Sílaba 4} & \multicolumn{3}{c}{Sílaba 3} & \multicolumn{3}{c}{Sílaba 2} & \multicolumn{3}{c}{Sílaba 1} \\
  \cmidrule(r){1-3}\cmidrule(lr){4-6}\cmidrule(lr){7-9}\cmidrule(l){10-12}
  At. & Nú. & Co. & At. & Nú. & Co. & At. & Nú. & Co. & At. & Nú. & Co. \\
  \midrule
  P  & A & R  & T & I &  & C & U &  & L  & A & R \\
  TR & A & NS & M & U &  & T & A &  & BL & E &   \\
  \bottomrule
\end{tabular}}

c. 1. V: \forma{te.\textbf{a}.tro}; 2. CV: \forma{va.\textbf{so}}; 3. VC:
\forma{\textbf{al}.to}; 4. CVC: \forma{\textbf{car}.tu.cho}; 5. CCV:
\forma{\textbf{bro}.ma}; 6. CCVC: \forma{\textbf{plan}.te.o}; 7. CCVCC:
\forma{\textbf{trans}.mi.tir}.
\end{solucion}

\begin{solucion}{\ref{ej:03-acento}}
b. El cambio de acento indica que se trata de dos clases de palabras diferentes:
un sustantivo y un verbo. c. El turco presenta acento fijo. Todas las palabras
del corpus son agudas u oxítonas.
\end{solucion}

\begin{solucion}{\ref{ej:03-puntuacion}}
a y b. Hola, ¿cómo estás? Te llamaba para consultarte, ¿pudiste conseguir el
libro de Lengua? Por favor, avisame. ¡No te olvides!
\end{solucion}

\begin{solucion}{\ref{ej:03-ambiguedad}}
2. \forma{No los aceptaron rápidamente} / \forma{No, los aceptaron
rápidamente}. En el primer caso, \forma{no} niega el verbo \forma{aceptar}, al
formar una única frase con el verbo. En el segundo caso, niega algo dicho
anteriormente y, al estar separado por una pausa, el verbo \forma{aceptar}
mantiene polaridad positiva. 3. En el primer caso, \forma{que tuvieron
problemas} forma una frase entonativa con \forma{clientes} y se interpreta como
un modificador especificativo que restringe el alcance del núcleo y hace
referencia a un grupo específico, «los que tuvieron problemas». En el segundo
caso, al formar una frase entonativa separada, se deben colocar las comas que
representan las pausas o los cambios de tono propios de la oralidad, y
\forma{que tuvieron problemas} se interpreta como un modificador explicativo que
se aplica a todos los clientes (no a un grupo específico) y puede omitirse. 4.
En el primer caso, \forma{Fernando} es el sujeto de la oración. En el segundo,
al separarse con una pausa representada por la coma en la escritura, se
interpreta como un vocativo y constituye una frase entonativa independiente.
\end{solucion}

\begin{solucion}{\ref{ej:03-humor}}
En la foto A, la omisión de la tilde genera un efecto ambiguo. Sin acento
ortográfico no se puede diferenciar en la escritura las formas del verbo en
1.ª persona del presente de indicativo (\forma{mato/asesino}) de las de
3.ª persona del pretérito perfecto simple (\forma{mató/asesinó}). Si bien la
noticia del medio de comunicación intentaba comunicar un hecho pasado («alguien
mató a seis personas ¡y también asesinó a su suegra!»), la omisión de la tilde
permite la lectura que aparece arriba de la foto. La imagen B muestra un chiste
empleado para destacar la importancia de las comas mediante el uso del humor
negro. El cartel contrasta una frase de uso común, \forma{vamos a comer,
niños}, en la cual la palabra \forma{niños} es un vocativo. Al eliminar la coma,
se realiza una lectura sin pausa en la cual \forma{niños} se interpreta como el
COD del verbo \forma{comer}. Esta reestructuración sintáctica, que nace de una
ambigüedad entonativa, resulta humorística.
\end{solucion}


% =====================================================================
%  capitulos/cap04.tex — Capítulo 4. La estructura interna de la palabra
%  Autor: Emiliano De Bin
%  Fuente: Capitulo_04_normalizado.docx
% =====================================================================
\capitulo{La estructura interna de la palabra}{Emiliano De Bin}

\section{Introducción}

En este capítulo presentamos y ejemplificamos los conceptos más relevantes de la
estructura interna de la palabra, de los que se ocupa la parte de la
lingüística que se conoce como morfología. Partimos del reconocimiento de la
existencia de unidades formales segmentables en el interior de las palabras,
establecemos luego las principales clases en las que se las agrupa y mostramos
cómo se organizan jerárquicamente estas unidades formando estructuras que
presentan relaciones de dependencia dentro de las piezas léxicas. Finalmente,
nos detenemos en los distintos procesos morfológicos, tanto los que intervienen
en la formación de palabras como los encargados de la manifestación de las
categorías gramaticales (género, número, tiempo, etc.).

\section{La segmentación morfológica}\label{sec:04-segmentacion}

En una primera aproximación intuitiva, a los hablantes del español nos resulta
evidente que la mayoría de las palabras de nuestra lengua no son un bloque
indivisible, sino que, por el contrario, son susceptibles de ser descompuestas
en unidades menores. Así, por ejemplo, en \forma{submarinas} podemos reconocer
cinco componentes: \forma{sub-}, \forma{-mar-}, \forma{-in-}, \forma{-a-} y
\forma{-s}, cada uno de los cuales contribuye con alguna parte de la
significación del ítem léxico en su conjunto. Estas unidades, que se denominan
\termino{morfemas}, se reconocen porque aparecen recurrentemente en otros ítems
léxicos de la lengua, realizando, por lo general, un aporte semántico semejante.
Encontramos a \forma{sub-} en \forma{subterráneo}, que comparte con
\forma{submarinas} la noción de \sig{estar debajo de}; \forma{mar}, por su
parte, puede conformar por sí solo una palabra, con el significado de
\sig{extensión de agua salada}; \forma{-in-}\footnote{Este segmento
\textit{-in-} es un sufijo adjetivador que aparece en palabras como
\textit{can-in-o} y \textit{cervantin-o}. No debe confundirse con el prefijo
\textit{in-} con valor negativo que se une a adjetivos: \textit{útil} >
\textit{inútil}.} reaparece en \forma{salino} con el valor de \sig{relativo a};
\forma{-a-} y \forma{-s}, como en \forma{alumna} y \forma{autos}, están
asociados a las ideas de \sig{género femenino} y \sig{pluralidad},
respectivamente. De modo análogo, \forma{descolorido} es susceptible de ser
descompuesto en \forma{des-}, que porta la idea de \sig{falta o privación} y
está presente en \forma{desnivel}; \forma{color} que, al igual que \forma{mar},
puede funcionar como una palabra independiente, referida a una propiedad física
de la materia; \forma{-id-}, que indica \sig{cualidad}, y reaparece en
\forma{dolorido}; y \forma{-o-}, vinculado a la noción de \sig{género
masculino}, como en \forma{alumno}.

La posibilidad de realizar esta descomposición habilita algunas preguntas. ¿Qué
realidad poseen los elementos que resultan de ella? ¿Se encuentran todos en el
mismo nivel jerárquico dentro de la palabra? ¿Juegan el mismo rol en su
funcionamiento interno? ¿De qué manera contribuyen estos formantes a la
predictibilidad del significado de los elementos léxicos de la lengua?

\subsection{Las dos ramas de la morfología: léxica y flexiva}\label{sec:04-ramas}

Interrogantes como los anteriores y otros que iremos presentando en este
capítulo constituyen el centro de interés de la morfología. Esta disciplina
gramatical, que estudia el funcionamiento interno de la palabra, está organizada
en dos ramas: la morfología léxica y la morfología flexiva.

Por un lado, la morfología léxica se pregunta por la estructura interna de la
palabra y por los procesos que permiten formar elementos léxicos nuevos sobre
la base de otros ya existentes. Por tanto, los procesos que permiten obtener
\forma{mar-in-o} o \forma{baja-mar} a partir de \forma{mar} son objeto de la
morfología léxica. La primera, \forma{mar-in-o}, se obtiene incorporando a la
forma simple de la palabra un sufijo, que también forma \forma{sal-in-o} a
partir de \forma{sal} o \forma{can-in-o} a partir de \forma{can}. Este proceso
se conoce como derivación. La segunda forma, \forma{bajamar}, proviene de la
unión de \forma{mar} con \forma{baja}, como también ocurre en
\forma{ple(n)a-mar} o \forma{alta-mar}. Se trata de un fenómeno de composición.
El español dispone de una gran variedad de procedimientos derivativos y
compositivos, como también de otros procesos menores de formación de palabras,
que iremos viendo a lo largo de este capítulo.

La morfología léxica también debe dar cuenta de cómo se encuentran
estructuradas por dentro las palabras. Por ejemplo, en el caso de
\forma{descolorido} está claro que el punto de partida para formarla es el
sustantivo \forma{color}, al que se une primero el sufijo \forma{-id-} para
formar \forma{colorido}, \sig{que tiene color}, y, luego, el prefijo
\forma{des-}, que construye el opuesto \forma{descolorido} (\sig{falto de
color}). Según vemos, en cada paso, la formación de la estructura se apoya en
una palabra existente. Sin embargo, hay casos un poco más complejos: así la
situación no es tan transparente en \forma{submarinas}. En el inicio de la
formación de esta palabra se encuentra el sustantivo \forma{mar}, pero, si bien
existe el adjetivo \forma{marino} (\sig{relativo al mar}) al cual podría
unírsele el prefijo \forma{sub-}, «tal análisis no es coherente con la
interpretación semántica del adjetivo:\footnote{En estos casos en que se produce
un desajuste entre la interpretación semántica y la segmentación morfológica,
se suele hablar de «paradojas de encorchetamiento». Dado el nivel introductorio
de esta obra, no profundizaremos en este tema, para lo que remitimos a
\textcite[4374-4376]{piera1999}.} \textit{submarino} no tiene el significado de
\sig{bajo lo relativo al mar} (interpretación que correspondería a la
yuxtaposición de \textit{sub-} y \textit{marino}: [\textit{sub} +
\textit{marino}]), sino \sig{relativo a lo que está bajo el mar} (interpretación
que implica que la estructura del adjetivo es [\textit{sub} + \textit{mar} +
\textit{ino}])» (Darmesteter, 1877, p.~129, cit. en
\cite[4734]{serranodolader1999}), que responde a un tipo de conformación
morfológica que veremos más adelante, en §\ref{sec:04-parasintesis}. Por tanto,
el desafío de la morfología léxica radica en poder explicar los pasos mediante
los que los morfemas se van relacionando entre sí, lo que, según vemos, puede
variar de una palabra a otra.

La morfología flexiva, por otro lado, da cuenta de las diversas formas que puede
adoptar una misma palabra para expresar distintas relaciones gramaticales y
adaptarse al funcionamiento sintáctico. Así, esta segunda vertiente explica, por
caso, que en el adjetivo \forma{submarinas} los morfemas \forma{-a-} y
\forma{-s} manifiestan dos categorías morfológicas que inciden en la combinatoria
sintáctica de la palabra para que pueda establecer relaciones con otras. Por
ejemplo, con un sustantivo masculino plural como \forma{recursos}, el adjetivo se
configura en una forma concordante para formar \forma{recursos submarinos}, la
cual replica el género y número del sustantivo al que modifica y, a la vez,
difiere de otra variante como \forma{submarina}, que es la utilizada para
concordar con un nombre femenino singular, como en \forma{fauna submarina}.
Además del género y el número, otras categorías que se expresan
morfológicamente en español son la persona, el tiempo, el modo y el aspecto.
Todas estas categorías no solo repercuten en la sintaxis: cada una de ellas
expresa algún contenido semántico muy general, como \sig{cantidad} o
\sig{tiempo}, tan relevante para la palabra que su omisión produce
agramaticalidad (v. \capref{2}): \juicio{*}{\forma{criaturas submarina}},
\juicio{*}{\forma{Llegará ayer}} (v. §\ref{sec:04-variantes}).

A lo largo de este capítulo abordaremos paulatinamente cada uno de estos
fenómenos. En todos los casos, el punto de partida del análisis morfológico es la
segmentación de la palabra en los componentes que la forman. A esta partición se
llega por medio de la comparación con otras palabras que exhiben los mismos
morfemas, esto es, observando su distribución dentro de las unidades léxicas.

\section{La unidad morfológica básica: el morfema}\label{sec:04-morfema}

Al igual que ocurre con otras unidades de la gramática, no resulta sencillo
definir con precisión qué es un morfema. En la tradición estructuralista, los
segmentos morfológicos fueron equiparados con el signo saussureano, tanto en la
vertiente norteamericana («forma lingüística con significado constante y
definido», \cite{bloomfield1933}) como en la europea, con la noción de monema
(«unidades dotadas de forma vocal y sentido», \cite{martinet1949}). Desde este
punto de vista, se consideró que era la asociación entre una representación
sonora y un contenido semántico lo que definía a los morfemas, los cuales
constituían «mínimas unidades significativas». Sin embargo, de la segmentación
morfológica de las palabras del español no siempre resultan unidades portadoras
de un contenido definido; así, por ejemplo, en los verbos, la llamada vocal
temática (VT) solo marca la clase conjugacional: \forma{-a} para los verbos de
primera conjugación (\forma{am-a-r}), \forma{-e} para los de segunda
(\forma{tem-e-r}), \forma{-i} para los de tercera (\forma{part-i-r}). Además,
algunos segmentos ni siquiera tienen una función clara en la estructura de la
palabra: se trata de componentes que funcionan como enlaces o que son resabios
de estados anteriores de la lengua. Por caso, en el análisis de una palabra como
\forma{inútil-i-dad-es}, el sufijo \forma{-dad} (\forma{mal-dad, bon-dad})
aparece precedido por una \forma{-i-} y, según veremos en el
ejercicio~\ref{ej:04-terraplanista}, en \forma{fals-e-dad-es}, el mismo sufijo
está precedido por una \forma{-e}, y en ambos casos la única función de estos
segmentos es vincular el sufijo \forma{-dad} con las respectivas bases
\forma{inútil-} y \forma{fals-}. Del mismo modo, también en el
ejercicio~\ref{ej:04-terraplanista}, la \forma{-a-} de \forma{plen-a-mente}
constituye un resto etimológico de la flexión de género que, en la actualidad,
no establece relación de concordancia ni recibe interpretación alguna.

Aceptado el hecho de que no todos los segmentos que se pueden obtener por medio
del análisis resultan significativos, otra objeción a la intuición que plantea a
los morfemas como unidades significativas es el hecho de que no siempre portan
un solo y único valor, sino que pueden aunar varios. Así, a veces un contenido
semántico se expresa de manera separada en un morfema independiente, como la
vocal temática en \forma{am-a-ban} y, en otros casos, ese contenido se presenta
unido en «amalgama» con otro segmento morfológico, como sucede con la vocal
temática en la primera persona singular del presente de indicativo, que se funde
con las otras categorías flexivas de la desinencia verbal (\forma{am-o, tem-o,
part-o}) (v.~§\ref{sec:04-vt}).

Por tanto, basándose en argumentos como los anteriores, \textcite[4318]{pena1999}
defiende una definición más amplia del morfema, como «la unidad mínima del
análisis morfológico y, en definitiva, del análisis gramatical», en el sentido
de que se trata de una unidad identificable dentro de la palabra, tenga o no
significado. En este último caso puede tener alguna de estas tres funciones:

\begin{itemize}
  \item distintiva o sin significado preciso (\forma{conocer / re-conocer}),
  \item clasificatoria (como la VT),
  \item formativa (como la \forma{-a} que aparece en algunos adverbios en
        \forma{-mente}).
\end{itemize}

\begin{ejercicio}[Terraplanistas]\label{ej:04-terraplanista}
Lea atentamente el siguiente fragmento extraído de una nota del diario español
\textit{El País} y luego realice las actividades indicadas más abajo.

% El fragmento del artículo se incorpora a mano desde el original (texto protegido).
\textoaincorporar{el fragmento de «No puedes convencer a un terraplanista y eso
debería preocuparte» (\textit{El País}, 27/02/2019,
\url{https://elpais.com/elpais/2019/02/27/ciencia/1551266455_220666.html}), con
las palabras subrayadas: \textit{convencer, terraplanista, negar, más,
falsedades, achatada, gigantesca, plenamente, representativos, explican,
entronización, pseudocientíficos}}

\textbf{Parte A}
\begin{consignas}
  \item Segmente en sus morfemas constitutivos las palabras subrayadas en el
        texto.
  \item Justifique el reconocimiento de cada morfema aportando al menos otras
        dos palabras en las que dicho morfema se encuentre presente con el mismo
        significado.
  \item Ofrezca una paráfrasis del significado y analice el aporte semántico que
        cada morfema realiza en los casos de \forma{falsedades},
        \forma{convencer}, \forma{plenamente} y \forma{explican}.
\end{consignas}

\begin{resuelto}
a) \forma{con-venc-e-r}. b) \forma{con-}: \forma{con-tribuir, con-tratar};
\forma{-venc-}: \forma{venc-iste, in-venc-ible}; \forma{-e-}: \forma{tem-e-rías,
respond-e-mos}; \forma{-r}:\footnote{El morfema \textit{-r} es el sufijo que
indica que se trata de una forma verbal de infinitivo.} \forma{ama-r,
parti-r}. c) \forma{fals-e-dad-es}. Significado: \sig{(dichos) que no se ajustan
a la verdad}. Aporte de cada morfema: \forma{fals-} \sig{que no se ajusta a la
verdad}; \forma{-dad-} \sig{propiedad de}; \forma{-es} \sig{plural}.
\end{resuelto}

\textbf{Parte B}

Decida cuáles de las generalizaciones expresadas más abajo describen con mayor
precisión lo observado.

\begin{consignas}
  \item Con respecto a la significatividad de los morfemas:
        \begin{exe}
          \ex Todos los morfemas que se pueden segmentar por medio del análisis
              distribucional realizan siempre un aporte semántico reconocible a
              la palabra.
          \ex No todos los morfemas que se pueden segmentar por medio del
              análisis distribucional realizan siempre un aporte semántico
              reconocible a la palabra: algunos de ellos tienen función de
              enlace entre otros morfemas o son vestigios de estadios anteriores
              de la lengua.
        \end{exe}
        \begin{resuelto}
        Opción correcta: 2.
        \end{resuelto}
  \item Con respecto a los límites del análisis formal:
        \begin{exe}
          \ex El análisis formal permite reconocer formantes morfológicos que se
              integraron a la palabra en otros estados de lengua anteriores y
              cuyo significado no es reconocible
              sincrónicamente.\footnote{\textit{Sincronía} y \textit{diacronía}
              son términos que provienen de \textcite{saussure1945}. La
              sincronía estudia un fenómeno tal como es en un momento
              específico, mientras que la diacronía estudia la evolución de ese
              mismo fenómeno a lo largo del tiempo.}
          \ex El análisis formal permite reconocer solamente formantes
              morfológicos integrados a la palabra sincrónicamente, y cuyo
              significado resulta siempre reconocible.
        \end{exe}
        Opción correcta:\completar[1]
  \item Con respecto a la relación entre los segmentos reconocidos y sus
        significados:
        \begin{exe}
          \ex Existen segmentos morfológicos que pueden expresar distintos
              significados y significados que pueden expresarse por medio de
              distintos segmentos morfológicos.
          \ex Todos los segmentos morfológicos expresan siempre los mismos
              significados, pero no todos los significados se expresan siempre
              en los mismos segmentos morfológicos.
          \ex Existen segmentos morfológicos que pueden expresar distintos
              significados, pero cada significado se expresa siempre en el mismo
              segmento morfológico.
          \ex Todos los segmentos morfológicos aportan siempre significados
              distintos.
        \end{exe}
        Opción correcta:\completar[1]
\end{consignas}
\end{ejercicio}

\subsection{La alomorfía}\label{sec:04-alomorfia}

Algunos morfemas del español están asociados a más de una manifestación
fonológica o morfo. Estas variantes de realización, que guardan entre sí un
parecido formal y se asocian todas de manera estable al mismo significado, se
denominan \termino{alomorfos}.

La alomorfía puede originarse en causas diversas. En algunos casos, la
alternancia puede estar motivada por factores externos al sistema gramatical.
Por ejemplo, por haber aparecido como resultado del cambio lingüístico de la
palabra. Así, las palabras \forma{apert-ura}, \forma{abert-ura} y
\forma{abr-i-r} tienen su origen etimológico en la raíz verbal latina
\forma{aper-}, que sufrió la sonorización de la consonante sorda intervocálica.
Eso produjo la coexistencia entre los alomorfos \forma{aper-}, \forma{aber-} y
\forma{abr-} (este último con caída de la segunda vocal). Dicho de otro modo,
estas variantes se explican, en última instancia, por razones léxicas.

Sin embargo, la motivación para la alomorfía también puede provenir del propio
sistema gramatical. En este sentido, hay morfemas cuya manifestación se ve
condicionada por factores vinculados con el contexto morfológico en el que
aparecen. Por ejemplo, por los morfemas que lo rodean: \forma{am-a-ble},
\forma{am-a-bil-ísim-o}, \forma{am-a-bil-idad} y \forma{am-a-ble-mente}
despliegan dos variantes en alternancia para un mismo sufijo: \forma{-ble-} y
\forma{-bil-}. La motivación para la aparición de una u otra forma radica en que
\forma{-ble} se emplea en posición final de palabra, mientras que \forma{-bil-}
aparece cuando le sigue(n) otro(s) morfema(s).

La alomorfía también puede venir motivada por razones fonológicas. Por caso, en
muchas palabras del español las vocales medias (\forma{e}, \forma{o}) diptongan
cuando reciben acento léxico. Es lo que ocurre en el contraste entre
\forma{hel-a-r} y \forma{hiel-o}. De la misma manera, el participio
\forma{abier-t-o} exhibe una cuarta forma de la raíz ejemplificada
anteriormente: \forma{abier-}. Lo que causa la aparición de esta última variante
es, en este caso, un factor de naturaleza fonológica que abre la vocal cuando
esta recibe el acento.

\begin{ejercicio}\label{ej:04-alomorfos}
Considere el siguiente corpus, que recoge palabras extraídas del texto del
comienzo del capítulo y lo completa con otras. Todas ellas manifiestan
alomorfos.

\begin{consignas}
  \item Identifique los morfemas que presentan alomorfos y describa su
        significado.
  \item Explique la causa que motiva la aparición de una u otra variante.
\end{consignas}

\begin{tabular}{@{}llll@{}}
  \forma{cuesta}        & \forma{terraplanistas} & \forma{culturales}  & \forma{irrumpido} \\
  \forma{terráqueos}    & \forma{roto}           & \forma{datos}       & \forma{costosa} \\
  \forma{terrícola}     & \forma{entierran}      & \forma{falsedades}  & \forma{ruptura} \\
  \forma{jóvenes}       & \forma{territorio}     & \forma{irrupción}   & \forma{disruptiva} \\
  \forma{maniqueas}     & \forma{aterrizar}      & \forma{rompió}      & \forma{expertos} \\
  \forma{Tierra}        & \forma{(las) crisis}   & \forma{costar}      & \forma{cuesten} \\
\end{tabular}

\begin{resuelto}
a) \forma{cuest-a, cost-a-r, cuest-e-n, cost-os-a}. En las palabras de este
grupo se identifica el morfema \textsc{cost-}, que significa \sig{precio o
esfuerzo demandado}. Los alomorfos son \fonol{ˈkwest} \textasciitilde{}
\fonol{kost}.\footnote{El signo \textasciitilde{} se lee \sig{alterna con}.} b)
La causa de la alomorfía es fonológica: \fonol{ˈkwest} aparece en sílaba tónica
y \fonol{kost}, en sílaba átona.
\end{resuelto}
\end{ejercicio}

\subsection{Asimilación y disimilación}

Aunque cada nivel lingüístico posee sus unidades y operaciones características,
no son pocos los fenómenos en los que estos entran en contacto. Algunos casos de
alomorfía fonológicamente motivada resultan de procesos fonológicos, como los de
asimilación y disimilación. El primero, que ya se describió en el \capref{3},
§\ref{sec:03-procesos}, consiste en que dos segmentos fónicos se vuelven
articulatoriamente más parecidos entre sí cuando entran en contacto. El caso más
típico es el de las consonantes nasales que, cuando preceden a otros sonidos
consonánticos, adaptan su punto de articulación al del segmento que les
sigue.\footnote{Es lo que se conoce como asimilación regresiva y es el tipo de
asimilación más frecuente en las lenguas (v. \capref{3},
§\ref{sec:03-procesos}).} De esto resulta que morfemas como el prefijo negativo
\forma{in-}, cuando se añaden a palabras iniciadas en consonante bilabial, se
realicen con una nasal bilabial: \forma{im-batible} \fon{imbaˈtiβle}, mientras
que en \forma{in-fum-a-ble} \fon{iɱfuˈmaβle}, \forma{in-toler-a-ble}
\fon{intoleˈɾaβle} e \forma{in-com-i-ble} \fon{iŋkoˈmiβle}, aunque la ortografía
no cambia e invisibiliza el fenómeno, también se pronuncian alófonos distintos
del archifonema nasal (v. \capref{3}, §\ref{sec:03-procesos}). El panorama de
posibilidades se completa ante palabras cuya primera consonante es un sonido
líquido: en este caso, el sonido nasal del prefijo cae y no se produce:
\forma{i-legal}, \forma{i-rracional}.\pendiente{El original dice «sonido
lateral», pero incluye \textit{irracional}}

Mientras que en la asimilación los sonidos se tornan formalmente más semejantes,
en la disimilación ocurre el proceso inverso: un segmento fonológico cambia su
forma habitual para diferenciarse de otro cercano. Las consonantes líquidas
suelen participar de este fenómeno. Eso explica por qué \forma{invern-al} y
\forma{circul-ar}, que contienen el mismo sufijo derivativo que expresa un valor
relacional (respectivamente \sig{relativo al invierno} y \sig{relativo al
círculo}), manifiestan, en el primer caso, un alomorfo terminado en lateral
(\forma{-al}) y, en el segundo, uno que finaliza en vibrante (\forma{-ar}) para
evitar la reiteración del segmento lateral: \juicio{*}{\forma{circulal}},
pronunciación que resultaría cacofónica en español.

\section{Clases de morfemas}

Los morfemas se pueden agrupar en clases según dos criterios principales: por
sus propiedades distribucionales y por su funcionamiento en el interior de la
palabra.

\subsection{Morfemas libres y ligados}

Desde un punto de vista distribucional, los morfemas se dividen en dos grandes
grupos: los \termino{morfemas libres}, que pueden aparecer solos ---o sea, tienen
capacidad para formar por sí solos una palabra--- y los \termino{morfemas
ligados}, que siempre aparecen en concurrencia con otros. Así, decimos que el
segmento \forma{-mar-}, que integra la palabra \forma{sub-mar-in-a-s}, pertenece
al primer grupo, dado que, como vimos, por sí mismo constituye una forma léxica
que puede funcionar dentro de una oración: \forma{Caminé hacia el mar}; en
cambio, ni el sufijo \forma{-in-} ni el prefijo \forma{sub-} ofrecen esta
posibilidad: su aparición siempre depende de que se encuentren junto con otros
segmentos morfológicos, como en \forma{bov-ino} o \forma{sub-desarrollado}.

\subsection{Raíces y afijos}

Más allá de su distribución, no todas las unidades morfológicas tienen la misma
función en la palabra. Las \termino{raíces}, que son las responsables del
significado léxico fundamental de la palabra, pueden ser libres, como
\forma{mar, árbol} o \forma{sol}, o ligadas, como \forma{convenc-} (en
\forma{convencer}) o \forma{chat-} (en \forma{achatado}). La mayor parte de las
raíces del español son ligadas y solo constituyen formas libres las que
integran palabras monomorfemáticas, como las de clase cerrada (categorías
funcionales, preposiciones, adverbios no terminados en \forma{-mente}) y un
puñado de adjetivos y nombres terminados en consonante, como \forma{árbol},
\forma{señor} o \forma{reloj}. Los miembros de las llamadas «familias de
palabras» comparten la raíz. Por ejemplo, \forma{submarino} se agrupa por medio
de este criterio con las ya mencionadas \forma{marino, bajamar, pleamar,
altamar} y también con \forma{marea, marear, marítimo, marinero, ultramarino,
marejada, maricultura, mareógrafo}, entre otras.

Los \termino{afijos} son unidades siempre ligadas entre las que cabe
diferenciar, por un lado, los derivativos, que aportan información léxica y
permiten formar palabras nuevas (como \forma{-in-} en \forma{mar-in-o} o
\forma{anti-} en \forma{anti-hongo-s}) y, por otro lado, los flexivos, que son
los que brindan información gramatical vinculada a las categorías lingüísticas
(como \forma{-o-} en \forma{marin-o}, que expresa género \sig{masculino}, y
\forma{-s} en \forma{antihongo-s}, que manifiesta \sig{plural}). Dependiendo de
su posición relativa respecto de la raíz, los afijos se dividen en prefijos y
sufijos. En español observamos que los flexivos se realizan exclusivamente como
sufijos, dado que nuestra lengua no posee prefijos
flexivos.\pendiente{El original dice «los derivativos se realizan exclusivamente
como sufijos»} En cuanto al orden relativo, en una palabra como
\forma{gigantesca} podemos observar que la raíz \forma{gigant-} primero se une
al sufijo derivativo \forma{-esc-} para formar el adjetivo \forma{gigantesc-}; y
luego entra en juego la morfología flexiva que incorpora el sufijo de género
femenino \forma{-a-}. Y, si quisiéramos pluralizar la palabra, el sufijo flexivo
plural \forma{-s} se agregaría después del de género. La secuencia que acabamos
de ejemplificar no es casual ni aislada: los sufijos flexivos son siempre más
externos que los derivativos en relación con la raíz, ya que, como demostró
\textcite{bybee1985}, los derivativos afectan el significado de la raíz,
mientras que los flexivos inciden en el comportamiento sintáctico de la palabra.
Asimismo, los elementos flexivos no solo son los últimos en incorporarse a la
palabra, sino que, además, lo hacen en un orden determinado. Por ejemplo, en una
forma verbal como \forma{estudi-a-ba-n}, a la raíz \forma{estudi-} le sigue la
vocal temática \forma{-a-}, a continuación se añade la manifestación amalgamada
de modo indicativo, tiempo pretérito y aspecto imperfectivo, manifestada por el
sufijo \forma{-ba-} y, finalmente, se une otra amalgama, que indica
conjuntamente tercera persona y número plural: \forma{-n}
(v.~§\ref{sec:04-variantes}).

Más allá de su evidente diferencia posicional, prefijos y sufijos poseen
características de funcionamiento distintas. Así, mientras que los primeros no
suelen ser selectivos con relación a la clase de la palabra a la que se
incorporan, los segundos son sensibles a la categoría léxica y, en general, la
modifican. Como ya hemos comentado, \forma{marino} se obtiene de \forma{mar} más
el sufijo \forma{-in}, que forma adjetivos relacionales a partir de nombres
(\forma{salino} < \forma{sal}, \forma{canino} < \forma{can}).\footnote{El signo <
se lee \sig{proviene de}.} Por contraste, el prefijo \forma{sub-} produce un
nombre a partir de otro en \forma{subgrupo} (< \forma{grupo}) y
\forma{subsistema} (< \forma{sistema}), pero también verbos derivados a partir
de verbos, como en \forma{subestimar} (< \forma{estimar}) y \forma{subarrendar}
(< \forma{arrendar}). Fonológicamente, el nivel de integración que exhiben con
la palabra también varía: suele ser mayor en el caso de los sufijos que en el de
los prefijos, entre los que existe incluso un grupo de formas separables, con
acento léxico propio, como \forma{ex-} en \forma{ex primer ministro} o
\forma{re-} en \forma{re me cae bien}.\footnote{Según la \textit{Ortografía de la
lengua española} \parencite[cap.~V, §\,2.2.2]{rae2010}, los prefijos deben
siempre escribirse unidos a la base, excepto, como en los casos de arriba,
cuando se coloquen delante de varias palabras.}

\begin{ejercicio}\label{ej:04-clases}
\textbf{Parte A}
\begin{consignas}
  \item Complete la siguiente tabla, clasificando los morfemas segmentados en el
        ejercicio~\ref{ej:04-terraplanista} a partir de su funcionamiento en la
        palabra.

        {\footnotesize
        \begin{tabular}{@{}l*{5}{c}@{}}
          \toprule
          Palabra & \makecell{Pref.\\flexivo} & \makecell{Pref.\\derivativo} & Raíz & \makecell{Suf.\\derivativo} & \makecell{Suf.\\flexivo} \\
          \midrule
          \forma{convencer}      & & con- & -venc- & -e- & -r \\
          \forma{falsedades}     & & & fals- & -e-dad- & -es \\
          \forma{terraplanistas} & & & & & \\
          \ldots & & & & & \\
          \bottomrule
        \end{tabular}}
  \item Identifique cuáles de los morfemas segmentados pueden clasificarse como
        libres y cuáles solo pueden funcionar como ligados.
\end{consignas}

\textbf{Parte B}

A partir de lo observado, decida cuáles de las siguientes generalizaciones
describen con mayor precisión para el español el funcionamiento de los morfemas
analizados.

\begin{consignas}
  \item Respecto de la relación entre las clases distribucionales y los tipos de
        morfemas:
        \begin{exe}
          \ex Las raíces son siempre morfemas libres, pero no todos los
              morfemas libres son siempre raíces.
          \ex Tanto las raíces como los afijos son siempre morfemas ligados.
          \ex Los afijos pueden ser morfemas libres o morfemas ligados.
          \ex Los afijos son siempre morfemas ligados, pero no todos los
              morfemas ligados son siempre afijos.
        \end{exe}
        \begin{resuelto}
        Opción correcta: 4.
        \end{resuelto}
  \item Respecto del aporte gramatical que realizan los sufijos y los prefijos
        del español a la construcción a la que se integran:
        \begin{exe}
          \ex Tanto los prefijos como los sufijos pueden expresar valores
              flexivos.
          \ex Los sufijos siempre forman palabras nuevas, pero no todas las
              palabras nuevas se forman por medio de la sufijación.
          \ex Solo los sufijos pueden expresar valores flexivos.
          \ex Solo los prefijos permiten formar palabras nuevas.
        \end{exe}
        Opción correcta:\completar[1]
\end{consignas}
\end{ejercicio}

\subsection{Otros afijos}

Existen otras formas de afijos que aparecen con menor frecuencia en la
morfología del español. Por un lado, hay formantes que se ubican entre raíces y
prefijos o entre raíces y sufijos (de los que ya hemos hablado en
§\ref{sec:04-morfema}), cuya función fundamental es actuar como enlaces
fonológicos, por lo que no portan significado. Es lo que ocurre con el segmento
\forma{-c-}, que se inmiscuye entre la raíz \forma{café-} y el sufijo
\forma{-it-} en \forma{cafe-c-ito}. Del mismo modo funciona el segmento
\forma{-t-} en \forma{café-t-ería} ante el sufijo \forma{-ería} y ante
\forma{-iv-} en \forma{representa-t-ivos}. A estos morfemas se los conoce como
\termino{interfijos}.

Por otro lado, existen casos muy marginales para el español en los que un
morfema puede partir la raíz e insertarse en su interior. A estos segmentos se
los denomina \termino{infijos}, y se limitan, en nuestra lengua, a la expresión
de algunos diminutivos, como en \forma{azuqu-ít-ar} y \forma{lej-it-os}.

Finalmente, existen afijos que constituyen morfemas discontinuos. Los
\termino{circunfijos}, por ejemplo, son formantes que se aplican rodeando la
palabra a la que se añaden. Puede pensarse que \forma{tras-papel-a-r} se
encuentra formado así, ya que ni \forma{tras-} ni \forma{-a-} se han incorporado
a la raíz \forma{papel} por separado, pues no se registran ni
\juicio{*}{\forma{traspapel}} ni \juicio{*}{\forma{papelar}}. Sin embargo, se
trata de formantes que funcionan de manera independiente en otras palabras,
como \forma{tras-nacion-al} o \forma{sal-a-r}, por lo que cabe mejor pensarlos
como prefijos y sufijos que se incorporan en un mismo paso en el proceso de
formación de la palabra. Analizaremos estos casos con más detalle más adelante,
al explicar el proceso de formación de palabras conocido como parasíntesis
(v.~§\ref{sec:04-parasintesis}).

\subsection{El concepto de \textit{tema}}\label{sec:04-tema}

Dentro de cada elemento léxico, la unidad a partir de la cual se forman las
variantes flexivas se denomina \termino{tema}. El tema es la unidad que codifica
el contenido léxico de la palabra y es común a todas sus formas flexionadas. Si
retomamos nuestro ejemplo de §\ref{sec:04-segmentacion}, el segmento
\forma{marin-} allí considerado constituye, precisamente, el tema de la palabra
\forma{marinos}, ya que es a \forma{marin-} que se unen los sufijos flexivos
\forma{-o-} de género y \forma{-s} de plural. Asimismo, el tema \forma{norma-}
es la unidad que permite formar tanto el singular \forma{norma} como el plural
\forma{norma-s}. A partir del tema \forma{norma-} se puede derivar
\forma{normal-} (\sig{relativo a la norma}) y, a partir del nuevo tema
\forma{normal-}, se producen tanto el plural \forma{normal-es} como el derivado
adverbial \forma{norm-al-mente}, mediante el agregado de un sufijo flexivo de
plural, en el primer caso, y de uno de derivación, en el segundo. Estos procesos
de derivación pueden representarse mediante una descripción estructural en la
que cada corchete encierra un paso en la formación de la palabra e indica la
clase de palabra resultante, como se muestra en \exref{ex:04-normalmente}:

\begin{exe}
  \ex\label{ex:04-normalmente} Descripción estructural (DE):
      [[[norma]\textsubscript{N} al]\textsubscript{A} mente]\textsubscript{Adv}\footnote{Las
      abreviaturas remiten a las clases de palabras resultantes de cada etapa:
      N: nombre; A: adjetivo; Adv: adverbio; V: verbo.}
\end{exe}

Entre los morfemas que integran el tema de la palabra existen relaciones
jerárquicas, lo que significa que expresan una dependencia de tipo estructural.
Esto implica, entre otras cosas, que en la forma de la palabra queda un
testimonio de los pasos que permitieron constituirla. Se considera que cada uno
de estos pasos debe conformar una palabra existente en la lengua (aunque más
adelante veremos que en algunos procesos este requisito no se mantiene,
v.~§\ref{sec:04-parasintesis}), que puede, a su vez, ser flexionada de acuerdo
con las posibilidades de su categoría.

Operativamente se puede decir que el tema resulta de quitarle a la palabra todos
sus afijos flexivos. En consecuencia, los afijos flexivos, que en español se
realizan exclusivamente como sufijos, no integran el tema. Esto no implica, no
obstante, que no pueda haber afijos derivativos en el tema, tanto sufijos como
prefijos. Por tanto, dado que la flexión es un proceso siempre posterior a la
derivación, cada paso en la historia derivacional de la palabra coincide con un
tema, y solo en la última etapa morfológica ---concluido ya el ciclo de
formación de la palabra--- se pueden adjuntar morfemas flexivos. Como se muestra
en la descripción estructural \exref{ex:04-marinos}, al sustantivo \forma{mar}
se le une el sufijo \forma{-in-} para formar un adjetivo y a este tema se le
añaden los elementos flexivos, que no involucran un paso más (otro corchete)
porque su agregado no involucra la formación de una palabra distinta.

\begin{exe}
  \ex\label{ex:04-marinos} DE: [[mar]\textsubscript{N} in]\textsubscript{A}-o-s
\end{exe}

\begin{ejercicio}\label{ej:04-tema}
\textbf{Parte A}

A partir de las palabras \forma{achatada} y \forma{pseudocientíficos}, que
fueron segmentadas en el ejercicio~\ref{ej:04-terraplanista}, e incorporando al
corpus \forma{entronización, culturales, bromear} y \forma{desprecio}, tomadas
del mismo texto, realice las actividades que se indican a continuación.

\begin{consignas}
  \item Segmente los ejemplos que no fueron analizados previamente en morfemas y
        proporcione una paráfrasis del significado de cada palabra.
  \item Indique los pasos mediante los que se ha ido formando la palabra y a qué
        categoría pertenece el tema resultante en cada uno de ellos.
  \item Realice la descripción estructural (DE) de la palabra.
\end{consignas}

\begin{resuelto}
\forma{entronización}: a) \forma{en-tron-iz-a-ción} = \sig{acción o efecto de
poner algo o a alguien en un trono}. b) \forma{trono}\textsubscript{N} >
\forma{entronizar}\textsubscript{V} > \forma{entronización}\textsubscript{N}; N >
V > N\textsubscript{abstracto}. c) [[en-[tron]\textsubscript{N}-iz-a]\textsubscript{V}-ción]\textsubscript{N}.
\end{resuelto}

\textbf{Parte B}

Decida cuáles de las siguientes afirmaciones describen con mayor precisión lo
observado en el análisis estructural realizado.

\begin{consignas}
  \item Con respecto al tema:
        \begin{exe}
          \ex Operativamente, el tema resulta de quitarle a la palabra todos sus
              afijos flexivos.
          \ex Operativamente, el tema resulta de quitarle a la palabra todos sus
              sufijos.
          \ex Operativamente, el tema resulta de quitarle a la palabra todos los
              afijos.
          \ex Operativamente, el tema coincide siempre con la raíz de la
              palabra.
        \end{exe}
        \begin{resuelto}
        Opción correcta: 1.
        \end{resuelto}
  \item Con respecto al proceso derivacional:
        \begin{exe}
          \ex Cada paso del proceso derivacional coincide con la incorporación
              de un morfema.
          \ex Cada paso del proceso derivacional coincide con la incorporación
              de un morfema derivativo.
          \ex Cada paso del proceso derivacional coincide con el tema de una
              palabra existente.
          \ex Cada paso del proceso derivacional coincide con una palabra
              flexionada.
        \end{exe}
        Opción correcta:\completar[1]
\end{consignas}
\end{ejercicio}

\subsection{Los temas cultos o semipalabras}\label{sec:04-temascultos}

Existe un segundo sentido según el cual se utiliza el término \forma{tema} en
español, que se diferencia del que hemos utilizado hasta ahora. Designa a una
clase de morfemas asociados a contenidos léxicos, los cuales por sí solos no
pueden recibir morfología flexiva ni funcionar como palabras, por lo que se los
conoce como \termino{semipalabras}, aunque también, por su origen en las lenguas
clásicas, se los llama \termino{temas cultos}. A diferencia de las palabras, que
siempre pertenecen a una clase determinada (N, V, A, etc.), los temas cultos no
se categorizan hasta que, para formar una palabra, se unen con otro elemento de
la lengua. Ese segundo formante puede ser un afijo derivativo (\forma{a-fonía},
\forma{fon-iatra}), otra palabra (\forma{fono-terapia}) u otro tema culto
(\forma{fono-logía}), como veremos en §\ref{sec:04-composicion}.

\begin{recuadro}[Para ampliar]
Aunque las gramáticas clásicas los asimilaban a los prefijos, los temas cultos
no siempre van al inicio de la palabra, como sí lo hacen los prefijos. En
realidad, los temas cultos no tienen posición fija en la palabra, en la que
pueden estar al principio o al final, como lo muestran pares como
\forma{fono-logía / telé-fono}, \forma{podo-logía / cefaló-podo},
\forma{logo-pedia / astró-logo}.
\end{recuadro}

\subsection{La vocal temática}\label{sec:04-vt}

Como ya adelantamos en §\ref{sec:04-morfema}, en algunas formas flexivas de los
verbos del español se advierte la presencia de un segmento que carece de
contenido léxico, pero se caracteriza por indicar la clase según la cual se
conjuga el verbo. Se trata de la \termino{vocal temática}, que aparece, por
ejemplo, precediendo a la flexión de infinitivo, donde contrastan las tres
conjugaciones: \forma{-a-} en \forma{amar}, \forma{-e-} en \forma{temer},
\forma{-i-} en \forma{partir}. Este morfema se manifiesta aislado o fusionado con
otros rasgos flexivos. Resulta segmentable en algunas formas, como la primera
persona del plural del pretérito imperfecto del indicativo
(\forma{am-á-ba-mos, tem-í-a-mos, part-í-a-mos});\footnote{Como puede apreciarse
en estos ejemplos, la vocal temática de la segunda y la tercera conjugación se
manifiesta por medio del mismo morfo.} en cambio, en otras, como la primera
persona del singular del presente del indicativo (\forma{am-o, tem-o, part-o}),
no tiene manifestación independiente. Estos morfemas, que contienen más de un
significado, se denominan \textit{portmanteau} (francés: \sig{maleta}) o
\termino{amalgama}. Una tercera posibilidad es que la vocal temática sea el
único elemento que acompaña al tema. Es un análisis posible para lo que ocurre
en los imperativos de segunda persona del singular, en los que las otras
categorías flexivas (tiempo, modo, aspecto, número y persona) no se manifiestan
(\forma{am-á, tem-é, part-í}).

Como se desprende de lo anterior, dado que señala la clase conjugacional del
verbo, la vocal temática permite predecir cómo se obtienen determinadas
variantes flexivas, como, por ejemplo, los pretéritos imperfectos, ya que la
clase conjugacional del verbo explica por qué no producimos formas como
\juicio{*}{\forma{amíamos}}, \juicio{*}{\forma{temábamos}} o
\juicio{*}{\forma{partábamos}}.

\begin{recuadro}[Para ampliar: la VT en los nombres]
El análisis anterior se puede extender al ámbito nominal, ya que en sustantivos
y adjetivos la formación del plural también se determina por la naturaleza de
la terminación del tema, según los siguientes casos:
\begin{itemize}
  \item Tema finalizado en consonante, salvo que esté seguido por un sufijo
        flexivo de género: el plural se forma con la variante \forma{-es}:
        \forma{ventanal-es} (y no \juicio{*}{\forma{ventanals}}).
  \item Tema finalizado en \fonol{s} (palabras graves y esdrújulas): las formas
        singulares y plurales son homónimas: \forma{la crisis / las crisis}.
  \item Tema finalizado en vocal o seguido por la marca de género: la flexión se
        realiza con \forma{-s} (\forma{casa-s, gat-a-s}, pero no
        \juicio{*}{\forma{casa-es}}, \juicio{*}{\forma{gat-a-es}}).
  \item Tema finalizado en vocales cerradas acentuadas: forma el plural como los
        temas finalizados en consonante, \forma{ñandú-es} y \forma{gurí-es}. No
        obstante, para estas terminaciones algunas variedades pueden admitir el
        alomorfo \forma{-s} (\forma{ñandú-s, gurí-s}) o incluso la variante
        \forma{-ses} (\forma{ñandú-ses, gurí-ses}).
\end{itemize}
\end{recuadro}

\section{Homonimia estructural en el nivel morfológico}\label{sec:04-homonimia}

Entre los morfemas que integran el tema existen relaciones de dependencia, por
lo que también algunas variantes interpretativas de las palabras pueden
explicarse como casos de homonimia estructural. Esto quiere decir que su
estructura se puede describir por medio de (al menos) dos análisis
alternativos. Por ejemplo, la palabra \forma{in-mov-il-iz-a-ble} puede
significar \sig{que se puede inmovilizar} o \sig{que no se puede movilizar}. La
interpretación que realizamos depende de en qué nivel de la formación de la
palabra consideramos que se adjunta el prefijo, si al tema \forma{móvil} o a
\forma{movilizable}, lo cual, según vemos, no es trivial, ya que cambia el
significado con que interpretamos la palabra. En la primera estructura, el tema
del verbo \forma{inmovilizar} se forma a partir del adjetivo \forma{inmóvil} y la
incorporación del sufijo \forma{-ble} produce el significado \sig{que se puede
inmovilizar}. En la segunda interpretación, \forma{in-} se une al ya formado
adjetivo \forma{movilizable} y genera la lectura \sig{no movilizable}, es decir,
\sig{que no se puede movilizar}. En definitiva, los segmentos morfológicos que
intervienen son los mismos, pero la trama de relaciones entre ellos varía y, en
consecuencia, también lo hace la interpretación.

\begin{ejercicio}\label{ej:04-homonimia}
Partiendo de que las siguientes palabras tienen por lo menos dos acepciones cuyo
origen puede explicarse estructuralmente, realice las actividades que se
solicitan a continuación:

\begin{consignas}
  \item Deduzca a partir de los ejemplos presentados las dos interpretaciones
        alternativas asociadas a cada una de estas palabras y realice una
        paráfrasis de cada uno de los significados propuestos.
  \item Presente los pasos de la derivación y ofrezca descripciones
        estructurales que den cuenta de los significados asignados a cada
        palabra.
\end{consignas}

\begin{exe}
  \ex \textbf{reconsiderado}
    \begin{xlist}
      \ex Ni siquiera me avisó que no iba. Jaja re piola, la mina,
          reconsiderada. Así estamos, hay gente para todo.
          (\url{https://www.facebook.com/karen.alma.90/posts/705446233192085})
      \ex Conicet: 47 aspirantes fueron reconsiderados e ingresarán a la carrera
          de investigador.
          (\url{https://www.clarin.com/sociedad/conicet-47-aspirantes-reconsiderados-ingresaran-carrera-investigador_0_r1NziPGxb.htm})
    \end{xlist}
  \ex \textbf{inutilizable}
    \begin{xlist}
      \ex Los celulares robados quedarán bloqueados y serán inutilizables.
          (\url{https://merloenlinea.com.ar/los-celulares-robados-quedaran-bloqueados-y-seran-inutilizables/})
      \ex El dirigible contratado por la empresa «La Serenísima» para fines
          publicitarios quedó inutilizable luego de un temporal en la localidad
          de Campo Andino, a 30 kilómetros de esta capital.
          (\url{https://www.lanacion.com.ar/sociedad/ahora-le-dicen-pinguino-porque-antes-volaba-nid170176})
    \end{xlist}
  \ex \textbf{desarmar}
    \begin{xlist}
      \ex Vasallo viene desarmando el discurso de la monogamia desde el 2013 y
          reflexiona sobre qué significa migrar de ese régimen.
          (\url{https://www.pagina12.com.ar/202982-los-malabares-del-deseo})
      \ex Al menos cinco ladrones rodearon y desarmaron este miércoles a
          Leonardo Sarmiento, un policía que trabajaba como encubierto en Villa
          Tranquila, Avellaneda\ldots{}
          (\url{https://www.infobae.com/sociedad/2018/11/02/menores-con-armas-de-juguete-amenazaron-a-una-periodista-en-vivo-al-policia-lo-queme-yo})
    \end{xlist}
\end{exe}

\begin{resuelto}
1. a. \forma{reconsiderado}: \sig{muy considerado o amable}.
\forma{considerar}\textsubscript{V} >
\forma{considerado}\textsubscript{Part/A}\footnote{El participio del verbo llega
a funcionar como adjetivo por un proceso derivativo conocido como conversión,
que estudiaremos en §\ref{sec:04-conversion}.} >
\forma{reconsiderado}\textsubscript{A};
[re-[[considera]\textsubscript{V}-d-o]\textsubscript{Part/A}]\textsubscript{A}; V >
Part/A > A. b. \sig{que se ha vuelto a considerar}:
\forma{considerar}\textsubscript{V} > \forma{reconsiderar}\textsubscript{V} >
\forma{reconsiderado}\textsubscript{V};
[[re-[considera]\textsubscript{V}]\textsubscript{V}-d-o]; V > V.
\textit{Aclaración}: en a., cuando el prefijo \forma{re-} se une a un A se
interpreta como intensivo (\sig{muy}); en b., cuando se une a un V, indica
repetición (\sig{volver a}).
\end{resuelto}
\end{ejercicio}

\section{Los procesos morfológicos}

Los morfemas de la lengua no se combinan unos con otros en forma aleatoria, sino
que lo hacen de manera sistemática y en un orden determinado. Cada uno de ellos
se incorpora a la palabra por medio de un proceso morfológico. Estos procesos se
dividen en dos grandes grupos: los \termino{de formación de palabras o
derivativos} y los \termino{flexionales}. Mientras que los primeros permiten
obtener palabras nuevas a partir de otras conocidas por el hablante, es decir,
crear nuevos \termino{lexemas}, por medio de los segundos se producen las
distintas \termino{variantes gramaticales} de un mismo lexema o unidad
vocabular de la lengua. Volviendo al ejemplo introducido en
§\ref{sec:04-segmentacion}, \forma{submarino}, \forma{submarina},
\forma{submarinos} y \forma{submarinas} son variantes gramaticales de un mismo
lexema adjetival, \forma{submarino}. Del mismo modo, \forma{mir-o, mir-á-s,
mir-a-ba-n, mir-a-r-ía-mos}, que cambian en los rasgos flexivos de tiempo,
persona y número, no son palabras distintas, sino formas que integran el
paradigma flexivo del mismo verbo \forma{mir-a-r}. Por contraste,
\forma{mir-a-dor}, que se forma añadiendo al verbo el sufijo derivativo
nominalizador \forma{-dor}, no es una variante del mismo lexema, sino otra
palabra diferente, que admite su propia variante flexiva plural
\forma{mir-a-dor-es}.

\subsection{La base}\label{sec:04-base}

La unidad sobre la cual se aplica un proceso morfológico se denomina
\termino{base}. Este concepto es relacional, puesto que la forma final de una
palabra puede obtenerse a partir de la aplicación ordenada de varios procesos
sucesivos. Retomando el ejemplo anterior, la base a partir de la cual se obtiene
\forma{mir-a-dor-es} es el tema del sustantivo \forma{mir-a-dor}, al cual se le
incorpora el sufijo flexivo \forma{-es}, que marca plural. Pero
\forma{mir-a-dor-} tiene como base, a su vez, el tema verbal \forma{mira-}, que
se convierte en nominal al adicionar el sufijo derivativo \forma{-dor}, que
forma un nuevo lexema.

\begin{exe}
  \ex\label{ex:04-mirador} \forma{mira(r)} > \forma{mira-dor} > \forma{mirador-es}
\end{exe}

\subsection{Productividad y transparencia}

Independientemente del hecho de que ambos grupos de procesos comparten el
recurso de la sufijación, poseen características bien diferenciadas, que
justifican una división del nivel morfológico en los dos ámbitos que
introdujimos en §\ref{sec:04-ramas}: la morfología léxica y la morfología
flexiva. Esta última se encuentra más próxima a la sintaxis y, por ello,
involucra procesos automáticos y obligatorios, es decir, que se pueden aplicar
sobre todas las palabras que poseen una determinada propiedad y no son
facultativos, sino que vienen determinados por los requisitos de buena formación
sintáctica de la oración en la que la palabra aparece. Así, todos los adjetivos
forman morfológicamente variantes de plural para concordar con los nombres a
los que modifican. La formación de palabras, en cambio, involucra procesos
optativos con distintos grados de productividad, que pueden no tener aplicación
sobre algunas bases. Por ejemplo, el sufijo \forma{-dad} produce sustantivos
abstractos como \forma{mal-dad, capac-i-dad, real-i-dad} a partir de adjetivos
como \forma{mal-o/a, capaz, real}, pero no se registra entre los hablantes una
forma como \juicio{*}{\forma{ternedad}} obtenida a partir de \forma{tiern-o/a}.
En cambio, se registra \forma{ternura}, obtenida de la incorporación del sufijo
\forma{-ura}, cercano semánticamente a \forma{-dad}, aunque menos productivo,
que también forma \forma{blancura}, \forma{dulzura} y \forma{verdura}. Podemos
decir, entonces, que la formación de un derivado abstracto en \forma{-ura}, en
el caso del adjetivo \forma{tierno}, bloquea o impide la formación de un
derivado abstracto en \forma{-dad}, que tendría el mismo significado.

Asimismo, por estar más cercana al léxico, esta área de la morfología tiende a
ser menos transparente semánticamente que la morfología flexiva y permite
interpretaciones más libres. Por caso, el significado de \forma{verdura}, en el
ejemplo anterior, no se interpreta habitualmente según su estructura
morfológica (\sig{cualidad de verde}), sino más bien como la designación de una
clase de alimentos, en muchos de los cuales predomina ese color. Ha ocurrido una
lexicalización por la que el significado de la palabra se ha vuelto más opaco.

\subsection{Un caso fronterizo: la morfología apreciativa}\label{sec:04-apreciativa}

La \termino{morfología apreciativa} constituye un caso especial, fronterizo entre
la derivación y la flexión. En español, comparte con esos dos procesos el
procedimiento de la sufijación, si bien, como señala \textcite{kornfeld2021},
también se expresa por medio de prefijos, como ocurre en usos ponderativos de
morfemas como \forma{re-} y otros análogos: \forma{re-leal, réquete-leal}. Al
igual que sucede con la morfología flexiva, la productividad de estos procesos
es muy alta (aunque no llega a ser automática). Por caso, casi todos los nombres
pueden formar diminutivos en \forma{-it-}, como \forma{cas-ita, airec-ito,
bosquec-ito}. Como se desprende de los ejemplos anteriores, los sufijos
apreciativos son selectivos respecto de la categoría de sus bases, pero no la
cambian, ni forman tampoco palabras nuevas al incorporarse. Presentan la
particularidad de que pueden repetirse, como en \forma{chiqu-it-it-o}. Sin
embargo, en consonancia con la formación de palabras, su aplicación no es
forzosa, no se encuentra determinada por la estructura sintáctica de la
oración, ya que estos sufijos no expresan información gramatical obligatoria. No
obstante, tampoco aportan significado léxico específico. Como se plantea en
Albano y Giammatteo (2000, p.~5),\pendiente{En la bibliografía figura Albano y
Giammatteo (2002)}

\begin{cita}
Por lo general, la semántica de los apreciativos se presenta relacionada con
parámetros tales como \sig{dimensión}, \sig{ponderación} y \sig{afectividad},
todos con la posibilidad de manifestar polaridad negativa tanto como positiva:
aumento/disminución (tamaño), valor/desvalor (ponderación), aprecio/desprecio
(afectividad).
\end{cita}

Es decir, los apreciativos se caracterizan por portar valores connotativos, que
codifican significados expresivos que quiere manifestar el hablante, lo cual
generalmente está muy vinculado al contexto en que se emiten tales expresiones.
Así, \textcite[128]{lang1997} reconoce lo siguiente:

\begin{cita}
\textit{-ito}, categorizado aisladamente como diminutivo y afectivo, puede ser
utilizado como peyorativo y despectivo en \textit{caudillo} >
\textit{caudillito}, o incluso como aumentativo en una construcción del tipo
\textit{comida calentita} (= comida muy caliente); \textit{-ote}, con claro
valor aumentativo, aparece en algunos contextos léxicos con un claro valor
diminutivo (cfr. \textit{isla} > \textit{islote}). La determinación del
significado concreto de los sufijos apreciativos en cada uno de los contextos
en que aparecen puede depender de factores externos a los sufijos mismos. El
valor peyorativo de \textit{caudillito} deriva del sentido negativo que se
aplica a la base \textit{caudillo}. O la intención del hablante, en función del
énfasis o el tono de voz, puede alterar el significado exacto ---\textit{¡qué
españolaza!} podría interpretarse como ponderación (\textit{¡qué españolaza más
guapa, más impresionante!}) o como desprecio (\textit{¡qué españolaza más
gorda, más fea!})---. Resultaría precipitado categorizar estos sufijos fuera de
su contexto.
\end{cita}

Distribucionalmente estos sufijos ocupan una posición intermedia entre la
morfología flexiva, que es más externa, y la léxica, que es más cercana a la
raíz, tal como se muestra en \forma{cas-er-it-o-s}, donde el sufijo apreciativo
\forma{-it-} se ubica entre el derivativo \forma{-er-}, que lo precede, y los
sufijos de flexión que indican género masculino \forma{-o} y número plural
\forma{-s}.\footnote{En algún sentido, esta generalización también parece
mantenerse para la prefijación, dado que, si bien en español no existen
prefijos flexivos, aquellos que participan del subsistema apreciativo resultan
más externos que los que aportan contenido léxico. Así, resulta aceptable
\textit{re-des-leal}, pero no *\textit{des-re-leal}.}

\begin{ejercicio}[El sufijo \textit{-ardo}]\label{ej:04-ardo}
Considere las siguientes imágenes, que incluyen palabras que han sido obtenidas
por un proceso que se popularizó en el léxico juvenil de las redes sociales en
los últimos años, y luego resuelva las consignas que se detallan más abajo:

\imagenpendiente{Imágenes del ejercicio: \textit{cafesardo, malardo, buenardo,
temardo, matesardo, quesardo}}

\textbf{Parte A}
\begin{consignas}
  \item Describa el proceso empleado en la formación de las palabras con el
        sufijo \forma{-ardo}, indicando la categoría de las bases y la del
        resultado.
  \item Realice una paráfrasis del significado del morfema en cuestión.
  \item Expanda el corpus con otros ejemplos de uso del mismo sufijo.
\end{consignas}

\begin{resuelto}[Parte A. Punto resuelto]
a) N\textsubscript{x}/A\textsubscript{y} + \forma{-ard-} >
N\textsubscript{x}/A\textsubscript{y}. (Los índices \textit{x} e \textit{y}
señalan que las bases no cambian su categoría).
\end{resuelto}

\textbf{Parte B}

Evalúe cuáles de las siguientes afirmaciones generalizan con mayor precisión las
conclusiones del análisis del ejercicio anterior.

\begin{consignas}
  \item Desde el punto de vista de la productividad del sufijo:
        \begin{exe}
          \ex El morfema se puede aplicar sobre cualquier base de una categoría,
              como los sufijos flexivos.
          \ex El morfema se puede aplicar sobre un conjunto acotado de bases,
              como los sufijos derivativos.
        \end{exe}
  \item Desde el punto de vista de la obligatoriedad:
        \begin{exe}
          \ex El proceso está determinado por la buena formación de la oración,
              como ocurre con la flexión.
          \ex El proceso es independiente de las condiciones de buena formación
              sintáctica, como ocurre con la formación de palabras.
        \end{exe}
  \item Desde el punto de vista del cambio categorial de la base:
        \begin{exe}
          \ex El proceso no cambia la categoría de la base.
          \ex El proceso cambia la categoría de la base.
        \end{exe}
  \item Desde el punto de vista del significado que aporta el sufijo a la
        formación:
        \begin{exe}
          \ex El morfema posee contenido léxico descriptivo, como los afijos
              derivativos.
          \ex El morfema posee contenido vinculado con las categorías
              lingüísticas, como los sufijos flexivos.
          \ex El morfema no tiene significado, como los interfijos.
          \ex El morfema tiene significado, pero su contenido no es descriptivo
              ni está vinculado con las categorías lingüísticas.
        \end{exe}
\end{consignas}

\textbf{Parte C}

Teniendo en cuenta los resultados obtenidos en su evaluación de los puntos
anteriores, ¿el sufijo \forma{-ardo} responde a las características de los
sufijos derivativos, a las de los flexivos o a las de los apreciativos?
\end{ejercicio}

\section{La formación de palabras}

Las lenguas distinguen entre palabras simples, que no provienen de otras, como
\forma{mar} o \forma{árbol}, y palabras complejas, formadas por procesos
morfológicos, como \forma{marítimo} y \forma{altamar}. La mayoría de las
palabras complejas se obtienen incorporando morfemas a una base, ya que la
morfología léxica es esencialmente aditiva o concatenativa. Según muestran los
ejemplos anteriores, los principales procesos de formación de palabras son la
\termino{derivación} y la \termino{composición}. La diferencia fundamental entre
ellos es la naturaleza de la unidad involucrada en el proceso morfológico:
mientras que en el caso de la derivación se añade un afijo, la composición une
dos o más palabras o temas. Mención aparte merece la \termino{parasíntesis}, en
la que ocurren de manera simultánea dos procesos, que pueden ser ambos de
naturaleza derivativa o uno derivativo y otro compositivo. A continuación,
presentaremos cada uno de estos procedimientos.

\subsection{La derivación}\label{sec:04-derivacion}

Una palabra derivada se obtiene por medio de la incorporación de un prefijo o un
sufijo a la base. Así, \forma{felic-i-dad} es una palabra derivada de
\forma{feliz} por medio del sufijo \forma{-dad} precedido del interfijo
\forma{-i-}; \forma{infeliz} también lo es, pero en este caso por acción del
prefijo \forma{in-}. La conversión y la formación regresiva se relacionan con
este proceso, con la diferencia de que no añaden segmentos morfológicos a la
base (v.~§\ref{sec:04-conversion} y §\ref{sec:04-regresiva}).

\begin{ejercicio}\label{ej:04-derivacion}
Considere el siguiente corpus y resuelva las consignas que se detallan a
continuación:

\begin{tabular}{@{}llll@{}}
  \forma{falsedades}  & \forma{famosas}        & \forma{chateaba}     & \forma{habilidad} \\
  \forma{bromear}     & \forma{deslealtad}     & \forma{cotidianidad} & \forma{amorosísimo} \\
  \forma{desinformó}  & \forma{bondades}       & \forma{desposteó}    & \forma{desagradables} \\
  \forma{grasosos}    & \forma{maldad}         & \forma{desamor}      & \forma{googlearon} \\
  \forma{ruidoso}     & \forma{disconformidad} & & \\
\end{tabular}

\textbf{Parte A}
\begin{consignas}
  \item Agrupe las palabras formadas por procesos comunes e indique de qué
        proceso se trata.
  \item Determine la categoría léxica de la base sobre la que se aplicó la regla
        y la categoría de la palabra resultante.
  \item Indique el valor semántico que los afijos aportan a la palabra.
  \item Realice la descripción fonológica del afijo incorporado a la estructura
        de la palabra.
\end{consignas}

\begin{resuelto}
Grupo 1. a) \forma{famosas, grasosos, amorosísimo, ruidoso}. Proceso común:
derivación con el sufijo \forma{-os-}. b) Categoría de la base: N. Categoría
resultante: A. c) Valor semántico: \sig{que tiene X}.\footnote{X = significado
de la base.} d) Descripción fonológica: \fonol{ˈos}.
\end{resuelto}

\textbf{Parte B}

A partir del análisis anterior, decida cuál de las siguientes generalizaciones
describe con mayor precisión lo observado en cuanto a la interacción entre los
morfemas derivativos y las categorías léxicas:

\begin{exe}
  \ex Todos los afijos derivativos son sensibles a la categoría de la base y
      cambian la clase o subclase de palabra.
  \ex Los sufijos derivativos son sensibles a la categoría de la base y
      típicamente cambian la clase o subclase de palabra; en cambio, los
      prefijos derivativos son típicamente insensibles a la categoría de la base
      y no cambian la clase o subclase de palabra.
  \ex Todos los sufijos derivativos son insensibles a la categoría de la base y
      no cambian la clase o subclase de palabra.
  \ex Los sufijos derivativos son insensibles a la categoría de la base y no
      cambian la clase o subclase de palabra; en cambio, los prefijos
      derivativos típicamente son sensibles a la categoría de la base y cambian
      la clase o subclase de palabra.
\end{exe}
\end{ejercicio}

\subsection{La composición}\label{sec:04-composicion}

Los compuestos se forman por la unión de dos o más palabras o temas. Se
denominan \termino{compuestos patrimoniales} aquellos que se forman con
palabras propias de la lengua, como en \forma{guardapolvo} (\forma{guarda} +
\forma{polvo}). Los \termino{compuestos grecolatinos}, en cambio, se producen
por la unión de los que se conocen como temas cultos
(v.~§\ref{sec:04-temascultos}), como en \forma{geógrafo} (\forma{geo} +
\forma{grafo}). También existen \termino{compuestos mixtos}, que se forman por
la unión de un tema culto y una palabra patrimonial, como \forma{hidroenergía} o
\forma{musicología}.

Según el grado de unión entre las bases, se reconocen \termino{compuestos
léxicos}, como \forma{sacacorchos, boquiabierto, agridulce}, y
\termino{compuestos sintagmáticos}, como \forma{ojo de buey, llave inglesa,
árabe-israelí}. En los compuestos léxicos las bases están plenamente soldadas y
unidas ortográficamente; en cambio, en los sintagmáticos, aunque tienen un
sentido unitario, los componentes no se escriben juntos e incluso pueden tener
una preposición entre ambos: \forma{máquina de coser, esmalte de uñas}.

La relación entre los elementos del compuesto puede ser de coordinación, como en
\forma{compraventa} y \forma{agridulce}, o de subordinación, como en
\forma{altiplanicie} o \forma{pelirrojo}. En los \termino{compuestos
coordinativos} ambos miembros pertenecen a la misma clase de palabra y pueden
parafrasearse por \forma{y} (\forma{agrio}\textsubscript{A} y
\forma{dulce}\textsubscript{A} a la vez); en cambio, en los \termino{compuestos
subordinativos} un miembro es modificador (\forma{pelirrojo}) o complemento
(\forma{marcapasos}) del otro.

\begin{ejercicio}\label{ej:04-composicion}
El siguiente corpus está integrado por distintos tipos de compuestos. Analícelos
y resuelva las consignas de abajo:

\begin{tabular}{@{}llll@{}}
  \forma{sacacorchos}      & \forma{sordomudo}      & \forma{aguardiente}      & \forma{lavarropas} \\
  \forma{socioeconómico}   & \forma{caradura}       & \forma{guardapolvos}     & \forma{blanquiceleste} \\
  \forma{pisapapeles}      & \forma{pelapapas}      & \forma{cazarrecompensas} & \forma{radioaficionado} \\
  \forma{dentoalveolar}    & \forma{matabichos}     & \forma{italoargentino}   & \forma{hierbabuena} \\
  \forma{cometierra}       & \forma{chupamedias}    & \forma{limpiaparabrisas} & \forma{agridulce} \\
\end{tabular}

\begin{consignas}
  \item Agrupe los que han sido obtenidos a partir de un esquema compositivo
        común.
  \item Segmente sus bases.
  \item A partir de uno de los ejemplos, realice la descripción estructural de
        cada esquema de palabra.
  \item Describa la regla que subyace a la formación de las palabras de cada
        grupo.
\end{consignas}

\begin{resuelto}
Grupo 1. a) \forma{aguardiente, caradura, radioaficionado, hierbabuena}. b)
\forma{agu(a)-ardiente; cara-dura; radio-aficionado; hierba-buena}. c)
[[agua]\textsubscript{N}[ardiente]\textsubscript{A}]\textsubscript{N}. d) N + A >
N.
\end{resuelto}
\end{ejercicio}

\begin{ejercicio}[Temas cultos]\label{ej:04-filos}
Considere el siguiente corpus integrado por elementos léxicos que comparten un
tema culto que proviene etimológicamente de la palabra del griego antiguo
\textit{philos} y resuelva las consignas de abajo:

\forma{filarmonía, filosofía, melófila, filonazi, hidrófilo, americanófila,
filología, tecnófilo, filántropo}

\begin{resuelto}[Ejemplo resuelto]
\forma{filarmonía}:
[[fil]\textsubscript{TC}[armonía]\textsubscript{N}]\textsubscript{N}; TC + N > N.
\end{resuelto}

\textbf{Parte A}

Realice la descripción estructural de cada palabra. Recuerde que los temas cultos
no pertenecen a ninguna clase de palabra, pero que las palabras que forman sí
están categorizadas.

\textbf{Parte B}

Compare el funcionamiento observado en los casos de arriba con el de los afijos
y el de las raíces y decida cuáles de las generalizaciones de abajo describen
con mayor precisión el comportamiento observado.

\begin{consignas}
  \item Con respecto a la distribución:
        \begin{exe}
          \ex al igual que los prefijos, los temas cultos son afijos que
              aparecen exclusivamente antes de la base.
          \ex al igual que los sufijos, los temas cultos son afijos que aparecen
              exclusivamente después de la base.
          \ex al igual que las raíces, los temas cultos son bases y pueden
              ocupar distintas posiciones en los esquemas compositivos.
          \ex los temas cultos son afijos, pero pueden ocupar distintas
              posiciones en los esquemas compositivos.
        \end{exe}
  \item Con respecto a la categoría gramatical:
        \begin{exe}
          \ex los temas cultos no poseen categoría gramatical propia ni
              determinan la categoría de la palabra resultante, al igual que los
              prefijos.
          \ex los temas cultos determinan la categoría de la palabra resultante
              al igual que los sufijos derivativos, pero no poseen categoría
              gramatical propia.
          \ex los temas cultos poseen categoría gramatical propia al igual que
              las raíces, pero no determinan la categoría de los compuestos.
          \ex los temas cultos poseen categoría gramatical propia y determinan
              la categoría de los compuestos al igual que las raíces.
        \end{exe}
\end{consignas}
\end{ejercicio}

\subsection{La parasíntesis}\label{sec:04-parasintesis}

Si bien se reconoce una tendencia a que cada paso de la historia derivacional de
las palabras esté asociado a un afijo derivativo, existen palabras que se forman
por medio de la incorporación simultánea de dos procesos, cada uno de los cuales
añade un morfema con contenido léxico a la base. Este fenómeno, que explica cómo
se obtiene la estructura de \forma{submarin-} sin haber pasado por
\forma{marin-} ni suponer un imposible \juicio{*}{\forma{submar}} (nuestro
interrogante en §\ref{sec:04-ramas}), se denomina \termino{parasíntesis}. La
peculiaridad de los procesos parasintéticos es que permiten suponer como paso
intermedio en la derivación una palabra inexistente.

En función del tipo de proceso que se aplica simultáneamente, existen dos tipos
de esquemas parasintéticos: \termino{por derivación} y \termino{por
composición}. El primero consiste en la combinación entre un proceso de
prefijación y otro de sufijación, como \forma{amontonar}, cuya base no es ni
\juicio{*}{\forma{montonar}} ni \juicio{*}{\forma{amontón}}, sino directamente el
sustantivo \forma{montón}, que incorpora en un mismo proceso el prefijo
\forma{a-} y el sufijo verbalizador formado por la vocal temática, al que se
añade la desinencia de infinitivo \forma{-r}, que es un morfema flexivo
\parencite[4702]{serranodolader1999}. Este primer tipo de parasíntesis en
español da lugar a verbos (\forma{enriquecer, agrandar, reciclar}) y a
adjetivos (\forma{anaranjado, aterrador, enrojecido}).

El segundo tipo de parasíntesis resulta de la formación de un derivado a partir
de un compuesto inexistente en la lengua al que se le añade un sufijo, como en
\forma{quinceañera}. Así, \forma{quince años}, que es una construcción
sintáctica y no una palabra compuesta existente en la lengua, igualmente actúa
como la base a la que se le aplica el sufijo \forma{-era} para formar el
derivado \forma{quinceañera}.

\begin{ejercicio}\label{ej:04-parasintesis}
Considere el siguiente corpus integrado por palabras parasintéticas y resuelva
las consignas de abajo:

\forma{desalmado, terraplanista, subterránea, mileuristas, salpimentar,
infumable}

\begin{consignas}
  \item Identifique en cada caso la base sobre la que se aplicó la parasíntesis.
  \item Realice la descripción estructural de la palabra y formalice los pasos
        de la derivación.
  \item Indique, en cada caso, de qué tipo de parasíntesis se trata (por
        derivación o por composición).
\end{consignas}

\begin{resuelto}
\forma{desalmado}. a) Base: \forma{alma-}. b)
[des-[alm]\textsubscript{N}-ad-o]\textsubscript{A}; \forma{des-} + N +
\forma{-ad-o} > A. c) Palabra parasintética por derivación.
\end{resuelto}
\end{ejercicio}

\section{Procesos menores de formación de palabras}

La morfología del español también dispone de mecanismos de creación léxica que
no agregan material fónico a la palabra. Entre ellos se pueden diferenciar, por
un lado, la \termino{conversión}, que no añade ningún afijo a la base, y la
\termino{formación regresiva} y los \termino{acortamientos}, que operan
sustrayendo formantes morfológicos a la base. Por el otro lado, otros procesos,
como los \termino{cruces léxicos} y los \termino{acrónimos}, directamente no
atienden a las fronteras morfológicas de la palabra y se explican a partir de
elisiones basadas en la estructura fonológica de la palabra. Por último, el
\termino{reanálisis} o \termino{etimología popular} es un proceso cuya
explicación recurre a la semántica, ya que lo que hace es reinterpretar el
significado de una parte de una palabra.

\subsection{La conversión}\label{sec:04-conversion}

Cuando una palabra cambia su categoría gramatical sin haber sufrido la
incorporación ni la sustracción de ningún segmento, se entiende que su
estructura morfológica sufrió una \termino{conversión}. La conversión se puede
considerar un proceso derivativo que, si bien no modifica los formantes
morfológicos, puede reanalizar, es decir, reinterpretar, alguno de los segmentos
de la palabra. En inglés es un proceso común que convierte sustantivos en
verbos: \forma{fish} (\sig{pez}) > \forma{(to) fish} (\sig{pescar});
\forma{house} (\sig{casa}) > \forma{(to) house} (\sig{alojar}). En español, por
ejemplo, los adverbios adjetivales, como \forma{rápido} en \forma{Llegué
rápido}, coinciden superficialmente con formas adjetivas flexionadas en
masculino singular. Sin embargo, su vocal final no puede considerarse un sufijo
flexivo porque, a diferencia de los adjetivos que flexionan en género y número,
los adverbios son palabras invariables; por tanto, la vocal final se ha
reanalizado como parte del tema adverbial. Así, en \exref{ex:04-rapido-a} y
\exref{ex:04-rapido-b}, \forma{rápido/rápida} son adjetivos predicativos que
concuerdan en género y número con sus respectivos sujetos (v. \capref{9}); en
cambio, \forma{rápido} en \exref{ex:04-rapidoadv} no manifiesta concordancia
porque es un adverbio equivalente a \forma{rápidamente}.

\begin{exe}
  \ex\label{ex:04-rapido}
    \begin{xlist}
      \ex\label{ex:04-rapido-a} Pedro es \forma{rápido} para hacer las cosas.
      \ex\label{ex:04-rapido-b} Margarita es \forma{rápida} para hacer las cosas.
    \end{xlist}
  \ex\label{ex:04-rapidoadv}
    \begin{xlist}
      \ex Él llegó \forma{rápido/rápidamente} a la estación.
      \ex Ella llegó \forma{rápido/rápidamente}.
    \end{xlist}
\end{exe}

Otro ejemplo de conversión en español puede considerarse el empleo de los
participios verbales \exref{ex:04-agotado-a} como adjetivos
\exref{ex:04-agotado-b} y \exref{ex:04-agotado-c}, y también el de los adjetivos
\exref{ex:04-terminal-a} como sustantivos \exref{ex:04-terminal-b}:

\begin{exe}
  \ex\label{ex:04-agotado}
    \begin{xlist}
      \ex\label{ex:04-agotado-a} Este trabajo nos \forma{ha agotado} a todos / a todas.
      \ex\label{ex:04-agotado-b} Juan quedó \forma{agotado} por el esfuerzo.
      \ex\label{ex:04-agotado-c} María quedó \forma{agotada} por el esfuerzo.
    \end{xlist}
  \ex\label{ex:04-terminal}
    \begin{xlist}
      \ex\label{ex:04-terminal-a} Llegamos temprano a la estación terminal.
      \ex\label{ex:04-terminal-b} La terminal estaba vacía.
    \end{xlist}
\end{exe}

Como vemos en \exref{ex:04-agotado-a}, el participio en los tiempos compuestos es
invariable, mientras que en el uso adjetivo debe concordar con el sustantivo al
que se refiere \exref{ex:04-agotado-b} y \exref{ex:04-agotado-c}. En
\exref{ex:04-terminal-a} tenemos un uso adjetivo de \forma{terminal} referido al
sustantivo \forma{estación}; en \exref{ex:04-terminal-b}, \forma{terminal} se
usa solo, precedido del determinante ---\forma{la terminal}---, por lo que
adquiere uso sustantivo.

\begin{ejercicio}\label{ej:04-conversion}
Las palabras que se destacan en los siguientes ejemplos, en ausencia de un
contexto sintáctico, podrían clasificarse en más de una
categoría.\footnote{Dado que el resultado de la conversión se vuelve perceptible
en el entorno sintáctico, este ejercicio entra en diálogo con los capítulos 5 y
11, sobre estructura del sintagma y perífrasis verbales. Recomendamos su
consulta para resolverlo.}

\begin{consignas}
  \item Segmente morfológicamente cada una de ellas, identificando raíz y tema.
  \item Analice los contextos sintácticos que corresponden a cada forma y
        permiten asignarla a una clase de palabra diferente en cada caso.
\end{consignas}

\begin{exe}
  \ex \textbf{cantar}
    \begin{xlist}
      \ex Aquí me pongo a cantar / Al compás de la vigüela, / Que el hombre que
          lo desvela / Una pena estrordinaria, / Como la ave solitaria / Con el
          cantar se consuela (José Hernández, \textit{El gaucho Martín Fierro},
          I).
      \ex Cantar de los cantares, el cual es de Salomón (\textit{Antiguo
          Testamento}, «Cantar de los cantares», 1:1).
    \end{xlist}
  \ex \textbf{comida}
    \begin{xlist}
      \ex El chef algoritmo: cómo es la comida hecha con inteligencia artificial
          que llegó a la Argentina.
          (\url{https://www.clarin.com/tecnologia/chef-algoritmo-comida-hecha-inteligencia-artificial-llego-argentina_0_znwOtd0bz.html})
      \ex Terror por una mujer que fue comida por un cocodrilo en Australia.
          (\url{https://www.minutouno.com/notas/3045423-terror-una-mujer-que-fue-comida-un-cocodrilo-australia/})
    \end{xlist}
  \ex \textbf{cursada}
    \begin{xlist}
      \ex Examen parcial es aquella evaluación que se toma en distintos momentos
          de la cursada de la materia, abarcando contenidos parciales.
          (\url{http://www.mdp.edu.ar/academica/ingreso/preguntas-uni.html})
      \ex [\ldots] a raíz de haber resultado infructuosa la notificación cursada
          en el domicilio que ahora denuncia la quejosa.
          (\url{https://www.abogados.com.ar/ratifican-validez-de-la-notificacion-del-traslado-de-la-demanda-practicada-en-el-domicilio-social-inscripto/16571})
    \end{xlist}
  \ex \textbf{privado}
    \begin{xlist}
      \ex Sin piso, el empleo privado cayó por sexto mes consecutivo en febrero.
          (\url{https://www.cronista.com/economiapolitica/Sin-piso-el-empleo-privado-cayo-por-sexto-mes-consecutivo-en-febrero-20190425-0040.html})
      \ex Artículo 61: El agente de la Administración Municipal no podrá ser
          privado de su empleo, ni objeto de sanciones disciplinarias sino por
          las causas y procedimientos determinados en este Estatuto.
          (\url{http://www.gob.gba.gov.ar/legislacion/legislacion/l-11757.html})
    \end{xlist}
  \ex \textbf{doctorando}
    \begin{xlist}
      \ex Me estoy doctorando, ergo, debo estar muy loca.
          (\url{http://lagataroma.blogspot.com/2017/02/me-estoy-doctorando-ergo-debo-estar-muy.html})
      \ex Un estudio asegura que los doctorandos son seis veces más propensos a
          desarrollar ansiedad o depresión en comparación con la población
          general.
          (\url{https://elpais.com/elpais/2018/03/15/ciencia/1521113964_993420.html})
    \end{xlist}
  \ex \textbf{cumbia}
    \begin{xlist}
      \ex \textoaincorporar{el fragmento de \textit{La luna y tú} (Ráfaga)}
      \ex Cumbia tu aldea y cumbiarás el mundo.
          (\url{https://www.pagina12.com.ar/200099-cumbia-tu-aldea-y-cumbiaras-el-mundo})
    \end{xlist}
\end{exe}

\begin{resuelto}
1. a) \forma{cant-a-r}. \forma{cantar}\textsubscript{a}: raíz \forma{cant-},
tema \forma{canta-}; \forma{cantar}\textsubscript{b}: \forma{cant-ar(-es)}, raíz
\forma{cant-}, tema \forma{cantar-}. b) \forma{cantar}\textsubscript{a}: se
trata de un verbo porque integra una perífrasis verbal (v. \capref{11}):
\forma{ponerse a cantar} (= \sig{empezar a cantar});
\forma{cantar}\textsubscript{b}: se trata de un sustantivo porque va precedido
del determinante \forma{el} y puede pluralizarse: \forma{los cantar-es}.
\end{resuelto}
\end{ejercicio}

\subsection{La formación regresiva}\label{sec:04-regresiva}

Hay algunas palabras que, contra lo que podríamos suponer, no se han formado por
un incremento afijal, sino por la sustracción de un morfema de la base, en un
proceso conocido como \termino{formación regresiva}. Así, por ejemplo, en la
historia de la lengua está atestiguado que, contra lo esperado, no es
\forma{deslizar} la forma que proviene de \forma{desliz}, sino a la inversa: el
nombre \forma{desliz} se formó a partir del verbo \forma{deslizar}. Lo mismo ha
sucedido con el verbo \forma{legislar}, que deriva de \forma{legislador} y no al
revés, es decir que para su formación se eliminó el sufijo \forma{-dor} de la
palabra base.

\begin{ejercicio}\label{ej:04-regresiva}
Considere los siguientes enunciados, que contienen palabras neológicas acuñadas
para un ámbito específico o para causar un efecto publicitario o humorístico, y
resuelva las siguientes consignas:

\begin{consignas}
  \item Identifique los temas que funcionan como base y resultado del proceso
        con sus categorías y el segmento morfológico afectado.
  \item Formule como una regla el proceso morfológico involucrado en la
        formación del resultado.
  \item Realice una paráfrasis del significado de la palabra resultante del
        proceso.
\end{consignas}

\begin{exe}
  \ex \textbf{televisar}: El canal que empezará a televisar la NBA en la
      Argentina y cómo se verá la jornada inaugural.
      (\url{https://www.lavoz.com.ar/deportes/basquet/el-canal-que-empezara-a-televisar-la-nba-en-la-argentina-y-como-se-vera-la-jornada-inaugural})
  \ex \textbf{espectar}: El modo espectador es casi igual al modo creativo,
      solo que, como el nombre lo dice, te da la posibilidad de espectar un mundo
      pero no interactuar con este\ldots{}
      (\url{https://aminoapps.com/c/mcft-s/page/item/espectador/KW0u_KIRYJlRb8pNnmb26bqbKj3QRBz})
  \ex \textbf{alfabeto}: La diferencia entre un alfabeto funcional y un
      analfabeto funcional radica en que el alfabeto funcional pondrá en
      práctica competencias, conocimientos y habilidades adquiridas durante su
      vida académica\ldots{}
      (\url{http://transformacion-educativa.com/2do-congreso/ponencias/Eje-1/L2-102.html})
  \ex \textbf{antitranspirar}\pendiente{Falta el ejemplo en el original}
\end{exe}

\begin{resuelto}
1. a) Base: \forma{televisión}. Categoría: N. Resultado: \forma{televis(a)-}.
Categoría: V. Segmento afectado: \forma{-sión}. b) \forma{televisar} <
\forma{televisión}: N -- \forma{-sión} > V. c) Significado: \sig{transmitir
imágenes a distancia}.
\end{resuelto}
\end{ejercicio}

\subsection{El acortamiento}

Una palabra acortada resulta de la sustracción de material fonológico de la
base. Como resultado de este proceso, se conserva o bien su parte inicial,
proceso que se conoce como \termino{apócope}, o bien la parte final, lo que se
denomina \termino{aféresis}. Son apócopes de uso cotidiano: \forma{celular} >
\forma{celu}; \forma{señorita} (= \sig{maestra}) > \forma{seño}/\forma{se};
\forma{bicicleta} > \forma{bici}. Como ejemplos de aféresis podemos mencionar:
\forma{violonchelo} > \forma{chelo} y \forma{autobús} > \forma{bus}. Un ejemplo
característico de este proceso lo constituyen los hipocorísticos, es decir, los
acortamientos de los nombres propios: \forma{Isa} por \forma{Isabel};
\forma{Maxi} por \forma{Maximiliano}. Muchos nombres permiten tanto el
acortamiento por apócope como por aféresis: \forma{Fer} y \forma{Nando}, por
\forma{Fernando}; \forma{Guada} y \forma{Lupe}, por \forma{Guadalupe}.

\begin{ejercicio}\label{ej:04-acortamiento}
Observe el siguiente corpus integrado por palabras acortadas y a partir de ellas
resuelva lo pedido más abajo:

\forma{cheto, profe, facu, tano, ñora, auto, moto, compu, ñero/ñeri, tamos,
taxi}

\begin{consignas}
  \item Agrupe las que se encuentran formadas por apócope o aféresis.
  \item Identifique sus bases y determine su categoría y la del resultado.
  \item Señale qué principio fónico se observa en la selección de los segmentos
        conservados.
  \item Evalúe qué modificación semántica produce la aplicación de estos
        procesos.
\end{consignas}

\begin{resuelto}
Grupo 1. a) Acortamientos por aféresis: \forma{cheto, tano, ñora, ñero/ñeri,
tamos}. b) \forma{cheto}\textsubscript{A} < \forma{concheto}\textsubscript{A};
\forma{tano}\textsubscript{A} < \forma{napolitano}\textsubscript{A};
\forma{ñora}\textsubscript{N} < \forma{señora}\textsubscript{N};
\forma{ñero/ñeri}\textsubscript{N} < \forma{compañero}\textsubscript{N};
\forma{tamos}\textsubscript{V} < \forma{estamos}\textsubscript{V}. c) Se
conservan las últimas dos sílabas de la base. La palabra resultante tiene
acentuación grave. d) El acortamiento por aféresis no produce cambio en el
significado de la base.
\end{resuelto}
\end{ejercicio}

\subsection{El cruce léxico o \textit{blending}}

Los \textit{blends} o «palabras macedonia» están formados a partir de los
extremos fónicos de otras palabras que se combinan entre sí. Estos extremos no
coinciden con morfemas, sino que pueden ser sílabas u otras secuencias de
segmentos tomadas del comienzo o del final de la palabra. Frecuentemente, el
punto de encuentro entre las bases coincide fónicamente. Por ejemplo, la palabra
\forma{estanflación} se obtiene a partir de fragmentos de las bases
\forma{estancamiento} e \forma{inflación} fusionando, respectivamente, los
extremos anterior (\forma{estan-}) y posterior (\forma{-flación}) al primer
sonido nasal que comparten (\forma{-n-}):
[\forma{estan(camiento)}]\textsubscript{N} +
[\forma{(i)nflación}]\textsubscript{N} >
[\forma{estanflación}]\textsubscript{N}. Por la naturaleza de las unidades
involucradas, el cruce léxico es un proceso morfológico que forma palabras
compuestas, cuya unión involucra reajustes fonológicos.

\begin{ejercicio}\label{ej:04-cruces}
El siguiente corpus está integrado por cruces léxicos de distinto tipo. Luego de
observar sus diferencias, realice las siguientes consignas:

\begin{consignas}
  \item Identifique las bases a partir de las cuales se formó cada compuesto.
  \item Determine la categoría de ambas bases y la del resultado.
  \item Analice cómo han sido seleccionados los segmentos de las bases
        utilizados en el compuesto.
\end{consignas}

\begin{exe}
  \ex Un \textbf{módem} es el dispositivo que permite a tus computadoras,
      teléfonos inteligentes, tabletas y otros dispositivos conectarse a
      Internet. (\url{https://www.lenovo.com/ar/es/glosario/que-es-modem/})
  \ex \textbf{Conmebol}
  \ex \textbf{Mercosur}
  \ex \textbf{vitalitud}
  \ex \textbf{deshamburguesar}
  \ex \textbf{artivismo}
  \ex \textbf{juernes}
\end{exe}

\imagenpendiente{Imágenes de los ítems 2 a 7}

\begin{resuelto}
1. \forma{módem}. a) Bases: \forma{modulador} + \forma{demodulador}. b)
Categoría de las bases: N + N. Categoría del resultado: N. c) Se seleccionaron
los primeros tres segmentos sonoros de cada una de las palabras, aprovechando el
fonema consonántico común: [mod(ulador) + dem(odulador)].
\end{resuelto}
\end{ejercicio}

\subsection{La acronimia}

La \termino{acronimia} resulta de la lexicalización, o conversión en palabra, de
una sigla. Las \termino{siglas} son formaciones que se producen a partir de las
iniciales de las palabras de clase mayor (sustantivos, adjetivos, verbos) de una
denominación compuesta por varias palabras. En algunos casos, también pueden
estar integradas por más de uno de los sonidos iniciales de cada palabra
(\forma{RENAPER}, por \forma{Registro Nacional de las Personas}) o incluir
palabras de clases menores (\forma{ADEBA}, por \forma{Asociación de Bancos de
Buenos Aires}). También existen siglas que resultan de una combinación entre
ambas estrategias, como \forma{EDENOR} (\forma{Empresa Distribuidora Norte}).
Pueden pronunciarse fonémicamente, es decir, de acuerdo con la realización
habitual de cada uno de sus segmentos sonoros (por ejemplo, \forma{ARCA}, para
\forma{Agencia de Recaudación y Control Aduanero}), o por medio de la mención de
los nombres de los grafemas que las integran (\forma{GNC}, pronunciado
\fonol{xeneˈse}, para \forma{Gas Natural Comprimido}).\pendiente{El original dice
«Gas Natural Licuado»}

Todas las siglas mencionadas anteriormente se encuentran lexicalizadas, por lo
que se emplean como si fueran palabras de la lengua; así, la \forma{Universidad
de Buenos Aires} es popularmente conocida como la \forma{UBA} y la
\forma{Organización de las Naciones Unidas}, como la \forma{ONU}. Las siglas
lexicalizadas, que pasan a funcionar en la lengua como palabras, incluso
permiten ser derivadas: por ejemplo, de \forma{SIDA} (\forma{síndrome de
inmunodeficiencia adquirida}) se ha formado \forma{sidoso} (\sig{la persona que
padece este síndrome}) y de \forma{OVNI} (\forma{objeto volador no
identificado}) se ha derivado \forma{ovnilogía} (disciplina que estudia los
ovnis).\footnote{También se emplea \textit{ufología}, derivado de la sigla en
inglés UFO (\textit{unidentified flying object}).}

\begin{ejercicio}\label{ej:04-acronimos}
El siguiente corpus está integrado por distintos acrónimos. En cada caso:

\begin{consignas}
  \item Identifique sus bases y determine su categoría.
  \item Determine la categoría de la construcción resultante.
  \item Observe si existen marcas morfológicas que evidencien la lexicalización
        del acrónimo.
  \item Evalúe qué modificación semántica produce la aplicación del proceso.
\end{consignas}

\begin{exe}
  \ex \textbf{TICs}: Una mirada institucional sobre la incorporación de las TICs.
      (\url{https://twitter.com/andynsane/status/1085360342513958914?lang=es})
  \ex \textbf{RAE}: La RAE vela, desde 1713, por el buen uso y la unidad de la
      lengua española en el mundo. (\url{https://www.instagram.com/laraeinforma/})
  \ex \textbf{ovnis}: La cantidad de personas que cree que los avistamientos de
      ovnis ofrecen una prueba probable de vida extraterrestre aumentó del 20\,\%
      en 1996 al 34\,\% en 2022.
      (\url{https://www.lanacion.com.ar/lifestyle/en-las-redes/la-peligrosa-idea-de-que-extraterrestres-visitaron-la-tierra-esta-cada-vez-mas-extendida-nid09092024/})
  \ex \textbf{FIFA / AFA}: La FIFA advirtió «severas sanciones» si el Gobierno
      interviene la AFA.\pendiente{El enlace del original repite el del ítem 1}
  \ex \textbf{UNA}\pendiente{Falta el ejemplo en el original}
\end{exe}

\begin{resuelto}
1. a) \forma{TICs}: tecnología\textsubscript{N}, información\textsubscript{N},
comunicación\textsubscript{N}. b) Categoría: N. c) El segmento final corresponde
con la pluralización, esperable para un nombre, si bien no se manifiesta por
medio del alomorfo correspondiente a la terminación del tema (\forma{-es}). d)
El acrónimo equivale referencialmente a la frase que abrevia.
\end{resuelto}
\end{ejercicio}

\subsection{Reanálisis o etimología popular}

El \termino{reanálisis} morfémico o simplemente reanálisis
\parencite{escandell2011}, también conocido como \termino{etimología popular}, es
el proceso por el que el hablante analiza las partes de una palabra y las
reinterpreta asignándoles significados que no tenían originalmente. Así, en
inglés, la palabra \forma{hamburger} originalmente era el gentilicio
correspondiente a los habitantes de la ciudad alemana de Hamburgo, pero la
primera parte se confundió con \forma{ham} (\sig{jamón}) y la palabra se
reinterpretó como un compuesto [[ham] [burger]], lo que dio origen a otros
compuestos, como \forma{cheeseburger} (\sig{hamburguesa de queso}),
\forma{fishburger} (\sig{hamburguesa de pescado}), \forma{beefburger}
(\sig{hamburguesa de carne}) y también a la creación, por formación regresiva,
de la palabra \forma{burger}, como denominación general de cualquier tipo de
hamburguesa. Otro caso muy conocido de reanálisis es el de la palabra
\forma{gate} (lit. \sig{portón}) en \forma{Watergate} (lit. \sig{portón de
agua}), nombre de las oficinas en las que, a principios de los 70, se produjo un
escándalo de espionaje que obligó a dimitir al presidente de los EE.\,UU.,
Richard Nixon. El segundo elemento del compuesto, \forma{-gate}, comenzó a
utilizarse en otras palabras con el nuevo significado de \sig{asunto
escandaloso}: \forma{Irangate} (venta de armas a Irán), \forma{sexgate}
(\textit{affaire} sexual del presidente Clinton en la Casa Blanca). El mecanismo
también ha sido adoptado en español: \forma{Yomagate} (caso de narcotráfico y
corrupción que involucraba a la familia Yoma), \forma{FIFAgate} (escándalo de
corrupción en la FIFA), etc.

El reanálisis también se puede producir con fines humorísticos, como se muestra
en la viñeta de la figura~\ref{fig:04-afortunado}, en la que, en el adjetivo
parasintético \forma{a-fortun-ad-o} (\sig{que tiene fortuna o suerte}), el
prefijo inicial se reinterpreta como si se tratara del prefijo negativo
\forma{a-} (\forma{a-moral, a-normal}, etc.) aplicado a \forma{fortunado},
produciendo el significado opuesto de \sig{desprovisto de fortuna}.

\figurapendiente{Reanálisis humorístico de \textit{afortunado}}{fig:04-afortunado}

\begin{ejercicio}[\textit{Siempre x, nunca in-x}]\label{ej:04-siempre}
La construcción sintáctica \forma{siempre x, nunca in-x}, de uso frecuente en las
redes sociales, se emplea para oponer una palabra al que se considera su
antónimo, para reafirmarla como propiedad caracterizadora del que la enuncia.
Esta construcción se popularizó como meme en las redes sociales a raíz de un tuit
de la modelo chilena Nicole Moreno, que tenía como una de sus frases
características \forma{Siempre fuerte, nunca infuerte} (\textit{sic}: «siempre
foerte nunca infoerte»). Por la productividad del proceso morfológico
subyacente, el esquema se puso de moda y produjo una catarata de neologismos
lúdicos y efímeros, como el de la propia Moreno. Más interesante aún, con la
misma calidad de fugaces y lúdicas, se acuñaron palabras como las que se
destacan en los ejemplos que, junto con otros de origen diverso, se reproducen a
continuación. A partir de ellas, resuelva las siguientes consignas:

\begin{consignas}
  \item Determine la categoría de las bases sobre las que se aplica esta
        operación y la de las palabras resultantes.
  \item Identifique el segmento que es reinterpretado en la formación de estas
        palabras e indique el significado que se le asigna.
  \item Explique qué particularidad presenta \forma{imberbe} > \forma{berbe}.
\end{consignas}

\begin{exe}
  \ex \textbf{bécil}: Siempre bécil, nunca imbécil.
      (\url{https://x.com/PabloGa27864933})
  \ex \textbf{pecable}: Siempre pecable, nunca impecable.
      (\url{https://x.com/teseodeseo_?lang=ar})
  \ex \textbf{berbe}: Siempre berbe, nunca imberbe.
      (\url{https://twitter.com/filozoofico})
  \ex \textbf{teligente}\pendiente{Falta el ejemplo en el original}
\end{exe}

\begin{resuelto}
a) 1. \forma{bécil}. Base: \forma{imbécil}. Categoría: A. Resultado:
\forma{bécil-}. Categoría: A. b) \forma{bécil} designa la propiedad opuesta a
\forma{imbécil}.\pendiente{El original agrega «c) Ejemplo: alfabeto: A --- an- >
A», que corresponde al ejercicio~\ref{ej:04-regresiva}}
\end{resuelto}
\end{ejercicio}

\begin{ejercicio}\label{ej:04-reanalisis}
¿Cómo se podría explicar la formación de las siguientes palabras? ¿Se trata de
casos de reanálisis? Justifique su respuesta teniendo en cuenta la palabra de la
que derivan y la interpretación del significado.

\begin{exe}
  \ex \textbf{flamable}\imagenpendiente[14mm]{Imagen del ítem 1}
  \ex \textsc{El periodista}.--- Mire, señor militar: me dicen que antes era
      dictadura y ahora \textbf{dictablanda}. ¿Qué les debo contestar?
      (\textit{Diario El Sol de España}, 12/03/1930;
      \url{https://cvc.cervantes.es/el_rinconete/anteriores/septiembre_12/17092012_02.htm})
  \ex {[\ldots]} \textit{Episodio IV: Una nueva esperanza} fue la primera
      película de Star Wars, el inicio de la saga como la conocemos. Mientras
      que los episodios V y VI se convirtieron en secuelas, los episodios I, II
      y III fueron las \textbf{precuelas}.
\end{exe}

\begin{resuelto}
1. \forma{flamable}. Es un reanálisis o reinterpretación humorística del prefijo
locativo \forma{in-}, presente en \forma{in-tubar} o \forma{in-filtrar}, como si
se tratara del prefijo negativo homónimo (\forma{in-sano, in-moral}). De ahí
que, paradójicamente, el significado de la palabra resultante no cambie:
\sig{que puede producir llamas} (quemarse).
\end{resuelto}
\end{ejercicio}

\section{Las variantes morfológicas de las palabras}\label{sec:04-variantes}

Uno de los modos posibles de clasificación de las palabras que forman el léxico
de la lengua permite separarlas en categorías variables e invariables. Esto
quiere decir que hay clases de palabras que oponen de forma morfológica
manifiesta determinados contenidos semánticos altamente gramaticalizados, como
la cantidad ---\forma{perro/perros}--- o el tiempo de ocurrencia de un suceso
---\forma{canté/canto/cantaré}---. Las manifestaciones lingüísticas de estos
significados generales se conocen como categorías lingüísticas. Las que se
reconocen para el español son las siguientes: género, número, persona, caso,
tiempo, modo, aspecto y voz (aunque esta última no tiene manifestación
morfológica en español). La expresión de estas categorías, hecha la excepción de
la voz, se da a través de la flexión, que Fábregas (2015, p.~175)\pendiente{En la
bibliografía figura Fábregas (2013)} define como el «conjunto de operaciones que
nos proporciona la forma que una palabra debe adoptar dado cierto contexto
sintáctico». Efectivamente, la flexión es relevante para la sintaxis porque en
algunos casos expresa de manera abierta la existencia de determinadas relaciones
gramaticales ---las distintas formas de la concordancia--- y es un factor
esencial en la buena formación de la oración. Así, ni \juicio{*}{\forma{el gato
blanca}} ni \juicio{*}{\forma{los gato blancas}} son sintagmas bien formados, ya
que, según se verá en el próximo capítulo, tanto el artículo como el adjetivo
deben concordar en género y número con el nombre núcleo del sintagma, de modo
que la forma gramatical sería \forma{el gato blanco}. La morfología flexiva del
español es fundamentalmente aditiva: se realiza mediante la anexión de sufijos a
la base de la palabra.

\subsection{Las categorías gramaticales y su manifestación en la morfología}\label{sec:04-categorias}

Las categorías lingüísticas son moldes a través de los cuales la mente
conceptualiza el mundo \parencite[6]{lakoff1987}. Se expresan, como hemos
visto, morfológicamente, pero también pueden hacerlo léxicamente. Por ejemplo,
para el género, además de contrastes como \forma{perr-o/perr-a}, que se
manifiestan en la oposición entre las desinencias \forma{-o} y \forma{-a},
también hay oposiciones léxicas (mediante cambio de palabra) ---\forma{vaca/toro}---
o sintácticas (mediante el cambio del determinante) ---\forma{el dentista / la
dentista}---.

Las categorías gramaticalizan ---es decir, manifiestan mediante los recursos
gramaticales de la lengua--- aspectos que los hablantes extraen de la situación
extralingüística. En \textcite{giammatteo2018} se propone una clasificación que
parte de esta vinculación de las categorías con la situación comunicativa. La
clasificación da cuenta de una gradación que va de aquellas categorías
fuertemente enraizadas en la situación comunicativa hasta las más abstractas y
aparentemente menos vinculadas al contexto de situación. En función de su grado
de relación con el entorno en el que se pone en uso la lengua, las categorías se
organizan en tres grupos: las de la situación comunicativa, las del evento y las
de la escena oracional.

En primer lugar, las \termino{categorías de la situación comunicativa} son las
que están estrechamente relacionadas con los participantes del coloquio y el
marco en el que se produce el intercambio lingüístico. En este grupo se incluyen
las categorías de persona y tiempo, que se consideran categorías deícticas
porque establecen «un tipo de vínculo referencial entre ciertas unidades o
expresiones lingüísticas y aquello que representan en el mundo o en el universo
del discurso» \parencite[932]{eguren1999}. La tercera categoría de este grupo,
el modo, si bien no es deíctica, establece una relación indirecta con la
situación comunicativa, ya que tiene que ver con el hablante y con la actitud
que asume respecto de lo que dice.

\begin{itemize}
  \item La \termino{persona} señala a los participantes de la comunicación: el
        emisor (\forma{yo} (sg.) / \forma{nosotros/as} (pl.)); el destinatario
        (\forma{vos / tú / Ud.} (sg.) // \forma{Uds. / vosotros/as}
        (pl.))\footnote{La forma \textit{vosotros/as} solo se usa en el español
        peninsular, mientras que \textit{tú} se utiliza en España y en otras
        regiones de Latinoamérica. \textit{Vos} es la forma característica del
        Río de la Plata, pero se emplea también en algunas otras partes de
        Latinoamérica. Mientras \textit{tú} y \textit{vos} indican confianza,
        \textit{Ud.} se emplea en el trato formal.} y un tercer participante que
        no forma parte de la relación de alocución (\forma{él / ella} (sg.) //
        \forma{ellos / ellas} (pl.)). Este tercer participante puede o no estar
        en la situación comunicativa y puede ser cualquier tipo de entidad,
        \sig{animada} o no. Por eso a la tercera persona también se la conoce
        como la «no persona» y es la que se emplea en todos los usos vinculados
        con la impersonalidad: \forma{Mañana llueve / Se habla mucho y se hace
        poco}. Además de en los pronombres, la persona se manifiesta en la
        desinencia del verbo para concordar con la persona del sujeto:
        \forma{Yo cant-o / Él cant-a}.
  \item El \termino{tiempo} relaciona un evento con un punto o intervalo
        temporal respecto del cual el hecho se sitúa. Este momento temporal
        puede ser extralingüístico y ubicarse directamente como anterior,
        simultáneo o posterior respecto del momento de enunciación, como en
        \forma{Ayer fui al cine}, que indica un evento sucedido con anterioridad
        al momento de la emisión; o bien puede orientarse indirectamente a partir
        de la referencia proporcionada por otro evento, como en \forma{Llegó}
        (Ev\textsubscript{principal}) \forma{cuando todos se habían ido}
        (Ev\textsubscript{2 anterior}).
  \item El \termino{modo} codifica la relación entre el enunciador y lo que
        dice, tanto en términos de la actitud que adopta en la enunciación
        (aseverar, preguntar, ordenar, etc.) como en cuanto al grado de realidad
        que le adjudica al evento (probable, obligatorio, etc.).
\end{itemize}

Tiempo y modo son dos categorías que se manifiestan en la desinencia verbal,
muchas veces en forma de amalgama (v.~§\ref{sec:04-vt}). Así, por ejemplo, en
\forma{cant-a-ba-mos} reconocemos la raíz \forma{cant-}, la vocal temática
\forma{-a-}, la amalgama de tiempo, modo y aspecto que se conoce como TAM y la
amalgama de persona y número \forma{-mos}.

En segundo lugar, las \termino{categorías del evento} son tres: el género, el
número y el aspecto. Estas categorías nos permiten clasificar, cuantificar o
delimitar las entidades de las que hablamos y los sucesos en los que intervienen
y que constituyen el tema de nuestras conversaciones.

\begin{itemize}
  \item El \termino{género} funciona esencialmente como un clasificador de las
        entidades, puesto que divide a los nombres en clases (femenino,
        masculino, neutro u otras, según las lenguas). En español solo se
        interpreta de forma sistemática, en términos de oposición de sexo, para
        los nombres animados (de persona y de animales superiores o domésticos):
        \forma{abogado/abogada; león/leona; gallo/gallina}; en cambio, en los
        nombres de entidades inanimadas resulta inmotivado o inherente (propio
        de la palabra): \forma{el árbol, la piedra}, etc. Aunque no todos los
        sustantivos terminan en vocal, existe una tendencia a que el morfema
        \fonol{-a} funcione como marca de femenino y \fonol{-o} marque el
        masculino, aunque hay excepciones como \forma{(la) mano} o \forma{(el)
        clima}. Además, la vocal \fonol{e} se distribuye complementariamente con
        el morfema \fonol{-o} como marca de masculino (\forma{elefante, jefe}),
        aunque también funciona como marca de ambos géneros en los adjetivos
        (\forma{elegante, amable}). La marcación de género femenino puede
        resultar suplida por otras desinencias (\forma{gallo/gall-ina,
        conde/cond-esa, actor/actr-iz}).
  \item El \termino{número} refiere a la cuantificación de las entidades (una o
        más de una) en los nombres ---\forma{gato/gatos}--- y a la de las
        personas intervinientes en el suceso, en el verbo ---(\forma{él})
        \forma{corre} / (\forma{ellos}) \forma{corren}---. En el verbo se
        manifiesta para concordar con el nombre núcleo del sujeto, junto con la
        categoría de persona: \forma{Yo cant-o / Nosotros cant-a-mos}.
  \item El \termino{aspecto} sirve tanto para expresar la forma interna del
        desarrollo de los eventos (si son dinámicos o estáticos, si poseen o no
        duración, si cuentan con una delimitación inherente o si pueden
        repetirse), como también para hacer referencia a sus alternativas de
        focalización (el punto de vista desde el que se los presenta:
        completados, en desarrollo o bien en alguna de sus fases, como el inicio
        o la conclusión). La principal, y tal vez la más extendida, oposición
        aspectual en las lenguas es la que distingue entre el aspecto
        imperfectivo, que muestra al evento en desarrollo, y el aspecto
        perfectivo, que lo concibe desde afuera, como un todo limitado. Este
        contraste se puede ilustrar para el español con la oposición entre
        \forma{cantó}, que indica un hecho puntual, y \forma{cantaba}, que
        señala un evento durativo.\footnote{Las nociones relativas al aspecto se
        amplían en el \capref{8}, donde se trata la estructura eventiva. También
        se puede consultar \textcite{debin2018}.}
\end{itemize}

Finalmente, las \termino{categorías de la escena oracional}, el caso y la voz o
diátesis (que no aparece en la flexión en español), se encargan de reproducir
dentro de la oración las relaciones que se establecen entre el evento y los
participantes que intervienen en él, manifestados, respectivamente, por el verbo
y sus argumentos.\footnote{Estas cuestiones se desarrollan en el \capref{7},
dedicado a la estructura argumental del verbo y a las funciones sintácticas.}

\begin{itemize}
  \item El \termino{caso} indica el tipo de relación que vincula a los sintagmas
        nominales, que representan los argumentos del verbo, con el predicado.
        Aunque en las lenguas clásicas, como el griego y el latín, y también en
        lenguas modernas, como el ruso y el alemán, estas diferencias se
        manifiestan morfológicamente, en español no tenemos una marca
        morfológica para los diferentes casos que indican las funciones
        sintácticas de los argumentos del verbo, excepto cuando se trata de
        algún pronombre. Los pronombres personales manifiestan caso nominativo,
        que es el que corresponde a la función sujeto ---\forma{yo}---, caso
        objetivo, para las funciones de COD y COI ---\forma{me}--- y caso oblicuo
        o terminal, para cuando siguen a una preposición ---(\forma{para})
        \forma{mí}---.\footnote{Este tema se amplía en el \capref{10}, §2.}
  \item La \termino{diátesis} es el nombre que reciben cada una de las
        construcciones gramaticales que permiten expresar de manera diversa los
        argumentos de un verbo; por ejemplo, la alternancia que se da entre la
        voz activa ---\forma{Yo leo un libro}--- y la pasiva ---\forma{Un libro
        es leído por mí}---. Como se muestra en el ejemplo, la diátesis no tiene
        manifestación morfológica en español, sino que se expresa mediante una
        construcción sintáctica.\footnote{La perífrasis de voz pasiva se trata
        en el \capref{11} y la pasiva cuasirrefleja, en el \capref{10}.}
\end{itemize}

\begin{ejercicio}\label{ej:04-genero}
Existe una percepción extendida entre los hablantes del español según la cual en
el entorno nominal la vocal \fonol{a} en posición final de palabra significa
\sig{género femenino} y la vocal \fonol{o} en esa misma posición se interpreta
como \sig{género masculino}. A partir del siguiente corpus de sintagmas
nominales del español resuelva lo pedido más abajo:

\begin{tabular}{@{}lllll@{}}
  1. la secretaria & 2. el sartén   & 3. la gata    & 4. la cebra    & 5. el naranjo \\
  6. la cuchilla   & 7. la araña    & 8. la mano    & 9. el clima    & 10. el alcalde \\
  11. el amante    & 12. el tilde   & 13. la bolsa  & 14. el calor   & \\
\end{tabular}

\textbf{Parte A}

Cambie el género gramatical de los nombres.

\begin{resuelto}
\forma{la secretari-a / el secretari-o}.
\end{resuelto}

\textbf{Parte B}

Decida cuáles de las siguientes generalizaciones se adecuan mejor a lo
observado:

\begin{consignas}
  \item Con respecto a la disponibilidad de la oposición en género gramatical:
        \begin{exe}
          \ex Todos los sustantivos expresan oposición de género gramatical a
              través de la flexión.
          \ex Solo los sustantivos que designan entidades sexuadas expresan
              oposición de género a través de la flexión.
          \ex Solo los sustantivos que designan entidades sexuadas expresan
              oposición de género, ya sea mediante la flexión o por la
              concordancia con el determinante.
          \ex No todos los sustantivos expresan oposición de género gramatical.
        \end{exe}
  \item Con respecto a la realización de las oposiciones en género:
        \begin{exe}
          \ex La oposición de género gramatical se vincula siempre a que los
              hablantes de distintas variedades del español adscriben los
              sustantivos a distintos géneros.
          \ex La oposición de género gramatical se interpreta en todos los
              sustantivos como un contraste de sexo.
          \ex La oposición en género gramatical puede interpretarse como un
              contraste de rasgos descriptivos de la entidad designada por el
              tema o puede deberse a la variación lingüística.
          \ex La oposición de género gramatical se interpreta siempre como un
              contraste de rasgos descriptivos de la entidad designada por el
              tema.
        \end{exe}
  \item Con respecto a la realización de las oposiciones en género:
        \begin{exe}
          \ex La vocal \fonol{a} siempre forma la variante del sustantivo que
              concuerda con la variante femenina del determinante y la vocal
              \fonol{o} siempre forma la variante que concuerda con la variante
              masculina del determinante.
          \ex No hay una correspondencia estricta entre la vocal con la que
              termina un sustantivo y su género gramatical.
          \ex La vocal \fonol{a} siempre forma la variante del sustantivo que
              concuerda con la variante femenina del determinante y la variante
              masculina se puede expresar tanto con la vocal \fonol{o} como por
              la ausencia de vocal.
          \ex En los sustantivos que designan entidades sexuadas la vocal
              \fonol{a} siempre forma la variante del sustantivo que concuerda
              con la variante femenina del determinante y la vocal \fonol{o}
              siempre forma la variante que concuerda con la variante masculina
              del determinante.
        \end{exe}
\end{consignas}
\end{ejercicio}

\begin{ejercicio}\label{ej:04-verbos}
Considere el siguiente corpus con formas verbales de cada una de las tres
conjugaciones del español y luego resuelva las consignas de abajo:

\begin{tabular}{@{}lll@{}}
  \textbf{A} & \textbf{B} & \textbf{C} \\
  1. \textbf{llegaba} & 1. \textbf{comía} & 1. \textbf{aburría} \\
  2. llega     & 2. come     & 2. aburre \\
  3. llegaban  & 3. comían   & 3. aburrían \\
  4. llegó     & 4. comió    & 4. aburrió \\
  5. llegabas  & 5. comías   & 5. aburrías \\
  6. llegara   & 6. comiera  & 6. aburriera \\
\end{tabular}

\begin{consignas}
  \item Segmente en morfemas todas las palabras del corpus.
  \item Evalúe qué categorías gramaticales se oponen entre la forma 1 de cada
        columna y cada una de las demás de su misma columna.
  \item Haga la descripción morfológica de cada forma verbal, indique el
        contenido semántico de cada uno de los morfemas que la integran y
        complete el cuadro de abajo como se muestra en el modelo.
\end{consignas}

\begin{resuelto}[Ejemplos resueltos]
a) 1. A. \forma{lleg-a-ba}; B. \forma{com-í-a}; C. \forma{aburr-í-a}. 2. A.
\forma{lleg-a}; B. \forma{com-e}; C. \forma{aburr-e}. b) Contraste 1/2: ambas
formas se oponen en tiempo: 1) pretérito / 2) presente. c) 1.
\forma{llegaba/comía/aburría}: 1.ª/3.ª persona singular del pretérito imperfecto
del indicativo. 2. \forma{llega/come/aburre}: 3.ª persona singular del presente
de indicativo.
\end{resuelto}

{\footnotesize
\begin{tabularx}{\linewidth}{@{}l L L L l@{}}
  \toprule
  & Raíz & VT & TAM (amalgama de tiempo, modo y aspecto) & Conc. (número y persona) \\
  \midrule
  1. \forma{llegaba/comía/aburría} & lleg- / com- / aburr- & -a- / -í- & -ba / -a & ∅ \\
  2. \forma{llega/come/aburre}     & lleg- / com- / aburr- & -a / -e   &          & ∅ \\
  3. \forma{llegaban/comían/aburrían} & & & & \\
  4. \forma{llegó/comió/aburrió}   & & & & \\
  5. \forma{llegabas/comías/aburrías} & & & & \\
  6. \forma{llegara/comiera/aburriera} & & & & \\
  \bottomrule
\end{tabularx}}
\end{ejercicio}

\begin{recuadro}[Reflexión final]
La morfología es el componente de la gramática que, a partir de nuestra
competencia lingüística, involucra el conocimiento que nos permite segmentar las
palabras en unidades menores jerárquicamente relacionadas, establecer
vinculaciones formales entre las entradas que forman parte de nuestro léxico
mental y producir tanto elementos léxicos nuevos como las variantes gramaticales
que cada uno de ellos necesita exhibir en la combinatoria sintáctica. En este
sentido, aunque el nivel morfológico interactúa estrechamente con otros
componentes de la lengua, como la fonología, la sintaxis y el léxico, puede
pensarse que constituye un módulo diferenciado, delimitado por la naturaleza de
las unidades y los procesos que intervienen en él.
\end{recuadro}

% --- Secciones finales -------------------------------------------------
\lecturasrecomendadas

\solucionario

\begin{solucion}{\ref{ej:04-terraplanista}}
\textit{Parte A}. a) \forma{terraplanista}: \forma{terra-plan-ist-a};
\forma{negar}: \forma{neg-a-r}; \forma{más}: \forma{más}; \forma{falsedades}:
\forma{fals-e-dad-es}; \forma{gigantesca}: \forma{gigant-esc-a};
\forma{plenamente}: \forma{plen-a-mente}; \forma{explican}:
\forma{ex-plic-a-n}; \forma{entronización}: \forma{en-tron-iz-a-ción};
\forma{pseudocientíficos}: \forma{pseudo-cient-ífic-o-s}.

b) Se ofrecen ejemplos de respuestas posibles a modo ilustrativo.
\forma{terraplanista}: \forma{terra-}: \forma{terrá-que-o, terri-torio};
\forma{-plan-}: \forma{plan-o, plan-icie}; \forma{-ista}: \forma{tax-ista,
conform-ista}. \forma{negar}: \forma{neg-}: \forma{neg-ación, in-neg-able};
\forma{-a-}: \forma{am-a-r, explic-a-n}; \forma{-r}: \forma{teme-r, parti-r}.
\forma{falsedades}: \forma{fals-}: \forma{fals-o, fals-ificar};
\forma{-e-dad-}: \forma{sol-e-dad, ebri-e-dad}; \forma{-es}: \forma{lápic-es,
tabú-es}. \forma{gigantesca}: \forma{gigant-}: \forma{gigant-es,
a-gigant-a-r}; \forma{-esc-}: \forma{dant-esc-a, carnaval-esc-o}; \forma{-a}:
\forma{león-a, maestr-a}. \forma{entronización}: \forma{en-}:
\forma{en-amorar, en-casillaron}; \forma{-tron-}: \forma{des-tron-aron,
en-tron-izar}; \forma{-iz-}: \forma{real-iz-ó, esclav-iz-a-ron}; \forma{-a-}:
VT; \forma{-ción}: \forma{conten-ción, nega-ción}. \forma{pseudocientífico}:
\forma{pseudo-}: \forma{pseudo-Aristóteles, pseudo-efedrina}; \forma{-cient-}:
\forma{cient-ificismo, cient-ología}; \forma{-ífic-}: \forma{magn-ífic-o,
terror-ífic-as}; \forma{-o}: \forma{amarill-o, segur-o-s}.

c) \forma{con-venc-e-r}. Significado: \sig{mover con razones a alguien a hacer o
pensar algo}. Aporte semántico de cada morfema: \forma{con-}: no posee un
significado preciso en la palabra; tiene valor distintivo porque permite
distinguir el verbo \forma{convencer} de \forma{vencer}; \forma{-venc-}:
\sig{prevalecer sobre algo o alguien}; \forma{-e-}: VT; \forma{-r}: desinencia
de infinitivo. \forma{plen-a-mente}. Significado: \sig{de manera plena o
completa}. Aporte semántico de cada morfema: \forma{plen-}: \sig{lleno,
completo}; \forma{-a-}: no posee significado; es un resto etimológico del morfema
de género femenino que en estados previos de lengua marcaba la concordancia con
el sustantivo femenino \forma{mente}; \forma{-mente}: \sig{de manera
X}.\footnote{X remite al contenido del adjetivo: \textit{rápidamente} se
interpreta como \sig{de manera rápida} y \textit{alegremente}, \sig{de manera
alegre}.} \forma{ex-plic-a-n}. Significado: \sig{hacen comprender algo}. Aporte
semántico de cada morfema: \forma{ex-}: \sig{fuera}; \forma{-plic-}:
etimológicamente remite a la misma raíz que \forma{pleg-a-r}; la presencia del
prefijo \forma{ex-} produce la interpretación de conjunto \sig{desplegar hacia
afuera}, es decir, \sig{hacer conocer o comprender}; \forma{-a-}: VT;
\forma{-n}: desinencia de número y persona (amalgama, v.~§\ref{sec:04-vt}).

\textit{Parte B}. b) Opción correcta: 1. c) Opción correcta: 1.
\end{solucion}

\begin{solucion}{\ref{ej:04-alomorfos}}
Además del ejemplo resuelto, en el corpus se reconocen tres grupos de palabras
más:

1. a) En \forma{i-rrump-ido, rot-o, rupt-ura, i-rrup-ción, dis-rupt-iva,
romp-ió} se identifica el morfema \textsc{romp-}, que significa \sig{separar las
partes de un todo deshaciendo su unión}. Los alomorfos son \fonol{romp}
\textasciitilde{} \fonol{rot} \textasciitilde{} \fonol{rupt} \textasciitilde{}
\fonol{rup} \textasciitilde{} \fonol{rump}. b) El corpus no permite identificar
un factor que determine la aparición de cada variante. La causa de la alomorfía
es léxica.

2. a) En \forma{terra-planistas, terrá-queos, terr-ícola, en-tierr-an,
terr-itorio, a-terr-izar, Tierra} se identifica el morfema \textsc{terr-}, que
significa \sig{planeta que habitamos}. Los alomorfos son \fonol{ˈtjer}
\textasciitilde{} \fonol{ter}. b) La causa de la alomorfía es fonológica:
\fonol{ˈtjer} aparece en sílaba acentuada y \fonol{ter}, en sílaba no acentuada.

3. a) En \forma{terraplanista-s, cultural-es, terráqueo-s, dato-s,
falsedad-es, joven-es, maniquea-s, experto-s, (las) crisis} se identifica el
morfema \forma{-s}, que significa \sig{plural}. Los alomorfos son \fonol{s}
\textasciitilde{} \fonol{es} \textasciitilde{} ∅ (alomorfo cero). b) La causa de
la alomorfía es fonológica: \fonol{s} se adjunta a bases terminadas en vocal, con
excepción de las vocales cerradas acentuadas; y \fonol{es}, a bases terminadas en
vocal cerrada acentuada o en consonante, con excepción de las palabras graves y
esdrújulas terminadas en \fonol{s}, que quedan igual, por lo que en ellas se
manifiesta el alomorfo cero, sin contenido fonológico.
\end{solucion}

\begin{solucion}{\ref{ej:04-clases}}
\textit{Parte A}. a)

{\footnotesize
\begin{tabular}{@{}l*{5}{c}@{}}
  \toprule
  Palabra & \makecell{Pref.\\flexivo} & \makecell{Pref.\\derivativo} & Raíz & \makecell{Suf.\\derivativo} & \makecell{Suf.\\flexivo} \\
  \midrule
  \forma{terraplanista}     & & & terra- -plan- & -ista & \\
  \forma{negar}             & & & neg- & -a- & -r \\
  \forma{más}               & & & más & & \\
  \forma{achatada}          & & a- & -chat- & -a- -d- & -a \\
  \forma{gigantesca}        & & & gigant- & -esc- & -a \\
  \forma{plenamente}        & & & plen- & -a- -mente & \\
  \forma{representativos}   & & re- & -present- & -a- -t- -iv- & -o- -s \\
  \forma{explican}          & & ex- & -plic- & -a- & -n \\
  \forma{entronización}     & & en- & -tron- & -iz- -a- -ción & \\
  \forma{pseudocientíficos} & & pseudo- & -cient- & -ífic- & -o- -s \\
  \bottomrule
\end{tabular}}

b) El único morfema libre en el corpus es \forma{más}. El resto de los formantes
son ligados. \textit{Parte B}. b) Opción correcta: 3.
\end{solucion}

\begin{solucion}{\ref{ej:04-tema}}
\textit{Parte A}. a) \forma{culturales}: \forma{cult-ur-al-es} = \sig{relativos a
la cultura}; \forma{bromear}: \forma{brom-e-a-r} = \sig{hacer bromas};
\forma{desprecio}: \forma{des-precio} = \sig{falta de aprecio}.

b) \forma{chato}\textsubscript{A} > \forma{achata(r)}\textsubscript{V} >
\forma{achatada}\textsubscript{A};\footnote{\textit{Achatada} es el participio
del V \textit{achatar} que, por un proceso de conversión, puede funcionar
también como adjetivo (v.~§\ref{sec:04-conversion}).} A > V > Part/A. Adjetivo:
\forma{La Tierra es una esfera achatada}. \forma{chato}\textsubscript{A} >
\forma{achata(r)}\textsubscript{V} > \forma{achatada}\textsubscript{V part.}; A
> V > V\textsubscript{part.}. Participio: \forma{La chapa fue achatada con una
maza por el herrero}. \forma{ciencia}\textsubscript{N} >
\forma{cien(cia)tífico}\textsubscript{A} >
\forma{pseudocientífico}\textsubscript{A}; N > A > A.
\forma{culto}\textsubscript{A} > \forma{cult(o)ura}\textsubscript{N} >
\forma{cultur(a)ales}\textsubscript{A}; A > N > A.
\forma{broma}\textsubscript{N} > \forma{brom(a)ea(r)}\textsubscript{V}; N > V.
\forma{precio}\textsubscript{N} > \forma{desprecio}\textsubscript{N}; N >
N\textsubscript{neg}.

c) DE:
[[a-[chat]\textsubscript{A}-a]\textsubscript{V}-d-a]\textsubscript{A} (como
adjetivo); [a-[chat]\textsubscript{A}-a-d-a]\textsubscript{V} (como participio
verbal);
[pseudo-[[cient]\textsubscript{N}-ífic-o]\textsubscript{A}]\textsubscript{A};
[[[cult]\textsubscript{A}-ur]\textsubscript{N}-al]\textsubscript{A};
[[brom]\textsubscript{N}-e-a-r]\textsubscript{V};
[des-[precio]\textsubscript{N}]\textsubscript{N}.

\textit{Parte B}. b) Opción correcta: 3.
\end{solucion}

\begin{solucion}{\ref{ej:04-homonimia}}
2. \forma{inutilizable}: a. \sig{que se puede inutilizar}; b. \sig{que no se
puede utilizar}. 3. \forma{desarmar}: a. \sig{revertir el proceso de armar (un
discurso)}; b. \sig{quitarle el arma a alguien}.

2. a. \forma{útil}\textsubscript{A} > \forma{inútil}\textsubscript{A} >
\forma{inutilizar}\textsubscript{V} > \forma{inutilizable}\textsubscript{A};
[[[in-[útil]\textsubscript{A}]\textsubscript{A}-iza]\textsubscript{V}-ble]\textsubscript{A};
A > A > V > A. b. \forma{útil}\textsubscript{A} >
\forma{utilizar}\textsubscript{V} > \forma{utilizable}\textsubscript{A} >
\forma{inutilizable}\textsubscript{A};
[in-[[[útil]\textsubscript{A}-iza]\textsubscript{V}-ble]\textsubscript{A}]\textsubscript{A};
A > V > A > A.

3. a. \forma{armar}\textsubscript{V} > \forma{desarmar}\textsubscript{V};
[des-[arm-a-r]\textsubscript{V}]\textsubscript{V}; V > V. b.
\forma{arma}\textsubscript{N} > \forma{armar}\textsubscript{V} >
\forma{desarmar}\textsubscript{V}; [des-[arm]\textsubscript{N}-a-r]\textsubscript{V};
N > V > V.
\end{solucion}

\begin{solucion}{\ref{ej:04-ardo}}
\textit{Parte A}. b) \forma{-ard-}: aunque el \textit{DLE} sostiene que este
sufijo forma sustantivos y adjetivos con valor aumentativo o despectivo, en el
uso juvenil actual su empleo asigna \sig{valoración positiva por parte del
hablante}. c) Ejemplos posibles sugeridos: \forma{epic-ard-o, amig-ard-o,
cervez-ard-a, vin-ard-o, mes-ard-a, min-ard-a, tip-ard-o, te-s-ard-o}.

\textit{Parte B}. a) 2; b) 2; c) 1; d) 4. El aporte semántico de \forma{-ard-} se
asocia a significados expresivos vinculados con valoraciones del hablante.

\textit{Parte C}. Las características evaluadas en los puntos anteriores
permiten incluir el sufijo \forma{-ardo} entre los que pertenecen a la
morfología apreciativa.
\end{solucion}

\begin{solucion}{\ref{ej:04-derivacion}}
\textit{Parte A}. Grupo 2. a) \forma{falsedades, habilidad, deslealtad,
cotidianidad, bondades, maldad, disconformidad}. Proceso común: derivación con
el sufijo \forma{-dad}. b) Categoría de la base: A. Categoría resultante: N. c)
Valor semántico: \sig{propiedad de ser X}. d) Descripción fonológica:
\fonol{ˈdad}, \fonol{ˈtad}. Las vocales \fonol{e} e \fonol{i} en \fonol{ˈedad} e
\fonol{ˈidad} pueden considerarse interfijos cuya aparición se explica por
razones fonológicas.

Grupo 3. a) \forma{chateaba, bromear, desposteó, googlearon}. Proceso común:
derivación con el sufijo \forma{-e(a)-}. b) Categoría de la base: N. Categoría
resultante: V. c) Valor semántico: \sig{hacer X}. d) Descripción fonológica:
\fonol{eˈa}.

Grupo 4. a) \forma{deslealtad, desinformó, desposteó, desagradables, desamor,
disconformidad}. Proceso común: derivación con el prefijo \forma{des-}. b)
Categorías de la base: A\textsubscript{x}, N\textsubscript{y},
V\textsubscript{z}. Categoría resultante: A\textsubscript{x},
N\textsubscript{y}, V\textsubscript{z} (los índices señalan que las bases no
cambian su categoría). c) Valor semántico: \sig{inversión / negación / privación
/ contrariedad respecto de X}. d) Descripción fonológica: \fonol{ˈdes},
\fonol{ˈdis}.

\textit{Parte B}. Opción correcta: 2.
\end{solucion}

\begin{solucion}{\ref{ej:04-composicion}}
Grupo 2. a) \forma{sacacorchos, lavarropas, guardapolvos, pisapapeles,
pelapapas, cazarrecompensas, matabichos, cometierra, chupamedias,
limpiaparabrisas}. b) \forma{saca-corchos; lava-ropas; guarda-polvos;
pisa-papeles; pela-papas; caza-recompensas; mata-bichos; come-tierra;
chupa-medias; limpia-parabrisas}. c)
[[saca]\textsubscript{V}[corchos]\textsubscript{N}]\textsubscript{N}. d) V + N >
N.

Grupo 3. a) \forma{sordomudo, socioeconómico, blanquiceleste, dentoalveolar,
italoargentino, agridulce}. b) \forma{sordo-mudo; socio-económico;
blanqu-i-celeste; dent-o-alveolar; ital-o-argentino; agri-dulce}. c)
[[sordo]\textsubscript{A}[mudo]\textsubscript{A}]\textsubscript{A}. d) A + A > A.
\end{solucion}

\begin{solucion}{\ref{ej:04-filos}}
\textit{Parte A}. \forma{filosofía}:
[[filo]\textsubscript{TC}[sofía]\textsubscript{TC}]\textsubscript{N}, TC + TC >
N; \forma{melófila}: [[meló]\textsubscript{TC}[fil]\textsubscript{TC}a]\textsubscript{A},
TC + TC > A; \forma{filonazi}:
[[filo]\textsubscript{TC}[nazi]\textsubscript{N}]\textsubscript{A}, TC + N > A;
\forma{hidrófilo}:
[[hidró]\textsubscript{TC}[fil]\textsubscript{TC}o]\textsubscript{A}, TC + TC >
A; \forma{americanófila}:
[[american]\textsubscript{N}ó[fil]\textsubscript{TC}a]\textsubscript{A}, N + TC >
A; \forma{filología}:
[[filo]\textsubscript{TC}[logía]\textsubscript{TC}]\textsubscript{N}, TC + TC >
N; \forma{tecnófilo}:
[[tecno]\textsubscript{TC}[fil]\textsubscript{TC}o]\textsubscript{A}, TC + TC >
A; \forma{filántropo}:
[[fil]\textsubscript{TC}[antrop]\textsubscript{TC}o]\textsubscript{N}, TC + TC >
N.

\textit{Parte B}. Opciones correctas: a) 3; b) 1.
\end{solucion}

\begin{solucion}{\ref{ej:04-parasintesis}}
\pendiente{El solucionario del original no corresponde a las consignas de este
ejercicio (analiza afijos y no bases parasintéticas); revisar}
Se ofrecen ejemplos de respuestas posibles a lo pedido a modo ilustrativo.
\forma{terra-plan-ista}. \forma{-ista}: a) \forma{tech-ista, social-ista}; b)
categoría de la base: N/A; c) categoría del resultado: N; N/A + \forma{-ista} >
N; d) valor semántico: \sig{que profesa / practica / cree en X}.
\forma{sub-terr-áne-a}. \forma{sub-}: a) \forma{sub-suelo, sub-valorar}; b)
categoría de la base: N/A/V; c) categoría del resultado: N/A/V;
X\textsubscript{i} + \forma{sub-} > X\textsubscript{i}; d) valor semántico:
\sig{debajo de, menos que X}. \forma{-áne-}: a) \forma{moment-áne-o,
cut-áne-a}; b) categoría de la base: N; c) categoría del resultado: A; N +
\forma{-áne-} > A; d) valor semántico: \sig{relacionado con / relativo a X}.
\forma{mil-eur-ista-s}. \forma{-ista} (v. \forma{terraplanista}).
\forma{sal-piment-a-r}. \forma{-a-}: a) \forma{lider-a-r, activ-a-r}; b)
categoría de la base: N/A; c) categoría del resultado: V; N/A + \forma{-a-} > V;
d) valor semántico: \sig{usar X}. \forma{in-fum-a-ble}. \forma{in-}: a)
\forma{in-moral, in-útil}; b) categoría de la base: A; c) categoría del
resultado: A; \forma{in-} + A > A; d) valor semántico: \sig{no X}.
\forma{-ble}: a) \forma{cont-a-ble, tem-i-ble-s}; b) categoría de la base: V; c)
categoría del resultado: A; V + \forma{-ble} > A; d) valor semántico: \sig{que
se puede X}.
\end{solucion}

\begin{solucion}{\ref{ej:04-conversion}}
En todos los casos las formas comparten la raíz y lo que varía es el tema.
2. \forma{com-i-d-a}. a) \forma{comida}\textsubscript{a}: tema \forma{comida-};
\forma{comida}\textsubscript{b}: tema \forma{comi-}. b)
\forma{comida}\textsubscript{a} es sustantivo porque va precedido del
determinante \forma{la}; \forma{comida}\textsubscript{b} es una forma verbal de
participio del verbo \forma{comer} que integra una perífrasis de voz pasiva.
3. \forma{curs-a-d-a}. a) \forma{cursada}\textsubscript{a}: tema
\forma{cursada-}; \forma{cursada}\textsubscript{b}: tema \forma{cursa-}. b)
\forma{cursada}\textsubscript{a} funciona como sustantivo porque va precedido
del determinante \forma{la}; \forma{cursada}\textsubscript{b} es la forma verbal
de participio del verbo \forma{cursar} empleada como adjetivo (por conversión)
que modifica al sustantivo \forma{notificación}.\pendiente{Revisar temas: el
original los invierte en a y b}
4. \forma{priv-a-d-o}. a) \forma{privado}\textsubscript{a}: tema
\forma{privad-}; \forma{privado}\textsubscript{b}: tema \forma{priva-}. b)
\forma{privado}\textsubscript{a} funciona como un adjetivo que modifica a
\forma{empleo}; \forma{privado}\textsubscript{b} es una forma verbal de
participio del verbo \forma{privar} que integra una perífrasis de voz pasiva.
5. \forma{doctor-a-ndo}. a) \forma{doctorando}\textsubscript{a}: tema
\forma{doctora-}; \forma{doctorando}\textsubscript{b}: tema
\forma{doctorando-}. b) \forma{doctorando}\textsubscript{a} es la forma verbal
de gerundio del verbo \forma{doctorarse} que integra una perífrasis verbal
progresiva de \forma{estar} + gerundio; \forma{doctorando}\textsubscript{b}
funciona como un sustantivo porque va precedido por el determinante \forma{los}.
6. \forma{cumbia}. a) \forma{cumbia}\textsubscript{a}: tema \forma{cumbia-};
\forma{cumbia}\textsubscript{b}: tema \forma{cumbia-}. b)
\forma{cumbia}\textsubscript{a} funciona como un sustantivo precedido por el
determinante \forma{esta}; \forma{cumbia}\textsubscript{b} ejemplifica un proceso
de conversión similar al del inglés, en el que los sustantivos pueden dar origen
a verbos de significado relacionado.
\end{solucion}

\begin{solucion}{\ref{ej:04-regresiva}}
2. \forma{espectar}. a) Base: \forma{espectador-}. Categoría: N.\pendiente{El
original dice «A»} Resultado: \forma{espect(a)-}. Categoría: V. Segmento
afectado: \forma{-dor}. b) \forma{espectar} < \forma{espectador}: N --
\forma{-dor} > V. c) Significado: \sig{participar de un juego como espectador}.
3. \forma{alfabeto}. a) Base: \forma{analfabet-}. Categoría: A. Resultado:
\forma{alfabet-}. Categoría: A. Segmento afectado: \forma{an-}. b)
\forma{alfabeto} < \forma{analfabeto}: A -- \forma{an-} > A. c) Significado:
\sig{que no es analfabeto}. 4. \forma{antitranspirar}. a) Base:
\forma{antitranspirante-}. Categoría: A. Resultado: \forma{antitranspir(a)-}.
Categoría: V. Segmento afectado: \forma{-nte}. b) \forma{antitranspirar} <
\forma{antitranspirante}: A -- \forma{-nte} > V. c) Significado: \sig{acción
contraria a transpirar}.
\end{solucion}

\begin{solucion}{\ref{ej:04-acortamiento}}
Grupo 2: acortamientos por apócope. a) \forma{profe, facu, auto, moto, compu,
taxi}. b) \forma{profe}\textsubscript{N} < \forma{profesor/a}\textsubscript{N};
\forma{facu}\textsubscript{N} < \forma{facultad}\textsubscript{N};
\forma{auto}\textsubscript{N} < \forma{automóvil}\textsubscript{N};
\forma{moto}\textsubscript{N} < \forma{motocicleta}\textsubscript{N};
\forma{compu}\textsubscript{N} < \forma{computadora}\textsubscript{N};
\forma{taxi}\textsubscript{N} < \forma{taxímetro}\textsubscript{N}. c) Se
conservan las primeras dos sílabas de la base. El resultado, al igual que en el
caso del grupo 1, tiene acentuación grave. d) El acortamiento por apócope
tampoco produce cambio alguno en el significado de la base.
\end{solucion}

\begin{solucion}{\ref{ej:04-cruces}}
2. \forma{Conmebol}. a) Bases: \forma{confederación + sudamericana + (de)
fútbol}. b) Categoría de las bases: N + A + N. Categoría del resultado: N. c)
Proceso de selección: primera sílaba de la primera base, tercera sílaba de la
segunda base, última sílaba de la última base: [con(federación) +
(suda)me(ricana) + (fút)bol]. 3. \forma{Mercosur}. a) Bases: \forma{mercado +
común + (del) sur}. b) Categoría de las bases: N + A + N. Categoría del
resultado: N. c) Proceso de selección: primera sílaba de cada base; la última de
las bases es monosilábica y se integra completa: [mer(cado) + co(mún) + (del)
sur]. 4. \forma{vitalitud}. a) Bases: \forma{vitalidad + actitud}. b) Categoría
de las bases: N + N. Categoría del resultado: N. c) Proceso de selección:
primeros seis segmentos sonoros de la primera base y cuatro segmentos finales de
la segunda, aprovechando el fonema vocálico común: [vital(idad) + (acti)tud]. 5.
\forma{deshamburguesar}. a) Bases: \forma{desaburguesar + hamburguesa}. b)
Categoría de las bases: V + N. Categoría del resultado: V. c) Proceso de
selección: primeros cuatro segmentos sonoros de la primera base y ocho segmentos
finales de la segunda, aprovechando la coincidencia en los últimos siete
fonemas: [desa(burguesar) + (hamburgues)a]. En la publicidad, el verbo
\forma{deshamburguesar} aparece en la forma imperativa de la 2.ª persona
singular, por lo que no se expresan morfemas flexivos, facilitando la
coincidencia entre ambas bases. En otras formas, \forma{deshamburguesar} debe
incorporar la morfología flexiva correspondiente a la 1.ª conjugación verbal,
ausente en el sustantivo \forma{hamburguesa}. 6. \forma{artivismo}. a) Bases:
\forma{arte + activismo}. b) Categoría de las bases: N + N. Categoría del
resultado: N. c) Proceso de selección: primeros tres segmentos sonoros de la
primera base y siete segmentos finales de la segunda, aprovechando el fonema
consonántico común: [art(e) + (ac)tivismo]. 7. \forma{juernes}. a) Bases:
\forma{jueves + viernes}. b) Categoría de las bases: N + N. Categoría del
resultado: N. c) Proceso de selección: primeros tres segmentos sonoros de la
primera base y cuatro segmentos finales de la segunda: [jue(ves) + (vie)rnes].
\end{solucion}

\begin{solucion}{\ref{ej:04-acronimos}}
En todos los casos el acrónimo equivale referencialmente al sintagma que
abrevia. 2. \forma{RAE}. a) Real\textsubscript{A} Academia\textsubscript{N}
Española\textsubscript{A}. b) Categoría de la construcción: N. c) No posee marcas
morfológicas adicionales. 3. \forma{ovnis}. a) objetos\textsubscript{N}
voladores\textsubscript{A} no\textsubscript{IP}\footnote{La sigla IP corresponde
a \textit{índice de polaridad}, en este caso, negativa. El concepto de
polaridad se desarrolla en el \capref{14}.} identificados\textsubscript{A}. b)
Categoría del resultado: N. c) El segmento final corresponde con la
pluralización, que se manifiesta por medio del alomorfo correspondiente a la
terminación del tema (\forma{-s}). 4. \forma{FIFA / AFA}. a)
Federación\textsubscript{N} Internacional\textsubscript{A} (de)
Fútbol\textsubscript{N} Asociado\textsubscript{A} / Asociación\textsubscript{N}
(del) Fútbol\textsubscript{N} Argentino\textsubscript{A}. b) Categoría de los
resultados: N. c) No poseen marcas morfológicas adicionales. 5. \forma{UNA}. a)
Universidad\textsubscript{N} Nacional\textsubscript{A} (de las)
Artes\textsubscript{N}. b) Categoría del resultado: N. c) No posee marcas
morfológicas adicionales.
\end{solucion}

\begin{solucion}{\ref{ej:04-siempre}}
a) Categoría de las bases: A. Categoría del resultado: A. b) El segmento
\forma{in-} con que comienzan todas estas palabras se reinterpreta como si fuera
el prefijo negativo \forma{in-}. c) En \forma{imberbe}, \forma{im-} representa
al prefijo negativo, ya que la palabra significa \sig{que no tiene barba}; no
obstante, \forma{berbe} constituye un caso de retroformación, dado que esta
palabra no existe en español y se formaría a partir de \forma{imberbe}.
\end{solucion}

\begin{solucion}{\ref{ej:04-reanalisis}}
Los tres son casos de reanálisis. 2. \forma{dictablanda}. Se trata de la
reinterpretación de la desinencia \forma{-dura} de \forma{dictadura}, formante de
sustantivos abstractos (\forma{gordura, dulzura}), como si se tratara del
adjetivo \forma{duro/a}, lo que permite la formación de un supuesto opuesto:
\forma{dictablanda}. El término se usó para hacer referencia al gobierno del
general Berenguer en España, ocurrido tras la dictadura de Primo de Rivera. Como
el gobierno de Berenguer se inició con algunas medidas liberalizadoras, se lo
denominó \forma{dictablanda}. 3. \forma{precuela}: versión española del inglés
\forma{prequel}, formado a partir de \forma{sequel} \sig{secuela}, derivado del
verbo latino \forma{sequor} (\sig{seguir}), sobre una base inexistente
\juicio{*}{\forma{-quel}} a la que se le añade el prefijo \forma{pre-} para
hacer referencia a una «obra literaria o cinematográfica que cuenta hechos que
preceden a los de otra obra ya existente» (\textit{DLE}, en línea).
\end{solucion}

\begin{solucion}{\ref{ej:04-genero}}
\textit{Parte A}. 2. \forma{el sartén / la sartén}. 3. \forma{la gat-a / el
gat-o}. 4. \forma{la cebra} / ∅. 5. \forma{el naranj-o / la naranj-a}. 6.
\forma{la cuchill-a / el cuchill-o}. 7. \forma{la araña} / ∅. 8. \forma{la
mano} / ∅. 9. \forma{el clima} / ∅. 10. \forma{el alcalde / la alcald-esa}. 11.
\forma{el amante / la amante}. 12. \forma{el tilde / la tilde}. 13. \forma{la
bols-a / el bols-o}. 14. \forma{el calor / la calor}.

\textit{Parte B}. a) Opción correcta: 4. Sustantivos como \forma{cebra},
\forma{araña}, \forma{mano} y \forma{clima} poseen género inherente. A los
sustantivos como \forma{cebra} y \forma{araña}, que abarcan indistintamente
individuos machos o hembras de la clase de animal denotada, se los conoce como
epicenos. b) Opción correcta: 3. Además de la oposición en términos de
macho/hembra para sustantivos que designan entidades sexuadas (como
\forma{secretario/secretaria}), la categoría también abarca otros contrastes de
significado, como árbol/fruto (\forma{naranjo/naranja}) o tamaño
(\forma{cuchillo/cuchilla, bolso/bolsa}). Asimismo, la selección de género
inherente en algunos nombres puede estar sujeta a variación socioestilística
(\forma{el/la sartén, el/la tilde}). c) Opción correcta: 2.
\end{solucion}

\begin{solucion}{\ref{ej:04-verbos}}
a) 3. A. \forma{lleg-a-ba-n}; B. \forma{com-í-a-n}; C. \forma{aburr-í-a-n}. 4.
A. \forma{lleg-ó}; B. \forma{com-i-ó}; C. \forma{aburr-i-ó}. 5. A.
\forma{lleg-a-ba-s}; B. \forma{com-í-a-s}; C. \forma{aburr-í-a-s}. 6. A.
\forma{lleg-a-ra}; B. \forma{com-ie-ra}; C. \forma{aburr-ie-ra}.

b) Contraste 1/3: oposición de número: 1) singular / 3) plural. Contraste 1/4:
oposición de aspecto: 1) imperfecto / 4) perfecto. Contraste 1/5: oposición de
persona: 1) 1.ª/3.ª / 5) 2.ª. Contraste 1/6: oposición de modo: 1) indicativo /
6) subjuntivo.

c) 3. 3.ª persona plural del pretérito imperfecto del indicativo. 4. 3.ª persona
singular del pretérito perfecto simple del indicativo. 5. 2.ª persona singular
del pretérito imperfecto del indicativo. 6. 1.ª/3.ª persona singular del
pretérito imperfecto del subjuntivo.

{\footnotesize
\begin{tabularx}{\linewidth}{@{}l L L L l@{}}
  \toprule
  & Raíz & VT & TAM & Conc. \\
  \midrule
  1. \forma{llegaba/comía/aburría}     & lleg- / com- / aburr- & -a- / -í- & -ba / -a & ∅ \\
  2. \forma{llega/come/aburre}         & lleg- / com- / aburr- & -a / -e   &          & ∅ \\
  3. \forma{llegaban/comían/aburrían}  & lleg- / com- / aburr- & -a- / -í- & -ba / -a & -n \\
  4. \forma{llegó/comió/aburrió}       & lleg- / com- / aburr- & ∅ / -i-   & -ó       & ∅ \\
  5. \forma{llegabas/comías/aburrías}  & lleg- / com- / aburr- & -a- / -í- & -ba / -a & -s \\
  6. \forma{llegara/comiera/aburriera} & lleg- / com- / aburr- & -a- / -ie- & -ra     & ∅ \\
  \bottomrule
\end{tabularx}}
\end{solucion}

% =====================================================================
%  capitulos/cap05.tex — Capítulo 5. Los constituyentes y la estructura
%  sintagmática
%  Autora: María Paula Bonorino
%  Fuente: Capitulo_05_normalizado.docx
% =====================================================================
\capitulo[Los constituyentes y la estructura sintagmática]%
  {Los constituyentes y la estructura sintagmática}{María Paula Bonorino}

\section{Introducción}

Este capítulo inicia el estudio de la sintaxis y para ello introduce una serie
de conceptos clave para reflexionar sobre las relaciones entre las palabras que
forman las oraciones, ya que, al igual que las unidades de la morfología, los
ítems léxicos no solo se ordenan de manera lineal, sino que también establecen
dependencias jerárquicas. En primer lugar, consideramos cómo se definen los
constituyentes sintácticos y presentamos las pruebas formales que permiten
identificarlos. En segundo lugar, nos centramos en las clases de palabras
léxicas, entendidas como aquellas que poseen significado descriptivo inherente,
y examinamos cómo se expanden en sintagmas, conformando así unidades intermedias
de análisis, que exhiben una organización interna. Finalmente, nos enfocamos en
las características de los diferentes tipos de sintagmas ---nominales (SSNN),
adjetivales (SSAA), adverbiales (SSADV) y preposicionales
(SSPP)\footnote{La duplicación de las siglas (SSNN, SSAA, etc.) indica
plural.}--- y consideramos sus principales propiedades sintácticas y
semánticas.

\section{El concepto de constituyente y la estructura jerárquica}\label{sec:05-constituyente}

Así como los morfemas se unen para formar palabras, también las palabras se
combinan en unidades mayores. Consideremos, por ejemplo, la siguiente oración:

\begin{exe}
  \ex\label{ex:05-ministro} El ministro insistió en su declaración.
\end{exe}

Si bien en \exref{ex:05-ministro} las palabras se vinculan linealmente, también
podemos reconocer ---al menos intuitivamente--- que la relación entre
\forma{ministro} y \forma{el} es más estrecha que la que existe entre
\forma{ministro} e \forma{insistió}, aunque estas dos últimas sean palabras
contiguas. La diferencia radica en que \forma{ministro} y \forma{el} están en el
mismo constituyente e \forma{insistió}, en cambio, se encuentra en otro. Pero
¿qué es un constituyente? Un \termino{constituyente} es un grupo de palabras que
funciona como una unidad «independientemente de si tiene una estructura interna
simple o compleja» \parencite[88]{camacho2018}. Así, en
\exref{ex:05-AB}, tenemos dos constituyentes mayores a los que llamaremos A y B:

\begin{exe}
  \ex\label{ex:05-AB} {}[El ministro]\etq{A} [insistió en su declaración]\etq{B}
\end{exe}

\begin{destacado}
El concepto de constituyente es relacional: «algo es un constituyente de una
determinada construcción si forma parte de esta\ldots{}»
\parencite[69]{ditullio2014}.
\end{destacado}

En línea con lo que plantea Di Tullio, un constituyente puede, a su vez, estar
integrado por constituyentes menores. El constituyente A está formado por el
determinante \forma{el} y el sustantivo \forma{ministro}.

\begin{exe}
  \ex {}[\etq{A} El ministro]
\end{exe}

El constituyente B, por su parte, plantea una estructura más compleja. Cabe
reconocer allí, en primer lugar, un constituyente formado por la combinación del
sustantivo \forma{declaración} con el posesivo \forma{su}. Este constituyente se
combina, a su vez, con la preposición \forma{en} para formar el constituyente
\forma{en su declaración}. Finalmente, este constituyente se combina con el verbo
\forma{insistió}, dando como resultado el constituyente B, \forma{insistió en su
declaración}.

\begin{exe}
  \ex {}[\etq{B} insistió [en [su declaración]]]
\end{exe}

\begin{destacado}
En síntesis, los constituyentes de una oración se relacionan entre sí de forma
jerárquica y no simplemente por contigüidad.
\end{destacado}

La estructura de constituyentes se puede representar mediante corchetes, como en
\exref{ex:05-corchetes}.

\begin{exe}
  \ex\label{ex:05-corchetes} {}[\etq{AB} [\etq{A} El ministro] [\etq{B} insistió
      [en [su declaración]]]]
\end{exe}

\begin{ejercicio}\label{ej:05-segmentar}
Segmente en constituyentes las siguientes oraciones, como se muestra en el
ejemplo resuelto~1.

\begin{exe}
  \ex Sus conclusiones asombraron a los científicos.
  \ex Nuestra tía traía regalos constantemente.
  \ex Muchos estudiantes prefieren la virtualidad.
  \ex Las medidas preocupan a la población.
  \ex Los jugadores ganaron la copa en 2022.
  \ex Aquellas profesoras incorporaron en sus clases herramientas tecnológicas.
\end{exe}

\begin{resuelto}
1. [[Sus conclusiones [sobre [el tema]]] [asombraron [a [la comunidad
científica]]]]\pendiente{La resolución no corresponde a la oración 1: «Sus
conclusiones asombraron a los científicos»}
\end{resuelto}
\end{ejercicio}

\subsection{Homonimia estructural}\label{sec:05-homonimia}

En algunos casos,\footnote{Este tema se presentó en el \capref{2},
§\ref{sec:02-ambiguedad}.} como se muestra en \exref{ex:05-salon}, la
segmentación en constituyentes puede ofrecer más de una alternativa
interpretativa: así, puede ser que \forma{el salón con flores} forme un solo
constituyente \exref{ex:05-salon-b} o que se trate de dos constituyentes
independientes, que pueden intercambiarse \exref{ex:05-salon-c}:

\begin{exe}
  \ex\label{ex:05-salon}
    \begin{xlist}
      \ex Redecoró el salón con flores.
      \ex\label{ex:05-salon-b} Redecoró [el salón [con flores]].
      \ex\label{ex:05-salon-c} Redecoró [el salón] [con flores]. / Redecoró [con
          flores] [el salón].
    \end{xlist}
\end{exe}

Las oraciones como \exref{ex:05-salon}, que admiten más de una segmentación en
constituyentes, son ejemplos de \termino{homonimia estructural}, el fenómeno
según el cual una oración presenta más de una descripción estructural (DE), cada
una asociada a una interpretación diferente de la oración. Así, en
\exref{ex:05-salon-b}, [\forma{el salón con flores}] constituye una unidad porque
hace referencia a \sig{un salón que ya tenía flores, por ejemplo, en el
empapelado}, y que fue redecorado, probablemente con otro diseño. En este caso,
[\forma{con flores}] es un constituyente menor que modifica a [\forma{el salón}],
por lo que, como constituyen una unidad, no admiten intercambiar sus posiciones.
En \exref{ex:05-salon-c}, en cambio, en la que interpretamos que \sig{puso flores
en el salón}, se puede considerar que [\forma{el salón}] y [\forma{con flores}]
son dos modificadores independientes del verbo \forma{redecoró} y que, por lo
tanto, sus posiciones son intercambiables.

\begin{destacado}
Dado que una misma oración puede recibir más de una descripción estructural,
según como se interpreten sus constituyentes, se considera que las secuencias
ambiguas establecen un argumento en favor de la teoría de que las construcciones
sintácticas poseen una estructura jerárquica.
\end{destacado}

\begin{ejercicio}\label{ej:05-homonimia}
Considere las siguientes oraciones y resuelva lo pedido, según se muestra en el
ejemplo resuelto~1.

\begin{consignas}
  \item Segmente en constituyentes.
  \item Proporcione las dos interpretaciones posibles de cada oración y explique
        a qué se deben.
  \item Luego de responder a los puntos a) y b), repase el
        ejercicio~\ref{ej:02-ambiguedad} del \capref{2} y considere qué efecto
        produce la doble posibilidad de segmentación que ofrecen estas
        oraciones.
\end{consignas}

\begin{exe}
  \ex Mi hermano trajo una foto de Egipto.
  \ex Llenó la bolsa de papel.
  \ex Encontraron el juguete roto.
\end{exe}

\begin{resuelto}
1. a) i) [[Mi hermano] [trajo [una foto] [de Egipto]]] / ii) [[Mi hermano]
[trajo [una foto de Egipto]]]. b) En i) se puede considerar que [una foto] y
[de Egipto] son dos modificadores independientes del verbo \forma{trajo}. En
ii), [de Egipto] modifica a [una foto], por lo que conforman un solo
constituyente.
\end{resuelto}
\end{ejercicio}

\subsection{Identificación de constituyentes}\label{sec:05-pruebas}

Para el reconocimiento de constituyentes se pueden aplicar pruebas que permiten
identificarlos. Las más habituales son las de sustitución, movimiento e
interrogación. Veamos en qué consiste cada una de ellas.

\subsubsection{Sustitución}\label{sec:05-sustitucion}

Toda unidad sintáctica compleja, como la secuencia subrayada en
\exref{ex:05-conicet-a}, puede concebirse como el desarrollo o expansión de otra
más simple por la que puede ser sustituida, como en \exref{ex:05-conicet-b} y
\exref{ex:05-conicet-c}, y toda unidad simple puede dar lugar a una expansión en
un constituyente más complejo, como se muestra en \exref{ex:05-conicet-d} y
\exref{ex:05-conicet-e}:

\begin{exe}
  \ex\label{ex:05-conicet}
    \begin{xlist}
      \ex\label{ex:05-conicet-a} \subrayado{Los investigadores del CONICET}
          analizaron con gran cuidado la nueva variante del virus.
      \ex\label{ex:05-conicet-b} \subrayado{Los investigadores} analizaron con
          gran cuidado la nueva variante del virus.
      \ex\label{ex:05-conicet-c} \subrayado{Ellos} analizaron con gran cuidado la
          nueva variante del virus.
      \ex\label{ex:05-conicet-d} \subrayado{Los investigadores más renombrados
          del CONICET} analizaron con gran cuidado la nueva variante del virus.
      \ex\label{ex:05-conicet-e} \subrayado{Los investigadores del CONICET
          (Consejo Nacional de Investigaciones Científicas y Técnicas)}
          analizaron con gran cuidado la nueva variante del virus.
    \end{xlist}
\end{exe}

\begin{ejercicio}\label{ej:05-expansion}
Proponga expansiones para los constituyentes subrayados, según se muestra en el
ejemplo~1.

\begin{exe}
  \ex \subrayado{Susana} firmó \subrayado{los papeles}.
  \ex \subrayado{El museo} recibió \subrayado{los ejemplares}.
  \ex \subrayado{Mis primas} \subrayado{llegaron}.
  \ex \subrayado{La temporada} trajo \subrayado{sorpresas}.
  \ex \subrayado{Los vehículos} circulaban \subrayado{lentamente}.
\end{exe}
\pendiente{Los subrayados no se conservan en el original; se reconstruyeron a
partir de las respuestas}

\begin{resuelto}[Ejemplo 1. Respuesta posible]
Susana, la directora del instituto, firmó los papeles necesarios para la
presentación del proyecto.
\end{resuelto}
\end{ejercicio}

\begin{ejercicio}\label{ej:05-sustitucion}
Sustituya en la oración los constituyentes subrayados. Proporcione para cada uno
de ellos una versión mediante sustitución por reducción y otra por
pronominalización, como se muestra para el ejemplo~1.

\begin{exe}
  \ex \subrayado{Los investigadores del CONICET} analizaron \subrayado{con gran
      cuidado} \subrayado{la nueva variante del virus}.
  \ex \subrayado{Varios ministros del gabinete} firmaron \subrayado{un acuerdo
      bilateral} \subrayado{con el embajador de la Unión Europea}.
\end{exe}

\begin{resuelto}
1. Reducción: Los investigadores analizaron con cuidado la variante.
Pronominalización: Ellos la analizaron así.
\end{resuelto}
\end{ejercicio}

\begin{ejercicio}\label{ej:05-encorchetar}
Considere las siguientes oraciones y resuelva las consignas a continuación,
según se muestra en el ejemplo resuelto~1.

\begin{consignas}
  \item Encorchete los constituyentes.
  \item Aplique la prueba de sustitución a todos los constituyentes
        encorchetados en a).
\end{consignas}

\begin{exe}
  \ex Sus tías enviaron los regalos por correo.
  \ex Los egresados buscaron con entusiasmo un salón de fiestas.
  \ex Nuestros candidatos presentaron sus propuestas a los contribuyentes.
  \ex Los fuertes temporales provocaron inundaciones en toda la ciudad.
\end{exe}

\begin{resuelto}
1. a) [[Sus tías] [enviaron [los regalos] [por correo]]]. b) Ellas enviaron los
regalos por correo. Sus tías los enviaron por correo. Sus tías enviaron los
regalos así. Ellas los enviaron así.
\end{resuelto}
\end{ejercicio}

\subsubsection{Interrogación}\label{sec:05-interrogacion}

Solo los constituyentes pueden funcionar como respuesta a preguntas formuladas
con una palabra interrogativa como \forma{qué, quién, cómo, cuándo}, etc. Estas
preguntas parciales\footnote{Las preguntas pueden ser totales o parciales:
mientras las parciales se responden mediante un constituyente de la oración,
como se muestra en §\ref{sec:05-interrogacion}, las preguntas totales son las
que interrogan por todo el enunciado y se responden afirmativa o negativamente:
\forma{¿Querés venir hoy al cine conmigo? Sí / No (quiero)}.} tienen que
responderse mediante algún constituyente específico de la oración, lo que
permite probar su condición de tal. Así, por ejemplo, respecto de
\exref{ex:05-conicet-a}, podríamos preguntar:

\begin{exe}
  \ex\label{ex:05-preguntas}
    \begin{xlist}
      \ex \forma{¿Quiénes} analizaron con gran cuidado la nueva variante del
          virus? Los investigadores del CONICET.
      \ex \forma{¿Qué} hicieron los investigadores del CONICET? Analizaron con
          gran cuidado la nueva variante del virus.
      \ex \forma{¿Cómo} analizaron los investigadores del CONICET la nueva
          variante del virus? Con gran cuidado.
      \ex \forma{¿Qué} analizaron los investigadores del CONICET con gran
          cuidado? La nueva variante del virus.
    \end{xlist}
\end{exe}

Como se advierte en los ejemplos anteriores, cada pregunta parcial recibe como
respuesta un constituyente de la oración.

\begin{ejercicio}\label{ej:05-preguntas}
¿Qué preguntas podríamos hacer para obtener cada uno de los constituyentes
encorchetados en la descripción estructural de las siguientes oraciones?

\begin{exe}
  \ex {}[El cartero] [trajo [a mi casa] [la nueva tarjeta de crédito]]
  \ex {}[Mis primos] [llegaron [ayer por la tarde] [a la ciudad] [para
      visitarnos]]
  \ex {}[La nueva profesora] es [muy exigente] [con los exámenes]
\end{exe}

\begin{resuelto}
1. ¿Quién trajo a mi casa la nueva tarjeta de crédito? / ¿Qué hizo el cartero? /
¿Adónde trajo el cartero la nueva tarjeta de crédito? / ¿Qué trajo a mi casa el
cartero?
\end{resuelto}
\end{ejercicio}

\begin{ejercicio}\label{ej:05-bioy}
Recurra a la prueba de la interrogación para verificar que las partes subrayadas
son constituyentes.

% Fragmento literario protegido: se incorpora a mano desde el original.
\textoaincorporar{el comienzo de «La salvación», de A.~Bioy Casares
(\textit{La invención y la trama}, Tusquets, 1991), con los constituyentes
subrayados: \textit{el escultor; con el tirano; por los jardines del palacio;
más allá del laberinto para los extranjeros ilustres, en el extremo de la
alameda de los filósofos decapitados; su última obra}}

\begin{resuelto}
¿Quién paseaba con el tirano por los jardines del palacio?
\end{resuelto}
\end{ejercicio}

\subsubsection{Movimiento}\label{sec:05-movimiento}

Los constituyentes se desplazan en bloque, por lo que solo un constituyente
completo puede moverse a otra posición dentro de la oración. En
\exref{ex:05-conicet-a}, que repetimos aquí como \exref{ex:05-mov-a}, los dos
constituyentes mayores pueden intercambiar sus posiciones, dando lugar a
\exref{ex:05-mov-b}:

\begin{exe}
  \ex\label{ex:05-mov}
    \begin{xlist}
      \ex\label{ex:05-mov-a} {}[[Los investigadores del CONICET] [analizaron con
          gran cuidado la nueva variante del virus]].
      \ex\label{ex:05-mov-b} {}[[Analizaron con gran cuidado la nueva variante
          del virus] [los investigadores del CONICET]].
    \end{xlist}
\end{exe}

Al realizar este movimiento, la oración sigue siendo gramatical. Asimismo,
también se puede hacer la misma operación con otros constituyentes presentes en
la oración:

\begin{exe}
  \ex {}[Con gran cuidado] [los investigadores del CONICET] [analizaron [la nueva
      variante del virus]].
\end{exe}

Sin embargo, si se delimita erróneamente un constituyente, el movimiento puede
dar lugar a una oración agramatical:

\begin{exe}
  \ex[*]{\label{ex:05-mov-agram} La nueva variante [los investigadores del
      CONICET] [analizaron [del virus] [con gran cuidado]].}
\end{exe}

En \exref{ex:05-mov-agram}, el constituyente \forma{la nueva variante del virus}
se ha fragmentado, lo que da como resultado una oración agramatical en español.

\begin{ejercicio}\label{ej:05-pares}
Compare los siguientes pares de oraciones y luego, según se muestra en el
ejemplo resuelto~1:

\begin{consignas}
  \item indique mediante encorchetamiento si las expresiones en cursiva
        conforman o no un único constituyente;
  \item fundamente la respuesta a partir de las pruebas de reconocimiento
        estudiadas.
\end{consignas}

\begin{exe}
  \ex
    \begin{xlist}
      \ex Le regalaron \textit{un pasaje a Daniel}.
      \ex Le regalaron \textit{un pasaje a México}.
    \end{xlist}
  \ex
    \begin{xlist}
      \ex Eliana encontró \textit{el auto de repente}.
      \ex Eliana encontró \textit{el auto de su primo}.
    \end{xlist}
  \ex
    \begin{xlist}
      \ex El hombre tocó \textit{al perro con la mano}.
      \ex El hombre tocó \textit{al perro con sarna}.
    \end{xlist}
  \ex
    \begin{xlist}
      \ex Juan abrió \textit{la puerta de una patada}.
      \ex Juan abrió \textit{la puerta de la habitación}.
    \end{xlist}
\end{exe}

\begin{resuelto}
1. a) i. Le regalaron [un pasaje] [a Daniel] / ii. Le regalaron [un pasaje a
México]. b) En i., \forma{un pasaje} y \forma{a Daniel} son dos constituyentes
que admiten las tres pruebas propuestas: se mueven separadamente (\forma{A Daniel
le regalaron un pasaje. Un pasaje le regalaron a Daniel}); responden a distintas
preguntas (\forma{¿Qué le regalaron a Daniel? ¿A quién le regalaron un pasaje?});
se pronominalizan independientemente (\forma{Le regalaron un pasaje. Se lo
regalaron a Daniel}).\footnote{V. la nota~\ref{fn:05-se}.} En
ii., \forma{un pasaje a México} es un solo constituyente porque se mueve
conjuntamente (\forma{Un pasaje a México le regalaron}); responde a una única
pregunta (\forma{¿Qué le regalaron?}); admite pronominalización conjunta
(\forma{Se lo regalaron}).
\end{resuelto}
\end{ejercicio}

\section{Los sintagmas}\label{sec:05-sintagmas}

Como vimos cuando aplicamos la prueba de sustitución, un constituyente puede
reducirse a un único elemento central: el núcleo. Por ejemplo, en
\exref{ex:05-pedro}, ambos constituyentes de la oración están formados por un
solo elemento:

\begin{exe}
  \ex\label{ex:05-pedro} {}[Pedro]\etq{A} [corre]\etq{B}
\end{exe}

Sin embargo, por lo general, los núcleos se expanden vinculándose con otras
palabras y conformando unidades mayores denominadas \termino{sintagmas}. Desde
esta perspectiva, \textcite[138]{bosque2009} definen los sintagmas como
«expansiones de un elemento central o nuclear, en torno al que se articulan
modificadores y complementos».

\begin{destacado}
Según vemos, entonces, todos los sintagmas tienen núcleo, es decir, son
endocéntricos y, en todos ellos, es el elemento nuclear el que da nombre al
sintagma.
\end{destacado}

\subsection{Las clases de palabras y su proyección sintagmática}\label{sec:05-clases}

Las palabras de la lengua se agrupan en dos clases mayores: las léxicas, por un
lado, y las gramaticales o funcionales, por el otro. Palabras como \forma{libro,
alto, leer, fácilmente} o \forma{para} se consideran léxicas porque su contenido
puede asociarse con ideas o conceptos \parencite[52]{rodriguezramalle2005}. Las
\termino{palabras léxicas} o de contenido descriptivo inherente incluyen el
sustantivo (N), el adjetivo (A), el verbo (V), el adverbio (Adv) y la
preposición (P).\footnote{Respecto de la preposición como palabra léxica, v. más
adelante §\ref{sec:05-sp}.} Pero, además de estas clases, también integran la
lengua las que Bosque (1991)\pendiente{En la bibliografía figura Bosque (1989)}
denominó «palabras engranaje», como \forma{y, mi, el} o \forma{que}, que no
tienen significado en sí mismas, sino que se definen fundamentalmente por la
función que cumplen en la oración. Son las llamadas \termino{palabras
funcionales}, constituidas por los determinativos (artículos, demostrativos,
posesivos, indefinidos, etc.) y las conjunciones (coordinantes y
subordinantes).

Aunque más adelante, en el \capref{6}, se verá que también las palabras y
categorías funcionales proyectan sintagmas en los que funcionan como núcleos,
aquí nos vamos a concentrar en los sintagmas formados por las categorías
léxicas. Todas las palabras léxicas se expanden o proyectan en sintagmas de los
que son núcleo, de modo que tenemos sintagmas nominales (SSNN), sintagmas
adjetivales (SSAA), sintagmas adverbiales (SSADV), sintagmas preposicionales
(SSPP) y sintagmas verbales (SSVV). En este capítulo nos ocupamos solo de las
cuatro primeras clases de sintagmas, ya que, dada la complejidad del SV, sus
particularidades se tratarán con más detalle en los capítulos
siguientes.\footnote{En el \capref{6} se trata la oración y su estructura; en el
\capref{7}, la estructura argumental y sintáctica, y en el \capref{8}, la
estructura eventiva.}

\begin{ejercicio}\label{ej:05-nucleo}
En los siguientes sintagmas, según se muestra en el ejemplo resuelto~1:

\begin{consignas}
  \item identifique el núcleo (tenga en cuenta cuál es la palabra central a la
        que las demás modifican o complementan);
  \item indique la categoría del núcleo de cada sintagma y señale qué clase de
        sintagma proyecta.
\end{consignas}

\begin{exe}
  \ex Ayer por la tarde.
  \ex Bueno con la gente.
  \ex Cerca de la esquina.
  \ex Directamente proporcional.
  \ex Para mayores de edad.
  \ex Profunda inquietud.
  \ex Armario antiguo.
  \ex Sin sentimientos negativos.
\end{exe}

\begin{resuelto}
1. a) El núcleo es \forma{ayer}. b) \forma{Ayer} es un adverbio y proyecta un
SAdv.
\end{resuelto}
\end{ejercicio}

\subsection{La estructura del sintagma}\label{sec:05-xbarra}

De acuerdo con el modelo de principios y parámetros propuesto por
\textcite{chomsky1981} (v. \capref{1}, §\ref{sec:01-estructura}), la teoría
X\prima{} o teoría de la X con barra, que se encarga de la buena formación de
los sintagmas, plantea que todos los sintagmas de todas las lenguas poseen la
misma estructura. No solo son siempre endocéntricos, es decir, poseen un núcleo
que da la categoría a todo el sintagma, sino que también adoptan la misma
configuración estructural. Veamos un ejemplo:

\begin{exe}
  \ex\label{ex:05-presidente} {}[El presidente de la república] inauguró las
      sesiones legislativas.
\end{exe}

Si consideramos el SN \forma{el presidente de la república} y tenemos en cuenta
cómo se conforman los constituyentes, podemos decir que, dado que \forma{ser
presidente} exige \forma{presidir algo}, el núcleo \forma{presidente} primero se
une a \forma{de la república}, que, al ser un constituyente exigido por el
núcleo, funciona como su complemento (Comp). El tercer elemento que integra el
SN es el artículo \forma{el}, también requerido para la buena formación de la
oración, ya que \juicio{*}{\forma{Presidente de la república inauguró las
sesiones legislativas}}\footnote{Dicha oración sería aceptable, pero en un
registro muy particular, como son los titulares de diarios.} no sería gramatical
en español. Por tanto, ambos, núcleo y complemento, deben ir precedidos del
artículo \forma{el}, que los determina y cumple la función de especificador
(Esp) dentro del sintagma. Podemos representar la estructura sintagmática que
acabamos de describir mediante encorchetamiento \exref{ex:05-presidente-a} o a
través de un diagrama arbóreo \exref{ex:05-presidente-b}:

\begin{exe}
  \ex\label{ex:05-presidente-de}
    \begin{xlist}
      \ex\label{ex:05-presidente-a} {}[\etq{SN} el\etq{D/Esp} [\etq{N\prima}
          presidente\etq{N/n} [\etq{SP/Comp} de la república]]]
      \ex\label{ex:05-presidente-b}
          \begin{forest} gramatica
            [SN
              [D\\\textit{el}]
              [N\prima
                [N\\\textit{presidente}]
                [SP [\textit{de la república}, roof]]]]
          \end{forest}
    \end{xlist}
\end{exe}

La estructura propuesta no es plana, es decir, no todos los elementos están en
el mismo nivel. Según se puede ver en \exref{ex:05-presidente-b}, la estructura
del sintagma es jerárquica e implica niveles que describen la relación entre los
constituyentes. En primer lugar, en el grado X se ubica el núcleo, que es el
elemento que no puede faltar, define la categoría del sintagma y selecciona los
complementos y modificadores. En el segundo nivel, X\prima, el núcleo se une a
su complemento, que es el elemento con el que tiene una relación más estrecha y,
por último, en el tercer nivel, el especificador es un constituyente que, según
su clase, determina ---\forma{el / este / mi presidente}--- o cuantifica
---\forma{algún / varios / todos los presidentes}--- conjuntamente al núcleo y a
su complemento. El sintagma es, pues, la máxima proyección del núcleo y
constituye el nivel X\prima\prima{} o SX (sintagma de X):

\begin{exe}
  \ex\label{ex:05-esquema}
      \begin{forest} gramatica
        [SX (X\prima\prima)
          [Esp]
          [X\prima
            [X\\(núcleo)]
            [Comp]]]
      \end{forest}
\end{exe}

\subsubsection{Complementos, especificadores y adjuntos}\label{sec:05-comp-esp-adj}

Dado que todas las categorías presentan la misma estructura, este esquema puede
aplicarse a cualquiera de ellas, si se sustituye la variable X por una etiqueta
categorial, como podrían ser N, V, A, P o Adv.\footnote{En este capítulo nos
centramos en los núcleos formados por palabras léxicas ---N, A, Adv y P---,
excepto el V, que se verá en el \capref{6}. A partir del \capref{7} nos iremos
adentrando en los sintagmas proyectados por palabras y categorías
funcionales.}

Es necesario remarcar que hay categorías transitivas, como algunos verbos y como
todas las preposiciones, que siempre exigen un complemento, y categorías
intransitivas, como los adverbios y algunos verbos. En las categorías
intransitivas, la posición de complemento queda vacía, como se muestra en
\exref{ex:05-sonrio}, en el que aparece el verbo intransitivo
\forma{sonreír}.\footnote{Las distintas clases de verbos se presentan en el
\capref{7}.} Asimismo, la posición de especificador también puede quedar vacía,
como se muestra en \exref{ex:05-irene}, donde el núcleo del SN sujeto es un
nombre propio, que no requiere especificador.\footnote{V. más adelante
§\ref{sec:05-sn}.}

\begin{exe}
  \ex\label{ex:05-sonrio} El muchacho [\etq{SV} sonrió\etq{V/n}].
  \ex\label{ex:05-irene} {}[\etq{SN} Irene\etq{N/n}] vendrá a visitarnos.
\end{exe}

Dentro de la estructura sintagmática, los \termino{complementos} son
constituyentes seleccionados o exigidos por el núcleo en función de sus
requerimientos léxicos y sintácticos. Así, por ejemplo, un verbo transitivo,
como \forma{comer}, \forma{escribir} o \forma{leer}, exige un SN vinculado con
su significado que lo complemente: \forma{comió torta, escribió una carta, leyó
un libro}. Desde el punto de vista estructural, también las preposiciones exigen
un complemento, ya que no pueden funcionar solas: \forma{de madera, por tu
culpa, sin aliento}.

Aunque los complementos son exigidos por el núcleo, esto no quiere decir que
siempre tengan que estar explicitados en la sintaxis, ya que pueden quedar
implícitos; por ejemplo, en \forma{Hoy comimos temprano} se sobreentiende que se
ingirió \sig{algo comestible}, que no se desea especificar.

Los \termino{especificadores} modifican conjuntamente al núcleo junto con sus
complementos y modificadores. La posición del especificador suele estar ocupada
por palabras funcionales o gramaticales, como los determinantes [D] (\forma{el /
este / su libro}) o los cuantificadores [Q] (\forma{varias / ambas / algunas
revistas}, \forma{bastante / muy desgastado}).\footnote{Esta información se
amplía más adelante, en §\ref{sec:05-sn}.}

Además, los sintagmas admiten constituyentes más externos, no exigidos por el
núcleo, aunque compatibles con él, como los \termino{adjuntos}, que pueden
modificar a cualquier tipo de unidad \parencite[49]{rodriguezramalle2005}. Con
la incorporación del adjunto, la estructura sintagmática queda como se muestra
en \exref{ex:05-adjunto}:

\begin{exe}
  \ex\label{ex:05-adjunto}
      \begin{forest} gramatica
        [SX
          [Esp]
          [X\prima
            [X\prima
              [X]
              [Comp]]
            [Adjunto]]]
      \end{forest}
\end{exe}

Los adjuntos son opcionales, es decir, su presencia no es obligatoria dentro del
sintagma, como es el caso de los circunstanciales dentro del SV (\forma{conversaron
en la oficina / llovió torrencialmente / se encendió con alcohol}) o de los
modificadores dentro del SN (\forma{camisa roja / mantel redondo / aulas pequeñas
/ luces de colores / fragancias de París}).

En \exref{ex:05-dibujo} el verbo \forma{hacer} está acompañado de dos
constituyentes, [\forma{un dibujo}] y [\forma{en su cuaderno}], pero solo
\forma{un dibujo}, que explicita \sig{lo hecho en el cuaderno}, es exigido por el
núcleo verbal, ya que su ausencia (\juicio{*}{\forma{Hizo en su cuaderno}}) torna
la secuencia agramatical. En cambio, si se omite \forma{en su cuaderno}, la
oración «pierde» información, pero se mantiene bien formada, lo que indica que
ese constituyente es un adjunto.

\begin{exe}
  \ex\label{ex:05-dibujo}
    \begin{xlist}
      \ex Hizo [un dibujo] [en su cuaderno].
      \ex Hizo [un dibujo].
      \ex[*]{Hizo [en su cuaderno].}
    \end{xlist}
\end{exe}

\begin{ejercicio}\label{ej:05-de}
Realice la descripción estructural (DE) de los siguientes sintagmas mediante
encorchetamiento, según se muestra en el ejemplo resuelto~1.

\begin{exe}
  \ex Muy afecto al baile.
  \ex Varios colaboradores expertos.
  \ex Arriba de la heladera.
  \ex Sin comentarios.
\end{exe}

\begin{resuelto}
1. [\etq{SA} muy\etq{Q/Esp} afecto\etq{A/n} [\etq{SP/Comp} a\etq{P/n}
[\etq{SN/Comp} el\etq{D/Esp} baile\etq{N/n}]]]
\end{resuelto}
\end{ejercicio}

\begin{ejercicio}\label{ej:05-comp-adj}
Considere los siguientes sintagmas y luego, según se muestra en el ejemplo
resuelto~1:

\begin{consignas}
  \item identifique el núcleo e indique el tipo de sintagma que proyecta;
  \item distinga si los constituyentes que acompañan a los núcleos son
        complementos o adjuntos. Justifique su respuesta de acuerdo con la clase
        de palabra a la que pertenece cada núcleo y sus requerimientos para que
        la secuencia quede bien formada.
\end{consignas}

\begin{exe}
  \ex Orgullosa de sus logros.
  \ex La feria del libro.
  \ex Similar a los demás casos.
  \ex Por la izquierda.
  \ex Bajo la mesada.
  \ex Antes de las diez.
\end{exe}

\begin{resuelto}
1. a) El núcleo del sintagma es \forma{orgullosa} y proyecta un SA. b) \forma{De
sus logros} es un SP exigido por \forma{orgullosa} y funciona como Comp.
\end{resuelto}
\end{ejercicio}

\begin{ejercicio}\label{ej:05-titulos}
Considere los títulos de las obras que se presentan a continuación y resuelva
las consignas siguientes, según se muestra en los ejemplos resueltos:

\begin{consignas}
  \item identifique el núcleo y clasifique cada sintagma;
  \item determine la función que cumplen los constituyentes que acompañan al
        núcleo (especificador, complemento o adjunto). Justifique su respuesta
        considerando lo visto anteriormente.
\end{consignas}

\begin{exe}
  \ex \textit{Caos} (Magalí Tajes)
  \ex \textit{Un mundo feliz} (Aldous Huxley)
  \ex \textit{El jardín de los cerezos} (Antón Chéjov)
  \ex \textit{Cien años de soledad} (Gabriel García Márquez)
  \ex \textit{Ágilmente} (Estanislao Bachrach)
  \ex \textit{Ese verano a oscuras} (Mariana Enriquez)
  \ex \textit{Terrenal} (Mauricio Kartun)
  \ex \textit{Desde la ventana} (Carmen Martín Gaite)
\end{exe}

\begin{resuelto}[Ejemplos resueltos]
1. a) Es un SN porque el núcleo es el sustantivo \forma{caos}. b) El núcleo no
presenta complementos ni modificadores. 2. a) Es un SN porque el núcleo es el
sustantivo \forma{mundo}. b) \forma{Feliz} es un modificador adjunto porque no
está seleccionado por el núcleo.
\end{resuelto}
\end{ejercicio}

\begin{ejercicio}\label{ej:05-vf}
Integre los conocimientos incorporados hasta ahora y decida si las siguientes
afirmaciones son verdaderas o falsas, según se muestra en el ejemplo
resuelto~1. Fundamente sus respuestas.

\begin{exe}
  \ex Las palabras léxicas ---N, A, Adv, V y P--- se expanden o proyectan en
      SSXX.
  \ex El núcleo de un sintagma es siempre una palabra gramatical.
  \ex Todos los constituyentes de un sintagma se vinculan de la misma forma con
      el núcleo.
  \ex Los complementos son requeridos por el núcleo y se diferencian de los
      adjuntos, que son opcionales.
\end{exe}

\begin{resuelto}
1. V(erdadero). Las palabras léxicas son los núcleos de los sintagmas que
proyectan.
\end{resuelto}
\end{ejercicio}

\section{Las clases de sintagmas}\label{sec:05-clasesintagmas}

En los apartados que siguen, analizamos con mayor detenimiento las
características de los distintos tipos de sintagmas.

\subsection{El sintagma nominal (SN) y su estructura interna}\label{sec:05-sn}

Como ya hemos visto, el SN es aquel que tiene por núcleo a un sustantivo. En
cuanto a su estructura interna, como sostiene \textcite[153]{ditullio2014},
«resulta muy compleja debido a las múltiples posibilidades de realización que
presentan sus dependencias y a la estrecha conexión entre estas y el tipo de
sustantivo que funciona como núcleo».

De acuerdo con la teoría X\prima, que presentamos en §\ref{sec:05-xbarra}, el SN
posee la siguiente estructura:

\begin{exe}
  \ex\label{ex:05-esquema-sn}
      \begin{forest} gramatica
        [SN
          [Esp\\D / Q]
          [N\prima
            [N\prima
              [N]
              [Comp\\SP]]
            [Mod\\SA / SP / SN]]]
      \end{forest}
\end{exe}

\paragraph{El núcleo.} Cualquier clase semántica de sustantivo puede ser núcleo
de un SN, tanto común \exref{ex:05-armario} como propio \exref{ex:05-martin},
colectivo \exref{ex:05-bandada}, individual concreto \exref{ex:05-armario} o
abstracto \exref{ex:05-ilusion}.\pendiente{El original remite a (23) para los
abstractos}

\begin{exe}
  \ex\label{ex:05-armario}
    \begin{xlist}
      \ex El \textit{armario} del fondo.
      \ex Un \textit{edificio} muy alto del centro de la ciudad.
    \end{xlist}
  \ex\label{ex:05-martin}
    \begin{xlist}
      \ex \textit{Martín}, mi primo de Catamarca.
      \ex \textit{Luisa}, la profesora de música.
    \end{xlist}
  \ex\label{ex:05-bandada}
    \begin{xlist}
      \ex Una \textit{bandada} de pájaros.
      \ex\label{ex:05-arboleda} Esa lejana \textit{arboleda} de mi infancia.
    \end{xlist}
  \ex\label{ex:05-ilusion}
    \begin{xlist}
      \ex Una \textit{ilusión} pasajera como el verano.
      \ex El \textit{consentimiento} mutuo de la pareja.
    \end{xlist}
\end{exe}

En el SN, el núcleo determina los rasgos de género y número de los modificadores
y especificadores. Así, en \exref{ex:05-arboleda}, tanto \forma{lejana} como
\forma{esa} están en femenino singular para concordar con el núcleo
\forma{arboleda}.

\paragraph{Los complementos.} Si bien la mayor parte de los nombres, como
\forma{mesa} o \forma{cuaderno}, no tienen como requisito semántico la exigencia
de un complemento que complete su significación, hay algunos sustantivos
específicos que requieren un SP como complemento, como los siguientes:

\begin{itemize}
  \item Los \termino{nombres eventivos}, que son aquellos que por su sentido
        implican la ocurrencia de un suceso, que se desarrolla en un tiempo y
        lugar. Muchos nombres eventivos son deverbales, es decir, provienen de
        verbos (\forma{construcción, agasajo, reunión}), pero otros no
        (\forma{fiesta, batalla, partido}). Se reconocen porque pueden ser
        sujeto de \forma{tener lugar} \exref{ex:05-lugar} y complemento de la
        preposición \forma{durante} \exref{ex:05-durante}:
\end{itemize}

\begin{exe}
  \ex\label{ex:05-lugar} {}[La \textit{formulación} de una nueva teoría] tuvo
      lugar en la Universidad de Buenos Aires.
  \ex\label{ex:05-durante} El suceso ocurrió durante [la \textit{guerra} de los
      seis días].
\end{exe}

\begin{itemize}
  \item Los \termino{nombres relacionales}, que son los que establecen un vínculo
        semántico con lo designado por el complemento, como ser \forma{padre de
        alguien, alumno de algún profesor} o \forma{especialista en alguna
        disciplina}. Son ejemplos de SSNN con nombres relacionales los
        siguientes:
\end{itemize}

\begin{exe}
  \ex {}[El \textit{hijo} de Susana] vino a visitarnos.
  \ex {}[El \textit{autor} de esa novela] vive en Francia.
\end{exe}

\begin{itemize}
  \item Los \termino{nombres de representación}, que son los que designan una
        imagen o descripción de lo manifestado por el complemento. Son ejemplos
        de SSNN de representación:
\end{itemize}

\begin{exe}
  \ex Esa es [la \textit{foto} de la abuela] cuando era joven.
  \ex {}[La \textit{imagen} de un héroe de Malvinas] me impactó.
\end{exe}

En todos los casos, los complementos seleccionados también pueden quedar
implícitos:

\begin{exe}
  \ex Puse en un marco plateado [el \textit{retrato} (de la
      abuela)\footnote{Convencionalmente, los elementos opcionales se colocan
      entre paréntesis.}].
\end{exe}

\paragraph{Los modificadores adjuntos.} Son constituyentes opcionales que
aportan características del nombre y permiten restringirlo a un subconjunto
determinado de la clase general de entidades que abarca. Los modificadores
típicos del nombre son los adjetivos, tanto los calificativos, que indican
propiedades, \exref{ex:05-hojas} y \exref{ex:05-fornidos}, como los
relacionales, que vinculan al nombre con un ámbito determinado. Dado que los
relacionales son derivados de sustantivos, la relación también puede expresarse
mediante un SP, \exref{ex:05-tambera} y \exref{ex:05-lunar}
(v.~§\ref{sec:05-sa}):

\begin{exe}
  \ex\label{ex:05-hojas} Las hojas \textit{rayadas} de mi cuaderno \textit{azul}.
  \ex\label{ex:05-fornidos} Los \textit{fornidos} trabajadores de las minas.
  \ex\label{ex:05-tambera} Una rica zona \textit{tambera / de tambos}.
  \ex\label{ex:05-lunar} El inmenso cráter de la superficie \textit{lunar / de la
      luna}.
\end{exe}

Otro tipo de modificador del nombre es la \termino{aposición}, formada por un
sustantivo que pertenece a la misma clase semántica del núcleo cuando está en
relación de adyacencia directa con él, como vemos en \exref{ex:05-gregorio} y
\exref{ex:05-gomez}.

\begin{exe}
  \ex\label{ex:05-gregorio} {}[Gregorio, \textit{mi gato},] está todo el día
      tirado en la cocina.
  \ex\label{ex:05-gomez} {}[El comisario \textit{Gómez}] ya está retirado del
      servicio activo.
\end{exe}

\begin{recuadro}[Para ampliar: modificadores restrictivos y explicativos]
Los modificadores adjetivos y las aposiciones pueden ser restrictivos o
explicativos. Los restrictivos, también llamados especificativos, acotan la
referencia del nombre \exref{ex:05-andino}; los explicativos, en cambio, agregan
información complementaria, por lo que van marcados prosódicamente por una
figura tonal propia, lo que se refleja en la escritura porque se colocan entre
comas y se pueden suprimir \exref{ex:05-emilce}.

\begin{exe}
  \ex\label{ex:05-andino}
    \begin{xlist}
      \ex {}[El paisaje \textit{andino}] me pareció espectacular.
      \ex {}[Los hermanos \textit{Grimm}] le agregaron un final \textit{feliz} a
          la historia. (Adaptado de A.~Litvin, \textit{El libro de los
          villanos}, V\&R, 2015)
    \end{xlist}
  \ex\label{ex:05-emilce}
    \begin{xlist}
      \ex {}[Emilce, \textit{agitada},] señaló con su dedito de tres años la
          puerta abierta de su cuarto. (M.~Di Marco, \textit{El fantasma del
          Reich}, Sudamericana, 1995)
      \ex {}[Jacobo, \textit{el niño tonto},] solía subirse a la azotea y espiar
          la vida de los vecinos. (E.~Anderson Imbert, \textit{El gato de
          Cheshire}, Losada, 1965)
    \end{xlist}
\end{exe}

También los SSPP pueden modificar al nombre y son, en general, modificadores
restrictivos:

\begin{exe}
  \ex
    \begin{xlist}
      \ex Ahora trabaja en [una escuela \textit{con ruralidad}].
      \ex {}[La entrevista \textit{de mañana}] va a ser fundamental para el
          proyecto.
    \end{xlist}
\end{exe}
\end{recuadro}

\paragraph{Los especificadores.} Los sustantivos comunes suelen ir precedidos
por un determinativo en la posición de especificador. Los determinativos son
palabras funcionales que se definen por su relación con otra u otras palabras de
las que delimitan su alcance o a las que cuantifican al indicar «cuáles o
cuántos de los elementos incluidos en la clase denotada [\ldots] se deben
considerar» \parencite[93]{rodriguezramalle2005}. La clase de los determinativos
está formada por los determinantes [D] y los cuantificadores [Q], que son
palabras que carecen de contenido descriptivo específico y no establecen una
relación semántica con el sustantivo al que preceden, por lo que pueden
anteponerse a cualquier tipo de sustantivo, del que anticipan el género y el
número:

\begin{exe}
  \ex {}[\textit{El / nuestro / este / algún} pueblo] no debe estar lejos de
      aquí.
  \ex {}[\textit{Mis / estas / varias} amigas] vendrán a visitarme para mi
      cumpleaños.
\end{exe}

La subclase de los determinantes está formada por los artículos (\forma{el, los,
la, las, lo}), los demostrativos (\forma{este, esos, aquella}) y los posesivos
(\forma{mi, nuestros, tus}). Los cuantificadores, por su parte, están conformados
por los numerales cardinales (\forma{cinco, ocho, cien}), los numerales ordinales
(\forma{primero, segundo, decimoquinto}) y los indefinidos (\forma{mucho/s,
poco/s, varios/as, ambos/as, sendos/as, cualquiera, todo/s, bastante,
demasiado}).

A diferencia de los sustantivos comunes, que requieren un determinativo para
funcionar oracionalmente, los propios, por lo general, se emplean solos como
núcleo del SN. Los nombres propios no necesitan especificador porque ya están
determinados; así, en español no es necesario especificar \forma{la Irene, el
Pedro}. Sin embargo, en varios dialectos los hablantes colocan un especificador
delante de los nombres propios para indicar familiaridad o que se trata de
alguien conocido:

\begin{exe}
  \ex
    \begin{xlist}
      \ex La Frida\ldots{} Qué linda estás. (R.~Cossa, \textit{Gris de
          ausencia}, 1981)
      \ex ¿Y vos sabés cómo sufren los yanquis por eso? ¡Uf! Todos los días se
          reúnen dos horas a llorar porque la Yoly de Lanús los tiene acá.
          (N.~Fernández Tiscornia, \textit{Made in Lanús}, Cántaro, 1986)
    \end{xlist}
\end{exe}

\begin{ejercicio}\label{ej:05-buenaformacion}
Considere la buena o mala formación de los SSNN siguientes. Tenga en cuenta la
clase semántica de los sustantivos que funcionan como núcleo y sus posibilidades
y exigencias combinatorias, según se muestra en el ejemplo resuelto~1.

\begin{exe}
  \ex Trajo un regalo para novio.
  \ex El Mario.
  \ex Los ingredientes de la receta.
  \ex Compraron arena.
  \ex Vive con alegría.
\end{exe}

\begin{resuelto}
1. Es agramatical porque \forma{novio} es un sustantivo común y, por lo tanto,
debe ir precedido por un determinativo.
\end{resuelto}
\end{ejercicio}

\begin{ejercicio}\label{ej:05-ratones}
Lea detenidamente el siguiente texto, considere los SSNN subrayados y señale sus
constituyentes mediante corchetes, según se muestra en el ejemplo resuelto.

% Texto de divulgación protegido: se incorpora a mano desde el original.
\textoaincorporar{«A los ratones no les gusta el queso», de A.~Litvin
(\textit{El libro de las mentiras}, V\&R, 2015), con los SSNN subrayados:
\textit{los ratones; el queso; una compañía de control de plagas; los alimentos
con alto contenido de azúcares}}

\begin{resuelto}
{}[\etq{SN} los\etq{D/Esp} ratones\etq{N/n}]
\end{resuelto}
\end{ejercicio}

\begin{ejercicio}\label{ej:05-titulares}
Lea con atención los siguientes titulares de diarios con sus respectivas
volantas\footnote{La volanta o copete es el texto que se ubica sobre el título
de un artículo o nota periodística.} y algunos fragmentos del cuerpo de la
noticia. Luego analice los SSNN subrayados, según se muestra en el ejemplo
resuelto.

\begin{consignas}
  \item Realice el encorchetamiento indicando el núcleo y las funciones de los
        restantes constituyentes.
  \item Indique si los modificadores son restrictivos o no.
\end{consignas}

\textbf{Noticia 1}

\begin{quote}\small
\textit{La diva vuelve a Sudamérica}\\
\textbf{(1)~\subrayado{Madonna} ofrecerá un concierto gratuito en
(2)~\subrayado{la playa de Copacabana}}\\
La reina del pop se presentaría en el mes de mayo. Este lunes, los organizadores
del evento anunciarán todos los detalles.

(3)~\subrayado{El columnista Lauro Jardim del diario O Globo} dijo este mes que
el concierto estaba previsto para el 4 de mayo, fecha aún no confirmada por las
autoridades. (4)~\subrayado{Madonna, de 65 años}, inició su gira
\textit{Celebration} el pasado mes de octubre.
\end{quote}

\textbf{Noticia 2}

\begin{quote}\small
\textit{Es (1)~\subrayado{la primera teatrista latinoamericana} y la segunda
mujer en obtenerlo}\\
\textbf{(2)~\subrayado{Lola Arias, ganadora del Premio Internacional Ibsen
2024}}\\
El jurado destacó la «dimensión ética» y (3)~\subrayado{el carácter
«transnacional» de la obra de la artista multidisciplinaria}.
(\url{www.pagina12.com.ar})
\end{quote}

\begin{resuelto}
Noticia 1. (1) a) [\etq{SN} Madonna\etq{N/n}]. b) El núcleo no presenta
modificadores ni complementos. (2) a) [\etq{SN} la\etq{D/Esp} playa\etq{N/n}
[\etq{SP/M} de\etq{P/n} [\etq{SN/Comp} Copacabana]]]. b) \forma{De Copacabana}
es un modificador restrictivo.
\end{resuelto}
\end{ejercicio}

\begin{ejercicio}\label{ej:05-lobos}
Lea atentamente el siguiente texto y resuelva las consignas que se presentan a
continuación, según se muestra en el ejemplo resuelto:

\begin{consignas}
  \item analice la estructura interna de los SSNN subrayados;
  \item determine en qué casos sería posible o no, dentro de este texto,
        alterar el orden de los constituyentes del sintagma;
  \item fundamente su respuesta teniendo en cuenta los efectos de esos posibles
        cambios (si se produce una modificación del significado o si el
        resultado es agramatical).
\end{consignas}

% Texto de divulgación protegido: se incorpora a mano desde el original.
\textoaincorporar{«Lobos. Dato real», adaptado de A.~Litvin (\textit{El libro de
los villanos}, V\&R, 2015), con los SSNN subrayados: \textit{el lobo, un
mamífero carnívoro; varias subespecies; el famoso aullido; los miembros de su
grupo; muy raras ocasiones}}

\begin{resuelto}
a) [\etq{SN} el\etq{D/Esp} lobo\etq{N/n} [\etq{SN/ApE} un\etq{D/Esp}
mamífero\etq{N/n} carnívoro\etq{SA/Mod}]]. b) Dado que el SN \forma{un mamífero
carnívoro} es una aposición explicativa, puede permutarse por el núcleo al que
modifica ---\forma{un mamífero carnívoro, el lobo}--- sin alterar ni la
gramaticalidad ni el sentido de la estructura que integra.
\end{resuelto}
\end{ejercicio}

\subsection{El sintagma adjetival (SA) y su estructura interna}\label{sec:05-sa}

La estructura del SA puede representarse de la siguiente manera:

\begin{exe}
  \ex\label{ex:05-esquema-sa}
      \begin{forest} gramatica
        [SA
          [Esp\\Q / SAdv]
          [A\prima
            [A\prima
              [A]
              [Comp\\SP]]
            [Mod\\SP]]]
      \end{forest}
\end{exe}

\paragraph{El núcleo.} El núcleo del SA es un adjetivo que puede ser de
distintas subclases: calificativo \exref{ex:05-calif}, relacional
\exref{ex:05-relac} o adverbial \exref{ex:05-adverb}:

\begin{exe}
  \ex\label{ex:05-calif}
    \begin{xlist}
      \ex Fue un accidente \textit{estremecedor}.
      \ex Esta semana se estrenó una \textit{interesante} película de aventuras.
    \end{xlist}
  \ex\label{ex:05-relac}
    \begin{xlist}
      \ex\label{ex:05-relac-a} La reforma \textit{educativa} impulsa un cambio
          de planes.
      \ex\label{ex:05-relac-b} Un escritor \textit{italiano} ganó el premio de
          novela este año.
    \end{xlist}
  \ex\label{ex:05-adverb}
    \begin{xlist}
      \ex Sus \textit{frecuentes} quejas me desesperan.
      \ex El \textit{actual} rey de Inglaterra es Carlos III.
    \end{xlist}
\end{exe}

Los adjetivos \termino{calificativos} conforman el grupo focal o central de la
clase. Son los que expresan propiedades del nombre (\forma{alto, linda, rojo,
amable}). Estos adjetivos admiten graduación (\forma{muy alto, bastante linda,
poco amable}) y suelen tener valores opositivos (\forma{lindo/feo; alto/bajo}).
En cuanto a su posición, pueden ubicarse antes del nombre o después de él
(\forma{un lindo recuerdo / un recuerdo lindo; esa amable mujer / esa mujer
amable}). A diferencia de los calificativos, los adjetivos \termino{relacionales}
no expresan propiedades, sino que restringen el alcance del sustantivo a un
ámbito determinado (\forma{consejo académico, reforma constitucional, educación
cívica}). Como no denotan propiedades, no admiten graduación
(\juicio{*}{\forma{reforma muy constitucional}}), ni manifiestan valores polares,
por lo que no suelen poseer un opuesto. Los relacionales solo pueden aparecer
pospuestos al sustantivo al que modifican (\juicio{*}{\forma{constitucional
reforma}}, \juicio{*}{\forma{cívica educación}}). Como se vio en
\exref{ex:05-relac}, algunos relacionales designan pertenencia a un ámbito
(\forma{político, económico, educativo, sexual}, etc.; v.~\exref{ex:05-relac-a})
o a un lugar determinado (\forma{francés, argentino, bonaerense};
v.~\exref{ex:05-relac-b}); estos últimos son los que se conocen como
\termino{gentilicios}.\pendiente{El original remite a (47) para los
relacionales} Por su parte, los adjetivos \termino{adverbiales} se vinculan con
determinados adverbios referidos a la localización o a la cuantificación
(\forma{antiguo, próximo, habitual, propio, único, simple, último}) y suelen
anteponerse al nombre (\forma{sus frecuentes reclamos / reclama
frecuentemente}). Al igual que los relacionales, no admiten graduación
(\juicio{*}{\forma{su muy / bastante último libro}}).

Los SSAA pueden estar formados por el núcleo solamente (\forma{estaba
\textbf{feliz}, seguía \textbf{triste}}), o bien combinarse con modificadores
(\forma{\textbf{triste} por la noticia}) o complementos (\forma{\textbf{carente}
de recursos, \textbf{apto} para el servicio}) y también pueden recibir
especificadores (\forma{plenamente \textbf{feliz} / muy \textbf{triste}}).

\paragraph{Los complementos.} Como vimos que sucede con los sustantivos, la
mayoría de los adjetivos (\forma{alto, grande, fugaz}, etc.) no necesita
complemento. No obstante, algunos, en función de su significado, exigen un
complemento que se manifiesta mediante un SP encabezado por una preposición
determinada, que puede ser:

\begin{itemize}
  \item \forma{a} (\forma{aficionado, equivalente, propicio, relativo});
  \item \forma{de} (\forma{deseoso, digno, orgulloso, cansado});
  \item \forma{con} (\forma{conforme, contradictorio, satisfecho});
  \item \forma{en} (\forma{experto, abundante, rico});
  \item \forma{para} (\forma{necesario, adecuado, disponible}).
\end{itemize}

Veamos algunos ejemplos:

\begin{exe}
  \ex Mi tío Juan es [muy \textit{aficionado a} la pesca deportiva].
  \ex Todos estaban [\textit{deseosos de} recibir al nuevo integrante de la
      familia].
  \ex El cliente quedó [sumamente \textit{conforme con} la atención].
  \ex Un profesional [\textit{experto en} informática] preparó la encuesta.
  \ex Lea detenidamente el instructivo con los requisitos [\textit{necesarios
      para} el proceso].
\end{exe}

\paragraph{Los modificadores.} Dado que, salvo en los casos que vimos más arriba,
la mayoría de los adjetivos no exige complemento, los SSPP que los acompañan se
consideran modificadores:

\begin{exe}
  \ex Lo vi [feliz \textit{por el encuentro}].
  \ex El vaso está [lleno \textit{hasta la mitad}].
  \ex Estaba [loco \textit{de amor}] por esa mujer.
  \ex Pedrito está [inmenso \textit{de alto}].
\end{exe}

Aunque no es exigido, el SP cuyo complemento es otro adjetivo o un sustantivo
usualmente restringe el alcance del núcleo del SA \parencite{ditullio2014}.

\paragraph{Los especificadores.} Los especificadores del núcleo del SA están
realizados por distintos tipos de cuantificadores, que pueden ser de grado,
\exref{ex:05-muy} y \exref{ex:05-bastante}, o algunos de los terminados en
\forma{-mente} con valor cuantificacional, \exref{ex:05-excesivamente} y
\exref{ex:05-sumamente}.

\begin{exe}
  \ex\label{ex:05-muy} Quedó [\textit{muy} herido por el accidente].
  \ex\label{ex:05-bastante} Es [\textit{bastante} inmaduro para resolver sus
      problemas].
  \ex\label{ex:05-excesivamente} Milena es [\textit{excesivamente} prolija para
      todo], pero su hermano es [\textit{terriblemente} desprolijo].
  \ex\label{ex:05-sumamente} Los inversores estaban [\textit{sumamente}
      interesados en el proyecto].
\end{exe}

También modifican al adjetivo otros adverbios en \forma{-mente}, que pueden
indicar punto de vista \exref{ex:05-pv}, modalidad \exref{ex:05-modalidad} o foco
\exref{ex:05-foco}:

\begin{exe}
  \ex\label{ex:05-pv} Algunos países son [\textit{lingüísticamente /
      étnicamente} heterogéneos].
  \ex\label{ex:05-modalidad} Muchos barrios periféricos son
      [\textit{lamentablemente / indudablemente} peligrosos].
  \ex\label{ex:05-foco} Ese es un riesgo [\textit{específicamente /
      principalmente} económico].
\end{exe}

\begin{ejercicio}\label{ej:05-slogans}
Analice la estructura interna del sintagma de estos viejos eslóganes
publicitarios, según se muestra en el ejemplo resuelto~1.

\begin{exe}
  \ex Demasiado rico para tomarlo despacito. (Nesquik)
  \ex Único en el mundo. / Único. (Fernet Branca)
  \ex Buenos de punta a punta. (Cigarrillos Imparciales)
  \ex Juntos. (Telefé)
\end{exe}

\begin{resuelto}
1. [\etq{SA} demasiado\etq{Q/Esp} [\etq{A\prima} rico\etq{A/n} [\etq{SP/M}
para\etq{P/n} [tomarlo despacito]]]]
\end{resuelto}
\end{ejercicio}

\begin{ejercicio}\label{ej:05-posicion}
Según se muestra en los ejemplos resueltos, en las siguientes series de SSNN
cuyos núcleos son modificados por SSAA:

\begin{consignas}
  \item cambie la posición del modificador;
  \item explique si el cambio produce agramaticalidad, altera el significado o
        si solo produce un efecto estilístico;
  \item en los casos en que cambia el significado, especifique en qué consiste
        la modificación.
\end{consignas}

\begin{exe}
  \ex el funcionario alto
  \ex la crisis económica
  \ex una tarde estupenda
  \ex el discurso político
  \ex ese pobre hombre
  \ex un rico chico
  \ex una enfermedad terrible
  \ex la educación sexual
\end{exe}

\begin{consignas}[resume]
  \item Teniendo en cuenta la resolución de los puntos anteriores, ¿puede
        establecer alguna relación entre lo observado y las subclases del
        adjetivo?
\end{consignas}

\begin{resuelto}[Ejemplos resueltos]
1. a) i) \forma{el funcionario alto} / ii) \forma{el alto funcionario}. b) El
cambio de posición no produce agramaticalidad, pero cambia el significado del A.
c) i) \sig{de elevada estatura} / ii) \sig{importante}. 2. a) i) \forma{la crisis
económica} / ii) \juicio{*}{\forma{la económica crisis}}. b) Si se cambia de
posición, la construcción se vuelve agramatical.
\end{resuelto}
\end{ejercicio}

\begin{ejercicio}\label{ej:05-expansionA}
\begin{consignas}
  \item Realice las expansiones de los siguientes núcleos adjetivos para formar
        sintagmas (considere las respuestas posibles que se muestran abajo para
        el ejemplo~1).
\end{consignas}

\begin{exe}
  \ex \textit{inteligente}
  \ex \textit{propenso}
  \ex \textit{habitual}
  \ex \textit{conforme}
  \ex \textit{deseoso}
\end{exe}

\begin{resuelto}[Respuestas posibles]
1. \forma{sumamente inteligente / muy inteligente para su edad}.
\end{resuelto}

\begin{consignas}[resume]
  \item Teniendo en cuenta la resolución del punto anterior, ¿qué particularidad
        presenta la expansión sintagmática de los adjetivos \forma{propenso,
        conforme} y \forma{deseoso}?
\end{consignas}
\end{ejercicio}

\subsection{El sintagma preposicional (SP) y su estructura interna}\label{sec:05-sp}

La estructura sintagmática del SP puede esquematizarse como sigue:

\begin{exe}
  \ex\label{ex:05-esquema-sp}
      \begin{forest} gramatica
        [SP
          [Esp\\Q]
          [P\prima
            [P]
            [Comp\\SN / SAdv / SP / O]]]
      \end{forest}
\end{exe}

\paragraph{El núcleo.} De todas las clases de palabras que hemos visto, sin duda
la preposición resulta «la clase más conflictiva, en el sentido de que se ha
discutido y se sigue discutiendo su carácter de palabra léxica»
\parencite[56]{giammatteo2009}. Aunque aquí adherimos a este último criterio, no
dejamos de reconocer, como señalan las mencionadas autoras, el carácter atípico
de esta clase de palabras. Para la gramática tradicional, que la preposición no
funcione en forma independiente, es decir, que siempre requiera un complemento,
fue una pauta determinante para considerar que no podía funcionar como núcleo de
un sintagma. Sin embargo, también algunos verbos transitivos y otras palabras,
como los sustantivos eventivos o relacionales, que ya hemos visto en
§\ref{sec:05-sn}, exigen un complemento. Así, por ejemplo, verbos como
\forma{tener} o \forma{regalar}, entre muchos otros, no admiten ir solos:
\juicio{*}{\forma{Yo tengo}}, \juicio{*}{\forma{Juan regaló}}.\footnote{La
omisión del complemento es posible como respuesta a una pregunta, que supone un
contexto previo en el que el complemento ya ha sido mencionado: \forma{¿Quién
trajo bebidas para la reunión? Luis trajo}. Igualmente, algunas preposiciones
permiten la omisión cuando el complemento ya ha sido nombrado: \forma{¿Compraste
té con hierbas? No, sin}.} Por tanto, el hecho de que la preposición siempre
deba ir acompañada por un complemento no impide que sea el núcleo del sintagma
que encabeza. Y esto es así porque, al menos en la mayoría de los casos, las
preposiciones tienen contenido léxico inherente ---\forma{por} (causa),
\forma{para} (fin), \forma{hacia} (dirección), etc.---, por lo que se ubican
dentro de las clases léxicas y funcionan como núcleos de SSPP. No obstante,
también debemos señalar que, en algunos usos que veremos más
adelante,\footnote{El funcionamiento gramatical de la preposición como marca de
función se trata en el \capref{7}.} las preposiciones tienen un funcionamiento
solamente gramatical y poseen escaso significado. En estos casos se considera
que son meras marcas de función que no actúan como núcleos ni determinan la
categoría del sintagma que introducen.

Por último, otra característica que destaca el carácter nuclear de la
preposición es que, al igual que los verbos, restringe semánticamente a su
complemento. Así, por ejemplo, como el complemento de un verbo como
\forma{comer} debe referir a \sig{algo comestible} (\forma{comí pan / postre /
verduras}), una preposición como \forma{hacia} exige que su complemento refiera
a lugar o tiempo ---\forma{hacia el río / el oeste / la noche}---.

Establecido el carácter léxico de la preposición, al menos para la mayoría de
ellas, podemos añadir que son palabras invariables y que, por otra parte, la
característica que las acerca a las palabras funcionales es su carácter
eminentemente relacional. Los SSPP sirven para establecer vínculos de
dependencia con distintas clases de palabras a las que modifican:

\begin{itemize}
  \item sustantivos: \forma{mesa de madera / planta sin flores / estante con
        libros};
  \item adjetivos: \forma{capaz de todo / bueno para nada / sobresaliente en
        matemáticas};
  \item adverbios: \forma{lejos de casa / antes de la madrugada / siempre en
        invierno};
  \item verbos: \forma{estudia en Madrid / piensa en su familia / depende de tu
        decisión / llegó con amigos};\footnote{En el \capref{7} se plantea que
        algunos verbos, como \forma{pensar} y \forma{depender}, exigen
        complementos encabezados por determinadas preposiciones, que se conocen
        como complementos de régimen (CR).}
  \item preposiciones: \forma{por sobre el mar / desde entre los matorrales}.
\end{itemize}

\paragraph{El complemento.} Según ya hemos dicho, dado que las preposiciones
poseen contenido léxico propio, el complemento que seleccionan debe estar acorde
con su significado. Además, la preposición también impone restricciones
formales a su complemento: este no solo es obligatorio, sino que ocupa una
posición fija dentro del sintagma, siempre inmediatamente después del núcleo:
\forma{de colores} / \juicio{*}{\forma{colores de}}.\footnote{Para las
particularidades de la acentuación de la mayor parte de las preposiciones puede
consultarse el \capref{3}, §\ref{sec:03-suprasegmentales}.} Asimismo, al igual
que el verbo exige caso\footnote{La categoría de caso se presentó en el
\capref{4}, §\ref{sec:04-categorias}.} acusativo a su COD, la preposición impone
caso oblicuo o terminal a su complemento cuando este está representado por un
pronombre (\forma{de mí} / \juicio{*}{\forma{de yo}}).

\paragraph{Los especificadores.} Además del complemento, opcionalmente los SSPP
admiten cuantificadores en la posición de especificador:

\begin{exe}
  \ex
    \begin{xlist}
      \ex Poné el cuadro [\textit{más} hacia la derecha].
      \ex {}[\textit{Muy} desde el fondo de mi corazón] debo decir gracias.
    \end{xlist}
\end{exe}

\begin{ejercicio}\label{ej:05-sp-modifica}
Identifique el núcleo al que modifica cada SP y señale a qué clase de palabra
pertenece, según se muestra en el ejemplo resuelto~1.

\begin{exe}
  \ex Plantas sin flores.
  \ex Lo eligieron por sobre los demás.
  \ex Cerca de la oficina.
  \ex Salir con viejos amigos.
  \ex Feliz por sus logros académicos.
\end{exe}

\begin{resuelto}
1. El SP encabezado por \forma{sin} modifica a \forma{plantas}, que es un N.
\end{resuelto}
\end{ejercicio}

\begin{ejercicio}\label{ej:05-de-sp}
Observe los SSPP subrayados de las siguientes oraciones y resuelva las consignas
de abajo, según se muestra en el ejemplo resuelto~1.

\begin{consignas}
  \item Haga la descripción estructural mediante encorchetamiento.
  \item Determine, en función de la palabra a la que modifica el SP, el
        funcionamiento y significado que adopta en cada caso la preposición
        \forma{de}.
\end{consignas}

\begin{exe}
  \ex La casa \subrayado{de mi abuela} es antigua.
  \ex El libro \subrayado{de gramática} resultó útil.
  \ex Me encantan los caramelos \subrayado{de menta}.
  \ex No supieron \subrayado{de aquella mujer} por varios días.
  \ex Los niños, ávidos \subrayado{de conocimiento}, concurrieron al instituto.
\end{exe}

\begin{resuelto}
1. a) [\etq{SP} de\etq{P/n} [\etq{SN/Comp} mi\etq{D/Esp} abuela\etq{N/n}]]. b)
\forma{De} modifica al N \forma{casa} e indica \sig{posesión}.\pendiente{El
original dice «cartera»}
\end{resuelto}
\end{ejercicio}

\begin{ejercicio}\label{ej:05-preposiciones}
Determine los significados de las preposiciones de los SSPP de la lista de
abajo. Use paráfrasis para completar la columna de la derecha, según se muestra
en el ejemplo resuelto.

\begin{center}
\begin{tabular}{@{}r l p{45mm}@{}}
  \toprule
  & SSPP & Significado de la preposición \\
  \midrule
  1.  & asombrados por la noticia & \respuesta{\sig{a causa de}} \\
  2.  & cien metros sobre el nivel del mar & \\
  3.  & el sapo entre los matorrales & \\
  4.  & los problemas sin solución & \\
  5.  & remedios caseros para el empacho & \\
  6.  & el camino hasta su casa & \\
  7.  & tu opinión sobre las instituciones & \\
  8.  & el edicto contra la guerra & \\
  9.  & un viaje hacia bellos lugares & \\
  10. & un pago mediante transferencia bancaria & \\
  11. & la verdad tras su reacción & \\
  \bottomrule
\end{tabular}
\end{center}
\end{ejercicio}

\subsection{El sintagma adverbial (SAdv) y su estructura interna}\label{sec:05-sadv}

De acuerdo con la teoría X\prima, podemos representar la estructura del SAdv
mediante un esquema arbóreo como el siguiente:

\begin{exe}
  \ex\label{ex:05-esquema-sadv}
      \begin{forest} gramatica
        [SAdv
          [Esp\\Q / SN]
          [Adv\prima
            [Adv\prima
              [Adv]
              [Comp\\SP]]
            [Mod\\SN / SP / O]]]
      \end{forest}
\end{exe}

\paragraph{El núcleo.} A diferencia de otras clases léxicas, como el
sustantivo, el adjetivo o el verbo, que manifiestan categorías morfológicas, los
adverbios, al igual que las preposiciones, son palabras invariables. Esto
implica que, cuando se proyectan en estructuras más complejas y funcionan como
núcleo del SAdv, no manifiestan concordancias. Por otra parte, mientras la
preposición, según hemos visto en §\ref{sec:05-sp}, siempre exige complemento
(\forma{por la ruta, sobre el tejado}), el adverbio, por lo general, no lleva
modificadores obligatorios, aunque, opcionalmente, algunos adverbios pueden
admitirlos (\forma{antes de la época actual, ayer por la tarde}).

Desde el punto de vista sintáctico y semántico, los adverbios se presentan como
una extensa y sumamente heterogénea clase, difícil de sistematizar por la
imposibilidad de determinar características comunes a todos sus miembros. Por
eso aquí presentamos clasificaciones que tienen en cuenta distintos criterios,
ya sean semánticos, categoriales o sintácticos.

Semánticamente se distinguen adverbios de:

\begin{itemize}
  \item lugar: \forma{alrededor, cerca, debajo};
  \item tiempo: \forma{siempre, ya, enseguida};
  \item modo: \forma{bien, peor, lentamente};
  \item afirmación: \forma{sí, cierto, obvio};
  \item negación: \forma{no, jamás, tampoco};
  \item duda:\footnote{\label{fn:05-cuantificadores}No incluimos en esta lista a
        los adverbios de cantidad porque hoy día se considera que palabras como
        \forma{muy, poco, bastante} o \forma{demasiado} integran la clase de los
        cuantificadores \parencite{leonetti2007,bosque2009}.}
        \forma{acaso, quizás, posiblemente}.
\end{itemize}

Desde el punto de vista de su naturaleza categorial, la \textit{NGLE}
\parencite[§\,30.2j]{rae2025} distingue dos subclases: adverbios de base léxica
y adverbios gramaticales.

\begin{itemize}
  \item Los \termino{adverbios de base léxica} incluyen la mayoría de los
        terminados en \forma{-mente} (\forma{rápidamente, estrepitosamente}) y
        otros no derivados, como \forma{bien} o \forma{temprano}. Los adverbios
        en \forma{-mente} provienen de adjetivos calificativos y admiten la
        paráfrasis mediante el SP \forma{de modo / de manera} + el A base:
        \forma{astuto > astutamente / de manera astuta}. Algunos adverbios en
        \forma{-mente} tienen una forma corta, sin el sufijo \forma{-mente}:
        \forma{trabajaron duramente / trabajaron duro; habló claramente / habló
        claro; aprendió rápidamente / aprendió rápido}. Estas variantes cortas
        se conocen también como \termino{adverbios adjetivales}, porque tienen la
        forma del adjetivo del que provienen, aunque se diferencian de ellos por
        la ausencia de flexión en género (\forma{Pedro habló claro / Julia habló
        claro}) (v. \capref{4}, §\ref{sec:04-conversion}). El empleo de estos
        adverbios cortos tiene algunas restricciones, ya que no aparecen con
        cualquier tipo de verbo, sino con algunos (\forma{Se portó genial,
        apuntó alto, jugaron sucio, hiló fino}). Varios de estos adverbios son
        propios del español de América: \forma{breve} (= \sig{pronto}),
        \forma{sabroso}, (\forma{oler}) \forma{lindo / rico},
        (\forma{equivocarse}) \forma{feo / fiero}; mientras que otros pertenecen
        al registro peninsular: (\forma{sentar algo}) \forma{fatal},
        (\forma{saberse algo}) \forma{fenomenal} (Bosque,
        1990).\pendiente{En la bibliografía figura Bosque (1989)}
  \item Los \termino{adverbios gramaticales} forman grupos cerrados. Se
        reconocen distintas subclases: \termino{demostrativos} (\forma{aquí,
        hoy, así}); \termino{identificativos} o \termino{referenciales}
        (\forma{antes / después, encima / debajo}); \termino{relativos}
        (\forma{cuando, como, donde}); \termino{enfáticos}\footnote{La subclase
        de adverbios enfáticos incluye los que la \textit{NGLE}
        \parencite[§\,30.2j]{rae2025} distribuye en dos subclases:
        interrogativos y exclamativos.} (\forma{cuándo, cómo, adónde}) y
        \termino{focales} (\forma{precisamente, también, solo, concretamente,
        incluso}).\footnote{Según lo explicado en la
        nota~\ref{fn:05-cuantificadores}, tampoco incluimos entre los adverbios
        gramaticales a los cuantificacionales.}\pendiente{El original da
        \textit{quien} entre los relativos}
\end{itemize}

Según su funcionamiento sintáctico, \textcite[724]{kovacci1999} distinguió dos
grupos de adverbios:

\begin{itemize}
  \item los \termino{adverbios nucleares}, que son los que modifican a un
        adjetivo \exref{ex:05-nucleares-a}, a otro adverbio
        \exref{ex:05-nucleares-b} o tienen funciones circunstanciales dentro del
        SV\footnote{Las funciones sintácticas dentro del SV se estudian en el
        \capref{7}.} \exref{ex:05-nucleares-c}:
\end{itemize}

\begin{exe}
  \ex\label{ex:05-nucleares}
    \begin{xlist}
      \ex\label{ex:05-nucleares-a} Es \textit{relativamente} alto.
      \ex\label{ex:05-nucleares-b} Vive \textit{extremadamente} lejos.
      \ex\label{ex:05-nucleares-c} Llegaron muy \textit{rápidamente}.
    \end{xlist}
\end{exe}

\begin{itemize}
  \item los \termino{adverbios periféricos}, que modifican a toda la oración o a
        la modalidad\footnote{La modalidad y el funcionamiento de los adverbios
        periféricos se tratan en el \capref{14}.} \exref{ex:05-perifericos}:
\end{itemize}

\begin{exe}
  \ex\label{ex:05-perifericos} \textit{Probablemente / Lamentablemente}, no
      esperaba esa noticia.
\end{exe}

\paragraph{Los complementos.} Pocos adverbios reciben complementos. Los que
derivan de adjetivos que exigen un complemento, por lo general, no lo heredan
(\forma{propenso a enfermarse} / \juicio{*}{\forma{propensamente a enfermarse}},
\forma{digno de admiración} / \juicio{*}{\forma{dignamente de admiración}})
\parencite{giammatteo2023a}. En cambio, los locativos que provienen de la unión
de un nombre y una preposición (\forma{de-bajo / a-bajo, de-lante / a-delante,
a-dentro, a-fuera, en-cima, a-rriba, al-rededor, en-frente, de-trás}) y algunos
temporales (\forma{antes, después}), que por su naturaleza categorial se ubican
entre los adverbios identificativos o referenciales que vimos más arriba, pueden
recibir complementos. La particularidad de estos adverbios, que son también
conocidos como \termino{adverbios preposicionales}, es que, al igual que las
preposiciones y a diferencia de otros adverbios, son transitivos
\parencite{giammatteo2013}:

\begin{exe}
  \ex
    \begin{xlist}
      \ex Unos amigos de mi hermana se mudaron [\textit{enfrente} de nuestra
          casa].
      \ex Llegaron muy puntuales, [\textit{antes} de la cena].
    \end{xlist}
\end{exe}

No obstante, la relación del adverbio con su complemento es más laxa que la que
mantiene una preposición, con la que el complemento no es omisible; con el
adverbio, en cambio, el complemento puede quedar implícito:

\begin{exe}
  \ex
    \begin{xlist}
      \ex El equipo quedó [afuera (del torneo)].
      \ex No sé lo que pasó porque yo llegué [después (del suceso)].
    \end{xlist}
\end{exe}

\paragraph{Los modificadores.} Un subgrupo de adverbios temporales, locativos y
algunos modales aceptan ser modificados por aposiciones nominales y
preposicionales (\forma{hoy, lunes; allá, en la esquina}), así como también por
oraciones subordinadas introducidas por un pronombre relativo\footnote{Las
subordinadas relativas, que son las introducidas por pronombres relativos como
\forma{que, cuando, donde}, etc., se introducen en el \capref{13}.}
(\forma{ayer, cuando nos vimos; allí, donde nos encontramos}).

\paragraph{Los especificadores.} Operan como especificadores de los adverbios
los cuantificadores de grado (\forma{muy lentamente, sumamente mal}), SSNN con
valor de grado o medida, que van siempre antepuestos (\forma{dos horas más tarde
de la reunión, tres metros delante de la marca}), y los focalizadores
(\forma{solo hoy, recién ahora}). Sin embargo, no todas las subclases de
adverbios admiten especificadores; así, por ejemplo, solo los graduables
aceptan cuantificadores de grado, de modo que podemos decir \forma{muy
lentamente}, pero no \juicio{*}{\forma{muy históricamente}}; \forma{bastante
lejos}, pero no \juicio{*}{\forma{bastante siempre}}.

\begin{ejercicio}\label{ej:05-adverbios}
Lea atentamente los siguientes fragmentos y resuelva las consignas a
continuación, según se muestra en el ejemplo resuelto:

\begin{consignas}
  \item identifique los adverbios, indique su estatus categorial y
        clasifíquelos semánticamente;
  \item indique si son núcleos o modificadores dentro de los sintagmas en los
        que aparecen.
\end{consignas}

% Fragmentos literarios protegidos: se incorporan a mano desde el original.
\textoaincorporar{1. fragmento de Saki, adaptado de \textit{Cuentos escogidos}
(Claridad, 2007), con los adverbios \textit{débilmente, definidamente, más bien,
elaboradamente}; 2. fragmento sobre el percebe de \textit{La naturaleza}
(Sopena, 1981), con \textit{sumamente, dentro}; 3. fragmento de
«Instrucciones para subir una escalera», de J.~Cortázar (\textit{Historias de
cronopios y de famas}, 1962), con \textit{de frente, atrás,
particularmente}}

\begin{resuelto}
a) \forma{débilmente}: adverbio de base léxica de modo. b) Modifica al adjetivo
\forma{iluminado} dentro del SA [\forma{débilmente iluminado}].
\end{resuelto}
\end{ejercicio}

\begin{ejercicio}\label{ej:05-p12}
Analice en forma completa mediante encorchetamiento los siguientes SSNN
extraídos de distintas noticias del diario \textit{Página/12} de Buenos Aires,
según se muestra en el ejemplo resuelto~1. Tenga en cuenta que los nombres
eventivos, relacionales y de representación llevan complemento.

\begin{exe}
  \ex El desecho de otros minerales potencialmente comercializables.
  \ex Una constante dolarización de excedentes económicos.
  \ex Una corriente de presión sobre este segmento del mercado de cambio.
  \ex El ministro de Turismo, un destacado empresario.
  \ex Un dólar financiero para el turista extranjero.
  \ex Un punto de inflexión en ese estilo de funcionamiento político de la
      coalición de gobierno.
\end{exe}

\begin{resuelto}
1. [\etq{SN} el\etq{D/Esp} desecho\etq{N/n} [\etq{SP/M} de\etq{P/n}
[\etq{SN/Comp} otros\etq{D/Esp} minerales\etq{N/n} [\etq{SA/M}
potencialmente\etq{Adv/M} comercializables\etq{A/n}]]]]
\end{resuelto}
\end{ejercicio}

\begin{recuadro}[Reflexión final]
En este capítulo vimos que las palabras, al igual que otras unidades de la
lengua, no solo se disponen de manera lineal, sino que se organizan en
estructuras jerárquicas que reflejan dependencias y relaciones internas. La
identificación de constituyentes sintácticos y la distinción entre palabras
léxicas y funcionales y sus posibilidades de proyección nos permitió
caracterizar a los sintagmas como unidades intermedias, así como también
describir y explicar su conformación estructural. Para ello, nos basamos en la
teoría X\prima, que sostiene que todos los sintagmas comparten una
configuración estructural endocéntrica, en la que el núcleo selecciona
complementos y admite adjuntos y especificadores, en función de factores tanto
léxicos como sintácticos. Nos detuvimos particularmente en la caracterización de
los SSNN, SSAA, SSPP y SSADV, lo cual se ampliará en los próximos capítulos con
la presentación del SV. El ámbito de la sintaxis, en conclusión, no solo aporta
criterios para definir, describir, analizar y fundamentar el funcionamiento y
la organización interna de sus unidades, sino que también ofrece un marco
teórico que permite establecer una conexión entre forma y significado, lo cual,
sin duda, enriquece nuestra comprensión del funcionamiento del sistema de la
lengua.
\end{recuadro}

% --- Secciones finales -------------------------------------------------
\lecturasrecomendadas

\solucionario

\begin{solucion}{\ref{ej:05-segmentar}}
2. [[Nuestra tía] [traía regalos constantemente]]. 3. [[Muchos estudiantes]
[prefieren [la virtualidad]]]. 4. [[Las medidas] [preocupan [a [la
población]]]]. 5. [[Los jugadores] [ganaron [la copa] [en 2022]]]. 6. [[Aquellas
profesoras] [incorporaron [en [sus clases]] [herramientas tecnológicas]]].
\end{solucion}

\begin{solucion}{\ref{ej:05-homonimia}}
En este ejercicio los constituyentes se identifican, pero no se analizan
exhaustivamente. a) 2. i) [Llenó [la bolsa] [de papel]]; ii) [Llenó [la bolsa
de papel]]. 3. i) [Encontraron [el juguete] [roto]]; ii) [Encontraron [el
juguete roto]]. b) 2. En i) se puede considerar que [la bolsa] y [de papel] son
dos constituyentes independientes. Se interpreta que \sig{la bolsa se llenó con
papel}. En ii), [de papel] modifica a [la bolsa], por lo que conforman un solo
constituyente. Se interpreta que \sig{la bolsa está hecha de papel}. 3. En i) se
puede considerar que [el juguete] y [roto] son dos constituyentes
independientes. Se interpreta que \sig{el juguete fue encontrado roto, pero
antes no lo estaba}. En ii), [roto] modifica a [el juguete], por lo que
conforman un solo constituyente. Se interpreta que \sig{se encontró el juguete
que ya estaba roto}. c) El efecto que produce la homonimia estructural es el de
ambigüedad, en este caso, de tipo sintáctico.
\end{solucion}

\begin{solucion}{\ref{ej:05-expansion}}
Respuestas posibles. 2. El renombrado museo de Ciencias Naturales de la ciudad
de La Plata recibió los nuevos ejemplares de restos fósiles encontrados por los
paleontólogos. / El museo de Arte Moderno de Buenos Aires, una institución
pública dependiente del Gobierno de la Ciudad de Buenos Aires, recibió los
ejemplares de las producciones artísticas más recientes. 3. Mis primas llegaron
el sábado a la noche muy cansadas por el viaje.\pendiente{El original dice
«cansados»} / Mis primas llegaron inmediatamente. 4. La última temporada de mi
serie preferida trajo sorpresas que no esperaba. / La temporada de invierno
trajo numerosas sorpresas referidas a la moda. 5. Los vehículos sin el
dispositivo de TelePASE circulaban muy lentamente. / Los vehículos antiguos
circulaban más lentamente.
\end{solucion}

\begin{solucion}{\ref{ej:05-sustitucion}}
2. Reducción: Varios ministros firmaron un acuerdo con el embajador.
Pronominalización: Ellos lo firmaron con él.
\end{solucion}

\begin{solucion}{\ref{ej:05-encorchetar}}
2. [[Los egresados] [buscaron [con entusiasmo] [un salón de fiestas]]]. Ellos
buscaron con entusiasmo un salón de fiestas. Los egresados lo buscaron con
entusiasmo. Los egresados buscaron así un salón de fiestas. Ellos lo buscaron
así. 3. [[Nuestros candidatos] [presentaron [sus propuestas] [a los
contribuyentes]]]. Ellos presentaron sus propuestas a los contribuyentes.
Nuestros candidatos las presentaron a los contribuyentes. Nuestros candidatos
les presentaron sus propuestas. Ellos se\footnote{\label{fn:05-se}La forma
\forma{se} delante de \forma{lo / la} es una variante de realización del clítico
\forma{le}, es decir, es un alomorfo (v. \capref{10}).} las presentaron. 4. [[Los
fuertes temporales] [provocaron [inundaciones] [en toda la ciudad]]].
Provocaron inundaciones en toda la ciudad. Aclaración: al tratarse de un sujeto
no animado, usualmente no se reemplaza por la forma pronominal \forma{ellos},
sino que se prefiere dejar el sujeto tácito.\footnote{La información sobre el
sujeto tácito se desarrolla en el \capref{6}.} Los fuertes temporales las
provocaron en toda la ciudad. Los fuertes temporales provocaron inundaciones
allí.
\end{solucion}

\begin{solucion}{\ref{ej:05-preguntas}}
2. ¿Quiénes llegaron ayer por la tarde a la ciudad para visitarnos? ¿Qué
hicieron mis primos? ¿Cuándo llegaron mis primos a la ciudad para visitarnos?
¿Adónde llegaron mis primos ayer por la tarde para visitarnos? ¿Para qué
llegaron mis primos a la ciudad ayer por la tarde? 3. ¿Quién es muy exigente con
los exámenes? ¿Cómo es la nueva profesora? ¿Con qué es muy exigente la nueva
profesora?
\end{solucion}

\begin{solucion}{\ref{ej:05-bioy}}
¿Con quién paseaba el escultor? ¿Por dónde paseaba el escultor con el tirano?
¿Dónde presentó el escultor su última obra? ¿Quién presentó su última obra más
allá del laberinto para los extranjeros ilustres, en el extremo de la alameda de
los filósofos decapitados? ¿Qué presentó el escultor más allá del laberinto para
los extranjeros ilustres, en el extremo de la alameda de los filósofos
decapitados?
\end{solucion}

\begin{solucion}{\ref{ej:05-pares}}
a) 2. i. Eliana encontró [el auto] [de repente] / ii. Eliana encontró [el auto
de su primo]. 3. i. El hombre tocó [al perro] [con la mano] / ii. El hombre tocó
[al perro con sarna]. 4. i. Juan abrió [la puerta] [de una patada] / ii. Juan
abrió [la puerta de la habitación].

b) 2. En i., \forma{de repente} y \forma{el auto} son dos constituyentes porque
se mueven separadamente (\forma{De repente Eliana encontró el auto. Eliana
encontró el auto de repente}); se pronominalizan independientemente
(\forma{Eliana lo encontró de repente. Eliana lo encontró así}); responden a
diferentes preguntas (\forma{¿Qué encontró Eliana de repente? ¿Cómo encontró
Eliana el auto?}). En ii., \forma{el auto de su primo} es un solo constituyente
porque se mueve conjuntamente (\forma{El auto de su primo encontró Eliana});
admite una única pronominalización (\forma{Eliana lo encontró}); responde a una
única pregunta (\forma{¿Qué auto encontró Eliana?}). 3. En i., \forma{al perro} y
\forma{con la mano} son dos constituyentes porque admiten las tres pruebas
propuestas: se mueven separadamente (\forma{El hombre tocó con la mano al perro.
Con la mano el hombre tocó al perro}); se pronominalizan independientemente
(\forma{El hombre lo tocó con la mano. El hombre tocó al perro así}); responden
a diferentes preguntas (\forma{¿A quién tocó el hombre con la mano? ¿Con qué tocó
el hombre al perro?}). En ii., \forma{al perro con sarna} es un solo
constituyente porque se mueve conjuntamente (\forma{Al perro con sarna tocó el
hombre}); admite una sola pronominalización (\forma{El hombre lo tocó}); responde
a una única pregunta (\forma{¿A quién tocó el hombre?}). 4. En i., \forma{la
puerta} y \forma{de una patada} son diferentes constituyentes porque admiten las
tres pruebas propuestas: se mueven separadamente (\forma{De una patada Juan abrió
la puerta. La puerta abrió Juan de una patada}); responden a distintas preguntas
(\forma{¿Qué abrió Juan de una patada? ¿Cómo abrió Juan la puerta?}); se
pronominalizan independientemente (\forma{Juan abrió la puerta así. Juan la
abrió de una patada}). En ii., \forma{la puerta de la habitación} es un
constituyente porque se mueve conjuntamente (\forma{La puerta de la habitación
abrió Juan}), admite una única pronominalización (\forma{Juan la abrió}) y
responde a una única pregunta (\forma{¿Qué abrió Juan?}).
\end{solucion}

\begin{solucion}{\ref{ej:05-nucleo}}
2. a) El núcleo es \forma{bueno}. b) \forma{Bueno} es un A y proyecta un SA. 3.
a) El núcleo es \forma{cerca}. b) \forma{Cerca} es un Adv y proyecta un SAdv. 4.
a) El núcleo es \forma{proporcional}. b) \forma{Proporcional} es un A y proyecta
un SA. 5. a) El núcleo es \forma{para}. b) \forma{Para} es una P y proyecta un
SP. 6. a) El núcleo es \forma{inquietud}. b) \forma{Inquietud} es un N y
proyecta un SN. 7. a) El núcleo es \forma{armario}. b) \forma{Armario} es un N y
proyecta un SN. 8. a) El núcleo es \forma{sin}. b) \forma{Sin} es una P y
proyecta un SP.
\end{solucion}

\begin{solucion}{\ref{ej:05-de}}
2. [\etq{SN} varios\etq{Q/Esp} colaboradores\etq{N/n} expertos\etq{A/M}]. 3.
[\etq{SAdv} arriba\etq{Adv/n}\footnote{\forma{Arriba} pertenece al subgrupo de
adverbios denominados identificativos o referenciales, que son transitivos y
llevan un complemento. Los tratamos en §\ref{sec:05-sadv}.} [\etq{SP/Comp}
de\etq{P/n} [\etq{SN/Comp} la\etq{D/Esp} heladera\etq{N/n}]]]. 4. [\etq{SP}
sin\etq{P/n} [\etq{SN/Comp} comentarios\etq{N/n}]].
\end{solucion}

\begin{solucion}{\ref{ej:05-comp-adj}}
2. a) El núcleo es \forma{feria} y proyecta un SN. b) \forma{Del libro} es un SP
no exigido por el núcleo, que funciona como M. 3. a) El núcleo es
\forma{similar} y proyecta un SA. b) \forma{A los demás casos} es un SP exigido
por el núcleo y funciona como Comp. 4. a) El núcleo es \forma{por} y proyecta un
SP. b) \forma{La izquierda} es un SN exigido por el núcleo y funciona como Comp.
5. a) El núcleo es \forma{bajo} y proyecta un SP. b) \forma{La mesada} es un SN
que funciona como Comp, ya que es exigido por el núcleo. 6. a) El núcleo es
\forma{antes} y proyecta un SAdv. b) \forma{De las diez} es un SP exigido por el
núcleo y funciona como Comp.
\end{solucion}

\begin{solucion}{\ref{ej:05-titulos}}
3. a) Es un SN porque el núcleo es el sustantivo \forma{jardín}. b) El D
\forma{el} es un especificador y \forma{de los cerezos} es un SP no requerido
por el núcleo, un modificador adjunto. 4. a) Es un SN porque el núcleo es el
sustantivo \forma{años}. b) El Q \forma{cien} es un especificador y \forma{de
soledad} es un SP no requerido por el núcleo, un modificador adjunto. 5. a) Es
un SAdv porque el núcleo es el adverbio \forma{ágilmente}. b) El núcleo no
presenta modificadores ni complementos. 6. a) Es un SN porque el núcleo es el
sustantivo \forma{verano}. b) El D \forma{ese} es un especificador y \forma{a
oscuras} es un SP no requerido por el núcleo, un modificador adjunto. 7. a) Es
un SA porque el núcleo es el adjetivo \forma{terrenal}. b) El núcleo no presenta
modificadores ni complementos. 8. a) Es un SP porque el núcleo es la preposición
\forma{desde}. b) \forma{La ventana} es un SN exigido por el núcleo, un Comp.
\end{solucion}

\begin{solucion}{\ref{ej:05-vf}}
2. F(also). Tanto las palabras léxicas como las gramaticales pueden ser núcleos
de sintagmas. 3. F(also). El núcleo mantiene diferentes relaciones con los
constituyentes del sintagma: los complementos son los exigidos por el núcleo,
pero los adjuntos son opcionales y los especificadores determinan o cuantifican
al núcleo y a sus complementos y modificadores. 4. V(erdadero). Los complementos
tienen una relación de dependencia estructural con el núcleo y, además,
completan su significado.
\end{solucion}

\begin{solucion}{\ref{ej:05-buenaformacion}}
2. Aunque los sustantivos propios no necesitan ir precedidos de un
determinativo, un sintagma como \forma{el Mario} puede ser aceptable en ciertas
variedades dialectales para referirse a alguien conocido por todos. 3. Es
gramatical porque el sustantivo común \forma{ingredientes} va precedido de un
determinante que lo especifica y del complemento que requiere por su carácter
relacional. 4. Es gramatical porque los nombres de materia o incontables no
requieren ir precedidos por un determinativo en ciertas posiciones sintácticas,
como la de COD. 5. Los nombres abstractos tienen un comportamiento sintáctico
que se asimila al de los de materia, de modo que pueden ir sin determinativo en
ciertas posiciones sintácticas, como la de complemento
preposicional.\footnote{Para ampliar sobre las diferencias de funcionamiento de
las distintas subclases de sustantivos se puede consultar
\textcite{giammatteo2009}.}
\end{solucion}

\begin{solucion}{\ref{ej:05-ratones}}
{}[\etq{SN} el\etq{D/Esp} queso\etq{N/n}] / [\etq{SN} una\etq{Q} compañía\etq{N/n}
[\etq{SP/Comp} de\etq{P/n} [\etq{SN/Comp} control\etq{N/n} [\etq{SP/Comp}
de\etq{P/n} [\etq{SN/Comp} plagas\etq{N/n}]]]]] / [\etq{SN} los\etq{D/Esp}
alimentos\etq{N/n} [\etq{SP/M} con\etq{P/n} [\etq{SN/Comp} alto\etq{A/M}
contenido\etq{N/n} [\etq{SP/Comp} de\etq{P/n} [\etq{SN/Comp} azúcares]]]]].
Aclaración: \forma{compañía} y \forma{contenido} exigen complemento porque son
nombres relacionales; \forma{control} exige un complemento porque es un nombre
eventivo.
\end{solucion}

\begin{solucion}{\ref{ej:05-titulares}}
a) Noticia 1. (3) [\etq{SN} el\etq{D/Esp} columnista\etq{N/n} [\etq{SN/ApR}
[\etq{N/n} Lauro Jardim]] [\etq{SP/Comp} de\etq{P/n} [\etq{SN/Comp}
el\etq{D/Esp} diario\etq{N/n} [\etq{SN/ApR} O Globo]]]]. \forma{Columnista} exige
complemento porque es un nombre relacional. (4) [\etq{SN} Madonna\etq{N/n},
[\etq{SP/M} de\etq{P/n} [\etq{SN/Comp} 65\etq{Q/Esp} años\etq{N/n}]]].
Noticia 2. (1) [\etq{SN} la\etq{D/Esp} primera\etq{A/M} teatrista\etq{N/n}
latinoamericana\etq{A/M}]. (2) [\etq{SN} [\etq{N/n} Lola Arias], [\etq{SN/ApE}
ganadora\etq{N/n} [\etq{SP/Comp} de\etq{P/n} [\etq{SN/Comp} el\etq{D/Esp}
Premio\etq{N/n} Internacional\etq{A/M} Ibsen\etq{SN/ApR}
2024\etq{SN/ApR}]]]].\footnote{La secuencia \forma{Premio Internacional Ibsen
2024} también podría considerarse un compuesto sintagmático que actúa como un
nombre propio (v. \capref{4}, §\ref{sec:04-composicion}).} (3) [\etq{SN}
el\etq{D/Esp} carácter\etq{N/n} transnacional\etq{SA/M} [\etq{SP/M} de\etq{P/n}
[\etq{SN/Comp} la\etq{D/Esp} obra\etq{N/n} [\etq{SP/M} de\etq{P/n}
[\etq{SN/Comp} la\etq{D/Esp} artista\etq{N/n} multidisciplinaria\etq{SA/M}]]]]].

b) \forma{Lauro Jardim}, \forma{del diario O Globo}: restrictivos; \forma{de 65
años}: explicativo; \forma{primera} y \forma{latinoamericana}: restrictivos;
\forma{ganadora del Premio Internacional Ibsen 2024}: explicativo;
\forma{transnacional}: restrictivo; \forma{multidisciplinaria}: restrictivo.
\end{solucion}

\begin{solucion}{\ref{ej:05-lobos}}
a) 2. [\etq{SN} varias\etq{Q/Esp} subespecies\etq{N/n}]. 3. [\etq{SN}
el\etq{D/Esp} famoso\etq{SA/M} aullido\etq{N/n}]. 4. [\etq{SN} los\etq{D/Esp}
miembros\etq{N/n} [\etq{SP/Comp} de\etq{P/n} [\etq{SN/Comp} su\etq{D/Esp}
grupo\etq{N/n}]]]. 5. [\etq{SN} [\etq{SA/M} muy\etq{Q/Esp} raras\etq{A/n}]
ocasiones\etq{N/n}].

b) 2. El Q \forma{varias} puede posponerse al núcleo \forma{subespecies} sin
alterar la gramaticalidad del sintagma, aunque aporta un ligero cambio
estilístico. 3. El A \forma{famoso} puede posponerse al N sin alterar la
gramaticalidad del sintagma, aunque pasa a tener valor restrictivo. 4. El D
\forma{los} no puede posponerse al núcleo por razones
fonológicas\footnote{V. \capref{3}, §\ref{sec:03-acento}.} y la anteposición del
complemento \forma{de su grupo} al núcleo produce un
hipérbaton.\footnote{Figura retórica en la que se altera la sintaxis habitual de
una oración.} 5. El SA puede posponerse al núcleo (\forma{ocasiones muy raras}),
adoptando valor restrictivo.
\end{solucion}

\begin{solucion}{\ref{ej:05-slogans}}
2. [\etq{SA} único\etq{A/n} [\etq{SP/M} en\etq{P/n} [\etq{SN/Comp} el\etq{D/Esp}
mundo\etq{N/n}]]] / [\etq{SA} único\etq{A/n}]. 3. [\etq{SA} buenos\etq{A/n}
[\etq{SP/M} de\etq{P/n} [\etq{SAdv/Comp} [\etq{LAdv/n} punta a
punta]]]].\footnote{\forma{De punta a punta} funciona como una unidad porque es
una locución adverbial. Las locuciones adverbiales se tratan en el
\capref{11}.}\pendiente{El análisis separa \textit{de} de la locución que la
nota declara unitaria} 4. [\etq{SA} juntos\etq{A/n}].
\end{solucion}

\begin{solucion}{\ref{ej:05-posicion}}
a) 3. i) \forma{una tarde estupenda} / ii) \forma{una estupenda tarde}. 4. i)
\forma{el discurso político} / ii) \juicio{*}{\forma{el político discurso}}. 5.
i) \forma{ese pobre hombre} / ii) \forma{ese hombre pobre}. 6. i) \forma{un rico
chico} / ii) \forma{un chico rico}. 7. i) \forma{una enfermedad terrible} / ii)
\forma{una terrible enfermedad}. 8. i) \forma{la educación sexual} / ii)
\juicio{*}{\forma{la sexual educación}}.

b) En 2, 4 y 8, el cambio de posición produce agramaticalidad; en 3 y 7, el
cambio de posición no produce ni agramaticalidad ni cambio de significado del A,
solo aporta un matiz estilístico; en 5 y 6, el cambio de posición no produce
agramaticalidad, pero cambia el significado del A.

c) 5. i) \sig{sin importancia, insignificante} / ii) \sig{sin
dinero}.\pendiente{El original invierte i y ii} 6. i) \sig{lindo} / ii)
\sig{adinerado}.

d) Los A calificativos pueden anteponerse o posponerse al núcleo al que
modifican. Por este motivo, en ninguna de las series que los contienen los
sintagmas se convierten en agramaticales. En 5 y 6 el cambio de posición altera
el significado del adjetivo, pero en 3 y 7 no lo altera y solo manifiesta un
matiz estilístico. Los A relacionales no admiten cambio de posición y siempre
deben posponerse al nombre.
\end{solucion}

\begin{solucion}{\ref{ej:05-expansionA}}
a) Respuestas posibles: 2. \forma{propenso a enfermarse / propenso a decir
mentiras}. 3. \forma{muy habitual por las mañanas / sumamente habitual}. 4.
\forma{conforme con su calificación / conforme a la ley}. 5. \forma{deseoso de
verte / deseoso de comer algo dulce}. b) Los adjetivos \forma{propenso,
conforme} y \forma{deseoso} exigen un complemento encabezado por una
preposición: respectivamente, \forma{a}; \forma{con} o \forma{a}; y
\forma{de}.\pendiente{El original dice «con o por» para
\textit{conforme}}
\end{solucion}

\begin{solucion}{\ref{ej:05-sp-modifica}}
2. El SP encabezado por \forma{por} modifica a \forma{eligieron}, que es un
verbo; el SP encabezado por \forma{sobre} es complemento de \forma{por}, que es
una preposición. 3. El SP encabezado por \forma{de} modifica a \forma{cerca},
que es un Adv. 4. El SP encabezado por \forma{con} modifica a \forma{salir}, que
es un V. 5. El SP encabezado por \forma{por} modifica a \forma{feliz}, que es un
A.
\end{solucion}

\begin{solucion}{\ref{ej:05-de-sp}}
a) 2. [\etq{SP} de\etq{P/n} [\etq{SN/Comp} gramática\etq{N/n}]]. 3. [\etq{SP}
de\etq{P/n} [\etq{SN/Comp} menta\etq{N/n}]]. 4. No supieron [\etq{SP/Comp}
de\etq{P/n} [\etq{SN/Comp} aquella\etq{D/Esp} mujer\etq{N/n}]] por varios días.
5. [\etq{SP} de\etq{P/n} [\etq{SN/Comp} conocimiento]].

b) 2. \forma{De} modifica al N \forma{libro} e indica \sig{tema}. 3. \forma{De}
modifica al N \forma{caramelos} e indica \sig{materia o sustancia}. 4.
\forma{De} es una preposición exigida por el V \forma{saber}. La P introduce un
complemento que modifica el significado del V: \forma{saber de (alguien)} es
\sig{tener información sobre esa persona}. 5. \forma{De} es una P exigida por el
adjetivo \forma{ávidos}, que introduce su complemento.
\end{solucion}

\begin{solucion}{\ref{ej:05-preposiciones}}
Respuestas posibles: 2. \sig{por encima de}. 3. \sig{en medio de}. 4.
\sig{carentes de}. 5. \sig{con el fin de}. 6. \sig{en dirección a}. 7.
\sig{acerca de}. 8. \sig{en oposición a}. 9. \sig{con destino a}. 10. \sig{con,
a través de}. 11. \sig{luego de}.
\end{solucion}

\begin{solucion}{\ref{ej:05-adverbios}}
1. \forma{definidamente}: a) adverbio de base léxica de modo; b) es núcleo del
SAdv [\forma{definidamente}], que modifica al V dentro del SV. \forma{más bien}:
a) locución adverbial de modo; b) especificador del Adv \forma{elaboradamente}
en el SAdv [\forma{más bien elaboradamente}]. \forma{elaboradamente}: a)
adverbio de base léxica de modo; b) modifica al A \forma{rígida} en el SA
[\forma{más bien elaboradamente rígida}]. 2. \forma{sumamente}: a) adverbio de
base léxica de modo; b) es especificador del A en el SA [\forma{sumamente
interesante}]. \forma{dentro}: a) adverbio gramatical referencial de lugar; b)
es núcleo del SAdv [\forma{dentro del caparazón del animal adulto}]. 3.
\forma{frente}: a) adverbio gramatical referencial de lugar; b) complemento de
la P \forma{de} en el SP [\forma{de frente}]. \forma{atrás}: a) adverbio
gramatical referencial de lugar; b) complemento de la P \forma{hacia} en el SP
[\forma{hacia atrás}]. \forma{particularmente}: a) adverbio de base léxica de
modo; b) es especificador del A en el SA [\forma{particularmente incómodas}].
\end{solucion}

\begin{solucion}{\ref{ej:05-p12}}
2. [\etq{SN} una\etq{Q/Esp} constante\etq{A/M} dolarización\etq{N/n} [\etq{SP/M}
de\etq{P/n} [\etq{SN/Comp} excedentes\etq{N/n} económicos\etq{A/M}]]].
3. [\etq{SN} una\etq{Q/Esp} corriente\etq{N/n} [\etq{SP/Comp} de\etq{P/n}
[\etq{SN/Comp} presión\etq{N/n} [\etq{SP/Comp} sobre\etq{P/n} [\etq{SN/Comp}
este\etq{D/Esp} segmento\etq{N/n} [\etq{SP/Comp} de\etq{P/n} [\etq{SN/Comp}
el\etq{D/Esp} mercado\etq{N/n} [\etq{SP/Comp} de\etq{P/n} [\etq{SN/Comp}
cambio]]]]]]]]].
4. [\etq{SN} el\etq{D/Esp} ministro\etq{N/n} [\etq{SP/Comp} de\etq{P/n}
[\etq{SN/Comp} Turismo]], [\etq{SN/Ap} un\etq{Q/Esp} destacado\etq{A/M}
empresario\etq{N/n}]].
5. [\etq{SN} un\etq{Q/Esp} dólar\etq{N/n} financiero\etq{A/M} [\etq{SP/M}
para\etq{P/n} [\etq{SN/Comp} el\etq{D/Esp} turista\etq{N/n}
extranjero\etq{A/M}]]].
6. [\etq{SN} un\etq{Q/Esp} punto\etq{N/n} [\etq{SP/M} de\etq{P/n} [\etq{SN/Comp}
inflexión\etq{N/n} [\etq{SP/Comp} en\etq{P/n} [\etq{SN/Comp} ese\etq{D/Esp}
estilo\etq{N/n} [\etq{SP/Comp} de\etq{P/n} [\etq{SN/Comp} funcionamiento\etq{N/n}
político\etq{A/M} [\etq{SP/Comp} de\etq{P/n} [\etq{SN/Comp} la\etq{D/Esp}
coalición\etq{N/n} [\etq{SP/M} de\etq{P/n} [\etq{SN/Comp}
gobierno\etq{N/n}]]]]]]]]]]].\pendiente{El original dice «su gobierno»}
\end{solucion}

% =====================================================================
%  capitulos/cap06.tex — Capítulo 6. La estructura de la oración
%  Autora: Mabel Giammatteo
%  Fuente: Capitulo_06_normalizado.docx
% =====================================================================
\capitulo{La estructura de la oración}{Mabel Giammatteo}

\section{Introducción}

Según se ha visto en los capítulos anteriores, la lengua es esencialmente un
mecanismo combinatorio. A partir de elementos mínimos sin significado, como los
fonemas y las sílabas, se van conformando unidades significativas cada vez más
amplias: primero morfemas y, luego, palabras, las que a su vez se proyectan en
unidades intermedias, como los sintagmas (SS). Ahora bien, ¿cómo ubicamos en esta
secuencia de complejidad creciente a la oración? En este capítulo mostramos que
la oración no es solo el resultado de la combinación de un SN, que funciona como
sujeto, y de un SV, que constituye el predicado, sino que también es necesario
dar cuenta de cómo se produce la vinculación entre estos dos constituyentes.
Tanto el tiempo (T), que ancla la oración en un momento dado, como la relación
de concordancia (Conc) entre sujeto y predicado resultan fundamentales para la
constitución de la oración. Es así que ambas categorías ---T y Conc---, reunidas
en la categoría funcional de flexión (F), constituyen el núcleo de la oración,
lo cual convierte a esta unidad en un sintagma de flexión (SF). Se explica
también en qué consiste la relación de predicación y por qué en algunas lenguas
se puede omitir el sujeto y en otras no. Por último, presentamos el concepto de
capa o nivel oracional, que permite dar cuenta de la estructura tripartita de
la oración.

\section{Hacia la definición de oración}\label{sec:06-definicion}

Como se muestra en \exref{ex:06-muchacho}, a pesar de que el contenido puede
resultar más o menos comprensible, la mera yuxtaposición de un SN y un SV no
constituye por sí misma una oración bien formada en español.

\begin{exe}
  \ex[*]{\label{ex:06-muchacho} Muchacho correr.}
\end{exe}

Para explicar la agramaticalidad de \exref{ex:06-muchacho}, podemos, por una
parte, recordar que, para funcionar adecuadamente dentro del SN sujeto de la
oración, el N común necesita estar especificado, es decir que debe ser
antecedido por algún tipo de determinante (D) ---\forma{el / este / tu
muchacho}--- o cuantificador (Q) ---\forma{un / algún / cualquier muchacho}---
(v. \capref{5}, §\ref{sec:05-sn}). Por otra parte, el V necesita manifestar el
tiempo (T) de ocurrencia del evento, como se muestra en \exref{ex:06-corre}.

\begin{exe}
  \ex\label{ex:06-corre} El muchacho corre / corrió / correrá.
\end{exe}

Además, ambos, N y V, deben estar vinculados entre sí mediante la concordancia
(Conc) en número y persona, como muestra la gramaticalidad de
\exref{ex:06-conc-a} frente a oraciones agramaticales como las de
\exref{ex:06-conc-b} y \exref{ex:06-conc-c}.

\begin{exe}
  \ex\label{ex:06-conc}
    \begin{xlist}
      \ex\label{ex:06-conc-a} El muchacho corrió.
      \ex[*]{\label{ex:06-conc-b} El muchacho corrieron.}
      \ex[*]{\label{ex:06-conc-c} Los muchachos corrió.}
    \end{xlist}
\end{exe}

\begin{destacado}
Según vemos, la vinculación entre el SN sujeto y el SV predicado se produce
gracias a las categorías de T y Conc, reunidas en la flexión (F), que se
manifiesta en la desinencia verbal.
\end{destacado}

Aunque se alojan en la desinencia verbal, tanto el T como la Conc trascienden al
verbo. El T es una categoría que afecta a toda la oración; por eso en
\exref{ex:06-conc-a} podemos decir que la oración (y no solo el verbo) está en
pasado.\pendiente{El original remite a (1)} En cuanto a la Conc, es la
categoría que pone en relación al verbo con el sujeto. La F es una categoría
funcional porque, al igual que las palabras funcionales que se estudiaron en el
capítulo anterior, transmite un significado gramatical: la categoría de tiempo
(T) expresa en qué momento ocurre el evento, y la de Conc manifiesta las
categorías de persona y número, que indican quiénes y cuántos intervienen en el
suceso expresado por el verbo.

Pero ¿qué clase de sintagma da como resultado la unión de un SN y un SV?
Recordemos que, según ya se ha visto (v. \capref{5}, §\ref{sec:05-sintagmas}),
todos los sintagmas tienen núcleo y que es el núcleo el que aporta la categoría
a todo el sintagma. Pues bien, dado que la F convierte a la unión de un SN y un
SV en oración (O), al aportarle los rasgos de tiempo (T) y concordancia (Conc),
la categoría funcional F se considera el núcleo de la O.

\begin{destacado}
En consecuencia, la O es un SF (sintagma de flexión) y se define formalmente como
una unidad de construcción que articula un \textbf{SN} y un \textbf{SV} en otro
sintagma de núcleo funcional: el \textbf{SF}.
\end{destacado}

Desde el punto de vista de la estructura sintagmática, el núcleo F toma como su
complemento al SV y aloja en su Esp al SN sujeto (aspecto que precisaremos más
adelante), lo que se puede representar mediante encorchetamiento
\exref{ex:06-sf-a} o con un diagrama arbóreo \exref{ex:06-sf-b}.

\begin{exe}
  \ex\label{ex:06-sf}
    \begin{xlist}
      \ex\label{ex:06-sf-a} {}[\etq{SF/O} [\etq{Esp/SN/Suj} \ldots{}]
          [\etq{F\prima} [\etq{F/n} T, Conc] [\etq{Comp/SV/Pred} \ldots{}]]]
      \ex\label{ex:06-sf-b}
          \begin{forest} gramatica
            [SF / O
              [SN\\(Esp / Suj)]
              [F\prima
                [{F\\T, Conc}]
                [SV\\(Comp / Pred)]]]
          \end{forest}
    \end{xlist}
\end{exe}

\begin{ejercicio}\label{ej:06-determinante}
Compare los sujetos de las oraciones de 1 con los de 2. Proponga alguna
hipótesis para explicar la falta del determinante (D) en los SSNN sujeto de las
oraciones de 2.

\textit{Dato adicional}: tenga en cuenta las subclases del nombre.

\begin{exe}
  \ex
    \begin{xlist}
      \ex Las ratas son las primeras en abandonar el barco.
      \ex Varios niños faltaron hoy a la escuela.
    \end{xlist}
  \ex
    \begin{xlist}
      \ex Pedro llegó cuando nadie lo esperaba.
      \ex Buenos Aires es una ciudad magnífica.
    \end{xlist}
\end{exe}

\textbf{Pista}: en 1, la omisión del D (\juicio{*}{\forma{Ratas son las primeras
en abandonar el barco}}) produce una oración agramatical.
\end{ejercicio}

\begin{ejercicio}\label{ej:06-oracion}
Considere las siguientes oraciones y explique lo pedido más abajo.

\begin{exe}
  \ex Mi padre me prometió un helado.
  \ex Me prometió un helado.
\end{exe}

\begin{consignas}
  \item Fundamente por qué 1 es una oración. Identifique el núcleo de 1 y haga
        la representación entre corchetes.
  \item ¿Considera que 2 también es una oración? ¿Qué particularidad presenta
        respecto de 1?
\end{consignas}

\begin{resuelto}[Punto resuelto]
a) 1 es una O porque contiene un SN sujeto y un SV predicado, vinculados por un
núcleo funcional: la F, que manifiesta T (pasado) y Conc (3.ª pers.~sg.).
Encorchetamiento: [\etq{SF/O} [\etq{SN/Suj} mi\etq{D/Esp} padre\etq{N/n}]
[\etq{F\prima} F\etq{n} [\etq{SV/Pred} me prometió un helado]]].
\end{resuelto}
\end{ejercicio}

\section{La relación sujeto-predicado}\label{sec:06-sujeto}

\subsection{Acerca del sujeto tácito o nulo}\label{sec:06-tacito}

Aunque, según acabamos de plantear, la buena formación de la oración requiere la
combinación de un SN sujeto con un SV predicado, como sabemos, en español el
sujeto puede elidirse, es decir que puede quedar tácito o no estar expresado
léxicamente. En el \capref{1} (v.~§\ref{sec:01-generativa}) vinculamos esta
posibilidad de omitir o no el sujeto que tienen las lenguas con uno de los
principios reconocidos por la gramática generativa: el \termino{principio de
proyección extendida} (PPE), que establece que todas las oraciones deben poseer
un sujeto del cual se predica.\footnote{Este principio es una extensión del
principio de proyección \parencite{chomsky1981}, según el cual la sintaxis debe
respetar los requisitos de selección de las piezas léxicas en todos los niveles
de representación.} Sin embargo, llama la atención que algunas lenguas puedan
elidirlo, según mostramos en \exref{ex:06-esp} y \exref{ex:06-ita}, pero otras
no, como se ejemplifica en \exref{ex:06-ing} y \exref{ex:06-fra}.

\begin{exe}
  \ex\label{ex:06-esp} Español: \forma{Llegaron tarde.}
  \ex\label{ex:06-ita} Italiano: \forma{Erano in ritardo.}
  \ex\label{ex:06-ing} Inglés: \forma{They arrived late.} /
      \juicio{*}{\forma{Arrived late.}}
  \ex\label{ex:06-fra} Francés: \forma{Ils étaient en retard.} /
      \juicio{*}{\forma{Étaient en retard.}}\pendiente{El original dice
      «étaint»}
\end{exe}

¿Cómo podemos explicar esta diferencia? Según ya hemos visto en el \capref{1},
los principios de carácter universal se realizan mediante parámetros que, por lo
general, presentan opciones binarias, lo cual significa que existen dos
posibilidades de realización, en este caso, del sujeto: la presencia obligatoria
en todas las oraciones, como sucede en inglés y en francés, frente a la
posibilidad facultativa de omitirlo, que se da en lenguas como el español y el
italiano, entre otras. Esta segunda opción se conoce como \termino{parámetro del
sujeto nulo} y las lenguas que permiten la omisión del sujeto se conocen como
\textit{pro-drop} (con caída del pronombre).

Ahora bien, ¿por qué en algunas lenguas es posible la omisión, mientras que
otras la rechazan? Las lenguas \textit{pro-drop} tienen una morfología verbal
rica, que permite reponer el sujeto a partir de las categorías morfológicas de
la desinencia verbal.\footnote{No obstante, hay lenguas como el alemán, con una
rica morfología, que no permiten la omisión, y otras como el chino, que no
expresan categorías de concordancia e igualmente permiten la omisión. No nos
extenderemos aquí sobre estas particularidades, que exceden el nivel inicial de
este texto.} Así, \forma{llegaron} \exref{ex:06-esp} y \forma{erano}
\exref{ex:06-ita} remiten a una tercera persona plural (\forma{ellos / ellas},
para el español, y \forma{loro}, femenino y masculino, para el italiano).

Por lo anteriormente expuesto, en la representación estructural de estas
oraciones se incluye un \textit{pro} (= pronombre) en la posición de sujeto,
como se muestra en \exref{ex:06-pro-a} y se representa en el diagrama arbóreo
\exref{ex:06-pro-b}.\footnote{El uso del triángulo en el esquema arbóreo
indica que, a fin de simplificar la descripción, se omite el análisis interno de
un constituyente. En el caso de \exref{ex:06-pro-b}, no se realiza el análisis
interno del SV \forma{llegaron tarde}.}

\begin{exe}
  \ex\label{ex:06-pro}
    \begin{xlist}
      \ex\label{ex:06-pro-a} {}[\etq{SF/O} \textit{pro}\etq{Esp/Suj}
          [\etq{F\prima} F\etq{n} [\etq{SV/Pred/Comp} llegaron tarde]]]
      \ex\label{ex:06-pro-b}
          \begin{forest} gramatica
            [SF / O
              [\textit{pro}]
              [F\prima
                [F]
                [SV [\textit{llegaron tarde}, roof]]]]
          \end{forest}
    \end{xlist}
\end{exe}

\begin{recuadro}[Para ampliar: ¿qué clase de pronombre es \textit{pro}?]
El \textit{pro} que hace las veces de sujeto en las oraciones de sujeto tácito se
caracteriza como una \termino{categoría vacía} o \termino{nula} que permite que
el PPE se respete. Se trata de un elemento no audible (sin rasgos fonológicos),
pero estructuralmente necesario para rellenar la posición de sujeto. Tal como se
muestra en \exref{ex:06-pro-b}, al igual que los sujetos expresos, \textit{pro}
se legitima en la posición de Esp de la F y se considera que sus rasgos
gramaticales de persona y número concuerdan con los de la desinencia verbal.
Una característica de \textit{pro} es que, debido a que carece de rasgos
fonológicos, no puede focalizarse (v.~§\ref{sec:06-capas}) o ponerse de relieve,
ya que este procedimiento requiere que el elemento destacado reciba acento
fuerte y, obviamente, un elemento que no se pronuncia no puede ser acentuado.
Así, mientras el sujeto léxico de \exref{ex:06-foco-a} puede destacarse mediante
la focalización, el sujeto tácito de \exref{ex:06-foco-b} no puede realzarse
entonativamente.

\begin{exe}
  \ex
    \begin{xlist}
      \ex\label{ex:06-foco-a} Ganó la beca ELLA, no su hermano.
      \ex[*]{\label{ex:06-foco-b} Ganó la beca ∅, no su hermano.}
    \end{xlist}
\end{exe}
\end{recuadro}

\subsubsection{La interpretación del sujeto tácito}

Según \textcite{bosque2009}, el contenido del sujeto tácito \textit{pro} se
recupera a partir del discurso precedente \exref{ex:06-andes1}, del
subsiguiente \exref{ex:06-andes2} o a partir del contexto extraoracional
\exref{ex:06-andes3}:

\begin{exe}
  \ex\label{ex:06-andes1} Los pueblos originarios de los Andes\etq{i} no
      conocían el caballo, aunque \textit{pro}\etq{i}\footnote{El subíndice que
      sigue a \textit{pro} indica cuál es el elemento del contexto al que
      \textit{pro} remite, el cual también va seguido de dicho índice; por eso
      se dice que ambos elementos están coindizados.} utilizaban llamas como
      animales de carga.
  \ex\label{ex:06-andes2} Aunque \textit{pro}\etq{i} no conocían el caballo, los
      pueblos originarios de los Andes\etq{i} utilizaban llamas como animales de
      carga.
  \ex\label{ex:06-andes3} \textit{pro}\etq{[1.ª pers. sg.]} Leí que los pueblos
      originarios utilizaban llamas como animales de carga.
\end{exe}

\subsection{Sujetos tácitos vs. sujetos explícitos}\label{sec:06-explicitos}

Para ver cuáles son las condiciones en que se usan los sujetos tácitos,
consideremos comparativamente las siguientes secuencias de oraciones:

\begin{exe}
  \ex\label{ex:06-nino} Descubrieron a un niño perdido en la montaña. Tenía las
      manos congeladas. Estaba muy asustado y no podía expresarse bien.
  \ex\label{ex:06-llegaron} Llegaron de noche. Estaban sumamente fatigados. Nadie
      los vio entrar. Traían las manos heladas y estaban muy desabrigados.
      Habían venido huyendo por leguas y, sin embargo, aún no tenían un destino
      seguro.
\end{exe}

Mientras que en \exref{ex:06-nino}, a partir de su mención en la primera
oración ---\forma{un niño perdido en la montaña}---, podemos fácilmente
recuperar el sujeto de la segunda y tercera oración, en \exref{ex:06-llegaron},
en cambio, aunque entendemos perfectamente lo que dice el texto, y también
podemos identificar que se está haciendo referencia a una tercera persona
plural, no podemos inferir de quiénes se está hablando en el fragmento. Tal como
plantea \textcite{camacho2018}, la interpretación de los sujetos nulos requiere
un \termino{tópico}\footnote{Aquí utilizamos \textit{tópico} o \textit{tema
discursivo} en el sentido general referido a aquello de lo que se está
hablando.} ya introducido en el discurso a partir del cual estos puedan ser
identificados. Además, los sujetos nulos también se emplean cuando el tópico
discursivo se mantiene; así, en \exref{ex:06-nino}, el tópico \forma{un niño
perdido en la montaña} permanece como sujeto en las dos oraciones siguientes.

Y ¿cuándo se utilizan los sujetos explícitos? Veamos lo que sucede en
\exref{ex:06-osos}.

\begin{exe}
  \ex\label{ex:06-osos} Una fotógrafa observó el comportamiento de los osos
      polares. Una cría de oso polar debe aprender las habilidades de
      supervivencia necesarias en el Ártico. Las madres se encargan de proteger
      a sus crías. Los machos, en cambio, no se hacen cargo del cuidado de los
      oseznos.
\end{exe}

El ejemplo de \exref{ex:06-osos} muestra que los sujetos explícitos se usan en
el comienzo de un texto o párrafo para identificar al referente del que estamos
hablando ---\forma{una fotógrafa}, en la primera oración--- y también se emplean
cuando se produce un cambio de tópico; así, del sujeto \forma{una fotógrafa} en
la primera oración se pasa a \forma{una cría de oso polar}, en la segunda; a
\forma{las madres}, en la tercera, y a \forma{los machos}, en la última.

\begin{ejercicio}\label{ej:06-peyrou}
Lea atentamente el siguiente fragmento adaptado de \textit{Se vuelven contra
nosotros}, de Manuel Peyrou, y luego realice las actividades pedidas, según se
muestra en los ejemplos resueltos.

\begin{consignas}
  \item Encorchete los sujetos de todas las oraciones y subraye los verbos a los
        que se vinculan.
  \item Inserte un \textit{pro} en el lugar correspondiente a los sujetos nulos y
        señale sus rasgos de concordancia (número y persona).
  \item Cuando sea posible, identifique el antecedente de cada sujeto nulo y
        proponga alguna hipótesis explicativa para aquellos casos en que no sea
        posible reconocer un antecedente en el texto.
\end{consignas}

% Fragmento literario protegido: se incorpora a mano desde el original.
\textoaincorporar{el fragmento adaptado de \textit{Se vuelven contra nosotros},
de M.~Peyrou, con las oraciones numeradas de (1) a (8): desde «Estos hombres
llevaban el nombre genérico de \textit{Proveedores}» hasta «\ldots{}y el hombre
dobló su fortuna en tres meses»}\pendiente{El original fecha el texto en 2020 y
escribe «Peyroú»}

\begin{resuelto}
a) (1) [\etq{Suj} Estos hombres]\etq{i} llevaban el nombre genérico de
«Proveedores». b) (2) A veces realmente [\etq{Suj}
\textit{pro}\etq{[3.ª pl.]}]\etq{i} vendían caballos y otros elementos
necesarios. c) El antecedente de \textit{pro} es \forma{estos hombres}, de la
oración anterior.
\end{resuelto}
\end{ejercicio}

\begin{recuadro}[Para ampliar: un caso de variación dialectal en el uso de los sujetos explícitos]
Aunque en el español general la reiteración del sujeto cuando no hay cambio de
tópico resulta redundante, en el español de la región de El Cibao, en República
Dominicana, los sujetos pronominales tienden a repetirse.

\begin{exe}
  \ex Cuando \textbf{tú} acabe \textbf{tú} me avisas.
      \parencite[ejemplo tomado de][393]{rodriguezramalle2005}
\end{exe}

Para Rodríguez Ramalle, esto sucede porque las marcas flexivas, que en el
español general permiten omitir el sujeto, se debilitan hasta no pronunciarse a
veces, lo que provoca la necesidad de reposición del sujeto.
\textcite{bullock2009} plantean que en esta variedad los pronombres explícitos de
sujeto son mucho más frecuentes que en el español general e incluso se usan,
contrariamente a lo esperado, para continuar el tópico discursivo, como muestra
el siguiente ejemplo tomado de los citados autores:

\begin{exe}
  \ex Entonces \textbf{él} fue a buscar el gallo, en busca el resto de las
      plumas, y \textbf{él} fue a buscarlo. Entonces, \textbf{él} vio un oso
      debajo de la mata de mango y vio el gallo a la parte de atrás del oso.
      \parencite[56; hablante dominicano de El Cibao]{bullock2009}
\end{exe}
\end{recuadro}

\begin{recuadro}[Diferencia entre oración y enunciado]
Según hemos planteado, la oración se define fundamentalmente como una unidad
propia de la gramática, formada por un SN sujeto y un SV predicado, articulados
por la F; se trata, pues, de una unidad estructural teórica, independientemente
de su empleo en situaciones concretas de habla. Para las emisiones efectivamente
proferidas en una situación concreta por un hablante determinado y dirigidas a
un oyente también determinado, se emplea la denominación de
\termino{enunciados}. Por tanto, mientras la oración es una unidad de la
gramática, el enunciado es una unidad del uso que no requiere tener la
estructura predicativa. Así, unidades como \forma{¡Socorro!, En la otra esquina,
Bien} o \forma{De ninguna manera} son enunciados perfectamente comprensibles y
que podríamos emplear en situaciones comunicativas concretas, aunque no pueda
reconocerse en ellos la estructura predicativa de sujeto y predicado.
\end{recuadro}

\section{Predicados no verbales y predicaciones sin sujeto}\label{sec:06-noverbales}

Para la \textit{Nueva gramática de la lengua española} \parencite[§\,1.13a]{rae2025},
las oraciones «son unidades mínimas de predicación organizadas en torno a un
verbo, que ponen en relación un sujeto con un predicado». El texto aclara que el
predicado puede ser verbal o no \parencite[§\,1.12d]{rae2025}. Al respecto
podemos tomar un caso como \forma{Interesante, la película}, en el que, aunque no
hay verbo, igualmente hay predicación. Se trata de una predicación de tipo
enfático, que se marca no solo por la ausencia del verbo cópula, que se podría
reponer ---\forma{Interesante es la película}---, sino por la alteración del
orden entre sujeto y predicado, que se vuelve significativa. En este caso
tenemos un predicado no verbal ---\forma{interesante}---, referido a un sujeto
pospuesto ---\forma{la película}---, junto al cual constituye una estructura
oracional de tipo enfático.

Ahora bien, en casos como \forma{Llueve} o \forma{Hace frío}, tenemos un verbo
conjugado, pero, dado que se trata de eventos climáticos, de naturaleza
impersonal, no hay un sujeto referencial con el que se los pueda relacionar. No
obstante, se trata de oraciones porque las formas verbales conjugadas
\forma{llueve} y \forma{hace (frío)} gramaticalmente manifiestan tercera persona
singular. Nuevamente, para el español, tendremos un sujeto \textit{pro}, ya que
en las lenguas \textit{pro-drop} los verbos de este tipo no pueden tener un
sujeto expreso. Y ¿qué sucede en estos casos con las lenguas que realizan la
primera opción del parámetro del sujeto nulo, como el inglés y el francés? Como
en estas lenguas la presencia de un sujeto es obligatoria, la posición de sujeto
es ocupada por un \termino{expletivo}, es decir, un elemento con rasgos
fonológicos, pero vacío de significado, como \forma{it} en inglés
\exref{ex:06-it} o \forma{il} en francés \exref{ex:06-il}.

\begin{exe}
  \ex\label{ex:06-it} \forma{It rains.} (\sig{Llueve})
  \ex\label{ex:06-il} \forma{Il pleut.} (\sig{Llueve})
\end{exe}

Las lenguas \textit{pro-drop}, en cambio, no necesitan un expletivo léxico, dado
que el pronombre sujeto puede ser tácito y estar representado por un
\textit{pro}\etq{expl}, tal como sucede en español \exref{ex:06-llueve} y en
italiano \exref{ex:06-piove}.

\begin{exe}
  \ex\label{ex:06-llueve} {}[\etq{SF/O} [\etq{Suj} \textit{pro}\etq{expl}]
      [\etq{SV/Pred} llueve\etq{V/n}]]
  \ex\label{ex:06-piove} {}[\etq{SF/O} [\etq{Suj} \textit{pro}\etq{expl}]
      [\etq{SV/Pred} piove\etq{V/n}]]
\end{exe}

En español también existe otro tipo de oraciones impersonales, con verbo en
tercera persona plural, como las siguientes:

\begin{exe}
  \ex\label{ex:06-arb}
    \begin{xlist}
      \ex\label{ex:06-arb-a} Te esperan en la puerta.
      \ex\label{ex:06-arb-b} Me llamaron por teléfono hoy temprano.
    \end{xlist}
\end{exe}

En oraciones como \exref{ex:06-arb-a} y \exref{ex:06-arb-b} podemos interpretar
que la referencia de la flexión remite a un sujeto plural mencionado
anteriormente y recuperable del contexto, tal como se planteó en
§\ref{sec:06-explicitos} ---por ejemplo, \forma{Vinieron tus amigos a buscarte.
Te esperan en la puerta}---. En este caso el sujeto plural corresponde a
entidades definidas e identificables: \forma{ellos / tus amigos}. Pero los
ejemplos de \exref{ex:06-arb} también admiten otra interpretación, muy común en
la lengua, en la que no hay una referencia a una entidad determinada y en la que
«el plural es una mera marca de indefinición, ya que puede referirse a un único
individuo» \parencite[109]{ditullio2014}. Este sujeto, que se interpreta como
indefinido e inespecífico, se conoce como \termino{pro arbitrario}:

\begin{exe}
  \ex
    \begin{xlist}
      \ex {}[\etq{SF/O} [\etq{Suj} \textit{pro}\etq{arb}] [\etq{SV/Pred} te
          esperan en la puerta]]
      \ex {}[\etq{SF/O} [\etq{Suj} \textit{pro}\etq{arb}] [\etq{SV/Pred} me
          llamaron por teléfono hoy temprano]]\pendiente{El original dice «Te
          llamaron»}
    \end{xlist}
\end{exe}

\begin{ejercicio}\label{ej:06-esoracion}
Decida cuáles de las siguientes son oraciones y fundamente su respuesta, según
se muestra en los ejemplos resueltos de abajo.

\begin{exe}
  \ex Mario es el barbero del barrio.
  \ex Un sinfín de problemas.
  \ex Cayó un meteorito en la autopista.
  \ex Vendrán más tarde.
  \ex Allá en el rancho grande.
  \ex Allá, solo ranchos y casuchas con techo de paja.
  \ex En el cuento había un ogro, un burro y una princesa.
  \ex Primero entraron los familiares cercanos y después, los demás parientes.
  \ex Granizó copiosamente en varias zonas del interior de la provincia.
  \ex Infaltable, el café de la sobremesa.
\end{exe}

\begin{resuelto}[Ejemplos resueltos]
1. O formada por el SN sujeto \forma{Mario} y el SV predicado \forma{es el
barbero del barrio}. 2. No es O, sino SN.
\end{resuelto}
\end{ejercicio}

\subsection{La relación predicativa}\label{sec:06-predicativa}

Ahora que ya hemos presentado la estructura oracional de acuerdo con el esquema
básico y universal de la teoría X\prima{} (v. \capref{5},
§\ref{sec:05-xbarra}), podemos también preguntarnos qué significa que el SN y el
SV se combinen para formar una oración. La relación que se establece entre los
dos constituyentes esenciales de la oración, el SN y el SV, se conoce como
\termino{predicación}. En ella el verbo asume un doble papel: por un lado, como
categoría léxica, se desempeña como núcleo del SV, soporte del predicado; por el
otro, en su desinencia, el verbo aloja a la flexión, categoría funcional que,
mediante la concordancia y la temporalización, instaura la relación
predicativa.

En tanto relación predicativa, toda oración manifiesta un evento o estado y, por
lo tanto, se organiza en torno al verbo que lo expresa. La mayor parte de los
verbos no se emplean solos, sino que tienen distintas exigencias de
complementación sintáctico-semánticas, que son satisfechas por los argumentos
que los acompañan. Como veremos más adelante, los argumentos del verbo remiten a
entidades del mundo extralingüístico que participan en el evento descripto por
el verbo. Por lo tanto, según su semántica, los verbos pueden requerir más o
menos argumentos. Así, por ejemplo, \forma{trabajar} solo exige un argumento,
referido a \sig{la persona que trabaja}\etq{x} \exref{ex:06-trabajar}; en
cambio, \forma{escribir} requiere dos argumentos: \sig{el que escribe}\etq{x} y
\sig{lo que se escribe}\etq{y} \exref{ex:06-escribir}; y hay también verbos que
piden tres argumentos, como, por ejemplo, \forma{enviar}, que debe manifestar
\sig{el que envía}\etq{x}, \sig{lo que se envía}\etq{y} y \sig{a quien se
envía}\etq{z} \exref{ex:06-enviar} (v. \capref{7}).

\begin{exe}
  \ex\label{ex:06-trabajar} \textit{trabajar}(x). Ej.: [Pedro]\etq{x} trabaja.
  \ex\label{ex:06-escribir} \textit{escribir}(x, y). Ej.: [Pedro]\etq{x}
      escribe [un correo electrónico]\etq{y}.
  \ex\label{ex:06-enviar} \textit{enviar}(x, y, z). Ej.: [Pedro]\etq{x} envía
      [un correo electrónico]\etq{y} [a sus alumnos]\etq{z}.
\end{exe}

\begin{ejercicio}\label{ej:06-argumentos}
Proporcione los ejemplos solicitados, según se muestra en el ejemplo de abajo.

\begin{exe}
  \ex Dos oraciones con verbos que exijan un solo argumento.
  \ex Dos oraciones con verbos que exijan dos argumentos.
  \ex Dos oraciones con verbos que exijan tres argumentos.
\end{exe}

\begin{resuelto}[Respuesta posible]
1. a) [Juan] trabaja en el bar de la otra cuadra. El V \forma{trabajar} exige un
solo argumento: \forma{Juan}. \forma{En el bar de la otra cuadra} es un
circunstancial no exigido por el V. b) [\etq{Suj} La pobre muchacha] bailó
frenéticamente hasta casi desfallecer. El V \forma{bailar} exige un solo
argumento. Tanto \forma{frenéticamente} como \forma{hasta casi desfallecer} son
circunstanciales no exigidos por el V.
\end{resuelto}
\end{ejercicio}

\begin{ejercicio}\label{ej:06-funciones}
Considere las siguientes oraciones y luego responda lo pedido, como se muestra
para el ejemplo resuelto~1.

\begin{exe}
  \ex La abogada me entregó los papeles en su oficina.
  \ex El primer ministro informó a la población sobre las nuevas medidas de
      gobierno.
  \ex El fiscal del caso guardó los expedientes en un sobre.
\end{exe}

\begin{consignas}
  \item ¿Cuántos argumentos exigen los verbos de las oraciones anteriores?
  \item ¿Qué funciones sintácticas cumplen los argumentos
        identificados?\footnote{Para las funciones sintácticas, puede consultar
        el \capref{7}.}
\end{consignas}

\begin{resuelto}
1. a) \forma{Entregó} exige tres argumentos: \forma{la abogada, los
certificados} y \forma{el nuevo dueño}. b) [\etq{Suj} La abogada] entregó
[\etq{COD} los certificados] [\etq{COI} al nuevo dueño].\pendiente{La resolución
no corresponde a la oración 1: «La abogada me entregó los papeles en su
oficina»}
\end{resuelto}
\end{ejercicio}

\begin{ejercicio}\label{ej:06-llover}
Observe la siguiente oración del español y responda lo pedido más abajo:

\begin{center}
\forma{Llovió torrencialmente en el fin de semana.}
\end{center}

\begin{consignas}
  \item ¿Cuántos argumentos exige \forma{llover} en la oración anterior?
  \item ¿Se le ocurren otros verbos que se construyan como \forma{llover}?
        Proporcione dos ejemplos oracionales.
  \item ¿Qué sucede en lenguas como el inglés o el francés con los verbos
        identificados en el punto anterior?
  \item ¿Qué nombre recibe el elemento que hace las veces de sujeto en estas
        lenguas?
  \item ¿Con qué parámetro lingüístico puede relacionar la diferencia existente
        en estas lenguas respecto de otras como el español y el italiano?
  \item ¿Qué nombre reciben estas dos últimas lenguas en función de esta
        característica?
\end{consignas}

\begin{resuelto}[Punto resuelto]
a) \forma{Llover} es un verbo impersonal que no exige ningún argumento, ni
siquiera sujeto. \forma{Torrencialmente} y \forma{en el fin de semana} son
circunstanciales no exigidos.
\end{resuelto}
\end{ejercicio}

\begin{ejercicio}\label{ej:06-cuatro}
Algunos autores, como \textcite[268]{bosque2009}, plantean la posibilidad de
verbos con hasta cuatro argumentos y ofrecen ejemplos como los siguientes:

\begin{exe}
  \ex Juan viaja desde Acassuso a Mar del Plata por la ruta dos.
  \ex María le compró su auto a su prima en 500\,000\,\$.
\end{exe}

\begin{consignas}
  \item Encorchete los argumentos de las oraciones anteriores, como se muestra
        en el ejemplo resuelto~1.
  \item Proporcione dos ejemplos más de verbos que admitan cuatro argumentos.
\end{consignas}

\begin{resuelto}
1. a) [Juan] viaja [desde Acassuso] [a Mar del Plata] [por la ruta dos].
\end{resuelto}
\end{ejercicio}

\subsection{Diferencias entre argumentos: argumento externo y argumentos internos}\label{sec:06-externo}

No todos los argumentos que exige el verbo están en el mismo nivel. Por
ejemplo, si bien un verbo transitivo como \forma{construir} exige dos argumentos
que dan cuenta, respectivamente, de \sig{el que construye} y de \sig{lo
construido} \exref{ex:06-arquitecto}, existe una diferencia entre ellos que lleva
a la distinción entre argumento externo y argumento interno
\parencite{williams1980}. El \termino{argumento externo}, que se identifica con
el sujeto por su relación estructural prominente, es el que mantiene con el
verbo la relación de predicación; en cambio, los \termino{argumentos internos} o
complementos verbales tienen una vinculación más estrecha con el verbo y, según
veremos más adelante, permanecen dentro del SV.

\begin{exe}
  \ex\label{ex:06-arquitecto} {}[\etq{Arg. ext.} El arquitecto] construyó
      [\etq{Arg. int.} un puente].
\end{exe}

\begin{recuadro}[Para ampliar: una relación muy estrecha]
Los argumentos internos, especialmente el COD, completan o complementan el
significado del V. Así, aunque se trate del mismo verbo, el significado de
\forma{servir} no es exactamente el mismo en cada una de las siguientes
oraciones, y esa variación se debe a la diferente naturaleza de los objetos que
lo acompañan en cada una de ellas:

\begin{exe}
  \ex El mozo \textit{sirvió} a los invitados (= \sig{los atendió}).
  \ex Ese hombre \textit{sirve} a Dios desde su juventud (= \sig{está consagrado
      a su servicio}).
  \ex Trajeron un toro campeón para \textit{servir} a las vacas (=
      \sig{fecundarlas, prestar servicios sexuales}).
\end{exe}

Por tanto, según estamos viendo, la vinculación entre el V y su objeto es
siempre más estrecha que la que existe entre el V y su sujeto. En español, hay
dos tipos de formaciones que muestran la integración entre el V y su OD, que los
lleva a constituir una unidad: los compuestos V + N y las locuciones verbales.
Dichas formaciones solo se dan entre el V y su OD y nunca entre el V y su
sujeto.

\begin{itemize}
  \item Los compuestos V + N del tipo de \forma{sacacorchos, abrelatas,
        sacapuntas, tirabombas, matamoscas, vendehumo, rompecorazones} (v.
        \capref{4}, §\ref{sec:04-composicion}).
  \item Diversas locuciones verbales, como \forma{hacer la cama} (=
        \sig{engañar}), \forma{dar fe} (= \sig{atestiguar}), \forma{echar
        raíces} (= \sig{radicarse}), \forma{escupir el asado} (=
        \sig{estropear}) (v. \capref{11}).
\end{itemize}
\end{recuadro}

\section{La organización tripartita de la oración: las capas oracionales}\label{sec:06-capas}

En tanto relación predicativa, toda oración manifiesta un evento o suceso y,
por lo tanto, se organiza en torno al verbo que lo expresa. En el denominado
\termino{enfoque cartográfico} \parencite{rizzi1995} se propone que la oración
se estructura en tres capas o niveles (\textit{shells}):

\begin{itemize}
  \item la \termino{capa léxica} o del SV, donde se produce la vinculación del V
        con sus argumentos (incluido el sujeto);
  \item la \termino{capa flexiva}, oracional o del SF, que pone en relación el
        sujeto con el predicado y en la que la oración se sitúa temporalmente;
  \item la \termino{capa informativa}, complementizadora o del SComp, que
        incluye la perspectiva del hablante y establece la articulación con el
        nivel del discurso.
\end{itemize}

Tal como sostiene \textcite{carnie2008}, se trata de

\begin{cita}
uno de los descubrimientos más interesantes de la lingüística del siglo veinte:
[ha sido] la naturaleza en tres capas de las oraciones. Con esto me refiero a la
forma en que un cierto número de enfoques sobre la estructura sintáctica han
convergido en la idea de que, translingüísticamente, las oraciones tienen en su
centro una predicación nuclear, que está dentro de una estructura de operadores
que ubican la predicación en tiempo y espacio respecto del tiempo de la emisión,
y finalmente están encerradas por una estructura que relaciona la proposición
con la perspectiva del hablante. (Traducción nuestra)
\end{cita}

Es importante destacar que esta organización abstracta se refleja en la forma
en que está estructurada la oración: desde el interior hacia el exterior, o sea,
de derecha a izquierda.

\begin{exe}
  \ex {}[\etq{SC} Quizás [\etq{SF} el profesor [\etq{SV} nos dé tarea para el fin
      de semana]]]
\end{exe}

Desde esta perspectiva, la capa léxica o del SV, que se organiza en torno al
verbo (aún no temporalizado), es la más interna. En ella la predicación se
despliega como un evento y el verbo, elemento central de la estructura
oracional, se relaciona con sus argumentos y con los modificadores no exigidos o
adjuntos.

\begin{exe}
  \ex\label{ex:06-castores} {}[\etq{SV} construir [\etq{Arg.~1} diques]
      [\etq{Arg.~2} los castores] [\etq{Adj} en ríos y arroyos]]
\end{exe}

Como se puede advertir en \exref{ex:06-castores}, \forma{construir} manifiesta
un evento de \sig{acción} (v. \capref{8}), que exige dos argumentos: \sig{lo
construido} ---\forma{diques}--- y \sig{quien construye} ---\forma{los
castores}---. La descripción del evento incluye un modificador no exigido por el
V, el adjunto de lugar ---\forma{en ríos y arroyos}---.

El siguiente nivel estructural se conoce como \termino{capa flexiva} o de la
oración porque, como ya adelantamos, mediante la flexión verbal ---que expresa T
y Conc--- el evento queda localizado respecto del momento de la emisión y se
establece la vinculación entre el SV y su sujeto. Es en este nivel en el que el
argumento sujeto debe desplazarse desde la posición interna que como argumento
del V ocupa dentro del SV, como se muestra en \exref{ex:06-viaje}, hacia el Esp
de la F, lo cual lo convierte en un «argumento externo», en tanto los demás
argumentos, que permanecen en el ámbito del SV, se identifican, según ya hemos
dicho, como «argumentos internos».

\begin{recuadro}[Para ampliar: el «viaje» del sujeto]
Llegados a este punto, ya podemos preguntarnos por qué se produce este
movimiento del sujeto hacia el exterior del SV, es decir, qué lo motiva. Además
de la temporalización y la concordancia, la conformación de la oración también
requiere que todos los argumentos del verbo «se vuelvan visibles» o manifiestos
para la sintaxis, para lo cual deben tener una marca que indique la función
sintáctica que cumplen. Ahora bien, ¿cuál es la categoría gramatical que marca
la función sintáctica de los SSNN? Según se vio en el \capref{4}
(v.~§\ref{sec:04-categorias}), el \termino{caso} es la categoría gramatical que
tiene esta función. La marca de caso «vuelve visibles» para su funcionamiento
oracional a los SSNN presentes en la oración. Recordemos que en las lenguas
clásicas ---griego y latín--- y también en algunas lenguas modernas, como el ruso
o el alemán, el caso es una categoría morfológica incluida en la desinencia de
la palabra: \forma{capra} (\sig{cabra}, caso nominativo o del sujeto);
\forma{capram} (\sig{a la cabra}, caso acusativo o del COD). En lenguas como el
español, en las que no existe el caso morfológico, se dice que esta marca es
abstracta, excepto en los pronombres personales, que varían en caso. Así, la
forma del pronombre de primera persona singular para la función de sujeto es
\forma{yo} ---\forma{\textbf{Yo} canto muy bien}---; para la función de objeto
es \forma{me} ---\forma{\textbf{Me} regalaron un hermoso cachorrito}---, y la
forma que se usa como complemento de preposición es \forma{mí} ---\forma{Lo
quiero para \textbf{mí}}---. En la sintaxis, el caso que le corresponde a un SN
determinado debe ser asignado. Así, por ejemplo, el V es la categoría léxica que
asigna caso acusativo al COD. El OI, en cambio, recibe su marca de caso de la
preposición \forma{a} que siempre lo precede: \forma{Les traje regalos
\textbf{a} los niños}. Y ¿qué sucede, entonces, con el sujeto? ¿Cómo recibe
caso? Para recibir caso nominativo, que es el que le corresponde, el sujeto,
originariamente un argumento interno a la capa léxica, se desplaza del SV al Esp
de la F. Por tanto, el sujeto recibe caso nominativo de una categoría funcional:
la F. Asimismo, el V también se desplaza del SV al núcleo de la F para cotejar
sus rasgos flexivos de tiempo y concordancia con el sujeto.

\begin{exe}
  \ex\label{ex:06-viaje}
    \begin{xlist}
      \ex {}[\etq{SF} [\etq{SN1} los castores]\etq{i} [\etq{F\prima} F
          [\etq{SV} [\etq{SN1} \sout{los castores}]\etq{i} [\etq{V\prima}
          construyen\etq{V} [\etq{SN2} diques]]]]]
      \ex
          \begin{forest} gramatica
            [SF
              [SN\etq{1}\\\textit{los castores}\etq{i}, name=suj]
              [F\prima
                [F]
                [SV
                  [SN\etq{1}\\{\color{gristexto}\textit{los castores}}\etq{i}, name=huella]
                  [V\prima
                    [V\\\textit{construyen}]
                    [SN\etq{2}\\\textit{diques}]]]]]
            \draw[->, densely dashed] (huella.south west) to[out=200, in=250]
              (suj.south);
          \end{forest}
    \end{xlist}
\end{exe}
\end{recuadro}

Dado que la oración ya queda conformada mediante la vinculación del SN sujeto y
del SV predicado dentro de la capa flexiva, resulta lícito preguntarnos para qué
se necesita proponer un nivel más o, mejor dicho, ¿qué aporta el tercer nivel
oracional? Este tercer nivel, al que se lo conoce como \termino{capa
informativa} o del \termino{complementante} (C), es el lugar en el que la
oración se conecta con la perspectiva del hablante, por lo que se ubican allí
distintos elementos que expresan la fuerza ilocutiva de la oración, tales como
los adverbios modales, el foco y el tópico, entre otros. Vayamos viendo de qué
se trata esto.

El C es, al igual que la F, un núcleo funcional y, así como la F contiene
valores de T y Conc, el C es de naturaleza modal. Si bien no vamos a
extendernos aquí sobre el complementante (v. capítulos 12 y 13), podemos
ilustrar su carácter modal mostrando la diferencia entre \exref{ex:06-que} y
\exref{ex:06-si}.

\begin{exe}
  \ex
    \begin{xlist}
      \ex\label{ex:06-que} Dice [\etq{SC} \textbf{que} [\etq{SF/O} Juan [\etq{SV}
          viene temprano]]].
      \ex\label{ex:06-si} Pregunta [\etq{SC} \textbf{si} [\etq{SF/O} Juan
          [\etq{SV} viene temprano]]].
    \end{xlist}
\end{exe}

En ambas oraciones, el complemento del verbo matriz o principal ---\forma{dice,
pregunta}--- no es un SN, como en \forma{Dice [la verdad]}, sino una oración
(SF) ---[\forma{Juan viene temprano}]---, por lo que la vinculación entre la
oración principal y la subordinada no es directa, sino mediada por una
conjunción, que puede ser \forma{que} o \forma{si}.

Según vemos, la distribución de ambas conjunciones ---\forma{que} y \forma{si}---
no es aleatoria. Pero ¿a qué se debe la alternancia? Sin duda no depende de la
oración subordinada, ya que en ambos casos ---\exref{ex:06-que} y
\exref{ex:06-si}--- es la misma ---\forma{Juan viene temprano}---. La diferencia
radica, entonces, en el verbo matriz: \forma{decir} en \exref{ex:06-que} y
\forma{preguntar} en \exref{ex:06-si}. Por tanto, es la modalidad de la oración
matriz, declarativa en \exref{ex:06-que} e interrogativa en \exref{ex:06-si}, la
que provoca la diferencia.

\begin{destacado}
\textbf{Una distribución complementaria}
\begin{itemize}
  \item Cuando el verbo de la oración principal es declarativo ---por ejemplo,
        \forma{decir, informar, afirmar}---, la oración seleccionada es
        introducida por el complementante \forma{que} \exref{ex:06-que}.
  \item Cuando el verbo de la oración principal es interrogativo o dubitativo
        ---por ejemplo, \forma{preguntar, ignorar, querer saber}---, la oración
        seleccionada es introducida por el complementante \forma{si}
        \exref{ex:06-si}.
\end{itemize}
\end{destacado}

El nivel del complementante siempre está presente en la oración, pero, a veces,
puede estar vacío, es decir, puede no haber ninguna marca explícita de la
modalidad en la oración, como normalmente sucede en las oraciones declarativas.

\begin{exe}
  \ex Llegó después de las ocho.
  \ex Ese día Juan estaba feliz.
  \ex\label{ex:06-periodico} Traía un periódico debajo del brazo.
\end{exe}

Otras modalidades oracionales, en cambio, llevan marcadores explícitos. Así, en
\exref{ex:06-talvez}, el adverbio \forma{tal vez} y el modo subjuntivo en el
verbo indican que se trata de una oración dubitativa; en \exref{ex:06-ojala}, el
adverbio \forma{ojalá} y el subjuntivo en el verbo son marcas de la modalidad
desiderativa de la oración; en \exref{ex:06-quetraia}, el índice de la modalidad
interrogativa de la oración es el pronombre interrogativo \forma{qué}, que
representa al constituyente ---\forma{un periódico} de
\exref{ex:06-periodico}--- desplazado en \exref{ex:06-quetraia} a la posición
inicial.

\begin{exe}
  \ex\label{ex:06-talvez} Tal vez llegue después de las ocho.
  \ex\label{ex:06-ojala} Ojalá esté feliz.
  \ex\label{ex:06-quetraia} ¿Qué traía debajo del brazo?
\end{exe}

\begin{recuadro}[Para ampliar: orden en el sintagma de complementante (SC)]
Como se advierte en los ejemplos de arriba, todos los marcadores léxicos
---\forma{tal vez, ojalá, qué}--- se ubican en el margen izquierdo de la
oración, que corresponde al SC. El núcleo del SC contiene la información sobre
la modalidad oracional y allí se coloca la conjunción ---\forma{que, si}---
cuando la oración contiene una subordinada (v. \exref{ex:06-que} y
\exref{ex:06-si}).

En el SC también se alojan otros elementos que sirven para enlazar la oración
con el discurso previo. En este nivel la organización de la oración se vincula
con la \termino{estructura informativa}, que establece la distinción entre tema
y rema, la cual no coincide exactamente con la partición sintáctica en sujeto y
predicado. El \termino{tema} da cuenta de la información ya conocida, previsible
o que el oyente puede inferir; el \termino{rema}, en cambio, aporta la
información nueva.

Si bien no todas las posturas teóricas coinciden, desde la perspectiva que
estamos considerando, un tema destacado en la posición inicial de la oración es
un \termino{tópico}. Como ya adelantamos, los tópicos manifiestan aquello de lo
que se habla en la oración. En \exref{ex:06-top-a}, por ejemplo, el elemento
topicalizado es el COD ---\forma{a Juan}---, que luego es retomado por el
clítico \forma{lo} dentro de la oración. En \exref{ex:06-top-b} y
\exref{ex:06-top-c}, los tópicos aparecen introducidos por marcadores de tópico o
topicalizadores ---\forma{en cuanto a, hablando de}---. Además, como se muestra
en \exref{ex:06-top-c}, el tópico es una categoría recursiva y puede haber más
de uno en la oración.

\begin{exe}
  \ex
    \begin{xlist}
      \ex\label{ex:06-top-a} {}[\etq{SC} [\etq{Top} A Juan\etq{i}] C [\etq{SF/Comp}
          F [\etq{SV/Pred} lo\etq{i} vi esta mañana en la estación]]]
      \ex\label{ex:06-top-b} {}[\etq{SC} [\etq{Top} En cuanto a Juan\etq{i}] C
          [\etq{SF/Comp} F [\etq{SV/Pred} lo\etq{i} vi esta mañana en la
          estación]]]
      \ex\label{ex:06-top-c} {}[\etq{SC} [\etq{Top1} Hablando de Juan\etq{i}],
          [\etq{Top2} a su hermano\etq{j}] C [\etq{SF/Comp} F [\etq{SV/Pred}
          lo\etq{j} vi esta mañana en la estación]]]
    \end{xlist}
\end{exe}

El \termino{foco}, otro de los elementos vinculados con la estructura
informativa incluidos en el SC, representa la información no presupuesta, es
decir, aquella que no forma parte del conocimiento compartido por el hablante y
el oyente. El foco se antepone estableciendo un contraste con las expectativas
ofrecidas por el contexto anterior, por lo que tiene que ser destacado por medio
de recursos enfáticos, como la entonación, \exref{ex:06-foc1} y
\exref{ex:06-foc2}, el orden de palabras \exref{ex:06-foc2}, estructuras de
relieve \exref{ex:06-foc3} o elementos léxicos de focalización, como
\forma{incluso} en \exref{ex:06-foc4}.

\begin{exe}
  \ex\label{ex:06-foc1} {}[\etq{SC} [\etq{Foc} EN LA ESQUINA]\footnote{Convencionalmente,
      la entonación enfática se indica en la escritura por el uso de
      mayúsculas.} fue el accidente] (relieve entonacional)
  \ex\label{ex:06-foc2} {}[\etq{SC} [\etq{Foc} ¿QUÉ] dijiste?] (relieve
      entonacional y orden marcado\footnote{La posición canónica es la esperada
      para un constituyente dentro de la oración, como, por ejemplo, la
      posición inicial para el sujeto o la posverbal para el COD. Cuando los
      elementos se encuentran en la posición canónica, la oración tiene un
      orden no marcado; en cambio, se considera orden marcado cuando se altera
      la disposición de alguno de los constituyentes, como sucede en las
      estructuras tematizadas y focalizadas.} OD -- V)
  \ex\label{ex:06-foc3} En la esquina fue donde ocurrió el accidente
      (estructura de relieve con el verbo \forma{ser})
  \ex\label{ex:06-foc4} {}[\etq{SC} [\etq{Foc} Incluso Juan] estaba enterado del
      asunto] (relieve léxico del foco \forma{Juan} mediante el focalizador
      \forma{incluso})
\end{exe}

También tienen su posición canónica en el margen izquierdo:

\begin{itemize}
  \item Los enlaces con el discurso precedente, como las dos instancias de
        \forma{pero}, que conectan con párrafos anteriores, o el uso inicial de
        \forma{y}, que enlaza con la oración previa, en \exref{ex:06-asimov}:
\end{itemize}

\begin{exe}
  \ex\label{ex:06-asimov}
      % Fragmento literario protegido: se incorpora a mano desde el original.
      \textoaincorporar{el diálogo de I.~Asimov, \textit{Orificios en la
      cabeza} (p.~46), con las dos instancias de \textit{pero} iniciales de
      párrafo y la \textit{y} que abre oración, en cursiva}
\end{exe}

\begin{itemize}
  \item Los \termino{vocativos}, que constituyen una apelación al oyente por
        parte del hablante.
\end{itemize}

\begin{exe}
  \ex \textit{Che}, pero yo realmente podría colaborar en \textit{La Nación} de
      los domingos. (J.~Cortázar, \textit{Rayuela}, p.~222)\pendiente{El
      original dice «Cortázar, R.»}
\end{exe}
\end{recuadro}

\begin{recuadro}[Reflexión final]
Hasta aquí hemos dado un panorama global de la organización oracional
basándonos en los planteos de \textcite{rizzi1995} sobre «la naturaleza en tres
capas o niveles de la oración», a las que hemos vinculado, respectivamente, con
el SV, el SF y el SC. Ahora, en los capítulos que siguen, vamos a adentrarnos en
la predicación, en relación con la cual hay que atender a tres aspectos
interrelacionados, que dan origen a tres estructuras conectadas entre sí:

\begin{itemize}
  \item la \termino{estructura temático-argumental}, que depende de la semántica
        del verbo ---ya sea de acción, de proceso, de estado, etc.--- y, como ya
        hemos anticipado, exige determinada cantidad de argumentos, los cuales,
        según se verá, desempeñan diferentes papeles temáticos (θ) en el evento,
        tales como agente, paciente, experimentante, etc.;
  \item la \termino{estructura sintáctica}, que constituye el enlace entre la
        semántica y la sintaxis, y manifiesta las diferentes funciones
        oracionales que cumplen los papeles temáticos, tales como sujeto, COD,
        COI, etc.;
  \item la \termino{estructura eventiva}, que depende de la naturaleza del
        evento o situación descripta por el verbo ---durativa, puntual, con o
        sin límite final, etc.---, se vincula con la categoría de aspecto, sobre
        todo el léxico (v. capítulos 4 y 8), y está determinada por las distintas
        categorías y constituyentes presentes en la predicación.
\end{itemize}

Asimismo, aunque por el carácter introductorio de este texto no sea posible
explayarnos, cuando la oración se vincula con una situación comunicativa
determinada en la que un emisor se dirige a un receptor, lo que la convierte en
una unidad de realización o enunciado, la manera en que es presentada la
información ---como nueva o conocida, como presupuesta o focalizada---
constituye lo que se conoce como \termino{estructura informativa}, que se
relaciona con los elementos que se ubican en el margen izquierdo de la oración,
correspondiente al nivel del SC.
\end{recuadro}

% --- Secciones finales -------------------------------------------------
\lecturasrecomendadas

\solucionario

\begin{solucion}{\ref{ej:06-determinante}}
Los NN comunes, como \forma{ratas} (1a) y \forma{niños} (1b), por su propia
naturaleza incluyen individuos en una clase ---\forma{Juan es doctor}---. Para
ser usados en referencia a alguna entidad determinada de la realidad
extralingüística deben construirse con un D que los delimite ---\forma{Juan es
el doctor que atiende a mi padre}---, o bien con un Q ---\forma{Llegó algún
amigo de tu padre}---. En cambio, los NN propios, como \forma{Pedro} (2a) o
\forma{Buenos Aires} (2b), tienen la referencia incluida; por eso se pueden usar
sin ningún especificador.
\end{solucion}

\begin{solucion}{\ref{ej:06-oracion}}
b) 2 es una O porque contiene un sujeto y un predicado vinculados por un núcleo
funcional: la F, que manifiesta T (pasado) y Conc (3.ª pers.~sg.). La
particularidad de 2 respecto de 1 es que no contiene un sujeto representado
léxicamente, sino que constituye un caso de lo que se conoce como sujeto tácito.
\end{solucion}

\begin{solucion}{\ref{ej:06-peyrou}}
(3) [\etq{Suj} Su abuelo] contó el caso de uno de ellos\etq{j}, que en una
ocasión [\textit{pro}\etq{[3.ª sg.]}]\etq{j} supo con anticipación la firma de
[\etq{Suj} un decreto] ordenando el pago de cinco años de sueldos atrasados. El
antecedente de \textit{pro} es \forma{uno de ellos}. (4) Inmediatamente
[\textit{pro}\etq{[3.ª sg.]}]\etq{j} visitó durante tres meses, en varias
provincias, las distintas guarniciones. El antecedente de \textit{pro} es
\forma{uno de ellos}, de la oración anterior. (5) A todos los soldados, cabos y
sargentos les [\textit{pro}\etq{[3.ª sg.]}]\etq{j} hizo la misma proposición:
(6) ¿A vos cuánto te [\textit{pro}\etq{[3.ª pl.]}] deben? El antecedente del
primer \textit{pro} es \forma{uno de ellos}; el del segundo \textit{pro} queda
indefinido. (7) Recibida [\etq{Suj} la respuesta],
[\textit{pro}\etq{[3.ª sg.]}]\etq{j} concluía: «[\textit{pro}\etq{[1.ª sg.]}] Te
doy la mitad si me [\textit{pro}\etq{[2.ª sg.]}] firmás recibo por todo». (8)
[\etq{Suj} Los hombres, a veces hambrientos y semidesnudos, y siempre
necesitados] aceptaban, y [\etq{Suj} el hombre] dobló su fortuna en tres meses.
El antecedente del primer \textit{pro} de (7) es \forma{uno de ellos}; el del
segundo \textit{pro} es el emisor; el del tercer \textit{pro}, uno de los
receptores.\pendiente{La numeración de las oraciones del solucionario no
coincide con la del texto; el original usa el índice \textit{i} para
\textit{uno de ellos}}
\end{solucion}

\begin{solucion}{\ref{ej:06-esoracion}}
3. O formada por el SN sujeto \forma{un meteorito} y el SV predicado
\forma{cayó en la autopista}. 4. O formada por un sujeto tácito y el SV
predicado \forma{vendrán más tarde}. 5. SAdv. 6. O formada por el SN sujeto
\forma{solo ranchos y casuchas con techo de paja} y el SAdv predicado
\forma{allá}. 7. O formada por un SV de V impersonal: \forma{En el cuento había
un ogro, un burro y una princesa}. 8. Hay dos oraciones coordinadas: la primera
está formada por el SN sujeto \forma{los familiares cercanos} y el SV predicado
\forma{primero entraron}; la segunda está formada por el SN sujeto \forma{los
demás parientes} y un SV con el V elidido, que puede reponerse porque es el mismo
que el de la primera: \forma{después (entraron)}. 9. O formada por un SV de V
impersonal de tipo meteorológico: \forma{Granizó copiosamente en varias zonas
del interior de la provincia}. 10. O formada por el SN sujeto \forma{el café de
la sobremesa} y el SA predicado \forma{infaltable}.

\textit{Aclaración}: los ejercicios que siguen son preparatorios para los del
capítulo siguiente, que ya especifican las diferencias entre argumentos. Aquí
solo se solicita identificar los argumentos y las funciones sintácticas que
cumplen.
\end{solucion}

\begin{solucion}{\ref{ej:06-argumentos}}
Respuestas posibles. 2. a) [\etq{Suj} La madre de Lucía] prepara [\etq{COD}
unas exquisitas tortas de chocolate]. b) [\etq{Suj} Algunos aspirantes]
aprobaron [\etq{COD} el examen]. 3. a) [\etq{Suj} El gerente de marketing]
presentó [\etq{COD} un informe] [\etq{COI} al director de la compañía]. b)
[\etq{Suj} Los organizadores del evento] enviaron [\etq{COD} invitaciones]
[\etq{COI} a los conferencistas].
\end{solucion}

\begin{solucion}{\ref{ej:06-funciones}}
a) 2. \forma{Informó} exige tres argumentos, que en este caso son \forma{el
primer ministro, la población} y \forma{las nuevas medidas de gobierno}. 3.
\forma{Guardar} exige tres argumentos: \forma{el fiscal del caso, los
expedientes} y \forma{en un sobre}. b) 2. [\etq{Suj} El primer ministro] informó
[\etq{COI} a la población] [\etq{COD} las nuevas medidas de gobierno]. 3.
[\etq{Suj} El fiscal del caso] guardó [\etq{COD} los expedientes] [\etq{CCL} en
un sobre]. Aclaración: \forma{en un sobre} es un argumento exigido por la
semántica del V, que sin su presencia quedaría incompleto
(\juicio{*}{\forma{El fiscal del caso puso los informes}}); por lo tanto,
se considera un argumento del verbo con la función sintáctica de complemento
circunstancial de lugar.\pendiente{En la prueba, el original pasa a «puso los informes»; en 2 el análisis omite \textit{sobre} ante \textit{las nuevas
medidas}}
\end{solucion}

\begin{solucion}{\ref{ej:06-llover}}
b) Respuestas posibles: \forma{Granizó muy fuertemente esta mañana. Tronaba
tanto que me asusté}. c) Dado que en inglés o francés no es posible omitir el
sujeto, los VV meteorológicos se construyen con un sujeto que se pronuncia, pero
no tiene significado. Ejemplos: inglés \forma{It snows} (\sig{nieva}); francés
\forma{Il neige}. d) Este elemento que «hace las veces de sujeto» se conoce como
\termino{expletivo}. e) La exigencia de un expletivo se relaciona con el
parámetro del sujeto, que en estas lenguas obligatoriamente debe estar expreso
en la oración, a diferencia del español o el italiano, que permiten la omisión.
Ejemplos: español \forma{Nieva}; italiano \forma{Nevica}. f) Las lenguas como el
español o el italiano, que permiten la omisión del sujeto, se conocen como
\textit{pro-drop}.
\end{solucion}

\begin{solucion}{\ref{ej:06-cuatro}}
a) 2. [María] le compró [su auto] [a su prima] [en 500\,000\,\$].
\textit{Aclaración}: \forma{le} y \forma{a su prima} son correferenciales y
expresan el mismo argumento. b) Respuestas posibles: [El profesor de Química]
viene todos los días [a la Facultad] [desde Tigre] [por la Panamericana]. / [Los
exportadores de soja] vendieron [el cereal] [al exterior] [a precio récord] este
año.
\end{solucion}

% =====================================================================
%  capitulos/cap07.tex — Capítulo 7. De la estructura argumental a las
%  funciones sintácticas
%  Autora: Laura Ferrari
%  Fuente: Capitulo_07_normalizado.docx
% =====================================================================
\capitulo[De la estructura argumental a las funciones sintácticas]%
  {De la estructura argumental a las funciones sintácticas}{Laura Ferrari}

\section{Introducción}

¿Cuáles son las relaciones entre la sintaxis y la semántica? ¿Qué incidencia
posee el léxico en la configuración sintáctica de la oración? ¿Qué tipo de
información está contenida en el nivel léxico? En este capítulo nos proponemos
dar respuesta a estas preguntas. Para ello, partimos del estudio de las
funciones semánticas para ir adentrándonos en las cuestiones sintácticas. En
primer lugar, mostramos cómo el contenido semántico de las palabras que
constituyen una oración determina mayormente sus propiedades combinatorias. Esto
significa que del contenido semántico depende el número de argumentos
seleccionados por un núcleo verbal dentro de su proyección sintagmática. Así, un
verbo como \forma{capturar} necesita de dos SSNN para la buena formación de la
oración que constituye, como en \forma{El policía capturó al ladrón}; en cambio,
\forma{caminar} solo selecciona un SN, como en \forma{Mario camina}. En segunda
instancia, presentamos la \termino{estructura temática}, es decir, las funciones
semánticas de los argumentos que el predicado selecciona, y damos cuenta de la
lista de papeles temáticos. En función de la distinción entre complementos y
adjuntos en el sintagma verbal, desarrollamos las nociones de transitividad
directa e indirecta y distinguimos dos tipos de verbos intransitivos:
inergativos e inacusativos.

\section{Relaciones entre la sintaxis y la semántica}\label{sec:07-sintsem}

\subsection{Predicados y argumentos}\label{sec:07-predicados}

La noción de predicación es semántica y tiene sus orígenes en el análisis
lógico-filosófico. Es un concepto que hace referencia a una condición relativa
al significado de las expresiones. Es decir, la estructura predicativa de las
oraciones es una parte sustancial de su significado. La relación de predicación
se establece entre un predicado y un argumento, como ya se ha señalado en el
capítulo anterior (§\ref{sec:06-predicativa}) respecto del vínculo que se
establece entre el SN sujeto y el SV que predica de él.

\begin{destacado}
Un \termino{predicado} es todo aquel elemento léxico que denota una acción,
proceso o estado y que selecciona un conjunto o red de argumentos en función de
su significado léxico \parencite{rodriguezramalle2005}.
\end{destacado}

Un \termino{argumento} se puede definir como el participante de una acción,
proceso o estado imprescindible para que esta pueda llevarse a cabo.

\begin{destacado}
Los argumentos son «expresiones referenciales que permiten identificar entidades
del mundo extralingüístico» \parencite[102]{ditullio2014}.
\end{destacado}

El ejemplo típico de predicado es el verbo, mientras que los argumentos suelen
ser SSNN.

\begin{exe}
  \ex\label{ex:07-visito} {}[\etq{Arg.~1} El mandatario de Canadá] [\etq{V}
      visitó] [\etq{Arg.~2} el Jardín Japonés].
\end{exe}

\forma{Visitar} en \exref{ex:07-visito} funciona como un predicado que debe
construirse con dos SSNN, que constituyen sus argumentos. El evento de
\forma{visitar} relaciona dos argumentos: x (= \forma{el mandatario de Canadá})
e y (= \forma{el Jardín Japonés}). Para que una oración esté completa, el
elemento seleccionador, en este caso el verbo, debe haberse combinado con todos
los argumentos que selecciona.

\begin{ejercicio}\label{ej:07-agramatical}
Explique la agramaticalidad o el carácter ambiguo de las siguientes oraciones,
según se muestra en los ejemplos resueltos.

\begin{exe}
  \ex María prestó a su hijo.
  \ex La comisión respaldó durante el receso parlamentario.
  \ex José colocó.
  \ex El viento provocó.
  \ex El vocero presidencial anunció gelatina.
  \ex El presidente reclamó a los bancos.
  \ex Los cóndores habitan.
  \ex Lucía va todas las semanas.
\end{exe}

\begin{resuelto}[Ejemplos resueltos]
1. Es ambigua: a) se puede interpretar que lo prestado por \forma{María} es
\forma{su hijo} y queda implícito \sig{a quién}; b) es agramatical porque falta
el argumento referido a \sig{lo que María le prestó a su hijo}. 2. Es
agramatical porque falta el argumento referido a \sig{lo respaldado por la
comisión}.
\end{resuelto}
\end{ejercicio}

\begin{ejercicio}\label{ej:07-seleccion}
Algunas de las oraciones que siguen, por diversas razones, no respetan los
requisitos de selección de los verbos que las conforman. Haga una lectura atenta
y responda lo pedido en las consignas, según se muestra en los ejemplos
resueltos de abajo.

\begin{consignas}
  \item Explique el porqué de la buena o mala formación de cada una de las
        oraciones.
  \item Luego especifique la estructura argumental de \forma{regresar} y
        \forma{llegar}.
\end{consignas}

\begin{exe}
  \ex El viajero regresó contento a su casa y trajo muchos regalos a sus
      familiares.
  \ex El viajero llegó contento a su casa y trajo muchos regalos.
  \ex Llegó contento a su casa y trajo muchos regalos.
  \ex El viajero llegó oblongo a su casa y trajo muchos regalos a sus familiares.
  \ex El viajero llegó contento su casa y trajo muchos regalos a sus familiares.
  \ex El viajero llegó contento a su aceituna y trajo muchos regalos a sus
      familiares.
\end{exe}

\begin{resuelto}[Ejemplos resueltos]
a) 1 es gramatical porque ambos verbos tienen todos sus argumentos. 4 es
agramatical porque \forma{oblongo} no es compatible desde el punto de vista
léxico-semántico con el sujeto ni con el verbo \forma{llegar}. b)
\forma{Llegar} [Arg.~1 SN, Arg.~2 SP].
\end{resuelto}
\end{ejercicio}

\subsubsection{La valencia verbal}\label{sec:07-valencia}

El número de argumentos que necesita un verbo para que esté completo o saturado
se denomina \termino{valencia verbal}. Pero ¿cuál es el número de argumentos de
un predicado? El número y el tipo de argumentos de un predicado dependen de cómo
se conceptualiza lingüísticamente el evento. Y ¿cuántos argumentos puede exigir
un verbo? Para \textcite{hale1998}:

\begin{cita}
Los verbos de una lengua natural, generalmente la más rica categoría al
respecto, son extremadamente limitados en la variedad y complejidad de las
estructuras argumentales que despliegan, y estas se conforman a una tipología
altamente restricta. \textbf{Pocos verbos tienen más de tres argumentos}\ldots{}
(Traducción nuestra)
\end{cita}

Según el número de argumentos que exigen, los verbos se clasifican en:

\begin{itemize}
  \item \termino{cero-ádicos} o \termino{avalentes}: \forma{Continúa granizando
        (en el sur)};\footnote{Los paréntesis indican los adjuntos.}
        \forma{Llovió (torrencialmente) (toda la noche)};
  \item \termino{monádicos} o \termino{monovalentes}: \forma{[El ministro de
        economía] viajó (ayer)}; \forma{Hay [café]};
  \item \termino{diádicos} o \termino{bivalentes}: \forma{[El juez] redactó [la
        sentencia]}; \forma{[Mi vecino] había perdido [las llaves]};
  \item \termino{triádicos} o \termino{trivalentes}: \forma{[El fiscal] entregó
        [sus informes] [a la justicia]}; \forma{[El periodista] solicitó
        [información] [a la embajada]}.
\end{itemize}

Los predicados seleccionan sus argumentos en función de su significado. El
conjunto de argumentos de un predicado representa un esqueleto de su propia
definición y constituye su \termino{estructura argumental}.

\begin{ejercicio}\label{ej:07-valencia}
Lea atentamente el siguiente fragmento de «Emma Zunz», de J.~L.~Borges, y luego
resuelva lo pedido en las consignas, según se muestra en el ejemplo resuelto de
abajo.

% Fragmento literario protegido: se incorpora a mano desde el original.
\textoaincorporar{el primer párrafo de «Emma Zunz», de J.~L.~Borges (\textit{El
Aleph}), desde «El catorce de enero de 1922» hasta «\ldots{}que se dirigía a la
hija del muerto»}

\begin{consignas}
  \item Según la cantidad de argumentos que selecciona, determine a qué clase
        pertenece cada verbo (monádico, triádico, etc.).
  \item Luego identifique los argumentos seleccionados por cada verbo.
\end{consignas}

\begin{resuelto}
\forma{Hallar}: a) es un V triádico porque selecciona tres argumentos; b) los
argumentos seleccionados son \forma{Emma Zunz} (sujeto), \forma{una carta,
fechada en el Brasil} (COD) y \forma{en el fondo del zaguán} (CCL).
\end{resuelto}
\end{ejercicio}

\subsection{Complementos y adjuntos}\label{sec:07-compadj}

Como ya se introdujo en el \capref{6}, §\ref{sec:06-externo}, existe una
diferencia entre el \termino{argumento externo} o sujeto, que por su relación
estructural prominente es el que mantiene con el verbo la relación de
predicación, y los argumentos internos o complementos verbales, que permanecen
dentro del SV \parencite{williams1980}.\pendiente{El original dice «Williams
1981»}

Los argumentos internos al SV seleccionados por el predicado se denominan
\termino{complementos}. Los complementos y el argumento externo (sujeto)
introducen información exigida por el significado de los predicados, como
\forma{el señor Maier} (argumento externo / sujeto) y \forma{una fuerte dosis de
veronal} (argumento interno / COD) en \exref{ex:07-maier-a}. Por otra parte, los
modificadores no seleccionados por el predicado se denominan
\termino{adjuntos}, como los que aparecen entre paréntesis en
\exref{ex:07-maier-a}. Generalmente los adjuntos pueden omitirse y la oración
sigue siendo gramatical \exref{ex:07-maier-b}. Los adjuntos también pueden
agregarse en una oración con gran libertad \exref{ex:07-maier-c}; en cambio,
cada complemento puede ser realizado exclusivamente por un único argumento. En
otras palabras, un verbo puede tener varios complementos, siempre y cuando
realicen diferentes funciones sintácticas. Como se puede apreciar,
\exref{ex:07-maier-d} es agramatical porque hay dos SSNN no coordinados que
realizan la función de OD. No obstante, algunos argumentos pueden quedar
implícitos, esto quiere decir que forman parte de la estructura argumental del
verbo, pero no se manifiestan necesariamente en la sintaxis
\exref{ex:07-fruta}.

\begin{exe}
  \ex\label{ex:07-maier}
    \begin{xlist}
      \ex\label{ex:07-maier-a} {}[El señor Maier] había ingerido (por error) [una
          fuerte dosis de veronal] y había fallecido (el tres del corriente) (en
          el hospital de Bagé).
      \ex\label{ex:07-maier-b} {}[El señor Maier] había ingerido [una fuerte
          dosis de veronal].
      \ex\label{ex:07-maier-c} {}[El señor Maier] había ingerido (por error) [una
          fuerte dosis de veronal] (con gran repugnancia) (esa mañana).
      \ex[*]{\label{ex:07-maier-d} [El señor Maier] había ingerido [una fuerte
          dosis de veronal] [veneno para ratas].}
    \end{xlist}
  \ex\label{ex:07-fruta} ---¿Comés fruta?\\
      ---No, no como. (El COD de \forma{comer} está implícito).
\end{exe}

\begin{ejercicio}\label{ej:07-compadj}
Lea atentamente el siguiente fragmento de «Emma Zunz» y resuelva lo pedido en
las consignas, según se muestra en el modelo de abajo.

% Fragmento literario protegido: se incorpora a mano desde el original.
\textoaincorporar{el fragmento de «Emma Zunz» que va de «En abril cumpliría
diecinueve años» hasta «\ldots{}se acostó y se obligó a dormir»}

\begin{consignas}
  \item Identifique los complementos internos al sintagma verbal y los adjuntos
        presentes en los predicados.
  \item Indique si alguno de los verbos posee un argumento implícito.
\end{consignas}

\begin{resuelto}
a) \forma{Cumpliría}: \forma{diecinueve años} (complemento) y \forma{en abril}
(adjunto). \forma{Inspiraban}: \forma{los hombres}, \forma{le} y \forma{un temor
casi patológico} (complementos); \forma{aún} (adjunto). b) \forma{Cumpliría}
tiene el argumento sujeto tácito o implícito.
\end{resuelto}
\end{ejercicio}

\section{La estructura temática}\label{sec:07-tematica}

Además de la información acerca de la cantidad y características de los
argumentos requeridos por un verbo, es necesario reconocer que los distintos
argumentos cumplen diferentes funciones semánticas, conocidas como papeles
temáticos, es decir, que los predicados asignan un valor semántico a los
argumentos que seleccionan. ¿Cómo definimos los papeles temáticos?

\begin{destacado}
Los \termino{papeles temáticos} dan cuenta de la relación establecida por cada
participante en el evento descripto por el verbo.
\end{destacado}

\subsection{La estructura temática como esqueleto del significado: los papeles temáticos}\label{sec:07-papeles}

La \termino{estructura temática} de un predicado es la variante enriquecida de
su estructura argumental, es decir, contiene información sobre la semántica de
los argumentos que el predicado selecciona.

Los predicados otorgan a sus argumentos distintos papeles temáticos que expresan
su participación en el evento. Si bien no hay acuerdo generalizado respecto de
cuántos y cuáles son los papeles temáticos, en la siguiente lista incluimos los
más reconocidos y en los que hay mayormente consenso.

\begin{itemize}
  \item \textbf{Agente}: el que voluntariamente causa o realiza una acción, como
        \forma{Emma} en \forma{Emma recogió el papel}. El agente es típicamente
        \sig{animado}. Cuando el causante de la acción no tiene el rasgo
        \sig{animado} se identifica como \textbf{Causa} o \textbf{Fuerza}, como
        \forma{el fuego} en \forma{El fuego destruyó la cabaña}, o como
        \textbf{Instrumento} ---el medio u objeto empleado para realizar la
        acción---, como \forma{el revólver} en \forma{Emma mató a Loewenthal con
        un revólver}.
  \item \textbf{Tema}: el objeto efectuado, afectado o trasladado, o cuya
        existencia o posición se predica, como \forma{la verja} en \forma{Emma
        empujó la verja}. Cuando es \sig{animado} suele denominarse
        \textbf{Paciente}, como \forma{el marinero} en \forma{Emma buscó al
        marinero}.
  \item \textbf{Experimentante}: el que experimenta un estado psicológico o
        cognitivo determinado, como \forma{Emma} en \forma{Emma sintió asco}.
  \item \textbf{Beneficiario}: el que recibe beneficio o daño
        ---\textbf{Maleficiario}--- de un suceso, como \forma{Perla} en
        \forma{Perla perdió su trabajo}.
  \item \textbf{Locativo}: el que designa una ubicación en relación con el
        evento, como \forma{en Almagro} en \forma{Pedro vivía en Almagro}.
  \item \textbf{Fuente} (u \textbf{Origen}): entidad o lugar a partir del cual se
        origina la actividad, como \forma{su socio} en \forma{Recibió el arma de
        su socio} o \forma{el Río de la Plata} en \forma{El viaje empieza en el
        Río de la Plata}.
  \item \textbf{Meta}: la entidad o lugar hacia el que se dirige la actividad,
        como \forma{la empleada} en \forma{Entregó los informes a la empleada} o
        \forma{la estación} en \forma{La avenida llegaba hasta la estación}.
  \item \textbf{Vía}: el lugar de tránsito del movimiento, como \forma{por la
        avenida 9 de Julio} en \forma{Llegó hasta el puerto por la avenida 9 de
        Julio}.
  \item \textbf{Cantidad} o \textbf{Medida}: la medida espacial o temporal del
        evento, como \forma{veinte centímetros} en \forma{Creció veinte
        centímetros}.
\end{itemize}

\begin{destacado}
\textbf{Criterios de reconocimiento}
\begin{itemize}
  \item Cada argumento recibe un solo papel temático.
  \item Cada papel temático se asigna a un solo argumento.
  \item Los papeles temáticos no coinciden con las funciones sintácticas
        tradicionales.
  \item Los adjuntos no forman parte de la estructura argumental.
\end{itemize}
\end{destacado}

\begin{ejercicio}\label{ej:07-papeles}
En el fragmento de «Emma Zunz» del ejercicio~\ref{ej:07-valencia}, señale los
papeles temáticos de los sintagmas subrayados e indique cuál es el predicado que
los selecciona, según se muestra en el ejemplo resuelto de abajo.

\textoaincorporar{los subrayados del primer párrafo de «Emma Zunz»: \textit{de la
fábrica de tejidos Tarbuch y Loewenthal; una carta; su padre; el sello y el
sobre; la letra desconocida; Emma; el señor Maier; un compañero de pensión de su
padre; la noticia; a la hija del muerto}}

\begin{resuelto}
\forma{De la fábrica de tejidos Tarbuch y Loewenthal} tiene el papel temático de
fuente y es seleccionado por el predicado \forma{volver}.
\end{resuelto}
\end{ejercicio}

\begin{recuadro}[Para ampliar]
\textcite{bosque2009} señalan algunos beneficios de trabajar con estructuras
temáticas, pero también advierten que este enfoque presenta algunos
inconvenientes.

\medskip\textbf{Ventajas de postular estructuras temáticas}
\begin{itemize}
  \item La estructura temática muestra la correspondencia entre papeles
        temáticos y construcciones estructurales canónicas; por ejemplo, SN para
        Agente y para Tema (encabezado por \forma{a}) y SP para Meta.
  \item La estructura temática permite deslindar la información sintáctica y la
        información semántica.
  \item Las estructuras temáticas permiten considerar clases semánticas
        naturales. Por ejemplo, podemos analizar los verbos del tipo [A, T, L]
        (Agente, Tema, Locativo) como una clase que agrupa a verbos tales como
        \forma{situar, colocar, ocultar, poner}, en lugar de asignarle
        individualmente una estructura temática a cada predicado.
  \item Las estructuras temáticas no informan sobre los adjuntos, sino sobre los
        argumentos. Algunas nociones pueden manifestarse como complementos con
        ciertos verbos y como adjuntos con otros predicados, lo que confirma que
        tiene importancia caracterizar a estos modificadores en relación con el
        predicado en el que participan. Por ejemplo, si comparamos
        \forma{Colocó los libros en la mesa del comedor} con \forma{Hizo la tarea
        en la mesa del comedor}, \forma{en la mesa del comedor} es complemento
        (argumento) en el primer caso y adjunto en el segundo.
\end{itemize}

\textbf{Inconvenientes de postular estructuras temáticas}
\begin{itemize}
  \item ¿Cuántos papeles temáticos tenemos exactamente? Lo que se pone en duda
        actualmente es que cada argumento de cada predicado tenga que recibir
        una etiqueta semántica que represente un concepto definido entre las
        nociones que articulan la teoría de la gramática.
  \item Algunos papeles temáticos son demasiado abarcadores. El ejemplo más
        claro lo representa el papel de Tema, que refiere tanto a OD canónicos,
        como en \forma{Comió una milanesa}, como a nociones menos específicas,
        como en \forma{Miró la pantalla de su celular}, donde el objeto no es ni
        afectado, ni efectuado, ni trasladado.
  \item ¿Cómo estar seguros de que identificamos el papel temático correcto? No
        siempre hay garantías demasiado claras y existen casos que pueden ser
        menos transparentes para su reconocimiento.
\end{itemize}
\end{recuadro}

En la estructura temática, los dos papeles de mayor centralidad son el agente y
el tema, que son los pilares de la oración transitiva, si bien también están
presentes en otros tipos de esquemas oracionales que revisaremos.

\subsection{La noción de agentividad}\label{sec:07-agentividad}

Las propiedades que permiten caracterizar a los agentes prototípicos son
\sig{intencionalidad}, \sig{control consciente sobre la acción}, \sig{origen,
causa o iniciador del proceso} y \sig{autonomía} (es decir, que el agente
preexiste al evento). Por otra parte, para identificar el agente se puede
aplicar una paráfrasis con el verbo \forma{hacer}. Así, en \exref{ex:07-juan},
esta prueba confirma que \forma{Juan} es agente.

\begin{exe}
  \ex\label{ex:07-juan}
    \begin{xlist}
      \ex {}[Juan] trabajó toda la tarde.
      \ex Lo que hizo Juan fue trabajar toda la tarde.
    \end{xlist}
\end{exe}

En cambio, en \exref{ex:07-telon}, si aplicamos esta prueba comprobamos que
\forma{el telón} no es agente, sino tema.

\begin{exe}
  \ex\label{ex:07-telon}
    \begin{xlist}
      \ex Cayó [el telón].
      \ex[*]{Lo que hizo el telón fue caer.}
    \end{xlist}
\end{exe}

Asimismo, solo una oración que posee un sujeto agentivo admite las paráfrasis
con construcciones finales, que indican la intencionalidad de quien realiza la
acción.

\begin{exe}
  \ex {}[Juan] trabajó para aumentar sus ingresos.
  \ex {}[Juan] camina para hacer ejercicio.
  \ex[*]{[El telón] cayó para terminar la función.}
  \ex[*]{[Juan] palideció para hacer ejercicio.}
\end{exe}

Los sujetos de los verbos de acción normalmente representan una «causa
externa», es decir que la transformación que puede sufrir el tema o paciente
tiene su origen en una causa que ejerce control sobre el evento y que puede ser
un agente que realiza deliberadamente esa acción. Por ejemplo, en \forma{Los
chicos arreglaron su cuarto}, el sujeto agente ejerce el control sobre el evento
y realiza deliberadamente la acción de \forma{arreglar el cuarto}. También puede
darse el caso de una causa desencadenante no animada que provoque una situación
determinada, como en \forma{El viento derribó el toldo}. O incluso un
instrumento que, por medio de un agente no manifiesto, origina un determinado
estado de cosas, como en \forma{El camión embistió el portón de hierro}, que
supone un conductor implícito.

\begin{ejercicio}\label{ej:07-agente}
Indique en el texto siguiente si los SSNN subrayados poseen el papel temático de
agente o no. Justifique su respuesta a partir de la aplicación de las pruebas
pertinentes, según se muestra en el ejemplo resuelto de abajo.

% Texto periodístico protegido: se incorpora a mano desde el original.
\textoaincorporar{el texto sobre la conferencia de David Hilbert en la Sorbona
(1900) y «los problemas de Hilbert», con los SSNN subrayados: \textit{un alemán
llamado David Hilbert; Hilbert; la legendaria lista}}

\begin{resuelto}
\forma{Un alemán llamado David Hilbert} es agente. Acepta la paráfrasis con
\forma{hacer} (\forma{Lo que hizo X fue ponerles a los asistentes la tarea de
matemáticas\ldots}) y la construcción final (\forma{Un alemán llamado David
Hilbert les puso a los asistentes la tarea de matemáticas (\ldots) para lograr
su propósito}).
\end{resuelto}
\end{ejercicio}

\section{El esquema transitivo}\label{sec:07-transitivo}

La noción de \termino{objeto directo} o \termino{complemento objeto directo}
(COD) es una herencia de la tradición gramatical, para la cual «el complemento
directo recibe la acción del verbo, e incluso [\ldots] esta pasa al objeto o se
manifiesta en él» \parencite[§\,34.1e]{rae2025}. Por ejemplo, en
\exref{ex:07-albaniles}, el SN \forma{la pared} constituye el objeto efectuado
por la acción expresada por el verbo y cumple la función sintáctica de COD.

\begin{exe}
  \ex\label{ex:07-albaniles} Los albañiles construyeron [la pared].
\end{exe}

Desde el punto de vista categorial, el objeto directo es un SN (\forma{El
cocinero preparó \textbf{la comida}}) o una oración
subordinada\footnote{La noción de oración subordinada se trata en el
\capref{13}.} (\forma{Los directores dijeron \textbf{que no les había gustado el
reglamento}}).

La definición tradicional de transitividad se basa en la existencia de acciones
realizadas por un agente (sujeto) que pasan o se trasladan a un objeto tema o
paciente. Sin embargo, existe una importante vaguedad e imprecisión en estos
conceptos, dado que no son aplicables a todas las estructuras transitivas
posibles. Por ejemplo, en una oración como \forma{María vio un avión}, el sujeto
posee el papel temático de experimentante y no ejerce ninguna acción que pueda
trasladarse al objeto.

Definimos, entonces, un \termino{verbo transitivo} como aquel que selecciona su
complemento y lo identifica con la marca de caso acusativo. Ese complemento
cumple la función sintáctica de COD. Léxicamente, los verbos transitivos
requieren la presencia de un complemento que complete su significado.

Desde el punto de vista sintagmático, la relación estrecha entre el V y el COD
puede observarse en el diagrama arbóreo \exref{ex:07-arbol-cod}.

\begin{exe}
  \ex\label{ex:07-arbol-cod}
      \begin{forest} gramatica
        [SV
          [SN\\(Suj)\\\textit{los albañiles}]
          [V\prima
            [V\\\textit{construyeron}]
            [SN\\(COD)\\\textit{la pared}]]]
      \end{forest}
\end{exe}

\subsection{Objeto efectuado y objeto afectado}

Existe otra distinción gramaticalmente relevante entre los casos en que el
objeto es \termino{efectuado} por la acción expresada por el verbo y aquellos en
que solamente es \termino{afectado} por ella, es decir, los casos en que la
acción se aplica a un objeto ya existente. Veamos la diferencia con un ejemplo
como \exref{ex:07-mesada}.

\begin{exe}
  \ex\label{ex:07-mesada}
    \begin{xlist}
      \ex\label{ex:07-mesada-a} Rubén construyó [la mesada].
      \ex\label{ex:07-mesada-b} Rubén pintó [la mesada].
    \end{xlist}
\end{exe}

En \exref{ex:07-mesada-a}, \forma{la mesada} refiere a un objeto que se
materializa a través de la acción del verbo \forma{construir}; en
\exref{ex:07-mesada-b}, en cambio, \forma{la mesada} señala un objeto
preexistente, que se modifica por la acción indicada por el verbo.

Los objetos directos admiten la presencia de un clítico acusativo (\forma{lo,
la, los, las}) que sustituye o dobla el objeto directo, como se muestra en la
paráfrasis siguiente:

\begin{exe}
  \ex Rubén [la] construyó.
\end{exe}

\subsection{El objeto directo animado y el marcado diferencial de objeto (MDO)}\label{sec:07-mdo}

En español, algunos objetos directos poseen una marca especial \forma{a}. Veamos
algunos ejemplos:

\begin{exe}
  \ex
    \begin{xlist}
      \ex Los docentes recibieron [a los alumnos] en el primer día de clase.
          (OD \sig{animado})
      \ex[*]{Los docentes recibieron [los alumnos] en el primer día de clase.}
    \end{xlist}
  \ex
    \begin{xlist}
      \ex Los empleados recibieron [la intimación]. (OD \sig{inanimado})
      \ex[*]{Los empleados recibieron [a la intimación].}
    \end{xlist}
  \ex\label{ex:07-perro}
    \begin{xlist}
      \ex Mariano encontró [un perro abandonado]. (OD \sig{indeterminado})
      \ex[?]{\label{ex:07-perro-b} Mariano encontró [a un perro abandonado].}
      \ex Mariano encontró [al perro abandonado].
    \end{xlist}
\end{exe}

En términos generales, el objeto directo lleva la marca de función \forma{a}
cuando refiere a una entidad animada, determinada y específica, aunque puede
haber excepciones, como \exref{ex:07-perro-b}. ¿Por qué \exref{ex:07-perro-b}
nos resulta extraña o poco aceptable? La oración resulta anómala porque lleva la
preposición \forma{a} un objeto directo inespecífico: \forma{un perro
abandonado}.

Algunos usos impersonales de verbos tales como \forma{haber} (\forma{No hay
plata}) y \forma{hacer} (\forma{Hace más de treinta años}) seleccionan un
complemento objeto directo de carácter marginal, ya que la sustitución por el
clítico acusativo (\forma{lo, la}) es solo posible en ciertos contextos
marcados.

\begin{ejercicio}\label{ej:07-mdo}
Señale en los textos siguientes las características de los objetos directos
subrayados (± animado, ± específico). Fundamente su respuesta, según se muestra
en los ejemplos resueltos de abajo.

% Textos protegidos: se incorporan a mano desde el original.
\textoaincorporar{Texto 1: el fragmento sobre Cantor y el infinito, con los COD
subrayados \textit{el significado del infinito; precisión matemática; un solo
infinito}. Texto 2: el fragmento de «Emma Zunz» que empieza «Recordó veraneos en
una chacra», con los COD subrayados \textit{veraneos; a su madre; la casita de
Lanús; los amarillos losanges de una ventana; el auto de prisión; el oprobio; los
anónimos}}

\begin{resuelto}[Ejemplos resueltos]
\forma{El significado del infinito}: COD \sig{inanimado} y \sig{específico}.
\forma{Un solo infinito}: COD \sig{inanimado} e \sig{inespecífico}, seleccionado
por un verbo impersonal de significado \sig{existencial}.
\end{resuelto}
\end{ejercicio}

\subsection{Los complementos preposicionales}\label{sec:07-regimen}

Si bien las gramáticas tradicionales solían incluir los complementos
preposicionales entre los circunstanciales, no todos ellos se comportan en
forma análoga. En el caso del \termino{complemento de régimen}, se trata de un
modificador seleccionado por el núcleo del predicado, es decir, tiene carácter
argumental. Por otra parte, el núcleo del predicado selecciona también la
preposición que lo encabeza. Veamos algunos ejemplos:

\begin{exe}
  \ex\label{ex:07-walmart} La marca Walmart se despide [del mercado argentino].
  \ex\label{ex:07-discutio} El ministro discutió [con los empresarios].
  \ex\label{ex:07-penso} Juan pensaba [en su madre].
  \ex\label{ex:07-quejo} Juana se quejó [de su madre].
\end{exe}

El complemento de régimen preposicional no puede ser sustituido ni duplicado por
un pronombre clítico acusativo o dativo. Algunos verbos, como \forma{despedir},
\forma{pensar} y \forma{discutir}, se comportan como transitivos en oraciones
tales como \forma{María despidió a su padre}, \forma{El ministro discutió sus
argumentos} o \forma{Juan pensaba la solución del problema}, en las que
seleccionan un COD, a diferencia de oraciones como \exref{ex:07-walmart},
\exref{ex:07-discutio} y \exref{ex:07-penso}, en las que se construyen con un
complemento de régimen. En cambio, el verbo \forma{quejar(se)}, en
\exref{ex:07-quejo}, no posee esa alternancia y siempre se construye con un
complemento encabezado por la preposición \forma{de}, aunque en ciertos
contextos el complemento puede quedar implícito: \forma{Juana se quejó todo el
tiempo}.

\subsection{Otros complementos}

Algunos verbos seleccionan complementos encabezados por una preposición cuyo
significado determina el papel temático del complemento que lo acompaña. Es el
caso de los \termino{complementos locativos} o \termino{de modo}:

\begin{exe}
  \ex\label{ex:07-volvia} Volvía [de Mar del Plata] por aquel entonces.
  \ex\label{ex:07-conduce} Conduce [a su manera].
\end{exe}

Los complementos señalados en \exref{ex:07-volvia} y \exref{ex:07-conduce} están
seleccionados por los verbos \forma{volver} y \forma{conducir}, respectivamente.
Se trata de complementos circunstanciales y no debemos confundirlos con los
adjuntos circunstanciales, que no son requeridos por el verbo, como, por
ejemplo, en:

\begin{exe}
  \ex\label{ex:07-trabaja} Trabaja (en Mar del Plata) (durante el verano).
  \ex\label{ex:07-pinto} Pintó la silla (a su manera).
\end{exe}

En \exref{ex:07-trabaja} y \exref{ex:07-pinto}, las construcciones entre
paréntesis funcionan como adjuntos circunstanciales de lugar y modo,
respectivamente.

\subsection{Esquemas ditransitivos}\label{sec:07-ditransitivos}

Son aquellos esquemas oracionales en los que el verbo selecciona dos argumentos
internos: uno directo y otro indirecto. En una oración como \forma{Los
socorristas pidieron [ayuda] [a los bomberos]}, los sintagmas entre corchetes
constituyen los dos complementos de objeto: el primero, el directo, y el
segundo, el indirecto. Típicamente, existen tres subclases de verbos que
constituyen esquemas ditransitivos: los verbos de comunicación
\exref{ex:07-anuncio}, los de transferencia \exref{ex:07-presto} y los de
influencia \exref{ex:07-solicito}. Veamos los ejemplos siguientes:

\begin{exe}
  \ex\label{ex:07-anuncio} El directorio anunció las medidas [a sus miembros].
  \ex\label{ex:07-presto} Prestó su bolsa de dormir [a su sobrino].
  \ex\label{ex:07-solicito} Solicitó [a sus jefes] una licencia.
\end{exe}

Los objetos indirectos admiten la presencia de un clítico dativo (\forma{le} o
\forma{les}), que sustituye o dobla el objeto indirecto, como se muestra en las
paráfrasis siguientes:

\begin{exe}
  \ex {}[Les] anunció las medidas.
  \ex {}[Le] prestó su bolsa de dormir.
  \ex {}[Les] solicitó una licencia.
\end{exe}

En el diagrama que presentamos a continuación se muestra la configuración
estructural del complemento objeto indirecto (COI) en el sintagma verbal y la
diferencia con la configuración estructural del COD.

\begin{exe}
  \ex\label{ex:07-arbol-coi} El directorio anunció las medidas a sus miembros.\\
      \begin{forest} gramatica
        [SV
          [V\prima
            [V\prima
              [V\\\textit{anunció}]
              [SN\\(COD)\\\textit{las medidas}]]
            [SP\\(COI)\\\textit{a sus miembros}]]]
      \end{forest}
\end{exe}

Otra clase de verbos ditransitivos (diádicos o bivalentes) son los llamados
verbos de «colocación», que seleccionan dos argumentos internos: un tema (COD) y
un locativo (CCL), como se expresa en los ejemplos que presentamos a
continuación:

\begin{exe}
  \ex Colocó estantes en la cocina.
  \ex Colgó los cuadros en la pared del dormitorio.
\end{exe}

\begin{ejercicio}\label{ej:07-seleccion2}
Lea atentamente el siguiente fragmento de «Emma Zunz» y responda lo pedido según
se muestra en los ejemplos resueltos de abajo.

\begin{consignas}
  \item Señale qué verbos seleccionan los sintagmas y pronombres clíticos
        subrayados.
  \item Relacione los sintagmas y pronombres clíticos subrayados con los verbos
        que los seleccionan.
\end{consignas}

% Fragmento literario protegido: se incorpora a mano desde el original.
\textoaincorporar{los dos fragmentos de «Emma Zunz» que van de «Emma, desde
1916, guardaba el secreto» hasta «\ldots{}el sabor de la victoria y de la
justicia» y de «Dio al fin con hombres del Nordstjärnan» hasta «\ldots{}a una
puerta que se cerró», con los subrayados: \textit{a nadie; a su mejor amiga,
Elsa Urstein; le} (dos veces)\textit{; hombres del Nordstjärnan; por otro; a una
puerta; a un turbio zaguán; a una escalera tortuosa; a un vestíbulo; a un
pasillo; a una puerta que se cerró}}

\begin{resuelto}[Ejemplos resueltos]
\forma{A nadie} y \forma{a su mejor amiga, Elsa Urstein} son seleccionados por
\forma{revelar}. a) El primer \forma{le} es seleccionado por \forma{inspirar}.
El segundo \forma{le} es seleccionado por el verbo \forma{deparar}. b)
\forma{Revelar} es un verbo de habla y selecciona un complemento objeto
indirecto, y \forma{deparar} es un verbo psicológico y selecciona un complemento
objeto indirecto.
\end{resuelto}
\end{ejercicio}

\subsection{Esquemas intransitivos: inergativos e inacusativos}\label{sec:07-intransitivos}

Los verbos intransitivos se han caracterizado tradicionalmente como aquellos que
carecen de objeto directo y solo disponen de un argumento sujeto. Sin embargo,
semánticamente, la estructura de los intransitivos es más compleja y hoy en día
se reconocen dos subclases de intransitivos: inergativos e inacusativos.

\subsubsection{Verbos inergativos}\label{sec:07-inergativos}

Son aquellos en los que el sujeto desempeña el papel temático de agente. Se
agrupan en diferentes clases:

\begin{enumerate}[label=\alph*)]
  \item Verbos que seleccionan un sujeto encargado de realizar la acción que
        tiene lugar en el propio individuo: \forma{bromear, gesticular,
        rezongar, temblar}.
        \begin{exe}
          \ex Los encargados de la seguridad gesticularon con muecas divertidas.
        \end{exe}
  \item Verbos que implican la emisión de un sonido o sustancia: \forma{gritar,
        ladrar, llorar, sudar}.
        \begin{exe}
          \ex El cachorro ladró durante toda la noche.
        \end{exe}
  \item Verbos que señalan la manera de moverse, que pueden tener un uso
        claramente inergativo: \forma{bailar, caminar, correr}.
        \begin{exe}
          \ex Los egresados bailaron hasta el amanecer.
        \end{exe}
\end{enumerate}

\subsubsection{Verbos inacusativos}\label{sec:07-inacusativos}

El segundo grupo de verbos intransitivos es el de los inacusativos; en ellos, el
sujeto es tema o paciente. Los verbos inacusativos se agrupan en diversos tipos:

\begin{enumerate}[label=\alph*)]
  \item Verbos de causa interna, que nunca admiten la existencia de un agente o
        causa externa, ya que refieren a un proceso en el que el sujeto puede
        cambiar de manera innata: \forma{arder, madurar, estallar, florecer,
        germinar}.
        \begin{exe}
          \ex La bomba molotov no estalló.
          \ex Los cerezos florecieron durante la primavera.
        \end{exe}
  \item Verbos que denotan cambio de estado psíquico o físico y que poseen
        también usos transitivos: \forma{adelgazar, engordar, hervir, encoger,
        calentarse, congelarse}.
        \begin{exe}
          \ex Se congelaron las cañerías durante el invierno.
          \ex La ropa encogió después del lavado.
        \end{exe}
  \item Verbos de movimiento de dirección inherente cuando se emplean con
        sujetos \sig{inanimados}: \forma{acercarse, alejarse, partir,
        aproximarse}.
        \begin{exe}
          \ex El auto blindado se alejó rápidamente.
          \ex El avión partió con retraso.
        \end{exe}
  \item Otros verbos, como \forma{moverse, bambolearse, temblar}, también
        presentan usos inacusativos.
        \begin{exe}
          \ex La silla se bamboleaba.
          \ex El gatito temblaba por la anestesia.
        \end{exe}
  \item Verbos de existencia o aparición, que tienden a dar lugar a
        construcciones de inversión locativa, es decir, el complemento de lugar
        aparece en posición inicial, como se puede observar en los ejemplos
        \exref{ex:07-bosque} y \exref{ex:07-alegria}.
        \begin{exe}
          \ex\label{ex:07-bosque} En el bosque aparecieron fantasmas.
          \ex\label{ex:07-alegria} En esta casa falta alegría.
        \end{exe}
\end{enumerate}

\paragraph{Comportamiento sintáctico de los verbos inacusativos.} Las pruebas
para el reconocimiento de los verbos inacusativos se basan en la idea de que el
sujeto de las estructuras inacusativas posee las mismas características
semánticas del argumento objeto directo de un verbo transitivo, aunque
sintácticamente adquiere rasgos de concordancia con el verbo que lo hacen
funcionar como sujeto de la oración. Algunas de las pruebas que se han
presentado para reconocer esta clase de verbos son las siguientes:

\begin{itemize}
  \item Los sujetos de los inacusativos, al igual que los objetos de los
        transitivos, pueden aparecer en construcciones de participio absoluto,
        \exref{ex:07-barco} y \exref{ex:07-avion}, lo cual no es posible para
        los inergativos, \exref{ex:07-ladrado} y \exref{ex:07-caminado}.
\end{itemize}

\begin{exe}
  \ex\label{ex:07-barco} {}[Llegado el barco al muelle], subimos rápidamente.
  \ex\label{ex:07-avion} {}[Una vez aterrizado el avión], me tranquilicé.
  \ex[*]{\label{ex:07-ladrado} [Ladrado el perro], advertimos que algo pasaba.}
  \ex[*]{\label{ex:07-caminado} [Caminado Pedro], nos fuimos con los demás.}
\end{exe}

\begin{itemize}
  \item La prueba anterior no es posible con los verbos que implican una
        realización gradual, como \forma{aumentar, mejorar, crecer}, porque las
        construcciones de participio absoluto exigen verbos que expresen el
        estado final.\footnote{Los verbos que indican el estado final se conocen
        como télicos y se estudian en el \capref{8}.}
\end{itemize}

\begin{exe}
  \ex[*]{[Aumentada la nafta], disminuyó el tráfico.}
  \ex[*]{[Crecidas las plantas], empezamos a disfrutar del jardín.}
\end{exe}

\begin{itemize}
  \item Los verbos inacusativos nunca admiten objetos directos
        cognados\footnote{Se denominan objetos o complementos cognados aquellos
        que comparten la base léxica del verbo, como \forma{soñar un sueño
        apacible} o \forma{vivir una vida armoniosa}.} \exref{ex:07-llegada}, a
        diferencia de los inergativos, que sí los aceptan \exref{ex:07-carrera}.
\end{itemize}

\begin{exe}
  \ex[*]{\label{ex:07-llegada} Llegó [una llegada fácil].}
  \ex\label{ex:07-carrera} Corrió [una carrera asombrosa].
\end{exe}

\begin{itemize}
  \item Los participios de los verbos transitivos e inacusativos pueden
        funcionar como modificadores cercanos a los adjetivos: \forma{una
        promesa rota, una obra acabada, una persona desaparecida, un bebé recién
        nacido}, pero no sucede lo mismo con los participios de los inergativos:
        \juicio{*}{\forma{una mujer sonreída}}. \textcite{bello2002} llamó
        \termino{participios deponentes} a los de los verbos
        inacusativos.\pendiente{El original cita «Bello (1928)»}
  \item Los verbos inacusativos tienden a ser incompatibles con el sufijo
        nominalizador \forma{-dor}, puesto que tiene valor agentivo:
        \juicio{*}{\forma{desaparecedor}}, \juicio{*}{\forma{crecedor}},
        \juicio{*}{\forma{moridor}}. Los inergativos, en cambio, aceptan bien
        este sufijo: \forma{corredor, trabajador, caminador}.
\end{itemize}

\subsection{Verbos con alternancia causativa-anticausativa incoativa}\label{sec:07-alternancia}

Existen verbos que denotan un cambio de estado físico o psicológico, que poseen
usos inacusativos y provienen de verbos transitivos. Veamos algunos ejemplos:

\begin{exe}
  \ex
    \begin{xlist}
      \ex La inundación estropeó las cañerías.
      \ex\label{ex:07-estropearon} Las cañerías se estropearon.
    \end{xlist}
  \ex
    \begin{xlist}
      \ex El calor derritió los helados.
      \ex\label{ex:07-derritieron} Los helados se derritieron.
    \end{xlist}
\end{exe}

En los ejemplos \exref{ex:07-estropearon} y \exref{ex:07-derritieron}, el
sujeto tiene el papel temático de tema. Se trata de usos inacusativos en los que
el propio objeto desencadena la acción que señala el verbo, sin que intervenga
ningún agente \parencite{rodriguezramalle2005}. A los verbos que intervienen en
este paradigma se los suele llamar verbos de \termino{alternancia
causativa-anticausativa incoativa}. Se los denomina incoativos porque estas
construcciones inacusativas señalan el comienzo de un proceso, aunque suelen
terminar en un estado final o resultante. Esto sucede especialmente cuando
proyectan un clítico \forma{se}.\footnote{Al incorporar el clítico \forma{se},
estos verbos se convierten en pronominales cuasirreflejos. Estas construcciones
se estudian en el \capref{10}.}

\begin{recuadro}[Para ampliar]
En otras lenguas romances, como el francés y el italiano, donde estos verbos
fueron primeramente identificados, existen algunos diagnósticos para reconocer
los verbos inacusativos. En primer lugar, para formar los tiempos compuestos,
los verbos inacusativos seleccionan el auxiliar \forma{essere / être}
(\sig{ser}). En cambio, los inergativos seleccionan \forma{avere / avoir}
(\sig{haber}).

\begin{exe}
  \ex \gll Mara è partita alle dieci.\\
           Mara es partida a.las diez\\
      \glt \sig{Mara partió a las diez.}
  \ex \gll I bimbi sono ormai cresciuti.\\
           los niños son ya crecidos\\
      \glt \sig{Los niños han crecido ya.}
\end{exe}

Ahora bien, cabe señalar que en italiano todos los verbos inacusativos
seleccionan \forma{essere}, pero no todos los que seleccionan el auxiliar
\forma{essere} son inacusativos. Por ejemplo, los reflexivos se construyen con
\forma{essere}, pero no son inacusativos:

\begin{exe}
  \ex \gll Giovanni si era guardato allo specchio a lungo.\\
           Juan se era mirado en.el espejo a largo\\
      \glt \sig{Juan se miró en el espejo durante largo rato.}
\end{exe}

En segundo lugar, otra propiedad típica del objeto directo que el sujeto
pospuesto de los verbos inacusativos presenta en italiano es la posibilidad de
usar el pronombre partitivo \forma{ne}:

\begin{exe}
  \ex \gll Ne sono arrivati gli autori.\\
           de.ellos son llegados los autores\\
      \glt \sig{Llegaron los autores.}
\end{exe}

Pero este reemplazo no es posible con los verbos inergativos:

\begin{exe}
  \ex \gll {*}Ne hanno telefonato gli autori.\\
           de.ellos han telefoneado los autores\\
      \glt \sig{Llamaron por teléfono los autores.}
\end{exe}

En síntesis, podemos concluir que, en lenguas como el italiano y el francés, los
diagnósticos de inacusatividad tienen un carácter claramente sintáctico. Esto no
es tan evidente en el español actual. No obstante, sabemos que el español
medieval seleccionaba el auxiliar \forma{ser} para los verbos inacusativos:

\begin{exe}
  \ex Los moços somos idos a comer y nos han dexado solos. (Juan de Valdés,
      \textit{Diálogo de la lengua}, hacia 1535)\\
      \sig{Los mozos se han ido a comer y nos han dejado solos.}
\end{exe}
\end{recuadro}

\begin{ejercicio}\label{ej:07-inacusativos}
Considere el siguiente texto (A) y las oraciones de (B). Luego, según se muestra
en el ejemplo resuelto de abajo:

\begin{consignas}
  \item clasifique los verbos subrayados en transitivos, inergativos e
        inacusativos. Fundamente su respuesta con las pruebas adecuadas;
  \item analice las alternancias propuestas en (B) y explique qué cambios se
        producen en relación con los esquemas transitivos, intransitivos e
        inacusativos.
\end{consignas}

\textbf{A}

% Fragmento literario protegido: se incorpora a mano desde el original.
\textoaincorporar{los fragmentos de «Emma Zunz» que van de «De pronto,
alarmada, se levantó» hasta «\ldots{}la empezó a leer y la rompió» y de «El
hombre la condujo a una puerta» hasta «\ldots{}las partes que los forman», con
los verbos subrayados: \textit{levantarse, correr, abrir, dejar, ver, leer,
romper, conducir, haber, cerrarse, formar}}

\textbf{B}

\begin{exe}
  \ex
    \begin{xlist}
      \ex Abrió el cajón de la cómoda.
      \ex El cajón de la cómoda se abrió de repente.
    \end{xlist}
  \ex
    \begin{xlist}
      \ex El déficit empeoró la balanza comercial.
      \ex Las condiciones meteorológicas empeoraron en las últimas horas.
    \end{xlist}
  \ex
    \begin{xlist}
      \ex Una puerta se cerró.
      \ex Cerró la puerta con llave.
    \end{xlist}
\end{exe}

\begin{resuelto}
a) \forma{Levantarse} y \forma{correr} son verbos intransitivos. El sujeto de
ambos es agente y el verbo pertenece al grupo de los que seleccionan un sujeto
encargado de realizar la acción que tiene lugar en el propio individuo. b) La
oración 1a es transitiva: el sujeto es agente y posee un complemento objeto
directo tema. 1b es un caso de alternancia causativa inacusativa en la que el
verbo es inacusativo, pero proviene de un verbo transitivo, \forma{abrir}.
\end{resuelto}
\end{ejercicio}

\begin{recuadro}[Reflexión final]
En este capítulo nos hemos concentrado en la relación entre la semántica y la
sintaxis. Hemos definido los conceptos de argumento y predicado; luego, nos
hemos ocupado de la estructura temática y hemos definido las funciones
semánticas o papeles temáticos. A partir de allí, presentamos las funciones
sintácticas, específicamente aquellas incluidas en el sintagma verbal. Para
finalizar, desarrollamos los esquemas verbales transitivo e intransitivo y,
dentro de los últimos, hemos distinguido entre los inergativos y los
inacusativos.
\end{recuadro}

% --- Secciones finales -------------------------------------------------
\lecturasrecomendadas

\solucionario

\begin{solucion}{\ref{ej:07-agramatical}}
3. Es agramatical porque faltan dos argumentos, referidos a \sig{lo colocado} y
\sig{el lugar de colocación}. 4. Es agramatical porque falta un argumento
referido a \sig{lo que provocó el viento}. 5. Hay una incompatibilidad
léxico-semántica entre el verbo \forma{anunciar} y el argumento
\forma{gelatina}. 6. Es ambigua: a) es gramatical en la interpretación de
\sig{hizo un reclamo a los bancos}; b) es agramatical en la interpretación en
que falta el argumento referido a \sig{lo que reclamó el presidente}. 7. Es
agramatical porque falta un argumento referido al lugar. 8. Es agramatical
porque falta un argumento referido al lugar.\pendiente{El original numera estas
respuestas como (8) y (9)}
\end{solucion}

\begin{solucion}{\ref{ej:07-seleccion}}
a) 2 es gramatical, aunque hay un argumento implícito. 3 es gramatical, aunque
tiene dos argumentos implícitos. 5 es gramatical.\pendiente{Falta la preposición
\textit{a} ante \textit{su casa}: revisar el juicio} 6 es agramatical porque
\forma{a su aceituna} no constituye un argumento de lugar. b) \forma{Regresar}
[Arg.~1 SN, Arg.~2 SP].
\end{solucion}

\begin{solucion}{\ref{ej:07-valencia}}
a) Son bivalentes porque seleccionan dos argumentos: \forma{engañar, inquietar,
ingerir, firmar} y \forma{dirigir(se)}. El único monovalente, que selecciona un
solo argumento, es \forma{fallecer}. b) Los argumentos de \forma{engañar} son
\forma{el sello y el sobre} (sujeto) y el pronombre personal en caso acusativo
\forma{la} (COD). Los argumentos de \forma{inquietar} son \forma{la letra
desconocida} (sujeto) y el pronombre personal en caso acusativo \forma{la}
(COD). Los argumentos de \forma{ingerir} son \forma{el señor Maier} (sujeto) y
\forma{una fuerte dosis de veronal} (COD). El argumento de \forma{fallecer} es el
sujeto tácito, correferencial con el sujeto de \forma{ingerir}. Los argumentos
de \forma{firmar} son \forma{un compañero de pensión de su padre} (sujeto) y
\forma{la noticia} (COD). Los argumentos de \forma{dirigir(se)} son el sujeto
tácito, correferencial con el sujeto de \forma{firmar}, y \forma{la hija del
muerto} (COD).\pendiente{Revisar: \textit{a la hija del muerto} no es COD de
\textit{dirigirse}}
\end{solucion}

\begin{solucion}{\ref{ej:07-compadj}}
a) \forma{Preparó}: \forma{una sopa de tapioca} y \forma{unas legumbres}
(complementos). \forma{Comió}: \forma{temprano} (adjunto). \forma{Se acostó} y
\forma{se obligó a dormir}: no tienen complementos ni adjuntos. b) Los cuatro
verbos tienen el sujeto tácito o implícito; \forma{comió} tiene también el COD
implícito.
\end{solucion}

\begin{solucion}{\ref{ej:07-papeles}}
\forma{Una carta}: tema seleccionado por \forma{hallar}. \forma{El sello y el
sobre}: causa seleccionada por \forma{engañar}.\pendiente{El original dice «Una
carta y el sobre»} \forma{Su padre}: paciente seleccionado por \forma{morir}.
\forma{La letra desconocida}: causa seleccionada por \forma{inquietar}.
\forma{Emma}: agente seleccionado por \forma{leer}. \forma{El señor Maier}:
agente seleccionado por \forma{ingerir}. \forma{Un compañero de pensión de su
padre}: agente seleccionado por \forma{firmar}. \forma{La noticia}: tema
seleccionado por \forma{firmar}. \forma{A la hija del muerto}: meta seleccionada
por \forma{dirigir(se)}.
\end{solucion}

\begin{solucion}{\ref{ej:07-agente}}
\forma{Hilbert} no es agente porque rechaza la paráfrasis con \forma{hacer}
(\juicio{*}{\forma{Lo que hizo Hilbert fue ser\ldots}}) y la construcción final
(\juicio{*}{\forma{Hilbert era uno de los ponentes del Congreso para realizar su
investigación}}). \forma{La legendaria lista} no es agente porque es inanimada y
no admite la paráfrasis con \forma{hacer} (\juicio{*}{\forma{Lo que hizo la
legendaria lista\ldots}}) ni la construcción final (\juicio{*}{\forma{La
legendaria lista definió las matemáticas de la era moderna para lograr su
propósito}}).
\end{solucion}

\begin{solucion}{\ref{ej:07-mdo}}
Texto 1. \forma{Precisión matemática}: COD \sig{inanimado} y
\sig{específico}.\pendiente{Como COD escueto, \textit{precisión matemática} es
más bien inespecífico} Texto 2. \forma{Veraneos}: COD \sig{inanimado} e
\sig{inespecífico}. \forma{A su madre}: \forma{a} es marca de función (MDO); COD
\sig{animado} y \sig{específico}. \forma{La casita de Lanús}: COD
\sig{inanimado} y \sig{específico}. \forma{Los amarillos losanges de una
ventana}: COD \sig{inanimado} y \sig{específico}. \forma{El auto de prisión}:
COD \sig{inanimado} y \sig{específico}. \forma{El oprobio}: COD \sig{inanimado}
y \sig{específico}. \forma{Los anónimos}: COD \sig{inanimado} y
\sig{específico}.
\end{solucion}

\begin{solucion}{\ref{ej:07-seleccion2}}
b) \forma{Dar con} = \sig{encontrar}; \forma{optar por} = \sig{elegir};
\forma{conducir a}. \forma{Hombres del Nordstjärnan} está seleccionado por el
verbo \forma{dar (con)}. \forma{Por otro} está seleccionado por el verbo de
\sig{elección} \forma{optar}. \forma{A una puerta, a un turbio zaguán, a una
escalera tortuosa, a un vestíbulo, a un pasillo} y \forma{a una puerta que se
cerró} están seleccionados por el verbo de \sig{movimiento} \forma{conducir}.
\end{solucion}

\begin{solucion}{\ref{ej:07-inacusativos}}
a) \forma{Abrir, dejar, ver, leer, romper} y \forma{conducir} son verbos
transitivos. Todos seleccionan dos argumentos: un agente y un tema o paciente.
\forma{Haber} es un verbo impersonal transitivo que se construye con un COD con
el papel temático de tema. \forma{Cerrarse} es la variante inacusativa
pronominal del verbo transitivo \forma{cerrar}. \forma{Formar} es un verbo
transitivo con un COD pronominal.

b) 2. \forma{Empeorar}: alternancia transitiva / inacusativa sin incorporación
del clítico \forma{se}. En a) el sujeto es causa y en b) es tema. 3. Alternancia
transitiva / inacusativa \forma{cerrarse / cerrar}. El sujeto del verbo
inacusativo de a) se convierte en COD del verbo transitivo en b).
\end{solucion}

% =====================================================================
%  capitulos/cap08.tex — Capítulo 8. La estructura eventiva
%  Autora: Mabel Giammatteo
%  Fuente: Capitulo_08_normalizado.docx
% =====================================================================
\capitulo{La estructura eventiva}{Mabel Giammatteo}

\section{Introducción}

Según hemos visto en el \capref{6}, la existencia misma de la oración depende de
su anclaje temporal respecto del momento de la emisión. Pero, además del tiempo,
la organización de la oración obedece a otra categoría que, desde el punto de
vista gramatical, en español se da casi siempre amalgamada morfológicamente con
el tiempo, pero que, como veremos, tiene una gran incidencia sintáctica: el
aspecto. A diferencia del tiempo, que es más externo, el aspecto tiene que ver
con la manera en que un evento se desarrolla u ocurre ---con o sin cambio; con o
sin límite; de forma única o repetida; en un período, acotado o no, o en un
instante; etc.---.

En este capítulo veremos cómo las propiedades más relevantes del aspecto
---duración, delimitación, dinamismo, reiteración--- han permitido identificar
los distintos subtipos de eventos que los verbos expresan. Esta clasificación
resulta de gran utilidad porque, como sostienen \textcite{bosque2009}, sirve
para explicar las propiedades combinatorias de los verbos con sus complementos
y adjuntos. Es necesario tener en cuenta, como plantean los mencionados autores,
que «[a]lgunos verbos pertenecen unívocamente a una sola clase, pero la mayor
parte puede pertenecer a varias si el contexto sintáctico lo permite»
\parencite[307]{bosque2009}. Por eso, también nos ocuparemos de mostrar cómo la
información aportada por el lexema verbal puede ser modificada por la que
aportan sus complementos y modificadores, o incluso por el tiempo y modo verbal
en que se encuentra.

\section{El evento}\label{sec:08-evento}

Informalmente, un \termino{evento} es cualquier clase de situación o suceso
denotado por un predicado \parencite[158]{pustejovsky2019}. A diferencia de los
objetos, que se ubican en el espacio, como \forma{el libro} en
\exref{ex:08-libro}, los eventos, como \forma{ir al cine} en \exref{ex:08-cine},
ocurren y transcurren en el tiempo.

\begin{exe}
  \ex\label{ex:08-libro} El libro está sobre la mesa.
  \ex\label{ex:08-cine} Fuimos al cine ayer a la noche.
\end{exe}

Si bien el tiempo y el aspecto son dos categorías que tradicionalmente se han
relacionado con el evento, ambas manifiestan nociones diferentes:

\begin{destacado}
Mientras el \termino{tiempo} relaciona al evento con una referencia externa que
da cuenta de su anterioridad, simultaneidad o posterioridad con respecto al
momento de la emisión, el \termino{aspecto} constituye una determinación más
interna, que provee información sobre el evento en sí, es decir, sobre cómo
ocurre o transcurre, si se repite o cuánto se extiende.
\end{destacado}

Actualmente se reconocen dos formas de expresar el aspecto:

\begin{itemize}
  \item El \termino{aspecto léxico} o \termino{de situación}, que transmite
        información sobre el tipo de evento de que se trata y se manifiesta en
        la raíz verbal.
  \item El \termino{aspecto gramatical} o \termino{flexivo}, que permite al
        hablante seleccionar el punto de vista desde el que presenta al evento y
        que es manifestado por el verbo en su desinencia.
\end{itemize}

\subsection{Las propiedades de los eventos}\label{sec:08-propiedades}

Si bien ambas formas de expresar la aspectualidad interactúan en la oración,
para la clasificación de los eventos, en principio, partimos del aspecto léxico.
Siguiendo a \textcite{pustejovsky2019}, algunas de las propiedades relevantes de
esta categoría, que inciden en la clasificación de eventos que veremos más
adelante, son las siguientes:

\begin{itemize}
  \item La \termino{duratividad}, que distingue entre eventos instantáneos o de
        duración mínima ---\forma{entrar, disparar} (un arma),
        \forma{encontrar}--- y eventos durativos, cuyo desarrollo implica un
        cierto tiempo ---\forma{estudiar, escribir, comer}---.
  \item La \termino{delimitación}, que diferencia entre eventos que terminan
        ---\forma{Juan construyó una casa en diez meses}--- y eventos que se
        extienden por un período de tiempo y que, en principio, no implican un
        límite ---\forma{Juan ama a su mujer}---.
  \item La \termino{dinamicidad}, que tiene en cuenta si los eventos producen
        cambios o no. Los estados, como \forma{amar, vivir} o \forma{ser alto},
        no involucran transformaciones y no pueden limitarse; así, no podemos
        decir \juicio{*}{\forma{Juan es alto hasta mañana}}. A diferencia de esto, los eventos dinámicos producen
        cambios a través del tiempo: \forma{Juan pintó la pared} (el evento
        produjo un cambio en el estado físico de la pared).
  \item La \termino{telicidad}, según la cual algunos eventos dinámicos tienen
        «una culminación natural» (un \textit{telos}\footnote{La palabra
        \textit{telos} proviene del griego τέλος, \sig{fin}, \sig{objetivo},
        \sig{propósito} o \sig{meta}.} interno), que impide que se prolonguen;
        así, por ejemplo, \forma{comer una manzana} o \forma{fumar un cigarrillo}
        son eventos télicos que no admiten continuarse: \forma{Juan comió una
        manzana (*y siguió haciéndolo toda la tarde)}.\footnote{Por convención,
        en los ejemplos gramaticales se marcan entre paréntesis las secuencias
        optativas. Si se coloca un asterisco antes del paréntesis inicial, la
        secuencia incluida se considera requerida y no puede omitirse. Si, en
        cambio, el asterisco aparece después del paréntesis inicial, se lee que
        agregar la secuencia delimitada por los paréntesis a la oración produce
        agramaticalidad.} En cambio, \forma{Juan leía} es un evento atélico (no
        télico) porque no tiene una culminación en sí mismo y acepta ser
        continuado: \forma{Juan leía y siguió haciéndolo toda la tarde}.
  \item La \termino{iteración}, que implica que algunos eventos son únicos y se
        producen una sola vez ---\forma{Ayer fui al cine} o \forma{Echó un
        tremendo bostezo}--- y otros se repiten ---\forma{Siempre voy al cine los
        fines de semana; Bostezaba sin parar} (\sig{daba varios bostezos})---.
  \item La \termino{intensidad}, que se refiere a que los eventos pueden tener
        diversos grados de fuerza o énfasis en su realización, como, por
        ejemplo, \forma{lloviznar / llover / diluviar} (\sig{llover
        intensamente}) o \forma{picotear} (algo comestible) / \forma{comer /
        devorar}.
\end{itemize}

\subsection{Las clases de eventos}\label{sec:08-clases}

La tradición gramatical adjudica a Aristóteles una primera diferenciación entre
verbos que describen eventos que tienen un punto final, a los que hemos
denominado télicos, y eventos que carecen de término final. En la actualidad, la
clasificación de eventos más reconocida proviene de otro filósofo,
\textcite{vendler1967}, quien distinguió cuatro subclases, a las que luego
\textcite{comrie1976} añadió una quinta. Veamos algunos ejemplos de cada clase
para luego aportar la caracterización de cada una de ellas, a partir de las
propiedades identificadas en §\ref{sec:08-propiedades}:

\begin{itemize}
  \item Estados: \forma{María tiene hambre; Luis es alto}.
  \item Actividades: \forma{Ana juega a las cartas; Corren alrededor de la
        plaza}.
  \item Realizaciones: \forma{Ana cruzó la calle; Comió un pedazo de torta}.
  \item Logros: \forma{La bomba explotó; Encontró lo que buscaba}.
  \item Semelfactivos: \forma{Bostezaba sin parar; Tosía estrepitosamente}.
\end{itemize}

Los \termino{estados} son sucesos durativos [+dur] y, por su carácter estático y
sin cambios [$-$din] ---\forma{ser alto, temer, estar triste}---, se diferencian
de las otras clases, que dan cuenta de situaciones con cambio o dinámicas
[+din], cuyo desarrollo puede implicar fases o momentos diferenciados. Una
característica adicional, relacionada con la falta de dinamismo de estas
predicaciones, es que su mantenimiento no implica ningún esfuerzo deliberado,
por lo que se consideran no agentivas.

Las \termino{actividades} o \termino{procesos} son eventos durativos [+dur] y
dinámicos (implican cambio) [+din] ---\forma{trabajar, correr, leer}---. Al
igual que los estados, su culminación no involucra un fin o \textit{telos}
interno [$-$tél], por lo que se distinguen de las realizaciones [+tél] y los
logros [+tél]. Las actividades no son télicas porque no necesitan llegar a su
culminación para haber tenido lugar o haberse realizado, ya que «se cumplen en
cada una de sus partes» \parencite[62]{garciafernandez1999}. Así, si decimos que
\forma{Ana juega a las cartas} y deja de hacerlo, igualmente puede decirse que
ha jugado a las cartas; en cambio, cuando decimos que \forma{Julia cruza la
calle}, si el evento es interrumpido, no se puede afirmar que cruzó la calle.

Las \termino{realizaciones} son eventos durativos [+dur] y dinámicos [+din], que
tienen una culminación inherente vinculada a la entidad afectada por el evento,
por lo general manifestada por el COD,\footnote{Ampliaremos esta caracterización
en §\ref{sec:08-complementos}.} y, además, se consideran eventos télicos [+tél]
---\forma{leer una carta, escribir un e-mail, construir un puente}---. Como se
mostró con el ejemplo \forma{Julia cruza la calle}, si la realización es
interrumpida, no puede decirse que efectivamente ocurrió.

Los \termino{logros} son eventos dinámicos [+din] y télicos [+tél], como las
realizaciones, pero, como son de cumplimiento instantáneo, se considera que
tienen una duración despreciable o que suceden instantáneamente [$-$dur]
---\forma{llegar, morir, hervir, encontrar}---.

Los \termino{semelfactivos} son eventos instantáneos [$-$dur] y dinámicos
[+din] que, a diferencia de los logros, no tienen \textit{telos} [$-$tél], lo
que permite que sean repetibles. «Usualmente expresan una ocurrencia única de
eventos que suelen ocurrir en secuencias» \parencite[161]{pustejovsky2019}
---\forma{toser, golpear, pestañear}---.

\begin{ejercicio}\label{ej:08-clases}
Lea atentamente el siguiente fragmento y luego resuelva lo pedido:

\begin{consignas}
  \item Según se muestra en los ejemplos resueltos de la tabla, identifique la
        clase eventiva a la que pertenece cada uno de los verbos.
  \item Complete la tabla indicando las características aspectuales de cada
        verbo.
\end{consignas}

\begin{quote}\small
Indiana Jones es un arqueólogo que viaja por el mundo buscando tesoros valiosos
para la humanidad. En \textit{Los cazadores del arca perdida}, que dirigió Steven
Spielberg, Indiana Jones vive una serie de aventuras hasta que finalmente
recupera el arca de la alianza. En sus andanzas, Jones enfrenta rivales a los
que golpea hasta que los derriba o atrapa con su lazo. Sin embargo, este
intrépido personaje tiene una terrible fobia que adquirió en su juventud: las
serpientes. Según aseguran, la saga del popular profesor Jones tiene un nuevo
guionista que ya ha escrito el libro para la próxima película.
\end{quote}
\pendiente{El original dice «Stephen Spilberg»}

\begin{center}\small
\begin{tabular}{@{}lcccl@{}}
  \toprule
  Verbo & Duratividad & Dinamicidad & Telicidad & Clase eventiva \\
  \midrule
  \forma{ser arqueólogo}      & + & $-$ & $-$ & estado \\
  \forma{viajar}              & + & + & $-$ & actividad \\
  \forma{buscar}              & & & & \\
  \forma{dirigir una película} & & & & \\
  \forma{vivir}               & & & & \\
  \forma{recuperar}           & & & & \\
  \forma{enfrentar}           & & & & \\
  \forma{golpear}             & & & & \\
  \forma{derribar}            & & & & \\
  \forma{atrapar}             & & & & \\
  \forma{tener}               & & & & \\
  \forma{adquirir}            & & & & \\
  \forma{asegurar}            & & & & \\
  \forma{escribir el libro}   & & & & \\
  \bottomrule
\end{tabular}
\end{center}
\end{ejercicio}

\section{Factores que intervienen en la delimitación del evento}\label{sec:08-factores}

Un aspecto importante a tener en cuenta en relación con la clasificación de los
eventos es que la información contenida en el lexema verbal puede ser
modificada por la que aportan otros constituyentes de la oración, como el
sujeto, los complementos y modificadores del verbo, así como también el punto de
vista manifestado por el aspecto flexivo o los valores temporales o modales
aportados por la desinencia verbal. Veamos las distintas posibilidades.

\subsection{Los complementos del V}\label{sec:08-complementos}

La presencia del COD resulta esencial para delimitar el evento, como lo muestra
la oposición entre \forma{escribir} (en el sentido de \sig{ser escritor})
\exref{ex:08-gm1} y \forma{escribir una novela} \exref{ex:08-gm2}. El primer
evento, al ser no delimitado, constituye una actividad; en cambio, el segundo,
al estar acotado en su realización al tiempo requerido para la escritura de la
novela, se convierte en una realización.

\begin{exe}
  \ex\label{ex:08-gm1} García Márquez \textit{escribió} durante toda su vida.
  \ex\label{ex:08-gm2} García Márquez \textit{escribió} [la novela \textit{Cien
      años de soledad}].
\end{exe}

Sin embargo, no todo COD delimita el evento: solo lo hacen aquellos que, como
\forma{la novela Cien años de soledad} en \exref{ex:08-gm2}, sirven para medir
el tiempo de desarrollo; en cambio, los objetos que no experimentan un proceso
de cambio o movimiento que los afecte no delimitan al evento. Así,
\exref{ex:08-cuadro} y \exref{ex:08-auto}, en que los objetos \forma{el cuadro} y
\forma{el auto de su padre} no sufren ninguna modificación debido al evento, se
consideran actividades.

\begin{exe}
  \ex\label{ex:08-cuadro} El estudiante \textit{contemplaba} [el cuadro]
      extasiado.
  \ex\label{ex:08-auto} Mariano \textit{manejó} [el auto de su padre] con gran
      precaución.
\end{exe}

Asimismo, tampoco en \exref{ex:08-carne} y \exref{ex:08-relatos} la presencia
del COD es suficiente para delimitar los respectivos predicados. Pero ¿por qué?
¿Qué sucede en estos casos?

\begin{exe}
  \ex\label{ex:08-carne} \textit{Vende} [carne] al exterior.
  \ex\label{ex:08-relatos} García Márquez \textit{escribió} [relatos] a lo largo
      de su vida.
\end{exe}

Los casos como los de \exref{ex:08-carne} y \exref{ex:08-relatos} siguen siendo
actividades porque los SSNN sin determinante, ya sean nombres de materia, como
\forma{carne} en \exref{ex:08-carne}, o plurales de nombres contables, como
\forma{relatos} en \exref{ex:08-relatos}, no producen el efecto de
particularización necesario para delimitar al evento. Para que este efecto se
produzca, el SN que funciona como COD debe estar acotado, ya sea mediante un
determinante, como vimos en \exref{ex:08-gm2}, o mediante un cuantificador,
como se muestra en \exref{ex:08-varias}.

\begin{exe}
  \ex\label{ex:08-varias} García Márquez \textit{escribió} [varias novelas] a lo
      largo de su vida.
\end{exe}

Por otra parte, los verbos de movimiento pueden ser delimitados por un
complemento preposicional que indique el límite de su realización.

\begin{exe}
  \ex Forrest Gump \textit{corrió} [de costa a costa del país / hasta quedar
      agotado].
  \ex \textit{Caminó} [hasta la otra cuadra].
\end{exe}

\begin{ejercicio}\label{ej:08-delimitar}
Transforme las siguientes actividades en realizaciones añadiéndoles un
complemento que las delimite y luego forme oraciones con ellas, según se muestra
en el ejemplo~1.

\begin{exe}
  \ex \textit{Leer}
  \ex \textit{Estudiar}
  \ex \textit{Diseñar}
  \ex \textit{Ir}
  \ex \textit{Construir}
  \ex \textit{Marchar}
\end{exe}

\begin{resuelto}[Respuesta posible]
1. Leyó [una novela de César Aira].
\end{resuelto}
\end{ejercicio}

\subsection{Incidencia del sujeto}\label{sec:08-sujeto}

Así como los complementos del V (por ejemplo, el COD) determinan la clase
aspectual del evento, también el carácter del SN que funciona como sujeto puede
incidir en su delimitación. Veamos cómo.

\begin{exe}
  \ex\label{ex:08-lluvia} {}[La lluvia] \textit{cayó / caía} sobre la ciudad
      durante toda la tarde.
  \ex\label{ex:08-rayo} {}[Un rayo] \textit{cayó / *caía} sobre el auto y lo
      destruyó.
  \ex\label{ex:08-rayos} {}[Rayos] \textit{cayeron / caían} sobre la ciudad
      durante toda la tarde.
\end{exe}

En \exref{ex:08-lluvia}, con un sujeto que refiere a un nombre de materia, como
\forma{la lluvia}, y que, por lo tanto, es incontable,\footnote{Los sustantivos
contables o discontinuos se refieren a objetos discretos y aislables, que no
pueden dividirse sin dejar de ser lo que son ---\forma{mesa, libro, coche}---.
Los sustantivos incontables o continuos tienen una constitución homogénea, por
lo que pueden dividirse sin perder su naturaleza ---\forma{agua, café, arena}---
\parencite[30]{giammatteo2009}.} el verbo \forma{caer} puede interpretarse como
un evento delimitado cuando está en perfecto simple ---\forma{cayó}--- y como no
delimitado y durativo cuando está en imperfecto ---\forma{caía}---. En cambio,
en \exref{ex:08-rayo}, con el sujeto \forma{un rayo}, que es un nombre contable,
el verbo solo admite perfecto simple y la única interpretación posible es la de
que se trata de un logro. Pero ¿qué sucede en \exref{ex:08-rayos}? ¿Por qué
ahora \forma{rayos} es compatible con el pretérito imperfecto? Según acabamos de
ver en relación con el COD, cuando los SSNN en plural no están determinados, el
evento no queda delimitado y admite interpretación durativa, lo que hace a
\forma{rayos} compatible con el pretérito imperfecto (v.~§\ref{sec:08-perfectivos}).

No solo la pluralidad morfológica del sujeto incide en la delimitación del
evento, ya que los nombres colectivos, que manifiestan pluralidad
semántica,\footnote{Los sustantivos individuales, en singular, se refieren a un
solo objeto y, si mencionan más de uno, deben pluralizarse ---\forma{tres libros,
muchos faroles, algunos ajíes}---. Los sustantivos colectivos, en forma
singular, designan varias entidades ---\forma{enjambre, rebaño, docena}---
\parencite[31]{giammatteo2009}.} también producen un efecto similar, de modo que
permiten una interpretación durativa de los eventos \exref{ex:08-multitud}, que
los individuales no admiten \exref{ex:08-profesor}.\pendiente{El original remite
dos veces a (16)}

\begin{exe}
  \ex\label{ex:08-multitud} {}[La multitud] \textit{salió} del estadio
      ordenadamente durante una hora.
  \ex\label{ex:08-profesor} {}[El profesor] \textit{salió} de su casa (*durante
      una hora).
\end{exe}

Asimismo, también el papel temático del sujeto puede alterar la interpretación
eventiva del predicado. Por ejemplo, \forma{rodear} en \exref{ex:08-parque} y
\forma{envolver} en \exref{ex:08-atmosfera}, con sujetos tema, se entienden como
verbos de estado, mientras que en \exref{ex:08-policia} y
\exref{ex:08-vendedor}, con sujeto agente, se interpretan como verbos de acción
del tipo eventivo de realización.

\begin{exe}
  \ex\label{ex:08-parque} {}[El parque] \textit{rodea} la casa.
  \ex\label{ex:08-policia} {}[La policía] \textit{rodeó} la manzana en 20
      segundos.
  \ex\label{ex:08-atmosfera} {}[La atmósfera] \textit{envuelve} a la Tierra.
  \ex\label{ex:08-vendedor} Con mucho cuidado [el vendedor] \textit{envolvió} el
      regalo para la Navidad.
\end{exe}

\begin{ejercicio}\label{ej:08-sujeto}
Compare las siguientes oraciones y explique los contrastes aspectuales que se
dan en cada grupo, según se muestra en el ejemplo resuelto~1.

\begin{exe}
  \ex
    \begin{xlist}
      \ex La luz del sol \textit{rebotó / rebotaba} en la pared toda la tarde.
      \ex La pelota \textit{rebotó / rebotaba} en la pared toda la tarde.
      \ex La bala \textit{rebotó} en la pared (*toda la tarde).
    \end{xlist}
  \ex
    \begin{xlist}
      \ex El soldado \textit{murió} durante la emboscada.
      \ex Juan \textit{murió} (*durante toda la tarde).
    \end{xlist}
  \ex
    \begin{xlist}
      \ex El agua líquida \textit{cubre} aproximadamente tres cuartas partes de
          la superficie terrestre.
      \ex Pedro \textit{cubrió} el auto para que el granizo no lo estropee.
    \end{xlist}
\end{exe}

\begin{resuelto}
1a. Con un sujeto que refiere a un nombre de materia incontable ---\forma{la luz
del sol}---, el verbo \forma{rebotar} en perfecto simple ---\forma{rebotó}--- se
interpreta como un evento delimitado, pero en imperfecto ---\forma{rebotaba}---
se considera no delimitado y durativo. 1b. Con el sujeto \forma{la pelota}, que
es un nombre común y contable, el verbo admite perfecto simple, con
interpretación de evento delimitado y puntual (= \sig{dio un rebote}); pero en
imperfecto el evento se considera reiterativo. 1c. Con el sujeto \forma{la
bala}, \forma{rebotar} se vuelve un evento semelfactivo, único e irrepetible,
por lo que rechaza el modificador durativo \forma{toda la tarde}.
\end{resuelto}
\end{ejercicio}

\subsection{Compatibilidad con distintas clases de modificadores temporales (MT)}\label{sec:08-mt}

Si comparamos los modificadores temporales (MMTT) de las siguientes oraciones,
vemos que en \exref{ex:08-estreno} y \exref{ex:08-regreso} refieren a un momento
de tiempo determinado respecto del que el evento se sitúa ---\forma{la semana
pasada, en 1994, a las nueve de la noche}---. En \exref{ex:08-duro} y
\exref{ex:08-tarea}, en cambio, el tiempo es concebido en su transcurrir
---\forma{dos horas, toda la tarde}---, sin atender a su culminación; así, por
ejemplo, \exref{ex:08-tarea-a} acepta continuarse con \forma{\ldots{}pero no la
terminamos}, como se muestra en \exref{ex:08-tarea-b}. Los modificadores
temporales de \exref{ex:08-estreno} y \exref{ex:08-regreso} se denominan
\termino{puntuales} (MTP), mientras que los de \exref{ex:08-duro} y
\exref{ex:08-tarea} son modificadores \termino{de duración} (MTD).

\begin{exe}
  \ex\label{ex:08-estreno} La película \textit{se estrenó} [la semana pasada / en
      1994].
  \ex\label{ex:08-regreso} \textit{Regresó} del cine [a las nueve de la noche].
  \ex\label{ex:08-duro} La proyección \textit{duró} [dos horas].
  \ex\label{ex:08-tarea}
    \begin{xlist}
      \ex\label{ex:08-tarea-a} \textit{Estuvimos haciendo} la tarea [toda la
          tarde].
      \ex\label{ex:08-tarea-b} \textit{Estuvimos haciendo} la tarea [toda la
          tarde], pero no la terminamos.
    \end{xlist}
\end{exe}

Por su parte, los modificadores que aparecen en \exref{ex:08-guion} y
\exref{ex:08-manejar}, si bien implican duración, delimitan un período de tiempo
y presentan al suceso como concluido. Se conocen como modificadores temporales
\termino{de período} (MTPer) y aceptan una paráfrasis o equivalencia con
\forma{tardar} + SN cuantificado + \forma{en} + infinitivo
\exref{ex:08-guion-b} o con \forma{llevarle (a alguien)} + SN cuantificado +
infinitivo \exref{ex:08-guion-c}.

\begin{exe}
  \ex\label{ex:08-guion}
    \begin{xlist}
      \ex \textit{Terminó} el guion de la película [en seis meses].
      \ex\label{ex:08-guion-b} \textit{Tardó} seis meses en terminar el trabajo.
      \ex\label{ex:08-guion-c} Le \textit{llevó} seis meses escribir el guion de
          la película.
    \end{xlist}
  \ex\label{ex:08-manejar} Aprendió a manejar [en un mes].
\end{exe}

Según veremos, las distintas subclases eventivas de verbos muestran afinidades e
incompatibilidades con los tres tipos de modificadores temporales mencionados.
Sin embargo, como sostiene \textcite[3047]{demiguel1999}, «al tratar de
información aspectual no se puede hablar de incompatibilidades o restricciones
auténticas, sino más bien de preferencias o tendencias», ya que, según veremos
en §\ref{sec:08-coercion}, las distintas subclases pueden cambiar su tipo
eventivo para compatibilizar con los modificadores u otros constituyentes
oracionales.

Las actividades [+din, +dur, $-$tél] y las realizaciones [+din, +dur, +tél]
presentan un comportamiento complementario respecto de los MTD y los MTPer: las
primeras típicamente se combinan con MTD \exref{ex:08-forrest}, mientras que las
segundas preferentemente seleccionan MTPer \exref{ex:08-recaudo}:

\begin{exe}
  \ex\label{ex:08-forrest} Forrest Gump \textit{corrió} de costa a costa (durante
      tres años / por mucho tiempo / *en una semana).
  \ex\label{ex:08-recaudo} La película \textit{recaudó} millones de dólares (en
      una semana / ?durante tres años / *por mucho tiempo).
\end{exe}

En cuanto a los estados [$-$din, +dur, $-$tél], una de sus características es su
duración potencialmente ilimitada, que impide que, por exceso, puedan emplearse
con MTD que los acoten \exref{ex:08-camarones}. Sin embargo, los estados
potencialmente finitos o transitorios son compatibles con MTD
\exref{ex:08-dan}:

\begin{exe}
  \ex[*]{\label{ex:08-camarones} Los camarones \textit{son} crustáceos [durante
      este año].}
  \ex\label{ex:08-dan} El teniente Dan \textit{era} el líder del pelotón de
      Forrest Gump [durante la guerra de Vietnam].
\end{exe}

\begin{recuadro}[Para ampliar]
En la bibliografía se suele explicar la diferencia entre estados de duración
ilimitada y estados finitos mediante la distinción que introduce
\textcite{carlson1977} entre \textit{individual-level predicates} (predicados de
nivel individual) y \textit{stage-level predicates} (predicados de estadio). Los
primeros son aquellos que persisten en el tiempo y pueden caracterizar una
entidad, como \forma{ser alto, tener ojos verdes} o \forma{ser cordobés}. Los
segundos son aquellos que pueden cambiar con el tiempo, como \forma{estar
enfermo, tener fiebre} o \forma{estar contento}. En español, esta distinción
entre \termino{estados de nivel individual} y \termino{estados de estadio}, en
líneas generales, se puede correlacionar con el uso de los verbos \forma{ser},
para los primeros ---\forma{Guillermo es enfermo desde su juventud}---, y
\forma{estar}, para los segundos ---\forma{Guillermo estuvo enfermo la semana
pasada}---.\footnote{Si bien la distinción entre \forma{ser} y \forma{estar}
tiene una base aspectual, existen varias excepciones en las que no podemos
extendernos aquí. Para ampliar, remitimos a \textcite[313 y
ss.]{bosque2009}.}
\end{recuadro}

Los logros [+din, $-$dur, +tél],\pendiente{El original dice aquí
«[+din, --dur, --tél]»} por su carácter no durativo, aceptan MTP referidos a un
momento muy breve \exref{ex:08-valla} o a un tiempo determinado
\exref{ex:08-1994}, mientras que los estados los rechazan
\exref{ex:08-ocho}.\pendiente{El original habla aquí de «MTPer»}

\begin{exe}
  \ex\label{ex:08-valla} Los refugiados \textit{saltaron} la valla [en un
      segundo].
  \ex\label{ex:08-1994} La película \textit{se estrenó} [en 1994].
  \ex[*]{\label{ex:08-ocho} Los camarones \textit{viven} en el mar [a las ocho].}
\end{exe}

Los semelfactivos típicamente se combinan con modificadores de repetición
(MRep) \exref{ex:08-tosio} y de frecuencia (MFr) \exref{ex:08-bostezaba}, pero
este tipo de modificadores son rechazados por los predicados de estado que, al
no tener una duración acotada, no admiten repetición ---\exref{ex:08-portero} y
\exref{ex:08-enferma-a}---, excepto que se trate de estados transitorios
\exref{ex:08-enferma-b}:

\begin{exe}
  \ex\label{ex:08-tosio} \textit{Tosió} [dos veces seguidas].
  \ex\label{ex:08-bostezaba} \textit{Bostezaba} [cada diez segundos].
  \ex[*]{\label{ex:08-portero} El portero de la escuela \textit{pertenece /
      perteneció} a una familia muy humilde [habitualmente / tres veces / en
      varias ocasiones].}
  \ex
    \begin{xlist}
      \ex[*]{\label{ex:08-enferma-a} [Habitualmente / Cada dos por tres] mi
          hermana \textit{es enferma} en invierno.}
      \ex\label{ex:08-enferma-b} {}[Habitualmente / Cada dos por tres] mi hermana
          \textit{está enferma} en invierno.
    \end{xlist}
\end{exe}

\subsection{La reinterpretación de la clase eventiva}\label{sec:08-coercion}

Cuando un verbo de una subclase eventiva se combina con un modificador que exige
una reinterpretación de su significado, se produce una «acomodación», para la
que \textcite{pustejovsky1995} utiliza la denominación de
\termino{coerción} o cambio del tipo eventivo.

\begin{destacado}
\textcite[330]{bosque2009} definen la \termino{coerción} o \termino{coacción}
como la «[o]peración léxica que cambia el tipo aspectual (una o más propiedades
aspectuales) de una pieza léxica. Sin este cambio la combinación resultante
sería no convergente».
\end{destacado}

Algunos de estos casos de coerción se producen:

\begin{itemize}
  \item cuando las realizaciones, que son télicas [+tél], se usan con MTD, por
        lo que estos modificadores pasan a referirse al desarrollo del evento
        \exref{ex:08-escribio}, mientras que si las actividades [$-$tél] se
        combinan con MTPer, estos se acomodan a hacer referencia al proceso o
        período de realización previo a la culminación del evento
        \exref{ex:08-vistio};
\end{itemize}

\begin{exe}
  \ex\label{ex:08-escribio} \textit{Escribió} el guion [durante un año].
  \ex\label{ex:08-vistio} \textit{Se vistió} [en quince minutos] y salió a la
      calle.
\end{exe}

\begin{itemize}
  \item cuando los logros se combinan con MTD, estos pueden hacer referencia al
        proceso previo \exref{ex:08-hirvio} o al estado resultante del evento
        \exref{ex:08-sefue}, o bien pueden tomar valor habitual
        \exref{ex:08-fuegos};
\end{itemize}

\begin{exe}
  \ex\label{ex:08-hirvio} La leche \textit{hirvió} [durante media hora].
  \ex\label{ex:08-sefue} \textit{Se fue} [por diez minutos].
  \ex\label{ex:08-fuegos} [Durante la semana de festejos] \textit{se lanzaron}
      fuegos artificiales.
\end{exe}

\begin{itemize}
  \item cuando logros y semelfactivos se combinan con MTPer o con MTD adquieren
        interpretación iterativa o habitual.
\end{itemize}

\begin{exe}
  \ex En China, Forrest \textit{ganó} varios partidos de pimpón [en una semana].
  \ex Un espectador de la tercera fila \textit{bostezó} sin parar [durante toda
      la película].
\end{exe}

\begin{ejercicio}\label{ej:08-construir}
Construya oraciones según lo pedido y luego indique la interpretación
resultante, según se muestra en los ejemplos resueltos de abajo.

\begin{exe}
  \ex Un verbo de estado con un MTD.
  \ex Un verbo de logro con un MTD.
  \ex Un verbo de realización con un MTPer.
  \ex Un verbo de estado transitorio con un MRep.
  \ex Un verbo de actividad con un MTD.
  \ex Un verbo semelfactivo con un MRep.
  \ex Un verbo de realización con un MTD.
  \ex Un verbo de logro con un MTPer.
\end{exe}

\begin{resuelto}[Respuestas posibles]
1. Los dinosaurios \textit{han existido} [por muchos millones de años]. El MTD
refuerza el valor durativo del estado. 2. \textit{Encontraron} muchos fósiles
[durante la excavación]. El MTD fuerza una interpretación durativa de repetición
(= \sig{hubo muchos hallazgos}) del logro.
\end{resuelto}
\end{ejercicio}

\subsection{Compatibilidad de los estados con otras clases de modificadores verbales}\label{sec:08-estados}

Dado su carácter homogéneo y que no indican cambio interno en su desarrollo, los
estados no son compatibles con modificadores de dinamicidad, como los de
manera, del tipo de \forma{rápidamente, con alegría, sin miedo}, ni con los
modificadores aspectuales que indican grado de completamiento del evento, como
\forma{poco a poco, de a poco, gradualmente, hasta la mitad}.

\begin{exe}
  \ex[*]{Las ballenas \textit{son} cetáceos [rápidamente / sin miedo].}
  \ex[*]{Indiana Jones \textit{teme} a las víboras [poco a poco / hasta la
      mitad].}
\end{exe}

Como los estados carecen de sujeto agente, tampoco son compatibles con los
adverbios orientados al sujeto, que describen el modo o manera en que el agente
realiza la acción ---\forma{cuidadosamente, esforzadamente,
inconscientemente, inteligentemente}, etc.---.

\begin{exe}
  \ex[*]{Los triángulos \textit{tienen} tres lados y tres ángulos
      [inteligentemente / deliberadamente].}
  \ex[*]{Algunos alumnos \textit{saben} inglés [cuidadosamente].}
\end{exe}

\subsection{Interacción del tipo eventivo (aspecto léxico) con el aspecto y el modo gramatical}\label{sec:08-interaccion}

\subsubsection{Tiempos perfectivos e imperfectivos}\label{sec:08-perfectivos}

Los \termino{tiempos imperfectivos} resaltan el carácter durativo y no
conclusivo de los eventos no télicos, como los estados \exref{ex:08-josefa} y
las actividades \exref{ex:08-mario}.

\begin{exe}
  \ex\label{ex:08-josefa} Josefa \textit{vivía} en Nueva York durante su
      juventud.
  \ex\label{ex:08-mario} Mario \textit{trabajaba} en una carpintería de
      Avellaneda.
\end{exe}

En cuanto a los eventos télicos ---realizaciones y logros---, la llamada
\termino{paradoja imperfectiva} hace referencia a que, cuando están en un tiempo
imperfectivo, no implican una culminación, \exref{ex:08-leia} y
\exref{ex:08-llegaba}, a diferencia de lo que sucede con los estados y
actividades, \exref{ex:08-vivian} y \exref{ex:08-corria}, que implican que el
suceso efectivamente ocurrió.

\begin{exe}
  \ex\label{ex:08-leia} Juan \textit{leía} una novela de misterio $\not\Rightarrow$
      Juan \textit{leyó} una novela de misterio.
  \ex\label{ex:08-llegaba} El doctor López \textit{llegaba} a su casa
      $\not\Rightarrow$ El doctor López \textit{llegó} a su casa.
  \ex\label{ex:08-vivian} Mis amigos \textit{vivían} en Hurlingham $\Rightarrow$
      Mis amigos \textit{vivieron} en Hurlingham.
  \ex\label{ex:08-corria} Forrest \textit{corría} por las rutas $\Rightarrow$
      Forrest \textit{corrió} por las rutas.
\end{exe}

Los tiempos imperfectivos son los preferidos por cualquier clase eventiva para
expresar hechos habituales o repetitivos.

\begin{exe}
  \ex \textit{Viajaba} todos los fines de semana al interior del país.
  \ex Estaba tan nervioso que \textit{entraba} y \textit{salía} continuamente.
\end{exe}

Por su parte, los \termino{tiempos perfectivos}, como el perfecto simple, se
emplean con cualquier clase de evento para poner de relieve su terminación.

\begin{exe}
  \ex Mi jefe \textit{leyó} atentamente el informe. (realización)
  \ex Marina \textit{entró} cansada a la oficina. (logro)
  \ex Forrest \textit{sufrió} mucho en su infancia. (estado)
  \ex El sol \textit{brilló} durante toda la semana. (actividad)
\end{exe}

También prefieren los tiempos perfectivos los eventos únicos:

\begin{exe}
  \ex Noam Chomsky \textit{nació} el 7 de diciembre de 1928.
\end{exe}

\subsubsection{La perífrasis progresiva}\label{sec:08-progresiva}

Dado que las perífrasis progresivas\footnote{Las perífrasis verbales se tratan
en el \capref{11}.} indican «el proceso de cambio asociado a una eventualidad»,
resultan incompatibles con los estados, que se caracterizan por no cambiar o no
avanzar en el tiempo \parencite[301]{bosque2009} ---\exref{ex:08-triangulo} y
\exref{ex:08-tigres}---. Sin embargo, los estados potencialmente delimitables
permiten la inclusión de un participante que se esfuerza en el sostenimiento de
la situación, que se presenta como más dinámica y compatible con un tiempo
expresado por una perífrasis progresiva ---\exref{ex:08-blanda} y
\exref{ex:08-gustando}---:

\begin{exe}
  \ex[*]{\label{ex:08-triangulo} Este triángulo \textit{está teniendo} tres
      lados.}
  \ex[*]{\label{ex:08-tigres} Los tigres \textit{están siendo} felinos.}
  \ex\label{ex:08-blanda} Marina \textit{está siendo} muy blanda con su hija.
  \ex\label{ex:08-gustando} Me \textit{está gustando} mucho la novela que me
      prestaste.
\end{exe}

Los logros tampoco aceptan los tiempos progresivos porque, al ser instantáneos,
no pueden indicar un desarrollo temporal del evento. Sin embargo, como ya hemos
señalado, admiten una referencia tanto al tiempo previo \exref{ex:08-resolviendo}
como al posterior \exref{ex:08-hirviendo}:

\begin{exe}
  \ex\label{ex:08-resolviendo} \textit{Estuvo resolviendo} el ejercicio un buen
      rato. (= \sig{Estuvo haciéndolo un buen rato hasta que lo resolvió})
  \ex\label{ex:08-hirviendo} La mezcla \textit{estuvo hirviendo} media hora más,
      como lo indicaba la receta. (= \sig{Continuó hirviendo media hora después
      del momento de ebullición})
\end{exe}

\subsubsection{Perífrasis fasales}\label{sec:08-fasales}

Las perífrasis fasales focalizan distintos momentos del evento: ya sea el inicio
---\forma{empezar / comenzar a}---, el fin ---\forma{dejar / acabar / terminar
de}--- o el desarrollo ---\forma{continuar / seguir}--- y, para hacerlo, exigen
que el verbo principal sea durativo y tenga delimitación. En consecuencia, los
verbos de logro no pueden emplearse con estas perífrasis \exref{ex:08-llegar},
excepto que tengan interpretación habitual \exref{ex:08-coleccionar}. Por la
misma razón interpretativa, los semelfactivos, que indican repetición del
evento, son posibles con estas perífrasis \exref{ex:08-bostezar}.

\begin{exe}
  \ex[*]{\label{ex:08-llegar} \textit{Empezó a / terminó de llegar / continuó
      llegando} toda la tarde.}
  \ex\label{ex:08-coleccionar} \textit{Empezó a / dejó de coleccionar / continuó
      coleccionando} objetos extraños.
  \ex\label{ex:08-bostezar} \textit{Empezó a / dejó de bostezar} enseguida /
      \textit{continuó bostezando} durante toda la función.
\end{exe}

Por otra parte, dado que los estados suelen carecer de delimitación, no pueden
utilizarse ni con las perífrasis inceptivas ---\forma{empezar / comenzar a}---
\exref{ex:08-salon} ni con las culminativas ---\forma{terminar de, acabar de}---
\exref{ex:08-contenta-a}, aunque admiten la perífrasis cesativa con \forma{dejar
de} \exref{ex:08-contenta-b}. La diferencia, para \textcite[95]{cinque1999},
estriba en que el \termino{aspecto cesativo} caracteriza una situación que ha
quedado inconclusa. Así, por ejemplo, en \forma{dejó de lavar los platos}, el
evento se interrumpe sin que necesariamente se haya completado, es decir, sin
que se hayan terminado de lavar todos los platos. En cambio, en el
\termino{aspecto culminativo}, el evento sí alcanza su punto final natural. Así,
en \forma{terminó de lavar los platos}, se interpreta que el evento culmina
cuando no quedan más platos por lavar.

\begin{exe}
  \ex[*]{\label{ex:08-salon} El salón \textit{empezó / comenzó a medir} 20
      metros.}
  \ex
    \begin{xlist}
      \ex[*]{\label{ex:08-contenta-a} \textit{Terminó de estar contenta} cuando
          supo la verdad.}
      \ex\label{ex:08-contenta-b} \textit{Dejó de estar contenta} cuando supo la
          verdad.
    \end{xlist}
\end{exe}

El hecho de que los estados no presenten fases de desenvolvimiento no favorece
la utilización de estos predicados con las perífrasis inceptivas; sin embargo,
no lo excluye totalmente \exref{ex:08-corpulento}. Además, los estados
transitorios admiten esta combinación sin problemas \exref{ex:08-recuperar}.

\begin{exe}
  \ex\label{ex:08-corpulento} \textit{Empezó a ser corpulento} a los doce años,
      cuando se desarrolló.
  \ex\label{ex:08-recuperar} \textit{Empezó a estar contenta} cuando supo que se
      iba a recuperar de su enfermedad.
\end{exe}

\subsubsection{Modo imperativo}\label{sec:08-imperativo}

La no agentividad de los estados, que deriva de su carácter estable y homogéneo
(sin fases), tiene como consecuencia que no se empleen en imperativo,
\exref{ex:08-ojos} y \exref{ex:08-alto}, excepto cuando representan un
comportamiento que puede concebirse como dependiente de la voluntad,
\exref{ex:08-caritativo} y \exref{ex:08-quieto}.

\begin{exe}
  \ex[*]{\label{ex:08-ojos} \textit{Tené} ojos claros.}
  \ex[*]{\label{ex:08-alto} \textit{Sé} alto.}
  \ex\label{ex:08-caritativo} \textit{Sé} caritativo y ayudá a quien lo necesita.
  \ex\label{ex:08-quieto} \textit{Quedate} quieto.
\end{exe}

\begin{ejercicio}\label{ej:08-naufrago}
Lea atentamente el siguiente fragmento adaptado de \textit{Relato de un
náufrago}, de G.~García Márquez (1970), y luego complete la tabla de abajo con
lo pedido, según se muestra en los ejemplos ya resueltos.

% Fragmento literario protegido: se incorpora a mano desde el original.
\textoaincorporar{el fragmento adaptado de la presentación de \textit{Relato de
un náufrago}, desde «El 28 de febrero de 1955 se conoció la noticia» hasta
«\ldots{}sus diez días en el mar», con los predicados que figuran en la tabla}

\begin{consignas}
  \item Determine la clase eventiva de todos los predicados del fragmento.
  \item Identifique, si corresponde, los modificadores temporales que reciben.
  \item Reconozca el aspecto gramatical de cada verbo y determine su incidencia
        en cada predicación.
  \item En todos los casos, especifique si la lectura o interpretación
        resultante es puntual o durativa.
  \item Determine si hay compatibilidad o coerción (cambio de tipo eventivo)
        entre los distintos factores considerados.
\end{consignas}

{\footnotesize
\begin{tabularx}{\linewidth}{@{}>{\raggedright\arraybackslash}p{19mm}
    >{\raggedright\arraybackslash}p{13mm} L L
    >{\raggedright\arraybackslash}p{10mm} L@{}}
  \toprule
  Predicado & Clase eventiva & MT & Aspecto verbal & Lectura &
    Compatibilidad / coerción \\
  \midrule
  \forma{se conoció} & estado & MTP \forma{el 28 de febrero de 1955} &
    perfectivo: presenta al evento como concluido & puntual & coerción del tipo
    eventivo por el MT y el aspecto verbal \\
  \forma{habían caído y desaparecido} & logro & --- & perfectivo: refuerza el
    valor conclusivo & puntual & compatibilidad entre tipo eventivo y aspecto
    verbal \\
  \forma{viajaba} & & & & & \\
  \forma{había sido sometida a reparaciones} & & & & & \\
  \forma{llegó} & & & & & \\
  \forma{se inició} & & & & & \\
  \forma{hacen oficios} & & & & & \\
  \forma{se desistió} & & & & & \\
  \forma{fueron declarados} & & & & & \\
  \forma{apareció} & & & & & \\
  \forma{había permanecido} & & & & & \\
  \forma{se llamaba} & & & & & \\
  \forma{es} & & & & & \\
  \forma{contó} & & & & & \\
  \forma{fue publicada} & & & & & \\
  \forma{tomaba notas} y \forma{soltaba preguntas} & & & & & \\
  \forma{pudimos reconstruir el relato} & & & & & \\
  \bottomrule
\end{tabularx}}
\end{ejercicio}

\begin{recuadro}[Reflexión final]
La estructura eventiva, dependiente tanto del significado léxico del verbo como
de la categoría de aspecto, se muestra como un fenómeno esencialmente
combinatorio y con múltiples aristas. Es así que, para todo predicado, hay que
tener en cuenta no solo la naturaleza del verbo, es decir, la clase aspectual a
la que pertenece (estado, actividad, realización, logro o semelfactivo), sino
también la estructura sintáctica que conforma con los distintos modificadores
que recibe, incluido el sujeto, y los distintos tiempos y formas verbales en que
se manifiesta el verbo. Si bien se trata de una cuestión compleja, su
reconocimiento resulta crucial para identificar, a partir de la naturaleza del
verbo y la configuración morfosintáctica en que se manifiesta, cuál será la
lectura o interpretación semántica que corresponde al predicado.
\end{recuadro}

% --- Secciones finales -------------------------------------------------
\lecturasrecomendadas

\solucionario

\begin{solucion}{\ref{ej:08-clases}}
\begin{center}\small
\begin{tabular}{@{}lcccl@{}}
  \toprule
  Verbo & Duratividad & Dinamicidad & Telicidad & Clase eventiva \\
  \midrule
  \forma{buscar}                         & + & + & $-$ & actividad \\
  \forma{dirigir una película}           & + & + & + & realización \\
  \forma{vivir una serie de aventuras}   & + & $-$ & +\textsuperscript{a} & actividad \\
  \forma{recuperar}                      & $-$ & + & + & logro \\
  \forma{enfrentar}                      & + & + & $-$ & actividad \\
  \forma{golpear}                        & $-$\textsuperscript{b} & + & $-$ & semelfactivo \\
  \forma{derribar}                       & $-$ & + & + & logro \\
  \forma{atrapar}                        & $-$ & + & + & logro \\
  \forma{tener}                          & + & $-$ & $-$ & estado \\
  \forma{adquirir}                       & $-$ & + & + & logro \\
  \forma{asegurar}                       & $-$ & + & + & logro \\
  \forma{escribir el libro}              & + & + & + & realización \\
  \bottomrule
\end{tabular}
\end{center}
\textsuperscript{a} El evento \forma{vivir una serie de aventuras} está
delimitado por el modificador temporal encabezado por \forma{hasta}.
\textsuperscript{b} Este evento se vuelve durativo por la repetición.
\pendiente{Los rasgos de \textit{vivir una serie de aventuras} ($-$din, +tél) no
corresponden a una actividad}
\end{solucion}

\begin{solucion}{\ref{ej:08-delimitar}}
Ejemplos de respuestas posibles. 2. Estudiamos [la voz pasiva] para rendir el
examen. 3. Diseñaron [un nuevo espacio verde] para la ciudad. 4. Fuimos [hasta
la plaza]. 5. ¿Ya construyeron [la nueva autopista]? 6. Marchó [hasta la
frontera].
\end{solucion}

\begin{solucion}{\ref{ej:08-sujeto}}
2a. El verbo de logro \forma{morir} es compatible con \forma{durante la
emboscada}, que limita la ocurrencia del evento a un período de tiempo. 2b.
\forma{Durante toda la tarde} indica duración y es incompatible con un verbo de
logro, puntual y delimitado, como \forma{morir}. 3. En 3a, con sujeto tema,
\forma{cubrir} se entiende como un verbo de estado, mientras que en 3b, con
sujeto agente, se interpreta como verbo de acción, del tipo eventivo de
realización.
\end{solucion}

\begin{solucion}{\ref{ej:08-construir}}
Ejemplos de respuestas posibles. 3. \textit{Planchó} toda la ropa [en una hora].
El MTPer mide la duración de la realización. 4. \textit{Estuvo} [varias veces] en
la ciudad. El MRep refiere a la reiteración de un estado transitorio. 5. Todas
las mañanas \textit{corre} [media hora] alrededor de la plaza. El MTD mide la
duración del evento de actividad. 6. \textit{Tosió} [tres o cuatro veces] durante
la conferencia. El MRep indica la repetición del evento semelfactivo. 7. Los
turistas \textit{recorrieron} la ciudad [durante la tarde]. El MTD mide el
tiempo del evento de realización. 8. \textit{Llegó} de su casa a la Facultad [en
dos horas]. El MTPer mide el tiempo previo al evento.
\end{solucion}

\begin{solucion}{\ref{ej:08-naufrago}}
{\footnotesize\setlength{\tabcolsep}{3pt}
\begin{longtable}{@{}>{\raggedright\arraybackslash}p{18mm}
    >{\raggedright\arraybackslash}p{14mm} >{\raggedright\arraybackslash}p{20mm}
    >{\raggedright\arraybackslash}p{20mm} >{\raggedright\arraybackslash}p{11mm}
    >{\raggedright\arraybackslash}p{32mm}@{}}
  \toprule
  Predicado & Clase eventiva & MT & Aspecto verbal & Lectura &
    Compatibilidad / coerción \\
  \midrule
  \endhead
  \forma{viajaba} & actividad & --- & imperfectivo: refuerza el valor durativo &
    durativa & compatibilidad entre tipo eventivo y aspecto verbal \\
  \forma{había sido sometida a reparaciones} & realización & --- &
    perfectivo\textsuperscript{a} & durativa & compatibilidad entre tipo eventivo
    y aspecto verbal \\
  \forma{llegó} & logro & MTP \forma{dos horas después} & perfectivo: refuerza
    el valor conclusivo & puntual & compatibilidad entre tipo eventivo, MTP y
    aspecto verbal \\
  \forma{se inició} & logro & MTP \forma{de inmediato} & perfectivo: refuerza el
    valor conclusivo & puntual & compatibilidad entre tipo eventivo, MTP y
    aspecto verbal \\
  \forma{hacen oficios} & actividad & --- & imperfectivo: refuerza el valor
    durativo & habitual & compatibilidad entre tipo eventivo y aspecto verbal;
    incidencia del COD plural en la lectura habitual \\
  \forma{se desistió} & logro & MTP \forma{al cabo de cuatro días} & perfectivo:
    refuerza el valor conclusivo & puntual & compatibilidad entre tipo eventivo,
    MTP y aspecto verbal \\
  \forma{fueron declarados} & logro & --- & perfectivo: refuerza el valor
    conclusivo & puntual & compatibilidad entre tipo eventivo y aspecto verbal \\
  \forma{apareció} & logro & MTP \forma{una semana más tarde} & perfectivo:
    refuerza el valor conclusivo & puntual & compatibilidad entre tipo eventivo,
    MTP y aspecto verbal \\
  \forma{había permanecido} & estado & MTD \forma{diez días} & perfectivo:
    presenta al evento como concluido & durativa & la compatibilidad entre tipo
    eventivo y MTD produce la lectura durativa; el aspecto perfecto señala
    anterioridad y conclusividad del evento, pero no coerciona el tipo eventivo
    del verbo \\
  \forma{se llamaba} & estado & --- & imperfectivo: refuerza el valor durativo &
    durativa & compatibilidad entre tipo eventivo y aspecto verbal \\
  \forma{es} & estado & --- & imperfectivo: refuerza el valor durativo & durativa
    & compatibilidad entre tipo eventivo y aspecto verbal \\
  \forma{contó} & realización (delimitada por el COD \forma{lo}) & --- &
    perfectivo: refuerza el valor conclusivo & durativa & compatibilidad entre
    tipo eventivo y aspecto verbal \\
  \forma{fue publicada} & realización & MTP \forma{un mes después} & perfectivo:
    refuerza el valor conclusivo & puntual & compatibilidad entre tipo eventivo,
    MTP y aspecto verbal \\
  \forma{tomaba notas} y \forma{soltaba preguntas} & logro & MTD \forma{durante
    las cuales} (veinte sesiones de seis horas diarias) & imperfectivo: refuerza
    el valor durativo & durativa & coerción del tipo eventivo de logro por el MTD
    y el aspecto verbal, que producen la lectura durativa \\
  \forma{pudimos reconstruir el relato} & realización & MTPer \forma{en veinte
    sesiones de seis horas diarias} & perfectivo: refuerza el valor conclusivo &
    durativa & la compatibilidad entre tipo eventivo y MTPer produce la lectura
    durativa; el tiempo perfecto destaca la conclusión del evento, pero no
    coerciona el tipo eventivo del verbo \\
  \bottomrule
\end{longtable}}

\textsuperscript{a} El sustantivo plural \forma{reparaciones} no permite
concluir la delimitación del evento.
\end{solucion}

% =====================================================================
%  capitulos/cap09.tex — Capítulo 9. Los predicativos
%  Autores: Pablo Zdrojewski y Fernando Carranza
%  Fuente: Capitulo_09_normalizado.docx
% =====================================================================
\capitulo{Los predicativos}{Pablo Zdrojewski y Fernando Carranza}

\section{Introducción}

El \termino{predicativo} es uno de los modificadores del núcleo verbal más
difíciles de caracterizar porque, en primer lugar, no está restringido
categorialmente, es decir, su núcleo puede pertenecer a más de una clase de
palabras y, en consecuencia, puede formar distintos tipos de sintagmas: SN, SA,
SP e, incluso, SAdv.

\begin{exe}
  \ex\label{ex:09-cat1}
    \begin{xlist}
      \ex\label{ex:09-cat1-a} Pedro es [\etq{SN} un arquitecto muy renombrado].
      \ex\label{ex:09-cat1-b} Varias torres de cristal se alzan [\etq{SA}
          gallardas] en Buenos Aires.
      \ex La abogada ofició [\etq{SP} de defensora].
      \ex La cosa es [\etq{SAdv} así].
    \end{xlist}
  \ex\label{ex:09-cat2}
    \begin{xlist}
      \ex\label{ex:09-cat2-a} Denominaron a este espécimen [\etq{SN}
          \textit{Patagosaurus fariasi}].
      \ex\label{ex:09-cat2-b} Guardé la ropa [\etq{SA} muy mojada].
      \ex Que no te traten [\etq{SP} de tonto].
      \ex A Juan lo vi [\etq{SAdv} bien].
    \end{xlist}
\end{exe}

Además, tal como iremos viendo, la relación que establece el predicativo con el
núcleo del predicado varía según el tipo de verbo de que se trate. Esto lleva a
que también cambie el tipo de contribución que hace el predicativo a la sintaxis
y a la semántica de las oraciones en las que aparece. En este capítulo vamos a
abordar los distintos tipos de predicativos que se reconocen en el español. En
primer lugar, partimos de una caracterización en relación con sus propiedades
sintácticas y semánticas e introducimos las nociones de \termino{predicación
primaria} y \termino{predicación secundaria}. Luego, en segundo lugar,
distinguimos las clases de predicativos en función de su orientación hacia el
sujeto o el objeto de la oración, y trazamos la diferencia entre obligatorios y
no obligatorios. Por último, señalamos cómo se marcan sintácticamente y
presentamos casos de homonimia estructural a que dan lugar.

\section{Los predicativos y la predicación}\label{sec:09-predicacion}

\subsection{Predicación semántica y predicación sintáctica}

La denominación \forma{predicativo} alude a que se trata de un constituyente que
expresa un tipo particular de predicación.

\begin{destacado}
Es importante tener en cuenta que, aunque están en el predicado, ninguno de los
tipos de predicativo que veremos en este capítulo se considera un argumento del
verbo; en consecuencia, tampoco les corresponde ningún papel temático, ya que
integran con el verbo la predicación principal o forman una predicación
secundaria.
\end{destacado}

Según ya se ha planteado en el \capref{7}, en términos semánticos, podemos
caracterizar la noción de predicación como una relación entre un predicado y sus
argumentos. Los argumentos denotan entidades del mundo extralingüístico,
mientras que los predicados atribuyen propiedades a los argumentos o denotan
relaciones entre argumentos. En este sentido, todas las clases de palabras y de
sintagmas pueden ser \termino{predicados semánticos}, como se aprecia en los
ejemplos de \exref{ex:09-clases}:

\begin{exe}
  \ex\label{ex:09-clases}
    \begin{xlist}
      \ex Juan [canta / construye casas]. \hfill [Verbo]
      \ex Juan [es cantante / arquitecto]. \hfill [Sustantivo]
      \ex Juan [es afinado / prolijo]. \hfill [Adjetivo]
      \ex Juan [está arriba / afuera / acá / bien]. \hfill [Adverbio]
      \ex Enzo [es de River / de Córdoba]. \hfill [Preposición]
    \end{xlist}
\end{exe}

Resulta conveniente, entonces, retomar la distinción entre las nociones
semántica y sintáctica de predicado, que se introdujeron en capítulos
anteriores. El \termino{predicado semántico} es la expresión que atribuye cierto
contenido descriptivo a los argumentos. El \termino{predicado sintáctico}, por
su parte, puede definirse como el constituyente de la oración que excluye al
sujeto, que contiene al predicado semántico y cuyo núcleo típicamente expresa
las categorías de tiempo, aspecto y modo, junto con la concordancia en número y
persona con el sujeto (los constituyentes delimitados por corchetes en
\exref{ex:09-clases}) (v. \capref{6}).\pendiente{El original remite a los
corchetes de (3) para ambos predicados} En consecuencia, solo los verbos pueden
ser a la vez predicados semánticos y núcleos de predicados sintácticos, con la
excepción de ciertas estructuras sintácticamente marcadas que involucran
predicados no verbales (v. \capref{6}, §\ref{sec:06-noverbales}) ---por ejemplo,
\forma{Excelente, ese parcial / Hermosas, tus zapatillas}---
\parencite{hernanz1999suner,ditullio2014}.\pendiente{El original remite a
«cap. 2: §4»}

\begin{ejercicio}\label{ej:09-predicados}
En las siguientes oraciones resuelva lo pedido, según se muestra en el ejemplo
resuelto~1.

\begin{consignas}
  \item Identifique el predicado sintáctico entre corchetes y subraye el
        predicado semántico.
  \item Determine a qué clase de sintagma pertenece el predicado semántico que
        identificó.
\end{consignas}

\begin{exe}
  \ex Eduarda Mansilla y Lucio Mansilla eran hermanos.
  \ex Eduarda Mansilla escribió varios libros.
  \ex Eduarda Mansilla ofició de traductora.
  \ex Eduarda Mansilla fue también compositora.
  \ex Margarita Abella Caprile nació en 1901.
  \ex La poesía de Margarita Caprile era muy íntima.
  \ex Margarita Caprile fue amiga de Alfonsina Storni.
  \ex María Luisa Carnelli era de La Plata.
  \ex María Luisa Carnelli firmaba sus letras de tango con seudónimo.
\end{exe}

\begin{resuelto}
1. Eduarda Mansilla y Lucio Mansilla [eran [\etq{SN} \subrayado{hermanos}]].
\end{resuelto}
\end{ejercicio}

\subsection{Predicación primaria y predicación secundaria}\label{sec:09-primaria}

La mayor parte de los verbos del español, como \forma{vivir}, \forma{combatir} y
\forma{dar}, tienen contenido léxico, es decir, aportan significado conceptual.
Esto los habilita a funcionar sintácticamente como verbos plenos. A diferencia
de estos, algunos verbos carecen total o parcialmente de contenido conceptual y
parecen cumplir una función meramente gramatical. Estos incluyen, por un lado, a
los \termino{verbos copulativos} \forma{ser, estar, parecer}, que reciben este
nombre precisamente porque funcionan como una cópula o nexo gramatical entre un
argumento y un predicado semántico.\footnote{Para una descripción de las clases
de palabras léxicas y funcionales en un nivel inicial, ver
\textcite{giammatteo2009}, y para ampliar estos conceptos, ver
\textcite{bosque1990}.}

\begin{exe}
  \ex\label{ex:09-copulativos}
    \begin{xlist}
      \ex María \textit{es} inteligente.
      \ex Juan \textit{está} enfermo.
      \ex Pedro \textit{parece} trabajador.
    \end{xlist}
\end{exe}

Por otro lado, esta clase abarca también a los \termino{verbos semicopulativos}
(o \termino{pseudocopulativos}), como \forma{resultar, terminar,
seguir},\footnote{No debe confundirse el uso semicopulativo de estos verbos, que
es el que exige un predicativo, con el uso pleno, en el que los verbos presentan
estructura argumental, como en \forma{terminó la tarea} o \forma{sigue por el
camino amarillo}.} que, si bien no están absolutamente desprovistos de contenido
conceptual, no funcionan como verbos plenos, pues en estos usos carecen de
estructura argumental y requieren el predicativo para completar su significado.

\begin{exe}
  \ex
    \begin{xlist}
      \ex Juan \textit{resultó} herido. / \juicio{*}{Juan resultó.}
      \ex Pedro \textit{terminó} loco.
      \ex Esteban \textit{sigue} enfermo.
    \end{xlist}
\end{exe}

Al perder su contenido léxico, estos dos tipos de verbos, si bien son núcleos de
predicados sintácticos, no constituyen predicados semánticos. La única
predicación de la oración es la que expresa el predicativo que estos verbos
exigen. Se dice, entonces, que el predicativo, como \forma{inteligente},
\forma{enfermo} y \forma{trabajador} en \exref{ex:09-copulativos}, expresa una
\termino{predicación primaria}.\pendiente{El original remite a (5)}

A diferencia de los predicativos con verbos copulativos y semicopulativos,
cuando el núcleo del predicado sintáctico es un verbo con significado léxico
pleno, este núcleo aporta la predicación principal de la oración ---es decir, la
predicación primaria---, mientras que el predicativo contribuye con una
predicación adicional ---es decir, una \termino{predicación secundaria}---, como
se ve en los ejemplos de \exref{ex:09-secundaria}.

\begin{exe}
  \ex\label{ex:09-secundaria}
    \begin{xlist}
      \ex Juan \textit{habló} dormido.\\
          Predicación primaria: \forma{Juan habló}.\\
          Predicación secundaria: \forma{Juan estaba dormido}.
      \ex María \textit{lee} muy concentrada.\\
          Predicación primaria: \forma{María lee}.\\
          Predicación secundaria: \forma{María está concentrada}.
      \ex San Martín \textit{cruzó} los Andes enfermo.\\
          Predicación primaria: \forma{San Martín cruzó los Andes}.\\
          Predicación secundaria: \forma{San Martín estaba enfermo}.
    \end{xlist}
\end{exe}

\begin{ejercicio}\label{ej:09-primaria}
Identifique en estas oraciones la predicación primaria y la secundaria (si la
hubiera), como se muestra en los ejemplos resueltos de abajo.

\begin{exe}
  \ex Carlos recibió la carta azorado.
  \ex Romina quedó perpleja.
  \ex La sala está llena.
  \ex El niño lloró desconsolado.
  \ex Las Malvinas son argentinas.
  \ex Ella respondió la llamada muy tranquila.
  \ex Martín se siente solo.
\end{exe}

\begin{resuelto}[Ejemplos resueltos]
1. Predicación primaria: \forma{Carlos recibió la carta}. Predicación
secundaria: \forma{Carlos estaba azorado}. 2. Predicación primaria:
\forma{Romina quedó perpleja}. Predicación secundaria: no hay.
\end{resuelto}
\end{ejercicio}

\section{Tipos de predicativos}\label{sec:09-tipos}

Tal como mencionamos en la introducción, es posible distinguir diferentes clases
de predicativos según dos criterios: por un lado, en función de su orientación
hacia el sujeto o hacia el objeto directo de la oración y, por el otro, en
relación con que su presencia sea obligatoria o facultativa. En lo que sigue,
presentamos ambos criterios y luego mostramos cómo se integran ambas
distinciones.

\subsection{Predicativos subjetivos y objetivos}\label{sec:09-subjetivos}

Por un lado, el predicativo es un constituyente interno del predicado verbal,
por lo que está asociado al núcleo. Por otra parte, se caracteriza por
establecer una relación de predicación con el sujeto \exref{ex:09-pedro} o con
el objeto de la oración \exref{ex:09-ciencia}.

\begin{exe}
  \ex\label{ex:09-pedro} Pedro es un arquitecto muy renombrado.
  \ex\label{ex:09-ciencia} La ciencia denominó a este espécimen
      \textit{Patagosaurus fariasi}.
\end{exe}

Representamos informalmente esta relación mediante flechas (Suj: sujeto; Pred:
predicado; V: verbo; OD: objeto directo; Pvo: predicativo): en
\exref{ex:09-flechas-a} se trata de un \termino{predicativo subjetivo}, mientras
que en \exref{ex:09-flechas-b} tenemos un \termino{predicativo objetivo}
\parencite{demonte1999}:

\begin{exe}
  \ex\label{ex:09-flechas}
    \begin{xlist}
      \ex\label{ex:09-flechas-a}
          \begin{tikzpicture}[baseline=(s.base), node distance=0pt,
              every node/.style={inner sep=1pt}]
            \node (s) {[\etq{Suj} Pedro]};
            \node[right=4pt of s] (v) {[\etq{Pred} \etq{V}es};
            \node[right=3pt of v] (p) {\etq{Pvo}[un arquitecto muy renombrado]]};
            \draw[->, >=Stealth] (p.south) .. controls +(0,-.5) and +(0,-.5) .. (s.south);
          \end{tikzpicture}
          \vspace{3ex}
      \ex\label{ex:09-flechas-b}
          \begin{tikzpicture}[baseline=(s.base), node distance=0pt,
              every node/.style={inner sep=1pt, font=\small}]
            \node (s) {[\etq{Suj} La ciencia]};
            \node[right=2pt of s] (v) {[\etq{Pred} \etq{V}denominó};
            \node[right=3pt of v] (o) {\etq{OD}[a este espécimen]};
            \node[right=3pt of o] (p) {\etq{Pvo}[\textit{Patagosaurus fariasi}]]};
            \draw[->, >=Stealth] (p.south) .. controls +(0,-.5) and +(0,-.5) .. (o.south);
          \end{tikzpicture}
          \vspace{3ex}
    \end{xlist}
\end{exe}

\begin{ejercicio}\label{ej:09-tabla}
\begin{consignas}
  \item Según se muestra en el ejemplo resuelto~1, determine si los predicativos
        señalados entre corchetes en la tabla son SN, SA o SP.
  \item Identifique en las oraciones de la tabla los sujetos y los objetos
        directos, si los hubiere.
  \item ¿Qué marca morfosintáctica visibiliza la relación que se produce entre
        el sujeto o el objeto y el predicativo, cuando este es un SN, un SA o un
        SP?
\end{consignas}

{\small
\begin{tabularx}{\linewidth}{@{}r L L@{}}
  \toprule
  & a. Predicativo subjetivo & b. Predicativo objetivo \\
  \midrule
  1. & [\etq{Suj} Santos Vega] fue [\etq{SN} un payador argentino]. &
       [\etq{Suj} Los demás payadores] consideraban [\etq{COD} a Santos Vega]
       [\etq{SA} invencible]. \\
  2. & Gabino Ezeiza murió [pobre] en el barrio de Flores. & Se suele
       considerar a José Betinotti [el nexo entre los payadores y los primeros
       tangueros]. \\
  3. & Para Gabino Ezeiza, Pancho Luna ofició [de maestro]. & José Betinotti
       convirtió [en exitosa] la línea de las payadas que exaltaban el amor
       materno. \\
  4. & Arturo de Nava solía vestir [elegante]. & El Mercosur declaró
       [Patrimonio Cultural] al arte de la payada. \\
  5. & José Betinotti siempre recordaba [agradecido] a su amigo Pablo Vázquez
       en cada función. & El alcalde acusó a Juan Moreira [de mentiroso]. \\
  \bottomrule
\end{tabularx}}
\end{ejercicio}

\begin{ejercicio}\label{ej:09-ps-po}
Según se muestra en los ejemplos resueltos, identifique cuáles de los
siguientes predicativos son subjetivos (PS) y cuáles objetivos (PO).

\begin{exe}
  \ex Por las noches cantaba ensimismada.
  \ex Rosita Quiroga interpretaba las canciones muy arrabaleras.
  \ex José Betinotti se declaraba [admirador de Alfredo Palacios].
  \ex Para una grabación particular, Mariano Villar Sáenz Peña convocó [de
      guitarrista] a Rosita Quiroga.
  \ex Las sucesivas crisis económicas dejaron al Café de los Angelitos [muy
      deteriorado].
  \ex Ángel Vargas tenía un pulmón [muy comprometido].
  \ex Alfredo Palacios y Juan B.~Justo solían discutir [enfervorizados] en el
      Café de los Angelitos.
  \ex Rosita Quiroga se retiró de su carrera artística [todavía joven].
  \ex En sus últimos años de vida a Enrique Santos Discépolo se lo veía [solo].
  \ex Azucena Maizani murió [casi olvidada] en 1970.
\end{exe}

\begin{resuelto}[Ejemplos resueltos]
1. Por las noches cantaba [\etq{PS} ensimismada]. 2. Rosita Quiroga
interpretaba las canciones [\etq{PO} muy arrabaleras].
\end{resuelto}
\end{ejercicio}

\subsection{Predicativos obligatorios y no obligatorios}\label{sec:09-obligatorios}

Además de la distinción entre predicativos subjetivos y objetivos que vimos en la
sección anterior, los predicativos también se pueden distinguir respecto de su
relación con el verbo. Si retomamos las oraciones de \exref{ex:09-cat1} y
\exref{ex:09-cat2}, se puede apreciar que la omisión del predicativo en
\exref{ex:09-cat1-b} y \exref{ex:09-cat2-b} no afecta la gramaticalidad de la
oración ni la interpretación del verbo, como se aprecia en
\exref{ex:09-opcionales}. Los paréntesis indican la opcionalidad de los
predicativos en estos casos:

\begin{exe}
  \ex\label{ex:09-opcionales}
    \begin{xlist}
      \ex Varias torres de cristal se alzan (gallardas) en Buenos Aires.
      \ex Guardé la ropa (muy mojada).
    \end{xlist}
\end{exe}

En contraste, los ejemplos de \exref{ex:09-obligatorios} muestran que en
\exref{ex:09-cat1-a} y \exref{ex:09-cat2-a} no es posible omitir el predicativo:

\begin{exe}
  \ex\label{ex:09-obligatorios}
    \begin{xlist}
      \ex[*]{Pedro es. (cfr. \forma{Pedro es un arquitecto muy renombrado})}
      \ex[*]{Nombraron a este dinosaurio.\footnote{La oración resulta gramatical
          si se interpreta \forma{nombrar} como \sig{mencionar su nombre}
          (\forma{Cuando tomaron lista, nombraron a Carlos}) o como \sig{designar
          a alguien en un cargo vacante} (\forma{Nombraron un nuevo ministro de
          Hacienda}).} (cfr. \forma{Nombraron a este dinosaurio
          \textit{Patagosaurus fariasi}})}
    \end{xlist}
\end{exe}

Esto nos permite distinguir dos tipos de predicativos: los \termino{predicativos
obligatorios}, que son aquellos que no pueden omitirse, y los
\termino{predicativos no obligatorios}, que son opcionales.

\begin{ejercicio}\label{ej:09-pso-psno}
Todos los ejemplos que aparecen a continuación tienen predicativos subjetivos.
Identifíquelos e indique en cada caso si son subjetivos obligatorios (PSO) o no
(PSNO), según se muestra en los ejemplos resueltos de abajo.

\begin{exe}
  \ex Pedro estaba furioso por haber perdido las llaves.
  \ex Pedro golpeó la mesa furioso.
  \ex Ernesto es biólogo.
  \ex Carlos salió despeinado de su casa.
  \ex Mi perro corre medio rengo cuando se asusta.
  \ex Ese reloj parece nuevo.
  \ex Rocío miró la foto asombrada.
  \ex Enrique escribió su novela muy entusiasmado.
  \ex Martina resultó la ganadora del concurso.
  \ex La actriz interpretó el papel muy segura.
\end{exe}

\begin{resuelto}[Ejemplos resueltos]
1. Pedro estaba [\etq{PSO} furioso] por haber perdido las llaves. 2. Pedro
golpeó la mesa [\etq{PSNO} furioso].
\end{resuelto}
\end{ejercicio}

\begin{ejercicio}\label{ej:09-sonono}
Observe los siguientes ejemplos y resuelva lo pedido más abajo, según se muestra
en el ejemplo resuelto~1.

\begin{exe}
  \ex
    \begin{xlist}
      \ex Anda [velozmente].
      \ex Anda [cansado].
    \end{xlist}
  \ex
    \begin{xlist}
      \ex Sigue [su instinto].
      \ex Sigue [mal].
      \ex Sigue [en Madrid].
      \ex Sigue [molesto] por lo que le dije.
    \end{xlist}
  \ex
    \begin{xlist}
      \ex Se puso [un sombrero].
      \ex Se puso [rojo].
    \end{xlist}
  \ex
    \begin{xlist}
      \ex Volvió [temprano].
      \ex Volvió [enfermo].
      \ex Volvió [loco] a su amigo.
      \ex Volvió [una noche].
      \ex Se volvió [loco].
    \end{xlist}
  \ex
    \begin{xlist}
      \ex Juan se sintió [olor a humo] después de hacer el asado.
      \ex Juan se sintió [cansado].
    \end{xlist}
  \ex
    \begin{xlist}
      \ex Esteban se hizo [el desayuno].
      \ex Esteban se hizo [el tonto].
    \end{xlist}
  \ex
    \begin{xlist}
      \ex Ese tipo cayó [mal]. Por eso se rompió el tobillo.
      \ex Ese tipo me cae [simpático].
    \end{xlist}
  \ex
    \begin{xlist}
      \ex Jugó [muy concentrado].
      \ex Jugó [el partido].
    \end{xlist}
\end{exe}

\begin{consignas}[label=\roman*)]
  \item Indique cuáles de los sintagmas delimitados entre corchetes son
        predicativos y cuáles no.
  \item En los casos en que haya identificado un predicativo, indique el tipo de
        sintagma de que se trata (SN o SA) y si es obligatorio (PO) o no (PNO).
  \item En los casos en que no hay predicativo, especifique el tipo de sintagma
        y la función sintáctica del constituyente delimitado.
\end{consignas}

\begin{resuelto}
1. i) a. \forma{Anda [velozmente]}: no es predicativo. b. \forma{Anda
[cansado]}: es predicativo. ii) b. Anda [\etq{SA/PO} cansado]. iii) a. Anda
[\etq{SAdv/ACM} velozmente].
\end{resuelto}
\end{ejercicio}

\subsection{Integración de los dos tipos de predicativos}\label{sec:09-integracion}

La combinación de las distinciones presentadas en las dos secciones anteriores
se puede presentar en una tabla de doble entrada como la
tabla~\ref{tab:09-tipos}.

\begin{table}[htbp]
\centering\small
\caption{Tipos de predicativos}\label{tab:09-tipos}
\begin{tabularx}{\linewidth}{@{}>{\bfseries\raggedright\arraybackslash}p{24mm} L L@{}}
  \toprule
  & \textbf{Predicativo subjetivo} & \textbf{Predicativo objetivo} \\
  \midrule
  Predicativo obligatorio & Pedro es [un arquitecto muy renombrado]. & La
    comisión nombró al postulante [presidente de la cámara]. \\
  Predicativo no obligatorio & Varias torres de cristal se alzan [gallardas] en
    Buenos Aires. & Guardé la ropa [muy mojada]. \\
  \bottomrule
\end{tabularx}
\end{table}

Siguiendo la tabla, se reconocen cuatro tipos de predicativos:

\begin{itemize}
  \item predicativo subjetivo obligatorio $\rightarrow$ PSO;
  \item predicativo subjetivo no obligatorio $\rightarrow$ PSNO;
  \item predicativo objetivo obligatorio $\rightarrow$ POO;
  \item predicativo objetivo no obligatorio $\rightarrow$ PONO.
\end{itemize}

Cada uno de estos tipos se puede describir mediante un grupo de características
que pueden utilizarse como pruebas de reconocimiento \parencite{ditullio2014},
según mostraremos en los apartados siguientes.

\subsubsection{Predicativo subjetivo obligatorio}\label{sec:09-pso}

\begin{itemize}
  \item Puede ser un sintagma de cualquier clase: SA, SN, SAdv, SP.
  \item Ocurre solo con verbos copulativos (\forma{ser, estar, parecer}) y
        semicopulativos (\forma{volverse, resultar, semejar}).
  \item Puede ser reemplazado por el pronombre invariable \forma{lo}:
        \forma{Rosa Guerra fue periodista} $\rightarrow$ \forma{Rosa Guerra lo
        fue}. Esta prueba sirve únicamente en el caso de que el verbo sea
        copulativo, por lo que no se puede aplicar a semicopulativos.
  \item Es importante señalar que este pronombre \forma{lo} es invariable y no
        debe ser confundido con el pronombre clítico que puede funcionar como
        COD y que varía en género y número (\forma{lo, la, los, las}), como se
        aprecia en el siguiente contraste:
\end{itemize}

\begin{exe}
  \ex ¿Rosa Guerra fue periodista? Sí, \textit{lo / *la} fue.
  \ex Compró el diario / la revista $\rightarrow$ \textit{Lo / La} compró.
\end{exe}

\begin{itemize}
  \item Si el predicativo es un SA, debe concordar en género y número con el
        sujeto de la oración: \forma{Rosa Guerra era argentina; Eugenio
        Cambaceres era argentino; Ellas son argentinas}.
\end{itemize}

\subsubsection{Predicativo subjetivo no obligatorio}\label{sec:09-psno}

\begin{itemize}
  \item Ocurre con verbos de contenido léxico y constituye una predicación
        secundaria.
  \item Acepta la paráfrasis por una coordinación entre la predicación primaria,
        que es la que aporta el verbo principal, y la secundaria, que requiere
        el agregado de un verbo copulativo: \forma{Pedro salió a trabajar
        despeinado} $\rightarrow$ \forma{Pedro salió a trabajar y estaba
        despeinado}. Esta prueba sirve para distinguir predicativos de
        circunstanciales, especialmente cuando se trata de un adverbio
        adjetival,\footnote{Los adverbios adjetivales se introdujeron en el
        \capref{5}, §\ref{sec:05-sadv}.} como en \forma{Juan camina rápido}
        $\rightarrow$ \juicio{*}{\forma{Juan camina y está rápido}}.\footnote{El
        diagnóstico no siempre resulta claro con predicativos conformados por
        SSPP. Por ejemplo, \forma{Juan fue al trabajo con un maletín} admite la
        coordinación \forma{Juan fue al trabajo y Juan estaba con un maletín}.}
  \item Si el predicativo es un SA, concuerda en género y número con el sujeto
        de la oración: \forma{Juan Manuel de Rosas leyó furioso / *furiosa el
        libro de Sarmiento; Manuelita Rosas leyó *furioso / furiosa el libro de
        Sarmiento}.
\end{itemize}

\subsubsection{Predicativo objetivo obligatorio}\label{sec:09-poo}

\begin{itemize}
  \item Ocurre con verbos que seleccionan un objeto directo y un predicativo,
        como, por ejemplo, \forma{nombrar, llamar, designar, notar, encontrar},
        etc. Estos verbos tienden a ser ambiguos con una lectura en la que no se
        selecciona un predicativo. Nótese que \forma{nombrar a alguien profesor}
        no significa lo mismo que \forma{nombrar a alguien}.
  \item Al ser un constituyente diferente del COD, puede anteponerse a este, es
        decir que acepta la prueba de movimiento:\footnote{La prueba de
        movimiento y las demás pruebas de reconocimiento de constituyentes se
        presentaron en el \capref{5}, §\ref{sec:05-pruebas}.} \forma{Nombraron
        al nuevo jugador capitán del equipo} $\rightarrow$ \forma{Nombraron
        capitán del equipo al nuevo jugador}.
  \item Asimismo, si el COD es reemplazado por el pronombre clítico
        correspondiente (\forma{lo, la, los, las}), el predicativo no queda
        incluido en la sustitución:
\end{itemize}

\begin{exe}
  \ex Considero \textit{esta novela} [valiosa] $\rightarrow$ \textit{La}
      considero [valiosa].
\end{exe}

\begin{itemize}
  \item Si el predicativo es un SA, concuerda en género y número con el objeto
        de la oración.
\end{itemize}

\begin{exe}
  \ex Considero \textit{esta novela} [muy valiosa]. / Considero \textit{estas
      novelas} [muy valiosas].
\end{exe}

\subsubsection{Predicativo objetivo no obligatorio}\label{sec:09-pono}

\begin{itemize}
  \item Ocurre solo con verbos transitivos.
  \item Puede anteponerse al COD (movimiento).
\end{itemize}

\begin{exe}
  \ex Encontró \textit{la mesa} [sucia] $\rightarrow$ Encontró [sucia]
      \textit{la mesa}.
\end{exe}

\begin{itemize}
  \item Si el COD es reemplazado por el pronombre clítico correspondiente
        (\forma{lo, la, los, las}), el predicativo no queda incluido en la
        sustitución:
\end{itemize}

\begin{exe}
  \ex Saludó \textit{a Pedro} [sentado en el parque] $\rightarrow$ \textit{Lo}
      saludó [sentado en el parque].
\end{exe}

\begin{itemize}
  \item Si el predicativo es un SA, concuerda en género y número con el objeto
        de la oración.
\end{itemize}

\begin{exe}
  \ex Encontró \textit{la mesa} [sucia]. / Encontró \textit{las mesas}
      [sucias].
\end{exe}

\begin{recuadro}[Para ampliar]
Los verbos que exigen predicativo objetivo obligatorio pertenecen a las
siguientes clases: verbos de entendimiento (\forma{considerar, creer, imaginar,
suponer, juzgar}), verbos de percepción sensible o intelectual (\forma{ver, oír,
notar, descubrir, encontrar}), verbos designativos (\forma{designar, elegir,
nombrar, denominar, llamar}) y verbos volitivos (\forma{querer})
\parencite[134]{ditullio2014}.
\end{recuadro}

\begin{ejercicio}\label{ej:09-pruebas}
Reconozca qué pruebas de las planteadas en §\ref{sec:09-integracion} se están
aplicando en los siguientes ejemplos con el fin de demostrar que los elementos
subrayados en la primera oración de cada par son predicativos, según se muestra
en el ejemplo resuelto~1.\pendiente{El original remite a «§5»}

\begin{exe}
  \ex
    \begin{xlist}
      \ex María sintió el pollo \subrayado{soso}.
      \ex María sintió soso el pollo.
    \end{xlist}
  \ex
    \begin{xlist}
      \ex Gaspar parecía \subrayado{distante}.
      \ex Gaspar lo parecía.
    \end{xlist}
  \ex
    \begin{xlist}
      \ex Consiguió \subrayado{caras} las entradas.
      \ex Consiguió las entradas y las entradas estaban caras.
    \end{xlist}
  \ex
    \begin{xlist}
      \ex La tarea le había resultado \subrayado{fácil}.
      \ex Las tareas le habían resultado fáciles.
    \end{xlist}
  \ex
    \begin{xlist}
      \ex Gaspar encontró las zapatillas \subrayado{baratas}.
      \ex Gaspar encontró baratas las zapatillas.
    \end{xlist}
\end{exe}

\begin{resuelto}
1. Prueba de movimiento: muestra que el predicativo no está incluido en el COD,
es decir que se trata de un constituyente diferente del COD.
\end{resuelto}
\end{ejercicio}

\begin{ejercicio}\label{ej:09-enriquez}
Lea atentamente el siguiente fragmento adaptado de «Alguien camina sobre tu
tumba», de Mariana Enriquez (2009), y resuelva las consignas que aparecen a
continuación, según se muestra en el ejemplo resuelto.

\begin{consignas}
  \item Identifique cuáles de los constituyentes destacados no son
        predicativos. Justifique su respuesta.
  \item Indique cuáles de las expresiones destacadas son predicativos y
        clasifíquelos. Explicite los criterios utilizados.
\end{consignas}

% Fragmento literario protegido: se incorpora a mano desde el original.
\textoaincorporar{el fragmento adaptado de «Alguien camina sobre tu tumba», de
M.~Enriquez, desde «El ferry que lleva a Paul hasta su trabajo» hasta
«\ldots{}en ocasiones, violenta», con los constituyentes destacados entre
corchetes: \textit{cerca del puerto; el uniforme; mi novio; australiano; de
mecánico de bicicletas; capital de Australia Occidental; a larga distancia;
insólitamente estable, permanente y feliz; la capital de estado más aislada del
mundo; la ciudad de más de un millón de habitantes más aislada del mundo;
comprobable y cierto; una localidad muy cool, algo estrafalaria, favorita de
mochileros y newagers; con jardines llenos de cristales curativos y hadas}}

\begin{resuelto}
a) \forma{El uniforme} no es un predicativo, sino un COD ([\etq{COD} el
uniforme]), ya que admite la sustitución por un pronombre acusativo. b)
\forma{Mi novio} es un predicativo subjetivo obligatorio, ya que ocurre como
complemento de un verbo copulativo y admite la sustitución por el pronombre
invariable \forma{lo} (\forma{Paul lo es}).
\end{resuelto}
\end{ejercicio}

\begin{ejercicio}\label{ej:09-significado}
Según se muestra en el ejemplo resuelto~1, responda las siguientes consignas.

\begin{consignas}[label=\roman*)]
  \item ¿Cuáles de los siguientes ejemplos presentan predicativos?
  \item ¿Cómo afecta la presencia de predicativos al significado de estos
        verbos? Explique en qué consiste la diferencia de significado.
\end{consignas}

\begin{exe}
  \ex
    \begin{xlist}
      \ex Paul encontró el uniforme sobre la mesa.
      \ex Paul encontró el uniforme sobre la mesa muy arrugado.
    \end{xlist}
  \ex
    \begin{xlist}
      \ex La gente permaneció estupefacta.
      \ex La gente permaneció en el lugar.
      \ex La gente permaneció estupefacta en el lugar.
    \end{xlist}
  \ex
    \begin{xlist}
      \ex Perth está aislada.
      \ex Perth está en Australia.
    \end{xlist}
  \ex
    \begin{xlist}
      \ex El testigo declaró en el juicio.
      \ex El juez declaró culpable al imputado.
    \end{xlist}
\end{exe}

\begin{resuelto}
i) 1a. No presenta ningún predicativo. 1b. PO [muy arrugado]. ii) En 1a,
\forma{encontrar} significa \sig{hallar algo cuyo paradero se desconocía}; en
1b, la predicación describe el estado en que el objeto ---el uniforme--- se
encontró.
\end{resuelto}
\end{ejercicio}

\begin{ejercicio}\label{ej:09-vf}
Revise las respuestas del ejercicio~\ref{ej:09-primaria} y determine cuáles de
estas afirmaciones son verdaderas y cuáles falsas, según se muestra en el
ejemplo resuelto i.

\begin{enumerate}[label=\roman*.]
  \item Los PSO expresan predicaciones primarias. \hfill V / F
  \item Los PSO expresan predicaciones secundarias. \hfill V / F
  \item Los PSNO expresan predicaciones primarias. \hfill V / F
  \item Los PSNO expresan predicaciones secundarias. \hfill V / F
\end{enumerate}

\begin{resuelto}
i. V.
\end{resuelto}
\end{ejercicio}

\begin{ejercicio}\label{ej:09-verbos}
\begin{consignas}
  \item Determine cuáles de los verbos que aparecen en la siguiente lista
        requieren predicativos en alguna lectura y complete la tabla.
  \item Siempre que sea posible, construya un ejemplo con un predicativo, tal
        como se ilustra en el ejemplo resuelto~1.
\end{consignas}

\begin{center}\small
\begin{tabular}{@{}r l c p{55mm}@{}}
  \toprule
  & Verbo & Requiere predicativo & Ejemplo \\
  \midrule
  1.  & \forma{parecer}   & sí & \respuesta{Pedro parece cansado.} \\
  2.  & \forma{escribir}  & & \\
  3.  & \forma{caer}      & & \\
  4.  & \forma{volverse}  & & \\
  5.  & \forma{construir} & & \\
  6.  & \forma{comer}     & & \\
  7.  & \forma{correr}    & & \\
  8.  & \forma{seguir}    & & \\
  9.  & \forma{resultar}  & & \\
  10. & \forma{haber}     & & \\
  11. & \forma{estar}     & & \\
  12. & \forma{ser}       & & \\
  13. & \forma{estudiar}  & & \\
  \bottomrule
\end{tabular}
\end{center}
\end{ejercicio}

\section{La sintaxis de los predicativos}\label{sec:09-sintaxis}

Como se presentó en el \capref{6}, en el marco de la gramática generativa, los
predicados sintácticos están compuestos por una proyección funcional que porta
la información flexiva (SF) y que contiene, a su vez, un sintagma verbal (SV),
que es el encargado de expresar la información eventiva.\footnote{Para más
detalles sobre las estructuras predicativas remitimos a
\textcite[211--250]{hernanz1987}.} Los predicativos, por su parte, son
semejantes a los predicados, pero carecen de toda esa articulación. Por eso, en
la bibliografía se los ha analizado en los términos de lo que se denomina
\termino{oración reducida} (OR), un tipo de estructura predicativa mínima que,
si bien expresa una predicación, carece de la manifestación de las categorías de
tiempo, aspecto y modo.\footnote{Para más detalles sobre las oraciones
reducidas, ver \textcite[425--429]{bosque2009}.} Dado que esa clase de análisis
requiere de una estructura más compleja, a los fines de este libro vamos a
marcar los predicativos simplemente del siguiente modo:

\begin{exe}
  \ex\label{ex:09-marcado}
    \begin{xlist}
      \ex Predicativo subjetivo obligatorio:\\
          {}[\etq{SF/O} [\etq{SN/Suj} Azucena Maizani] [\etq{SV/Pred}
          fue\etq{V/n} [\etq{SN/PSO} una\etq{Q/Esp} cantante\etq{N/n}]]]
      \ex Predicativo subjetivo no obligatorio:\\
          {}[\etq{SF/O} [\etq{SN/Suj} Carlos] [\etq{SV/Pred} volvió\etq{V/n}
          [\etq{SA/PSNO} cansado]]]
      \ex Predicativo objetivo obligatorio:\\
          {}[\etq{SF/O} [\etq{SN/Suj} la\etq{D/Esp} justicia\etq{N/n}]
          [\etq{SV/Pred} declaró\etq{V/n} [\etq{SN/COD} a\etq{MF} la\etq{D/Esp}
          acusada\etq{N/n}] [\etq{SA/POO} culpable]]]
      \ex Predicativo objetivo no obligatorio:\\
          {}[\etq{SF/O} [\etq{SN/Suj} Alfredo] [\etq{SV/Pred} compró\etq{V/n}
          [\etq{SN/COD} el\etq{D/Esp} disco\etq{N/n}] [\etq{SA/PONO} barato]]]
    \end{xlist}
\end{exe}

\begin{recuadro}[Para ampliar: análisis de oración reducida]
El análisis en términos de OR asume que el predicativo y el argumento al que se
aplica conforman una cláusula que no flexiona en tiempo, aspecto y modo. Para
obtener una oración completa, este análisis requiere algún mecanismo que
garantice que el sujeto de la oración reducida sea correferente con el sujeto u
objeto de la oración principal, según se trate de un predicativo subjetivo u
objetivo, respectivamente. Por ejemplo, en el caso del predicativo subjetivo
obligatorio, el sujeto se mueve desde la cláusula reducida, como se ve en el
siguiente árbol:

\begin{exe}
  \ex\label{ex:09-or1}
    \begin{xlist}
      \ex Azucena Maizani fue una cantante de tangos.
      \ex {}[\etq{SF/O} [\etq{SN/Suj} Azucena Maizani]\etq{i} [\etq{SV/Pred}
          fue\etq{V/n} [\etq{ORed/Comp} \sout{Azucena Maizani}\etq{i} [\etq{SN}
          una\etq{Q/Esp} cantante\etq{N/n} [\etq{SP/M} de\etq{P/n}
          [\etq{SN/Comp} tangos\etq{N/n}]]]]]]
      \ex
          \begin{forest} gramatica
            [SF / O
              [SN\etq{i}\\\textit{Azucena Maizani}]
              [SV
                [V\\\textit{fue}]
                [ORed
                  [SN\etq{i}\\{\color{gristexto}\textit{Azucena Maizani}}]
                  [SN [\textit{una cantante de tangos}, roof]]]]]
          \end{forest}
    \end{xlist}
\end{exe}

En este análisis, los dos elementos que aparecen en la oración reducida, la
huella de \forma{Azucena Maizani}, por un lado, y \forma{una cantante de
tangos}, por otro, no se etiquetan en relación con la función sintáctica que
cumplen, ya que la oración reducida carece de articulación interna. Los
predicativos subjetivos no obligatorios, como \exref{ex:09-or2}, se pueden
analizar también como oraciones reducidas.\pendiente{El original remite a (20)}
Sin embargo, a diferencia de lo que sucede con los obligatorios, en los no
obligatorios el argumento del cual se predica en la oración reducida no es un
elemento que luego se mueva, sino un pronombre nulo controlado, es decir, que es
obligatoriamente correferente con el sujeto de la oración principal. Llamaremos
a este pronombre PRO (no confundir con \textit{pro} en minúscula):

\begin{exe}
  \ex\label{ex:09-or2}
    \begin{xlist}
      \ex Carlos volvió cansado.
      \ex {}[\etq{SF/O} [\etq{SN/Suj} Carlos]\etq{i} [\etq{SV/Pred}
          volvió\etq{V/n} [\etq{ORed/Comp} PRO\etq{i} cansado\etq{SA}]]]
      \ex
          \begin{forest} gramatica
            [SF / O
              [SN\etq{i}\\\textit{Carlos}]
              [SV
                [V\\\textit{volvió}]
                [ORed
                  [PRO\etq{i}]
                  [SA\\\textit{cansado}]]]]
          \end{forest}
    \end{xlist}
\end{exe}

Análisis semejantes se pueden postular para los predicativos objetivos, tal como
se ilustra, respectivamente, en los siguientes ejemplos de predicativo objetivo
obligatorio y no obligatorio:

\begin{exe}
  \ex\label{ex:09-or3}
    \begin{xlist}
      \ex La justicia declaró a la acusada culpable.
      \ex {}[\etq{SF/O} [\etq{SN/Suj} la\etq{D/Esp} justicia\etq{N/n}]
          [\etq{SV/Pred} declaró\etq{V/n} [\etq{ORed/Comp} [\etq{SN/COD}
          a\etq{MF} la\etq{D/Esp} acusada\etq{N/n}] [\etq{SA} culpable]]]]
      \ex
          \begin{forest} gramatica
            [SF / O
              [SN [\textit{la justicia}, roof]]
              [SV
                [V\\\textit{declaró}]
                [ORed
                  [SN [\textit{a la acusada}, roof]]
                  [SA\\\textit{culpable}]]]]
          \end{forest}
    \end{xlist}
  \ex\label{ex:09-or4}
    \begin{xlist}
      \ex Alfredo compró el disco barato.
      \ex {}[\etq{SF/O} [\etq{SN/Suj} Alfredo] [\etq{SV/Pred} compró\etq{V/n}
          [\etq{SN/COD} el\etq{D/Esp} disco\etq{N/n}]\etq{i} [\etq{ORed}
          PRO\etq{i} barato\etq{SA}]]]
      \ex
          \begin{forest} gramatica
            [SF / O
              [SN\\\textit{Alfredo}]
              [SV
                [V\prima
                  [V\\\textit{compró}]
                  [SN\etq{i} [\textit{el disco}, roof]]]
                [ORed
                  [PRO\etq{i}]
                  [SA\\\textit{barato}]]]]
          \end{forest}
    \end{xlist}
\end{exe}
\end{recuadro}

\begin{ejercicio}\label{ej:09-analisis}
Haga el análisis sintáctico de las siguientes oraciones, según se muestra en el
ejemplo resuelto~1.

\begin{exe}
  \ex Ahora oficia de mecánico de bicicletas.
  \ex Mi hermano siempre toma el mate amargo.
  \ex Josefina desayuna callada.
  \ex Los niños escondieron culposos el jarrón roto.
  \ex En cada picadito de Argentina, un patadura se cree Messi.
  \ex Los violinistas callejeros se hacen los Paganini.
  \ex Me desperté muerto de hambre.
  \ex Las nuevas medidas dejaron desfinanciada la educación.
\end{exe}

\begin{resuelto}
1. [\etq{SF/O} [\etq{SN/Suj} \textit{pro}] [\etq{SV/Pred} [\etq{SAdv/ACT}
ahora\etq{Adv/n}] oficia\etq{V/n} [\etq{SP/PSO} de\etq{P/n} [\etq{SN/Comp}
mecánico\etq{N/n} [\etq{SP/M} de\etq{P/n} [\etq{SN/Comp}
bicicletas\etq{N/n}]]]]]]
\end{resuelto}
\end{ejercicio}

\section{Predicativos y homonimia estructural}\label{sec:09-homonimia}

Tal como se observó previamente, los predicativos no están limitados
categorialmente ---es decir, pueden ser sintagmas de cualquier clase--- y
establecen diferentes relaciones estructurales respecto del verbo que funciona
como núcleo del predicado y del argumento con el que se vinculan, sea el sujeto
o el objeto.\footnote{La homonimia estructural se introdujo en el \capref{2},
§\ref{sec:02-ambiguedad}, y en el \capref{5}, §\ref{sec:05-homonimia}.} Estas
dos características hacen que los predicativos tiendan a aparecer en oraciones
que son ambiguas por razones estructurales. El ejemplo de
\exref{ex:09-sentados} tiene tres lecturas posibles porque puede recibir tres
análisis sintácticos diferentes: en \exref{ex:09-sentados-a}, \forma{sentados}
forma parte del complemento objeto directo; en \exref{ex:09-sentados-b},
\forma{sentados} es un predicativo objetivo (los que están sentados son los
jugadores), y en \exref{ex:09-sentados-c} se trata de un predicativo subjetivo
(los que están sentados son los preparadores).

\begin{exe}
  \ex\label{ex:09-sentados} Los preparadores físicos observaban a los jugadores
      sentados.
    \begin{xlist}
      \ex\label{ex:09-sentados-a} {}[\etq{SF/O} [\etq{SN/Suj} los preparadores
          físicos] [\etq{SV/Pred} observaban [\etq{SN/COD} a los jugadores
          sentados]]]
      \ex\label{ex:09-sentados-b} {}[\etq{SF/O} [\etq{SN/Suj} los preparadores
          físicos] [\etq{SV/Pred} observaban [\etq{SN/COD} a los jugadores]
          [\etq{SA/PONO} sentados]]]
      \ex\label{ex:09-sentados-c} {}[\etq{SF/O} [\etq{SN/Suj} los preparadores
          físicos] [\etq{SV/Pred} observaban [\etq{SN/COD} a los jugadores]
          [\etq{SA/PSNO} sentados]]]
    \end{xlist}
\end{exe}
\pendiente{En los análisis, el original pone «Los entrenadores» como sujeto}

\begin{ejercicio}\label{ej:09-ambiguas}
Las siguientes oraciones son ambiguas. Explique mediante una paráfrasis cuáles
son las interpretaciones posibles y presente los análisis correspondientes para
cada una de esas lecturas, según se muestra en el ejemplo resuelto~1.

\begin{exe}
  \ex Trajeron a los perros enfermos.
  \ex Les entregaron los platos sucios.
  \ex Pedro pidió un postre borracho.
\end{exe}

\begin{resuelto}
1. a. [\etq{SF/O} [\etq{SN/Suj} \textit{pro}\etq{Pr/n}] [\etq{SV/Pred}
trajeron\etq{V/n} [\etq{SN/COD} a\etq{MF} los\etq{D/Esp} perros\etq{N/n}]
[\etq{SA/PO} enfermos\etq{A/n}]]]: \sig{Quienes estaban enfermos trajeron a los
perros}. b. [\etq{SF/O} [\etq{SN/Suj} \textit{pro}\etq{Pr/n}] [\etq{SV/Pred}
trajeron\etq{V/n} [\etq{SN/COD} a\etq{MF} los\etq{D/Esp} perros\etq{N/n}]
[\etq{SA/PSNO} enfermos\etq{A/n}]]]: \sig{Cuando trajeron a los perros, estos
estaban enfermos}. c. [\etq{SF/O} [\etq{SN/Suj} \textit{pro}\etq{Pr/n}]
[\etq{SV/Pred} trajeron\etq{V/n} [\etq{SN/COD} a\etq{MF} los\etq{D/Esp}
perros\etq{N/n} enfermos\etq{A/M}]]]: \sig{Trajeron a los perros, que estaban
enfermos}.\pendiente{Las etiquetas de a y b están intercambiadas respecto de las
paráfrasis (a es PSNO; b, PONO); el original rotula las lecturas a, b, b}
\end{resuelto}
\end{ejercicio}

\begin{ejercicio}\label{ej:09-vinas}
Lea atentamente el siguiente fragmento de \textit{Un dios cotidiano}, de
D.~Viñas. Luego identifique todos los predicativos y señale de qué tipo de
predicativo se trata en cada caso.

% Fragmento literario protegido: se incorpora a mano desde el original.
\textoaincorporar{el fragmento de \textit{Un dios cotidiano}, de D.~Viñas, desde
«Que-lo-nios, dije tres o cuatro veces» hasta «\ldots{}la voz me salía pastosa,
la matizaba»}
\end{ejercicio}

\begin{recuadro}[Reflexión final]
En este capítulo hemos abordado un tipo particular de función del sintagma
verbal: el predicativo. Al respecto, hemos observado que los predicativos pueden
corresponder a diferentes categorías: SN, SA, SP e incluso SAdv. Hemos
explorado también la tipología básica de los predicativos, que se desarrolla en
torno a dos parámetros: i) si están orientados al sujeto o al objeto y ii) si
son obligatorios o no obligatorios. Esto lleva a establecer cuatro tipos de
predicativos: subjetivo obligatorio (PSO), subjetivo no obligatorio (PSNO),
objetivo obligatorio (POO) y objetivo no obligatorio (PONO). Hemos también
presentado cómo se los analiza sintácticamente y qué propiedades sintácticas
tiene cada subclase.
\end{recuadro}

% --- Secciones finales -------------------------------------------------
\lecturasrecomendadas

\solucionario

\begin{solucion}{\ref{ej:09-predicados}}
2. Eduarda Mansilla [\etq{SV} \subrayado{escribió varios libros}]. 3. Eduarda
Mansilla [\etq{SV} \subrayado{ofició de traductora}]. 4. Eduarda Mansilla [fue
también [\etq{SN} \subrayado{compositora}]]. 5. Margarita Abella Caprile
[\etq{SV} \subrayado{nació en 1901}]. 6. La poesía de Margarita Caprile [era
[\etq{SA} \subrayado{muy íntima}]]. 7. Margarita Caprile [fue [\etq{SN}
\subrayado{amiga de Alfonsina Storni}]]. 8. María Luisa Carnelli [era [\etq{SP}
\subrayado{de La Plata}]]. 9. María Luisa Carnelli [\etq{SV}
\subrayado{firmaba sus letras de tango con seudónimo}].
\end{solucion}

\begin{solucion}{\ref{ej:09-primaria}}
3. Predicación primaria: \forma{La sala está llena}. Predicación secundaria: no
hay. 4. Predicación primaria: \forma{El niño lloró}. Predicación secundaria:
\forma{El niño estaba desconsolado}. 5. Predicación primaria: \forma{Las
Malvinas son argentinas}. Predicación secundaria: no hay. 6. Predicación
primaria: \forma{Ella respondió la llamada}. Predicación secundaria: \forma{Ella
estaba muy tranquila}. 7. Predicación primaria: \forma{Martín se siente solo}.
Predicación secundaria: no hay.
\end{solucion}

\begin{solucion}{\ref{ej:09-tabla}}
a y b. 1a. [\etq{Suj} Santos Vega] fue [\etq{SN} un payador argentino]. 1b.
[\etq{Suj} Los demás payadores] consideraban [\etq{COD} a Santos Vega]
[\etq{SA} invencible]. 2a. [\etq{Suj} Gabino Ezeiza] murió [\etq{SA} pobre] en
el barrio de Flores. 2b. Se suele considerar [\etq{COD} a José Betinotti]
[\etq{SN} el nexo entre los payadores y los primeros tangueros]. 3a. Para Gabino
Ezeiza, [\etq{Suj} Pancho Luna] ofició [\etq{SP} de maestro]. 3b. [\etq{Suj}
José Betinotti] convirtió [\etq{SP} en exitosa] [\etq{COD} la línea de las
payadas que exaltaban el amor materno]. 4a. [\etq{Suj} Arturo de Nava] solía
vestir [\etq{SA} elegante]. 4b. [\etq{Suj} El Mercosur] declaró [\etq{SN}
Patrimonio Cultural] [\etq{COD} al arte de la payada]. 5a. [\etq{Suj} José
Betinotti] siempre recordaba [\etq{SA} agradecido] [\etq{COD} a su amigo Pablo
Vázquez] en cada función. 5b. [\etq{Suj} El alcalde] acusó [\etq{COD} a Juan
Moreira] [\etq{SP} de mentiroso].

c. La marca morfosintáctica que visibiliza la relación entre el predicativo y el
sujeto o el objeto es la concordancia en género y número, en los casos en que el
núcleo del predicativo es un adjetivo. En los casos en que el predicativo es un
SN o un SP, la relación que se establece es de compatibilidad de rasgos
semánticos, pero no una concordancia propiamente dicha. En el contexto de los
SSPP, esa compatibilidad se da con el complemento de la preposición.
\end{solucion}

\begin{solucion}{\ref{ej:09-ps-po}}
3. José Betinotti se declaraba [\etq{PS} admirador de Alfredo Palacios]. 4. Para
una grabación particular, Mariano Villar Sáenz Peña convocó [\etq{PO} de
guitarrista] a Rosita Quiroga. 5. Las sucesivas crisis económicas dejaron al
Café de los Angelitos [\etq{PO} muy deteriorado]. 6. Ángel Vargas tenía un
pulmón [\etq{PO} muy comprometido]. 7. Alfredo Palacios y Juan B.~Justo solían
discutir [\etq{PS} enfervorizados] en el Café de los Angelitos. 8. Rosita Quiroga
se retiró de su carrera artística [\etq{PS} todavía joven]. 9. En sus últimos
años de vida, a Enrique Santos Discépolo se lo veía [\etq{PS}
solo].\pendiente{Con \textit{se lo veía}, \textit{solo} se predica del objeto
(\textit{lo})} 10. Azucena Maizani murió [\etq{PS} casi olvidada] en 1970.
\end{solucion}

\begin{solucion}{\ref{ej:09-pso-psno}}
3. Ernesto es [\etq{PSO} biólogo]. 4. Carlos salió [\etq{PSNO} despeinado] de su
casa. 5. Mi perro corre [\etq{PSNO} medio rengo] cuando se asusta. 6. Ese reloj
parece [\etq{PSO} nuevo]. 7. Rocío miró la foto [\etq{PSNO} asombrada]. 8.
Enrique escribió su novela [\etq{PSNO} muy entusiasmado]. 9. Martina resultó
[\etq{PSO} la ganadora del concurso]. 10. La actriz interpretó el papel
[\etq{PSNO} muy segura].
\end{solucion}

\begin{solucion}{\ref{ej:09-sonono}}
i) Hay predicativos solo en 1b, 2b, 2d, 3b, 4b, 4c, 4e, 5b, 6b, 7b y 8a.
\pendiente{El original omite 1b de la lista, aunque lo resuelve, y analiza 2b como
predicativo en iii)}

ii) 2d. Sigue [\etq{SA/PO} molesto] por lo que le dije. 3b. Se puso [\etq{SA/PO}
rojo]. 4b. Volvió [\etq{SA/PNO} enfermo]. 4c. Volvió [\etq{SA/POO} loco] a su
amigo. 4e. Se volvió [\etq{SA/PO} loco]. 5b. Juan se sintió [\etq{SA/PO}
cansado]. 6b. Esteban se hizo [\etq{SN/PO} el tonto]. 7b. Ese tipo me cae
[\etq{SA/PSO} simpático]. 8a. Jugó [\etq{SA/PNO} muy concentrado].

iii) 2a. Sigue [\etq{SN/COD} su instinto]. 2b. Sigue [\etq{SAdv/PO} mal]. 2c.
Sigue [\etq{SP/ACL} en Madrid]. 3a. Se puso [\etq{SN/COD} un sombrero]. 4a.
Volvió [\etq{SAdv/ACT} temprano]. 4d. Volvió [\etq{SN/ACT} una noche]. 5a. Juan
se sintió [\etq{SN/COD} olor a humo] después de hacer el asado. 6a. Esteban se
hizo [\etq{SN/COD} el desayuno]. 7a. Ese tipo cayó [\etq{SAdv/ACM} mal]. Por eso
se rompió [\etq{SN/COD} el tobillo]. 8b. Jugó [\etq{SN/COD} el partido].
\end{solucion}

\begin{solucion}{\ref{ej:09-pruebas}}
2. Prueba de reemplazo por \forma{lo} invariable. 3. Prueba de coordinación
entre la predicación primaria y la secundaria. 4. Prueba de concordancia. 5.
Prueba de movimiento: misma explicación que en el ejemplo 1.
\end{solucion}

\begin{solucion}{\ref{ej:09-enriquez}}
[\etq{CCL} cerca del puerto] / [\etq{PSO} australiano] / [\etq{PSNO} de mecánico
de bicicletas] / [\etq{ApE} capital de Australia Occidental] / [\etq{M} a larga
distancia] / [\etq{PSO} insólitamente estable, permanente y feliz] / [\etq{PSO}
la capital de estado más aislada del mundo] / [\etq{PSO} la ciudad de más de un
millón de habitantes más aislada del mundo] / [\etq{PSO} comprobable y cierto] /
[\etq{PSO} una localidad muy cool, algo estrafalaria, favorita de mochileros y
\textit{newagers}] / [\etq{M} con jardines llenos de cristales curativos y
hadas].
\end{solucion}

\begin{solucion}{\ref{ej:09-significado}}
2. i) a y c: [\etq{PSO} estupefacta]; b no presenta ningún predicativo. ii) En
el primer y el tercer uso, \forma{permanecer} refiere a que \sig{la gente no
salió de su asombro}; en el segundo caso, se dice que \sig{la gente no se
trasladó espacialmente}. 3. i) a: [\etq{PSO} aislada]; b no presenta ningún
predicativo. ii) En a, el predicativo describe el estado de Perth; en b, se
indica su ubicación espacial. 4. i) a no presenta ningún predicativo; b:
[\etq{PO} culpable]. ii) En a, el verbo \forma{declarar} significa \sig{hizo una
declaración}; en b, la predicación indica que \sig{el juez condenó al imputado}.
\end{solucion}

\begin{solucion}{\ref{ej:09-vf}}
ii. F. iii. F. iv. V.
\end{solucion}

\begin{solucion}{\ref{ej:09-verbos}}
a. No requieren predicativo: \forma{escribir, caer, construir, comer, correr,
haber, estudiar}; requieren predicativo: \forma{volverse, seguir, resultar,
estar, ser}. b. Respuestas posibles: 4. El profesor se volvió loco. 8. Mi
hermano sigue en cama / enfermo. 9. La clase de gramática resultó muy
interesante. 11. La sopa está caliente. 12. Comer frutas por la mañana es
verdaderamente importante.
\end{solucion}

\begin{solucion}{\ref{ej:09-analisis}}
2. [\etq{SF/O} [\etq{SN/Suj} mi\etq{D/Esp} hermano\etq{N/n}] [\etq{SV/Pred}
[\etq{SAdv/ACT} siempre\etq{Adv/n}] toma\etq{V/n} [\etq{SN/COD} el\etq{D/Esp}
mate\etq{N/n}] [\etq{SA/PONO} amargo\etq{A/n}]]].

3. [\etq{SF/O} [\etq{SN/Suj} Josefina\etq{N/n}] [\etq{SV/Pred} desayuna\etq{V/n}
[\etq{SA/PSNO} callada\etq{A/n}]]].

4. [\etq{SC=SF/O} [\etq{SN/Suj} los\etq{D/Esp} niños\etq{N/n}] [\etq{SV/Pred}
escondieron\etq{V/n} [\etq{SA/PSNO} culposos\etq{A/n}] [\etq{SN/COD}
el\etq{D/Esp} jarrón\etq{N/n} roto\etq{A/M}]]].

5. [\etq{SF/O} [\etq{SP/ACL} en\etq{P/n} [\etq{SN/Comp} cada\etq{Q/Esp}
picadito\etq{N/n} [\etq{SP/M} de\etq{P/n} [Argentina\etq{N/n}]]]]\etq{i},
[\etq{SN/Suj} un\etq{Q/Esp} patadura\etq{N/n}] [\etq{SV/Pred} se\etq{IPCR}
cree\etq{V/n} [\etq{SN/PSO} Messi\etq{N/n}] h\etq{i}]]. \textit{Aclaración}: el
constituyente \forma{En cada picadito de Argentina} se halla desplazado de la
posición en la que se origina, interna al predicado.

6. [\etq{SF/O} [\etq{SN/Suj} los\etq{D/Esp} violinistas\etq{N/n}
callejeros\etq{A/M}] [\etq{SV/Pred} se\etq{IPCR} hacen\etq{V/n} [\etq{SN/PSO}
los\etq{D/Esp} Paganini\etq{N/n}]]]. \textit{Aclaración}: \forma{Paganini}
aparece recategorizado como nombre común.

7. [\etq{SF/O} [\etq{SN/Suj} \textit{pro}] [\etq{SV/Pred} me\etq{IPCR}
desperté\etq{V/n} [\etq{SA/PSNO} muerto\etq{A/n} [\etq{SP/M} de\etq{P/n}
[\etq{SN/Comp} hambre\etq{N/n}]]]]].

8. [\etq{SF/O} [\etq{SN/Suj} las\etq{D/Esp} nuevas\etq{A/M} medidas\etq{N/n}]
[\etq{SV/Pred} dejaron\etq{V/n} [\etq{SA/POO} desfinanciada\etq{A/n}]
[\etq{SN/COD} la\etq{D/Esp} educación\etq{N/n}]]].
\end{solucion}

\begin{solucion}{\ref{ej:09-ambiguas}}
En cada caso incluimos una paráfrasis de modelo. Eso no implica que no haya
otras paráfrasis posibles.

2. a. [\etq{SF/O} [\etq{SN/Suj} \textit{pro}] [\etq{SV/Pred} les\etq{Pr/COI}
entregaron\etq{V/n} [\etq{SN/COD} los\etq{D/Esp} platos\etq{N/n}] [\etq{SA/PO}
sucios\etq{A/n}]]]: \sig{Les entregaron sucios los platos}. b. [\etq{SF/O}
[\etq{SN/Suj} \textit{pro}] [\etq{SV/Pred} les\etq{Pr/COI} entregaron\etq{V/n}
[\etq{SN/COD} los\etq{D/Esp} platos\etq{N/n} sucios\etq{A/M}]]]: \sig{Les
entregaron los platos que estaban sucios}.

3. a. [\etq{SF/O} [\etq{SN/Suj} Pedro\etq{N/n}] [\etq{SV/Pred} pidió\etq{V/n}
[\etq{SN/COD} un\etq{D/Esp} postre\etq{N/n}] [\etq{SA/PSNO}
borracho\etq{A/n}]]]: \sig{Pedro pidió el postre estando borracho}. b.
[\etq{SF/O} [\etq{SN/Suj} Pedro\etq{N/n}] [\etq{SV/Pred} pidió\etq{V/n}
[\etq{SN/COD} un\etq{D/Esp} postre borracho\etq{N/n}]]]:\footnote{\forma{Postre
borracho} en 3b funciona como un sustantivo compuesto.} \sig{Pedro pidió un
postre llamado «postre borracho»}.
\end{solucion}

\begin{solucion}{\ref{ej:09-vinas}}
Era una palabra [\etq{PSO} divertida]. / La ballena me salió [\etq{PSO}
estupenda]. / Parecía [\etq{PSO} el comienzo de una firma ampulosa, una «Q»
inflada, potente]. / El ornitorrinco provocó asombro: era [\etq{PSO} un bicho
ridículo y conmovedor]. / Sentía la mano [\etq{POO} ligera]. / Yo era [\etq{PSO}
eficaz], sabía, podía dibujar\ldots{} la voz me salía [\etq{PSO} pastosa], la
matizaba.
\end{solucion}

% =====================================================================
%  capitulos/cap10.tex — Capítulo 10. Las construcciones pronominales
%  Autoras: Mabel Giammatteo y Mariana Cuñarro
%  Fuente: Capitulo_10_normalizado.docx
% =====================================================================
\capitulo{Las construcciones pronominales}{Mabel Giammatteo y Mariana Cuñarro}

\section{Introducción}

En este capítulo nos ocupamos de los casos en que verbos como \forma{lavar,
divertir, abrazar} o \forma{acercar} pueden también emplearse en forma
pronominal ---\forma{lavarse, divertirse, abrazarse} o \forma{acercarse}---, lo
cual implica que el verbo y el pronombre deben coincidir en número y persona.
Para ello, en primer lugar, revisamos el funcionamiento de los pronombres
personales a fin de distinguir las características de aquellos que intervienen
en estas construcciones. Luego consideramos los dos grandes grupos en que se
dividen las construcciones pronominales: las \termino{reflejas}, en las que el
pronombre manifiesta un argumento del verbo, por lo que tiene papel temático y
desempeña la función sintáctica de COD o de COI, y las
\termino{cuasirreflejas}, en las que, aunque el pronombre comparte los rasgos
flexivos de número y persona del verbo, no representa un papel temático ni tiene
función sintáctica. No obstante, en todos los casos, la incorporación del
pronombre altera en diferentes sentidos el significado del verbo, ya sea porque
el evento que describe revierte sobre quien o quienes lo realizan (\forma{me
lavo, nos abrazamos}), o bien porque se producen cambios en la estructura
argumental (\forma{El viento voló los techos} frente a \forma{Los techos se
volaron}), la estructura eventiva (\forma{Comió una manzana} frente a \forma{Se
comió una manzana}) o la estructura sintáctica (\forma{Olvidé la clave} frente a
\forma{Me olvidé de la clave}).

\section{Los pronombres personales}\label{sec:10-pronombres}

Los pronombres personales son \termino{deícticos}: su referencia es variable de
acuerdo con los participantes del coloquio. «La primera persona se refiere al
hablante o a los hablantes, la segunda, al oyente o a los oyentes, y la tercera,
a las personas o cosas de las que se habla» \parencite[§\,4.4a]{rae2025}.

\begin{destacado}
En español, además de la categoría lingüística de \termino{persona} (primera,
segunda y tercera), los pronombres personales manifiestan \termino{número}
(singular y plural), \termino{género} (masculino y femenino) y \termino{caso}
(nominativo o caso del sujeto; acusativo o caso del OD; dativo o caso del OI, y
preposicional o terminal, que es el caso del pronombre cuando sigue a una
preposición).
\end{destacado}

\begin{table}[htbp]
\centering\footnotesize\setlength{\tabcolsep}{3pt}
\caption{Formas pronominales del español}\label{tab:10-pronombres}
\resizebox{\linewidth}{!}{%
\begin{tabular}{@{}l l c c c c c c c c@{}}
  \toprule
  & & \multicolumn{2}{c}{\textbf{Primera}} & \multicolumn{4}{c}{\textbf{Segunda}} &
    \multicolumn{2}{c}{\textbf{Tercera}} \\
  \cmidrule(lr){3-4}\cmidrule(lr){5-8}\cmidrule(l){9-10}
  & & & & \multicolumn{2}{c}{$-$cortesía\textsuperscript{a}} &
    \multicolumn{2}{c}{+cortesía} & & \\
  \cmidrule(lr){5-6}\cmidrule(lr){7-8}
  \textbf{Caso} & & Sg. & Pl. & Sg. & Pl. & Sg. & Pl. & Sg. & Pl. \\
  \midrule
  Nominativo & & yo & nosotros/-as & vos / tú & vosotros/-as\textsuperscript{b} &
    usted & ustedes & él / ella & ellos / ellas \\
  \addlinespace
  Objetivo\textsuperscript{c} & oblicuas: acusativo & \multirow{2}{*}{me} &
    \multirow{2}{*}{nos} & \multirow{2}{*}{te} & \multirow{2}{*}{os} & lo / la &
    los / las & lo / la & los / las \\
  & oblicuas: dativo & & & & & le & les & le & les \\
  & reflexiva & & & & & \multicolumn{2}{c}{se} & \multicolumn{2}{c}{se} \\
  \addlinespace
  Preposicional & oblicuas & mí & nosotros/-as & vos / ti & vosotros/-as & usted
    & ustedes & él / ella & ellos / ellas \\
  & reflexiva & & & & & \multicolumn{2}{c}{sí} & \multicolumn{2}{c}{sí} \\
  \bottomrule
\end{tabular}}

\medskip
\begin{minipage}{\linewidth}\scriptsize
\textsuperscript{a} Como señala \textcite[169]{ditullio2014}, además de las
categorías de persona, número, género y caso, hay una dimensión adicional,
presente en «los pronombres que designan al destinatario, la cortesía, que opone
las formas de tratamiento \textit{vos} o \textit{tú} (de acuerdo con el dialecto)
a \textit{usted} y, en español peninsular, \textit{vosotros} a \textit{ustedes}».
\textsuperscript{b} Las formas \textit{vosotros/-as} y \textit{os} se usan solo
en la variedad peninsular. \textit{Tú} se emplea en España y en muchas regiones
de Latinoamérica. \textit{Vos} es la forma típica de la variedad rioplatense,
pero también se utiliza en otras regiones de América, aunque con distinto grado
de aceptación social. \textsuperscript{c} Con la denominación de caso objetivo
reunimos al acusativo y al dativo, que solo tienen manifestación diferenciada en
la tercera persona, singular y plural, y también en las formas de segunda
persona \textit{usted} y \textit{ustedes}.
\end{minipage}
\end{table}

\subsection{Pronombres tónicos y pronombres átonos o clíticos}\label{sec:10-tonicos}

Los pronombres personales se pueden agrupar en dos series:

\begin{itemize}
  \item \termino{Tónicos}: llevan acento fuerte y son todos los que están en
        caso nominativo o preposicional.
  \item \termino{Átonos} o \termino{clíticos}: no tienen acento propio y se
        apoyan acentualmente en otra palabra a la que se anteponen
        (\termino{proclíticos}: \forma{Juan se cansó; ¿Me ayudás?}) o a la que
        se posponen (\termino{enclíticos}: \forma{Juan suele cansarse con
        facilidad; ¡Dame eso ya!}). Están en caso acusativo o dativo, y la
        primera y la segunda persona emplean la misma forma para ambos casos.
\end{itemize}

\begin{exe}
  \ex Después de dos horas de espera, el médico [\etq{Pr/COD} \textit{me}]
      atendió.
  \ex {}[\etq{Pr/COI} \textit{Me}] recetó unas pastillas para el dolor de
      cabeza.
\end{exe}

En la tercera persona, en cambio, las formas de acusativo y dativo difieren: se
emplean \forma{lo, la, los, las} para acusativo, y \forma{le} o \forma{les} para
dativo.

\begin{exe}
  \ex {}[\etq{Pr/COD} \textit{Lo}] encontré esta mañana en el bar de enfrente.
  \ex {}[\etq{Pr/COI} \textit{Le}] habían dado un premio en la Academia.
\end{exe}

\begin{ejercicio}\label{ej:10-caso}
En las siguientes oraciones señale el caso de los pronombres y justifique su
respuesta indicando la función sintáctica que cumplen, según se muestra en el
ejemplo resuelto~1.

\begin{exe}
  \ex \textit{Ud.} no sabe nada de lo sucedido.
  \ex ¿Cuántas veces \textit{lo} tengo que repetir?
  \ex Varios amigos \textit{me} comentaron del asunto.
  \ex Mis padres \textit{me} regalaron un coche nuevo.
  \ex \textit{Te lo} dije varias veces, pero no \textit{me} entendiste.
\end{exe}

\begin{resuelto}
1. \forma{Ud.}: caso nominativo, función sintáctica sujeto.
\end{resuelto}
\end{ejercicio}

\subsection{Clíticos oblicuos y reflexivos}\label{sec:10-oblicuos}

Cuando los clíticos difieren de la persona señalada por el sujeto ---o por la
flexión verbal, si el sujeto es tácito--- se dice que son \termino{oblicuos}
\exref{ex:10-medico}; en cambio, cuando son correferenciales con el sujeto, el
clítico es \termino{reflexivo} \exref{ex:10-yomisma}.

\begin{exe}
  \ex\label{ex:10-medico} {}[\etq{3.ª pers.} El médico] [\etq{1.ª pers.}
      \textit{me}] examinó las mamas cuidadosamente.
  \ex\label{ex:10-yomisma} {}[\etq{1.ª pers.} Yo misma] [\etq{1.ª pers.}
      \textit{me}] examiné las mamas cuidadosamente.
\end{exe}

Para que el pronombre sea reflexivo no solo debe coincidir con la persona del
sujeto, sino que ambos ---sujeto y objeto--- deben remitir al mismo referente
extralingüístico. En \exref{ex:10-paciente1}, aunque sujeto y objeto refieren a
la tercera persona, no hay reflexividad porque \forma{el médico} y \forma{la
paciente} son dos personas (es decir, referentes) distintas. Para indicar
correferencia en tercera persona se utiliza la forma reflexiva del pronombre:
\forma{se} \exref{ex:10-paciente2}.

\begin{exe}
  \ex\label{ex:10-paciente1} La paciente llegó al consultorio y [\etq{SN/Suj
      3.ª pers.} el médico] [\etq{Pr Oblic/COD 3.ª pers.} \textit{la}] examinó
      cuidadosamente.
  \ex\label{ex:10-paciente2} {}[\etq{SN/Suj 3.ª pers.} La paciente] [\etq{Pr
      Refl/COI 3.ª pers.} \textit{se}] examinó las mamas cuidadosamente.
\end{exe}

En el SV, los pronombres átonos representan a los argumentos internos del verbo;
por tanto, tienen papel temático en la estructura argumental (EA) y,
sintácticamente, cumplen las funciones de COD, \exref{ex:10-paciente1} y
\exref{ex:10-testigo}, o de COI, \exref{ex:10-paciente2} y
\exref{ex:10-vacuna}.

\begin{exe}
  \ex\label{ex:10-testigo} El testigo [\etq{Pr Refl/COD} \textit{se}] inculpó
      por el crimen.
  \ex\label{ex:10-vacuna} El investigador [\etq{Pr Refl/COI} \textit{se}]
      inyectó [\etq{SN/COD} la vacuna] en su propio brazo.
\end{exe}

\begin{ejercicio}\label{ej:10-funcion}
Determine la función sintáctica y el papel temático de los pronombres en
negrita, según se muestra en el ejemplo resuelto~1.

\begin{exe}
  \ex El jugo de naranja \textbf{le} desagrada.
  \ex Elena \textbf{le} dijo que esta vez llegaría temprano.
  \ex Elena \textbf{lo} dijo una vez más.
  \ex Elena \textbf{le} dijo a su hermano la verdad.
  \ex\label{ex:10-selo} Elena \textbf{se lo} dijo a su hermano.
\end{exe}

\begin{resuelto}
1. \forma{Le}: función sintáctica COI; papel temático experimentante.
\end{resuelto}
\end{ejercicio}

\subsection{El clítico \textit{se} como alomorfo de \textit{le}}\label{sec:10-alomorfo}

A partir de los resultados del ejercicio anterior, ¿cómo se puede explicar lo que
sucede en el ejemplo~5, en el que el pronombre que precede a \forma{lo}, con
función de COI, es \forma{se} en vez de \forma{le}? En este caso, \forma{se} no
es correferente del sujeto ---no remite a \forma{Elena}---, sino al COI \forma{a
su hermano} y, por lo tanto, no es un pronombre reflexivo, ya que se trata de
una variante de realización del clítico \forma{le}, es decir, de un
alomorfo\footnote{Este concepto se presentó en el \capref{4},
§\ref{sec:04-alomorfia}.}. Cuando \forma{le} precede a un clítico acusativo de
tercera persona (\forma{lo, la, los, las}), en su lugar aparece la forma
\forma{se}, que reemplaza a \forma{le} cuando tanto el COD como el COI son
pronominales; de modo que decimos \forma{Le di un regalo esta mañana}, pero
\forma{Se lo di esta mañana} y no \juicio{*}{\forma{Le lo di esta
mañana}}.\footnote{Según \textcite[173, n.~9]{ditullio2014}, «No se trata de un
condicionamiento fonológico, como lo prueba la existencia de secuencias como
\textit{le lavé la ropa}, \textit{le lotearon el terreno}». Se trata, pues, de un
condicionamiento morfológico propio de la concatenación entre los morfemas
pronominalizados de OD y OI.\pendiente{El original cita «Di Tullio (2010, 173,
nota 9)»}}

\section{Las construcciones pronominales}\label{sec:10-construcciones}

La aparición de los pronombres personales clíticos en la sintaxis no es un
fenómeno transparente y unívoco. Por el contrario, se trata de uno de los temas
más complejos y debatidos en las gramáticas. En primer lugar, puede haber más de
una forma para una misma función, como vimos en el ejercicio anterior, en el que
el dativo de tercera persona presenta dos variantes, \forma{se} y \forma{le},
cuya aparición depende del contexto oracional. En segundo lugar, una forma
pronominal no siempre se corresponde con una única función sintáctica, ya que,
como también vimos, puede manifestar caso acusativo o dativo, y cumplir función
de COD o de COI, según la oración. En tercer lugar, según veremos un poco más
adelante (§\ref{sec:10-cuasirreflejas}), los clíticos no siempre tienen función
sintáctica ni desempeñan papel temático.

\subsection{Prueba de sustitución}\label{sec:10-sustitucion}

Para reconocer la función sintáctica de los clíticos de primera y segunda
persona, que no poseen formas diferenciadas para acusativo y
dativo,\footnote{Para esta prueba se tienen en cuenta los datos propios de la
variedad del español rioplatense, en la que la pronominalización distingue entre
las formas \forma{lo, la, los, las} para el COD y las formas \forma{le, les}
para el COI. Existen variedades leístas, en las que se utilizan \forma{le, les}
tanto para el COI como para el COD (\forma{Le veo con frecuencia; Ese cuadro le
compramos el año pasado}). Asimismo, en las variedades laístas de ciertas zonas
del centro y norte de España, ocurre el fenómeno inverso: se utilizan las formas
\forma{la, las} para el COI femenino (\forma{A Luisa la dije que viniera; A la
pared la dieron una mano de pintura}) y, en las regiones loístas, se emplea
\forma{lo, los} para el COI masculino (\forma{A sus socios los mintió; Al vino
lo echan de todo}) \parencite{bosque2019}.} se puede aplicar una
\termino{prueba de sustitución}, mediante la cual se los reemplaza por un
pronombre de tercera persona, que tiene formas diferenciadas para acusativo y
dativo. Si al aplicar la prueba se cambia el clítico por \forma{la(s)} o
\forma{lo(s)}, se trata de un COD; en cambio, si se reemplaza por \forma{le(s)},
es un COI.

\begin{exe}
  \ex \textit{Me} encontró dormida. $\rightarrow$ \textit{La / *Le} encontró
      dormida. $\rightarrow$ \forma{Me} es COD.
  \ex \textit{Te} mostraron las fotos del viaje. $\rightarrow$ \textit{Le / *Lo}
      mostraron las fotos del viaje. $\rightarrow$ \forma{Te} es
      COI.\pendiente{El original dice «Me es COI»}
\end{exe}

\begin{ejercicio}\label{ej:10-sustitucion}
Según se muestra en el ejemplo resuelto~1:

\begin{consignas}
  \item aplique la prueba de sustitución para determinar qué función cumplen
        los pronombres en cursiva;
  \item luego, explique el resultado obtenido.
\end{consignas}

\begin{exe}
  \ex \textit{Me} arreglé frente al espejo.
  \ex \textit{Me} arreglé el pelo frente al espejo.
  \ex \textit{Se} vistió de gala.
  \ex ¿\textit{Te} ayudaron con la tarea?
  \ex \textit{Nos} dimos la vacuna para viajar.
\end{exe}

\begin{resuelto}
1. a) Sustitución por pronombre de tercera persona singular: \forma{Lo arreglé
frente al espejo}. b) El pronombre manifiesta caso acusativo; por lo tanto, su
función sintáctica es COD.
\end{resuelto}
\end{ejercicio}

\begin{ejercicio}\label{ej:10-papeles}
En los siguientes ejemplos analice el papel temático y la función sintáctica de
los pronombres en cursiva, según se muestra en el ejemplo resuelto~1.

\begin{exe}
  \ex \textit{Me} encontró dormida y \textit{te lo} contó.
  \ex \textit{Te} vimos cuando \textit{nos} comprabas los helados.
  \ex \textit{Se lo} llevé a su casa.
  \ex \textit{Me la} entregó en mi casa.
  \ex Los chicos \textit{se las} regalaron a sus amiguitos.
\end{exe}

\begin{resuelto}
1. \forma{Me}: paciente, COD; \forma{te}: meta, COI; \forma{lo}: tema, COD.
\end{resuelto}
\end{ejercicio}

\subsection{Transitividad y reflexividad}\label{sec:10-transitividad}

Como se planteó en el \capref{7}, la base de la \termino{transitividad} reside en
un sujeto que «actúa» (física o verbalmente) sobre un objeto, que siempre
representa una entidad diferente del sujeto.

\begin{exe}
  \ex {}[\etq{SN/Suj} El pastelero] preparó [\etq{SN/COD} varios postres] para la
      boda.
  \ex {}[\etq{SN/Suj} El automóvil] embistió [\etq{SN/COD} al peatón].
  \ex {}[\etq{SN/Suj} Los camilleros] trasladaron [\etq{SN/COD} al paciente].
\end{exe}

En \exref{ex:10-refl1} y \exref{ex:10-refl2} (que retoman \exref{ex:10-yomisma} y
\exref{ex:10-paciente2}), el panorama respecto de la transitividad cambia: sujeto
y objeto no solo coinciden en la persona gramatical, sino que corresponden al
mismo referente extralingüístico. Cuando la entidad a la que remite el objeto es
la misma que la del sujeto, hablamos de \termino{reflexividad}.

\begin{exe}
  \ex\label{ex:10-refl1} {}[\etq{1.ª pers.} Yo misma] [\etq{Pr Refl 1.ª pers.}
      \textit{me}] examiné las mamas cuidadosamente.
  \ex\label{ex:10-refl2} {}[\etq{3.ª pers.} La paciente] [\etq{Pr Refl 3.ª
      pers.} \textit{se}] examinó \textit{a sí misma} las mamas cuidadosamente.
\end{exe}

Además del pronombre reflexivo, las oraciones anteriores incluyen un COD que no
remite a alguien diferente, sino que es parte (\forma{las mamas}) del mismo
sujeto (\forma{yo} en \exref{ex:10-refl1}; \forma{la paciente} en
\exref{ex:10-refl2}), es decir, se trata de lo que se denomina una
\termino{posesión inalienable}.\footnote{\label{fn:10-inalienable}Posesión
inalienable es la que no puede ser separada de quien la posee porque es parte
integrante de él o de ella, como las partes del cuerpo o incluso, en algunos
casos, las pertenencias de su esfera personal.} En \exref{ex:10-refl1}, la
reflexividad se refuerza en el sujeto por el agregado del modificador enfático
\forma{misma} y, en \exref{ex:10-refl2}, mediante la forma \forma{a sí misma},
que duplica al COI.\pendiente{El original dice «a mí misma»}

\begin{recuadro}[Para ampliar: ¿qué sucede en las oraciones reflexivas con la transitividad?]
Aunque la forma transitiva se mantiene en la conformación estructural de estas
oraciones, que formalmente tienen un sujeto y un COD, ambos remiten a la misma
entidad del mundo extralingüístico, ya que el sujeto interviene en un evento que
revierte sobre sí mismo y no sobre un objeto externo a él. Por tanto, según
\textcite[499--500]{morenocabrera1991}, la reflexividad es «[u]no de los
procesos típicos que conduce a la intransitividad. La situación reflexiva típica
es aquella en la que el sujeto y el objeto son referencialmente idénticos
[\ldots] de modo que, aunque la oración sea sintácticamente transitiva, desde el
punto de vista semántico, se considera que es intransitiva».
\end{recuadro}

\section{Construcciones pronominales reflejas}\label{sec:10-reflejas}

Las \termino{construcciones pronominales} están formadas por un verbo y un
pronombre no oblicuo que coinciden en los rasgos flexivos de número y persona.
Dentro de estas construcciones, se consideran \termino{reflejas} aquellas en las
que, además de que la referencia del pronombre coincide con la del sujeto, el
pronombre representa uno de los argumentos del verbo y tiene la función
sintáctica de COD o de COI, según los casos, como ya hemos visto. Entre las
reflejas se distinguen dos grupos principales: \termino{reflexivas} y
\termino{recíprocas}. En las reflexivas, el sujeto interviene en un evento que
revierte sobre sí mismo (\forma{bañarse, peinarse}) o sobre alguna parte de sí
(\forma{cortarse las uñas, peinarse el flequillo}). En las recíprocas, el sujeto
es siempre plural y el evento implica una realización cruzada entre los
participantes (\forma{saludarse, darse las manos, abrazarse}). A continuación,
veremos las principales características de los distintos subgrupos.

\subsection{Construcciones pronominales reflexivas}\label{sec:10-reflexivas}

Dentro de las construcciones pronominales reflexivas distinguimos tres
subgrupos: reflexivas puras, reflexivas accidentales y reflexivas factitivas.

\subsubsection{Reflexivas puras}\label{sec:10-puras}

Son las reflejas más típicas y ocurren en todas las personas del paradigma
verbal. En ellas, el evento que realiza el sujeto recae sobre sí mismo
(\forma{bañarse, cortarse las uñas, rascarse}). El pronombre cumple la función
sintáctica de COD \exref{ex:10-marta} o de COI \exref{ex:10-sofia}. Como vimos en
\exref{ex:10-refl2}, como prueba de reconocimiento se puede incorporar el
sintagma con valor reflexivo \forma{a sí mismo/a}, que duplica el COD o el COI y
confirma que el pronombre es correferencial con el sujeto.

\begin{exe}
  \ex\label{ex:10-marta}
    \begin{xlist}
      \ex Marta \textit{se} baña con jabón biodegradable.
      \ex {}[\etq{SC=SF/O} [\etq{SN/Suj} Marta] [\etq{SV/Pred} [\etq{Pr Refl/COD}
          \textit{se}] baña\etq{V/n} [\etq{SN/COD} \textit{a sí misma}]
          [\etq{SP/ACI} con jabón biodegradable]]]
    \end{xlist}
  \ex\label{ex:10-sofia}
    \begin{xlist}
      \ex Sofía \textit{se} cortó las uñas esta mañana.
      \ex {}[\etq{SC=SF/O} [\etq{SN/Suj} Sofía] [\etq{SV/Pred} [\etq{Pr Refl/COI}
          \textit{se}] cortó\etq{V/n} [\etq{SN/COD} las uñas] [\etq{SN/COI}
          \textit{a sí misma}]]]
    \end{xlist}
\end{exe}

Típicamente, las reflexivas puras se construyen con verbos de acción que tienen
sujeto agente, \exref{ex:10-marta} y \exref{ex:10-sofia}, pero también admiten
verbos con sujeto experimentante, \exref{ex:10-adolescentes} y
\exref{ex:10-genio}:

\begin{exe}
  \ex\label{ex:10-adolescentes} Los adolescentes aún no se conocen bien a sí
      mismos.
  \ex\label{ex:10-genio} Ese escritor se cree un genio.
\end{exe}

\subsubsection{Reflexivas accidentales}\label{sec:10-accidentales}

En las reflejas que hemos estado considerando se presentan acciones realizadas
deliberadamente por los sujetos, que mayormente cumplen el papel temático de
agente. Pero veamos ahora qué sucede en \exref{ex:10-dedo}: ¿se trata también
de una acción deliberada?

\begin{exe}
  \ex\label{ex:10-dedo} Sofía \textit{se} cortó el dedo con un cuchillo.
\end{exe}

El ejemplo de \exref{ex:10-dedo} puede tener dos interpretaciones: la
«masoquista» o intencional, en la que Sofía se cortó un dedo voluntariamente, o
aquella en que el suceso le ocurrió en forma accidental. En el primer caso,
tenemos una construcción reflexiva pura, en la que Sofía es el agente que
deliberadamente realiza una acción que recae sobre ella misma. Sin embargo, la
interpretación «más natural» de esta oración es la segunda, en la que el corte
se ha producido de manera accidental, de modo que Sofía no interviene como
agente ejecutor, sino que, dado que el evento le sucede de modo involuntario, le
corresponde el papel temático de paciente.

Las construcciones que aceptan una interpretación de este tipo se conocen como
\termino{reflexivas accidentales} y presentan semejanzas y diferencias respecto
de las reflexivas puras. Al igual que las reflexivas puras
\exref{ex:10-paradigma-a}, las accidentales son paradigmáticas (es decir, de
toda persona) \exref{ex:10-paradigma-b} y, además de la construcción reflexiva
con un clítico de persona coincidente con el sujeto, también admiten clíticos
oblicuos (no coincidentes) con interpretación de acción deliberada o accidental
\exref{ex:10-oblicuo}.

\begin{exe}
  \ex\label{ex:10-paradigma}
    \begin{xlist}
      \ex\label{ex:10-paradigma-a} Reflexiva pura: \forma{Me peino. / Te peinás.
          / Se peina(n). / Nos peinamos\ldots}
      \ex\label{ex:10-paradigma-b} Reflexiva accidental: \forma{Me lastimé. / Te
          lastimaste. / Se lastimó (lastimaron). / Nos lastimamos\ldots}
    \end{xlist}
  \ex\label{ex:10-oblicuo}
    \begin{xlist}
      \ex Reflexiva pura: \forma{Me peino} (clítico coincidente con el sujeto) /
          \forma{Lo peino} (clítico no coincidente).
      \ex Reflexiva accidental: \forma{Me lastimé} (clítico coincidente) /
          \forma{Lo lastimé en la pelea} (clítico no coincidente con
          interpretación deliberada) / \forma{Lo lastimé sin darme cuenta}
          (clítico no coincidente con interpretación accidental).
    \end{xlist}
\end{exe}

Sin embargo, las accidentales también presentan diferencias respecto de las
reflexivas puras, algunas de las cuales las acercan a las cuasirreflejas medias,
que veremos más adelante (v.~§\ref{sec:10-medias}), como, por ejemplo, no tener
un sujeto agente ni admitir los refuerzos típicos de las reflexivas puras, como
\forma{a mí mismo/a, a vos mismo/a}, etc. Las reflexivas accidentales admiten
refuerzos del tipo de \forma{involuntariamente, accidentalmente} o \forma{por
casualidad}.

\begin{exe}
  \ex
    \begin{xlist}
      \ex {}[\etq{agente} Iván] se quebró (a sí mismo) una pierna para no ir a
          la guerra.\footnote{En este ejemplo, la presencia de un adjunto
          circunstancial de fin es un claro diagnóstico de la agentividad de
          Iván.}
      \ex\label{ex:10-ivan-b} {}[\etq{paciente} Iván] se quebró una pierna en el
          partido de rugby (*a sí mismo / accidentalmente).
    \end{xlist}
\end{exe}

Según vemos en \exref{ex:10-ivan-b}, aunque se acercan a las cuasirreflejas
medias, que dan cuenta de procesos no agentivos (\forma{El árbol se quebró por la
tormenta / solo}), las reflexivas accidentales no se intransitivizan, sino que
pueden presentar un COD referido a la parte del sujeto en que se produjo la
lesión, de modo que el papel temático que le corresponde es el de locativo, como
si dijéramos \sig{Sofía tuvo un corte en el dedo} o \sig{Juan sufrió una
fractura en la pierna}.

El carácter de reflexivas de las accidentales se manifiesta en que la acción,
aunque involuntaria, revierte sobre el sujeto y el clítico tiene función
sintáctica de COD \exref{ex:10-espina-a} o de COI \exref{ex:10-espina-b}.

\begin{exe}
  \ex
    \begin{xlist}
      \ex\label{ex:10-espina-a} Se pinchó con una espina de cactus.
      \ex\label{ex:10-espina-b} Se pinchó el dedo con una espina de cactus.
    \end{xlist}
\end{exe}

\subsubsection{Reflexivas factitivas}\label{sec:10-factitivas}

Comparemos estas dos oraciones:

\begin{exe}
  \ex\label{ex:10-selfi} Ayer \textit{me} saqué una selfi.
  \ex\label{ex:10-foto} El martes fui a un estudio fotográfico y \textit{me}
      saqué una foto para el registro de conducir.
\end{exe}

Aunque en ambas oraciones encontramos la construcción \forma{me saqué} referida a
una fotografía, solo la primera permite la inclusión del refuerzo \forma{a mí
mismo/a}; en la segunda, en cambio, es obvio que, si alguien va a un estudio
fotográfico, es para que otra persona le saque la foto. ¿Cómo podemos explicar
estas diferencias? En \exref{ex:10-selfi} encontramos las características de las
reflexivas puras que ya hemos visto en §\ref{sec:10-puras}, pero
\exref{ex:10-foto} implica que el sujeto no es el agente ejecutor de la acción,
que queda a cargo de otro ---la persona que en la casa de fotografía
efectivamente saca la foto---, quien es el verdadero agente, aunque este suele
quedar implícito o representado por un adjunto: \forma{Me saqué una foto con un
fotógrafo famoso}. Y ¿qué papel temático le corresponde al sujeto? Para
\textcite{kovacci1992}, quien promueve que otro realice la acción se considera un
\termino{causante instigador}. Las oraciones como \exref{ex:10-foto} se incluyen
con las reflexivas porque la acción que el causante instiga a otro a realizar
revierte sobre él mismo.

Las construcciones en las que el sujeto es distinto del ejecutor de la acción se
denominan \termino{reflexivas factitivas}. En ellas el pronombre cumple función
de COD o de COI. No admiten el refuerzo \forma{a mí mismo/a}, aunque aceptan una
paráfrasis con \forma{hacer}: \sig{Hice que otro (el fotógrafo) me sacara una
fotografía}.

\begin{exe}
  \ex {}[\etq{SC=SF/O} [\etq{Suj} \textit{pro}] [\etq{SV/Pred} [\etq{SN/ACT} el
      martes] [\etq{Pr Refl/COI} me] saqué\etq{V/n} [\etq{SN/COD} una foto]
      [\etq{SP/ACF} para el registro de conducir]]]
\end{exe}

Las factitivas son frecuentes en el uso cotidiano, tal como muestran los
siguientes ejemplos:

\begin{exe}
  \ex\label{ex:10-barberia} Hoy en la barbería \textit{me} corté el pelo. (=
      \sig{Hice que me cortaran el pelo})
  \ex Mi primo \textit{se} hizo una hermosa casa en un barrio privado de zona
      norte. (= \sig{Hizo que le hicieran / construyeran una hermosa casa\ldots})
  \ex \textit{Me} hice un piercing y ahora me sale pus de la nariz. (= \sig{Hice
      que me hicieran un piercing})
  \ex\label{ex:10-lorena} Lorena \textit{se} vacunó contra el sarampión esta
      mañana. (= \sig{Hizo que la vacunaran\ldots})
\end{exe}

\begin{ejercicio}\label{ej:10-factitivas}
En las oraciones anteriores, \exref{ex:10-barberia} a \exref{ex:10-lorena},
identifique el agente y el causante instigador. Señale la función sintáctica que
cumple el pronombre, según se muestra en el ejemplo resuelto de abajo.

\begin{resuelto}
\exref{ex:10-barberia}: agente: el barbero; causante instigador: el sujeto
tácito \forma{yo}. El pronombre \forma{me} es COI.\pendiente{El original
intercambia agente y causante instigador, también en el solucionario}
\end{resuelto}
\end{ejercicio}

\subsection{Construcciones pronominales recíprocas}\label{sec:10-reciprocas}

Las construcciones recíprocas responden a la caracterización de las reflejas: el
pronombre es correferencial con el sujeto y cumple función sintáctica de COD
\exref{ex:10-matias} o de COI \exref{ex:10-novios}.

\begin{exe}
  \ex\label{ex:10-matias} Matías y Francisco \textit{se} esperaron a la salida
      del colegio.
  \ex\label{ex:10-novios} Durante la ceremonia, los novios \textit{se} pusieron
      los anillos.
\end{exe}

Ahora bien, ¿qué diferencia estas construcciones de las reflexivas? Aquí también
el evento «transita» de un sujeto a un objeto que remite al mismo referente, pero
se diferencian porque se trata de acciones cruzadas: en \exref{ex:10-matias},
Matías esperó a Francisco y, recíprocamente, Francisco esperó a Matías, y, en
\exref{ex:10-novios}, la novia le puso el anillo al novio mientras él, a su vez,
se lo colocaba a ella. Los eventos recíprocos ocurren simultáneamente (o casi),
ya que la acción realizada por cada agente revierte en otro, que hace las veces
de paciente \exref{ex:10-recip-a} o meta \exref{ex:10-recip-b}.

\begin{exe}
  \ex
    \begin{xlist}
      \ex\label{ex:10-recip-a} {}[\etq{SC=SF/O} [\etq{SN/Suj} Matías y Francisco]
          [\etq{SV/Pred} [\etq{Pr Recipr/COD} se] esperaron\etq{V/n}
          [\etq{SP/ACT} a la salida del colegio]]]
      \ex\label{ex:10-recip-b} {}[\etq{SC=SF/O} [\etq{SN/Suj} los novios]
          [\etq{SV/Pred} [\etq{Pr Recipr/COI} se] pusieron\etq{V/n} [\etq{SN/COD}
          los anillos]]]
    \end{xlist}
\end{exe}

Como se trata de una acción que va de uno a otro, estas construcciones siempre
exigen un sujeto que exprese pluralidad, ya sea morfológicamente
\exref{ex:10-novios}, sintácticamente, mediante la coordinación de SSNN
\exref{ex:10-matias}, o semánticamente, por medio de un sustantivo colectivo
(\forma{La tripulación se abrazaba feliz por haber llegado a destino}).

Para reconocer estas construcciones se pueden aplicar las siguientes pruebas:

\begin{itemize}
  \item agregar adverbios como \forma{mutuamente} o \forma{recíprocamente}, el SP
        \forma{entre sí} o construcciones como \forma{uno al otro}:
\end{itemize}

\begin{exe}
  \ex {}[\etq{SC=SF/O} [\etq{SN/Suj} los novios] [\etq{SV/Pred} [\etq{Pr
      Recipr/COI} se] pusieron\etq{V/n} [\etq{SN/COD} los anillos]
      [\etq{SAdv/ACM} mutuamente]]]\pendiente{El original rotula aquí el
      pronombre como COD}
\end{exe}

\begin{itemize}
  \item dado que el sujeto es plural, también se pueden descomponer sus miembros
        en dos oraciones coordinadas por \forma{y}:
\end{itemize}

\begin{exe}
  \ex Matías esperó a Francisco y Francisco esperó a Matías.
\end{exe}

\begin{itemize}
  \item pronominalizar separadamente los dos OODD:
\end{itemize}

\begin{exe}
  \ex Matías \textit{lo} esperó a Francisco y Francisco \textit{lo} esperó a
      Matías.\pendiente{El original dice «Federico» en estas dos pruebas}
\end{exe}

\begin{destacado}
\textbf{Para tener en cuenta.} Cuando el sujeto es plural, se puede producir
ambigüedad entre la interpretación reflexiva y la recíproca, y será el contexto
el que permita desambiguar. Así, en \forma{Luisa y Mara se miraron en el espejo
antes de entrar en el salón}, podemos interpretar que \sig{cada una se miró a sí
misma en el espejo} (reflexiva pura) o que \sig{a través del espejo, una miró a
la otra, es decir, que se miraron mutuamente} (recíproca). El contexto suele
contener «pistas» que nos permiten desambiguar ambas lecturas:

\begin{exe}
  \ex
    \begin{xlist}
      \ex Luisa y Mara se miraron en el espejo y cada una revisó su peinado.
          (reflexiva pura)
      \ex Luisa y Mara se miraron en el espejo y ambas se sonrieron en forma
          cómplice. (recíproca)
    \end{xlist}
\end{exe}
\end{destacado}

\begin{ejercicio}\label{ej:10-ambiguas}
Teniendo en cuenta que las siguientes oraciones son ambiguas (admiten más de una
interpretación), resuelva lo pedido más abajo.

\begin{exe}
  \ex Ambos luchadores \textit{se torcieron} la muñeca durante la competencia.
  \ex Ana y su hermano \textit{se tiñeron} el pelo.
  \ex Los dos estudiantes \textit{se copiaron} en el examen.
\end{exe}

\begin{consignas}
  \item Provea dos contextos que permitan desambiguarlas, según se muestra en el
        ejemplo resuelto~1.
  \item En 2, además de la reflexiva pura y la recíproca, ¿qué otra
        interpretación es posible?
\end{consignas}

\begin{resuelto}
1. Reflexiva pura: \forma{Uno se torció su propia muñeca al caer y el otro la
suya, contra las sogas}. Recíproca: \forma{Se torcieron uno al otro la muñeca
como resultado de una toma peligrosa en la que ambos participaron}.
\end{resuelto}
\end{ejercicio}

\begin{ejercicio}\label{ej:10-tiporefleja}
Analice los siguientes ejemplos y decida de qué tipo de refleja se trata en cada
caso.

\begin{exe}
  \ex Juanita \textit{se} vistió de rojo para la fiesta y \textit{se} pintó los
      ojos de dorado.
  \ex Silvina decidió tatuar\textit{se} el nombre de su hijo en el brazo.
  \ex \textit{Se} depiló las cejas para emprolijar su mirada.
  \ex El pintor holandés Vincent van Gogh \textit{se} cortó una oreja.
  \ex Hipólito Yrigoyen y Lisandro de la Torre \textit{se} enfrentaron en un
      duelo.
\end{exe}

\begin{resuelto}
1. Ambos pronombres \forma{se} forman construcciones reflexivas puras.
\end{resuelto}
\end{ejercicio}

\begin{ejercicio}\label{ej:10-danna}
Analice todas las construcciones subrayadas en el siguiente fragmento;
identifique las que son pronominales y aporte pruebas que sustenten sus
respuestas, según se muestra en el ejemplo resuelto de abajo.

% Texto periodístico protegido: se incorpora a mano desde el original.
\textoaincorporar{el fragmento adaptado de una nota de \textit{Infobae}
(14/04/2020) sobre Danna Paola, con las construcciones subrayadas: \textit{se
había bañado; se toca el cabello; le arregle las pestañas y las cejas; me estoy
agarrando el pelo; me vean bien el pelo; me las tengo que hacer yo; me las haga;
se comió todas las uñas de la mano derecha}}

\begin{resuelto}
\forma{Se había bañado}: construcción pronominal refleja. El pronombre \forma{se}
es correferencial con el sujeto. Admite el refuerzo \forma{a sí misma}.
\end{resuelto}
\end{ejercicio}

\section{Pronominales sí, pero no reflexivas: las construcciones pronominales cuasirreflejas}\label{sec:10-cuasirreflejas}

No todas las construcciones pronominales son reflexivas: hay un grupo en las que
el verbo y el pronombre coinciden en número y persona, pero no comparten el resto
de las características de las reflexivas. Para ir diferenciando las reflexivas
de las otras, a las que \textcite{bello2002} denominó \termino{cuasirreflejas},
comencemos observando las siguientes frases extraídas a) de una pintada callejera
y b) de un afiche del gobierno de la provincia de Buenos Aires.

\imagenpendiente{a) pintada callejera; b) afiche del gobierno de la provincia de
Buenos Aires}

\begin{exe}
  \ex\label{ex:10-feliz} No sabía qué ponerme y \textit{me} puse feliz.
  \ex\label{ex:10-turnos} \textit{Se} están enviando nuevos turnos de vacunación.
  \ex\label{ex:10-vacunate} Buenos Aires, vacuná\textit{te}.
\end{exe}

Aunque todas las construcciones pronominales están formadas por un verbo y un
pronombre que comparten los rasgos flexivos de número y persona, las reflejas,
que hemos visto en §\ref{sec:10-reflejas}, deben cumplir otros dos requisitos
más: el pronombre debe desempeñar algún papel temático y la función sintáctica
de COD o de COI. En cambio, en las \termino{cuasirreflejas}, como las que
aparecen en las imágenes de arriba, aunque verbo y pronombre comparten los rasgos
de flexión, no cumplen los otros dos requisitos: el pronombre no expresa ningún
papel temático y tampoco tiene función sintáctica. Para
\textcite[§\,759]{bello2002}, a diferencia de las reflexivas, en las que agente y
paciente coinciden, en las cuasirreflejas «gramaticalmente parece decirse que el
sujeto obra en sí mismo», lo que significa que las construcciones cuasirreflejas
refieren a un proceso en el que el sujeto no es el agente que lo realiza; por lo
tanto, los refuerzos que habíamos visto con las reflexivas,
\exref{ex:10-desvesti} y \exref{ex:10-odian}, no se aplican,
\exref{ex:10-despidio} y \exref{ex:10-jactaban}:

\begin{exe}
  \ex\label{ex:10-desvesti} Cuando llegué a mi casa, \textit{me} desvestí (a mí
      misma) y \textit{me} di una ducha (a mí misma).
  \ex\label{ex:10-odian} Mi prima y su hermana \textit{se} odian (mutuamente /
      una a la otra).
  \ex\label{ex:10-despidio} Ya era tarde cuando Karina \textit{se} despidió (*a
      sí misma) de nosotros.
  \ex\label{ex:10-jactaban} Ambos \textit{se} jactaban de sus hazañas (*a sí
      mismos / *mutuamente).
\end{exe}

Las construcciones cuasirreflejas son sumamente complejas porque hay distintos
subgrupos cuyos límites «son inestables y, por decirlo así, porosos»
\parencite[17]{sanchezlopez2002}; por tanto, es difícil clasificarlas siguiendo
criterios que las abarquen a todas. En líneas generales, se trata de
construcciones que producen algún «cambio» en la oración. Pero ¿de qué tipo de
cambio se trata y a qué afecta? En las cuasirreflejas, el clítico no cumple
función sintáctica ni de COD ni de COI y se interpreta como una marca de función
que recibe el nombre de \termino{índice pronominal cuasirreflejo} (IPCR). La
introducción de este pronombre incide en la conformación de alguna de las tres
estructuras del predicado que presentamos en el \capref{6}: la estructura
argumental (EA), la estructura sintáctica (ES) o la estructura eventiva (EE), por
lo que dividiremos estas construcciones en tres grupos:

\begin{itemize}
  \item Construcciones cuasirreflejas que producen cambio en la estructura
        argumental. La introducción del clítico elimina la manifestación de
        algún argumento de la estructura sintáctica. Dentro de estas
        construcciones se encuentran las cuasirreflejas medias, las
        medio-pasivas, las pasivas y las impersonales.
  \item Construcciones cuasirreflejas que producen cambio en la estructura
        sintáctica. El clítico produce una modificación en la forma en que se
        manifiesta la transitividad. Las construcciones de este tipo son las
        cuasirreflejas alternantes y las inherentes.
  \item Construcciones cuasirreflejas que producen cambio en la estructura
        eventiva. La presencia del clítico ocasiona cambios en la aspectualidad.
        Se ubican en este grupo las construcciones cuasirreflejas que delimitan
        o intensifican el evento.
\end{itemize}

\subsection{Construcciones cuasirreflejas que producen cambio en la estructura argumental}\label{sec:10-ea}

Las cuasirreflejas que cambian la estructura argumental comprenden cuatro
subgrupos: medias, medio-pasivas, pasivas e impersonales.

Todas estas construcciones tienen en común que alteran la manifestación
sintáctica del argumento agente. Esto, en las construcciones medias,
medio-pasivas y pasivas, produce la promoción del tema/paciente a la función de
sujeto e intransitiviza la oración. En las impersonales, en cambio, ningún
argumento se promueve como sujeto y, como veremos, no todas las construcciones
se intransitivizan. Otra diferenciación importante entre estos subgrupos es que
solo las construcciones medias son paradigmáticas, es decir, el pronombre puede
aparecer en cualquiera de las personas gramaticales (primera, segunda o tercera)
y tanto en singular como en plural (\forma{Me entristecí / Te entristeciste / Se
entristecieron con las noticias}, etc.). En las medio-pasivas, pasivas e
impersonales, \forma{se} no es paradigmático: está fijado en tercera persona,
singular o plural en las medio-pasivas (\forma{Este error se subsana fácilmente}
/ \forma{Estos errores se subsanan\ldots}) y en las pasivas (\forma{Se presentó
un informe} / \forma{Se presentaron varios informes}), y en tercera persona
singular en las impersonales (\forma{Aquí no se fuma} /
\juicio{*}{\forma{Aquí no se fuman}}). En lo que sigue describiremos las
características de cada subgrupo.

\subsubsection{Construcciones cuasirreflejas medias}\label{sec:10-medias}

En las construcciones de este subgrupo se elimina el agente, lo que motiva la
promoción o ascenso del argumento tema/paciente, que se convierte en sujeto de
la oración.

\begin{exe}
  \ex {}[\etq{agente} Pedro] abrió [\etq{tema} la puerta]. $\rightarrow$
      [\etq{tema} La puerta] se abrió (sola).
\end{exe}

Como en todas las cuasirreflejas, en estas construcciones el pronombre no
desempeña ninguna función sintáctica y se considera una marca de función.
Asimismo, se destaca que el sujeto «sufre un proceso que escapa a su control
efectivo» \parencite[77]{sanchezlopez2002}; de ahí que se considere que estas
oraciones se corresponden con la \termino{voz media}.

\begin{destacado}
Según \textcite[1636]{mendikoetxea1999}, la voz media «se interpreta como la
forma de expresar que la acción que denota el verbo “afecta”, en mayor o menor
grado, al sujeto (incide sobre el sujeto, interesa al sujeto, indica un cambio en
el sujeto, etc.). En otras palabras, las construcciones medias se caracterizan
por tener sujetos “afectados”».
\end{destacado}

Este tipo de construcciones ocurre con verbos transitivos de \sig{cambio de
estado}, que se intransitivizan por la pronominalización. Según el tipo de cambio
involucrado, se reconocen tres subgrupos: con verbos que indican \sig{cambio de
estado físico} (\forma{romperse, secarse, hundirse}), con verbos que indican
\sig{cambio de estado mental} (\forma{entristecerse, asombrarse, divertirse}) y
con verbos que indican \sig{ubicación o posición} (\forma{acercarse, moverse,
levantarse}).

\paragraph{Medias con verbos de cambio de estado físico.} Prestemos atención a
los siguientes pares de oraciones:

\begin{exe}
  \ex\label{ex:10-pared}
    \begin{xlist}
      \ex\label{ex:10-pared-a} El albañil rompió la pared.
      \ex\label{ex:10-pared-b} La pared se rompió.
    \end{xlist}
  \ex\label{ex:10-ropa}
    \begin{xlist}
      \ex\label{ex:10-ropa-a} El viento secó la ropa.
      \ex\label{ex:10-ropa-b} La ropa se secó con el viento.
    \end{xlist}
\end{exe}

Tanto \exref{ex:10-pared-a} como \exref{ex:10-ropa-a} son oraciones transitivas
en las que el primer argumento es un agente (\forma{el albañil}) o una causa
(\forma{el viento}), y el segundo es un tema (\forma{la pared} o \forma{la
ropa}). Pero ¿qué sucede en \exref{ex:10-pared-b} y \exref{ex:10-ropa-b}? El tema
ha sido promovido a la función más prominente en la estructura oracional, la de
sujeto, lo que en \exref{ex:10-pared-b} produce la eliminación total del agente
(\juicio{*}{\forma{La pared se rompió por el albañil}}) o, como sucede en
\exref{ex:10-ropa-b}, la causa es relegada a una posición marginal, en la que
cumple la función de adjunto circunstancial.

¿Qué efecto producen estos cambios? En \exref{ex:10-belgrano} se muestra el
contraste entre una oración activa, de lectura agentiva \exref{ex:10-belgrano-a},
y una oración media \exref{ex:10-belgrano-b}, en la que la eliminación del agente
(o su democión a una función marginal) promueve una \termino{lectura de proceso}
que, en algunos casos, permite la introducción del refuerzo \forma{solo/a} o
\forma{por sí solo/a}.

\begin{exe}
  \ex\label{ex:10-belgrano}
    \begin{xlist}
      \ex\label{ex:10-belgrano-a} Los ingleses hundieron el \textit{General
          Belgrano} durante la guerra de Malvinas.\\
          {}[\etq{SC=SF/O} [\etq{SN/Suj/agente} los ingleses] [\etq{SV/Pred}
          hundieron\etq{V/n} [\etq{SN/COD/tema} el \textit{General Belgrano}]
          [\etq{SP/ACT} durante la guerra de Malvinas]]]
      \ex\label{ex:10-belgrano-b} El \textit{Titanic} se hundió (solo) / (por un
          iceberg / porque chocó con un iceberg).\\
          {}[\etq{SC=SF/O} [\etq{SN/Suj/tema} el \textit{Titanic}] [\etq{SV/Pred}
          [\etq{Pr/IPCR} se] hundió\etq{V/n} [\etq{SA/ACM} solo] / [\etq{SP/ACC}
          por un iceberg]]]
    \end{xlist}
\end{exe}

Los verbos que intervienen en estas construcciones se pueden caracterizar como
\sig{verbos de cambio} o que producen un cambio en el objeto. Muchos permiten lo
que se conoce como \termino{alternancia causativa} y ofrecen dos posibles
manifestaciones:

\begin{enumerate}[label=\alph*)]
  \item la \termino{construcción transitiva causativa} típica, en la que el
        cambio en el objeto es producido por un agente, \exref{ex:10-pared-a} y
        \exref{ex:10-belgrano-a}, o por una causa \exref{ex:10-ropa-a};
  \item la \termino{construcción media anticausativa}, en la que el agente se
        elimina \exref{ex:10-pared-b} o se demueve a una función sintáctica
        opcional, \exref{ex:10-ropa-b} y \exref{ex:10-belgrano-b}, y el tema se
        promueve a la función de sujeto. Esta transformación, que elimina al COD,
        produce una intransitivización o inacusativización, por lo que se dice
        que el pronombre cumple función intransitivizadora.
\end{enumerate}

\begin{recuadro}[Para ampliar]
No todos los verbos de cambio transitivos entran en la alternancia causativa.
Para que la alternancia sea posible, el evento tiene que poder ser interpretado
como un proceso que le ocurre al tema sin intervención de un agente. Así, verbos
transitivos típicamente agentivos como \forma{cortar, construir, leer,
asesinar}, por lo general, no admiten la alternancia, ya que siempre requieren un
sujeto agente.

\begin{exe}
  \ex
    \begin{xlist}
      \ex El pastelero cortó el postre en porciones.
      \ex[*]{El postre se cortó (solo) en porciones.}
    \end{xlist}
  \ex
    \begin{xlist}
      \ex Niemeyer construyó los principales edificios de Brasilia.
      \ex[*]{Los principales edificios de Brasilia se construyeron (solos).}
    \end{xlist}
\end{exe}

Por otra parte, los verbos de proceso tampoco se comportan todos igual. Para
explicar estas diferencias, \textcite{mendikoetxea2011} compara \forma{romper} y
\forma{estallar}, dos verbos que implican \sig{la desintegración material de un
objeto}, que representa al tema afectado por el evento. La autora encuentra que,
a pesar de su aparente similitud semántica, las construcciones que ambos verbos
permiten no son las mismas:

\begin{exe}
  \ex
    \begin{xlist}
      \ex La tormenta rompió el vidrio.
      \ex[*]{La tormenta estalló el vidrio.}
    \end{xlist}
  \ex
    \begin{xlist}
      \ex[*]{El vidrio rompió.}
      \ex El vidrio estalló.
    \end{xlist}
  \ex
    \begin{xlist}
      \ex El vidrio se rompió.
      \ex[*]{El vidrio se estalló.}
    \end{xlist}
\end{exe}

¿Cómo se pueden explicar estas diferencias? Si bien tanto \forma{romper} como
\forma{estallar} son verbos de cambio de estado, se diferencian porque pertenecen
a dos subclases distintas, según el evento sea producido por una causa externa o
por una causa interna:

\begin{enumerate}[label=\alph*)]
  \item \termino{Verbos de causa externa}, como \forma{romper, secar, quemar} o
        \forma{abrir}: refieren una acción que tiene una causa que la produce.
        Estos verbos admiten que el tema afectado por el evento aparezca tanto en
        la posición de objeto como en la de sujeto, pero en este segundo caso
        requieren la introducción del clítico \forma{se}:
        \begin{exe}
          \ex
            \begin{xlist}
              \ex Juana secó la ropa.
              \ex La ropa \textit{se} secó.
            \end{xlist}
          \ex
            \begin{xlist}
              \ex El incendio está quemando el bosque.
              \ex El bosque \textit{se} está quemando (debido al incendio).
            \end{xlist}
        \end{exe}
  \item \termino{Verbos de causa interna}, como \forma{estallar, arder} o
        \forma{envejecer}: no involucran causa ni agente externo, es decir que
        nadie o nada ocasiona el evento, sino que este simplemente se da u ocurre
        por un proceso interno, que se produce o tiene lugar en el objeto. Estos
        verbos solo admiten la construcción con el tema en función de sujeto y
        rechazan \forma{se}:
        \begin{exe}
          \ex
            \begin{xlist}
              \ex La leña verde no arde bien.
              \ex[*]{La leña verde no \textit{se} arde bien.}
            \end{xlist}
          \ex
            \begin{xlist}
              \ex Algunas personas envejecen muy bien.
              \ex[*]{Algunas personas \textit{se} envejecen muy bien.}
            \end{xlist}
        \end{exe}
\end{enumerate}
\end{recuadro}

\paragraph{Medias con verbos de cambio de estado mental (verbos psicológicos).}
Los verbos psicológicos, como \forma{entristecerse, enojarse, divertirse,
sorprenderse}, también se incluyen entre los de cambio de estado y señalan un
cambio interno de tipo mental que alguien experimenta ---\forma{la abuela}
\exref{ex:10-abuela-a} y \forma{la población} \exref{ex:10-abuela-b}---.

\begin{exe}
  \ex
    \begin{xlist}
      \ex\label{ex:10-abuela-a} El nieto alegró a su abuela con una sonrisa.\\
          {}[\etq{SC=SF/O} [\etq{SN/Suj/agente} el nieto] [\etq{SV/Pred}
          alegró\etq{V/n} [\etq{SN/COD/experimentante} a su abuela] [\etq{SP/ACI}
          con una sonrisa]]]
      \ex\label{ex:10-abuela-b} Un nuevo virus preocupa a la población.
    \end{xlist}
\end{exe}

Al igual que los verbos de cambio de estado físico, los que indican cambio
psicológico admiten alternancia entre:

\begin{enumerate}[label=\alph*)]
  \item la construcción transitiva, con agente \exref{ex:10-abuela-a} o causa
        \exref{ex:10-abuela-b} y un experimentante;
  \item la construcción media, en la que el agente o la causa se elimina o
        demueve a una función sintáctica opcional de adjunto y el experimentante
        es promovido a la función de sujeto mediante un proceso de
        intransitivización.
\end{enumerate}

\begin{exe}
  \ex
    \begin{xlist}
      \ex La abuela se alegró con la sonrisa de su nieto.\\
          {}[\etq{SC=SF/O} [\etq{SN/Suj/experimentante} la abuela]
          [\etq{SV/Pred} [\etq{Pr/IPCR} se] alegró\etq{V/n} [\etq{SP/ACI} con la
          sonrisa de su nieto]]]
      \ex La población se preocupó por el nuevo virus.
    \end{xlist}
\end{exe}

\paragraph{Medias con verbos de cambio de ubicación o posición.} Dentro del
amplio grupo de verbos de movimiento, en este apartado nos centramos en los que
son transitivos y admiten la alternancia causativa. Más adelante, dentro de las
construcciones cuasirreflejas que producen cambio en la estructura eventiva,
tratamos algunos verbos de movimiento intransitivos que presentan una variante
pronominal (v.~§\ref{sec:10-delimitacion}).

Si bien la mayoría de los autores incluyen estos verbos con los de cambio
físico, los de ubicación o posición presentan características que los
distinguen. Cotejemos los siguientes ejemplos para ver sus particularidades.

\begin{exe}
  \ex
    \begin{xlist}
      \ex\label{ex:10-lobo-a} Los cazadores alejaron al lobo.
      \ex El lobo se alejó.
    \end{xlist}
  \ex\label{ex:10-sento}
    \begin{xlist}
      \ex\label{ex:10-sento-a} La enfermera sentó al paciente.
      \ex\label{ex:10-sento-b} El paciente se sentó por sí solo.
    \end{xlist}
  \ex
    \begin{xlist}
      \ex\label{ex:10-chofer-a} El chofer detuvo el coche.
      \ex El coche se detuvo (porque el chofer paró el motor / porque se paró el
          motor) $\rightarrow$ el motor del automóvil se detuvo.
    \end{xlist}
\end{exe}

En casos como los anteriores, con verbos transitivos, las construcciones
cuasirreflejas con verbos de cambio de ubicación o posición se comportan como las
de cambio físico o mental que ya hemos visto, es decir que alternan entre:

\begin{enumerate}[label=\alph*)]
  \item una construcción transitiva causativa, con sujeto agente (\forma{los
        cazadores, la enfermera, el chofer}) y un objeto tema o paciente
        (\forma{el lobo, el paciente, el coche});
  \item una construcción media anticausativa, formada mediante la inclusión de
        \forma{se}, en la que se elimina al agente, promoviendo al tema o
        paciente a la función de sujeto.
\end{enumerate}

Pero ¿qué particularidad presenta el nuevo sujeto de la construcción pronominal?
¿Es un tema/paciente como en los verbos de cambio de estado? Observemos el
contraste que presenta \exref{ex:10-muneca} respecto de \exref{ex:10-sento}.

\begin{exe}
  \ex\label{ex:10-muneca}
    \begin{xlist}
      \ex\label{ex:10-muneca-a} La niña sentó a su muñeca.
      \ex[?]{\label{ex:10-muneca-b} La muñeca se sentó (sola).}
    \end{xlist}
\end{exe}

La construcción transitiva de \exref{ex:10-muneca-a} es paralela a la de
\exref{ex:10-sento-a}; sin embargo, aunque podríamos pensar en una situación
particular en que \exref{ex:10-muneca-b} fuera aceptable, fuera de contexto
resulta anómala. Para \textcite[200]{kovacci1992}, en estas construcciones,
cuando el sujeto es animado, como en \exref{ex:10-sento-b}, realiza
simultáneamente al agente y al paciente, ya que «[\ldots] dos participantes
pueden ser correferenciales y manifestarse en la misma función». También para
\textcite[1593]{mendikoetxea1999}, dado que estos son verbos de causa externa,
«el sujeto se interpreta como causa y tema»,\footnote{Para
\textcite[495]{martinezrojas2021}, a diferencia de las anticausativas, que
refieren eventos espontáneos con sujeto paciente, las que se forman con verbos
de movimiento del cuerpo (\textit{body motion verbs}) (autocausativas o
endorreflexivas) refieren a «eventos en los que el sujeto animado es agente».}
ya que podemos tener contrastes como \forma{Yo me levanto solo (voz media) y no
espero que otros me levanten (voz activa)} frente a \forma{Yo me levanto solo y
no espero ser levantado (por otros) (voz pasiva)}. Asimismo, estas oraciones
admiten la construcción final con \forma{para}, que se considera un diagnóstico
de agentividad: \forma{Yo me levanto para no seguir escuchando lo que dicen / El
paciente se sentó para comer / El cazador se alejó para perseguir al lobo}.

Y ¿qué encontramos en oraciones como \exref{ex:10-tormenta} a
\exref{ex:10-findeano}, en las que estos verbos se construyen con sujetos no
animados? Se trata de entidades que por sí solas desencadenan un proceso. En
\exref{ex:10-findeano}, aunque el sujeto tácito \forma{nosotros} sea
\sig{animado} y \sig{humano}, no posee voluntariedad para realizar un evento como
\sig{acercarse a fin de año}, que se presenta como un proceso vinculado con el
transcurso temporal. En estos casos, los sujetos solo manifiestan el papel
temático de tema o paciente, según corresponda.

\begin{exe}
  \ex\label{ex:10-tormenta} La tormenta se aproxima a la ciudad.
  \ex La solución se asentó en el fondo del recipiente.
  \ex\label{ex:10-findeano} Nos acercamos a fin de año.
\end{exe}

\begin{destacado}
\textbf{Para tener en cuenta.} Las diferencias presentadas a través de los
ejemplos incluidos en este apartado muestran que los verbos tienen siempre
múltiples posibilidades de construcción y de combinación con diferentes
complementos (incluido el sujeto) y adjuntos y que, por eso, la interpretación de
cada oración debe hacerse teniendo en cuenta todos estos factores.
\end{destacado}

\begin{ejercicio}\label{ej:10-levantar}
Compare las oraciones de los siguientes grupos y resuelva lo pedido, según se
muestra en el ejemplo resuelto~1.

\begin{consignas}
  \item Señale semejanzas y diferencias.
  \item Aplique alguna prueba para corroborar la agentividad de los sujetos.
  \item ¿Cuál es el papel temático de \forma{las sillas} en 1a y en 1c?
  \item ¿Qué refuerzo puede aplicar en 2b?
\end{consignas}

\begin{exe}
  \ex
    \begin{xlist}
      \ex Los asistentes levantaron las sillas en señal de protesta.
      \ex Los asistentes levantaron la mano para hacer preguntas.
      \ex Los asistentes se levantaron de sus sillas para aplaudir de pie al
          cantante.
    \end{xlist}
  \ex
    \begin{xlist}
      \ex El nuevo gobierno extendió los límites de su dominio.
      \ex El incendio se extendió por el bosque nativo.
    \end{xlist}
\end{exe}

\begin{resuelto}
1. Semejanzas entre 1a y 1b: en ambas, \forma{levantar} tiene empleo transitivo
con sujeto agente ---\forma{los asistentes}--- y COD tema ---\forma{las sillas}
(1a) y \forma{la mano} (1b)---, y aceptan la paráfrasis con \forma{hacer}:
\forma{¿Qué hicieron los asistentes?} Diferencia entre 1a y 1b: en 1b el COD es
una posesión inalienable, \forma{la mano}. La diferencia con 1c es que se trata
de una construcción media, en la que el sujeto es a la vez agente y paciente. El
papel temático de \forma{las sillas} en 1a es tema y en 1c es
origen.\pendiente{El original dice «en (b)» en la última oración}
\end{resuelto}
\end{ejercicio}

\begin{ejercicio}\label{ej:10-cambio}
Compare las estructuras de cada par, en las que el verbo aparece con y sin
pronombre. Luego resuelva lo pedido en las consignas, según se muestra en el
ejemplo resuelto~1.

\begin{consignas}
  \item Identifique la clase del verbo según el tipo de cambio que produce (de
        estado físico, de estado mental, de ubicación o posición).
  \item Teniendo en cuenta la estructura argumental de los verbos, trate de
        explicar qué cambios produce la introducción del pronombre.
\end{consignas}

\begin{exe}
  \ex
    \begin{xlist}
      \ex Juan abrió la puerta con la llave.
      \ex Se abrió la puerta con el viento.
    \end{xlist}
  \ex
    \begin{xlist}
      \ex Levanté las cortinas para que entre la luz.
      \ex Me levanto temprano y estudio.
    \end{xlist}
  \ex
    \begin{xlist}
      \ex Acuesto a los niños luego de bañarlos.
      \ex Se acostó mientras los niños dormían.
    \end{xlist}
\end{exe}

\begin{resuelto}
1. Verbo de cambio de estado físico. 1a: construcción transitiva con agente y
tema (\forma{la puerta}). 1b: construcción cuasirrefleja con promoción del tema
(\forma{la puerta}) a sujeto y democión de la causa a ACC (\forma{con el
viento}).
\end{resuelto}
\end{ejercicio}

\subsubsection{Construcciones cuasirreflejas medio-pasivas}\label{sec:10-mediopasivas}

Al igual que en las construcciones medias, también en las medio-pasivas se
demueve o elimina el agente, pero, a diferencia de lo que sucede con el grupo
anterior, en las medio-pasivas el agente eliminado tiene interpretación
indefinida. Se trata de oraciones como las siguientes:

\begin{exe}
  \ex\label{ex:10-problema} El problema de matemáticas se soluciona en dos
      minutos / en dos patadas.
  \ex\label{ex:10-libros} Los libros de divulgación científica se leen
      fácilmente / rápidamente / de corrido.
  \ex\label{ex:10-pastillas} Estas pastillas se toleran bien / sin problemas.
\end{exe}

Una lectura atenta de estas oraciones pone en evidencia que comparten
características con las pasivas (v.~§\ref{sec:10-pasivas}) y con las
impersonales (v.~§\ref{sec:10-impersonales}) con \forma{se}: en primer lugar,
solo se pueden construir en tercera persona, es decir que el único pronombre que
aparece es \forma{se}. En segundo lugar, como las pasivas, admiten concordancia
en singular \exref{ex:10-problema} y en plural, \exref{ex:10-libros} y
\exref{ex:10-pastillas}, y, en ambas construcciones, el tema o paciente se
promueve a la posición de sujeto. También como las pasivas, estas construcciones
implican la eliminación del agente de la correspondiente oración activa, pero se
diferencian porque en las medio-pasivas el agente eliminado tiene interpretación
indefinida: \forma{Este problema se soluciona en dos minutos} $\rightarrow$
\sig{Cualquiera lo soluciona / lo puede solucionar}. Otra característica propia
de estas construcciones es que, a diferencia de las pasivas y de las
impersonales, que son eventivas, las medio-pasivas constituyen predicados de
estado\footnote{La diferencia entre estados y otros tipos de predicados se
desarrolló en el \capref{8}.} que manifiestan una cualidad inherente y
caracterizadora del tema; así, por ejemplo, en \exref{ex:10-problema},
interpretamos que, si el problema se soluciona en dos minutos, se trata de un
problema \sig{fácil} o \sig{solucionable}. Dado que la propiedad referida se
presenta como atemporal, estas oraciones no se construyen en tiempos
perfectivos, que denotan que el evento efectivamente ocurrió, sino en tiempos
imperfectivos (presente o imperfecto), que indican que es posible que el evento
ocurra, de donde se desprende su valor modal\footnote{Sobre la lectura modal de
estas construcciones, v. \capref{14}.} \parencite{sanchezlopez2002}. Tanto los
tiempos imperfectivos como la presencia de adverbios o SSPP modales
(\forma{fácilmente, bien, con dificultad}) son inductores que promueven la
\termino{lectura caracterizadora} atemporal de estas construcciones.

\subsubsection{Construcciones cuasirreflejas pasivas (pasivas cuasirreflejas o pasivas con \textit{se})}\label{sec:10-pasivas}

Observemos las relaciones entre las construcciones de las columnas A, B y C.

\begin{center}\small
\begin{tabularx}{\linewidth}{@{}r L L L@{}}
  \toprule
  & A & B & C \\
  \midrule
  \refstepcounter{exx}(\theexx)\label{ex:10-traductor} & Un traductor novato tradujo el último libro
    premiado en el concurso Planeta. & El último libro premiado en el concurso
    Planeta fue traducido (por un traductor novato). & Se tradujo el último
    libro premiado en el concurso Planeta. \\
  \refstepcounter{exx}(\theexx) & Uno de los abogados filtró a la prensa las declaraciones del imputado. &
    Las declaraciones del imputado fueron filtradas a la prensa (por uno de los
    abogados). & Se filtraron a la prensa las declaraciones del imputado. \\
  \refstepcounter{exx}(\theexx) & Una empresa extranjera construirá una nueva represa en el sur del país. &
    Una nueva represa será construida en el sur del país (por una empresa
    extranjera). & Se construirá una nueva represa en el sur del país. \\
  \bottomrule
\end{tabularx}
\end{center}

Mientras las oraciones de la columna A son todas transitivas, ¿qué sucede con las
de B y C? En todos los casos el verbo sufre un proceso de intransitivización que
promueve el tema o paciente a la posición de sujeto y produce una
\termino{lectura de afectación}, porque el tema o paciente, que es el afectado
por el evento, pasa a ser más prominente. Sin embargo, debido a la
inacusativización de la estructura, el sujeto suele aparecer pospuesto, sobre
todo cuando es indeterminado. En cuanto al agente, en las oraciones de B puede
omitirse o aparecer marginalmente como complemento agente; en cambio, las de C
no admiten su presencia.\footnote{Aunque la presencia del agente es excepcional
con la pasiva con \forma{se}, en algunos textos jurídicos puede aparecer.}
Tradicionalmente, las oraciones de B se conocen como \termino{pasivas
perifrásticas}\footnote{Los conceptos de perífrasis verbal y de verbo auxiliar
se desarrollan en el \capref{11}.} porque el verbo adopta una forma compuesta
por el verbo auxiliar \forma{ser} (en el tiempo correspondiente de la voz activa
y concordado en número y persona con el sujeto tema) seguido del participio del
verbo principal (también concordado con el sujeto, en este caso, en género y
número). En las construcciones de C, que son \termino{pasivas cuasirreflejas},
el verbo está en forma activa y el clítico es un IPCR.

La eliminación sintáctica del agente no debe hacernos pensar que estas oraciones
no tienen un agente implícito que representa a \sig{quien realiza la acción}.
Una oración como \juicio{?}{\forma{La nueva represa se construirá sola}} resulta
extraña porque las pasivas, tanto perifrásticas como cuasirreflejas, involucran
eventos de acción transitivos que son realizados por un agente al que se puede
hacer referencia mediante un adverbio \exref{ex:10-deliberadamente} o una
construcción final \exref{ex:10-convenio}.

\begin{exe}
  \ex\label{ex:10-deliberadamente} Fue adelantada / Se adelantó la fecha de la
      reunión \textit{deliberadamente}. $\rightarrow$ Alguien la adelantó.
  \ex\label{ex:10-convenio} Se garantizó un nuevo convenio de precios \textit{para
      favorecer a los sectores de menores ingresos}. $\rightarrow$ Alguien lo
      garantizó.
\end{exe}

Sin embargo, no existe un paralelismo absoluto entre la pasiva perifrástica y la
pasiva con \forma{se}:

\begin{itemize}
  \item la primera admite todas las personas gramaticales, tanto singulares
        \exref{ex:10-estafado} como plurales \exref{ex:10-convocados}; la
        segunda está restringida a la tercera persona, tanto singular
        \exref{ex:10-prorrogo} como plural \exref{ex:10-prorrogaron}:
\end{itemize}

\begin{exe}
  \ex\label{ex:10-estafado} Fui estafado por uno de mis acreedores. /
      \juicio{*}{\forma{Se me estafó.}}\pendiente{En el español rioplatense,
      \textit{Se me estafó} es una impersonal aceptable}
  \ex\label{ex:10-convocados} Uds. serán convocados en el próximo llamado.
  \ex\label{ex:10-prorrogo} Se prorrogó la convocatoria.
  \ex\label{ex:10-prorrogaron} Se prorrogaron las fechas de examen.
\end{exe}

\begin{itemize}
  \item a diferencia de la pasiva perifrástica, que generalmente expresa eventos
        únicos \exref{ex:10-napoleon}, la pasiva con \forma{se} suele utilizarse
        en oraciones de sentido general \exref{ex:10-oportunidades}:
\end{itemize}

\begin{exe}
  \ex\label{ex:10-napoleon} Napoleón Bonaparte fue derrotado en 1815 en la
      batalla de Waterloo.
  \ex\label{ex:10-oportunidades} Se encuentran más oportunidades con una buena
      educación.
\end{exe}

Como la estructura pasiva exige ambos papeles ---tema y agente---, ni la forma
perifrástica ni la cuasirrefleja admiten verbos intransitivos, ya que los
inacusativos carecen de agente y los inergativos, de tema
\parencite{giammatteo2010}.

\begin{destacado}
\textbf{Para tener en cuenta.} \textcite{ditullio2014} plantea que, cuando el
sujeto de la pasiva cuasirrefleja es singular, se puede producir ambigüedad entre
la interpretación pasiva y la impersonal. En oraciones como \forma{Se convocó un
plebiscito} o \forma{Se requiere una auxiliar para tareas contables}, si se
considera que el SN acepta pronominalizarse, como sucede en el español
rioplatense (\forma{Se lo convocó} o \forma{Se la requiere}), la interpretación
es impersonal y el SN es COD. En cambio, si, dado que el SN concuerda con el
verbo en singular, se lo considera sujeto de la oración, la construcción tiene
interpretación pasiva.
\end{destacado}

\subsubsection{Construcciones cuasirreflejas impersonales}\label{sec:10-impersonales}

En español se consideran impersonales las oraciones que carecen de sujeto
explícito, ya sea porque el evento descripto excluye esa función, como en los
verbos meteorológicos (\forma{llover, tronar, granizar}), o porque el argumento
exigido carece de referencia definida, como en las \termino{construcciones
cuasirreflejas impersonales} con \forma{se}.

\begin{exe}
  \ex\label{ex:10-bebe} Aquí no \textit{se bebe}.
  \ex\label{ex:10-senado} En el Senado \textit{se debatió} durante largas horas.
  \ex\label{ex:10-sufre} \textit{Se sufre} si el amor no es correspondido.
  \ex\label{ex:10-esta} \textit{Se está} bien en este lugar.
  \ex\label{ex:10-cima} \textit{Se llega} a la cima muchas veces agotado.
  \ex\label{ex:10-descansa} Ya no \textit{se descansa} tan bien como antes en
      este lugar.
\end{exe}

¿Qué cambio se reconoce en la estructura argumental de estas oraciones? En todos
los ejemplos, los verbos suponen un sujeto ---agente en \exref{ex:10-bebe},
\exref{ex:10-senado}, \exref{ex:10-cima} y \exref{ex:10-descansa}; experimentante
en \exref{ex:10-sufre}, o paciente en \exref{ex:10-esta}---, aunque este no se
explicita y su referencia permanece indeterminada, ya que su mención se
considera innecesaria o irrelevante. Por tanto, el sujeto se puede representar
por un pronombre arbitrario.\footnote{La noción de \textit{pro} arbitrario se
introdujo en el \capref{6}, §\ref{sec:06-noverbales}.} La interpretación del
sujeto admite ya sea:

\begin{enumerate}[label=\alph*)]
  \item una \termino{lectura genérica}, parafraseable por \sig{todo el mundo /
        nadie / cualquiera}; por ejemplo, para \exref{ex:10-bebe}: \sig{Aquí
        nadie bebe}, o
  \item una \termino{lectura existencial}, parafraseable por \sig{alguno(s) /
        alguien / cualquiera}, como en \exref{ex:10-esta}: \sig{Cualquiera está
        bien en este lugar}.
\end{enumerate}

Los ejemplos muestran que no solo los verbos transitivos, \exref{ex:10-bebe} y
\exref{ex:10-senado}, admiten esta construcción, como sucede con la pasiva, sino
que prácticamente todos los tipos de verbos la aceptan, incluso los estativos,
\exref{ex:10-sufre} y \exref{ex:10-esta}, y los intransitivos inacusativos
\exref{ex:10-cima} e inergativos \exref{ex:10-descansa}.

\begin{recuadro}[Para ampliar]
Aunque algunos gramáticos han planteado que en las construcciones cuasirreflejas
impersonales el clítico \forma{se} representa el sujeto, esta hipótesis ha sido
descartada dentro del marco de la gramática generativa. Se postula, en cambio,
que el sujeto es un pronombre nulo que se especifica convencionalmente como
[\textit{pro}\etq{arb}] (con el subíndice \textit{arb}, \sig{arbitrario}) para
indicar que su referencia no es específica. Se considera que \forma{se} es el
elemento que desencadena la interpretación impersonal de la construcción, ya que
«absorbe» los rasgos definidos de persona y número de la flexión, que está
fijada en la forma no marcada de tercera persona singular, la más apta para
remitir a un sujeto inespecífico.\footnote{Para ampliar las nociones acerca de
la categoría gramatical de persona se puede consultar
\textcite{cunarro2018}.}
\end{recuadro}

Si bien en todas las construcciones cuasirreflejas impersonales ejemplificadas se
puede considerar que se ha producido un proceso de intransitivización, también
existen cuasirreflejas impersonales transitivas en las que, a diferencia de lo
que sucede en las cuasirreflejas medias, medio-pasivas y pasivas, el argumento
interno del verbo no es promovido a sujeto. La característica común de los
objetos de estas impersonales es que tienen el rasgo \sig{humano} o, al menos,
\sig{animado} y, para bloquear su interpretación como sujetos, reciben la marca
de función \forma{a} y pueden, además, ser pronominalizados.

\begin{exe}
  \ex Se arrestó a los sospechosos del triple crimen. $\rightarrow$ Se
      \textit{los} arrestó.
  \ex En la comisaría se metió a los imputados en una celda. $\rightarrow$ Se
      \textit{los} metió en una celda.
  \ex Al día siguiente, se convocó al abogado. $\rightarrow$ Se \textit{lo}
      convocó.
\end{exe}

\begin{recuadro}[Para ampliar: concordancias anómalas en las impersonales cuasirreflejas]
A pesar de que la cuasirrefleja impersonal está fijada en tercera persona
singular, en el uso surgen construcciones que «las gramáticas del español
deploran de forma unánime» \parencite[1689]{mendikoetxea1999}, con
concordancias «anómalas» desde el punto de vista normativo, como \forma{En el
ingreso al estadio se cachean a todos los que van entrando} o \forma{Se
encerraron a los culpables}. Pero ¿qué sucede en estas oraciones? Estos casos
pueden ponerse en línea con otros tan frecuentes como \forma{Hubieron fuertes
tormentas en el sur} o \forma{Habían muchas personas en la manifestación de
ayer}, en los que el hablante hace concordar un verbo impersonal como
\forma{haber} con su COD. Si bien para la normativa estas oraciones son
incorrectas, muchos hispanohablantes no podrían fácilmente detectar el error.
Según se plantea en \textcite{giammatteo2010}, en todas estas construcciones
está operando una tendencia regularizadora que lleva al hablante a «rellenar el
hueco estructural con lo que dentro de la oración mejor califica para sujeto»,
que es el único SN disponible: el COD. Volviendo al caso de las impersonales con
\forma{se} con objeto animado plural, un mecanismo interpretativo similar admite
la presencia de una doble señal: por un lado, la marca funcional \forma{a}, que
bloquea la interpretación agentiva del objeto y, al mismo tiempo, el plural en
el verbo, que señala que el paciente es reinterpretado por el hablante como
sujeto sintáctico de la oración.
\end{recuadro}

\begin{ejercicio}\label{ej:10-cuadro}
Complete el siguiente cuadro con las características comunes y diferenciales de
los cuatro subgrupos de construcciones cuasirreflejas que producen cambio de
estructura argumental, según se muestra en la primera columna.

{\small
\begin{tabularx}{\linewidth}{@{}L c c c c@{}}
  \toprule
  Características & Medias & Medio-pasivas & Pasivas & Impersonales \\
  \midrule
  Omisión del agente (o democión a función periférica) & + & & & \\
  Promoción del tema a sujeto de la oración & + & & & \\
  Tipo de lectura promovida & proceso & & & \\
  Persona de la concordancia & todas\textsuperscript{a} & & & \\
  Número de la concordancia & sg. y pl. & & & \\
  Presencia de inductores & $-$ & & & \\
  Naturaleza del predicado resultante & proceso & & & \\
  \bottomrule
\end{tabularx}}

\textsuperscript{a} \pendiente{El original dice «3.ª persona», aunque el texto
afirma que las medias son paradigmáticas}
\end{ejercicio}

\begin{ejercicio}\label{ej:10-meme}
Preste atención al siguiente meme y responda lo pedido en las consignas de
abajo.

\imagenpendiente{Meme con la construcción \textit{se veía}}

\begin{consignas}
  \item La construcción pronominal \forma{se veía}, que aparece en el meme,
        ofrece dos lecturas posibles. ¿Podría explicarlas y decir qué tipo de
        construcción pronominal es cada una de ellas?\\
        \textit{Lectura 1}: \forma{En el espejo se veía a sí mismo venir}.
        Construcción pronominal refleja pura.\\
        \textit{Lectura 2}: \completar[3]
  \item Luego explique la incidencia de ambas interpretaciones en el efecto
        humorístico del meme.
\end{consignas}
\end{ejercicio}

\begin{ejercicio}\label{ej:10-criolipolisis}
Analice las construcciones con \forma{se} de la siguiente publicidad y resuelva
lo pedido en las consignas, según se muestra en el ejemplo resuelto de abajo.

\begin{consignas}
  \item Discuta si es posible más de una lectura y aporte pruebas que lo
        fundamenten.
  \item ¿Qué funcionamiento tienen estas estructuras con \forma{se} en el texto?
\end{consignas}

\begin{quote}\small
\textbf{Criolipólisis}\\
Se ofrece un tratamiento para eliminar adiposidad localizada. Es no invasivo. Se
realiza una vez por mes y se requiere de 4 a 6 sesiones. Cada zona se trata
durante 15 minutos. Precio promocional 2x1 (la segunda persona se atiende sin
cargo).
\end{quote}

\begin{resuelto}
\forma{Se ofrece}: pasiva cuasirrefleja. No puede interpretarse como impersonal,
dado que acepta la versión en plural: \forma{Se ofrecen tratamientos}.
\end{resuelto}
\end{ejercicio}

\begin{ejercicio}\label{ej:10-shua}
Lea el texto «Ejemplares raros», de \textit{Botánica del caos}, de Ana María
Shua (2000). Observe las estructuras subrayadas y luego complete el cuadro de
abajo según se muestra en los casos resueltos.

% Texto literario protegido: se incorpora a mano desde el original.
\textoaincorporar{el microrrelato «Ejemplares raros», de A.~M.~Shua
(\textit{Botánica del caos}, 2000), con las estructuras subrayadas que figuran
en el cuadro}

{\small
\begin{tabularx}{\linewidth}{@{}>{\raggedright\arraybackslash}p{24mm}
    >{\raggedright\arraybackslash}p{13mm} L
    >{\raggedright\arraybackslash}p{15mm} >{\raggedright\arraybackslash}p{16mm}@{}}
  \toprule
  Construcción & Tipo de verbo & Cambio de la estructura argumental & Tipo de
    construcción pronominal & Tipo de lectura \\
  \midrule
  \forma{Se les frota el cuerpo} & acción & El agente es eliminado y el tema es
    promovido a sujeto. Acepta plural: \forma{Se les frotan los cuerpos}. &
    pasiva & afectación \\
  \forma{Se aplica (\ldots) una pomada antiséptica} & & & & \\
  \forma{Se las cubre} (a las llagas) & acción & El sujeto agente es eliminado,
    pero el tema no es promovido a sujeto y se mantiene en la función de COD. No
    acepta plural: \juicio{*}{\forma{Se las cubren a las llagas}}. & impersonal
    & existencial (\sig{alguien las cubre}) \\
  \forma{Se cierran} (las llagas) \forma{sin dejar marcas} & & & & \\
  \forma{las llagas se inflaman} & & & & \\
  \forma{Su movimiento se debe solo a la velocidad\ldots} & & & & \\
  \forma{Se las rocía} & & & & \\
  \forma{Se cubren} & & & & \\
  \forma{Las heridas se cierran} & & & & \\
  \bottomrule
\end{tabularx}}
\end{ejercicio}

\subsection{Construcciones cuasirreflejas que producen cambio en la estructura sintáctica}\label{sec:10-es}

En algunas construcciones pronominales cuasirreflejas de toda persona, el
pronombre indica un cambio en la estructura sintáctica que afecta a la
manifestación de la transitividad. En otras palabras, se cambia la relación
sintáctica que establece el verbo con su argumento interno, que pasa de ser COD a
complemento preposicional de régimen (CPR). Se distinguen dos subgrupos:

\begin{itemize}
  \item \termino{construcciones cuasirreflejas alternantes}, con verbos que
        permiten las dos construcciones, con COD y con CPR (\forma{olvidar
        (algo)} y \forma{olvidarse de (algo)});
  \item \termino{construcciones cuasirreflejas puras} o \termino{inherentes},
        con verbos que solo admiten la construcción pronominal (\forma{quejarse}
        / \juicio{*}{\forma{quejar}}; \forma{jactarse} /
        \juicio{*}{\forma{jactar}}).
\end{itemize}

\subsubsection{Construcciones cuasirreflejas alternantes}\label{sec:10-alternantes}

Observemos lo que sucede si comparamos las estructuras siguientes:

\begin{exe}
  \ex\label{ex:10-olvide}
    \begin{xlist}
      \ex Olvidé lo que quería decirte.
      \ex[*]{Olvidé de lo que quería decirte.}
      \ex Me olvidé de lo que quería decirte.
    \end{xlist}
  \ex\label{ex:10-ocupa}
    \begin{xlist}
      \ex Ocupa dos estantes de la librería.
      \ex Me ocupo de dos estantes de la librería.
    \end{xlist}
\end{exe}

Las construcciones cuasirreflejas que producen un cambio en la manifestación de
la transitividad tienen un comportamiento sintáctico homogéneo: la incorporación
del elemento pronominal siempre convierte el COD en un CPR. Sin embargo, desde la
perspectiva semántica, no todas las construcciones se comportan igual: algunas
mantienen su significado y otras lo alteran. Por ejemplo, en
\exref{ex:10-olvide}, el cambio sintáctico no modifica el significado del verbo,
mientras que, en \exref{ex:10-ocupa}, la alternancia no solo implica un cambio de
transitividad (de directa a indirecta), sino que también ocasiona un cambio en el
significado del verbo: en (a), \forma{ocupar} quiere decir \sig{llenar un
espacio} y, en (b), \forma{ocuparse de} significa \sig{asumir la responsabilidad
de un asunto}.\pendiente{El original remite a (88) y (89)}

\begin{ejercicio}\label{ej:10-alternantes}
Contraste las formas transitivas con las formas pronominales de los verbos
cuasirreflejos alternantes que figuran en la siguiente tabla y luego resuelva lo
pedido abajo.

{\small
\begin{tabularx}{\linewidth}{@{}l L l L@{}}
  \toprule
  Forma transitiva & Significado & Forma pronominal & Significado \\
  \midrule
  \forma{acordar}    & \sig{establecer un acuerdo} & \forma{acordarse de} &
    \sig{recordar} \\
  \forma{cruzar}     & & \forma{cruzarse con} & \\
  \forma{hallar}     & & \forma{hallarse en} & \\
  \forma{aprovechar} & & \forma{aprovecharse de} & \\
  \forma{empeñar}    & & \forma{empeñarse en} & \\
  \bottomrule
\end{tabularx}}

\begin{consignas}
  \item Según se muestra en el ejemplo resuelto, complete la tabla con el
        significado de cada uno de los verbos y luego redacte oraciones con cada
        uno de los significados de cada verbo.
  \item ¿Qué diferencia observa en el cambio de significado de los pares de 1
        frente a los de 2?
        \begin{enumerate}[label=\arabic*.]
          \item \forma{acordar / acordarse de; ocupar / ocuparse de; empeñar /
                empeñarse en}
          \item \forma{lamentar / lamentarse de; compadecer / compadecerse de;
                arriesgar / arriesgarse a}
        \end{enumerate}
\end{consignas}
\end{ejercicio}

\subsubsection{Construcciones cuasirreflejas puras o inherentes}\label{sec:10-inherentes}

Algunos verbos son inherentemente pronominales y no presentan alternancia. Solo
admiten la construcción pronominal, de modo que el clítico \forma{se} forma
parte del lexema. Dentro de estos verbos, algunos carecen de régimen verbal
(\forma{rebelarse, vanagloriarse, enfadarse, suicidarse, acurrucarse,
portarse}, etc.) y otros exigen un CPR (\forma{arrepentirse de, atreverse a,
dignarse a, quejarse de, abstenerse de, esforzarse en / por, jactarse de},
etc.).

\begin{exe}
  \ex
    \begin{xlist}
      \ex Se queja de todo permanentemente.
      \ex[*]{Queja de todo permanentemente.}
    \end{xlist}
  \ex
    \begin{xlist}
      \ex Estaba deprimido y se suicidó.
      \ex[*]{Estaba deprimido y suicidó.}
    \end{xlist}
\end{exe}

\begin{ejercicio}\label{ej:10-inherentes}
Analice las siguientes estructuras y determine qué cambios introduce el
pronombre, según se muestra en el ejemplo resuelto~1.

\begin{exe}
  \ex
    \begin{xlist}
      \ex Se atrevieron a decir la verdad.
      \ex[*]{Atrevieron a decir la verdad.}
    \end{xlist}
  \ex
    \begin{xlist}
      \ex La profesora despedía a sus alumnos.
      \ex La profesora se despedía de sus alumnos.
    \end{xlist}
  \ex
    \begin{xlist}
      \ex Cruzaron el bosque para llegar al lago.
      \ex Se cruzaron con ciervos de los pantanos.
    \end{xlist}
  \ex
    \begin{xlist}
      \ex Se dignaron venir por la mañana.
      \ex[*]{Dignaron venir por la mañana.}
    \end{xlist}
  \ex
    \begin{xlist}
      \ex Aseguró el clavo en la pared.
      \ex Se aseguró de la verdad antes de venir.
    \end{xlist}
\end{exe}

\begin{resuelto}
1. La agramaticalidad de 1b muestra que se trata de un caso de cuasirrefleja
pura o inherente: solo es posible la variante pronominal (1a), por lo que se
entiende que el pronombre forma parte del lexema.
\end{resuelto}
\end{ejercicio}

\subsection{Construcciones cuasirreflejas que producen cambio en la estructura eventiva}\label{sec:10-ee}

Para considerar los casos de este tercer grupo, podemos partir de los siguientes
ejemplos:

\begin{exe}
  \ex\label{ex:10-vino} ¿Quién (se) ha tomado todo el vino? (canción de La Mona
      Jiménez)
  \ex\label{ex:10-telon} (Se) cayó el telón.
  \ex\label{ex:10-duermo} Siempre (me) duermo temprano.
  \ex\label{ex:10-fui} (Me) fui hasta tu casa, pero no te encontré.
  \ex\label{ex:10-pipa} (Se) fumaron hasta la pipa de la paz.
\end{exe}

En todas estas construcciones se puede eliminar el pronombre y las oraciones
siguen siendo gramaticales. De ahí que se haya considerado que el pronombre es
un elemento espurio o expletivo, y que se haya buscado asimilar su función a la
del denominado \termino{dativo ético}.\footnote{Se denomina \textit{dativo ético}
al pronombre átono no reflexivo que se refiere al individuo implicado
afectivamente por la acción que denota el predicado
\parencite[102]{bosque2019}.} Sin embargo, mientras el dativo ético no concuerda
con el sujeto (\forma{Este chico no me come; No me vaya a llorar ahora};
ejemplos de \cite[102]{bosque2019}), el pronombre que aparece en las oraciones
de arriba concuerda en persona y número con el sujeto.

Por otra parte, desde el punto de vista semántico, ¿significan lo mismo las
versiones con y sin pronombre de las oraciones anteriores? Por ejemplo, cuando
decimos \forma{Cayó el telón} entendemos que el telón se bajó porque la obra
terminó; en cambio, \forma{Se cayó el telón} supone una caída accidental. Sin
embargo, en otros casos, las diferencias de significado entre la construcción con
y sin pronombre pueden ser más escurridizas.

El factor común en las oraciones de arriba es que la inclusión del pronombre no
elimina ni promueve argumentos, aunque puede incidir en su aceptabilidad. En
este tipo de construcciones pronominales, que pueden darse tanto con verbos
transitivos, \exref{ex:10-vino} y \exref{ex:10-pipa}, como con intransitivos,
\exref{ex:10-telon} a \exref{ex:10-fui}, la presencia del pronombre se asocia
con el aspecto léxico.\footnote{Las diferencias entre aspecto léxico y aspecto
gramatical y otras nociones relacionadas se tratan en el \capref{8}. Asimismo,
para ampliar estos conceptos, remitimos a \textcite{debin2018}.}

\begin{destacado}
\textbf{Para recordar.} El \termino{aspecto} es una categoría gramatical que
tiene que ver con la manera en que el evento se desarrolla y que tiene gran
incidencia sintáctica. Así, por un lado, circunscribe o delimita la estructura
interna del evento, lo que se conoce como \termino{aspecto cualitativo}, y, por
el otro, afecta su duración e intensidad, lo que corresponde al \termino{aspecto
cuantitativo}.
\end{destacado}

Pero ¿cómo inciden estas características del aspecto en las construcciones
pronominales? A partir de la distinción entre aspecto cualitativo, vinculado con
la delimitación, y aspecto cuantitativo, relacionado con la intensificación, se
reconocen dos subgrupos: construcciones cuasirreflejas de delimitación y
construcciones cuasirreflejas de intensificación.

\subsubsection{Construcciones cuasirreflejas de delimitación}\label{sec:10-delimitacion}

La diferencia de significado entre los verbos, tanto transitivos como
intransitivos, que se construyen con y sin pronombre que trataremos en este
apartado ha sido ampliamente debatida en la bibliografía, y el consenso al que
se ha llegado es que la variante pronominal restringe o delimita al verbo. Pero
¿en qué sentido? La incorporación del pronombre tiene valor aspectual y convierte
un verbo \termino{atélico} (no delimitado) en \termino{télico} (delimitado), con
un momento inicial y un estado resultante. Sin embargo, el efecto que esta
delimitación produce no es el mismo en todos los casos. Según vimos en los
ejemplos de arriba, estas construcciones pueden manifestar varios matices, todos
ellos vinculados con la delimitación, de modo que reconocemos dos subgrupos:
construcciones culminativas con verbos transitivos, como \exref{ex:10-vino} y
\exref{ex:10-pipa},\pendiente{El original incluye también (93)} y construcciones
ingresivas con verbos intransitivos, como \exref{ex:10-duermo} y
\exref{ex:10-fui}.

\paragraph{Construcciones culminativas con verbos transitivos.} Como ya hemos
señalado en el \capref{8},\pendiente{El original dice «capítulo 9»} el aspecto
es composicional porque no está dado por el verbo solamente, sino que múltiples
factores ---el tipo de sujeto, los complementos o algunos modificadores--- pueden
delimitar la realización del evento. Así, por ejemplo, a un verbo lo podemos
acotar con un COD determinado o cuantificado, como muestra el contraste entre
\exref{ex:10-manzanas-a} y \exref{ex:10-manzanas-b}, por un lado, y entre
\exref{ex:10-abuelo-a} y \exref{ex:10-abuelo-b}, por el otro:

\begin{exe}
  \ex
    \begin{xlist}
      \ex\label{ex:10-manzanas-a} Pedro come manzanas.
      \ex\label{ex:10-manzanas-b} Pedro comió una manzana / la manzana que estaba
          en la frutera / tres manzanas.
    \end{xlist}
  \ex
    \begin{xlist}
      \ex\label{ex:10-abuelo-a} El abuelo toma medicación para dormir.
      \ex\label{ex:10-abuelo-b} El abuelo tomó la / su medicación para dormir.
    \end{xlist}
\end{exe}

Solo en las construcciones en que el verbo aparece delimitado por el objeto
---\exref{ex:10-manzanas-b} y \exref{ex:10-abuelo-b}--- es posible el uso del
denominado \forma{se} \termino{aspectual}, que denota acciones «completas y
delimitadas» \parencite[111]{sanchezlopez2002}, lo que explica la
agramaticalidad de \exref{ex:10-secome-a} y \exref{ex:10-setoma-a} frente a
\exref{ex:10-secome-b} y \exref{ex:10-setoma-b}, que resultan
gramaticales.\pendiente{El original remite a «(108b)»}

\begin{exe}
  \ex
    \begin{xlist}
      \ex[*]{\label{ex:10-secome-a} Pedro se come manzanas.}
      \ex\label{ex:10-secome-b} Pedro se comió una manzana / la manzana que
          estaba en la frutera.
    \end{xlist}
  \ex
    \begin{xlist}
      \ex[*]{\label{ex:10-setoma-a} El abuelo se toma medicación para dormir.}
      \ex\label{ex:10-setoma-b} El abuelo se tomó su medicación / una pastilla
          para dormir.
    \end{xlist}
\end{exe}

La exigencia de un punto final que señale la culminación del evento hace que la
construcción pronominal no se emplee ni con actividades ni con estados.
Típicamente, estas construcciones se combinan con verbos de consumo
(\forma{comer, beber, fumar, gastar}), aunque también son compatibles con otros
verbos de realización que indiquen un estado final resultante. Asimismo, el
sujeto de estos verbos puede ser agente o experimentante (\forma{Facundo se sabe
el himno de memoria}), pero no paciente (\juicio{*}{\forma{Facundo se vive a la
vuelta de mi casa}}).

\paragraph{Construcciones ingresivas con verbos intransitivos.} Algunos verbos
intransitivos de movimiento (\forma{ir, venir, salir, llegar, correr, caer,
quedar, aparecer, volver}, etc.) y otros intransitivos de cambio de estado
(\forma{dormir, morir, despertar, estar}) presentan una construcción pronominal
(\forma{irse, venirse, salirse, llegarse, correrse, subirse, bajarse, caerse,
quedarse, aparecerse, volverse, dormirse, morirse, despertarse, estarse}, etc.).
Aunque todos estos verbos tienen en común el carácter intransitivo, las
construcciones pronominales que forman no muestran un comportamiento uniforme;
por eso vamos a ver ambos grupos separadamente.

\paragraph{Ingresivas con verbos intransitivos de movimiento.} Nuevamente vamos a
partir de la observación de algunos ejemplos contrastivos:

\begin{exe}
  \ex\label{ex:10-voy}
    \begin{xlist}
      \ex\label{ex:10-voy-a} Voy ([\etq{Origen} de la clase de Griego])
          [\etq{Meta} a la de Gramática].
      \ex[*]{\label{ex:10-voy-b} Voy [\etq{Origen} de la clase de Griego].}
      \ex Me voy [\etq{Origen} de la clase de Griego] ([\etq{Meta} a la de
          Gramática]).
    \end{xlist}
  \ex\label{ex:10-corrio}
    \begin{xlist}
      \ex\label{ex:10-corrio-a} Corrió ([\etq{Origen} de aquí]) [\etq{Meta} a la
          esquina].
      \ex[*]{\label{ex:10-corrio-b} Corrió [\etq{Origen} de aquí].}
      \ex Se corrió [\etq{Origen} de aquí].
      \ex[*]{Se corrió colina abajo.}
    \end{xlist}
\end{exe}

¿Qué efecto produce la inclusión del pronombre en las oraciones de arriba?
Prestemos atención a los papeles temáticos. Cuando estos verbos se construyen sin
pronombre, \exref{ex:10-voy-a} y \exref{ex:10-corrio-a}, el evento no se
encuentra delimitado y pueden aparecer ambos complementos, origen y meta, o solo
la meta, pero las oraciones resultan agramaticales cuando se construyen solo con
el argumento de origen, \exref{ex:10-voy-b} y \exref{ex:10-corrio-b}. La
inaceptabilidad de estas construcciones con un SP de origen muestra que, para
que el límite inicial pueda explicitarse, el verbo tiene que estar acompañado
por una forma pronominal, que «va a ser el elemento encargado de delimitar el
evento de \textit{ir} y hacerlo compatible con el SP delimitador»
\parencite[2986]{demiguel1999}. Para este subgrupo de verbos, el pronombre es
obligatorio cuando el verbo tiene lo que \textcite{demiguel1999} denomina
«presunción de origen» e indica abandono del lugar donde se encuentra el
hablante. Por tanto, al marcar el inicio del movimiento, estas construcciones se
consideran \termino{ingresivas} o \termino{incoativas}.

Según muestran los ejemplos anteriores, la presencia del pronombre no solo
influye en la aceptación de los complementos, sino que también incide en el
significado del verbo: mientras \forma{ir} es «moverse de un lugar hacia otro»
(\textit{DLE}), \forma{irse} significa «dejar un lugar (para ir a otro)»
\parencite[2986]{demiguel1999};\pendiente{El original cita aquí «De Miguel,
2009»} y frente a \forma{correr}, que significa \sig{ir deprisa},
\forma{correrse} (no registrado con esta acepción por el \textit{DLE}) es
\sig{salirse de un lugar} (Giammatteo, 2019).\pendiente{Giammatteo (2019) no
figura en la bibliografía}

No todos los verbos de movimiento presentan un claro contraste entre la forma
con y sin pronombre. Existen oposiciones más difusas, como \forma{Vino de la clase
de Griego a la de Gramática} y \forma{Se vino de la clase de Griego a la de
Gramática}, que suelen explicarse en función del grado de involucramiento del
sujeto, que sería mayor en la construcción pronominal. Incluso un mismo verbo
puede no presentar un comportamiento uniforme. Así, por ejemplo, son posibles
tanto \forma{Cayó de un quinto piso} como \forma{Se cayó de un quinto piso};
pero difiere la interpretación entre \forma{Cayó el pelotón de infantería}
(\sig{se rindió}) y \forma{Se cayó el pelotón de infantería (en la fosa)}
(\sig{se hundió}).

\begin{ejercicio}\label{ej:10-correr}
Observe las siguientes oraciones y resuelva lo solicitado en las consignas,
según se muestra en el ejemplo resuelto~1.

\begin{consignas}
  \item Explique los contrastes entre los casos a y b.
  \item Señale los papeles temáticos de todos los argumentos.
\end{consignas}

\begin{exe}
  \ex
    \begin{xlist}
      \ex ¡Corré!
      \ex ¡Correte de acá!
    \end{xlist}
  \ex
    \begin{xlist}
      \ex Cayó gente en la fiesta.
      \ex Se cayó gente en la fiesta.
    \end{xlist}
  \ex
    \begin{xlist}
      \ex La lluvia cae (al piso).
      \ex[*]{La lluvia se cae (al piso).}
    \end{xlist}
  \ex
    \begin{xlist}
      \ex Me cayó una maceta del balcón.
      \ex Se (me) cayó una maceta del balcón.
    \end{xlist}
  \ex
    \begin{xlist}
      \ex Se (me) cayó la silla.
      \ex Me caí de la silla.
    \end{xlist}
\end{exe}

\begin{resuelto}
1. a) En 1b, el pronombre \forma{te} tiene un valor aspectual y hace foco en el
inicio del evento. Se trata de una construcción ingresiva o incoativa que
produce cambio de significado: \forma{correr} es \sig{trasladarse con cierta
velocidad hacia algún otro lugar}, pero \forma{correrse} se interpreta como
\sig{abandonar un sitio para dirigirse a otro}. b) \forma{De acá} tiene el papel
temático de origen.
\end{resuelto}
\end{ejercicio}

\begin{ejercicio}\label{ej:10-cocina}
Observe las estructuras subrayadas del siguiente texto y determine los cambios
que produce la aparición del pronombre \forma{se}, según se muestra en el
ejemplo resuelto de abajo.

% Texto periodístico protegido: se incorpora a mano desde el original.
\textoaincorporar{el copete de la nota «Antes y después. Una cocina se actualiza
en blanco y en azul», con las estructuras subrayadas: \textit{se actualiza;
ampliarse; mudarse; quedarse}}

\begin{resuelto}
\forma{Una cocina se actualiza en blanco y en azul}. \forma{Se} produce cambios
en la estructura argumental del verbo: elimina al agente, que tiene
interpretación genérica (\forma{Cualquiera (sabe) que una cocina se actualiza en
blanco y azul}), y el tema es promovido a sujeto. Se pasa de una construcción
transitiva a una cuasirrefleja medio-pasiva con valor caracterizador de lo que el
texto plantea que \sig{debe ser una cocina actualizada}. El inductor del sentido
genérico de la construcción es el presente con valor modal deóntico (de
obligación) y, si bien no hay en el contexto un adverbio modal, por el sentido
puede reponerse \forma{siempre} $\rightarrow$ \forma{Una cocina debe ser
(siempre) actualizada en blanco y azul / Para ser actualizada, una cocina debe
estar en blanco y azul}.
\end{resuelto}
\end{ejercicio}

\paragraph{Ingresivas con verbos intransitivos de cambio de estado.} ¿Qué
diferencia hay entre \forma{dormir} y \forma{dormirse}? \forma{Dormir} es quizás
el caso más claro de los verbos de este grupo. Mientras que la forma no
pronominal refiere una actividad no delimitada, \forma{dormirse} es un verbo de
logro que marca el inicio del suceso, por lo que tiene valor ingresivo. Por
tanto, en uno y otro caso difiere la compatibilidad con modificadores
temporales.\footnote{La compatibilidad de las distintas clases de verbos con
modificadores temporales se presentó en el \capref{8},
§\ref{sec:08-mt}.}

\begin{exe}
  \ex
    \begin{xlist}
      \ex Estaba muy cansada y dormí profundamente toda la noche / *en un
          segundo.
      \ex Estaba muy cansada y me dormí en un segundo / *durante toda la noche.
    \end{xlist}
\end{exe}

Si bien existe cierto acuerdo en que la variante pronominal restringe los
valores del verbo y resulta aspectualmente marcada al denotar eventos
perfectivos, ¿cómo se aplica esto a verbos como \forma{morir} o
\forma{despertar}, que ya son télicos? Algo en común que tienen todos los verbos
intransitivos que admiten construcciones pronominales, tanto de movimiento como
de cambio de estado, es que solo las variantes pronominales permiten la
incorporación de un dativo no concordado con el sujeto, que señala al afectado
por el evento.

\begin{exe}
  \ex\label{ex:10-gato}
    \begin{xlist}
      \ex Murió el / mi gato. / \juicio{*}{\forma{Me murió el gato.}}
      \ex\label{ex:10-gato-b} Se murió el / mi gato. / Se me murió el gato.
    \end{xlist}
  \ex\label{ex:10-lapiz}
    \begin{xlist}
      \ex Cayó el lápiz. / \juicio{*}{\forma{Me cayó el lápiz.}}
      \ex\label{ex:10-lapiz-b} Se cayó el lápiz. / Se me cayó el lápiz.
    \end{xlist}
\end{exe}

Varios de los verbos pronominales inacusativos admiten estas
\termino{construcciones doblemente pronominales} o \termino{bipronominales}, que
incorporan un dativo no concordado como segunda forma pronominal, cuya presencia
altera el orden de constituyentes porque el dativo se coloca entre el pronombre
concordado y el verbo \exref{ex:10-anteojos}:

\begin{exe}
  \ex\label{ex:10-anteojos}
      \begin{tabular}[t]{@{}llll@{}}
        Se & me & cayeron & los anteojos. \\
        \footnotesize clít. concordado & \footnotesize + clít. dat. no
          concordado & \footnotesize + V & \footnotesize + sujeto \\
      \end{tabular}
\end{exe}

En la forma no marcada de las construcciones doblemente pronominales, el sujeto
se pospone y el dativo, que representa al afectado por el evento, puede estar
reduplicado por un SP que suele colocarse al inicio:

\begin{exe}
  \ex A Juan se le cayó el cuaderno.
\end{exe}

\begin{recuadro}[Para ampliar: el dativo afectado]
En estas construcciones el dativo se considera el argumento más prominente, ya
que refiere a una entidad animada, mientras que el sujeto pospuesto suele ser no
animado. Para \textcite[9]{maldonado1994}, este \termino{dativo de afectación}
es «una construcción intermedia cuyo análisis ha sido problemático a lo largo de
la historia lingüística del español». Según \textcite[103]{sanchezlopez2002},
«aunque dativos, estos argumentos no pueden considerarse benefactivos, sino que
representan un argumento que es en cierto modo responsable del proceso sin
intervenir en él directamente». Otros autores
\parencite{masullo1992,fernandezsoriano1999,mendikoetxea2011} los consideran
sujetos en dativo, ya que pueden ocupar la posición canónica inicial y hacer las
veces de «sujeto psicológico» de la oración.\pendiente{El original cita
«Fernández Soriano, 1990»} El dativo afectado se combina con distintos tipos de
verbos que tienen en común el hecho de no poseer un sujeto que cause el evento,
como sucede con:

\begin{itemize}
  \item las formas pronominales de los verbos de cambio de estado de causa
        externa (\forma{quemar, estropear, romper}), con los que el dativo
        representa al responsable indirecto del evento;
        \begin{exe}
          \ex El jarrón se me rompió.
        \end{exe}
  \item las formas pronominales de verbos intransitivos inacusativos
        (\forma{morirse} \exref{ex:10-gato-b}, \forma{caerse}
        \exref{ex:10-lapiz-b}, \forma{irse, salirse}) con «presunción de
        origen». Como muestra \exref{ex:10-diente}, la incorporación del
        pronombre puede implicar un cambio de significado en el
        verbo;\pendiente{El original remite a (107)}
        \begin{exe}
          \ex\label{ex:10-diente}
            \begin{xlist}
              \ex Le salió un diente nuevo al bebé. (\sig{apareció})
              \ex Se le salió un diente al abuelo. (\sig{se desprendió de la
                  encía})
            \end{xlist}
        \end{exe}
  \item predicaciones impersonales (\forma{llover, anochecer}).
        \begin{exe}
          \ex Nos llovió toda la semana y tuvimos que volver de las vacaciones.
        \end{exe}
\end{itemize}

Con los inacusativos pronominales y los impersonales meteorológicos, el dativo
representa al experimentante afectado psicológicamente por el suceso.
\textcite{company2017} plantea que las construcciones con dativo afectado
constituyen una estrategia de subjetivización del hablante para interpretar los
sucesos desde su óptica personal y de manera afectiva. Así, aunque con
\forma{quemar} la construcción agentiva es perfectamente gramatical, una
oración como \exref{ex:10-asado-a} sería poco probablemente emitida por un
hablante. Por lo general, se prefiere la formulación anticausativa
\exref{ex:10-asado-b}, que presenta al evento como no causado y omite al
responsable. La aparición del dativo \exref{ex:10-asado-c} posibilita que el
causante se presente como un afectado por el hecho, en una lectura de proceso
accidental, en la que el sujeto se muestra como responsable indirecto o
involuntario del suceso, aunque involucrado en él.

\begin{exe}
  \ex
    \begin{xlist}
      \ex\label{ex:10-asado-a} Yo quemé el asado. ⟨agente, tema⟩
      \ex\label{ex:10-asado-b} El asado se quemó. ⟨tema⟩
      \ex\label{ex:10-asado-c} [Me distraje y] el asado se me quemó. ⟨tema,
          experimentante⟩
    \end{xlist}
\end{exe}

Los distintos casos que hemos presentado permiten ver el complejo estatus de
estos dativos, tradicionalmente considerados superfluos, pero que, sin embargo,
son determinantes a la hora de fijar la interpretación de la oración.
\end{recuadro}

Asimismo, como ya vimos con algunos verbos intransitivos de movimiento, también
con los de cambio la diferencia entre la construcción con y sin pronombre suele
estar dada porque la primera hace hincapié en el esfuerzo o la intencionalidad
del sujeto en la realización del evento. Así, en algunos casos se otorga
agentividad a sujetos que en las construcciones no pronominales se considerarían
pacientes, tal como muestran los contrastes siguientes:

\begin{exe}
  \ex
    \begin{xlist}
      \ex Vino a mi casa el domingo.
      \ex Se vino a mi casa desde José C. Paz en tren.
    \end{xlist}
  \ex
    \begin{xlist}
      \ex Apareció la niña que estaban buscando hacía varios días
          (casualmente).
      \ex Se apareció en la fiesta, aunque no estaba invitado (deliberadamente).
    \end{xlist}
  \ex
    \begin{xlist}
      \ex Trajo una torta para festejar.
      \ex Se trajo una torta para festejar.
    \end{xlist}
\end{exe}

\subsubsection{Construcciones cuasirreflejas de intensificación}\label{sec:10-intensificacion}

Estas construcciones se dan tanto con verbos transitivos como con intransitivos
inergativos y, según \textcite{ditullio2014}, son propias de la variedad
rioplatense. La incorporación del pronombre, apoyada por la presencia de algún
elemento intensificador, señala «[\ldots] un grado extremo en que se realiza una
actividad y enfatiza, así, la capacidad o habilidad del agente para llevarla a
cabo» \parencite[181]{ditullio2014} \exref{ex:10-traviata}, aunque el sujeto
también puede ser un experimentante, \exref{ex:10-floreros} y
\exref{ex:10-sabian}.

\begin{exe}
  \ex\label{ex:10-traviata} Se cantó todo / la vida / hasta \textit{La
      traviata}.
  \ex\label{ex:10-floreros} Después del viaje teníamos tanta sed que nos tomamos
      todo / hasta el agua de los floreros.
  \ex\label{ex:10-sabian} Se sabían todo / vida y milagros del cantante.
\end{exe}

\begin{ejercicio}\label{ej:10-intensificacion}
Analice los cambios producidos en las siguientes estructuras a partir de la
incorporación del pronombre, según se muestra en el ejemplo resuelto~1.

\begin{exe}
  \ex a) Tomó café. / b) *Se tomó café. / c) Se tomó un café. / d) Se tomó el
      café de su vida.
  \ex Se sabe todo, pero no lo cuenta a nadie.
  \ex a) Se comió toda la carne. / b) Se comió hasta las migas.
  \ex a) Fuma desde chico. / b) ¡Hoy se fumó la vida! / c) Se fumó varios
      cigarrillos durante el examen.
  \ex a) Salió con sus hijos al parque. / b) Se salió el botón de la camisa. /
      c) Me salí de la dieta.
  \ex Se iban a la iglesia a pedir perdón. (Videograf del canal televisivo
      \textit{Crónica})
\end{exe}

\begin{resuelto}
1. a) Oración transitiva. b) Oración agramatical: el COD es un nombre de materia
no determinado ni cuantificado que no delimita el evento, por lo que no es
posible la incorporación del \forma{se} aspectual delimitador. c) Construcción
cuasirrefleja pronominal culminativa de delimitación: el COD está delimitado y
la incorporación del \forma{se} aspectual es gramatical. d) Construcción
cuasirrefleja pronominal de intensificación que se interpreta como \sig{el mejor
café que tomó en su vida}, con valor superlativo.
\end{resuelto}
\end{ejercicio}

\begin{ejercicio}\label{ej:10-diferencias}
Identifique en cada caso la construcción pronominal involucrada y explique las
diferencias entre las versiones a y b de cada ejemplo, según se muestra en el
ejemplo resuelto~1.

\begin{exe}
  \ex
    \begin{xlist}
      \ex Se buscan especialistas en manejo de herramientas informáticas.
      \ex Se busca a los especialistas interesados en tomar becas.
    \end{xlist}
  \ex
    \begin{xlist}
      \ex Para las fiestas se come mucho.
      \ex Se comieron la entrada, el plato principal y el postre.
    \end{xlist}
  \ex
    \begin{xlist}
      \ex Este tema se aprende con facilidad.
      \ex Se aprendió de memoria la lección.
    \end{xlist}
  \ex
    \begin{xlist}
      \ex Se están enviando turnos de vacunación.
      \ex Los alumnos se enviaron mensajes anónimos.
    \end{xlist}
\end{exe}

\begin{resuelto}
1. La incorporación del pronombre \forma{se} en ambas construcciones provoca un
cambio en la estructura argumental: el agente es eliminado. Las diferencias
estriban en que en a el tema (\forma{especialistas\ldots}) es promovido a
sujeto, lo que permite una lectura pasiva; en b, en cambio, la marca de función
\forma{a} impide la promoción del tema (\forma{los especialistas
interesados\ldots}) a la función de sujeto, lo que produce una lectura
impersonal.
\end{resuelto}
\end{ejercicio}

\begin{ejercicio}[Actividad de integración]\label{ej:10-akasha}
Observe las estructuras subrayadas en el siguiente fragmento del cuento
«Akasha», de S.~Fernández (2019), y resuelva las consignas que se presentan a
continuación del texto, según se muestra en los ejemplos resueltos de abajo.

% Texto literario protegido: se incorpora a mano desde el original.
\textoaincorporar{el epígrafe y el fragmento de «Akasha», de S.~Fernández
(2019), desde «Cierro mis ojos para descansar» hasta «\ldots{}esta nueva
sensación de desamparo», con las estructuras pronominales subrayadas}

\begin{consignas}
  \item En cada uno de los casos, identifique si la referencia del pronombre
        coincide con la del sujeto y determine si la construcción es pronominal
        o no.
  \item En las construcciones pronominales, evalúe si el pronombre personal
        involucrado manifiesta un argumento del verbo al que acompaña o no:
        \begin{enumerate}[label=\roman*)]
          \item en caso afirmativo, determine el papel temático y la función
                sintáctica del clítico y proporcione pruebas para identificar de
                qué tipo de pronombre reflejo se trata (reflejo puro, recíproco
                o factitivo);
          \item en caso negativo, cuando el pronombre involucrado no manifieste
                ningún argumento del verbo, indique qué tipo de cambio introduce
                (en la estructura sintáctica, en la argumental o en la
                eventiva), de qué tipo de construcción cuasirrefleja se trata y
                aporte pruebas que justifiquen su decisión.
        \end{enumerate}
\end{consignas}

\begin{resuelto}[Ejemplos resueltos]
\forma{(El aire) me molesta}. a) La referencia del pronombre no coincide con la
del sujeto y la construcción no es pronominal. \forma{(El poder espiritual de
carácter omnipresente) se encuentra\ldots}: a) la referencia del pronombre
coincide con la del sujeto y la construcción es pronominal. b) El pronombre no
manifiesta ningún argumento del verbo. ii) El pronombre produce cambio en la
estructura argumental y la construcción es cuasirrefleja pasiva porque admite
tercera persona singular y plural: \forma{Los poderes espirituales de carácter
omnipresente se encuentran}.\pendiente{El original dice aquí «estructura
sintáctica»; el solucionario, «estructura argumental»}
\end{resuelto}
\end{ejercicio}

\begin{recuadro}[Las construcciones pronominales en el texto]
En este fragmento del cuento se observa que la narradora construye una escena en
la que los eventos que se suceden «le suceden» al personaje. Aun en los casos en
los que a la protagonista le corresponde el papel de agente, parece actuar
dejándose llevar por las circunstancias (\forma{cierro (mis ojos), vuelo,
extiendo (mis brazos), rozo (algo), vuelvo (al tren), digo, viajo (lo hago por
primera vez)}) o a la expectativa de lo que le viene del exterior (\forma{miro,
escucho}). Para lograr el clima de despersonalización que la lleva de la
ensoñación y la sensación de libertad inicial a la soledad y el desamparo final
---un pasaje de la vida a la muerte, como prefigura el epígrafe---, la
protagonista cada vez asume una actitud más pasiva, que deja ver un personaje
inseguro, al cual «las cosas le ocurren». Este efecto se logra porque predominan
los casos en que ella no es el agente que lleva adelante la acción, sino que
aparece ya sea como objeto paciente de verbos transitivos (\forma{(el aire)
\textbf{me} rodea, \textbf{me} eleva; (la puerta) \textbf{me} saca del ensueño;
(mi acompañante) \textbf{me} mira; \textbf{me} despierta (una inmensa
curiosidad); \textbf{me} hace sentir; \textbf{me} pregunto}), como sujeto
paciente de verbos de estado (\forma{estoy sentada, estoy en un tren, no estoy
tan sola}) o como experimentante de verbos psicológicos (\forma{siento,
\textbf{me} molesta, creo, no sé, vuelvo a pensar}). Otro recurso de
despersonalización empleado en el fragmento es el uso de construcciones
pronominales medias, que favorecen la presencia de sujetos no agentivos
(\forma{se deslizan (mis dedos), se borra (mi sonrisa), me levanto}). Junto con
las construcciones verbales, refuerzan la actitud de inseguridad de la
protagonista las preguntas y los índices de duda y negación que indican
desconocimiento de la situación (\forma{¿Cuánto hace que estoy sentada? No lo
sé. ¿Cómo lo sé? No tengo idea. ¿A dónde voy? ¿Acaso importa?}, \forma{quizás
sea quien respondió, no estoy tan sola, quizás, a dónde voy, a dónde vamos}).
Asimismo, también los elementos léxicos acompañan esta presentación y la
protagonista pasa de una \forma{ensoñación} en la que se siente \forma{libre} a
un estado en que \forma{se le borra la sonrisa} y \forma{se siente incómoda, con
malestar}, hasta finalmente sentir \forma{la soledad y el desamparo}.
\end{recuadro}

\begin{recuadro}[Reflexión final]
El recorrido realizado en este capítulo nos ha mostrado que el campo de las
construcciones pronominales es sumamente amplio y que las distintas subclases no
siempre presentan bordes nítidos, lo que las convierte en uno de los temas más
complejos de la gramática del español, aunque también, según creemos, en uno de
los más interesantes. En el capítulo hemos distinguido dos grupos principales:
las reflejas, en las que el pronombre tiene función sintáctica de COD o de
COI,\pendiente{El original dice «de COD o de COD»} y las cuasirreflejas, que se
parecen a las primeras en la forma, pero en las que el pronombre no cumple
ninguna función sintáctica. Si bien en las reflejas, dentro de las que hemos
distinguido entre las reflexivas (puras, accidentales y factitivas) y las
recíprocas, se da un paso hacia la intransitivización, en tanto el agente y el
paciente coinciden, es en las cuasirreflejas en las que se advierte este proceso
como completado, debido a que se produce el borrado o la democión del agente
dentro de la estructura oracional. Con fines didácticos hemos adoptado para estas
construcciones cuasirreflejas una triple distinción, que tiene en cuenta si la
introducción del pronombre afecta a la estructura argumental (cuasirreflejas
medias, medio-pasivas, pasivas e impersonales), a la estructura sintáctica
(cuasirreflejas alternantes e inherentes) o a la estructura eventiva
(cuasirreflejas de delimitación y de intensificación).
\end{recuadro}

Por último, a modo de recapitulación, en la tabla~\ref{tab:10-sintesis} ofrecemos
una síntesis de las distinciones reconocidas en el capítulo.

{\footnotesize\setlength{\tabcolsep}{3pt}
\begin{longtable}{@{}>{\raggedright\arraybackslash}p{30mm}
    >{\raggedright\arraybackslash}p{38mm} >{\raggedright\arraybackslash}p{52mm}@{}}
  \caption{Sistematización de construcciones pronominales}\label{tab:10-sintesis}\\
  \toprule
  Grupo & Subtipo & Ejemplo \\
  \midrule
  \endfirsthead
  \toprule
  Grupo & Subtipo & Ejemplo \\
  \midrule
  \endhead
  \multicolumn{3}{@{}l}{\textbf{Reflejas}: el pronombre es COD o COI} \\
  \addlinespace
  Reflexivas (Pr Refl) & reflexivas puras & \textit{Juan se vio (a sí mismo).} \\
  & reflexivas accidentales & \textit{Juan se rompió el dedo
    (accidentalmente).} \\
  & reflexivas factitivas & \textit{Juan se cortó el pelo en la peluquería} (=
    \textit{Se hizo cortar el pelo en la peluquería}). \\
  Recíprocas (Pr Recipr) & --- & \textit{María y Pedro se miraron (el uno al
    otro / recíprocamente).} \\
  \midrule
  \multicolumn{3}{@{}p{120mm}}{\textbf{Cuasirreflejas}: el pronombre no tiene
    función sintáctica (es IPCR); producen cambios en la oración} \\
  \addlinespace
  Cambio en la estructura sintáctica & alternantes (con cambio de significado o
    sin él) & \textit{Olvidé eso / Me olvidé de eso. Ocupa dos estantes / Se
    ocupa de sus hijos.} \\
  & inherentes & \textit{Juan se arrepintió de lo que dijo.} \\
  \addlinespace
  Cambio en la estructura argumental: de toda persona & medias: verbos de cambio
    de estado físico (\textit{romper / romperse}), de cambio de estado mental
    (\textit{molestar / molestarse}) o de cambio de posición (\textit{sentar /
    sentarse}) & \textit{La puerta se cerró (sola). María se divirtió. El bebé ya
    se sienta (solo).} \\
  Cambio en la estructura argumental: de tercera persona & medio-pasivas &
    \textit{Esta carne se corta fácilmente.} \\
  & pasivas (el verbo concuerda con el sujeto en 3.ª persona, sg. o pl.) &
    \textit{Se tomó una medida polémica / Se tomaron medidas polémicas.} \\
  & impersonales (el verbo está fijado en 3.ª persona sg.) & \textit{Se habla
    mucho y se trabaja poco.} \\
  \addlinespace
  Cambio en la estructura eventiva o aspectual & de delimitación & \textit{Se
    comió toda la torta. Se fue a su casa. Se durmió temprano.} \\
  & de intensificación & \textit{Se comió hasta las migas.} \\
  \bottomrule
\end{longtable}}

% --- Secciones finales -------------------------------------------------
\lecturasrecomendadas

\solucionario

\begin{solucion}{\ref{ej:10-caso}}
2. \forma{Lo}: caso acusativo, función sintáctica COD. 3. \forma{Me}: caso
dativo, función sintáctica COI. 4. \forma{Me}: caso dativo, función sintáctica
COI. 5. \forma{Te}: caso dativo, función sintáctica COI; \forma{lo}: caso
acusativo, función sintáctica COD; \forma{me}: caso acusativo, función
sintáctica COD.
\end{solucion}

\begin{solucion}{\ref{ej:10-funcion}}
2. \forma{Le}: función sintáctica COI, papel temático meta. 3. \forma{Lo}:
función sintáctica COD, papel temático tema. 4. \forma{Le}: función sintáctica
COI, papel temático meta (\forma{le} duplica al COI \forma{a su hermano}). 5.
\forma{Se}: función sintáctica COI, papel temático meta. A diferencia de 4, aquí
aparecen ambos objetos pronominalizados, por lo que en lugar de \forma{le}
aparece el alomorfo \forma{se}. \forma{Lo}: función sintáctica COD, papel
temático tema.
\end{solucion}

\begin{solucion}{\ref{ej:10-sustitucion}}
2. a) Sustitución: \forma{Le arreglé el pelo frente al espejo}. b) El pronombre
sustituto \forma{le} manifiesta caso dativo; por lo tanto, la función sintáctica
de \forma{me} es COI. 3. a) Sustitución: \forma{Lo vistió de gala}. b) El
pronombre sustituto \forma{lo} manifiesta caso acusativo; por lo tanto, la
función sintáctica de \forma{se} es COD. 4. a) Sustitución: \forma{Lo ayudaron
con la tarea}. b) El pronombre sustituto \forma{lo} manifiesta caso acusativo;
por lo tanto, la función sintáctica de \forma{te} es COD. 5. a) Sustitución:
\forma{Les dieron la vacuna para viajar}. b) El pronombre sustituto \forma{les}
manifiesta caso dativo; por lo tanto, la función sintáctica de \forma{nos} es
COI.
\end{solucion}

\begin{solucion}{\ref{ej:10-papeles}}
2. \forma{Te}: paciente, COD; \forma{nos}: beneficiario, COI. 3. \forma{Se}:
beneficiario, COI; \forma{lo}: tema, COD. 4. \forma{Me}: beneficiario, COI;
\forma{la}: tema, COD. 5. \forma{Se}: beneficiario, COI; \forma{las}: tema,
COD.
\end{solucion}

\begin{solucion}{\ref{ej:10-factitivas}}
Mi primo \forma{se} hizo una hermosa casa: agente: la empresa o persona que
construye; causante instigador: el sujeto, \forma{mi primo}. El pronombre
\forma{se} es COI. \forma{Me} hice un piercing: agente: la persona que colocó el
piercing; causante instigador: el sujeto tácito \forma{yo}. El pronombre
\forma{me} es COI. Lorena \forma{se} vacunó: agente: la persona o enfermera que
colocó la vacuna; causante instigador: \forma{Lorena}. El pronombre \forma{se}
es COI.\pendiente{El original intercambia «agente» y «causante instigador» y
analiza \textit{se} como COI en \textit{se vacunó} (es COD)}
\end{solucion}

\begin{solucion}{\ref{ej:10-ambiguas}}
Respuestas posibles. a) 2. Reflexiva: \forma{Cada uno de ellos se pasó la
tintura por el propio pelo}. Recíproca: \forma{Ana le pasó la tintura a su
hermano mientras él le hacía algunos retoques de tintura a ella}. 3. Refleja
pura: \forma{Cada uno se copió usando su propio celular}. Recíproca: \forma{Cada
uno se copió del otro / Se copiaron mutuamente}. b) En 2 también es posible la
interpretación como refleja factitiva: \forma{Ana y su hermano se hicieron teñir
el pelo en la peluquería de la esquina de su casa}.
\end{solucion}

\begin{solucion}{\ref{ej:10-tiporefleja}}
2. Reflexiva factitiva. 3. Reflexiva pura / factitiva. 4. Reflexiva pura. Esta
oración es ambigua entre una interpretación accidental y la reflexiva pura, pero
se desambigua por conocimiento de mundo, ya que, según se sabe, el famoso pintor
se automutiló. 5. Refleja recíproca.
\end{solucion}

\begin{solucion}{\ref{ej:10-danna}}
\forma{Se toca el cabello}: construcción pronominal refleja. El pronombre
\forma{se} es correferencial con el sujeto. Admite el refuerzo \forma{a sí
misma}. \forma{Le arregle las pestañas y las cejas}: no es construcción
pronominal. El pronombre \forma{le} no es correferencial con el sujeto.
\forma{Me estoy agarrando el pelo}: construcción pronominal refleja. El
pronombre \forma{me} es correferencial con el sujeto. Admite el refuerzo
\forma{a mí misma}. \forma{Me vean bien el pelo}: no es construcción pronominal.
El pronombre \forma{me} no es correferencial con el sujeto. \forma{Me las tengo
que hacer yo}: construcción pronominal refleja. El pronombre \forma{me} es
correferencial con el sujeto. Admite el refuerzo \forma{a mí misma}. \forma{Me
las haga}: no es construcción pronominal. El pronombre \forma{me} no es
correferencial con el sujeto. \forma{Se comió todas las uñas de la mano
derecha}: construcción pronominal refleja. El pronombre \forma{se} es
correferencial con el sujeto. Admite el refuerzo \forma{a sí misma}.
\end{solucion}

\begin{solucion}{\ref{ej:10-levantar}}
En 2a, el verbo \forma{extender} se comporta como transitivo: se construye con
sujeto agente (\forma{el nuevo gobierno}) y COD tema (\forma{los límites}), y
acepta la paráfrasis con \forma{hacer}: \forma{¿Qué hizo el gobierno?} En 2b
aparece la variante pronominal de \forma{extender}, referida a un suceso que
afecta a un tema. Se puede añadir el refuerzo \forma{solo} o \forma{por sí
solo}.
\end{solucion}

\begin{solucion}{\ref{ej:10-cambio}}
2. Verbo de cambio de posición. 2a: construcción transitiva con agente (tácito)
y tema (\forma{las cortinas}); 2b: construcción cuasirrefleja con sujeto
(tácito) simultáneamente agente y tema. 3. Verbo de cambio de posición. 3a:
construcción transitiva con agente (tácito) y paciente (\forma{los niños}); 3b:
construcción cuasirrefleja con sujeto (tácito) simultáneamente agente y tema.
\end{solucion}

\begin{solucion}{\ref{ej:10-cuadro}}
{\small
\begin{tabularx}{\linewidth}{@{}L c >{\centering\arraybackslash}p{20mm} c
    >{\centering\arraybackslash}p{18mm}@{}}
  \toprule
  Características & Medias & Medio-pasivas & Pasivas & Impersonales \\
  \midrule
  Omisión del agente (o democión a función periférica) & + & + & + & + \\
  Promoción del tema a sujeto de la oración & + & + & + & $-$ \\
  Tipo de lectura promovida & proceso & caracterizadora & afectación & genérica
    o existencial \\
  Persona de la concordancia & todas & 3.ª & 3.ª & 3.ª \\
  Número de la concordancia & sg. y pl. & sg. y pl. & sg. y pl. & sg. \\
  Presencia de inductores & $-$ & tiempos imperfectivos y adverbios modales &
    $-$ & $-$ \\
  Naturaleza del predicado resultante & proceso & estado & acción & sin
    restricciones \\
  \bottomrule
\end{tabularx}}
\end{solucion}

\begin{solucion}{\ref{ej:10-meme}}
a) Lectura 2: \forma{Se veía (en el espejo) esto: que él venía}. Construcción
pronominal cuasirrefleja pasiva. b) El efecto humorístico surge del contraste
entre ambas interpretaciones: en la lectura 1 se da a entender que el hombre
muere al chocar (él mismo) contra el espejo; en la 2, lo que se plantea es que
la muerte era previsible (= \forma{se (la) veía venir}).
\end{solucion}

\begin{solucion}{\ref{ej:10-criolipolisis}}
\forma{Se realiza}, \forma{se trata} y \forma{se atiende} son pasivas
cuasirreflejas. No pueden interpretarse como impersonales porque pueden
pluralizarse: \forma{Se realizan tratamientos; Todas las zonas se tratan durante
15 minutos; Las dos personas se atienden sin cargo}. \forma{Se requiere} es una
cuasirrefleja impersonal. No puede interpretarse como pasiva porque no acepta el
plural: \juicio{*}{\forma{Se requieren de 4 a 6 sesiones}}. Función textual de
estas construcciones: tanto las cuasirreflejas pasivas como las impersonales, al
promover el «borrado» del agente, resultan aptas para textos en los que cobra
mayor relevancia el tema que el agente.
\end{solucion}

\begin{solucion}{\ref{ej:10-shua}}
{\footnotesize\setlength{\tabcolsep}{3pt}
\begin{longtable}{@{}>{\raggedright\arraybackslash}p{24mm}
    >{\raggedright\arraybackslash}p{16mm} >{\raggedright\arraybackslash}p{44mm}
    >{\raggedright\arraybackslash}p{14mm} >{\raggedright\arraybackslash}p{18mm}@{}}
  \toprule
  Construcción & Tipo de verbo & Cambio de la EA & Tipo de construcción
    pronominal & Tipo de lectura \\
  \midrule
  \endhead
  \forma{Se aplica (\ldots) una pomada antiséptica} & acción & El agente es
    eliminado y el tema es promovido a sujeto. Acepta plural: \forma{Se aplican
    pomadas}. & pasiva & afectación \\
  \forma{Se cierran (las llagas) sin dejar marcas} & cambio de estado físico & El
    agente es eliminado y el tema es promovido a sujeto. & media & proceso \\
  \forma{Las llagas se inflaman} & cambio de estado físico & El agente es
    eliminado y el tema es promovido a sujeto. & media & proceso \\
  \forma{Su movimiento se debe solo a la velocidad (\ldots)} & acción & El agente
    es demovido a ACC y el tema es promovido a sujeto. Acepta plural: \forma{Sus
    movimientos se deben a (\ldots)}. & pasiva & afectación \\
  \forma{Se las rocía} & acción & El agente es eliminado, pero el tema no es
    promovido a sujeto y se mantiene en la función de COD. No acepta plural:
    \juicio{*}{\forma{Se las rocían (a las raicillas)}}. & impersonal &
    existencial (\sig{alguien las rocía}) \\
  \forma{Se cubren} & acción & El agente es eliminado y el tema es promovido a
    sujeto. Acepta singular: \forma{Se cubre}. & pasiva & afectación \\
  \forma{Las heridas se cierran} & cambio de estado físico & El agente es
    eliminado y el tema es promovido a sujeto. & media & proceso \\
  \bottomrule
\end{longtable}}
\pendiente{Revisar el análisis de \textit{se debe}: \textit{deberse} es más bien
cuasirreflejo inherente}
\end{solucion}

\begin{solucion}{\ref{ej:10-alternantes}}
\textit{Aclaración}: los ofrecidos en el cuadro son ejemplos posibles de uso de
los verbos.

a)

{\small
\begin{tabularx}{\linewidth}{@{}l L l L@{}}
  \toprule
  Forma transitiva & Significado & Forma pronominal & Significado \\
  \midrule
  \forma{cruzar} & \sig{pasar de un lado a otro opuesto}: \forma{Cruzamos el
    puente para llegar a Zárate.} & \forma{cruzarse con} & \sig{encontrarse con
    alguien viniendo de lugares opuestos}: \forma{Me crucé con mi vecino en el
    supermercado.} \\
  \forma{hallar} & \sig{encontrar}: \forma{No pudieron hallar la solución.} &
    \forma{hallarse en} & \sig{estar en algún lugar o situación}: \forma{En el
    sueño me hallaba en un lugar maravilloso.} \\
  \forma{aprovechar} & \sig{utilizar algo útilmente, obtener beneficio de
    algo}: \forma{Aprovechamos las vacaciones para conocer Córdoba.} &
    \forma{aprovecharse de} & \sig{sacar provecho de algo o de alguien}:
    \forma{Se aprovecharon de su ingenuidad y lo estafaron.} \\
  \forma{empeñar} & \sig{dejar en prenda algo}: \forma{Empeñó su palabra y todos
    confiaron en que cumpliría lo pactado.} & \forma{empeñarse en} &
    \sig{esforzarse, obstinarse en algo}: \forma{Está empeñado en aprender a
    nadar.} \\
  \bottomrule
\end{tabularx}}

b) Los verbos de 1 son cuasirreflejos alternantes con cambio de significado; los
verbos de 2 son alternantes, pero no presentan cambio de significado.
\end{solucion}

\begin{solucion}{\ref{ej:10-inherentes}}
En 2, 3 y 5, la introducción del pronombre produce cambios en la estructura
sintáctica: la variante sin pronombre exige un COD (\forma{sus alumnos, el
bosque, el clavo}), mientras que la variante pronominal exige un CPR. Se trata
de construcciones cuasirreflejas alternantes. Significado: en 2a, quienes se
retiran son \forma{los alumnos}; en 2b, quien se retira y saluda o se despide es
\forma{la profesora}. En 3a, \forma{cruzar} se interpreta como \sig{atravesar};
en 3b, \forma{cruzarse con} se interpreta como \sig{encontrarse con (de
casualidad), toparse}. En 5a, \forma{asegurar} se interpreta como \sig{afirmar,
reforzar, fortalecer}; en 5b, \forma{asegurarse} se interpreta como
\sig{cerciorarse, confirmar, estar seguro de algo}. En 4, la variante sin
pronombre es agramatical, lo que muestra que se trata de un caso de cuasirrefleja
pura o inherente, en la que el pronombre forma parte del lexema.
\end{solucion}

\begin{solucion}{\ref{ej:10-correr}}
En todos los casos, la introducción del pronombre produce cambios en la
estructura eventiva y afecta la delimitación del evento. 2. a) En 2b, \forma{se}
tiene valor aspectual y hace foco en el inicio del evento: se trata de una
construcción ingresiva o incoativa que produce cambio de significado: en 2a,
\forma{caer} refiere a \sig{llegar de improviso a un lugar}, pero en 2b
\forma{caerse} se interpreta como \sig{desplomarse en el suelo}. b) En ambas
oraciones, \forma{gente} tiene el papel temático de paciente. 3. a) En 3b, el
clítico no es posible porque el evento de llover se presenta como no delimitado:
se interpreta como \sig{movimiento descendente, hacia abajo} y no implica
\sig{desprendimiento de ningún lugar}. b) En 3a, \forma{la lluvia} es tema. 4. a)
En 4a no hay coincidencia entre la persona del pronombre \forma{me} y el verbo
\forma{cayó}: no hay construcción pronominal; 4b es una construcción pronominal y
tiene significado accidental. b) En 4a, \forma{me} tiene el papel temático de
meta y en 4b \forma{me} indica al afectado por el evento: es un benefactivo
negativo omisible. 5. a) En 5a no hay coincidencia entre la persona del pronombre
\forma{me} y el verbo \forma{cayó}: no hay construcción pronominal; 5b es una
construcción pronominal cuasirrefleja. b) En 5a, \forma{me} indica al afectado
por el evento, es un benefactivo negativo omisible, y \forma{la silla} es tema;
en 5b, \forma{la silla} es origen.
\end{solucion}

\begin{solucion}{\ref{ej:10-cocina}}
\forma{Ampliarse}: uso reflexivo del verbo transitivo \forma{ampliar}, que se
construye con agente y tema: \sig{alguien amplía algo}. En el texto se plantea
que H.~Milberg y su familia se amplían a sí mismos al ampliar su vivienda, que es
considerada una posesión inalienable (v. nota~\ref{fn:10-inalienable}).
\forma{Se} funciona como COD. \forma{Mudarse}: se trata de una construcción
cuasirrefleja alternante, que produce un cambio en la estructura sintáctica: el
COD se convierte en CPR y se pasa de la construcción transitiva \forma{mudar
casa} (registrada en el \textit{DLE}, aunque no utilizada en la variedad
rioplatense) a otra de transitividad indirecta, \forma{mudar(se) de casa}.
\forma{Quedarse}: \forma{se} produce un cambio en la estructura eventiva: la
variante pronominal del verbo tiene interpretación ingresiva debido al efecto
del \forma{se} delimitador.
\end{solucion}

\begin{solucion}{\ref{ej:10-intensificacion}}
2. Construcción cuasirrefleja pronominal de intensificación producida por la
incorporación de \forma{se} y el intensificador \forma{todo}. Se interpreta que
el evento se lleva a cabo hasta un punto extremo. 3. a) Construcción
cuasirrefleja pronominal culminativa de delimitación. El núcleo del COD, que es
un nombre de materia (\forma{carne}), está delimitado por el cuantificador
\forma{todo} y habilita la incorporación del \forma{se} aspectual. b)
Construcción cuasirrefleja pronominal de intensificación producida por la
incorporación de \forma{se} y el COD intensificado por \forma{hasta}.\footnote{En
estas construcciones, \forma{hasta} no es preposición, sino un focalizador, que
equivale a \forma{incluso}. Los focalizadores pueden resaltar o enfatizar
distintos constituyentes de la oración.} Se interpreta que el evento se llevó a
cabo hasta un punto extremo. 4. a) El COD de \forma{fumar} queda implícito y se
interpreta como \sig{alguna sustancia fumable, por lo general, cigarrillos}. b)
Construcción cuasirrefleja pronominal de intensificación producida por la
incorporación de \forma{se} y el añadido del COD con sentido superlativo,
\forma{la vida}. c) Construcción cuasirrefleja pronominal culminativa de
delimitación. El COD está cuantificado y permite la incorporación del
\forma{se} aspectual. 5. a) Oración intransitiva. El verbo significa \sig{ir
afuera}. b) Construcción cuasirrefleja de delimitación ingresiva. El verbo
cambia el significado a \sig{desprenderse}, con presunción de origen. El
argumento de origen es \forma{de la camisa}. c) \forma{Salirse} es una variante
pronominal de \forma{salir}. \forma{Me} es omisible y \forma{de la dieta} indica
origen (el lugar figurado de donde se sale). 6. \forma{Ir} es un verbo
intransitivo de movimiento que presupone un origen no explicitado y una meta
(\forma{a la iglesia}). Se trata de una construcción cuasirrefleja de
delimitación ingresiva.
\end{solucion}

\begin{solucion}{\ref{ej:10-diferencias}}
2. En a, el verbo transitivo tiene lectura genérica y la construcción es
impersonal; en b, la enumeración de los diferentes platos del menú se interpreta
como \sig{se comieron todo} y la construcción es cuasirrefleja de
intensificación. 3. En a, el pronombre \forma{se} produce un cambio en la
estructura argumental: se elimina el agente. Dado que el verbo está en presente y
lo acompaña el SP \forma{con facilidad}, se produce una lectura genérica y la
construcción es medio-pasiva. En b, la incorporación del pronombre produce un
cambio en la estructura eventiva: el \forma{se} delimitador convierte el evento
durativo \forma{aprender} en uno delimitado. 4. En a, el pronombre \forma{se}
provoca un cambio en la estructura argumental: se elimina el agente y se
promueve el tema (\forma{turnos de vacunación}) a sujeto de la oración. La
construcción es cuasirrefleja pasiva. En b, el pronombre \forma{se} materializa
uno de los argumentos del verbo, la meta, y cumple función de COI. Se trata de
una estructura refleja recíproca, que admite el refuerzo \forma{mutuamente /
unos a otros}.
\end{solucion}

\begin{solucion}{\ref{ej:10-akasha}}
\forma{(El aire) me rodea; (el aire) me eleva; (mi acompañante) me mira; (un
tren moderno) me despierta una inmensa curiosidad; no me importa (eso); (la
puerta) me saca del ensueño; me encanta ese lugar} y \forma{el silencio me hace
sentir incómoda}: a) la referencia del pronombre no coincide con la del sujeto y
la construcción no es pronominal.

\forma{Me pregunto (\ldots)}: a) la referencia del pronombre coincide con la del
sujeto y la construcción es pronominal. b) El pronombre manifiesta al paciente y
cumple la función de COI. i) La construcción es refleja pura porque acepta el
refuerzo \forma{a mí mismo/a}.

\forma{(El poder espiritual de carácter omnipresente) se encuentra\ldots}: a) la
referencia del pronombre coincide con la del sujeto y la construcción es
pronominal. b) El pronombre no manifiesta ningún argumento del verbo. ii) El
pronombre produce cambio en la estructura argumental (elimina al agente y
promueve al tema a sujeto de la oración con lectura de afectación). La
construcción es cuasirrefleja pasiva porque admite tercera persona singular y
plural: \forma{Los poderes espirituales de carácter omnipresente se encuentran}.

\forma{Se forma el espíritu; mis dedos se deslizan por el espacio; (mis brazos)
se convierten en alas; la puerta se abre} y \forma{(la puerta) se abrió; no me
levanto; mi sonrisa se borra; el tren apenas se mueve}: a) la referencia del
pronombre coincide con la del sujeto y se trata de construcciones pronominales.
b) El pronombre no manifiesta ningún argumento del verbo. ii) El pronombre
produce cambio en la estructura argumental (elimina al agente y promueve al tema
a sujeto de la oración, con lectura de proceso). Las construcciones son
cuasirreflejas medias porque admiten el refuerzo \forma{solo/a/os/as} o \forma{por
sí solo/a/os/as}.

\textit{Aclaraciones}: 1. \forma{Convertir} es un verbo de transformación, que
exige tres papeles temáticos ---agente, tema y meta---: \forma{Yo convierto mis
brazos en alas} (voz activa) $\rightarrow$ \forma{Mis brazos se convierten en
alas} (voz media). La meta es un CPR. 2. Con los verbos de cambio de posición o
ubicación, como \forma{levantar}, cuando el sujeto es animado se ha considerado
que es a la vez agente y paciente.
\end{solucion}

% =====================================================================
%  capitulos/cap11.tex — Capítulo 11. Las perífrasis verbales
%  Autores: Mabel Giammatteo, Ana Marcovecchio y Andrés Kaller
%  Fuente: Capitulo_11_normalizado.docx
% =====================================================================
\capitulo{Perífrasis y locuciones verbales}{Mabel Giammatteo, Ana Marcovecchio y Andrés Kaller}

\section{Introducción}

Hasta aquí, la estructura oracional se ha presentado articulada en torno a una
única forma verbal, en la que coexisten la información léxica, aportada por la
raíz, y el contenido gramatical, expresado por la flexión. En este capítulo nos
ocupamos de los casos en que las unidades verbales son más complejas y están
constituidas por un verbo que funciona como centro léxico-semántico de la
predicación y otro que lo auxilia, manifestando la información gramatical
(tiempo, modo y aspecto) y las categorías de concordancia (persona y número).
Esas unidades que funcionan conjuntamente como núcleo del SV reciben la
denominación de \termino{perífrasis verbales} (PPVV) y, según el aporte de
significado que hace el auxiliar, se especializan en manifestar diferentes
valores relacionados con el evento: aspectuales (\forma{empezar a leer, seguir
luchando}), modales (\forma{tener que estudiar, poder descifrar}), temporales
(\forma{ir a comer, haber disfrutado}) o de diátesis pasiva (\forma{ser
premiado}). En este capítulo nos ocupamos también de los casos más complejos, en
los que se forman cadenas de auxiliares que preceden al verbo principal o
auxiliado, así como de distinguir las PPVV de las \termino{locuciones verbales}
(LLVV), que también son unidades verbales complejas, pero cuyo significado no
deriva composicionalmente del de sus constituyentes, por lo que su
interpretación no es literal (así, por ejemplo, \forma{dar fe} significa
\sig{atestiguar} y \forma{echar de menos} se interpreta como \sig{extrañar}).

\section{Concepto de perífrasis verbal}\label{sec:11-concepto}

Según \textcite[3425]{gomeztorrego1999}:

\begin{cita}
Una perífrasis verbal es la unión de dos o más verbos que constituyen un solo
núcleo del predicado. El primer verbo, llamado «auxiliar», comporta las
informaciones morfológicas de número y persona, y se conjuga en todas (o en
parte de) las formas y tiempos de la conjugación. El segundo verbo, llamado
«auxiliado» o «principal», debe aparecer en infinitivo, gerundio o participio,
es decir, es una forma no personal. Según se trate de una u otra forma, hablamos
de perífrasis de infinitivo, de gerundio o de participio. A veces contiene
partículas de enlace como \textit{a}, \textit{de}, \textit{por} y \textit{que}.
\end{cita}

Las perífrasis verbales se constituyen por un proceso de unificación de dos
verbos: el \termino{auxiliar} y el verbo \termino{principal} o
\termino{auxiliado}, que transmite el contenido léxico de la predicación. Los
auxiliares pueden ser verbos con escaso contenido semántico (por ejemplo, los
copulativos\footnote{La noción de verbo copulativo se introdujo en el
\capref{9}.} \forma{ser} y \forma{estar} y otros, como \forma{haber} y
\forma{soler}) o pueden originarse en verbos de significado léxico pleno (por
ejemplo, \forma{deber, ir, andar}) que, en su empleo como auxiliares, sufren
alguna reducción de su contenido y pasan a interpretarse con un significado
gramatical más abstracto, característico de los auxiliares. El verbo principal,
en cambio, desprovisto de flexión y en una forma no finita (infinitivo,
participio o gerundio), actúa como verbo pleno, ya que mantiene sus valores
semánticos y requisitos argumentales regulares. Así, en \exref{ex:11-fiesta}, es
\forma{comer} el que selecciona el tema ---\forma{sándwiches y bocaditos}--- y
no el auxiliar \forma{estar}; en \exref{ex:11-mozos}, ambos argumentos, tema
---\forma{los alimentos}--- y locativo ---\forma{en las mesas}---, son exigidos
por el verbo de colocación \forma{poner} y no por el auxiliar \forma{empezar}.
Además, con ambos verbos, son los auxiliados los que exigen un agente como
sujeto.

\begin{exe}
  \ex\label{ex:11-fiesta} Todos los invitados a la fiesta \textit{estaban
      comiendo} sándwiches y bocaditos.
  \ex\label{ex:11-mozos} Los mozos \textit{empezaron a poner} los alimentos en
      las mesas muy temprano.
\end{exe}

Según vemos en \exref{ex:11-mozos}, entre las dos formas verbales que integran la
perífrasis puede haber otro elemento, que puede ser una preposición, como
\forma{a} ---\forma{empezaron a correr}--- o \forma{de} ---\forma{terminaron de
decir}---, y también la conjunción \forma{que} ---\forma{tener que salir}---.

\subsection{Criterios de reconocimiento}\label{sec:11-criterios}

\subsubsection{La referencia a un evento único}\label{sec:11-evento}

Para que la unión de los dos verbos, auxiliar y auxiliado, constituya una
perífrasis, ambos deben referir a un único evento. Veamos algunos ejemplos:

\begin{exe}
  \ex\label{ex:11-juan}
    \begin{xlist}
      \ex\label{ex:11-juan-a} Juan \textit{está estudiando} en su habitación.
      \ex\label{ex:11-juan-b} Juan \textit{estudia} en su habitación.
    \end{xlist}
  \ex\label{ex:11-alumnos}
    \begin{xlist}
      \ex\label{ex:11-alumnos-a} Los alumnos \textit{deben asistir} con uniforme.
      \ex\label{ex:11-alumnos-b} Los alumnos \textit{asistirán} con uniforme.
    \end{xlist}
\end{exe}

En \exref{ex:11-juan-a}, \forma{está estudiando} equivale a la forma simple del
verbo, \forma{estudia} \exref{ex:11-juan-b}, a la que añade un matiz de
\sig{duración}; en \exref{ex:11-alumnos-a}, \forma{deben asistir} es equivalente
a \forma{asistirán} \exref{ex:11-alumnos-b}, al que agrega el valor de
\sig{obligación}. En ambos casos hay un único evento.

\subsubsection{La conmutación estructural}\label{sec:11-estructural}

Sin embargo, no siempre que existen dos verbos seguidos, uno conjugado y otro en
una forma no finita, ya sea participio, infinitivo o gerundio, se trata de una
perífrasis. Para establecer las diferencias podemos comparar las oraciones de
\exref{ex:11-problema} y \exref{ex:11-profesor}:

\begin{exe}
  \ex\label{ex:11-problema}
    \begin{xlist}
      \ex\label{ex:11-problema-a} El problema \textit{fue resuelto}.
      \ex\label{ex:11-problema-b} El problema \textit{se resolvió}. $\rightarrow$
          Un solo evento.
      \ex[*]{\label{ex:11-problema-c} El problema fue que resuelto.}
          Rechazo de la conmutación estructural.
    \end{xlist}
  \ex\label{ex:11-profesor}
    \begin{xlist}
      \ex\label{ex:11-profesor-a} El profesor \textit{decidió posponer} el
          examen.
      \ex\label{ex:11-profesor-b} El profesor \textit{decidió (tomó la decisión) y
          pospuso} el examen. $\rightarrow$ Dos eventos.
      \ex\label{ex:11-profesor-c} El profesor \textit{decidió que posponía} el
          examen. Aceptación de la conmutación estructural.
    \end{xlist}
\end{exe}

Una prueba que podemos aplicar para distinguir estos casos es la de
\termino{conmutación}\footnote{Se entiende por \textit{conmutación} la operación
que consiste en cambiar un elemento por un equivalente funcional para comprobar
si produce o no cambio de significado. Por ejemplo, \forma{Llegó triste} puede
ser reemplazado por \forma{Llegó con tristeza} o \forma{Llegó apenado}. En el
\capref{3} esta prueba se aplicó para contrastar sonidos y corroborar si se
trataba de fonemas diferentes o de alófonos del mismo fonema.}
\termino{estructural}, que consiste en cambiar la estructura sintáctica para ver
si se trata de dos formas unificadas, como en las perífrasis, o de dos formas
independientes, cada una de las cuales remite a un evento diferente. En
\exref{ex:11-problema-a} tenemos una perífrasis porque ambas formas verbales
remiten a un único evento, de modo que es posible cambiar \forma{fue resuelto}
por otra forma simple equivalente, \forma{se resolvió} \exref{ex:11-problema-b};
pero un cambio estructural, como el que aparece en \exref{ex:11-problema-c},
produce un resultado agramatical. El caso de \exref{ex:11-profesor-a}, en cambio,
no tiene una interpretación de evento único, como vemos en
\exref{ex:11-profesor-b}, y, por lo tanto, acepta el cambio por otra estructura
\exref{ex:11-profesor-c}, lo que nos permite comprobar que no se trata de una
perífrasis.

En las PPVV, auxiliar y auxiliado forman una unidad compleja que funciona como
núcleo del SV \exref{ex:11-problema-a}, como vemos en \exref{ex:11-de-pv}:

\begin{exe}
  \ex\label{ex:11-de-pv} {}[\etq{SF/O} [\etq{SN/Suj} El problema] [\etq{SV/Pred}
      [\etq{V/n} fue resuelto] [\etq{SP/CAg} por los estudiantes]]]
\end{exe}

En cambio, en el caso de \exref{ex:11-profesor-a}, cada uno de los verbos
pertenece a una oración diferente: la forma conjugada es el núcleo del predicado
verbal de la oración principal y la forma no finita constituye el verbo núcleo de
otra oración, subordinada,\footnote{El concepto de oración subordinada se
desarrolla en el \capref{13}.} tal como vemos en \exref{ex:11-de-nopv}:

\begin{exe}
  \ex\label{ex:11-de-nopv} {}[\etq{SF/O} [\etq{SN/Suj} El profesor]
      [\etq{SV/Pred} decidió\etq{V/n} [\etq{SC=SF/COD} posponer el
      examen]]]\pendiente{El original analiza aquí «resolvió aplazar»}
\end{exe}

En algunos casos de lectura ambigua, como \exref{ex:11-perdon-a}, la conmutación
estructural permite anclar una interpretación: en \exref{ex:11-perdon-b} se
interpreta \forma{volvieron a pedir} como un evento que se repite (=
\sig{pidieron de nuevo}), por lo que se trata de una perífrasis; pero en
\exref{ex:11-perdon-c}, con \forma{volver} acompañado de un complemento locativo
---\forma{volvieron al lugar}---, se disuelve la posibilidad de lectura como
perífrasis, ya que admite conmutación estructural: el infinitivo se puede
reemplazar por un SP encabezado por \forma{para} ---\forma{volvieron para pedir
perdón}---.

\begin{exe}
  \ex
    \begin{xlist}
      \ex\label{ex:11-perdon-a} \textit{Volvieron a pedir} perdón.
      \ex\label{ex:11-perdon-b} Reiteraron sus disculpas / Pidieron perdón de
          nuevo. $\rightarrow$ Un solo evento.
      \ex\label{ex:11-perdon-c} Volvieron al lugar y pidieron perdón.
          $\rightarrow$ Dos eventos.
    \end{xlist}
\end{exe}

\begin{ejercicio}\label{ej:11-reconocer}
En las siguientes oraciones, determine si se trata o no de PPVV o si hay
ambigüedad. Según se muestra en los ejemplos resueltos de abajo, fundamente sus
respuestas aplicando las pruebas de reconocimiento suministradas (referencia a
un evento único, admisión o no de conmutación estructural). En los casos
ambiguos, para las PPVV ofrezca una paráfrasis adecuada del significado.

\begin{exe}
  \ex La sentencia \textit{ha salido} del juzgado hoy.
  \ex Los deportistas \textit{vienen corriendo} más rápido.
  \ex Mi abuela no \textit{salió a pasear} en la pandemia.
  \ex \textit{Volvió a hacer} el trabajo práctico de la materia.
  \ex El domingo \textit{empezó a llover} desde temprano.
\end{exe}

\begin{resuelto}[Ejemplos resueltos]
1. Es PV porque señala un único evento y no admite conmutación estructural:
\juicio{*}{\forma{La sentencia ha y salió del juzgado hoy}}. 2. Es ambigua y
posee dos interpretaciones: i) no es PV porque indica dos eventos y admite
conmutación estructural: \forma{Los deportistas vienen (al club) y en el
trayecto apuran su marcha}; ii) es PV y no admite conmutación estructural.
Paráfrasis: \forma{Los deportistas mejoran progresivamente sus habilidades para
la carrera}.
\end{resuelto}
\end{ejercicio}

\subsubsection{La conmutación léxica}\label{sec:11-lexica}

Aunque no admiten conmutación estructural, las PPVV aceptan \termino{conmutación
léxica}. Esta prueba, que consiste en cambiar una palabra por otra de la misma
categoría, se aplica a las PPVV mediante la conmutación del verbo principal por
otro de significado equivalente:

\begin{exe}
  \ex
    \begin{xlist}
      \ex Ana \textit{debe comer} mejor porque está anémica.
      \ex Ana \textit{debe alimentarse} mejor porque está anémica.
    \end{xlist}
  \ex
    \begin{xlist}
      \ex El famoso actor \textit{fue premiado} por su labor en la nueva serie
          de televisión.
      \ex El famoso actor \textit{fue galardonado} por su labor en la nueva
          serie de televisión.
    \end{xlist}
  \ex
    \begin{xlist}
      \ex Todos \textit{continuaron mirando} apasionadamente la nueva temporada
          de la serie.
      \ex Todos \textit{continuaron viendo} apasionadamente la nueva temporada
          de la serie.
    \end{xlist}
\end{exe}

\begin{ejercicio}\label{ej:11-series}
Lea atentamente el texto siguiente y resuelva lo pedido en las consignas que se
presentan debajo.

% Texto periodístico protegido: se incorpora a mano desde el original.
\textoaincorporar{el fragmento adaptado de \textit{Infobae} (08/11/2020) sobre
dos tipos de series, con las formas en cursiva: \textit{nos hacen sentir; son
devoradas; puede distraer; comenzarlas a ver; permiten conectarnos; suelen
proveer}}

\begin{consignas}
  \item Determine si las formas en cursiva constituyen o no PPVV. Aporte pruebas
        que fundamenten sus afirmaciones (referencia o no a un evento único,
        conmutación léxica y conmutación estructural), según se muestra en el
        ejemplo resuelto~1.
  \item Reescriba el fragmento incorporando los sinónimos obtenidos en la
        conmutación léxica, sin alterar el significado del texto.
\end{consignas}

\begin{resuelto}
1. \forma{(Nos) hacen sentir} no es PV. Se trata de dos eventos; admite
conmutación léxica, \forma{(nos) hacen experimentar}, y conmutación estructural,
\forma{hacen que sintamos}.
\end{resuelto}
\end{ejercicio}

\section{Criterios de clasificación}\label{sec:11-clasificacion}

Las perífrasis verbales suelen clasificarse teniendo en cuenta dos criterios
complementarios: a) la forma no finita que interviene en la perífrasis y b) el
aporte semántico del auxiliar.

\begin{itemize}
  \item Según la forma no finita del verbo auxiliado,\pendiente{El original dice
        «la forma no finita del auxiliar que acompaña al verbo auxiliado»}
        tenemos:
        \begin{itemize}
          \item \termino{perífrasis de infinitivo}: \forma{volver a leer, empezar a
                escribir, poder entender}, etc.;
          \item \termino{perífrasis de gerundio}: \forma{estar jugando, seguir
                estudiando}, etc.;
          \item \termino{perífrasis de participio}: \forma{haber leído, ser
                vacunado}, etc.
        \end{itemize}
\end{itemize}

\begin{destacado}
\textbf{Para tener en cuenta.} El comportamiento del participio en cuanto a la
concordancia no es el mismo en todas las perífrasis: mientras que con los tiempos
compuestos es invariable \exref{ex:11-primos}, en la perífrasis de voz pasiva el
participio concuerda en género y número con el sujeto de la oración
\exref{ex:11-francia} y, en las perífrasis con verbos como \forma{tener} o
\forma{llevar}, concuerda con el objeto \exref{ex:11-libros}.

\begin{exe}
  \ex\label{ex:11-primos} Participio invariable: Mis primos / mis primas
      \textit{han visitado} Francia / varios países.
  \ex\label{ex:11-francia} Participio concordado con el sujeto:
    \begin{xlist}
      \ex Francia \textit{fue visitada} por mis primos.
      \ex El país de Luis XV \textit{fue visitado} por mis primos.
    \end{xlist}
  \ex\label{ex:11-libros} Participio concordado con el objeto:
    \begin{xlist}
      \ex \textit{Tengo / Llevo leídos} más de 20 libros sobre el tema.
      \ex \textit{Tengo / Llevo leídas} más de 20 novelas de este autor.
    \end{xlist}
\end{exe}
\end{destacado}

\begin{itemize}
  \item Según el aporte semántico del auxiliar, se reconocen:
        \begin{itemize}
          \item perífrasis de diátesis: \forma{ser leído};
          \item perífrasis aspectuales: \forma{empezar a leer, continuar
                escuchando};
          \item perífrasis temporales: \forma{ir a comer, haber dicho};
          \item perífrasis modales: \forma{deber escribir, tener que mirar, poder
                ser}.
        \end{itemize}
\end{itemize}

\begin{recuadro}[Para ampliar]
Con respecto a las categorías verbales (tiempo, aspecto y modo), las formas no
finitas solo manifiestan aspecto, no de forma flexional, sino inherente, es
decir, por su propia naturaleza. Como señala \textcite{hernanz1999}, el
infinitivo es neutro en relación con las otras dos formas no finitas, que marcan
valores específicos: el gerundio presenta el proceso en desarrollo, por lo
tanto, es durativo, y el participio indica el proceso como terminado; de ahí su
sentido perfectivo-resultativo. Del valor aspectual de cada forma no finita se
deriva su utilización en las PPVV:

\begin{itemize}
  \item el \termino{infinitivo}, que es neutro, es el que más combinaciones
        admite y puede manifestar valores tales como el inceptivo
        ---\forma{empezó a cantar}---, el terminativo ---\forma{dejó de
        cantar}---, el habitual ---\forma{suele cantar}--- o el de repetición
        ---\forma{volvió a cantar}---;
  \item el \termino{gerundio} se emplea en perífrasis con valor durativo, como
        la progresiva ---\forma{estoy cantando}--- y la continuativa
        ---\forma{siguió cantando}---;
  \item el \termino{participio}, que es perfectivo-resultativo, se utiliza en los
        tiempos compuestos ---\forma{había cantado}--- y en la voz pasiva
        ---\forma{fue cantado}--- \parencite[50--51]{giammatteo2009}.
\end{itemize}
\end{recuadro}

\begin{ejercicio}\label{ej:11-tato}
Considere los siguientes fragmentos y resuelva lo pedido en las consignas, según
se muestra en el ejemplo resuelto~1.

\begin{consignas}
  \item Determine si las combinaciones de palabras subrayadas son perífrasis
        verbales o no.
  \item En las que son perífrasis, haga la descripción estructural de sus
        constituyentes.
\end{consignas}

% Textos protegidos (monólogos de Tato Bores): se incorporan a mano desde el original.
\textoaincorporar{los seis fragmentos adaptados de monólogos de Tato Bores
(cultura.gob.ar, 12/01/2022), con las combinaciones subrayadas: 1. \textit{se
tuvieron que aguantar; (no) crea que le miento; voy a detallar}; 2.
\textit{iremos a comprar; van a tener que entregar; dicen ser; están corriendo};
3. \textit{estamos viviendo; estamos diciendo}; 4. \textit{estoy aburrida; dicen
defenderme}; 5. \textit{iremos a pagarle; van a tener que entregar}; 6.
\textit{estamos pagando; vamos a pagar}}

\begin{resuelto}
1. \forma{Se tuvieron que aguantar} es PV. Descripción estructural (DE): V
auxiliar \forma{tener} + \forma{que} + V auxiliado en infinitivo
\forma{aguantarse}. \forma{(No) crea que le miento}: no es PV. \forma{Voy a
detallar}: es PV. DE: V auxiliar \forma{ir} + \forma{a} + V auxiliado en
infinitivo \forma{detallar}.
\end{resuelto}
\end{ejercicio}

\section{Principales características de los auxiliares de perífrasis}\label{sec:11-auxiliares}

Los auxiliares de perífrasis verbales presentan las siguientes características:

\paragraph{Carácter delimitado de la clase.} Los verbos auxiliares que componen
las perífrasis verbales son un conjunto relativamente acotado de pocos miembros
que se repiten en la mayoría de las lenguas. Según se plantea en
\textcite[205]{giammatteo2018marcovecchio}:

\begin{cita}
[\ldots] la comparación entre lenguas no necesariamente emparentadas entre sí
nos muestra que no cualquier verbo deviene auxiliar. Así, no esperamos, ni
encontramos en las distintas lenguas del mundo, verbos con significados
específicos como, por ejemplo, \sig{comer}, \sig{leer} o \sig{cocinar}, que
hayan adoptado una función auxiliar. Por lo general, los auxiliares son verbos
de significados más generales (Heine, 1991), como \textit{do} (ingl.),
\textit{hacer} (esp.); o bien que señalan \sig{posesión}, como \textit{tener}
(esp.), \textit{have} (ingl.) y \textit{avoir} (fr.). Este tipo de verbos son
buenos candidatos para la auxiliarización.
\end{cita}

Además de los mencionados, otro grupo de verbos que frecuentemente se emplean
como auxiliares en las distintas lenguas son los de movimiento ---\forma{ir,
llevar, andar, venir, seguir}---, que pasan de señalar \sig{desplazamiento} en el
orden espacial a referir al \sig{desarrollo} del evento, ya sea en el tiempo
\parencite{cuenca1999} ---\forma{voy a comer} (valor futuro: \forma{comeré})---
o en relación con diferentes valores aspectuales ---\forma{voy / vengo / ando /
sigo / continúo estudiando}---.

\paragraph{Desemantización.} Algunos de los auxiliares, como \forma{ser, estar} y
\forma{haber}, son verbos con muy poco significado léxico inherente; otros, como
\forma{ir, venir} o \forma{seguir}, son verbos de significado pleno, que en sus
usos auxiliares se considera que han sufrido un proceso de
\termino{desemantización}. Según lo explica \textcite[21]{garciafernandez2006},
«[d]esde el punto de vista semántico, las construcciones perifrásticas suelen ser
descritas por el proceso de desemantización que experimenta el auxiliar y que
puede ilustrarse fácilmente con la oración \textit{Va a nevar}, donde la forma
\textit{va} no denota movimiento alguno». En el caso de estos auxiliares
desemantizados, es necesario destacar que en sus otros usos continúan
funcionando como verbos plenos y que pierden (parte de) su significado solo
cuando se combinan con otra forma verbal no finita, es decir que «la pérdida de
su significado viene provocada por factores colocacionales\ldots»
\parencite[124]{tsutahara2022}.

Para determinar si una secuencia de verbo conjugado más otro no finito constituye
una perífrasis verbal o si se trata de una secuencia de dos verbos, es
indispensable evaluar si el verbo conjugado mantiene su significado léxico pleno
o si sufrió algún tipo de modificación o pérdida. Por ejemplo, en
\exref{ex:11-mariano}, \forma{va} no indica \sig{desplazamiento}, sino
\sig{acción futura} y, en \exref{ex:11-sigue}, \forma{seguir} perdió el valor
espacio-temporal de \sig{ir detrás o después} para pasar a significar
\sig{continuar}.

\begin{exe}
  \ex\label{ex:11-mariano} Mariano \textit{va a publicar} su tesis.
  \ex\label{ex:11-sigue} \textit{Sigue trabajando} en una gran empresa.
\end{exe}

Sin embargo, en otros casos, la oposición entre verbos plenos y verbos
auxiliares «puede ser demasiado radical, ya que no deja lugar para los verbos que
se comportan parcialmente como auxiliares» \parencite{bosque2019}. Por ejemplo,
se considera que \forma{querer} en \exref{ex:11-salsa-a} no constituye una
perífrasis verbal porque no experimenta una interpretación distinta de su
significado de \sig{desear}; de ahí que estemos frente a un empleo de
\forma{querer} como verbo pleno y que, por lo tanto, esperaríamos que acepte la
prueba de conmutación estructural. Sin embargo, como muestra
\exref{ex:11-salsa-b}, el resultado de la conmutación estructural es una oración
agramatical.

\begin{exe}
  \ex\label{ex:11-salsa}
    \begin{xlist}
      \ex\label{ex:11-salsa-a} Yo \textit{quiero bailar} salsa.
      \ex[*]{\label{ex:11-salsa-b} Yo \textit{quiero} que (yo) \textit{baile}
          salsa.}
    \end{xlist}
\end{exe}

No obstante, lo que realmente impide aplicar la prueba de conmutación estructural
en \exref{ex:11-salsa-b} es la coincidencia de sujeto entre \forma{querer} y
\forma{bailar}: ambos verbos tienen el mismo sujeto, \forma{yo}. Cuando no hay
coincidencia de persona entre el sujeto de la forma personal y el de la forma no
finita, como se muestra en \exref{ex:11-ana}, \forma{querer} + infinitivo acepta
conmutación estructural.

\begin{exe}
  \ex\label{ex:11-ana} Yo \textit{quiero} que Ana \textit{baile} salsa.
\end{exe}

En \exref{ex:11-ana}, como se trata de dos personas gramaticales distintas para
cada verbo ---\forma{yo} para \forma{quiero} y \forma{ella} para
\forma{baile}---, la conmutación estructural es posible, por lo que queda claro
que se trata de dos eventos distintos y que no hay perífrasis. Pero ¿qué sucede
en \exref{ex:11-llover}?

\begin{exe}
  \ex\label{ex:11-llover} Parece que \textit{quiere llover}.
\end{exe}

En \exref{ex:11-llover} no le podemos asignar el deseo a ningún individuo. Por
tanto, asumimos que \forma{querer} ya no actúa como un verbo pleno que selecciona
a su sujeto, sino que ha adquirido un significado diferente al de \sig{desear}:
el de sugerir \sig{la inminencia de la lluvia}. Ambas características hacen que
en su empleo con verbos impersonales \forma{querer} se considere un auxiliar.
Como plantea \textcite[241]{ditullio2014}, «los auxiliares son semánticamente
transparentes» y, por eso, a diferencia de los verbos plenos, que imponen
restricciones a sus sujetos \exref{ex:11-escuela}, pueden construirse con verbos
impersonales y aceptar cualquier tipo de SN como sujeto
\exref{ex:11-complicado}.

\begin{exe}
  \ex\label{ex:11-escuela} Juan / *el año \textit{va} temprano a la escuela.
  \ex\label{ex:11-complicado} Juan / el año \textit{va a continuar} complicado.
\end{exe}

\paragraph{Integración con el verbo principal.} El grado de integración entre el
verbo auxiliar y el auxiliado se desprende de ciertas propiedades que actúan
sobre el conjunto, como el alcance de los adjuntos. Cuando tenemos dos verbos
independientes que no forman perífrasis, cada uno de ellos puede recibir
modificadores adjuntos \exref{ex:11-mariana1}; en las perífrasis, en cambio,
auxiliar y auxiliado deben ser modificados conjuntamente \exref{ex:11-mariana2}.

\begin{exe}
  \ex\label{ex:11-mariana1} Hoy Mariana decidió \textit{ir} mañana a la reunión en
      la facultad.
  \ex\label{ex:11-mariana2}
    \begin{xlist}
      \ex Hoy Mariana \textit{va a ir} a la facultad.
      \ex[*]{Hoy Mariana \textit{va a ir} mañana a la facultad.}
    \end{xlist}
\end{exe}

Asimismo, a diferencia de lo que sucede con la secuencia de dos verbos
\exref{ex:11-neg-a}, el auxiliar no admite ser negado en forma independiente del
verbo principal \exref{ex:11-neg-c}.\footnote{No obstante, según veremos en
§\ref{sec:11-modales}, \forma{poder} epistémico admite doble negación: en el
verbo principal y en el auxiliar.}

\begin{exe}
  \ex
    \begin{xlist}
      \ex\label{ex:11-neg-a} Mariana no decidió no \textit{ir} mañana a la
          facultad.
      \ex Mariana no \textit{va a ir} a la facultad.
      \ex[*]{\label{ex:11-neg-c} Mariana no \textit{va a} no \textit{ir} a la
          facultad.}
    \end{xlist}
\end{exe}

Otra muestra del grado de integración entre el auxiliar y el verbo principal es
que, a diferencia de lo que sucede cuando se trata de una secuencia de dos verbos
independientes \exref{ex:11-imputado}, las perífrasis no aceptan la omisión del
verbo principal ni en la interrogación \exref{ex:11-suele-b} ni en la
focalización con estructuras de relieve\footnote{El concepto de foco y los
recursos mediante los que se manifiesta se presentaron en el \capref{6},
§\ref{sec:06-capas}.} \exref{ex:11-suele-d}. No obstante, si bien en las PPVV el
verbo principal no puede elidirse, puede ser reemplazado por \forma{hacer}, tanto
en la interrogación \exref{ex:11-suele-c} como en la focalización
\exref{ex:11-suele-e}.

\begin{exe}
  \ex\label{ex:11-imputado}
    \begin{xlist}
      \ex El imputado \textit{prometió} decir la verdad.
      \ex ¿Qué \textit{prometió} el imputado? $\rightarrow$ Interrogación con
          omisión del verbo principal.
      \ex Lo que el imputado \textit{prometió} es decir la verdad. $\rightarrow$
          Focalización con omisión del verbo principal.
    \end{xlist}
  \ex\label{ex:11-suele}
    \begin{xlist}
      \ex El imputado \textit{suele mentir} con gran facilidad.
      \ex[*]{\label{ex:11-suele-b} ¿Qué \textit{suele} el imputado con gran
          facilidad?} $\rightarrow$ Interrogación sin aceptar la omisión del verbo
          principal.
      \ex\label{ex:11-suele-c} ¿Qué \textit{suele hacer} el imputado con gran
          facilidad? $\rightarrow$ Interrogación con reemplazo del verbo principal
          por el proverbo \forma{hacer}.
      \ex[*]{\label{ex:11-suele-d} Lo que el imputado \textit{suele} es mentir con
          gran facilidad.} $\rightarrow$ Focalización sin aceptar la omisión del
          verbo principal.
      \ex\label{ex:11-suele-e} Lo que el imputado \textit{suele hacer} es mentir
          con gran facilidad. $\rightarrow$ Focalización con reemplazo del verbo
          principal por el proverbo \forma{hacer}.
    \end{xlist}
\end{exe}

No obstante la cohesión que existe entre los dos verbos, es posible la
intercalación del sujeto \exref{ex:11-oferta-a} y de determinados adverbios,
\exref{ex:11-oferta-b} y \exref{ex:11-oferta-c}, entre ambos
\parencite[243]{bosque2019}.

\begin{exe}
  \ex
    \begin{xlist}
      \ex\label{ex:11-oferta-a} ¿\textit{Puede} Juan \textit{rechazar} una oferta
          tan atractiva?
      \ex\label{ex:11-oferta-b} ¿Juan \textit{puede} realmente \textit{rechazar}
          una oferta tan atractiva?
      \ex\label{ex:11-oferta-c} \textit{Terminaron} recientemente \textit{de
          elaborar} el nuevo proyecto.
    \end{xlist}
\end{exe}

\paragraph{Comportamiento respecto de los clíticos.} En las oraciones con un
núcleo verbal simple, los clíticos de COD (\forma{lo, la, los, las}) y de COI
(\forma{le, les, se}) adoptan la posición preverbal.

\begin{exe}
  \ex
    \begin{xlist}
      \ex Ana \textit{cortó} la torta. / Ana la \textit{cortó}.
      \ex Ana le \textit{dio} una porción a Juan. / Ana se la \textit{dio}.
    \end{xlist}
\end{exe}

Las perífrasis de infinitivo y de gerundio, al actuar como una unidad verbal,
aceptan que el clítico se ubique antes del verbo conjugado, pero también admiten
que se ubique después de la forma no finita, \exref{ex:11-torta} y
\exref{ex:11-examenes}. Las perífrasis de participio, en cambio, solo aceptan la
anteposición de los clíticos \exref{ex:11-premio}.

\begin{exe}
  \ex\label{ex:11-torta} Ana \textit{comenzó a cortar} la torta. / Ana la
      \textit{comenzó a cortar}. / Ana \textit{comenzó a cortarla}.
  \ex\label{ex:11-examenes} Pedro \textit{está preparando} sus exámenes. / Pedro
      los \textit{está preparando}. / Pedro \textit{está preparándolos}.
  \ex\label{ex:11-premio} El director de la compañía \textit{ha recibido} un
      premio. / El director de la compañía lo \textit{ha recibido}. /
      \juicio{*}{\forma{El director de la compañía ha recibídolo.}}
\end{exe}

A diferencia de las PPVV, las secuencias verbales no perifrásticas no tienen un
comportamiento uniforme: algunos verbos, como \forma{querer, saber, necesitar,
intentar, tratar de}, aceptan anteposición y posposición \exref{ex:11-turistas},
pero otros, como \forma{decidir, lamentar, creer, pensar}, solo admiten
pronombres enclíticos \exref{ex:11-lamento}.

\begin{exe}
  \ex\label{ex:11-turistas} Los turistas \textit{quieren conocer} la ciudad. /
      Los turistas la \textit{quieren conocer}. / Los turistas \textit{quieren
      conocerla}.
  \ex\label{ex:11-lamento} Juan \textit{lamentó comprar} ese auto. / Juan
      \textit{lamentó comprarlo}. / \juicio{*}{\forma{Juan lo lamentó comprar.}}
\end{exe}

Para \textcite{ditullio2014}, la posición de los clíticos sirve como prueba
negativa para reconocer una perífrasis, porque si una forma finita seguida de un
infinitivo no permite la anteposición, no es perífrasis \exref{ex:11-lamento};
pero si la permite, puede tanto ser \exref{ex:11-torta} como no ser
\exref{ex:11-turistas} una verdadera perífrasis.

\paragraph{Pasivización.} Del mismo modo que con los verbos simples
correspondientes, las perífrasis verbales conformadas por un verbo principal que
tenga un sujeto y un objeto directo pueden pasivizarse:

\begin{exe}
  \ex
    \begin{xlist}
      \ex Juan \textit{suele bailar} \textit{La cumparsita}.
      \ex \textit{La cumparsita} \textit{suele ser bailada} por Juan.
    \end{xlist}
\end{exe}

En cambio, en el caso de las oraciones con un verbo conjugado más un verbo no
finito que no conforman perífrasis verbal, puede suceder que, si el verbo exige
un sujeto animado, la pasiva resulte agramatical:

\begin{exe}
  \ex
    \begin{xlist}
      \ex Juan \textit{quiere bailar} \textit{La cumparsita}.
      \ex[*]{\textit{La cumparsita} \textit{quiere ser bailada} por Juan.}
    \end{xlist}
\end{exe}

\begin{ejercicio}\label{ej:11-vf}
Indique si las siguientes afirmaciones son verdaderas o falsas y justifique sus
respuestas, según se muestra en el ejemplo resuelto~1.

\begin{exe}
  \ex En la mayoría de las lenguas, los verbos auxiliares que componen
      perífrasis verbales pertenecen a un conjunto relativamente cerrado de pocos
      miembros.
  \ex Los auxiliares de perífrasis verbales son verbos desemantizados, que en
      otros contextos pueden mantener su significado y actuar como verbos
      plenos.
  \ex Las perífrasis verbales de participio permiten la posposición del clítico.
  \ex Los auxiliares de perífrasis nunca pueden construirse con verbos
      impersonales y aceptan cualquier tipo de SN como sujeto.
\end{exe}

\begin{resuelto}
1. Verdadero. La mayoría de las lenguas suele recurrir a verbos con significados
generales, como \forma{ser, tener} o \forma{hacer}, y resulta poco probable que
utilicen como auxiliares a aquellos con significados más específicos, como
\forma{estudiar, saltar} o \forma{construir}.
\end{resuelto}
\end{ejercicio}

\begin{ejercicio}\label{ej:11-caracteristicas}
Teniendo en cuenta las principales características de los auxiliares de
perífrasis verbales señaladas en §\ref{sec:11-auxiliares}, explique por qué las
combinaciones de verbos presentes en las siguientes oraciones son o no PPVV,
según se muestra en el ejemplo resuelto~1.

\begin{exe}
  \ex Juan no querrá no ir a la escuela mañana.
  \ex ¿Qué prometió Juan? (Juan prometió viajar a Córdoba).
  \ex ¿Qué suele hacer Juan todas las mañanas?
  \ex Juan pudo recientemente comprar su casa.
  \ex María necesita saberlo todo.
\end{exe}

\begin{resuelto}
1. No es PV porque no existe integración entre los verbos y cada uno de ellos
aparece negado en forma independiente.
\end{resuelto}
\end{ejercicio}

\begin{recuadro}[Para ampliar: los auxiliares de perífrasis y los procesos de gramaticalización]
En las lenguas, los verbos más frecuentes suelen sufrir un «desgaste» que va
erosionando sus significados; así, por ejemplo, el verbo \forma{habere}, que en
latín indicaba \sig{posesión}, en las lenguas romances ha perdido parcialmente
ese valor y se ha convertido en un auxiliar para los tiempos compuestos: esp.
\forma{haber}, fr. \forma{avoir}, it. \forma{avere}, rum. \forma{avea}. En estos
casos se dice que el verbo ha sufrido un proceso de
\termino{gramaticalización}, entendido como «un cambio lingüístico por el que una
voz con significado léxico adquiere valor gramatical en su evolución histórica»
\parencite[147]{bosque2019}. En el caso de \forma{haber}, en su paso del latín
al español, su significado de \sig{posesión} se ha perdido completamente.
Asimismo, también puede suceder que, como resultado de un proceso de
gramaticalización, un verbo ya de muy poco contenido, como \forma{ser}, que,
según vimos en el \capref{9}, exige un predicativo para completar su
significado, «adquiera una función aún más gramatical»
\parencite{company2003}, como cuando se convierte en el soporte de las
categorías de tiempo, modo, aspecto y concordancia (número y persona) en la voz
pasiva. En otros casos, lo que sucede es que el verbo mantiene usos con
significado pleno (\forma{Anda por toda la ciudad en moto} = \sig{se desplaza})
y otros en los que se ha desemantizado, es decir, ha perdido parte de su
significado léxico (\forma{Anda rumoreando que habrá cambios en el gobierno}).
En este último ejemplo, \forma{andar} ya no implica \sig{desplazamiento}, pero
mantiene un valor referido a que se trata de un evento continuativo-reiterativo
\parencite{yllera1999}.

Según vemos, entonces, no en todas las perífrasis el auxiliar muestra el mismo
grado de pérdida de significado léxico; por eso las características de las
perífrasis verbales que se han venido señalando dependen en gran medida del grado
de consolidación de la estructura perifrástica, condicionada, a su vez, por el
estadio de gramaticalización que haya alcanzado el auxiliar. Es decir, nos
enfrentamos a un \textit{continuum} en el que en un extremo encontramos verbos
que solo pueden tener empleo auxiliar, como \forma{soler}; luego, otros como
\forma{ir} o \forma{deber}, que funcionan como auxiliares, pero también retienen
su uso con valor pleno, y, en el polo opuesto, se encuentran ciertos verbos
plenos que comienzan a mostrar indicios de desemantización. Así, por ejemplo,
\forma{prometer} se comporta como verbo pleno en \forma{El deportista promete
llegar temprano a los entrenamientos}, donde la unión del verbo conjugado y la
forma no finita no constituye perífrasis porque admite conmutación estructural:
\forma{El deportista promete que llegará temprano}. Sin embargo, en \forma{Este
deportista promete ser la estrella del equipo}, el mismo verbo se muestra en vías
de auxiliarización, ya que no expresa ninguna promesa por parte del deportista,
sino que ha adquirido un nuevo significado por el cual el hablante expresa una
deducción anticipatoria sobre el lugar de prominencia que posiblemente alcance el
deportista en el futuro. En este segundo caso, la oración se interpreta como
\sig{Este deportista posiblemente será la estrella del equipo}.

Considerado en perspectiva histórica, el proceso por el que un verbo pleno se
convierte en un auxiliar de PV puede analizarse de la siguiente manera. Tomemos,
por ejemplo, un verbo originariamente léxico como \forma{volver}, que supone un
\sig{desplazamiento físico de retorno a un lugar}, lo que se manifiesta por el
verbo seguido de un SP con valor de dirección, como en \forma{Volvió a la
escuela}. Paulatinamente, al comenzar a emplearse con un infinitivo como
complemento de la preposición, ha ido sumando una nueva interpretación, como en
\forma{María vuelve a leer el libro}, donde es posible entender tanto que
\sig{María vuelve / regresa a un lugar determinado para leer el libro}
(\forma{María volvió a la biblioteca para leer el libro}) como que la reiteración
no se aplica al movimiento de regreso, sino a la repetición del evento
(\forma{María leyó el libro nuevamente}). Cuando se produce una nueva lectura, la
consolidación de esa interpretación tiene como contraparte la rigidización de la
estructura sintáctica. Hay una nueva correspondencia entre la forma y el
significado, y la estructura queda cristalizada como \textit{verbo auxiliar +
verbo auxiliado}. En otras palabras, mientras el verbo se comporte como el núcleo
léxico-sintáctico de la predicación, posee mayor amplitud combinatoria:
\forma{María vuelve \{a la escuela / a leer / del extranjero\}}; en cambio,
cuando se convierte en auxiliar, la estructura se vuelve más rígida y solo acepta
ir seguida del infinitivo: \forma{El árbol volvió a \{florecer / perder las
hojas\}}. Por tanto, si bien hay ciertas perífrasis consolidadas, en otros
casos, según hemos visto, puede producirse ambigüedad y, en otros, puede tratarse
de un cambio aún en proceso. De ahí que resulte clave el análisis que realicemos
al enfrentarnos a una secuencia de verbo finito + verbo no finito, ya que no solo
las distintas combinaciones, según las pruebas que acepten, pueden ser o no
perífrasis, sino que un mismo verbo, según la interpretación que habilite el
contexto, puede tener una u otra lectura. Y aún más, como vimos con
\forma{prometer}, puede haber verbos que en su sentido más general no admitan las
pruebas de auxiliarización, pero que en algunos usos «más innovadores» comiencen
a mostrar cierto desplazamiento de su sentido original.
\end{recuadro}

\begin{ejercicio}\label{ej:11-corpes}
Considere las combinaciones verbales subrayadas en las siguientes oraciones y
resuelva las consignas según se muestra en el ejemplo resuelto~1.

\begin{consignas}
  \item Determine si las secuencias que conforman constituyen o no una perífrasis
        verbal.
  \item En cada caso, fundamente su respuesta.
\end{consignas}

\begin{exe}
  \ex Alguien le comentó alguna vez que la salida más fácil para terminar una
      historia \subrayado{era matando} al protagonista.
  \ex Ves que tu hermana no \subrayado{era descabellada} del todo.
  \ex Me pareció que el Ojo \subrayado{era sacudido} por un escalofrío.
  \ex El Bando Azul de Matanzas aparece como fundado en 1910 y vinculado al
      cabildo arará. Según se infiere de los datos compilados, \subrayado{comenzó
      siendo} un coro o bando de clave y derivó a uno de guaguancó.
  \ex Él es el verdadero autor de su propia gloria y uno no puede atribuirse más
      que la alegría, pero sería egoísta si no reconociera que yo lo ayudé a
      nacer. Cuando \subrayado{comenzó pude adivinar} que tenía enormes
      condiciones. Me maravilló su actitud iconoclasta.
  \ex Lo destapó, \subrayado{comenzó a beberlo} directamente de la botella y
      apagó el monitor de vigilancia.
\end{exe}

\noindent{\footnotesize Textos extraídos de CORPES XXI.}
\pendiente{Los subrayados no se conservan en el original; se reconstruyeron a
partir del solucionario}

\begin{resuelto}
1. \forma{Era matando}: no se trata de una PV porque el gerundio puede ser
reemplazado por el adverbio \forma{así}.
\end{resuelto}
\end{ejercicio}

\section{Contenido gramatical de las perífrasis verbales}\label{sec:11-contenido}

Según adelantamos, más allá de la conformación estructural que puedan adoptar,
las perífrasis verbales también se clasifican teniendo en cuenta el aporte de
significado del auxiliar, que manifiesta valores específicos vinculados ya sea
con el tiempo \exref{ex:11-valores-a}, el modo \exref{ex:11-valores-b}, el
aspecto \exref{ex:11-valores-c} o bien con la diátesis, que es la categoría que
se relaciona con la configuración que adoptan los argumentos del verbo en la
oración \exref{ex:11-valores-d}.

\begin{exe}
  \ex
    \begin{xlist}
      \ex\label{ex:11-valores-a} María \textit{va a leer} un libro en el verano.
          $\rightarrow$ PV de temporalidad: indica futuridad.
      \ex\label{ex:11-valores-b} María \textit{puede leer} un libro durante el
          verano. $\rightarrow$ PV de modalidad: indica posibilidad.
      \ex\label{ex:11-valores-c} María \textit{comenzó a leer} un libro durante
          el verano. $\rightarrow$ PV de aspectualidad: indica inicio.
      \ex\label{ex:11-valores-d} Un libro \textit{fue leído} por María durante el
          verano. $\rightarrow$ PV de diátesis: indica voz pasiva.
    \end{xlist}
\end{exe}

En función de los distintos valores que pueden manifestar, presentamos cuatro
grupos principales de perífrasis verbales: de diátesis, aspectuales, temporales y
modales.\footnote{La clasificación de las PPVV que presentamos, basada en
\textcite{giammatteo2009marcovecchio,giammatteo2010marcovecchio}, no pretende
ser exhaustiva, sino dar cuenta de las PPVV más representativas de los distintos
valores semánticos manifestados por estas unidades.}

\subsection{Perífrasis verbal de diátesis: la voz pasiva}\label{sec:11-pasiva}

Como se vio en el \capref{7}, los verbos determinan la estructura argumental del
predicado, es decir, seleccionan los argumentos y les asignan los papeles
temáticos que les corresponden. Sin embargo, la forma en que se manifiestan estos
argumentos puede variar. En gramática recibe el nombre de \termino{diátesis} la
categoría que

\begin{cita}
establece (i) qué argumentos verbales se seleccionan para una predicación, y
(ii) qué funciones sintácticas le corresponden a cada argumento. Así,
dependiendo de las posibles formas de disponer los argumentos del verbo en una
oración, se considera que existen \textit{distintas formas de diátesis}; puede
decirse que cada diátesis constituye un patrón de correspondencia específico
entre los argumentos de un verbo y un conjunto de funciones sintácticas
\parencite[222]{munoz2018}.
\end{cita}

Para manifestar las oposiciones de diátesis, las lenguas acuden a recursos
morfológicos (flexión y derivación) o sintácticos (fundamentalmente,
perífrasis). Estas diferencias formales que expresa la diátesis, como, por
ejemplo, \forma{despertó} frente a \forma{fue despertado}, se denominan
\termino{voces} \parencite{rae2025}. Así, tanto \exref{ex:11-vidrios-a} como
\exref{ex:11-vidrios-b} expresan los mismos argumentos, pero mediante
configuraciones oracionales diferentes:

\begin{exe}
  \ex
    \begin{xlist}
      \ex\label{ex:11-vidrios-a} Juan rompió los vidrios. $\rightarrow$ Voz
          activa.
      \ex\label{ex:11-vidrios-b} Los vidrios fueron rotos (por Juan).
          $\rightarrow$ Voz pasiva.
    \end{xlist}
\end{exe}

El cambio que manifiestan \exref{ex:11-vidrios-a} y \exref{ex:11-vidrios-b} se
conoce como \termino{alternancia activa-pasiva}. Mientras la voz activa se
expresa por la forma simple del verbo \exref{ex:11-vidrios-a}, la pasiva lo hace
mediante la perífrasis formada por el auxiliar \forma{ser} y el participio del
verbo principal \exref{ex:11-vidrios-b}. El participio de la voz pasiva
concuerda en género y número con el sujeto.

\begin{exe}
  \ex
    \begin{xlist}
      \ex Un investigador de la universidad \textit{fue premiado} por su
          trayectoria académica.
      \ex Varias investigadoras de la universidad \textit{fueron premiadas} por
          sus trayectorias académicas.
    \end{xlist}
\end{exe}

La voz activa es la forma no marcada característica de los verbos transitivos.
Solo los verbos transitivos, que poseen dos argumentos ---un agente, que funciona
como sujeto, y un tema o paciente, en la función de COD---, permiten la
alternancia activa-pasiva. La forma pasiva reestructura la configuración
oracional promoviendo el tema o paciente a la función oracional más prominente,
la de sujeto, al mismo tiempo que desplaza el agente a la función marginal y
omisible de complemento agente.

Además de la pasiva perifrástica, en español también tenemos una pasiva
cuasirrefleja, formada por el clítico \forma{se} seguido de la forma activa del
verbo. La pasiva cuasirrefleja también promueve el tema o paciente a la función
de sujeto, pero no admite la manifestación sintáctica del agente.

\begin{exe}
  \ex \textit{Se rompieron} los vidrios (*por Juan).\footnote{En el \capref{10},
      §\ref{sec:10-pasivas}, se presentaron comparativamente las diferencias
      entre la pasiva perifrástica y la cuasirrefleja pasiva.}\pendiente{El
      original remite al «capítulo 9»}
\end{exe}

Asimismo, también existe otra perífrasis de diátesis, que «expresa un estado
resultante de una acción pasiva o media\ldots» \parencite[3429]{yllera1999}. Se
forma con el verbo \forma{estar} + participio concordado con el sujeto y puede
incluir un complemento agente\footnote{La adscripción de \forma{estar} como
auxiliar de perífrasis no es unívoca, dado que en este verbo conviven el
comportamiento de auxiliar, como se muestra en los ejemplos
\exref{ex:11-escombros} a \exref{ex:11-plantel}, y el de verbo copulativo, como
en \forma{Está angustiado por todo lo que le pasó}, donde el participio puede
ser conmutado por un adjetivo: \forma{Está triste por todo lo que le
pasó}.} referido al participio, que tiene interpretación pasiva. Según sugieren
\textcite[105]{garciafernandez2019}, a diferencia de las pasivas con
\forma{ser}, estas perífrasis de resultado que se construyen con \forma{estar}
«se pueden formar regularmente sobre verbos intransitivos, especialmente sobre
aquellos que tienen un incremento clítico e interpretación
media».\pendiente{El original cita «Fernández García y Krivochen»}

\begin{exe}
  \ex\label{ex:11-escombros} Los escombros \textit{estaban esparcidos} por toda
      la zona.
  \ex La reunión \textit{estuvo presidida} por el actual jefe de gobierno.
  \ex\label{ex:11-plantel} Ud. ya \textit{está incorporado} en el nuevo plantel
      docente.
\end{exe}

\subsection{Perífrasis verbales aspectuales}\label{sec:11-aspectuales}

Las perífrasis aspectuales tienen que ver con «la configuración temporal interna
(o contorno) de las situaciones (acciones, sucesos, estados, etc.)»
\parencite[346]{lyons1997}. Se reconocen dos subgrupos: las que se refieren al
desarrollo del evento y dan cuenta de las fases de su desenvolvimiento, y las que
se vinculan con su cuantificación y consideran la pluralización de los sucesos o
de las ocasiones en que ocurren, o bien su intensificación. Veamos más
detenidamente cada subtipo.

\paragraph{Perífrasis aspectuales de fase.} Manifiestan las etapas o «fases» de
desarrollo interno del evento e incluyen:

\begin{itemize}
  \item Las \termino{perífrasis incoativas}, que señalan el inicio del evento
        ---\forma{empezar / comenzar / ponerse / echar / romper a} +
        infinitivo--- o su inminencia ---\forma{estar a punto de / estar por} +
        infinitivo---.
\end{itemize}

\begin{exe}
  \ex
    \begin{xlist}
      \ex\label{ex:11-inc-a} María \textit{empezó / comenzó / se puso a
          estudiar}.
      \ex\label{ex:11-inc-b} María \textit{echó / rompió a llorar / a correr /
          *a estudiar / *a pensar}.
    \end{xlist}
  \ex\label{ex:11-parir} María \textit{estaba a punto de / por parir} y llamaba
      la atención su serenidad.
\end{exe}

Mientras que las perífrasis de \exref{ex:11-inc-a} solo señalan el comienzo del
evento de \forma{estudiar}, las de \exref{ex:11-inc-b} añaden que este se produce
de manera abrupta y repentina; por lo tanto, estas perífrasis tienen más
restricciones para combinarse con distintos verbos. Las perífrasis verbales de
\exref{ex:11-parir} describen una fase previa al inicio efectivo del evento.

\begin{itemize}
  \item Las \termino{perífrasis terminativas} indican la conclusión efectiva del
        evento ---\forma{terminar / concluir / acabar de} + infinitivo--- o su
        interrupción ---\forma{cesar / dejar de} + infinitivo---.
\end{itemize}

\begin{exe}
  \ex
    \begin{xlist}
      \ex\label{ex:11-term-a} María \textit{terminó / concluyó / acabó de
          estudiar} Derecho.
      \ex\label{ex:11-term-b} María \textit{dejó de estudiar} Derecho, pero
          retomará el año próximo.
    \end{xlist}
\end{exe}

Mientras las perífrasis de \exref{ex:11-term-a} indican la terminación del
evento, en \exref{ex:11-term-b} se señala que se suspendió la realización, sin
que ello implique que el hecho haya concluido.

\begin{itemize}
  \item Las \termino{perífrasis continuativas} marcan la continuación de un
        evento iniciado con anterioridad ---\forma{seguir / continuar} +
        gerundio---.
\end{itemize}

\begin{exe}
  \ex María \textit{siguió / continuó estudiando}.
\end{exe}

\paragraph{Perífrasis aspectuales de cuantificación.} Manifiestan un valor
incremental referido al desarrollo del evento o a su repetición, y forman dos
subgrupos:

\begin{itemize}
  \item Las \termino{perífrasis progresivas} marcan el progreso efectivo
        ---\forma{estar / andar} + gerundio--- o gradual ---\forma{ir} +
        gerundio--- del desarrollo de un evento ya comenzado.
\end{itemize}

\begin{exe}
  \ex
    \begin{xlist}
      \ex\label{ex:11-prog-a} María \textit{estaba recorriendo} el territorio. (=
          \sig{Recorría efectivamente el territorio})
      \ex\label{ex:11-prog-b} María \textit{iba recorriendo} el territorio. (=
          \sig{Recorría gradualmente el territorio})
    \end{xlist}
\end{exe}

Mientras en \exref{ex:11-prog-a} se describe un evento en su desarrollo
continuado, en \exref{ex:11-prog-b} se indica que el evento se realiza de manera
gradual.

\begin{itemize}
  \item Las \termino{perífrasis reiterativas} indican la repetición
        ---\forma{volver a} + infinitivo--- o el desarrollo interrumpido de uno o
        varios eventos ---\forma{andar} + gerundio---.
\end{itemize}

\begin{exe}
  \ex
    \begin{xlist}
      \ex\label{ex:11-reit-a} María \textit{volvió a recorrer} el territorio. (=
          \sig{Recorrió nuevamente el territorio})
      \ex\label{ex:11-reit-b} María \textit{anda recorriendo} el territorio. (=
          \sig{Realiza una recorrida que se extiende durante un cierto tiempo o
          recorre frecuentemente el territorio})
    \end{xlist}
\end{exe}

Mientras que \exref{ex:11-reit-a} indica la repetición de un evento que ya ha
concluido (\forma{volvió a leer el libro} $\rightarrow$ \sig{lo había leído y lo
leyó nuevamente}), \exref{ex:11-reit-b} «describe situaciones que se desarrollan
con interrupciones o de modo intermitente» \parencite[§\,28.14a]{rae2025},
aunque pueden involucrar un solo proceso ---\forma{anda terminando la tesis}--- o
varios ---\forma{andan buscando al sospechoso por todos lados}---.\pendiente{El
original remite a (47a) y (47b)}

\subsection{Perífrasis verbales temporales}\label{sec:11-temporales}

Dentro de las perífrasis verbales temporales incluimos la que forma los tiempos
compuestos y la que indica futuridad.

\paragraph{Perífrasis de tiempos compuestos.} Se forman con \forma{haber} +
participio no concordado. Se consideran témporo-aspectuales porque señalan un
valor temporal de anterioridad respecto de otro evento y un valor aspectual
perfectivo, que indica \sig{completamiento del evento}. Consideremos los
siguientes ejemplos:

\begin{exe}
  \ex
    \begin{xlist}
      \ex\label{ex:11-llegue-a} Cuando yo \textit{llegué}, Juan \textit{se
          retiró}.
      \ex\label{ex:11-llegue-b} Cuando yo \textit{llegué}, Juan ya \textit{se
          había retirado}.
    \end{xlist}
\end{exe}

Mientras que en \exref{ex:11-llegue-a} las formas verbales están marcando un
valor temporal de simultaneidad entre ambos eventos ---\forma{llegué} y \forma{se
retiró}---, en \exref{ex:11-llegue-b}, en cambio, lo que se está señalando es que
el retiro de Juan fue anterior a mi llegada. Además de marcar anterioridad con
respecto a otro tiempo, las perífrasis de tiempos compuestos señalan
\termino{aspecto perfecto}. \textit{Perfecto} viene del latín
\textit{perfectum}, que significa \sig{acabado o terminado}, lo que quiere decir
que, aspectualmente, los tiempos compuestos refieren un evento ya concluido y del
cual se muestran los resultados \parencite{garciafernandez1998}.

En las perífrasis que conforman los tiempos compuestos, la morfología del verbo
\forma{haber} reproduce el paradigma flexivo de los tiempos simples y constituye
formas verbales que marcan anterioridad a otro tiempo de referencia.

\begin{exe}
  \ex
    \begin{xlist}
      \ex\label{ex:11-comp-a} Ya les \textit{han dicho} la verdad. (anterioridad
          al presente $\rightarrow$ \sig{antes de ahora})
      \ex\label{ex:11-comp-b} Cuando \textit{sonó} la alarma, los ladrones ya se
          \textit{habían escapado}. (anterioridad a un evento del pasado)
      \ex\label{ex:11-comp-c} Cuando María \textit{llegue},\footnote{Después de
          \forma{cuando}, el presente de subjuntivo \forma{llegue} tiene valor
          futuro.} la fiesta ya \textit{habrá terminado}. (anterioridad a un
          evento futuro)
    \end{xlist}
\end{exe}

El pretérito perfecto compuesto \forma{han dicho} señala un evento anterior al
presente \exref{ex:11-comp-a};\footnote{En el español de la Argentina, el
pretérito perfecto también admite uso deíctico equivalente al perfecto simple:
\forma{Ayer he ido al cine} (= \forma{Ayer fui al cine}).} el pretérito
pluscuamperfecto \forma{se habían escapado} \exref{ex:11-comp-b} señala un
tiempo anterior a otro del pasado, \forma{cuando sonó la
alarma},\footnote{En el español actual, el pretérito anterior (\forma{hubo
entrado}) está prácticamente desaparecido del habla coloquial y muy restringido
en la lengua escrita \parencite{rojo1999}.} y el futuro perfecto \forma{habrá
terminado} \exref{ex:11-comp-c} indica anterioridad a un hecho futuro, \forma{la
llegada de María}.

\paragraph{Perífrasis de futuro.} Se forman con \forma{ir a} + infinitivo.
Originariamente indicaban posterioridad inmediata, es decir, un futuro próximo,
pero el sentido de inmediatez se fue perdiendo y la perífrasis pasó a significar
futuridad en general, tanto próxima \exref{ex:11-fut-a} como remota
\exref{ex:11-fut-b}.

\begin{exe}
  \ex
    \begin{xlist}
      \ex\label{ex:11-fut-a} En cinco minutos \textit{vamos a salir} para el
          campo.
      \ex\label{ex:11-fut-b} El año próximo \textit{vamos a veranear} en Brasil.
    \end{xlist}
\end{exe}

\subsection{Perífrasis verbales de modalidad}\label{sec:11-modales}

La \termino{modalidad} es la categoría que señala la actitud del hablante con
respecto a su enunciado.\footnote{Este tema se desarrolla en el
\capref{14}.\pendiente{El original remite al «capítulo 12»}}
\textcite{palmer1986} distingue entre \termino{modalidad deóntica} y
\termino{modalidad epistémica} como dos valores presentes en todas las lenguas.
El primero está vinculado con distintos matices dentro del ámbito de la
obligación y el segundo, con valores ligados a la posibilidad. Ambas modalidades
pueden ser expresadas por perífrasis verbales formadas por los auxiliares
\forma{poder, deber, tener que} + infinitivo, mientras que \forma{haber que} +
infinitivo solo admite valor deóntico.

\begin{itemize}
  \item Las \termino{perífrasis deónticas} están orientadas hacia el sujeto de la
        oración. En ellas, \forma{poder} + infinitivo adopta valor de
        \sig{capacidad} \exref{ex:11-correr} o de \sig{permiso}
        \exref{ex:11-volver}; y \forma{deber} o \forma{tener / haber que} +
        infinitivo manifiestan \sig{obligación} \exref{ex:11-bertha} o
        \sig{necesidad} \exref{ex:11-agua}.
\end{itemize}

\begin{exe}
  \ex\label{ex:11-correr} Juan \textit{puede correr} los 100 metros en 9
      segundos.
  \ex\label{ex:11-volver} María \textit{puede volver} a casa a la una.
  \ex\label{ex:11-bertha} Bertha \textit{debe entregar} la solicitud antes de las
      11.
  \ex\label{ex:11-agua} \textit{Tenemos que beber} más agua para estar sanos.
\end{exe}

\begin{itemize}
  \item Las \termino{perífrasis epistémicas} utilizan los mismos auxiliares que
        las deónticas (excepto \forma{haber que}), pero con significados
        relacionados con la \sig{posibilidad} o la \sig{probabilidad} de la
        ocurrencia del evento.
\end{itemize}

\begin{exe}
  \ex María \textit{puede llegar} a las 8~h. $\rightarrow$ Es posible que María
      llegue a las 8~h.
  \ex María \textit{debe (de) llegar} cerca de las 8~h. $\rightarrow$ Supongo que
      María llegará cerca de las 8~h.
\end{exe}

Las perífrasis modales presentan algunas particularidades. Según dijimos en
§\ref{sec:11-auxiliares}, dado que en las PPVV auxiliar y auxiliado forman una
unidad compleja, el auxiliar no admite ser negado en forma independiente del
verbo principal:

\begin{exe}
  \ex[*]{Mariana no va a no ir a la facultad.}
\end{exe}

Sin embargo, las perífrasis de modalidad (salvo \forma{tener que}), además de la
negación del auxiliar \exref{ex:11-negmod-a}, también permiten la de ambos
verbos \exref{ex:11-negmod-b}, con lectura predominantemente deóntica; en
cambio, cuando se niega solo el auxiliado, la interpretación resulta epistémica
\exref{ex:11-negmod-c}. Estas posibilidades muestran que la unión entre auxiliar
y auxiliado es más laxa en estas perífrasis verbales.

\begin{exe}
  \ex
    \begin{xlist}
      \ex\label{ex:11-negmod-a} \{\textit{No puede / no debe}\} \textit{venir}
          temprano. (No está autorizado).
      \ex\label{ex:11-negmod-b} \{\textit{No puede / no debe}\} no \textit{venir}
          temprano. (Está obligado a venir temprano).
      \ex\label{ex:11-negmod-c} \{\textit{Puede / debe}\} no \textit{venir}
          temprano. (Tiene la posibilidad de no venir / Supongo que no vendrá
          temprano).
    \end{xlist}
\end{exe}

Como vimos en \exref{ex:11-negmod-c}, las perífrasis epistémicas admiten la
negación del auxiliado \exref{ex:11-camino1}, pero cuando la negación precede a
\forma{poder}, este se interpreta como \sig{es imposible}
\exref{ex:11-camino2}:

\begin{exe}
  \ex\label{ex:11-camino1} Pudo no encontrar el camino. $\rightarrow$ Es posible
      que no haya encontrado el camino / Sospecho que no ha encontrado el camino.
  \ex\label{ex:11-camino2} No pudo perderse en ese camino. $\rightarrow$ Es
      imposible / No creo que se haya perdido en ese camino.
\end{exe}

\begin{recuadro}[Para ampliar: las perífrasis de habitualidad en la intersección entre modalidad y aspectualidad]
La perífrasis verbal de habitualidad se forma con \forma{soler, acostumbrar} y,
en el Río de la Plata, también con \forma{saber} + infinitivo. En el caso de
\forma{soler}, presenta la particularidad de que nunca actúa como verbo pleno,
sino que únicamente funciona como auxiliar seguido de infinitivo.

\begin{exe}
  \ex [L]as águilas harpías son aves muy escurridizas, que \textit{suelen volar}
      bajo entre las copas de los árboles y son difíciles de avistar. (\textit{El
      Ciudadano}, 17/02/2025)
  \ex Es por esto que te digo que no \textit{acostumbro a plantar} dalias en los
      macizos del jardín. (\textit{Revista Mi Jardín}, 2020)
  \ex Ya de vuelta de esas tierras, cuando me encontraba en la proximidad de los
      puertos, me \textit{sabía venir} la tentación de interrogar a los marinos
      que volvían de viaje. (J.~J.~Saer, \textit{El entenado})
\end{exe}

Si bien las perífrasis de habitualidad generalmente se han ubicado entre las
aspectuales, según la \textit{NGLE} \parencite[§\,28.9b]{rae2025}, «\textit{soler}
+ infinitivo ocupa un lugar intermedio entre las perífrasis modales y las
aspectuales». Anteriormente, también \textcite[40]{comrie1985} había sostenido
que «la expresión gramatical de la habitualidad siempre estará integrada en el
sistema aspectual o modal de una lengua». Pero ¿a qué se debe esta
particularidad? El valor aspectual de estas perífrasis se relaciona con que hacen
referencia a la reiteración de un evento a lo largo de un período en el que el
suceso típicamente ocurre \parencite{olbertz1998}. En cuanto a su valor modal,
deriva de que dicho evento debe cumplirse en la mayoría, pero no necesariamente
en todas las ocasiones, lo que implica que su ocurrencia tiene una cierta
regularidad en el período involucrado, aunque la reiteración puede no ser
siempre efectiva.

\begin{exe}
  \ex Juan \textit{suele llegar} (frecuentemente llega) a estas horas, pero hoy
      todavía no ha llegado (aún no está aquí).
  \ex Don Gómez \textit{sabe venir} por acá los sábados, pero este fin de semana
      no lo hemos visto.
\end{exe}

En \textcite{giammatteo2010marcovecchio} y \textcite{giammatteo2011} se plantea
que este carácter virtual de la reiteración de un evento vincula habitualidad y
posibilidad. De este modo, \forma{poder / deber (de) / tener que}, por un lado, y
\forma{soler / saber / acostumbrar}, por el otro, manifiestan grados dentro de
la modalidad epistémica. Así, \forma{soler} presupone la realización efectiva
del evento, aunque le permite al hablante evitar un compromiso absoluto con el
valor de verdad de su aseveración para todas las ocasiones; \forma{poder}, en
cambio, indica un grado aún menor de compromiso del hablante con sus dichos, ya
que solo hace referencia a una realización eventual o posible del evento, que no
confiere ninguna certeza a su realización.

\begin{exe}
  \ex Suele ir al cine los domingos. (Generalmente lo hace, aunque algún domingo
      puede no ir).
  \ex Puede ir al cine este domingo. (Existe una posibilidad de que el hecho
      suceda, aunque también podría no ocurrir).
\end{exe}
\end{recuadro}

A modo de recapitulación, en la tabla~\ref{tab:11-perifrasis} ofrecemos una
sistematización de las distintas clases de perífrasis reconocidas en el
capítulo.

{\footnotesize\setlength{\tabcolsep}{3pt}
\begin{longtable}{@{}>{\raggedright\arraybackslash}p{24mm}
    >{\raggedright\arraybackslash}p{30mm} >{\raggedright\arraybackslash}p{30mm}
    >{\raggedright\arraybackslash}p{36mm}@{}}
  \caption{Sistematización de las perífrasis verbales}\label{tab:11-perifrasis}\\
  \toprule
  Clase & Subclase & Descripción estructural & Ejemplos \\
  \midrule
  \endfirsthead
  \toprule
  Clase & Subclase & Descripción estructural & Ejemplos \\
  \midrule
  \endhead
  \multicolumn{4}{@{}p{126mm}}{\textbf{De diátesis}: modifican la relación entre
    los argumentos del verbo y las funciones sintácticas que se manifiestan en la
    oración.} \\
  \addlinespace
  Pasiva & --- & ⟨\textit{ser} + participio concordado⟩ & \textit{La casa fue
    comprada por un nuevo vecino.} \\
  Resultativa & --- & ⟨\textit{estar} + participio concordado⟩ & \textit{La casa
    está ocupada solo por muebles viejos.} \\
  \midrule
  \multicolumn{4}{@{}p{126mm}}{\textbf{Aspectuales}: aportan información sobre el
    desarrollo interno del evento.} \\
  \addlinespace
  De fase (I): etapas de desarrollo del evento & incoativas: marcan el inicio del
    evento & ⟨\textit{empezar / comenzar / echar / romper / ponerse a} + inf.⟩;
    ⟨\textit{estar a punto de / por} + inf.⟩ (inminencial) & \textit{Ana empezó a
    dibujar el plano. El niño rompió a llorar desconsolado. Los jueces están por
    dictar sentencia.} \\
  & terminativas: marcan el fin o la interrupción del evento &
    ⟨\textit{terminar / acabar / concluir de} + inf.⟩; ⟨\textit{llegar /
    alcanzar a} + inf.⟩; ⟨\textit{cesar / dejar de} + inf.⟩ (cesativa) &
    \textit{Pedro terminó de pintar la pared. Raúl dejó de leer el libro.} \\
  & continuativas: marcan la continuación de un evento en desarrollo &
    ⟨\textit{continuar / hallarse / llevar / seguir / venir} +
    ger.⟩\textsuperscript{a} & \textit{Los alumnos se hallan buscando
    información. El gato continúa persiguiendo al ratón.} \\
  De cuantificación (II): incremento en el desarrollo del evento & reiterativas:
    marcan la repetición del evento & ⟨\textit{volver a} + inf.⟩;
    ⟨\textit{andar} + ger.⟩ (reiterativa-frecuentativa) & \textit{Los libros
    volvieron a ser disfrutados por todos. Mariana anda bostezando por todas
    partes.} \\
  & progresivas: marcan el desarrollo ininterrumpido de uno o varios eventos &
    ⟨\textit{estar} + ger.⟩; ⟨\textit{ir} + ger.⟩ (progresiva gradual) &
    \textit{Margarita está traduciendo un libro al inglés. Juanita va aprendiendo
    a contar.} \\
  \midrule
  \multicolumn{4}{@{}p{126mm}}{\textbf{Temporales}: aportan información sobre el
    tiempo del evento.} \\
  \addlinespace
  De tiempos compuestos & expresan anterioridad a otro evento e indican aspecto
    perfecto & ⟨\textit{haber} + part.⟩ & \textit{Todos hemos aprobado
    Gramática.} \\
  De futuro & expresan posterioridad al tiempo presente & ⟨\textit{ir a} + inf.⟩
    & \textit{El próximo mes voy a disfrutar de mis vacaciones.} \\
  \midrule
  \multicolumn{4}{@{}p{126mm}}{\textbf{Modales}: manifiestan información sobre la
    actitud del hablante ante el contenido de su enunciado.} \\
  \addlinespace
  Deónticas & expresan obligación, necesidad o capacidad del sujeto de la oración
    & ⟨\textit{deber (de) / poder / tener que / haber que} + inf.⟩ &
    \textit{Deben estudiar este tema para el examen.} \\
  Epistémicas & expresan el grado de posibilidad que el hablante adjudica a la
    ocurrencia de un evento\textsuperscript{b} & (los mismos auxiliares, salvo
    \textit{haber que}) & \textit{Deben (de) haber estudiado mucho porque
    aprobaron.} \\
  De habitualidad (intersección entre aspecto y modalidad) & señalan que en un
    período extendido de tiempo algún suceso típicamente ocurre &
    ⟨\textit{soler} + inf.⟩; ⟨\textit{acostumbrar a} + inf.⟩; ⟨\textit{saber}
    + inf.⟩ & \textit{En el pueblo suelen festejar el 25 de Mayo con locro y
    empanadas. Sabe estar siempre en su casa tomando mate.} \\
  \bottomrule
\end{longtable}}

{\scriptsize\noindent\textsuperscript{a} Dada la índole introductoria de este
libro, omitimos detallar los diferentes matices aspectuales que manifiestan
estas PPVV. \textsuperscript{b} Los distintos grados de posibilidad o
probabilidad de la ocurrencia de un evento dependen del grado de certeza con que
el hablante lo comunica (información de primera mano, inferencia directa,
indirecta, etc.). Estos temas se amplían en el \capref{14}, dedicado a la
modalidad.\par}

\begin{ejercicio}\label{ej:11-pleno}
En las siguientes oraciones, diferencie el significado del verbo en su uso pleno
del que aporta como auxiliar de perífrasis, según se muestra en el ejemplo
resuelto~1.

\begin{exe}
  \ex
    \begin{xlist}
      \ex Juana va todos los días a la facultad a estudiar Derecho.
      \ex Juana va a estudiar Derecho el año próximo.
    \end{xlist}
  \ex
    \begin{xlist}
      \ex El iracundo rockero rompió la guitarra en el escenario.
      \ex El iracundo rockero rompió a llorar desconsoladamente.
    \end{xlist}
  \ex
    \begin{xlist}
      \ex Al final del período cretácico, ya no había dinosaurios.
      \ex Antes de su muerte habían ocurrido algunos hechos sospechosos.
    \end{xlist}
  \ex
    \begin{xlist}
      \ex Mi abuelo era pelado.
      \ex Mi abuelo era conocido por todos en el pueblo.
    \end{xlist}
  \ex
    \begin{xlist}
      \ex El jugador está en el área de penales.
      \ex El jugador está por salir al campo de juego.
    \end{xlist}
\end{exe}

\begin{resuelto}
1a. \sig{dirigirse a}. 1b. Futuridad.
\end{resuelto}
\end{ejercicio}

\begin{ejercicio}\label{ej:11-significados}
Determine si las siguientes oraciones contienen perífrasis verbales y, en caso de
que lo hagan, señale los significados que manifiestan, según se muestra en los
ejemplos resueltos.

\begin{exe}
  \ex María necesita aprender gramática para aprobar su examen.
  \ex María tiene que aprender gramática para aprobar su examen.
  \ex María está aprendiendo gramática para aprobar su examen.
  \ex María aprende cantando.
  \ex María va aprendiendo gramática.
  \ex María continúa aprendiendo gramática.
  \ex María es dominada por su pasión por la gramática.
  \ex María está sentada en primera fila.
  \ex María suele estudiar gramática.
  \ex María empezó a estudiar gramática.
  \ex María no volvió a rendir el examen de gramática porque lo aprobó.
  \ex María no vuelve a estudiar para el examen de gramática, sino para
      descansar.
\end{exe}

\begin{resuelto}[Ejemplos resueltos]
1. No es PV. \textit{Aclaración}: en \forma{necesita aprender} no es posible
aplicar la prueba de conmutación estructural, según se explicó en
§\ref{sec:11-auxiliares}, porque el sujeto de ambos verbos es el mismo,
\forma{María}; pero la aplicación de la prueba es posible si cambiamos la persona
del sujeto del primer verbo: \forma{Yo necesito que María aprenda gramática}. 2.
PV de modalidad deóntica.
\end{resuelto}
\end{ejercicio}

\begin{recuadro}[Para ampliar: equivalencias semánticas entre perífrasis verbales y adverbios]
Algunos de los significados que expresan las perífrasis verbales pueden también
manifestarse a través de adverbios, como puede observarse en los siguientes pares
de oraciones, en los que las versiones de (a) expresan respectivamente futuridad,
continuidad y posibilidad mediante adverbios, y las de (b) lo hacen a través de
perífrasis verbales.

\begin{exe}
  \ex Futuridad:
    \begin{xlist}
      \ex \textit{Mañana} voy a Buenos Aires.
      \ex \textit{Voy a ir} a Buenos Aires.
    \end{xlist}
  \ex Continuidad:
    \begin{xlist}
      \ex \textit{Todavía} Juan está en Buenos Aires.
      \ex Juan \textit{sigue estando} en Buenos Aires.
    \end{xlist}
  \ex Posibilidad:
    \begin{xlist}
      \ex \textit{Posiblemente / Tal vez} Juan esté en Buenos Aires.
      \ex Juan \textit{puede estar} en Buenos Aires.
    \end{xlist}
\end{exe}
\end{recuadro}

\begin{ejercicio}\label{ej:11-reformular}
Reformule las siguientes oraciones con una perífrasis verbal de manera que
conserven el mismo significado, según se muestra en los ejemplos resueltos.
Tenga en cuenta los segmentos subrayados.\pendiente{Los subrayados no se
conservan en el original}

\begin{exe}
  \ex Se compraron algunas provisiones.
  \ex Tal vez te guste la nueva novela.
  \ex Mañana le doy a mi compañero la contraseña.
  \ex Le ordeno que se detenga.
  \ex Aquí nieva con frecuencia en invierno.
  \ex Posiblemente, Ricardo se ha dormido.
  \ex Enseguida inicio el armado del equipo.
  \ex El escritor recién terminó su última novela.
  \ex El escritor vivió durante muchos años en San Juan, pero ahora vive en
      Santiago de Chile.
  \ex Mariela nuevamente será mamá.
  \ex El profesor les pidió a sus alumnos que mantengan el ritmo de estudio.
  \ex Por favor, ya no peleen más.
  \ex Todos sabemos que Miguel tuvo la capacidad de vencer a su rival.
  \ex Finalmente, se dieron las condiciones para que Rosa duerma.
\end{exe}

\begin{resuelto}[Ejemplos resueltos]
1. Algunas provisiones fueron compradas. 2. La nueva novela te puede gustar.
\end{resuelto}
\end{ejercicio}

\section{Agrupaciones o cadenas perifrásticas}\label{sec:11-cadenas}

Las perífrasis verbales también pueden combinarse entre sí para formar unidades
verbales complejas que reúnen más de un auxiliar. En estos casos, solo el primer
auxiliar es una forma conjugada; los demás mantienen la forma no finita que
corresponde a la clase de perífrasis que integran ---\forma{pueden empezar a
comer, había sido designado, iban a estar conversando}---.

\begin{ejercicio}\label{ej:11-cadenas}
Observe los ejemplos siguientes. Luego realice una descripción estructural de
las perífrasis que integran cada núcleo verbal e indique el significado que
transmite cada uno de sus constituyentes, según se muestra en el ejemplo
resuelto~1.

\begin{exe}
  \ex María \textit{debe seguir estudiando} en la universidad.
  \ex María \textit{pudo terminar de estudiar} Derecho el año pasado.
  \ex María \textit{va a poder empezar a estudiar} en la universidad el próximo
      año.
  \ex Juan y María \textit{deberían haber empezado a estudiar} para los exámenes
      hace rato.
  \ex Juan \textit{tiene que poder dejar de comer} grasas saturadas.
\end{exe}

\begin{resuelto}
1. \forma{Debe} (obligatoriedad) + \forma{seguir} (continuidad) + gerundio
\forma{estudiando}.
\end{resuelto}
\end{ejercicio}

\subsection{Ámbito oracional de las perífrasis verbales}\label{sec:11-ambito}

Ahora bien, ¿qué sucede con el orden entre perífrasis verbales cuando se
combinan entre sí formando cadenas? Observemos los siguientes casos:

\begin{exe}
  \ex\label{ex:11-chomsky}
    \begin{xlist}
      \ex[*]{\label{ex:11-chomsky-a} El libro de Chomsky \textit{empezó a haber
          sido leído} por los estudiantes.}
      \ex\label{ex:11-chomsky-b} El libro de Chomsky \textit{pudo ser leído} por
          los estudiantes.
      \ex[*]{\label{ex:11-chomsky-c} El libro \textit{fue debido de ser leído}
          por los estudiantes.}
      \ex\label{ex:11-chomsky-d} Los estudiantes \textit{habían vuelto a leer} el
          libro de Chomsky antes del examen.
      \ex[*]{\label{ex:11-chomsky-e} Los estudiantes \textit{continuaron habiendo
          comprado} libros.}
    \end{xlist}
\end{exe}

Según vemos en los ejemplos de arriba, algunos órdenes entre PPVV son posibles,
mientras que otros producen resultados agramaticales. Pero ¿qué determina la
gramaticalidad o agramaticalidad de las oraciones anteriores? Aunque podría
suponerse que la combinación de perífrasis verbales se da en un orden aleatorio,
su ordenamiento secuencial tiene que ver con los niveles o capas de la estructura
oracional que hemos visto en el \capref{6}.

\begin{destacado}
\textbf{Para recordar.} Según la propuesta de \textcite{rizzi1995}, la oración se
organiza en tres capas o niveles (\textit{shells}): la \termino{capa léxica} o
del SV, donde se produce la vinculación del V con sus argumentos y
modificadores; la \termino{capa oracional} o del SF, que pone en relación el
sujeto con el predicado y en la que la oración se temporaliza, y la
\termino{capa informativa} o del SComp, que incluye la perspectiva del hablante y
establece la articulación con el nivel del discurso.
\end{destacado}

Dentro del dominio del SV se ubica la perífrasis de voz pasiva,\pendiente{El
original dice «voz activa»} que es la más interna y afecta a la estructura
argumental del verbo.

\begin{exe}
  \ex {}[\etq{PAS} Fue reconocido] por su obra literaria.
\end{exe}

También dentro del mismo nivel del SV se sitúan, en orden correlativo, las
perífrasis aspectuales de fase, que preceden a la pasiva \exref{ex:11-asp-a};
luego, las cuantificacionales, que anteceden a las de fase y a la pasiva
\exref{ex:11-asp-b}, y, por último, las de modalidad deóntica, que preceden a
todas las anteriores, \exref{ex:11-deont-a} y \exref{ex:11-deont-b}:

\begin{exe}
  \ex
    \begin{xlist}
      \ex\label{ex:11-asp-a} {}[\etq{ASP I} Empezó a [\etq{PAS} ser reconocido]]
          por su obra literaria en años recientes.
      \ex\label{ex:11-asp-b} {}[\etq{ASP II} Está [\etq{ASP I} empezando a
          [\etq{PAS} ser reconocido]]] por su obra literaria recién ahora.
    \end{xlist}
  \ex
    \begin{xlist}
      \ex\label{ex:11-deont-a} {}[\etq{DEÓNT} Debe [\etq{ASP I} empezar a
          [\etq{PAS} ser premiado]]] por su obra literaria.
      \ex\label{ex:11-deont-b} {}[\etq{DEÓNT} Puede [\etq{ASP II} volver a
          inscribirse]] en el curso, si completa este formulario.
          (\forma{Puede} = \sig{se le permite}).
    \end{xlist}
\end{exe}

En el nivel del SF se ubican las perífrasis temporales, que se vinculan con el
tiempo gramatical de la oración. Estas perífrasis anteceden a todas las que se
sitúan en el nivel previo del SV (pasivas, aspectuales de fase y de
cuantificación, y modales deónticas).

\begin{exe}
  \ex {}[\etq{TEMP} Había [\etq{ASP I} terminado de estudiar]] ese año.
  \ex {}[\etq{TEMP} Va a [\etq{ASP II} volver a [\etq{PAS} ser premiado]]] por su
      trayectoria.
\end{exe}

Finalmente, en la capa más externa de la oración, correspondiente al SComp,
donde se asientan los elementos que conectan el contenido proposicional de la
oración con la fuerza ilocutiva del hablante, se ubican las perífrasis de
modalidad epistémica:

\begin{exe}
  \ex {}[\etq{EPIST} Deben (de) [\etq{TEMP} haber [\etq{PAS} sido capturados]]]
      por la policía. (\forma{Deben (de)} = \sig{es posible que hayan sido
      capturados}).
  \ex Los manifestantes [\etq{EPIST} pueden [\etq{ASP II} estar [\etq{PAS} siendo
      dispersados]]]. (\forma{Pueden} = \sig{es posible que estén siendo
      dispersados}).
\end{exe}

Por tanto, según acabamos de ver, en la combinación de perífrasis el orden de los
auxiliares que se observa es el siguiente:

\begin{center}
[EPIST [TEMP [DEÓNT [ASP II (cuant.) [ASP I (de fase) [PAS]]]]]]
\end{center}

En el caso de los auxiliares de habitualidad, como \forma{soler, saber} y
\forma{acostumbrar}, una de las pruebas que se aportan en
\textcite{giammatteo2009marcovecchio,giammatteo2010marcovecchio} para demostrar
que también corresponden al ámbito del SComp es que estos auxiliares no aceptan
combinarse secuencialmente con los modales epistémicos, por lo que estarían
ocupando la misma posición estructural, lo que hace incompatible su
combinación.

\begin{exe}
  \ex
    \begin{xlist}
      \ex[*]{\textit{Suele deber de llegar} temprano.}
      \ex[*]{\textit{Puede soler venir} los viernes.}
    \end{xlist}
\end{exe}

Al igual que \forma{poder} y los demás auxiliares epistémicos, \forma{soler}
ocupa siempre la posición más externa en la estructura oracional, por lo que
antecede al resto de los auxiliares que pueden conformar perífrasis:

\begin{exe}
  \ex
    \begin{xlist}
      \ex María \textit{suele tener que trabajar} los sábados.
      \ex[*]{María \textit{tiene que soler trabajar} los sábados.}
    \end{xlist}
\end{exe}

\begin{recuadro}[Para ampliar: ámbito y posibilidades de desplazamiento]
Si bien la concatenación de perífrasis responde a restricciones no aleatorias, su
reflejo en las lenguas puede variar. Mientras en inglés, que tiene orden fijo, la
secuencia de auxiliares permanece siempre invariable, en español, lengua de orden
más libre, la secuencialidad entre auxiliares no es tan rígida y permite algunos
desplazamientos. No obstante, los movimientos permitidos se circunscriben a los
límites que imponen las capas o niveles de estructuración en los que se ubican
las distintas perífrasis. Así, aquellas que se encuentran dentro del ámbito del
SV (pasivas, aspectuales de fase, cuantificacionales y deónticas) admiten más de
una posición entre ellas.

\begin{exe}
  \ex Este artículo [\etq{PAS} fue [\etq{ASP I} empezado a / terminado de
      revisar]] la semana pasada.
  \ex Este artículo [\etq{ASP II} volvió a [\etq{DEÓNT} tener que [\etq{PAS} ser
      leído]]] entre los académicos.
  \ex {}[\etq{ASP II} Está [\etq{DEÓNT} pudiendo comprar]] por Internet
      nuevamente. (\forma{Poder} = \sig{capacidad o permiso}).
\end{exe}

Sin embargo, ni la pasiva ni las aspectuales o las deónticas pueden ocurrir fuera
del límite del SV, por lo que no pueden preceder ni a las temporales,
\exref{ex:11-desp1} y \exref{ex:11-desp2}, ni a las epistémicas, excepto que
admitan cambio de significado. Así, \forma{ir a} + infinitivo pierde su valor de
futuridad cuando se pospone a una PV aspectual o de modalidad deóntica
\exref{ex:11-desp2}, y \forma{poder, deber} y \forma{tener que}, cuando van
precedidos por perífrasis aspectuales, solo admiten interpretación deóntica
\exref{ex:11-desp3}.\pendiente{El original remite aquí a (86) y (88)}

\begin{exe}
  \ex[*]{\label{ex:11-desp1} [\etq{ASP I} Empezó a [\etq{TEMP} haber
      conseguido]] otra oportunidad.}
  \ex[*]{\label{ex:11-desp2} [\etq{ASP I/II} Volvió a / Terminó de [\etq{TEMP} ir
      a escribir]] un nuevo ensayo.}
  \ex\label{ex:11-desp3} {}[\etq{ASP I} Empezó a [\etq{DEÓNT} tener que
      resolver]] solo sus problemas.
\end{exe}

Asimismo, las perífrasis témporo-aspectuales, que se alojan en el nivel de la
capa flexiva, no tienen alcance sobre el ámbito más externo, correspondiente a la
modalidad oracional, es decir que el orden de los auxiliares permite determinar
el valor deóntico o epistémico de \forma{poder} y \forma{deber}: a) pospuestos a
la perífrasis temporal con \forma{haber}, se interpretan con valor deóntico
(dentro de la capa del SV) \exref{ex:11-ha-podido}; b) antepuestos a la
perífrasis temporal con \forma{haber}, quedan fuera del ámbito del SF y su única
interpretación posible es con valor epistémico \exref{ex:11-pudo-haber}.

\begin{exe}
  \ex\label{ex:11-ha-podido} Ella [\etq{TEMP} ha [\etq{DEÓNT} podido traducir]]
      el libro. $\rightarrow$ Lectura deóntica: \sig{Ella ha tenido la capacidad
      de traducir el libro}.
  \ex\label{ex:11-pudo-haber} Ella [\etq{EPIST} pudo [\etq{TEMP} haber traducido]]
      el libro. $\rightarrow$ Lectura epistémica: \sig{Es posible que fuera ella
      quien lo tradujo}.
\end{exe}
\end{recuadro}

\begin{ejercicio}\label{ej:11-orden}
Luego de haber reflexionado sobre el orden secuencial al que responden las
cadenas de PPVV, según se muestra en el ejemplo resuelto \exref{ex:11-chomsky-a},
¿podría decir en qué consiste la gramaticalidad o agramaticalidad de las
oraciones de \exref{ex:11-chomsky-b} a \exref{ex:11-chomsky-e}?

\begin{resuelto}
\exref{ex:11-chomsky-a}: es agramatical porque la PV de inicio con
\forma{empezar} no puede anteceder a la de tiempo compuesto.
\end{resuelto}
\end{ejercicio}

\begin{ejercicio}\label{ej:11-tabla}
Teniendo en cuenta los significados que manifiestan, ubique en la tabla de abajo
los verbos auxiliares y auxiliados de las combinaciones de perífrasis de las
oraciones siguientes, según se muestra en el ejemplo resuelto.

\begin{exe}
  \ex Los recién llegados deben comenzar a aprender las costumbres del lugar.
  \ex Woody Allen puede haber tenido que dejar de ser atendido por su
      psicoanalista.
  \ex Los alumnos suelen tener que volver a leer los textos de gramática.
  \ex Los libros han tenido que ser restaurados por especialistas.
  \ex Las páginas del libro deben (de) haber sido deterioradas por el paso del
      tiempo.
  \ex Mañana vamos a comenzar a practicar para el examen.
\end{exe}

\begin{center}
\resizebox{\linewidth}{!}{%
\begin{tabular}{@{}c *{6}{c} c@{}}
  \toprule
  & \multicolumn{6}{c}{Verbos auxiliares} & \\
  \cmidrule(lr){2-7}
  Ej. & \makecell{Mod. epistémica /\\habitualidad} & Temporalidad &
    \makecell{Mod.\\deóntica} & \makecell{Aspect.\\tipo II} &
    \makecell{Aspect.\\tipo I} & \makecell{Voz\\pasiva} & \makecell{Verbo\\auxiliado} \\
  \midrule
  1 & --- & --- & \textit{deben} & --- & \textit{comenzar a} & --- &
    \textit{aprender} \\
  2 & & & & & & & \\
  3 & & & & & & & \\
  4 & & & & & & & \\
  5 & & & & & & & \\
  6 & & & & & & & \\
  \bottomrule
\end{tabular}}
\end{center}
\end{ejercicio}

\begin{ejercicio}\label{ej:11-valores}
Identifique los valores gramaticales que manifiestan los auxiliares de las
perífrasis verbales subrayadas en las siguientes oraciones, según se muestra en
el ejemplo resuelto~1.

\begin{exe}
  \ex Los profesores \subrayado{van a tener que entregar} los programas con
      varios cambios significativos.
  \ex Tranquilamente, Susana \subrayado{puede seguir siendo considerada} la mejor
      alumna del colegio.
  \ex Jorge regresó al hospital para \subrayado{poder volver a ser atendido} por
      la enfermera.
  \ex Camilo \subrayado{ha vuelto a dejar de tomar} su medicina.
\end{exe}

\begin{resuelto}
1. Valores gramaticales: temporal de futuridad (\forma{van a}) + modal deóntico
(\forma{tener que entregar}). La PV de modalidad tiene valor deóntico porque
está precedida por la temporal.
\end{resuelto}
\end{ejercicio}

\begin{ejercicio}\label{ej:11-entrar}
Observe el funcionamiento de los complejos verbales encabezados por el verbo
\forma{entrar} en los siguientes fragmentos y luego resuelva lo pedido en las
consignas, según se muestra en el ejemplo resuelto~1.

\begin{consignas}
  \item Aplique las pruebas correspondientes para determinar si se trata de
        perífrasis verbales o no.
  \item Identifique el significado que transmite \forma{entrar} + infinitivo en
        sus usos perifrásticos.
\end{consignas}

\begin{exe}
  \ex Me levanto y voy a bañarme. \textit{Entro a trabajar} a las nueve, pero me
      he levantado media hora antes de lo acostumbrado; son las siete.
  \ex (\ldots) De allí, despertada su vocación militar por la guerra,
      \textit{entró a formar} parte del Ejército del Rin, distinguiéndose en la
      batalla de Wissembourg, donde fue herido en un brazo.
  \ex (\ldots) La astucia imperceptible de Antonio \textit{entró a funcionar}
      ---trataba de borrar el reciente traspié y congraciarse---.
  \ex Treparon una reja para \textit{entrar a robar} a una peluquería de Palermo
      y los detuvieron \textit{in fraganti}.
  \ex Luego se \textit{entró a absorber} el ahorro nacional que se volcaba al
      Mercado de Capitales. Las empresas no tenían posibilidad de competir con
      los bonos del Estado.
  \ex En un primer momento nos quedamos anonadados; unos segundos después el
      resto de los pibes se \textit{entró a reír} a carcajadas, algunas minas
      también, Adriana lució su mejor cara de perra triunfante y Alicia se apuró
      en ir a consolar a Patri.
\end{exe}

\noindent{\footnotesize Las oraciones 1--3, 5 y 6 fueron extraídas de M.~Davies
(2016--), \textit{Corpus del Español: Web/Dialects}
(\url{http://www.corpusdelespanol.org/web-dial/}), y la 4, de \textit{Infobae}
(17/06/2024).}

\begin{resuelto}
1. No es PV: admite conmutación estructural (\forma{entro y trabajo / entro para
trabajar}) y \forma{entrar} actúa como verbo pleno.
\end{resuelto}
\end{ejercicio}

\section{Locuciones verbales}\label{sec:11-locuciones}

Las \termino{locuciones verbales} son unidades complejas en las que el primer
constituyente, como sucede en las perífrasis verbales, es un verbo; pero, a
diferencia de estas últimas, el segundo constituyente puede ser tanto una forma
verbal (\forma{echar a perder}) como otra clase de palabra (\forma{tomar el
pelo, dar fe, echar de menos}). Además, en las locuciones verbales los
constituyentes pueden incluir más de dos palabras (\forma{hacer la cama} =
\sig{engañar}; \forma{dejar sin aliento} = \sig{impactar}). Otra diferencia
fundamental respecto de las perífrasis es que las locuciones verbales son
conjuntos de palabras con un significado unitario que, en líneas generales, no
deriva composicionalmente de los significados que individualmente aportan sus
constituyentes, es decir, las locuciones no pueden ser interpretadas
literalmente. Podemos ejemplificar estos conceptos con los siguientes casos:

\begin{exe}
  \ex\label{ex:11-entender} María \textit{dio a entender} sus pretensiones.
  \ex Pedro \textit{mete la pata} todo el tiempo.
  \ex La profesora \textit{echa de menos} su tierra natal.
  \ex\label{ex:11-toalla1} La estudiante \textit{tiró la toalla} con la lectura
      del libro de Chomsky.
\end{exe}

Los significados de las locuciones en \exref{ex:11-entender} a
\exref{ex:11-toalla1}, que pueden glosarse como \sig{sugerir}, \sig{equivocarse},
\sig{extrañar} y \sig{rendirse}, no se desprenden de los de las palabras que las
integran, sino, más bien, de un proceso de reinterpretación por el que llegan a
adquirir un significado unitario. El grado de integración de las unidades que
componen las locuciones se corrobora con ciertas propiedades que actúan sobre el
conjunto y no sobre cada unidad por separado. Así, como muestran los siguientes
ejemplos, los componentes no pueden ser modificados \exref{ex:11-loc-a} ni
negados \exref{ex:11-loc-b} en forma independiente, aunque admiten adverbios
adjuntos \exref{ex:11-pata-a} e interpolados \exref{ex:11-pata-b}.

\begin{exe}
  \ex
    \begin{xlist}
      \ex[*]{\label{ex:11-loc-a} María \textit{dio} hoy \textit{a conocer} ayer
          sus pretensiones.}
      \ex[*]{\label{ex:11-loc-b} María no \textit{dio a} no \textit{conocer} sus
          pretensiones.}
    \end{xlist}
  \ex
    \begin{xlist}
      \ex\label{ex:11-pata-a} Pedro \textit{mete la pata} sistemáticamente.
      \ex\label{ex:11-pata-b} Pedro \textit{mete} sistemáticamente \textit{la
          pata}.
    \end{xlist}
\end{exe}

Al igual que las perífrasis verbales, las locuciones conforman unidades
complejas que pueden funcionar como núcleo del predicado.

\begin{exe}
  \ex
    \begin{xlist}
      \ex El artista [\etq{SV/Pred} [\etq{V/n} \textit{ha terminado}] con el
          espectáculo]. $\rightarrow$ PV.
      \ex El artista [\etq{SV/Pred} [\etq{V/n} \textit{dio por terminado}] el
          espectáculo]. $\rightarrow$ LV.
    \end{xlist}
\end{exe}

\subsection{Reconocimiento}\label{sec:11-reconocimiento-lv}

A diferencia de las perífrasis verbales, que, según hemos visto, aceptan
conmutación léxica (§\ref{sec:11-lexica}), pero no estructural
(§\ref{sec:11-estructural}), las locuciones son complejos verbales más
solidificados y no admiten ninguno de estos cambios, por lo que la aplicación de
estas dos pruebas puede permitirnos distinguir entre ambas estructuras.

\begin{exe}
  \ex\label{ex:11-gobernador} Perífrasis verbales:
    \begin{xlist}
      \ex El gobernador \textit{no pudo regresar} a su provincia.
      \ex\label{ex:11-gobernador-b} Conmutación léxica: El gobernador no pudo
          volver / retornar a su provincia.
      \ex\label{ex:11-gobernador-c} Conmutación estructural: *El gobernador no
          pudo que regresaba a su provincia.
    \end{xlist}
  \ex\label{ex:11-partido} Locuciones verbales:
    \begin{xlist}
      \ex María \textit{echó a perder} el partido.
      \ex\label{ex:11-partido-b} Conmutación léxica: *María echó a ganar el
          partido.
      \ex\label{ex:11-partido-c} Conmutación estructural: *María echó y perdió el
          partido.
    \end{xlist}
\end{exe}

Según vemos, en \exref{ex:11-gobernador} es posible la conmutación léxica
\exref{ex:11-gobernador-b}, pero no la conmutación estructural
\exref{ex:11-gobernador-c}, por lo que se trata de una PV; en cambio, en
\exref{ex:11-partido}, el complejo \forma{echó a perder} no acepta ni conmutación
léxica \exref{ex:11-partido-b} ni conmutación estructural
\exref{ex:11-partido-c}, por lo que constituye una locución verbal (LV).

No obstante, hay que tener en cuenta que en la interpretación del significado de
las locuciones verbales y en su reconocimiento como tales incide el contexto en
el que se encuentran. Por ejemplo, en \exref{ex:11-toalla-a}, \forma{tiró la
toalla} admite una interpretación literal, en la que \forma{tiró} es un verbo
pleno, que actúa como núcleo del predicado y va seguido de un COD y de un CCL; en
cambio, en \exref{ex:11-toalla-b}, se trata de una locución porque ambos
constituyentes se interpretan como \sig{rendirse}. El caso de
\exref{ex:11-toalla-c} es interesante porque admitiría ambas interpretaciones: el
entrenador puede haber tirado literalmente la toalla con que le seca el sudor al
boxeador (sentido pleno de \forma{tirar}), pero en el boxeo ese gesto se
interpreta mayoritariamente como \sig{rendirse}, es decir, como una locución
verbal.

\begin{exe}
  \ex
    \begin{xlist}
      \ex\label{ex:11-toalla-a} María \textit{tiró la toalla} en el cesto de la
          ropa sucia.
      \ex\label{ex:11-toalla-b} María \textit{tiró la toalla} después de
          desaprobar el examen.
      \ex\label{ex:11-toalla-c} El entrenador del boxeador \textit{tiró la toalla}
          en mitad de la pelea.
    \end{xlist}
\end{exe}

\begin{recuadro}[Para ampliar: locuciones verbales frente a verbos de apoyo]
Los denominados \termino{verbos de apoyo} o \termino{de soporte} (\forma{dar un
paseo, hacer una pregunta, formar parte}), que son estructuras semilexicalizadas
formadas por un verbo y un SN que funciona como su complemento, poseen puntos de
contacto con las locuciones verbales. Los verbos de apoyo, que muchas veces
tienen un equivalente verbal aproximado (\forma{dar un paseo} = \forma{pasear};
\forma{hacer una pregunta} = \forma{preguntar}; \forma{formar parte} =
\forma{participar}), difieren de las locuciones en que su significado suele ser
composicional y se deduce de la suma de sus partes. Además, en el caso de los
verbos de apoyo, las estructuras que forman con el SN que los acompaña están
menos lexicalizadas, por lo que los sustantivos de las construcciones admiten
variantes: \forma{dar un paseo / una caminata / unas vueltas}. Además, a partir
del sustantivo se pueden formar oraciones subordinadas\footnote{Las oraciones
subordinadas, tanto las de relativo como las completivas, se estudian en el
\capref{13}.} de relativo con el verbo de apoyo como núcleo del predicado de la
subordinada (\forma{el paseo que di, la pregunta que hice}), y el sustantivo
admite modificadores: \forma{di un paseo muy largo; hice una pregunta bastante
difícil; formé parte fundamental del directorio de la empresa}.
\end{recuadro}

\begin{ejercicio}\label{ej:11-locuciones}
Responda las siguientes consignas, según se muestra en el ejemplo resuelto~1.

\begin{consignas}
  \item Determine si las siguientes oraciones contienen PPVV o LLVV aplicando las
        pruebas estudiadas: conmutación léxica y conmutación estructural.
  \item En el caso de las locuciones, proponga una forma verbal de significado
        equivalente.
\end{consignas}

\begin{exe}
  \ex Juan le \textit{dio una mano} a Ricardo con su nuevo negocio.
  \ex Todos \textit{hemos aprendido} mucho con los últimos acontecimientos.
  \ex María \textit{echó de menos} a sus padres durante el viaje a Mendoza.
  \ex Carlos \textit{hizo añicos} su auto deportivo.
  \ex El mes próximo \textit{vamos a disfrutar} de una nueva puesta teatral de
      \textit{Hamlet}.
  \ex El político \textit{dio a conocer} a sus votantes su programa de gobierno.
  \ex Nadie \textit{tuvo en cuenta} sus opiniones sobre esta locución.
  \ex Todos \textit{cayeron en la cuenta} de que la caída del régimen se
      aproximaba.
  \ex Los autores \textit{dan por terminado} este capítulo del libro.
  \ex Siempre \textit{puede haber} algo más para agregar.
\end{exe}

\begin{resuelto}
1. a) LV: no acepta ni conmutación léxica ni conmutación estructural. b)
Equivale a \sig{ayudar}.
\end{resuelto}
\end{ejercicio}

\begin{recuadro}[Reflexión final]
En este capítulo hemos presentado un panorama de conjunto sobre un grupo de
estructuras de las que el hablante se sirve para indicar, de modo más específico
que con las formas verbales simples, distintos matices respecto del evento: si
ocurre antes de otro suceso, si ocurrirá en el futuro, si acaba de terminar o
está próximo a suceder, si se repite, etc. Las perífrasis y locuciones verbales
estudiadas son unidades sintácticas, ya que funcionan conjuntamente como núcleo
del predicado, pero también son unidades semánticas: tienen significado
unitario. No obstante, difieren porque, en el caso de las perífrasis, el
significado se interpreta composicionalmente, como resultado de lo que expresa
cada uno de sus constituyentes, mientras que, en el caso de las locuciones, se
trata de unidades más fosilizadas cuya interpretación no es literal.
\end{recuadro}

% --- Secciones finales -------------------------------------------------
\lecturasrecomendadas

\solucionario

\begin{solucion}{\ref{ej:11-reconocer}}
3. No es PV porque indica dos eventos y admite conmutación estructural:
\forma{Mi abuela ni salió ni paseó durante la pandemia}. 4. Es ambigua y posee
dos interpretaciones: i) no es PV porque admite conmutación estructural:
\forma{Volvió a la facultad para hacer el trabajo práctico}; ii) es PV y no
admite conmutación estructural. Acepta la paráfrasis \forma{Hizo el trabajo
práctico nuevamente}. 5. Es PV porque señala un único evento y no admite
conmutación estructural: \juicio{*}{\forma{El domingo empezó que llovía desde
temprano}}.
\end{solucion}

\begin{solucion}{\ref{ej:11-series}}
a) 2. \forma{Son devoradas} es PV. Se trata de un evento único; admite
conmutación léxica, \forma{son tragadas}, pero no conmutación estructural,
\juicio{*}{\forma{son y se devoran}}. 3. \forma{Puede distraer} es PV. Se trata
de un evento único; admite conmutación léxica, \forma{puede entretener}, pero no
conmutación estructural, \juicio{*}{\forma{puede y distrae}}. 4.
\forma{Comenzarlas a ver} es PV. Se trata de un evento único; admite conmutación
léxica, \forma{comenzarlas a mirar}, pero no conmutación estructural,
\juicio{*}{\forma{comenzar y verlas}}. 5. \forma{Permiten conectarnos} no es PV.
Se trata de dos eventos; admite conmutación léxica, \forma{permiten
vincularnos}, y conmutación estructural, \forma{permiten que nos conectemos}. 6.
\forma{Suelen proveer} es PV. Se trata de un evento único; admite conmutación
léxica, \forma{suelen suministrar}, pero no conmutación estructural,
\juicio{*}{\forma{suelen y proveen}}.

b) El fragmento reescrito reemplaza las formas en cursiva por \forma{nos hacen
experimentar, son tragadas, puede entretener, comenzarlas a mirar, permiten
vincularnos} y \forma{suelen suministrar}.\pendiente{El original transcribe el
fragmento completo reescrito}
\end{solucion}

\begin{solucion}{\ref{ej:11-tato}}
2. \forma{Iremos a comprar}: es PV. DE: V auxiliar \forma{ir} + \forma{a} + V
auxiliado \forma{comprar} en infinitivo. \forma{Tener que entregar}: es PV. DE: V
auxiliar \forma{tener} + \forma{que} + V auxiliado \forma{entregar} en
infinitivo. \forma{Dicen ser}: no es PV. \forma{Están corriendo}: es PV. DE: V
auxiliar \forma{estar} + V auxiliado \forma{correr} en gerundio. 3.
\forma{Estamos viviendo}: es PV. DE: V auxiliar \forma{estar} + V auxiliado
\forma{vivir} en gerundio. \forma{Estamos diciendo}: es PV. DE: V auxiliar
\forma{estar} + V auxiliado \forma{decir} en gerundio. 4. \forma{Estoy
aburrida}: no es PV. \forma{Dicen defenderme}: no es PV. 5. \forma{Iremos a
pagarle}: no es PV. \forma{Van a tener que entregar}:\footnote{En
§\ref{sec:11-cadenas} se describe el funcionamiento de las agrupaciones o
cadenas perifrásticas.} es PV. DE: V auxiliar\etq{1} \forma{ir} + \forma{a} + V
auxiliar\etq{2} \forma{tener} en infinitivo + \forma{que} + V auxiliado en
infinitivo. 6. \forma{Estamos pagando}: es PV. DE: V auxiliar \forma{estar} + V
auxiliado \forma{pagar} en gerundio. \forma{Vamos a pagar}: es PV. DE: V auxiliar
\forma{ir} + \forma{a} + V auxiliado en infinitivo.\pendiente{En 2,
\textit{iremos a comprar} figura como PV y, en 5, \textit{iremos a pagarle}, como
no PV: revisar}
\end{solucion}

\begin{solucion}{\ref{ej:11-vf}}
2. Verdadero. Los auxiliares de perífrasis solo pierden (parte de) su significado
cuando se combinan con una forma verbal no finita, pero en sus otros usos
continúan funcionando como verbos plenos. Ejemplos: \forma{Va todos los días a
la escuela} (V pleno con significado de movimiento); \forma{Va aprendiendo cada
día más} (V auxiliar que indica progreso del evento y no desplazamiento). 3.
Falso. Las perífrasis de participio solo aceptan la anteposición de los clíticos.
Ejemplo: \forma{El candidato ha rechazado el premio / lo ha rechazado /
*ha rechazádolo}. 4. Falso. Los auxiliares de perífrasis pueden construirse con
verbos impersonales. Ejemplo: \forma{Mañana puede llover todo el día}.
\end{solucion}

\begin{solucion}{\ref{ej:11-caracteristicas}}
2. No es PV porque se puede omitir el verbo principal en la interrogación. 3. Es
PV porque el verbo principal está reemplazado en la interrogación por el proverbo
\forma{hacer}. 4. Es PV porque es posible intercalar el adverbio
\forma{recientemente} entre los verbos que la componen. 5. No es PV porque
\forma{necesitar} mantiene su significado de V pleno, es decir que no está
desemantizado.
\end{solucion}

\begin{solucion}{\ref{ej:11-corpes}}
2. No se trata de una PV porque el participio \forma{descabellada} actúa como
adjetivo y funciona sintácticamente como el predicativo del V \forma{ser}. 3. Se
trata de una PV de \forma{ser} + participio que manifiesta la voz pasiva. 4. No
es una PV porque el gerundio funciona como un adverbio de modo y puede ser
reemplazado por \forma{así}. 5. \forma{Comenzó} y \forma{pude} no forman PV
porque son dos formas finitas; \forma{pude} forma PV con \forma{adivinar}. 6. Se
trata de una PV de \forma{comenzar} + \forma{a} + infinitivo que manifiesta el
inicio del evento.
\end{solucion}

\begin{solucion}{\ref{ej:11-pleno}}
2a. \sig{destruir}; 2b. \sig{comenzar abruptamente}. 3a. existencia; 3b.
anterioridad a otro evento del pasado. 4a. posesión de un atributo; 4b. voz
pasiva. 5a. ubicación; 5b. comienzo inminente.
\end{solucion}

\begin{solucion}{\ref{ej:11-significados}}
3. PV aspectual de cuantificación progresiva. 4. No es PV. 5. Es ambigua y posee
dos interpretaciones: i) no es PV porque indica dos eventos y admite conmutación
estructural: \forma{María va por el camino (mientras lee un libro) y aprende
gramática}; ii) PV aspectual de cuantificación progresiva. 6. PV aspectual de
fase continuativa. 7. PV de diátesis pasiva. 8. No es PV. 9. PV de modalidad
habitual. 10. PV aspectual de fase incoativa. 11. El contexto (la oposición entre
\forma{no volvió a rendir} y \forma{aprobó}) determina la interpretación de
\forma{volver} como PV aspectual de cuantificación reiterativa. 12. El contexto
(la oposición entre \forma{no vuelve a estudiar\ldots, sino para descansar})
determina la interpretación de \forma{volver} como V pleno de movimiento.
\end{solucion}

\begin{solucion}{\ref{ej:11-reformular}}
3. Mañana le voy a dar a mi compañero la contraseña. 4. Debe detenerse. 5. Aquí
suele nevar en invierno. 6. Ricardo se debe (de) haber dormido / puede haberse
dormido. 7. Enseguida me pongo / comienzo / empiezo a armar el equipo. 8. El
escritor acaba de terminar su última novela. 9. El escritor ha vivido durante
muchos años en San Juan, pero ahora vive en Santiago de Chile. 10. Mariela
volverá a ser mamá. 11. El profesor les pidió a sus alumnos que sigan / continúen
estudiando. 12. Por favor, ya dejen de pelear. 13. Todos sabemos que Miguel pudo
vencer a su rival. 14. Finalmente, Rosa pudo dormirse.
\end{solucion}

\begin{solucion}{\ref{ej:11-cadenas}}
2. \forma{Pudo} (posibilidad) + \forma{terminar de} (terminación) +
\forma{estudiar}. 3. \forma{Va a} (futuridad) + \forma{poder} (posibilidad) +
\forma{empezar a} (inicio) + \forma{estudiar}. 4. \forma{Deberían} (obligación) +
\forma{haber} (anterioridad) + \forma{empezado a} (inicio) + \forma{estudiar}. 5.
\forma{Tiene que} (obligación) + \forma{poder} (capacidad) + \forma{dejar de}
(interrupción / término) + \forma{comer}.
\end{solucion}

\begin{solucion}{\ref{ej:11-orden}}
\exref{ex:11-chomsky-b}: es gramatical porque la PV de modalidad precede a la de
voz pasiva. \exref{ex:11-chomsky-c}: es agramatical porque la PV de voz pasiva no
puede anteceder a ninguna otra en la oración. \exref{ex:11-chomsky-d}: es
gramatical porque la PV de tiempo compuesto precede a la aspectual reiterativa
con \forma{volver}. \exref{ex:11-chomsky-e}: es agramatical porque la PV
aspectual de fase continuativa no puede preceder a la de tiempo
compuesto.\pendiente{El original numera estas respuestas (70) a (73)}
\end{solucion}

\begin{solucion}{\ref{ej:11-tabla}}
\begin{center}
\resizebox{\linewidth}{!}{%
\begin{tabular}{@{}c *{6}{c} c@{}}
  \toprule
  & \multicolumn{6}{c}{Verbos auxiliares} & \\
  \cmidrule(lr){2-7}
  Ej. & \makecell{Mod. epistémica /\\habitualidad} & Temporalidad &
    \makecell{Mod.\\deóntica} & \makecell{Aspect.\\tipo II} &
    \makecell{Aspect.\\tipo I} & \makecell{Voz\\pasiva} & \makecell{Verbo\\auxiliado} \\
  \midrule
  2 & \textit{puede} & \textit{haber} & \textit{tenido que} & --- & \textit{dejar
    de} & \textit{ser} & \textit{atendido} \\
  3 & \textit{suelen} & --- & \textit{tener que} & \textit{volver a} & --- & ---
    & \textit{leer} \\
  4 & --- & \textit{han} & \textit{tenido que} & --- & --- & \textit{ser} &
    \textit{restaurados} \\
  5 & \textit{deben (de)} & \textit{haber} & --- & --- & --- & \textit{sido} &
    \textit{deterioradas} \\
  6 & --- & \textit{vamos a} & --- & --- & \textit{comenzar a} & --- &
    \textit{practicar} \\
  \bottomrule
\end{tabular}}
\end{center}
\end{solucion}

\begin{solucion}{\ref{ej:11-valores}}
2. Significado: modal deóntico (\forma{puede}) + aspectual continuativo
(\forma{seguir}) + de voz pasiva (\forma{siendo considerada}). La PV modal
indica que se admite o permite que Susana siga siendo considerada\ldots, por lo
que se considera deóntica. 3. Significado: modal epistémica (\forma{poder}) +
aspectual reiterativa (\forma{volver}) + pasiva (\forma{ser atendido}). La PV
modal se interpreta como \sig{tener la posibilidad de ser atendido\ldots}, por lo
que se considera epistémica.\pendiente{Según la jerarquía del capítulo,
\textit{poder} antepuesto a una aspectual con lectura de \textit{tener la
posibilidad} es deóntico: revisar} 4. Significado: témporo-aspectual
(\forma{ha}) + aspectual reiterativa (\forma{vuelto a}) + aspectual de
interrupción (\forma{dejar de tomar}).
\end{solucion}

\begin{solucion}{\ref{ej:11-entrar}}
2, 3, 5 y 6 son PPVV: no admiten conmutación estructural. La PV de
\forma{entrar} + infinitivo tiene valor aspectual de inicio. 4 no es PV: admite
conmutación estructural (\forma{entraron y robaron / entraron para robar}) y
\forma{entrar} actúa como verbo pleno.
\end{solucion}

\begin{solucion}{\ref{ej:11-locuciones}}
a) 2, 5 y 10 son PPVV y no aceptan conmutación estructural. Aceptan conmutación
léxica: 2. \forma{hemos asimilado}; 5. \forma{vamos a gozar}; 10. \forma{puede
existir}. 3, 4, 6, 7, 8 y 9 son LLVV: no aceptan ni conmutación léxica ni
conmutación estructural.\pendiente{El original numera «(1) hemos asimilado»} b)
Equivalencias: 3. \sig{extrañar}; 4. \sig{destrozar}; 6. \sig{informar}; 7.
\sig{considerar}; 8. \sig{reconocer}; 9. \sig{finalizar}.
\end{solucion}

% =====================================================================
%  capitulos/cap12.tex — Capítulo 12. Las relaciones interoracionales:
%  coordinación
%  Autora: Natalia Giollo
%  Fuente: Capitulo_12_normalizado.docx
% =====================================================================
\capitulo{Las relaciones interoracionales: coordinación}{Natalia Giollo}

\section{Introducción}

En este capítulo comenzamos el estudio de las relaciones que vinculan distintas
oraciones, también conocidas como \termino{relaciones interoracionales}. Para
ello, partimos de mostrar cómo una predicación simple puede ir ampliándose para
conformar estructuras más complejas, que incluyan otra u otras predicaciones, ya
sea secundarias, o incluidas, como en la subordinación,\footnote{La subordinación
se aborda en el \capref{13}.} o conservando cada una su independencia
sintáctica, como en la coordinación. Luego, nos centramos en la coordinación y
abordamos las particularidades de los distintos subtipos (copulativa, disyuntiva
y adversativa). Asimismo, también introducimos el concepto de marcador discursivo
y nos ocupamos de la identificación de las instancias de elipsis asociadas a la
coordinación.

\section{El pasaje de la oración simple a la oración compuesta}\label{sec:12-pasaje}

La oración, tal como se la definió en el \capref{6}, §\ref{sec:06-definicion},
constituye una unidad sintáctico-semántica, en la que un SN, como \forma{mi
hermano} en \exref{ex:12-hermano1}, y un SV, como \forma{lavó los platos}
\exref{ex:12-hermano1}, establecen una relación predicativa articulada a través
de las categorías funcionales de Tiempo y Concordancia para conformar un
sintagma de flexión (SF): \forma{Mi hermano lavó los platos}.

\begin{exe}
  \ex\label{ex:12-hermano1} {}[\etq{SF/O} [\etq{SN/Suj} Mi hermano]
      [\etq{SV/Pred} lavó los platos]]
\end{exe}

Ahora bien, además de expresar una predicación simple, la oración también puede
incorporar una segunda instancia predicativa, como cuando contiene algún
predicativo no obligatorio, que aporta una predicación adicional o secundaria
\exref{ex:12-hermano2} (véase el \capref{9}, §\ref{sec:09-primaria}):

\begin{exe}
  \ex\label{ex:12-hermano2} {}[\etq{SF/O} [\etq{SN/Suj} Mi hermano]
      [\etq{SV/Pred} lavó los platos [\etq{SA/PSNO} contento]]]
\end{exe}

En \exref{ex:12-hermano1} el SF contiene una única predicación, organizada a
partir de \forma{lavar}, un verbo transitivo que expresa una acción realizada
por un sujeto agente ---\forma{mi hermano}--- sobre un tema (el SN/COD \forma{los
platos}). En \exref{ex:12-hermano2} se mantiene la predicación de
\exref{ex:12-hermano1}, pero aparece una segunda instancia predicativa: la
\termino{predicación secundaria}, formada por el adjetivo \forma{contento},
referido al mismo sujeto de \exref{ex:12-hermano1}. Según vimos en el
\capref{9}, §\ref{sec:09-sintaxis}, el predicativo forma una oración reducida,
articulada sin flexión, ya que la oración reducida carece de la manifestación de
las categorías de tiempo, aspecto y modo. Asimismo, el sujeto en esta segunda
instancia predicativa ---\forma{(estar) contento}--- tiene otro papel temático:
es un \textsc{paciente}.

Si bien en \exref{ex:12-hermano2} hay dos predicaciones, el SF se sigue
considerando una oración simple, debido a que la segunda predicación no es
autónoma o completa, puesto que carece de flexión y se temporaliza a partir del
tiempo de la predicación principal. En efecto, en \exref{ex:12-hermano2}, la
única interpretación posible del evento \forma{(estar) contento} es que se
encuentra en el mismo tiempo pasado establecido por el evento de \forma{lavar}.
En este sentido, \exref{ex:12-hermano2} se podría parafrasear como se muestra en
\exref{ex:12-hermano3}.

\begin{exe}
  \ex\label{ex:12-hermano3} {}[\etq{SF/O} Mi hermano lavó los platos].
      [\etq{SF/O} (\textit{pro}) Estaba contento].
\end{exe}

En \exref{ex:12-hermano3}, hay dos sintagmas de flexión independientes y cada uno
constituye una oración simple. La estructura argumental que manifiestan es la
siguiente: \forma{lavar} ⟨alguien⟩ ⟨algo⟩ y \forma{estar contento} ⟨alguien⟩. En
la primera oración, los argumentos del verbo son ⟨\forma{mi hermano}⟩ y
⟨\forma{los platos}⟩. En la segunda, como ya adelantamos, el argumento
correspondiente al sujeto está representado sintácticamente por un sujeto tácito
o desinencial que, de acuerdo con lo desarrollado en el \capref{6},
§\ref{sec:06-tacito}, se corresponde con una forma pronominal abstracta, un
\textit{pro}, es decir, un pronombre fonéticamente vacío, pero semánticamente
interpretable y recuperable a partir del contexto de situación o de lo dicho
anteriormente. En este caso, \textit{pro} recoge el valor referencial del SN
\forma{mi hermano}, el sujeto de la oración inmediatamente anterior. A partir de
esa referencia entendemos que estas oraciones están vinculadas: lo que era el
contenido de una única oración simple \exref{ex:12-hermano2} ahora aparece
distribuido en dos oraciones, simples también, una contigua a la otra, en una
relación de \termino{yuxtaposición}.

\begin{destacado}
\textbf{Para recordar.} \textcite[334]{bosque2019} define la yuxtaposición como

\begin{cita}
Unión sintáctica de dos o más elementos en la cual se omite la marca formal (una
conjunción o una preposición), que expresa la relación gramatical que se
establece entre ellos [\ldots].
\end{cita}

Los elementos que se unen mediante yuxtaposición pueden ser sintagmas nominales,
adjetivales, preposicionales, adverbiales y verbales, además de oracionales. Al
igual que en la coordinación, los segmentos yuxtapuestos son elementos contiguos
del mismo nivel. La relación semántica que se da entre ellos es también semejante
a la que se establece en la coordinación (adición, exclusión, etc.).
\end{destacado}

Siguiendo esta línea de razonamiento, podríamos parafrasear
\exref{ex:12-hermano3} con una estructura como la presentada en
\exref{ex:12-hermano4}, en la que se emplea una conjunción para explicitar la
vinculación entre ambas oraciones:

\begin{exe}
  \ex\label{ex:12-hermano4} Mi hermano lavó los platos y estaba contento.
\end{exe}

En \exref{ex:12-hermano4}, a partir de la definición de oración que hemos
adoptado, se puede considerar que tenemos una \termino{oración compuesta} (OC)
constituida por dos oraciones simples (OS) coordinadas entre sí, cada una con su
sujeto: el SN \forma{mi hermano}, para la primera oración, y el pronombre vacío
\textit{pro}, para la segunda. Es cierto que, en términos referenciales, ambas
formas denotan la misma entidad del mundo extralingüístico y que comparten,
necesariamente, los rasgos de persona y número (tercera persona del singular)
responsables de la concordancia con los verbos \forma{lavó} y \forma{estaba},
respectivamente. Sin embargo, no se trata del mismo sujeto en términos
semánticos, porque cada instancia aporta un papel temático distinto. En este
sentido, en [\forma{mi hermano lavó los platos}] el sujeto es la entidad que
realiza la acción, el agente de \forma{lavar}, mientras que en [\forma{pro
estaba contento}], el sujeto es la entidad a la que se le atribuye un estado, con
el papel paciente de \forma{estar contento}. Es importante destacar que, tal como
se ha desarrollado en el \capref{7}, cada predicado tiene una estructura
argumental y temática propia, independientemente de que sus respectivos sujetos
compartan (o no) valores referenciales, sintácticos y semánticos.

Por último, también es posible reformular \exref{ex:12-hermano4} incluyendo una
oración en otra, es decir, utilizando una estructura incluida o subordinada a la
principal, como se muestra en \exref{ex:12-hermano5}:

\begin{exe}
  \ex\label{ex:12-hermano5}
    \begin{xlist}
      \ex\label{ex:12-hermano5-a} Mi hermano, [que estaba contento], lavó los
          platos.
      \ex\label{ex:12-hermano5-b} Mi hermano, [que lavó los platos], estaba
          contento.
      \ex\label{ex:12-hermano5-c} [Cuando lavó los platos] mi hermano estaba
          contento.
    \end{xlist}
\end{exe}

En los ejemplos de arriba se encuentran las mismas dos predicaciones
(\forma{lavar} ⟨alguien⟩ ⟨algo⟩, por un lado, y \forma{estar contento}
⟨alguien⟩, por el otro), ambas realizadas como predicados independientes. Pero
ahora una de ellas es una oración subordinada (que aparece entre corchetes), es
decir, una oración incluida en otra principal de la que forma parte
constituyendo una estructura mayor: una oración compuesta, en la que hay una
oración principal (predicación principal) y dentro de esta una oración
subordinada (predicación secundaria). Como se desprende del contraste entre
\exref{ex:12-hermano5-a}, \exref{ex:12-hermano5-b} y \exref{ex:12-hermano5-c},
el hablante puede crear distintas configuraciones sintácticas, preservando ---en
líneas generales--- el valor de cada predicación.

\begin{recuadro}[Sobre el concepto de oración compuesta]
Si bien en las gramáticas tradicionales se distinguía entre oración compuesta,
formada por oraciones coordinadas, y oración compleja, constituida por una
oración principal y una o varias subordinadas, en este libro seguimos los
planteos actuales de la \textit{NGLE} \parencite[§\,1.13ñ]{rae2025}, que adopta el
término \textit{oración compuesta} para todos los casos de combinación de
oraciones. Igualmente, en \textcite{bosque2019} se plantea que las oraciones que
forman una oración compuesta pueden estar vinculadas por coordinación
(\forma{María estudia y su hermana trabaja; He visto la película, pero no me ha
gustado}) o por subordinación (\forma{Me prometiste que vendrías; El contrato
que firmó no era válido}). Las oraciones compuestas también se pueden vincular
por yuxtaposición (\forma{Tu padre está cansado; pasa a buscarlo con el
coche}).
\end{recuadro}

En resumen, lo que se muestra de \exref{ex:12-hermano1} a \exref{ex:12-hermano5}
es que una oración simple, como \exref{ex:12-hermano1} \forma{Mi hermano lavó los
platos}, se puede ir transformando en estructuras que incorporen otras
predicaciones, en grados de complejidad creciente. Por ejemplo, en
\exref{ex:12-hermano2} \forma{Mi hermano lavó los platos contento}, la
predicación se expande con la adición del predicativo, que se añade a la
predicación principal. Asimismo, ambas predicaciones pueden ser recogidas en dos
oraciones independientes entre las que medie una mera yuxtaposición
\exref{ex:12-hermano3}: \forma{Mi hermano lavó los platos. Estaba contento}. O
bien, como hemos visto en \exref{ex:12-hermano4} y \exref{ex:12-hermano5},
podemos vincular ambas predicaciones, ya sea mediante conjunciones coordinantes
o subordinantes: la conjunción \forma{y}, que coordina dos predicaciones
principales, en \exref{ex:12-hermano4}, y los relativos \forma{que} y
\forma{cuando}, que establecen la relación entre la oración subordinada y la
principal, en \exref{ex:12-hermano5}. Si bien las predicaciones son las mismas,
lo que va variando es la manera en que se articulan entre sí, de modo que se
están relacionando oraciones (cláusulas o también proposiciones, según las
denominaciones dadas en las distintas tradiciones gramaticales) mediante
diferentes relaciones interoracionales.

En los siguientes apartados, nos centramos en analizar el complejo campo de las
relaciones interoracionales que se establecen mediante el mecanismo de la
coordinación, con énfasis en el análisis de los distintos tipos de conjunciones
coordinantes. Queda para el próximo capítulo el tratamiento más extenso de los
distintos tipos de oraciones subordinadas.

\begin{ejercicio}\label{ej:12-predicaciones}
En las siguientes oraciones, distinga la predicación primaria y la secundaria.
Luego, proponga una versión por yuxtaposición, otra por coordinación y otra por
subordinación, según se muestra en el ejemplo resuelto~1.

\begin{exe}
  \ex Los niños caminaron hacia el parque felices.
  \ex Los transeúntes miraban la escena asombrados.
  \ex María colgó el cuadro bien centrado.
\end{exe}

\begin{resuelto}
1. Predicación primaria: \forma{Los niños caminaron hacia el parque}.
Predicación secundaria: \forma{Los niños estaban felices}. Yuxtaposición:
\forma{Los niños caminaron hacia el parque, estaban felices}. Coordinación:
\forma{Los niños caminaron hacia el parque y estaban felices}. Subordinación: a)
\forma{Los niños, que estaban felices, caminaron hacia el parque}; b) \forma{Los
niños, que caminaron hacia el parque, estaban felices}; c) \forma{Cuando
caminaron hacia el parque, los niños estaban felices}.
\end{resuelto}
\end{ejercicio}

\section{La coordinación}\label{sec:12-coordinacion}

La \termino{coordinación} puede definirse como «el proceso gramatical por el cual
dos o más constituyentes sintácticos (los "coordinandos") se unen con una
conjunción formando una frase compleja» \parencite[301]{camacho2016}. Mediante la
coordinación se pueden vincular distintos tipos de constituyentes, que
generalmente son sintagmas de la misma categoría, de modo tal que la
coordinación se caracteriza como \termino{homocategorial}.

\begin{exe}
  \ex\label{ex:12-juan-pedro} {}[\etq{SSNN} Juan \textit{y} Pedro] fueron al río.
  \ex\label{ex:12-rio} El río estaba [\etq{SSAA} calmo, \textit{pero} muy frío].
  \ex\label{ex:12-ahora} Lo reclamamos [\etq{SSAdv} ahora \textit{o} nunca].
  \ex\label{ex:12-animo} Juan se encuentra [\etq{SSPP} de buen ánimo \textit{y} en
      buen estado de salud].
  \ex\label{ex:12-maria} {}[[\etq{SF} María llegó a la oficina] \textit{y}
      [\etq{SF} encendió la computadora]].
\end{exe}

La conjunción no solo coordina constituyentes que pertenecen a la misma clase,
sino que también une constituyentes de distinta clase, pero con una función
equivalente, por lo que también existe la coordinación
\termino{heterocategorial}.

\begin{exe}
  \ex\label{ex:12-cansado} Entró [[\etq{SA} cansado] \textit{y} [\etq{SP} sin
      prisa]]. $\rightarrow$ Coordinación de un SA y un SP.
  \ex\label{ex:12-rapido} Llegó [[\etq{SAdv} rápidamente] \textit{y} [\etq{SP}
      con buena disposición]]. $\rightarrow$ Coordinación de un SAdv y un SP.
  \ex\label{ex:12-bebida} Compramos [\etq{SN} la bebida] \textit{y} [\etq{SF}
      todo lo que estaba en oferta]. $\rightarrow$ Coordinación de un SN y un
      SF.
\end{exe}

\subsection{La estructura de la coordinación}\label{sec:12-estructura}

En relación con la coordinación, tanto la gramática tradicional como la
estructural planteaban una estructura plana en la que los constituyentes unidos
por la conjunción estaban en el mismo nivel y no había un núcleo que diera la
categoría al sintagma resultante de dicha unión:

\begin{exe}
  \ex {}[[\etq{SN} Los estudiantes] \textit{y} [\etq{SN} sus profesores]]
      cantaron el himno nacional en la ceremonia.
  \ex {}[[\etq{SF} El gobierno aumentó los impuestos], \textit{pero} [\etq{SF}
      la ciudadanía se manifestó en contra]].
\end{exe}

Sin embargo, una segmentación como la de arriba no se adecua a los planteos de la
teoría sintagmática propuesta por la gramática generativa (y presentada en el
\capref{5}, §\ref{sec:05-xbarra}), ya que no permite identificar un núcleo que
otorgue la categoría a toda la construcción. Al respecto, recordemos que, según
el esquema de la teoría X\prima, todos los sintagmas son endocéntricos, es decir,
tienen un núcleo, que puede estar solo o ir acompañado de un complemento en el
nivel X\prima, y que ambos, núcleo y complemento, pueden recibir un especificador
que incide sobre ellos, tal como se muestra en \exref{ex:12-xbarra}.

\begin{exe}
  \ex\label{ex:12-xbarra}
      \begin{forest} gramatica
        [SX
          [Esp]
          [X\prima
            [X\\(núcleo)]
            [Comp]]]
      \end{forest}
\end{exe}

Siguiendo, entonces, estos planteos, necesitamos recurrir a un nuevo sintagma de
naturaleza funcional: el \termino{sintagma de conjunción} (SConj). Su núcleo es
una palabra funcional, la \termino{conjunción coordinante} (Conj), que toma como
complemento al segundo constituyente de la coordinación y, como especificador,
al primer constituyente, según el esquema de \exref{ex:12-sconj}.

\begin{exe}
  \ex\label{ex:12-sconj}
      \begin{forest} gramatica
        [SConj
          [Esp\\(1.\textsuperscript{er} coordinando)]
          [Conj\prima
            [Conj\\(núcleo)]
            [Comp\\(2.\textsuperscript{o} coordinando)]]]
      \end{forest}
\end{exe}

Es importante resaltar que el SConj se proyecta en toda instancia en que aparece
una conjunción, más allá de la naturaleza de los elementos coordinados. En todos
los casos, en términos estructurales, el elemento que aparece a la derecha de la
conjunción coordinante ocupa la posición de complemento, mientras que el
elemento a la izquierda se ubica en la posición de especificador. Veamos, en
\exref{ex:12-an-juan}--\exref{ex:12-an-maria}, cómo se aplica este análisis a las
oraciones de \exref{ex:12-juan-pedro} a \exref{ex:12-maria}.

\begin{exe}
  \ex\label{ex:12-an-juan} {}[\etq{SConj} [\etq{SN/Esp} Juan] \textit{y}\etq{Conj/n}
      [\etq{SN/Comp} Pedro]] fueron al río.
  \ex\label{ex:12-an-rio} El río estaba [\etq{SConj} [\etq{SA/Esp} calmo],
      \textit{pero}\etq{Conj/n} [\etq{SA/Comp} muy frío]].
  \ex Lo reclamamos [\etq{SConj} [\etq{SAdv/Esp} ahora] \textit{o}\etq{Conj/n}
      [\etq{SAdv/Comp} nunca]].
  \ex Juan se encuentra [\etq{SConj} [\etq{SP/Esp} de buen ánimo]
      \textit{y}\etq{Conj/n} [\etq{SP/Comp} en buen estado de salud]].
  \ex\label{ex:12-an-maria} {}[\etq{SConj} [\etq{SF/Esp} María llegó a la
      oficina] \textit{y}\etq{Conj/n} [\etq{SF/Comp} encendió la computadora]].
\end{exe}

Ahora bien, analizar dos (o más) constituyentes como los de la coordinación, con
la misma o equivalente naturaleza categorial, como una estructura jerárquica con
diferentes niveles puede, en principio, parecer contraintuitivo. Por eso,
dedicamos el próximo apartado a fundamentar estos planteos.

\subsection{Sobre la estructura jerárquica de la coordinación}\label{sec:12-jerarquica}

El mecanismo de la coordinación tiene la particularidad de presentar al mismo
tiempo características de equivalencia categorial entre sus miembros, pero
también, desde la perspectiva estructural o de conformación sintagmática, se
reconocen relaciones jerárquicas entre los coordinandos. En principio, aunque los
coordinandos compartan la misma naturaleza categorial, a partir de la estructura
propuesta en \exref{ex:12-sconj}, asumimos una relación jerárquica o asimétrica
entre ellos, en tanto el complemento y el especificador no se encuentran en el
mismo nivel, sino que ocupan diferentes posiciones estructurales.

Asimismo, la vinculación entre ambas oraciones coordinadas requiere una operación
sintáctica básica: el \termino{ensamble} o
\termino{fusión}.\pendiente{El original dice que el ensamble «ya [fue]
mencionado en el capítulo 6», pero no aparece allí} Como se explica en
\textcite[219]{bosque2009}, el ensamble es «el proceso fundamental por el que
formamos nuevos constituyentes». Cuando dos expresiones se ensamblan, se produce
un nuevo constituyente que hereda los rasgos categoriales de uno de los
elementos fusionados. En el caso de las estructuras coordinadas, primero se
ensambla la conjunción coordinante, que actúa como núcleo de la construcción, con
su complemento, que es el segundo miembro de la coordinación. Luego, el
constituyente resultante (Conj\prima) se ensambla con el primer miembro de la
coordinación, que está en la posición de especificador, dando lugar a la
proyección máxima del sintagma (SConj). En \exref{ex:12-arbol-maria} presentamos
la estructura arbórea de \exref{ex:12-maria}.

\begin{exe}
  \ex\label{ex:12-arbol-maria}
      \begin{forest} gramatica
        [SConj
          [SF\\Esp [\textit{María llegó a la oficina}, roof]]
          [Conj\prima
            [Conj\\\textit{y}]
            [SF\\Comp [\textit{encendió la computadora}, roof]]]]
      \end{forest}
\end{exe}

Ahora bien, asumir que el segundo coordinando es el complemento de la conjunción
implica plantear que entre este y la conjunción existe una relación más estrecha
---de complementación--- que la que se da entre la conjunción y el primer
coordinando. ¿Cómo podemos fundamentar esta diferencia estructural? Con el fin de
recoger evidencia, prestemos atención a los datos que nos puede proporcionar un
simple intercambio entre dos hablantes, A y B. En \exref{ex:12-dialogo-a}, el
hablante A hace una afirmación y su interlocutor B la continúa en una típica
réplica \exref{ex:12-dialogo-b}, que puede ser tanto enunciativa como
interrogativa y está conformada solo por la conjunción y su complemento. Según
vemos en la réplica de \exref{ex:12-dialogo-b}, la estructura formada por la
conjunción seguida del segundo coordinando (su complemento), pero sin incluir al
primer constituyente (que ya fue expresado por el primer emisor), resulta
perfectamente gramatical. En cambio, si la réplica del interlocutor fuera como
\exref{ex:12-dialogo-c} y estuviera formada por el primer constituyente seguido
de la conjunción (sin el complemento), no resultaría gramatical.

\begin{exe}
  \ex\label{ex:12-dialogo}
    \begin{xlist}
      \ex\label{ex:12-dialogo-a} A: ---Mañana vendrá María.
      \ex\label{ex:12-dialogo-b} B: ---\textit{Y} Pedro. / ---¿\textit{Y} Pedro?
      \ex[*]{\label{ex:12-dialogo-c} B\prima: ---Juan \textit{y}.}
    \end{xlist}
\end{exe}

Veamos otros dos ejemplos, que nos permitirán reforzar la argumentación
anterior:

\begin{exe}
  \ex\label{ex:12-juguete}
    \begin{xlist}
      \ex Todavía no sabemos si fue un accidente, \textit{pero} rompimos el
          juguete en mil pedazos. (H.~Casciari, \textit{Teníamos un juguete})
      \ex[*]{\label{ex:12-juguete-b} Todavía no sabemos si fue un accidente,
          \textit{pero}.}
      \ex\label{ex:12-juguete-c} \textit{Pero} rompimos el juguete en mil
          pedazos.
    \end{xlist}
  \ex\label{ex:12-padres}
    \begin{xlist}
      \ex Algunos padres todavía se amontonan frente a las puertas \textit{y}
          gritan los nombres de sus hijos. (S.~Schweblin, \textit{Mariposas})
      \ex[*]{\label{ex:12-padres-b} Algunos padres todavía se amontonan frente a
          las puertas \textit{y}.}
      \ex\label{ex:12-padres-c} \textit{Y} gritan los nombres de sus hijos.
    \end{xlist}
\end{exe}

A partir del contraste de gramaticalidad entre los casos (b) y (c) de los
ejemplos \exref{ex:12-juguete} y \exref{ex:12-padres}, podemos nuevamente
comprobar que la conjunción coordinante establece una relación muy estrecha con
su complemento y que solamente es gramatical la estructura en que va seguida de
su complemento, con ausencia de especificador, como en \exref{ex:12-juguete-c} y
\exref{ex:12-padres-c}; mientras que ni \exref{ex:12-juguete-b} ni
\exref{ex:12-padres-b} resultan gramaticales debido, justamente, a la ausencia
del complemento de la conjunción, que no puede ser omitido.\footnote{El
complemento puede omitirse si queremos generar suspenso: \forma{Algunos padres
gritan los nombres de sus hijos y\ldots}, caso en que se considera que el
complemento queda implícito.}

Por otra parte, también debemos señalar que existe una equivalencia funcional
entre los coordinandos cuando se toman por separado, es decir que ambos podrían
cumplir independientemente la misma función sintáctica. En \exref{ex:12-rio},
reproducida aquí como \exref{ex:12-equiv-a}, tanto el SA \forma{calmo} como el
SA \forma{muy frío} podrían funcionar como PSO en un contexto sin coordinación. A
su vez, en \exref{ex:12-equiv-b}, tanto el SA \forma{cansado} como el SP
\forma{sin prisa} podrían actuar como PSNO: \forma{Entró cansado; Entró sin
prisa}.

\begin{exe}
  \ex
    \begin{xlist}
      \ex\label{ex:12-equiv-a} El río estaba [\etq{SA/PSO} calmo]. El río estaba
          [\etq{SA/PSO} muy frío].
      \ex\label{ex:12-equiv-b} Entró [\etq{SA/PSNO} cansado]. Entró
          [\etq{SP/PSNO} sin prisa].
    \end{xlist}
\end{exe}

Ahora bien, cuando los elementos están coordinados, asumimos que el SConj es el
que cumple la función sintáctica que ambos tenían por separado. Al mismo tiempo,
al establecer entre ambos miembros una relación jerárquica dentro del SConj, cuyo
núcleo es la conjunción, el primer coordinando pasa a funcionar como Esp y el
segundo, como Comp:

\begin{exe}
  \ex
    \begin{xlist}
      \ex El río estaba [\etq{SConj/PSO} [\etq{SA/Esp} calmo],
          \textit{pero}\etq{Conj/n} [\etq{SA/Comp} muy frío]].
      \ex Entró [\etq{SConj/PSNO} [\etq{SA/Esp} cansado] \textit{y}\etq{Conj/n}
          [\etq{SP/Comp} sin prisa]].
    \end{xlist}
\end{exe}

\begin{ejercicio}\label{ej:12-desantis}
En las siguientes oraciones de \textit{La sexta lámpara}, de Pablo De Santis,
señale entre corchetes los dos miembros de cada coordinación y la estructura
resultante, como se muestra en el ejemplo resuelto~1.

\begin{exe}
  \ex El viejo Treviso supo aprovechar nuestro fanatismo y empezó a pedirnos
      pequeños préstamos a cambio de una información siempre escatimada y a
      menudo contradictoria.
  \ex Hicimos una colecta para pagarle el legado, sin saber si comprábamos un
      tesoro o un montón de diarios viejos.
  \ex Habíamos subido el baúl por la escalera casi a los saltos, pero la visión
      de aquellos papeles nos paralizó.
  \ex Balestri sigue siendo un desconocido, y de su torre no poseemos más que la
      sombra de una sombra.
\end{exe}

\begin{resuelto}
1. [\etq{SConj/OC} [\etq{SF/Esp} El viejo Treviso supo aprovechar nuestro
fanatismo] y\etq{Conj/n} [\etq{SF/Comp} empezó a pedirnos pequeños préstamos a
cambio de una información [\etq{SConj/M} [\etq{SV(participio)/Esp} siempre
escatimada] y\etq{Conj/n} [\etq{SA/Comp} a menudo contradictoria]]]].
\end{resuelto}
\end{ejercicio}

\begin{ejercicio}\label{ej:12-enriquez}
Señale entre corchetes las instancias de coordinación del siguiente fragmento de
\textit{Éste es el mar}, de Mariana Enriquez, tal como se observa en el ejemplo
resuelto.

\begin{textolectura}
Tenían una madre que rara vez era vista o convocada; la madre Hécate. En el
Enjambre se creía que era una Leyenda, una mentira, pero las Luminosas se reían
de tanta ignorancia. Y tenían hermanas mayores: Perséfone bajo la tierra, con
flores en el pelo; y las Imago, las que se escondían en lugares secos y oscuros.
\end{textolectura}

\begin{resuelto}
[\etq{Oración 1} Tenían una madre que rara vez era
[\etq{SConj}\footnote{En este caso, la conjunción une dos formas verbales: los
dos auxiliados de la perífrasis de voz pasiva \forma{era vista o convocada}.
Como la coordinación se da dentro de la perífrasis verbal, el SConj no tiene
función sintáctica. Las perífrasis verbales se tratan en el
\capref{11}.\pendiente{El original dice «los dos auxiliares» y etiqueta
\textit{vista} y \textit{convocada} como Vaux}} [\etq{V/Esp} vista]
o\etq{Conj/n} [\etq{V/Comp} convocada]]; la madre Hécate].
\end{resuelto}
\end{ejercicio}

\begin{destacado}
\textbf{En síntesis.} Estructuralmente, la coordinación es una relación que une
dos constituyentes (generalmente, de la misma categoría) entre los que establece
una relación jerárquica. Por tanto, podemos afirmar que, al igual que el resto de
los sintagmas estudiados en otros capítulos, la coordinación se adecua al esquema
de X\prima, respetando el principio de endocentricidad: posee un núcleo funcional
---la conjunción--- que toma un complemento obligatorio (el segundo coordinando) y
un especificador que puede, eventualmente, omitirse (el primer coordinando).
\end{destacado}

\subsection{Las conjunciones coordinantes}\label{sec:12-conjunciones}

Para analizar la relación que se establece entre los coordinandos es fundamental
identificar la naturaleza del núcleo de la coordinación, es decir, de las
distintas conjunciones coordinantes. En principio, las conjunciones coordinantes
pueden ser de dos tipos:

\begin{itemize}
  \item \termino{conjunciones coordinantes simples}: \forma{y}, con su variante
        \forma{e}; \forma{ni} (= \forma{y no}); \forma{o}, con su variante
        \forma{u}; \forma{pero}; \forma{sino}; \forma{mas};
  \item \termino{conjunciones coordinantes correlativas}: \forma{ni\ldots\
        ni\ldots; tanto\ldots\ como\ldots; tanto\ldots\ cuanto\ldots;
        así\ldots\ como\ldots; sea\ldots\ sea\ldots; ya\ldots\ ya\ldots;
        ora\ldots\ ora\ldots; bien\ldots\ bien\ldots; o\ldots\ o\ldots}
\end{itemize}

Es importante destacar que, más allá del hecho de que las conjunciones
coordinantes sean simples o compuestas, la coordinación puede ser

\begin{itemize}
  \item \termino{binaria}, cuando se establece entre solo dos coordinandos;
  \item \termino{múltiple}, cuando vincula varios coordinandos.
\end{itemize}

Algunas conjunciones pueden usarse de forma binaria o múltiple. La conjunción
\forma{y}, por ejemplo, es simple, pero puede participar tanto en coordinaciones
binarias, como las de \exref{ex:12-ajedrez-a} y \exref{ex:12-ajedrez-b}, como en
múltiples, como las de \exref{ex:12-ajedrez-c} y \exref{ex:12-ajedrez-d}:

\begin{exe}
  \ex
    \begin{xlist}
      \ex\label{ex:12-ajedrez-a} Juan juega al ajedrez \textit{y} a las damas.
      \ex\label{ex:12-ajedrez-b} Juan juega al ajedrez \textit{y} Pedro juega a
          las damas.
      \ex\label{ex:12-ajedrez-c} Juan juega al ajedrez, las damas, el backgammon
          \textit{y} el póker.
      \ex\label{ex:12-ajedrez-d} Juan juega al ajedrez, Pedro prefiere las damas
          \textit{y} su hermana María ama el póker.
    \end{xlist}
\end{exe}

Las conjunciones que pueden aparecer en coordinaciones múltiples son: \forma{y,
ni, tanto\ldots\ como, o, o sea, bien\ldots\ bien}. En cambio, las conjunciones
\forma{pero} y \forma{sino} solo pueden participar en coordinaciones binarias,
como se aprecia en los ejemplos de \exref{ex:12-binarias}:\pendiente{\textit{o
sea} figura entre las conjunciones, pero en la tabla~\ref{tab:12-marcadores} se
clasifica como marcador explicativo}

\begin{exe}
  \ex\label{ex:12-binarias}
    \begin{xlist}
      \ex Juan juega al ajedrez, \textit{pero} Pedro prefiere jugar a las damas.
      \ex[*]{Juan juega al ajedrez, \textit{pero} Pedro prefiere jugar a las
          damas, \textit{pero} su hermana María ama jugar al póker.}
      \ex Juan no juega al ajedrez, \textit{sino} a las damas.
      \ex[*]{Juan no juega al ajedrez, \textit{sino} a las damas, \textit{sino}
          al póker.}
    \end{xlist}
\end{exe}

Por otra parte, según su funcionamiento sintáctico-semántico, las conjunciones se
clasifican en tres grupos:

\begin{itemize}
  \item \termino{copulativas}: \forma{y, e, ni, tanto\ldots\ como\ldots};
  \item \termino{disyuntivas}: \forma{o, u, bien\ldots\ bien\ldots,
        sea\ldots\ sea\ldots};
  \item \termino{adversativas}: \forma{pero, sino, mas}.
\end{itemize}

Cuando la conjunción coordinante es copulativa, la relación de coordinación se
interpreta como una suma o adición. En cambio, cuando es disyuntiva, se entiende
que hay una elección o alternancia. En el caso de las conjunciones adversativas,
se plantea una relación de oposición entre los coordinandos. En los próximos
apartados, describimos los distintos tipos de conjunciones coordinantes y las
características de las estructuras de coordinación que conforman.

\subsubsection{La coordinación copulativa}\label{sec:12-copulativa}

En la coordinación copulativa se establece entre los coordinandos un vínculo de
adición. La conjunción coordinante copulativa simple \forma{y} constituye el
recurso más empleado para la coordinación. Otras conjunciones copulativas son:
\forma{ni; tanto\ldots\ como\ldots};\footnote{\textcite[194]{jimenezjulia2023}
explica que el valor de \forma{tanto\ldots\ como\ldots} se solapa con el de la
conjunción copulativa \forma{y}, en tanto ambas formas permiten combinar acciones
que se producen en paralelo: según el autor, para expresar la acción paralela el
español, como gran parte de las lenguas indoeuropeas, gramaticalizó la forma
\forma{tanto X como Y}, heredera de una estructura latina que también originó
esquemas comparativos. Señala también que esta forma gramaticalizada tiene un
correlato negativo: \forma{ni X ni Y}.} \forma{ni\ldots\ ni\ldots;
tanto\ldots\ cuanto\ldots; así\ldots\ como}. Cuando se trata de una
coordinación múltiple, como ya hemos anticipado, la conjunción \forma{y}
normalmente se coloca solo ante el último elemento, pero se interpreta un
vínculo de suma entre todos los elementos del conjunto.

\begin{destacado}
\textbf{Para tener en cuenta.} La omisión de la conjunción ante el último
elemento o bien su repetición ante cada elemento coordinado genera efectos que
dan lugar a dos figuras retóricas. Como se explica en la \textit{NGLE}
\parencite[§\,31.2b]{rae2025}, el proceso de coordinar expresiones sin
conjunción manifiesta se denomina \termino{asíndeton} y constituye una de las
formas de yuxtaposición; la suspensión de la conjunción final deja a menudo en
suspenso la enumeración. En este caso, se interpreta que la enumeración no es
exhaustiva: \forma{Compramos cuadernos, hojas, lápices, lapiceras\ldots} En
cambio, se denomina \termino{polisíndeton} la figura retórica que consiste en el
proceso contrario, es decir, la repetición de la conjunción copulativa ante cada
miembro coordinado: se trata de un recurso enfático dirigido a lograr mayor
energía en la expresión, como en \forma{Se lo dije una, y dos, y tres veces, pero
no me hacía caso} \parencite[§\,31.2c]{rae2025}.
\end{destacado}

\paragraph{Lectura distributiva y colectiva de la coordinación copulativa.}\label{sec:12-distributiva}
En general, cuando se construyen predicados en los que participan sintagmas
nominales unidos por una conjunción copulativa, emerge una \termino{lectura
distributiva}, de modo que a cada coordinando le corresponde la misma
predicación, por lo que \exref{ex:12-teatro-a} puede parafrasearse como
\exref{ex:12-teatro-b}.

\begin{exe}
  \ex
    \begin{xlist}
      \ex\label{ex:12-teatro-a} María \textit{y} Ana fueron al teatro.
      \ex\label{ex:12-teatro-b} María fue al teatro \textit{y} Ana fue al
          teatro.
    \end{xlist}
\end{exe}

Por otro lado, cuando tenemos predicados de naturaleza colectiva (véase el
recuadro \textit{Para ampliar}), como en \exref{ex:12-parecen} y
\exref{ex:12-cunadas}, la presencia de una coordinación requiere de una
interpretación conjunta o acumulativa (\termino{lectura colectiva}), y se rechaza
la lectura distributiva.

\begin{exe}
  \ex\label{ex:12-parecen}
    \begin{xlist}
      \ex Pedro y Joaquín se parecen.
      \ex[*]{\label{ex:12-parecen-b} Pedro se parece \textit{y} Joaquín se
          parece.}
    \end{xlist}
  \ex\label{ex:12-cunadas}
    \begin{xlist}
      \ex María \textit{y} Ana son cuñadas.
      \ex[*]{\label{ex:12-cunadas-b} María es cuñada \textit{y} Ana es cuñada.}
    \end{xlist}
\end{exe}

\begin{recuadro}[Para ampliar]
De acuerdo con la \textit{NGLE} \parencite[§\,31.6c]{rae2025}, los predicados de
naturaleza colectiva «se suelen llamar simétricos y se caracterizan por
predicarse de grupos, manifestados mediante plurales o elementos coordinados».
Estos predicados pueden ser verbos (\forma{casarse, conversar, juntar,
negociar}), adjetivos (\forma{coincidente, compatible, diferente}) o sustantivos
(\forma{adversario, amigo, cuñado, pariente, socio}). También tienen
características de predicados simétricos un SP con \forma{entre} como núcleo y
algunos adverbios o locuciones adverbiales, como \forma{lejos, cerca,
conjuntamente, al lado, en colaboración}.
\end{recuadro}

Como resulta claro a partir de la agramaticalidad de \exref{ex:12-parecen-b} y
\exref{ex:12-cunadas-b}, el predicado no puede repetirse para cada elemento del
SConj, por lo que la única lectura posible es la colectiva: se predica de los
coordinandos en forma conjunta. Tampoco hay ambigüedad cuando se emplean, en
cualquier predicado, conjunciones copulativas compuestas, también llamadas
\textit{discontinuas} o \textit{correlativas}, porque impulsan una
interpretación distributiva. En este sentido, la agramaticalidad de una oración
como \exref{ex:12-tanto} se debe a la tensión que se produce entre la lectura
colectiva del predicado \forma{ser cuñadas} y la lectura distributiva que
requiere la conjunción \forma{tanto\ldots\ como}:

\begin{exe}
  \ex[*]{\label{ex:12-tanto} Tanto María como Ana son cuñadas entre sí.}
\end{exe}

Sin embargo, en predicados no simétricos, es posible encontrar ejemplos de
ambigüedad entre la interpretación colectiva y la distributiva de \forma{y}. Por
ejemplo, una oración como \exref{ex:12-pastel-a} puede interpretarse como
\exref{ex:12-pastel-b} o como \exref{ex:12-pastel-c}:

\begin{exe}
  \ex
    \begin{xlist}
      \ex\label{ex:12-pastel-a} María \textit{y} Ana cocinaron un pastel de papas
          y un budín.
      \ex\label{ex:12-pastel-b} María \textit{y} Ana cocinaron un pastel de papas
          y un budín \textit{juntas}. (Lectura colectiva).
      \ex\label{ex:12-pastel-c} María cocinó un pastel de papas y un budín,
          \textit{y} Ana cocinó un pastel de papas y un budín. (Lectura
          distributiva).
      \ex\label{ex:12-pastel-d} María \textit{y} Ana cocinaron un pastel de papas
          y un budín, \textit{respectivamente}. (Lectura ordenada).
    \end{xlist}
\end{exe}

Asimismo, como vemos en \exref{ex:12-pastel-d}, la ambigüedad puede evitarse
mediante la incorporación de \forma{respectivamente}.

\begin{ejercicio}\label{ej:12-lecturas}
Lea las siguientes oraciones y proponga las lecturas posibles, según se muestra
en el ejemplo resuelto~1.

\begin{exe}
  \ex Frodo y Sam compraron frutas y galletas para el viaje.
  \ex Frodo y Sam se llevan bien desde la infancia.
  \ex Frodo y Sam viajan a Mordor.
  \ex Sam y Gollum se odian.
  \ex Después de la batalla, Frodo y Aragorn viajan hacia Mordor y hacia el
      Bosque de Fangorn.
\end{exe}

\begin{resuelto}
1. a) Lectura distributiva: \forma{Frodo compró frutas y galletas para el viaje
y Sam también compró frutas y galletas para el viaje} / \forma{Tanto Frodo como
Sam compraron frutas y galletas}. b) Lectura colectiva: \forma{Frodo y Sam
compraron frutas y galletas juntos}. c) Lectura ordenada: \forma{Frodo y Sam
compraron frutas y galletas, respectivamente}.
\end{resuelto}
\end{ejercicio}

\paragraph{Coordinación simétrica y asimétrica.}\label{sec:12-simetrica}
Por último, retomemos un aspecto de los coordinandos que ya habíamos mencionado
antes, por lo que repetiremos los ejemplos de \exref{ex:12-hermano4},
\exref{ex:12-maria}, \exref{ex:12-cansado} y \exref{ex:12-bebida}, reformulados
como \exref{ex:12-sim-a}--\exref{ex:12-sim-d}:

\begin{exe}
  \ex
    \begin{xlist}
      \ex\label{ex:12-sim-a} Mi hermano lavó los platos y estaba contento.
      \ex María llegó a la oficina y encendió la computadora.
      \ex Entró cansado y sin prisa.
      \ex\label{ex:12-sim-d} Compramos la bebida y todo lo que estaba en oferta.
    \end{xlist}
\end{exe}

En principio, la coordinación copulativa establece una relación simétrica entre
los coordinandos, que se manifiesta no solo en la homogeneidad categorial y
funcional, que hemos visto en §\ref{sec:12-coordinacion}, sino también en el
orden: si la relación es simétrica, los coordinandos son intercambiables. Ahora
bien, como afirma \textcite[282]{ditullio2014}, cuando tal intercambio no es
posible, «es porque se asocian al coordinante otros valores semánticos» y se
produce lo que se conoce como \termino{coordinación asimétrica}, en la que los
coordinandos no son intercambiables. Veamos en \exref{ex:12-inv} los resultados
del cambio de orden de las oraciones anteriores.

\begin{exe}
  \ex\label{ex:12-inv}
    \begin{xlist}
      \ex\label{ex:12-inv-a} Mi hermano estaba contento \textit{y} lavó los
          platos.
      \ex[*]{\label{ex:12-inv-b} María encendió la computadora \textit{y} llegó
          a la oficina.}
      \ex\label{ex:12-inv-c} Entró sin prisa \textit{y} cansado.
      \ex\label{ex:12-inv-d} Compramos todo lo que estaba en oferta \textit{y} la
          bebida.
    \end{xlist}
\end{exe}

Notamos que, luego del cambio de orden, la gramaticalidad se mantiene
indiscutiblemente en \exref{ex:12-inv-a}, \exref{ex:12-inv-c} y
\exref{ex:12-inv-d}, por lo que validamos que se trata de coordinaciones
simétricas. Tal vez podríamos señalar que, con la inversión,
\exref{ex:12-inv-a} adquiere cierto matiz de causa, en tanto podría entenderse
como equivalente a \forma{Como mi hermano estaba contento, lavó los platos}. En
cambio, \exref{ex:12-inv-b} es agramatical,\pendiente{Con la lectura de
sucesión, \exref{ex:12-inv-b} es gramatical, aunque describe otra secuencia de
eventos; quizá corresponda \# en lugar de *} lo que nos lleva a afirmar que los
coordinandos no son simétricos en su configuración, es decir, no se puede
alterar el orden de aparición. En este caso, notamos que el orden de los
constituyentes se encuentra asociado a una relación temporal que no permite la
inversión: María primero llega a la oficina y luego enciende la computadora para
trabajar.

Además del valor temporal, la conjunción copulativa puede adquirir otros valores
semánticos: condicional, concesivo, adversativo, ilativo. Tal como se explica en
la \textit{NGLE} \parencite[§\,31.2m]{rae2025}, cuando uno de los dos predicados
coordinados no expresa un suceso, sino una situación o un estado de cosas, no se
obtiene la interpretación temporal. En esos casos son comunes la interpretación
adversativa (\forma{Tenía examen y no era capaz de recordar nada}, es decir,
\sig{y sin embargo\ldots}), la ilativa (\forma{Estaba cansado y me fui a
dormir}, es decir, \sig{luego\ldots}) o la condicional, frecuente en los
imperativos retóricos (\forma{Cásate y verás}).

Respecto de estos valores asociados a la conjunción, algunos autores consideran
que se trata, en realidad, de un fenómeno vinculado con aspectos pragmáticos, es
decir, con el uso de la lengua. Por ejemplo, \textcite[195]{jimenezjulia2023}
afirma que la posibilidad de interpretar contextualmente otros valores «no es una
consecuencia del valor de \textit{y}, sino de nuestro conocimiento de la
realidad» y de la relación lógica esperable cuando se encadenan, en cierto
orden, los eventos expresados.

\begin{ejercicio}\label{ej:12-copulativas}
Lea las siguientes oraciones, adaptadas del artículo «\textit{Homo sapiens}, el
animal que todo lo extingue»,\footnote{P. Esteban, «\textit{Homo sapiens}, el
animal que todo lo extingue», \textit{Página/12}, 12/10/2024
(\url{https://www.pagina12.com.ar/774085-homo-sapiens-el-animal-que-todo-lo-extingue}).}
y luego resuelva lo pedido, según se muestra en el ejemplo resuelto~1.

\begin{consignas}
  \item Encierre entre corchetes las coordinaciones copulativas e indique, en cada
        caso, la función sintáctica del SConj y la naturaleza categorial de los
        elementos coordinados.
  \item Distinga entre coordinación simétrica y asimétrica, cambiando el orden de
        los coordinandos. En caso de que haya asimetría, identifique el valor
        semántico asociado a la conjunción.
\end{consignas}

\begin{exe}
  \ex Las especies exóticas introducidas por los humanos suelen ser
      determinantes para su eliminación porque, como están en islas, no pueden
      escapar ni pueden asentarse con éxito en otras latitudes.
  \ex Meteoritos y volcanes, en el pasado, provocaron las extinciones masivas
      documentadas a lo largo de la historia.
  \ex El dodo vivió en la Isla Mauricio y se extinguió hacia el siglo XVII.
  \ex La compañía aislará y cultivará óvulos y células productoras de esperma
      para luego editar las secuencias de ADN que coincidan con las del dodo.
  \ex Estas células darían lugar a un animal mezcla de dodo y paloma.
  \ex El mamut lanudo habitó Asia y Europa hace 6 mil años, y podría volver a
      existir gracias a la ciencia.
\end{exe}

\begin{resuelto}
1. a) Las especies exóticas introducidas por los humanos suelen ser
determinantes para su eliminación porque, como están en islas,
[\etq{SConj/Pred} [\etq{SV/Esp} no pueden escapar] ni\etq{Conj/n}
[\etq{SV/Comp} pueden asentarse con éxito en otras latitudes]]. b) La
coordinación es asimétrica. No es posible cambiar el orden de los coordinandos,
debido a la presencia de un valor temporal de sucesión asociado a la
conjunción.\pendiente{La inversión (\textit{no pueden asentarse\ldots\ ni pueden
escapar}) parece aceptable: revisar el carácter asimétrico}
\end{resuelto}
\end{ejercicio}

\subsubsection{La coordinación disyuntiva}\label{sec:12-disyuntiva}

La coordinación disyuntiva se expresa, principalmente, a partir del uso de la
conjunción simple \forma{o}, que puede unir dos o más elementos que se
interpretan como \sig{opciones}. Además, se pueden utilizar conjunciones
discontinuas: \forma{o\ldots\ o\ldots, bien\ldots\ bien\ldots, ya\ldots\
ya\ldots, ora\ldots\ ora\ldots, sea\ldots\ sea\ldots, fuera\ldots\ fuera\ldots}
Excepto \forma{o\ldots\ o}, las restantes, debido a su alto grado de formalidad,
en la actualidad tienen uso restringido y se emplean sobre todo en la lengua
escrita.

Tomando, entonces, como ejemplo típico el de la conjunción \forma{o}, podemos
señalar que se establece un primer matiz en la interpretación de la disyunción en
función de si la elección entre las opciones es obligatoria o no. En el primer
caso, hablamos de una \termino{interpretación exclusiva}, como en
\exref{ex:12-disy-a} y \exref{ex:12-disy-b}. En cambio, cuando la elección no se
impone necesariamente, hablamos de una \termino{interpretación inclusiva}, en la
que las opciones presentadas se interpretan como instancias posibles, como en
\exref{ex:12-disy-c} y \exref{ex:12-disy-d}:

\begin{exe}
  \ex
    \begin{xlist}
      \ex\label{ex:12-disy-a} La azafata preguntó a los pasajeros si querían
          carne \textit{o} pasta.
      \ex\label{ex:12-disy-b} ¿Compramos el pasaje ahora \textit{o} esperamos a
          que baje el precio?
      \ex\label{ex:12-disy-c} Podemos ir al cine mañana, el sábado \textit{o} el
          domingo.
      \ex\label{ex:12-disy-d} Mañana podemos ir al cine, \textit{o} visitar unos
          amigos \textit{o} preparar una rica cena.
    \end{xlist}
\end{exe}

\begin{recuadro}[Para ampliar]
Con respecto a la interpretación inclusiva y la exclusiva de la conjunción
\forma{o} del español, \textcite[195]{jimenezjulia2023} establece un vínculo con
los valores expresados en latín por las conjunciones \forma{aut} y \forma{uel}.
De acuerdo con lo señalado por este autor, en latín \forma{aut} expresaba el
valor exclusivo y \forma{uel} (o \forma{siue / seu}), el inclusivo. El español
heredó la conjunción \forma{o} de la primera, pero no conservó \forma{uel}. Ello,
sin embargo, no implicó la pérdida de los dos valores: se produjo una
reconversión, de modo que la conjunción \forma{o} pasó a ser la general para
ambos valores. Asimismo, Jiménez Juliá añade que la repetición de la conjunción,
que la convierte en correlativa (\forma{o X o Y}), solo expresa valor exclusivo.
\end{recuadro}

En algunos casos, la conjunción \forma{o} tiene otras interpretaciones. Por un
lado, puede usarse con un sentido denominativo para introducir una aclaración
del valor de una expresión. Como vemos en los ejemplos de
\exref{ex:12-carpincho}, en estas instancias se trata de un uso metalingüístico:

\begin{exe}
  \ex\label{ex:12-carpincho}
    \begin{xlist}
      \ex El carpincho \textit{o} capibara es el roedor más grande del mundo.
      \ex Pekín \textit{o} Beijing es la capital de la República Popular China.
    \end{xlist}
\end{exe}

Por otro lado, la disyunción puede usarse con un valor correctivo, con el
objetivo de introducir una reformulación que, a su vez, incorpora algún matiz
respecto del primer coordinando:

\begin{exe}
  \ex
    \begin{xlist}
      \ex Nueva Constitución, \textit{o} el cuento de la inversión extranjera.
          (\textit{Página/12}, 2/9/2024)
      \ex Tu obra fue muy corta \textit{o}, mejor dicho, el tiempo se pasó
          volando.
    \end{xlist}
\end{exe}

Además, la conjunción \forma{o} puede usarse de forma discontinua o
distributiva, es decir, antepuesta a ambos elementos de la coordinación:

\begin{exe}
  \ex
    \begin{xlist}
      \ex Podemos ir al concierto \textit{o} el jueves \textit{o} el viernes
          (pero no los dos días).
      \ex \textit{O} nos unimos al mundo globalizado \textit{o} quedamos
          aislados.
    \end{xlist}
\end{exe}

Si bien hay otras conjunciones disyuntivas discontinuas, como \forma{ya\ldots\
ya\ldots, bien\ldots\ bien\ldots, ora\ldots\ ora\ldots}, su frecuencia de uso es
mucho menor. Respecto de estas conjunciones correlativas,
\textcite[196]{jimenezjulia2023} señala que no son un tipo específico de
conjunciones, sino que se les da ese nombre a estructuras coordinadas
copulativas o disyuntivas expresadas con nexos correlativos. Como señala este
autor, las conjunciones correlativas pueden tener un valor disyuntivo, como en
\exref{ex:12-corr-a}, o copulativo, como en \exref{ex:12-corr-b}:

\begin{exe}
  \ex
    \begin{xlist}
      \ex\label{ex:12-corr-a} Todos sabemos que \textit{prefijo} es lo que
          \sig{figura o se coloca delante} de un número o una palabra, ya sea el
          \textit{34} para llamar a España desde el extranjero, ya sea el
          elemento compositivo \textit{auto-}. (Fundéu RAE)
      \ex\label{ex:12-corr-b} La perfección absoluta, bien en las artes, bien en
          las demás cosas, no es de aquellas que el cielo quiso conceder a la
          tierra. (M.~Menéndez Pelayo, \textit{Historia})
    \end{xlist}
\end{exe}

Como puede observarse en \exref{ex:12-corr-b}, cuando el valor es copulativo,
las conjunciones correlativas son equivalentes a la expresión
\forma{tanto\ldots\ como\ldots}. En caso de que su valor sea disyuntivo, solo se
establece la interpretación de una disyunción exclusiva, es decir, se impone
necesariamente una decisión por uno de los dos elementos coordinados.

\begin{ejercicio}\label{ej:12-disyuntivas}
Lea las siguientes oraciones y luego resuelva lo pedido según se muestra en el
ejemplo resuelto~1.

\begin{consignas}
  \item Encierre entre corchetes las coordinaciones disyuntivas y determine la
        función sintáctica del SConj y la naturaleza categorial de los
        coordinandos.
  \item Identifique el valor con el que se interpreta la coordinación.
\end{consignas}

\begin{exe}
  \ex El carpincho o capibara es el roedor más grande del mundo.
  \ex Los capibaras son dóciles y usualmente los humanos los acarician o les dan
      de comer.
  \ex Dependiendo de la subespecie, el periodo de gestación dura unos 110 días o
      unos 150 días.
  \ex Sus depredadores naturales más importantes son félidos, como jaguares o
      pumas.
  \ex En la lengua tupí, el término es \textit{ka'apiûara}, que significa \sig{el
      que come hojas delgadas}, o \sig{comedor de hierba}.
  \ex En 2020, los carpinchos invadieron Nordelta o, mejor dicho,\footnote{Como
      se verá en §\ref{sec:12-marcadores}, \forma{mejor dicho} se considera un
      conector discursivo, que se combina con la conjunción, pero no es parte de
      ella.} volvieron a su territorio autóctono en los humedales del río
      Paraná.
\end{exe}

\begin{resuelto}
1. a) [\etq{SConj/Suj} [\etq{SN/Esp} El carpincho] o\etq{Conj/n}
[\etq{SN/Comp} capibara]] es el roedor más grande del mundo. b) Disyunción
metalingüística denominativa: introduce una segunda denominación del término
\forma{carpincho}.
\end{resuelto}
\end{ejercicio}

\subsubsection{La coordinación adversativa}\label{sec:12-adversativa}

En la coordinación adversativa se establece entre los coordinandos una relación
de \sig{oposición}. Las conjunciones adversativas principales son \forma{pero} y
\forma{sino}. También es adversativa la conjunción
\forma{mas},\footnote{\textcite[197]{jimenezjulia2023} comenta que \forma{mas}
proviene del latín \forma{magis}, origen de la conjunción adversativa dominante
en lenguas románicas como el francés (\forma{mais}) o el portugués
(\forma{mas}); en cambio, tanto en el español europeo como en las variedades
americanas parece relegada a registros cultos.} de características similares a
\forma{pero}, aunque su uso está más restringido por motivos de registro. Dado
que las dos conjunciones principales tienen características formales bien
diferenciadas, las presentaremos por separado.

\paragraph{La conjunción \textit{pero}.}\label{sec:12-pero}
Se puede utilizar para coordinar sintagmas de distinto tipo, aunque es más
frecuentemente utilizada para coordinar oraciones:

\begin{exe}
  \ex
    \begin{xlist}
      \ex\label{ex:12-depto-a} El departamento es [\etq{SA} chico], \textit{pero}
          [\etq{SA} luminoso].
      \ex El departamento tiene [\etq{SN} muchas puertas], \textit{pero}
          [\etq{SN} pocas ventanas].
      \ex\label{ex:12-depto-c} {}[\etq{SF/O} El departamento está a buen precio],
          \textit{pero} [\etq{SF/O} aún no apareció ningún comprador].
    \end{xlist}
\end{exe}

Con respecto a los valores de la relación adversativa establecida con la
conjunción \forma{pero}, \textcite[284]{ditullio2014} señala la existencia, por
un lado, de un valor de \termino{contraexpectativa} y, por el otro, de una
\termino{oposición valorizadora}. En el primer caso, se establece una relación de
contraposición entre lo expresado por ambos coordinandos, pero a partir de la
interpretación de una inferencia del hablante. Por ejemplo, en
\exref{ex:12-depto-c}, de lo dicho por el primer coordinando (\forma{el
departamento está a buen precio}) se podría inferir que se debería vender
rápido, pero esa inferencia se contrapone con lo manifestado por el segundo
coordinando (\forma{aún no apareció ningún comprador}), en el que se afirma un
evento o un estado de cosas contrario al inferido. En palabras de Di Tullio, la
conjunción \forma{pero} cancela la expectativa generada por el primer
coordinando.

Con respecto a la lectura valorizadora, la autora señala que tal interpretación
se produce cuando los elementos coordinados se pueden relacionar con una
jerarquía, de modo que la oposición permite priorizar lo expresado en el segundo
coordinando. Por ejemplo, en \exref{ex:12-depto-a} se contraponen dos
características de un departamento: \forma{ser chico} y \forma{ser luminoso}.
Pero, por el orden dado a los elementos, se interpreta una valorización de la
segunda característica y, por ende, una visión positiva del departamento. Si
alteramos el orden de los coordinandos (\forma{El departamento es luminoso, pero
chico}), la característica ahora valorizada es \forma{ser chico}, la cual suele
asociarse con un rasgo negativo, por lo que la oración se interpreta como un
juicio negativo sobre la propiedad.

En ambos casos, se destaca el hecho de que la coordinación adversativa se vincula
necesariamente con las inferencias que el hablante puede establecer a partir
tanto de su conocimiento del mundo como del contexto discursivo. Es en este
sentido que, como se afirma en la \textit{NGLE} \parencite[§\,31.10i]{rae2025},
la coordinación adversativa obliga al oyente a establecer relaciones de causa y
efecto entre oraciones mediante mecanismos discursivos que no se basan
únicamente en el contenido de las palabras. A modo de ejemplo, veamos cómo
funciona la coordinación adversativa en las siguientes oraciones de \textit{Éste
es el mar}, de Mariana Enriquez:

\begin{exe}
  \ex
    \begin{xlist}
      \ex\label{ex:12-helena} Helena las miraba atenta y tímida, \textit{pero}
          finalmente se atrevió a hablar.
      \ex\label{ex:12-humanos} Cuando los humanos investigaban los Misterios de
          los Últimos Días, notaban la rareza, \textit{pero} no comprendían su
          enormidad.
      \ex\label{ex:12-enjambre} El Enjambre era vasto, eran muchas sus
          integrantes, \textit{pero} no podían cambiar de rostro ni de cuerpo.
    \end{xlist}
\end{exe}

En \exref{ex:12-helena} se muestra el valor de contraexpectativa, dado que el
segundo coordinando (\forma{finalmente se atrevió a hablar}) cancela la
expectativa que genera el primero, a partir de la inferencia de que la timidez
no iba a dejar que hable. De manera similar, en \exref{ex:12-humanos}, a partir
del primer coordinando se puede inferir que los humanos, dado que pueden notar la
rareza, podrían encontrarle una explicación; sin embargo, tal expectativa se ve
cancelada por el contenido del segundo coordinando. Finalmente, en
\exref{ex:12-enjambre}, predomina una interpretación valorizadora, dado que los
coordinandos expresan dos características de las integrantes del Enjambre: ser
numerosas y no poder cambiar de rostro ni de cuerpo. Dado el orden elegido, en
este fragmento la evaluación se interpreta como negativa, en tanto se realza
aquello que las integrantes no pueden hacer. Dicha valoración podría volverse
positiva con simplemente alterar el orden de los coordinandos: \forma{No podían
cambiar de rostro ni de cuerpo, pero el Enjambre era vasto, eran muchas sus
integrantes}.

\paragraph{La conjunción \textit{sino}.}\label{sec:12-sino}
Esta conjunción, al igual que \forma{pero}, establece una coordinación binaria
opositiva, con la particularidad de que exige la presencia de una negación en el
primer coordinando y el segundo coordinando se interpreta como el reemplazo del
elemento negado, por lo que la coordinación constituye una \termino{oposición
exclusiva}. La conjunción \forma{sino} puede coordinar distintos tipos de
sintagmas, desde sintagmas léxicos hasta oraciones independientes o
subordinadas:

\begin{exe}
  \ex\label{ex:12-sino}
    \begin{xlist}
      \ex Juan llegó \textit{no} a la mañana, \textit{sino} a la tarde.
      \ex Pedro \textit{no} tiene dos perros, \textit{sino} dos gatos.
      \ex El animal que trajiste \textit{no} es macho, \textit{sino} hembra.
      \ex Ana \textit{no} dijo que compró el libro, \textit{sino} que lo había
          visto a la venta.
    \end{xlist}
\end{exe}

Cabe destacar que, en muchas ocasiones, la negación del primer término exigida
por \forma{sino} se usa de forma combinada con el adverbio \forma{solo},
correlativo del adverbio \forma{también}, lo cual cambia la interpretación de
oposición exclusiva, de modo que el segundo término de la coordinación ya no se
interpreta como un reemplazo, sino que adquiere valor de adición. Para que tal
interpretación sea posible, es fundamental que haya compatibilidad entre los
predicados de cada coordinando. Como vemos, si agregamos \forma{solo} y
\forma{también} a los ejemplos de \exref{ex:12-sino}, podemos obtener secuencias
gramaticales, como \exref{ex:12-solo-b}, o dudosas, como \exref{ex:12-solo-d},
ambas con cambio de significado: en \exref{ex:12-solo-b}, Pedro pasa a tener
cuatro mascotas en vez de dos y, en \exref{ex:12-solo-d}, se pierde el sentido
exclusivo de la adversativa. Por último, en \exref{ex:12-solo-a} y
\exref{ex:12-solo-c} las secuencias que se obtienen son incoherentes, dado que
los significados de los elementos léxicos coordinados ---\forma{mañana / tarde} y
\forma{macho / hembra}--- son incompatibles y no pueden añadirse.

\begin{exe}
  \ex
    \begin{xlist}
      \ex[\#]{\label{ex:12-solo-a} Juan llegó \textit{no solo} a la mañana,
          \textit{sino también} a la tarde.}
      \ex\label{ex:12-solo-b} Pedro \textit{no solo} tiene dos perros,
          \textit{sino también} dos gatos.
      \ex[\#]{\label{ex:12-solo-c} El animal que trajiste \textit{no solo} es
          macho, \textit{sino también} hembra.}
      \ex[?]{\label{ex:12-solo-d} Ana \textit{no solo} dijo que compró el libro,
          \textit{sino también} que lo había visto a la venta.}
    \end{xlist}
\end{exe}

\begin{ejercicio}\label{ej:12-adversativas}
Lea las siguientes oraciones, adaptadas del artículo «\textit{Homo sapiens}, el
animal que todo lo extingue», y luego resuelva lo pedido según se muestra en el
ejemplo resuelto~1.

\begin{consignas}
  \item Encierre entre corchetes las coordinaciones adversativas y determine la
        función sintáctica del SConj y la naturaleza categorial de los
        coordinandos.
  \item Analice el valor con el que se interpreta la coordinación.
\end{consignas}

\begin{exe}
  \ex Las sociedades no solo barren con el número de especies, sino que arrancan
      de raíz ecosistemas enteros.
  \ex Estos números representan una pérdida ya consolidada del 5 por ciento de
      los linajes, es decir, de especies distintas, pero con un ancestro en
      común.
  \ex No solo se achicará la biodiversidad de la fauna, sino que también se
      perderán funciones ecosistémicas cruciales, como la polinización y el
      control de plagas.
  \ex La paloma migratoria era la especie de ave más abundante en Norteamérica,
      pero ahora está completamente extinta.
  \ex El ser humano no se siente parte del ecosistema, sino que está totalmente
      desconectado de la naturaleza de la que forma parte.
  \ex Hay medidas para frenar este conflicto, pero se tienen que tomar muy
      rápido.
\end{exe}

\begin{resuelto}
1. a) [\etq{SConj/OC} [\etq{SF/Esp} Las sociedades no solo barren con el número
de especies], sino que\etq{Conj/n} [\etq{SF/Comp} arrancan de raíz ecosistemas
enteros]]. b) La coordinación adversativa con \forma{sino} no se interpreta como
oposición exclusiva, sino como adición contrastiva, por la presencia del
adverbio \forma{solo} en el primer coordinando.
\end{resuelto}
\end{ejercicio}

\subsection{Diferencias entre conjunciones y marcadores discursivos}\label{sec:12-marcadores}

Las conjunciones coordinantes a menudo se utilizan combinadas con elementos como
\forma{además, sin embargo, encima}. Si bien estos elementos hacen su aporte para
establecer la conexión entre los constituyentes, no pertenecen a la clase de las
conjunciones coordinantes, sino que se los conoce como \termino{marcadores} o
\termino{conectores discursivos}.

En el \textit{Glosario de términos gramaticales}, se define el marcador
discursivo como una «unidad léxica, formada por una o varias palabras, que
establece vínculos lógicos, discursivos o argumentativos entre los segmentos de
un texto» \parencite[74]{bosque2019}. La mayor parte de los marcadores
discursivos son locuciones adverbiales (\forma{sin embargo, por si fuera poco,
no obstante}), pero dentro del paradigma también encontramos adverbios
(\forma{encima, además}), interjecciones (\forma{claro, bueno}) y locuciones
conjuntivas (\forma{ya que, puesto que}). Como se explica en la \textit{NGLE}
\parencite[§\,30.12a]{rae2025}, el concepto de conector o marcador discursivo
«no constituye una clase sintáctica de palabras análoga a \textit{verbo,
conjunción} o \textit{adverbio}, sino un grupo establecido con criterios
textuales».

Para dar cuenta de la necesidad de diferenciar las conjunciones de los conectores
discursivos, debemos tener en cuenta que las conjunciones tienen una posición fija
entre los coordinandos, como se observa en \exref{ex:12-fija}, mientras que los
conectores pueden desplazarse, tal como se muestra en \exref{ex:12-movil}:

\begin{exe}
  \ex\label{ex:12-fija}
    \begin{xlist}
      \ex Juan \textit{y} Pedro vendrán mañana.
      \ex[*]{Juan Pedro \textit{y} vendrán mañana.}
    \end{xlist}
  \ex\label{ex:12-movil}
    \begin{xlist}
      \ex Juan salió temprano y, \textit{sin embargo}, no llegó a tiempo.
      \ex Juan salió temprano y no llegó a tiempo, \textit{sin embargo}.
    \end{xlist}
\end{exe}

Una segunda diferencia importante entre marcadores y conjunciones radica en que,
cuando están al inicio de la oración, los primeros van seguidos de pausa,
mientras que las conjunciones generalmente forman una unidad entonacional con el
constituyente al que preceden. En las oraciones de \exref{ex:12-pausa} puede
apreciarse el contraste:

\begin{exe}
  \ex\label{ex:12-pausa}
    \begin{xlist}
      \ex \textit{Además}, Juan todavía no llegó.
      \ex[*]{\textit{Pero}, Juan todavía no llegó.}
      \ex \textit{Pero} Juan todavía no llegó.
    \end{xlist}
\end{exe}

Otro aspecto que muestra la necesidad de diferenciar entre conjunciones
coordinantes y marcadores discursivos es la posibilidad de combinación: al
asociar dos conjunciones coordinantes se obtiene una secuencia agramatical, como
se muestra en \exref{ex:12-comb-a}, mientras que la combinación de una conjunción
y un conector discursivo produce secuencias que sí son gramaticales, como la de
\exref{ex:12-comb-b}:

\begin{exe}
  \ex
    \begin{xlist}
      \ex[*]{\label{ex:12-comb-a} Juan todavía no llegó, \textit{pero y} podemos
          empezar.}
      \ex\label{ex:12-comb-b} Juan todavía no llegó, \textit{pero, sin embargo},
          podemos empezar.
    \end{xlist}
\end{exe}

En definitiva, los marcadores discursivos sirven para orientar «la manera en que
la oración o el fragmento oracional sobre el que inciden han de ser
interpretados» en relación con el contexto precedente o con las inferencias que
de él se obtienen \parencite[§\,30.12i]{rae2025}, pero no tienen las mismas
propiedades sintácticas que las conjunciones. De hecho, un contraste adicional
puede encontrarse en el inventario de ítems de cada grupo. Como ya vimos, las
conjunciones coordinantes, si bien pueden presentar distintos valores
(copulativo, disyuntivo, adversativo), constituyen un inventario bastante
acotado: \forma{y, ni, o, pero, sino} (y las variantes correlativas). En cambio,
existe un amplio repertorio de marcadores discursivos, tan abundante que
normalmente se presenta en subconjuntos, agrupados por su función discursiva. En
la tabla~\ref{tab:12-marcadores} presentamos la propuesta de clasificación de
los grupos fundamentales, tomada de la \textit{NGLE}
\parencite[§\,30.13a]{rae2025}.

{\small
\begin{longtable}{@{}>{\raggedright\arraybackslash}p{38mm}
    >{\raggedright\arraybackslash\itshape}p{85mm}@{}}
  \caption{Marcadores discursivos adverbiales}\label{tab:12-marcadores}\\
  \toprule
  Grupo & \normalfont Marcadores \\
  \midrule
  \endfirsthead
  \toprule
  Grupo & \normalfont Marcadores \\
  \midrule
  \endhead
  1. Aditivos y de precisión o particularizadores & a decir verdad, además,
    análogamente, aparte, asimismo, de hecho, encima, en el fondo, en realidad,
    es más, por añadidura, por otro lado, por si fuera poco, sobre todo \\
  2. Adversativos y contraargumentativos & ahora bien, al contrario, antes bien,
    después de todo, empero, en cambio, eso sí, no obstante, por el contrario,
    sin embargo, todo lo contrario \\
  3. Concesivos & así y todo, aun así, con todo, de cualquier manera, de todas
    formas, de todas maneras, de todos modos, en cualquier caso \\
  4. Consecutivos e ilativos & así pues, consiguientemente, de ese modo, en
    consecuencia, entonces, por consiguiente, por ende, por lo tanto, por tanto,
    pues \\
  5. Explicativos & a saber, es decir, esto es, o sea \\
  6. Reformuladores & dicho con otras palabras, dicho en otros términos, dicho de
    otra forma, dicho de otra manera, de otro modo, más claramente, más
    llanamente, hablando claro \\
  7. Ejemplificativos & así, así por ejemplo, así tenemos, por ejemplo,
    verbigracia \\
  8. Rectificativos & más bien, mejor dicho, por mejor decir \\
  9. Recapitulativos & a fin de cuentas, al fin y al cabo, en conclusión, en
    definitiva, en fin, en resumen, en resumidas cuentas, en síntesis, en suma,
    en una palabra, resumiendo, total \\
  10. Ordenadores & a continuación, antes que nada, de una parte, de otra parte,
    en primer lugar, en primer término, finalmente, para empezar, para
    terminar \\
  11. De apoyo argumentativo & así las cosas, dicho esto, en vista de ello, pues
    bien \\
  12. De digresión & a propósito, a todo esto, dicho sea de paso, entre
    paréntesis, por cierto \\
  \bottomrule
\end{longtable}}

\subsection{Coordinación y elipsis}\label{sec:12-elipsis}

Un fenómeno estrechamente vinculado con la coordinación de sintagmas y oraciones
es la \termino{elipsis}, es decir, la posibilidad de eliminar elementos del
segundo coordinando, en función de que pueden interpretarse a partir del primer
coordinando. En algunos casos, la elipsis es fácilmente detectable: cuando se
omite el verbo del segundo coordinando, como en \exref{ex:12-elip-a}, donde se
puede observar la presencia de una coma que marca la elisión; cuando se elide el
verbo y algún complemento, como en \exref{ex:12-elip-b}, o cuando en el segundo
coordinando se emplean los adverbios \forma{también, tampoco, sí} o \forma{no},
que dan lugar a la elipsis de todo el sintagma verbal, como en
\exref{ex:12-elip-c} y \exref{ex:12-elip-d}:\pendiente{En
\exref{ex:12-elip-b} se elide el verbo con su COD, pero el remanente es un
adjunto}

\begin{exe}
  \ex\label{ex:12-elip}
    \begin{xlist}
      \ex\label{ex:12-elip-a} Juan estudia Letras \textit{y} Pedro, Filosofía.
      \ex\label{ex:12-elip-b} Ana leyó el libro en su casa \textit{y} María, en la
          plaza.
      \ex\label{ex:12-elip-c} Juan no estudia Filosofía, \textit{pero} Pedro sí.
      \ex\label{ex:12-elip-d} Ana no leyó el libro en la biblioteca, \textit{y}
          María tampoco.
    \end{xlist}
\end{exe}

Siguiendo a \textcite[278]{ditullio2014}, diferenciamos las instancias de
coordinación con elipsis, como las de \exref{ex:12-elip}, de otros casos en los
que no es necesario plantear una elisión, sino que se puede establecer una
coordinación a nivel sintagmático, como en los ejemplos de
\exref{ex:12-sinelip}:

\begin{exe}
  \ex\label{ex:12-sinelip}
    \begin{xlist}
      \ex\label{ex:12-sinelip-a} Pedro estudia Letras \textit{y} Filosofía.
      \ex\label{ex:12-sinelip-b} Ana leyó el libro en su casa \textit{y} en la
          plaza.
      \ex\label{ex:12-sinelip-c} Ana \textit{y} María son compañeras de trabajo.
    \end{xlist}
\end{exe}

En \exref{ex:12-sinelip-a}, se coordinan dos sintagmas nominales dentro del COD.
En \exref{ex:12-sinelip-b}, los que se coordinan son dos sintagmas
preposicionales, que funcionan como adjuntos dentro del predicado. Por último, en
\exref{ex:12-sinelip-c} se coordinan dos sintagmas nominales, pero dentro del
sujeto y, en este caso en particular, no sería admisible pensar una estructura
alternativa con elipsis, dado que \forma{ser compañeras}, al ser un predicado
simétrico, requiere un sujeto plural. En este sentido, si bien existiría la
posibilidad de plantear versiones completamente redundantes para los casos de
\exref{ex:12-sinelip-a} y \exref{ex:12-sinelip-b}, no existe tal opción para
\exref{ex:12-sinelip-c},\footnote{Si bien es cierto que
\exref{ex:12-sinelip-c} es aceptable si se interpreta que Ana y María son
compañeras de trabajo de una tercera persona, esa lectura no sería simétrica.}
tal como se puede apreciar en los siguientes ejemplos:

\begin{exe}
  \ex
    \begin{xlist}
      \ex Pedro estudia Letras \textit{y (también)} estudia Filosofía.
      \ex Ana leyó el libro en su casa \textit{y (también)} lo leyó en la plaza.
      \ex[*]{Ana es compañera de trabajo \textit{y} María es compañera de
          trabajo.}
    \end{xlist}
\end{exe}

\begin{ejercicio}\label{ej:12-elipsis}
Identifique los elementos que se eliden en el segundo coordinando de las
siguientes oraciones, según se muestra en el ejemplo resuelto~1.

\begin{exe}
  \ex Frodo quiere destruir el anillo y Sam también.
  \ex Aragorn está enamorado de Arwen y Sam, de Rosita.
  \ex Frodo trata bien a Gollum, pero Sam no.
  \ex Boromir no puede ser rey de Gondor, y Faramir tampoco.
  \ex Gandalf no es un mago ambicioso, pero Saruman sí.
  \ex Los elfos desconfían de los hombres y de los enanos.
\end{exe}

\begin{resuelto}
1. Se elide \forma{quiere destruir el anillo}.
\end{resuelto}
\end{ejercicio}

\begin{ejercicio}[Integración]\label{ej:12-meta}
Lea el siguiente fragmento de la nota periodística «¿Puede Meta usar tus fotos
para entrenar su IA?» y resuelva lo pedido respecto de las coordinaciones
señaladas entre corchetes, según se muestra en el ejemplo resuelto~1.

\begin{consignas}
  \item Identifique el tipo de conjunción coordinante.
  \item Analice la estructura del SConj.
  \item Determine si se trata de una coordinación homocategorial o
        heterocategorial.
  \item Considere si hay valores asociados a la conjunción que inciden en la
        interpretación de la coordinación.
\end{consignas}

\begin{textolectura}
(1) La última operación de Meta ha suscitado cierta preocupación entre los
usuarios de [Facebook e Instagram]. La empresa ha anunciado recientemente que
utilizará los datos que estas personas han colgado en sus plataformas para
entrenar su nueva inteligencia artificial (IA), a no ser que expliciten lo
contrario. (2) Un movimiento que ha hecho emerger todo tipo de preguntas sobre
[la protección de datos o su impacto social]. (3) A su vez, ha abierto una
oportunidad para reflexionar sobre la falta de consciencia que tienen los
usuarios [sobre el uso de sus datos en general, pero también sobre las maniobras
normalizadas y ocultas del poder]. [\ldots] (4) (A) [Desde Meta se ha afirmado
que buscan cumplir con (B) [las normativas de privacidad y protección de datos],
pero la desconfianza persiste]. (5) (A) [A su juicio, el problema no gira tanto
en torno a si de aquí en adelante la foto de tu hijo pequeño aparecerá en un
futuro en anuncios de pañales, sino que la cuestión se sitúa en torno a (B) [si
basta con que Meta dé la oportunidad a los usuarios de oponerse a que sus datos
sean utilizados con ese fin o es necesario un consentimiento específico]].
\end{textolectura}
\pendiente{Falta la referencia de la nota (medio, fecha, URL)}

\begin{resuelto}
1. a) Conjunción coordinante copulativa. b) [\etq{SConj/Comp} [\etq{SN/Esp}
Facebook] \textbf{e}\etq{Conj/n} [\etq{SN/Comp} Instagram]]. c) Coordinación
homocategorial entre dos SSNN. d) No hay valores asociados a la conjunción, más
allá del aditivo.
\end{resuelto}
\end{ejercicio}

\begin{recuadro}[Reflexión final]
Como hemos visto en este capítulo, la coordinación se basa en un mecanismo
sumamente productivo, que permite conformar una unidad sintagmática de distintos
grados de complejidad, a partir de la incorporación de una conjunción
coordinante, que se desempeña como núcleo funcional del nuevo sintagma. Mientras
que en la coordinación de sintagmas léxicos (SV, SN, SA, SAdv, SP) el fenómeno se
inscribe dentro de la sintaxis de la oración simple, la coordinación de oraciones
nos lleva a la esfera de las relaciones interoracionales, campo que comparte con
la subordinación.
\end{recuadro}

% --- Secciones finales -------------------------------------------------
\lecturasrecomendadas

\solucionario

\begin{solucion}{\ref{ej:12-predicaciones}}
2. Predicación primaria: \forma{Los transeúntes miraban la escena}. Predicación
secundaria: \forma{Los transeúntes estaban asombrados}. Yuxtaposición:
\forma{Los transeúntes miraban la escena; estaban asombrados}. Coordinación:
\forma{Los transeúntes miraban la escena y estaban asombrados}. Subordinación:
\forma{Los transeúntes, que miraban la escena, estaban asombrados}. 3.
Predicación primaria: \forma{María colgó el cuadro}. Predicación secundaria:
\forma{El cuadro quedó bien centrado}. Yuxtaposición: \forma{María colgó el
cuadro; el cuadro quedó bien centrado}. Coordinación: \forma{María colgó el
cuadro y el cuadro quedó bien centrado}. Subordinación: \forma{María colgó el
cuadro, que quedó bien centrado}.
\end{solucion}

\begin{solucion}{\ref{ej:12-desantis}}
2. Hicimos una colecta para pagarle el legado, sin saber si comprábamos
[\etq{SConj/COD} [\etq{SN/Esp} un tesoro] o\etq{Conj/n} [\etq{SN/Comp} un
montón de diarios viejos]]. 3. [\etq{SConj/OC} [\etq{SF/Esp} Habíamos subido el
baúl por la escalera casi a los saltos], pero\etq{Conj/n} [\etq{SF/Comp} la
visión de aquellos papeles nos paralizó]]. 4. [\etq{SConj/OC} [\etq{SF/Esp}
Balestri sigue siendo un desconocido], y\etq{Conj/n} [\etq{SF/Comp} de su torre
no poseemos más que la sombra de una sombra]].
\end{solucion}

\begin{solucion}{\ref{ej:12-enriquez}}
[\etq{Oración 1} Tenían una madre que rara vez era [\etq{SConj} [\etq{V/Esp}
vista] o\etq{Conj/n} [\etq{V/Comp} convocada]]; la madre Hécate].
[\etq{SConj/Oración 2} [\etq{SF/Esp} En el Enjambre se creía que era una
Leyenda, una mentira], pero\etq{Conj/n} [\etq{SF/Comp} las Luminosas se reían
de tanta ignorancia]]. [\etq{SConj/OS} Y\etq{Conj/n} [\etq{SF/Comp} tenían
hermanas mayores: [\etq{SConj/ApE} [\etq{SN/Esp} Perséfone bajo la tierra, con
flores en el pelo]; y\etq{Conj/n} [\etq{SN/Comp} las Imago, las que se escondían
en lugares [\etq{SConj/M} [\etq{SA/Esp} secos] y\etq{Conj/n} [\etq{SA/Comp}
oscuros]]]]]].\pendiente{La oración que empieza con \textit{Y} se etiqueta como
SConj sin especificador expreso: conviene explicitar el análisis}
\end{solucion}

\begin{solucion}{\ref{ej:12-lecturas}}
2. a) Lectura distributiva: \juicio{*}{\forma{Frodo se lleva bien desde la
infancia y Sam también se lleva bien desde la infancia}}. b) Lectura colectiva:
\forma{Frodo y Sam se llevan bien entre sí desde la infancia}. c) Lectura
ordenada: \juicio{*}{\forma{Frodo y Sam se llevan bien desde la infancia,
respectivamente}}.

3. a) Lectura distributiva: \forma{Frodo viaja a Mordor y Sam también viaja a
Mordor} / \forma{Tanto Frodo como Sam viajan a Mordor (separados)}. b) Lectura
colectiva: \forma{Frodo y Sam viajan a Mordor juntos}. c) Lectura ordenada:
\juicio{*}{\forma{Frodo y Sam viajan a Mordor, respectivamente}}. No es posible
la lectura porque no hay un segundo elemento en el predicado que permita
realizar el ordenamiento.

4. a) Lectura distributiva: se puede armar la secuencia por separado, pero no
mantiene el mismo sentido porque las predicaciones se vuelven reflexivas:
\forma{Sam se odia (a sí mismo) y Gollum también se odia (a sí mismo)} /
\forma{Tanto Sam como Gollum se odian (a sí mismos)}. Al tratarse de un verbo
recíproco, se podría reponer la lectura distributiva si se eliminara el pronombre
\forma{se} y se asignaran las funciones de sujeto y objeto directo
respectivamente a Sam y a Gollum: \forma{Sam odia a Gollum y Gollum odia a Sam}.
b) Lectura colectiva: \forma{Sam y Gollum se odian entre sí / mutuamente}. Es la
única posible, por tratarse de un predicado simétrico, de naturaleza colectiva,
que se predica sobre un grupo. c) Lectura ordenada: no es posible, porque en el
predicado no hay elementos que puedan ordenarse.

5. a) Lectura distributiva: es posible estructuralmente: \forma{Frodo viaja hacia
Mordor y hacia el Bosque de Fangorn, y Aragorn también} / \forma{Tanto Frodo como
Aragorn viajan hacia Mordor y hacia el Bosque de Fangorn}. Podría admitirse tal
lectura si consideramos que primero van a un lugar y después al otro, cada uno
por su cuenta (aunque tal sentido no se corresponde con lo que ocurre en el
relato). b) Lectura colectiva: \forma{Frodo y Aragorn viajan juntos hacia Mordor
y el Bosque de Fangorn}. Es posible si consideramos que van juntos primero a un
lugar y después al otro (aunque tal sentido no se corresponde con lo que ocurre
en el relato). c) Lectura ordenada: \forma{Frodo y Aragorn viajan hacia Mordor y
hacia el Bosque de Fangorn, respectivamente} / \forma{Frodo viaja hacia Mordor y
Aragorn viaja hacia el Bosque de Fangorn}. Este es el sentido compatible con lo
que ocurre en el relato, donde Frodo viaja con Sam hacia el este, mientras que
Aragorn, con otra compañía, viaja hacia el oeste.
\end{solucion}

\begin{solucion}{\ref{ej:12-copulativas}}
2. a) [\etq{SConj/Suj} [\etq{SN/Esp} Meteoritos] y\etq{Conj/n} [\etq{SN/Comp}
volcanes]], en el pasado, provocaron las extinciones masivas documentadas a lo
largo de la historia. b) Coordinación simétrica: es posible cambiar el orden de
los coordinandos. 3. a) El dodo [\etq{SConj/Pred} [\etq{SV/Esp} vivió en la Isla
Mauricio] y\etq{Conj/n} [\etq{SV/Comp} se extinguió hacia el siglo XVII]]. b)
Coordinación asimétrica: se reconoce un valor temporal de sucesión asociado a la
conjunción.\pendiente{El original etiqueta los coordinandos como SF} 4. a) La
compañía [\etq{SConj1/Pred} [\etq{SV/Esp} aislará] y\etq{Conj/n} [\etq{SV/Comp}
cultivará [\etq{SConj2/COD} [\etq{SN/Esp} óvulos] y\etq{Conj/n} [\etq{SN/Comp}
células productoras de esperma]] para luego editar las secuencias de ADN que
coincidan con las del dodo]]. b) SConj1: coordinación asimétrica: se reconoce un
valor temporal de sucesión asociado a la conjunción. SConj2: coordinación
simétrica: es posible cambiar el orden de los coordinandos. 5. a) Estas células
darían lugar a un animal mezcla de [\etq{SConj/Comp} [\etq{SN/Esp} dodo]
y\etq{Conj/n} [\etq{SN/Comp} paloma]]. b) Coordinación simétrica: es posible
cambiar el orden de los coordinandos. 6. a) El mamut lanudo [\etq{SConj1/Pred}
[\etq{SV/Esp} habitó [\etq{SConj2/COD} [\etq{SN/Esp} Asia] y\etq{Conj/n}
[\etq{SN/Comp} Europa]] hace 6 mil años], y\etq{Conj/n} [\etq{SV/Comp} podría
volver a existir gracias a la ciencia]]. b) SConj1: coordinación asimétrica: se
reconoce un valor temporal de sucesión asociado a la conjunción. SConj2:
coordinación simétrica: es posible cambiar el orden de los
coordinandos.\pendiente{El original etiqueta SConj2 como CCL, \textit{Asia} como
SV/Esp y lo califica de «asimétrica: es posible cambiar el orden»}
\end{solucion}

\begin{solucion}{\ref{ej:12-disyuntivas}}
2. a) Los capibaras son dóciles y usualmente los humanos [\etq{SConj/Pred}
[\etq{SV/Esp} los acarician] o\etq{Conj/n} [\etq{SV/Comp} les dan de comer]].
b) Disyunción inclusiva: ambas opciones son posibles. 3. a) Dependiendo de la
subespecie, el periodo de gestación dura [\etq{SConj/CCT} [\etq{SN/Esp} unos 110
días] o\etq{Conj/n} [\etq{SN/Comp} unos 150 días]]. b) Disyunción exclusiva: la
duración del periodo de gestación es distinta para cada subespecie. 4. a) Sus
depredadores naturales más importantes son félidos, como
[\etq{SConj/Comp}\footnote{En (4), consideramos que el núcleo de \forma{como
jaguares o pumas} es \forma{como}, dado que introduce la secuencia que sirve de
ejemplo.} [\etq{SN/Esp} jaguares] o\etq{Conj/n} [\etq{SN/Comp} pumas]]. b)
Disyunción inclusiva: ambas opciones son posibles. 5. a) En la lengua tupí, el
término es \textit{ka'apiûara}, que significa [\etq{SConj/PSO} [\etq{SF/Esp}
\sig{el que come hojas delgadas}] o\etq{Conj/n} [\etq{SN/Comp} \sig{comedor de
hierba}]]. b) Disyunción metalingüística denominativa: introduce una segunda
denominación posible para \textit{ka'apiûara}. 6. a) En 2020, los carpinchos
[\etq{SConj/Pred} [\etq{SV/Esp} invadieron Nordelta] o\etq{Conj/n}, mejor dicho,
[\etq{SV/Comp} volvieron a su territorio autóctono en los humedales del río
Paraná]]. b) Disyunción con valor correctivo: introduce una reformulación que, a
su vez, incorpora la opinión del emisor, que no considera que los carpinchos
hayan invadido la zona, sino que volvieron porque era su territorio autóctono.
\end{solucion}

\begin{solucion}{\ref{ej:12-adversativas}}
2. a) Estos números representan una pérdida ya consolidada del 5 por ciento de
los linajes, es decir, de especies [\etq{SConj/M} [\etq{SA/Esp} distintas],
pero\etq{Conj/n} [\etq{SP/Comp} con un ancestro en común]]. b) La coordinación
adversativa con \forma{pero} establece una contraposición entre los coordinandos,
y el segundo explicita una contraexpectativa: aunque la afirmación del primer
coordinando de que las especies son distintas lleva a inferir que no tienen
ancestros en común, estas especies sí lo tienen (y por eso son del mismo linaje).
3. a) [\etq{SConj/OC} [\etq{SF/Esp} No solo se achicará la biodiversidad de la
fauna], sino que\etq{Conj/n} [\etq{SF/Comp} también se perderán funciones
ecosistémicas cruciales, como la polinización y el control de plagas]]. b) La
coordinación adversativa con \forma{sino} no se interpreta como una oposición
exclusiva, sino como una adición contrastiva, por la presencia del adverbio
\forma{solo} en el primer coordinando y de \forma{también} en el segundo. 4. a)
La paloma migratoria [\etq{SConj/Pred} [\etq{SV/Esp} era la especie de ave más
abundante en Norteamérica], pero\etq{Conj/n} [\etq{SV/Comp} ahora está
completamente extinta]]. b) La coordinación adversativa con \forma{pero} se
interpreta como valorizadora del segundo coordinando. Los coordinandos se
relacionan en una jerarquía de valores (antes era abundante, ahora está extinta),
y se prioriza el segundo. El orden de los coordinandos podría cambiarse
(\forma{La paloma migratoria ahora está completamente extinta, pero era la
especie más abundante en Norteamérica}); en ese caso, cambiaría también la
lectura valorizadora, dado que siempre se le da prominencia al segundo
elemento.\pendiente{En el original, las respuestas 3 y 4 están intercambiadas} 5.
a) [\etq{SConj/OC} [\etq{SF/Esp} El ser humano no se siente parte del
ecosistema], sino que\etq{Conj/n} [\etq{SF/Comp} está totalmente desconectado de
la naturaleza de la que forma parte]]. b) La coordinación adversativa con
\forma{sino} se interpreta como oposición exclusiva, en la que el segundo
coordinando reemplaza al elemento negado en primer lugar. 6. a) [\etq{SConj/OC}
[\etq{SF/Esp} Hay medidas para frenar este conflicto], pero\etq{Conj/n}
[\etq{SF/Comp} se tienen que tomar muy rápido]]. b) La coordinación adversativa
con \forma{pero} establece una contraposición entre los coordinandos, aunque, en
realidad, la contraposición se da entre las inferencias que se realizan sobre
cada coordinando: hay medidas para frenar el conflicto, de lo cual se infiere
cierto optimismo, pero las medidas se deben tomar muy rápido, de lo cual se
infiere que el panorama optimista requiere algo más que posibles medidas. Esta
interpretación se valida en una afirmación que hace el autor unas líneas más
adelante, según la cual el optimismo debe alimentarse de opciones concretas.
\end{solucion}

\begin{solucion}{\ref{ej:12-elipsis}}
2. Se elide \forma{está enamorado}. 3. Se elide \forma{trata bien a Gollum}. 4.
Se elide \forma{puede ser rey de Gondor}. 5. Se elide \forma{es un mago
ambicioso}. 6. No hay elipsis, sino una coordinación copulativa de dos SSPP. El
SConj funciona como el complemento preposicional de régimen del verbo
\forma{desconfiar}.
\end{solucion}

\begin{solucion}{\ref{ej:12-meta}}
2. a) Conjunción coordinante disyuntiva \forma{o}. b) [\etq{SConj/Comp}
[\etq{SN/Esp} la protección de datos] \textbf{o}\etq{Conj/n} [\etq{SN/Comp} su
impacto social]]. c) Coordinación homocategorial entre dos SSNN. d) La
disyunción se interpreta como inclusiva: no se impone la elección de una u otra
opción.

3. a) Conjunción coordinante adversativa \forma{pero}. b) [\etq{SConj/CPR}
[\etq{SP/Esp} sobre el uso de sus datos en general], \textbf{pero}\etq{Conj/n}
también [\etq{SP/Comp} sobre las maniobras normalizadas y ocultas del poder]].
c) Coordinación homocategorial entre dos SSPP. d) La conjunción está acompañada
por el adverbio \forma{también}, que funciona como marcador discursivo. La
coordinación se interpreta con una lectura valorizadora, en la que se da
prominencia al segundo constituyente.\pendiente{El original incluye en el primer
coordinando \textit{sobre la falta de consciencia que tienen los usuarios}}

4. (A) a) Conjunción coordinante adversativa \forma{pero}. b) [\etq{SConj/OC}
[\etq{SF/Esp} Desde Meta se ha afirmado que buscan cumplir con las normativas de
privacidad y protección de datos], \textbf{pero}\etq{Conj/n} [\etq{SF/Comp} la
desconfianza persiste]]. c) Coordinación homocategorial entre dos oraciones. d)
Los dos coordinandos son sintagmas de flexión, por lo que el SConj se constituye
como una oración compuesta. Se interpreta un valor de contraexpectativa porque
Meta afirma que cumple con la normativa pertinente y se infiere que eso debería
generar confianza en los usuarios, pero eso no sucede. (B) a) Conjunción
coordinante copulativa \forma{y}. b) [\etq{SConj/Comp} [\etq{SN/Esp} privacidad]
\textbf{y}\etq{Conj/n} [\etq{SN/Comp} protección de datos]]. c) Coordinación
homocategorial entre dos SSNN. d) No hay valores asociados a la conjunción.

5. (A) a) Conjunción coordinante adversativa \forma{sino}. b) [\etq{SConj/OC}
[\etq{SF/Esp} A su juicio, el problema no gira tanto en torno a si de aquí en
adelante la foto de tu hijo pequeño aparecerá en un futuro en anuncios de
pañales], \textbf{sino que}\etq{Conj/n} [\etq{SF/Comp} la cuestión se sitúa en
torno a si basta con que Meta dé la oportunidad a los usuarios de oponerse a que
sus datos sean utilizados con ese fin o es necesario un consentimiento
específico]]. c) Coordinación homocategorial entre dos oraciones. d) La
coordinación se interpreta con valor de oposición exclusiva, en la que el
segundo elemento reemplaza al primero. (B) a) Conjunción coordinante disyuntiva
\forma{o}. b) [\etq{SConj} [\etq{SF/Esp} si basta con que Meta dé la
oportunidad a los usuarios de oponerse a que sus datos sean utilizados con ese
fin] \textbf{o}\etq{Conj/n} [\etq{SF/Comp} (si) es necesario un consentimiento
específico]]. c) Coordinación homocategorial entre dos oraciones. d) La
disyunción se interpreta como exclusiva: se impone la elección de una u otra
opción. En el segundo coordinando hay una elipsis del subordinante \forma{si},
que ya está mencionado en el primer constituyente.\pendiente{El original
etiqueta este SConj como OComp, pero está incluido como término de \textit{en
torno a}}
\end{solucion}

% =====================================================================
%  capitulos/cap13.tex — Capítulo 13. Las relaciones interoracionales:
%  subordinación
%  Autores: Augusto M. Trombetta y Mabel Giammatteo
%  Fuente: Capitulo_13_normalizado.docx
%
%  Siglas (decisión editorial, ver HANDOFF): Ocomp = oración compuesta;
%  OS = oración subordinada (genérica); OC = oración completiva;
%  OR = oración relativa.
% =====================================================================
\capitulo{Las relaciones interoracionales: subordinación}{Augusto M. Trombetta y Mabel Giammatteo}

\section{Introducción}

En este capítulo se completa el estudio de las relaciones interoracionales con el
tratamiento de la subordinación y sus principales subclases. En primer lugar, se
plantea en qué consiste la subordinación oracional (la oración carente de
autonomía sintáctica, esto es, incluida en una estructura
mayor)\pendiente{El original dice «no incluida en una estructura mayor»} y,
luego, se hace foco en la estructura interna de las oraciones subordinadas. Para
ello se tiene en cuenta la presencia (o no) de subordinantes léxicos y su
naturaleza o bien conjuntiva o bien relativa, así como el tipo de flexión verbal
(conjugada o no finita) involucrada. La consideración de estos aspectos
estructurales se completa con el tratamiento de las funciones sintácticas que
cumplen las oraciones subordinadas. La tradición gramatical clasificó a las
subordinadas teniendo en cuenta su equivalencia funcional con algunas clases de
palabras, por lo que se distinguía entre subordinadas sustantivas, adjetivas y
adverbiales. Sin embargo, aquí se adopta una clasificación que, según el tipo de
encabezador que las introduce, distingue entre subordinadas completivas y
relativas y, dentro de estas últimas, diferencia las que se ligan con un elemento
nominal o adverbial previo, que funciona como su antecedente, y las que carecen
de antecedente léxico y se conocen como relativas libres y semilibres.
Finalmente, desde la perspectiva semántico-pragmática se consideran las
construcciones que incluyen subordinadas que transmiten valores
lógico-relacionales (causales, finales, concesivas y condicionales) o bien
valores cuantificacionales-intensivos (comparativas y consecutivas).

\section{La relación de subordinación}\label{sec:13-relacion}

Según plantean \textcite[731]{bosque2009}, «[l]a subordinación es una relación de
dependencia sintáctica», en la que se establece una relación jerárquica entre dos
constituyentes sintácticos, uno de los cuales complementa o modifica al otro, que
funciona como núcleo. Así entendida, la relación de subordinación se aplica por
igual a cualquier vinculación entre un núcleo y sus complementos o adjuntos, con
independencia de la configuración sintáctica que adopten:

\begin{cita}
la relación sintáctica que existe entre el verbo y su \textit{complemento} es la
misma en \textit{Lamento que las cosas estén así} que en \textit{Lamento la
situación}, aun cuando tradicionalmente se entendía que la oración que encabeza
\textit{que} estaba \sig{subordinada}, mientras que esa noción no se aplicaba al
grupo nominal \textit{la situación} \parencite[§\,1.13m]{rae2025}.
\end{cita}

\begin{destacado}
\textbf{Para tener en cuenta.} En sintonía con lo anterior, se puede señalar que
la ampliación del concepto de \termino{nuclearidad} a todas las categorías
gramaticales (incluidas las preposiciones; véase el \capref{5},
§\ref{sec:05-sp}), junto con la universalización de la
\termino{endocentricidad} (que sostiene que toda construcción gramatical posee un
núcleo; véase el \capref{5}, §\ref{sec:05-sintagmas}), ha permitido entender que
el núcleo de una construcción gramatical es el constituyente al que se
subordinan los restantes constituyentes.
\end{destacado}

\subsection{La subordinación oracional}\label{sec:13-oracional}

Centrándonos ahora en las relaciones entre oraciones, la
\termino{subordinación} (o \termino{hipotaxis}, según la denominación de las
gramáticas clásicas) es la relación estructural entre una oración subordinada y
una oración matriz (también llamada \textit{oración principal}), a la que se
integra y de la que forma parte. En consecuencia, la subordinada queda incluida
en un constituyente más extenso, la oración principal
\parencite[214--215]{bosque2019}, y forma con ella una estructura mayor conocida
como \termino{oración compuesta}.\footnote{El concepto de oración compuesta se
presentó en el \capref{12}, §\ref{sec:12-pasaje}.} Según iremos viendo, la
subordinada «desempeña una función dentro de la oración "principal" (sujeto,
complemento directo, etc.) o dentro de un constituyente de ella»
\parencite[731]{bosque2009}. El esquema de \exref{ex:13-diag1} representa la
relación.

\begin{exe}
  \ex\label{ex:13-diag1}
      \begin{forest} gramatica
        [Ocomp\\(oración matriz)
          [SN\\Suj]
          [SV\\Pred
            [V]
            [OS\\(argumento o adjunto) [\ldots, roof]]]]
      \end{forest}
\end{exe}

La oración subordinada (OS) se encuentra inserta en una estructura mayor
(oración matriz) de la que forma parte; junto con ella, forma una oración
compuesta (Ocomp). Dentro de esa estructura, la subordinada establece una
relación de dependencia y puede cumplir una función argumental (sujeto,
complemento) o no argumental (adjunto). En corchetes: [\etq{Ocomp} \ldots\
[\etq{OS} \ldots] ].\pendiente{Diagrama 1 rehecho («cambiar diagrama»)}

La subordinación involucra una diferencia sintáctica entre la oración subordinada
y la oración matriz a la que se integra, ya que la subordinada es una estructura
carente de autonomía sintáctica. En cambio, la oración compuesta, que resulta de
la relación matriz-subordinada, constituye una estructura sintácticamente
autónoma:

\begin{exe}
  \ex\label{ex:13-gobierno}
    \begin{xlist}
      \ex\label{ex:13-gobierno-a} {}[\etq{Ocomp} El Gobierno decidió [\etq{OC}
          aumentar los sueldos de las fuerzas de seguridad]].
      \ex\label{ex:13-gobierno-b} {}[\etq{Ocomp} El Gobierno decidió [\etq{OC}
          que aumenten los sueldos de las fuerzas de seguridad]].
    \end{xlist}
\end{exe}

Tanto \exref{ex:13-gobierno-a} como \exref{ex:13-gobierno-b} son oraciones
compuestas, que constituyen estructuras completas en términos
sintáctico-semánticos y pueden ser inicio de discurso. Ambas cualidades son
aplicables a cualquier oración simple por la que se puedan intercambiar:

\begin{exe}
  \ex\label{ex:13-simples}
    \begin{xlist}
      \ex El Gobierno decidió [\etq{SN} el aumento de los sueldos de las fuerzas
          de seguridad].
      \ex El Gobierno decidió [\etq{SN} un aumento de sueldo para las fuerzas de
          seguridad].
    \end{xlist}
\end{exe}

En contraste, aun cuando constituye una estructura completa en términos
sintáctico-semánticos, puesto que contiene una predicación, la subordinada no
puede ser inicio de discurso:

\begin{exe}
  \ex
    \begin{xlist}
      \ex[??]{aumentar los sueldos de las fuerzas de seguridad.}
      \ex[??]{que aumenten los sueldos de las fuerzas de seguridad.}
    \end{xlist}
\end{exe}

En cuanto a la oración matriz, puede ser una estructura completa, como las de
\exref{ex:13-simples}, o la presencia de la subordinada puede ser obligatoria
porque cumple una función argumental (sujeto, complemento) en ella. En cambio,
cuando la subordinada no es argumental, funciona como un modificador adjunto
dentro de la oración matriz:

\begin{exe}
  \ex\label{ex:13-decidio}
    \begin{xlist}
      \ex\label{ex:13-decidio-a} El Gobierno decidió [\etq{OC} que aumenten los
          sueldos de las fuerzas armadas].
      \ex\label{ex:13-decidio-b} El Gobierno decidió [\etq{OC} aumentar los
          sueldos de las fuerzas armadas].
      \ex\label{ex:13-decidio-c} El Gobierno decidió algo inesperado, [\etq{OC}
          que aumenten los sueldos de las fuerzas armadas].
      \ex\label{ex:13-decidio-d} El Gobierno decidió algo inesperado, [\etq{OC}
          aumentar los sueldos de las fuerzas armadas].
    \end{xlist}
\end{exe}

En \exref{ex:13-decidio-a} y \exref{ex:13-decidio-b}, la oración subordinada
cumple la función argumental de complemento objeto directo del verbo transitivo
\forma{decidir} ⟨alguien, algo⟩. En ese sentido, la subordinada contribuye a la
buena formación (gramaticalidad; véase el \capref{2},
§\ref{sec:02-aceptabilidad}) de la oración matriz, ya que
\juicio{*}{\forma{El Gobierno decidió}} forma una secuencia incompleta. En
\exref{ex:13-decidio-c} y \exref{ex:13-decidio-d}, en cambio, la misma oración
subordinada cumple la función no argumental de aposición explicativa de un
núcleo nominal (el pronombre neutro \forma{algo}; véase el \capref{5},
§\ref{sec:05-sn}) y su eventual ausencia no afectaría la buena formación
oracional, ya que \forma{El Gobierno decidió algo inesperado} es una oración
completa, tanto desde el punto de vista sintáctico como semántico.

Asimismo, en un sentido formal, los ejemplos de \exref{ex:13-decidio} dejan ver
que, internamente, las oraciones subordinadas pueden constituirse alrededor de un
verbo conjugado en forma personal (\forma{El Gobierno decidió [que
\textbf{aumenten} los sueldos de las fuerzas armadas]}) o de un verbo en forma no
personal (\forma{El Gobierno decidió [\textbf{aumentar} los sueldos de las
fuerzas armadas]}). En forma correlativa, las subordinadas con verbos finitos
están encabezadas por un subordinante que puede ser conjuntivo o relativo,
mientras que las subordinadas con verbos no finitos carecen de un encabezador que
destaque la relación subordinativa.

Recordemos que en el \capref{6}, §\ref{sec:06-capas}, se introdujo la idea de que
la estructura de la oración se puede entender como un conjunto de tres capas
formales (complementante (SC), flexiva (SF) y léxica (SV)) que se corresponden
con valores discursivos, deícticos y semánticos, respectivamente. El esquema de
\exref{ex:13-diag2} muestra que el núcleo complementante (C) puede estar
realizado léxicamente por una conjunción \forma{que / si} o puede no tener
realización léxica (Ø)\footnote{El símbolo Ø indica una categoría
fonológicamente vacía, es decir, que no está realizada por un elemento léxico.}
en la oración subordinada.

\begin{exe}
  \ex\label{ex:13-diag2}
      \begin{forest} gramatica
        [{SC = OS}
          [C\\\textit{que / si / Ø}]
          [SF [\ldots, roof]]]
      \end{forest}
\end{exe}

Dentro de la oración subordinada, los subordinantes conjuntivos (expresos o
nulos) encabezan el sintagma complementante. Como corresponde a su función
nuclear, el subordinante determina la naturaleza de la oración subordinada
(conjuntiva; véase §\ref{sec:13-encabezador}). El sintagma complementante es la
máxima proyección oracional (es decir, SC = O) y su complemento es el SF.

Así, es importante destacar que la relación de subordinación puede estar presente
aun en ausencia de encabezadores que hagan explícito el vínculo entre principal y
subordinada:

\begin{exe}
  \ex\label{ex:13-xi}
    \begin{xlist}
      \ex\label{ex:13-xi-a} Xi Jinping pregunta [\etq{OC} cuánta presión de
          Estados Unidos soportará China].
      \ex\label{ex:13-xi-b} Xi Jinping busca [\etq{OC} ejercer presión sobre
          Estados Unidos].
    \end{xlist}
\end{exe}

En los ejemplos de \exref{ex:13-xi}, ni la subordinada interrogativa
\exref{ex:13-xi-a} ni la subordinada de infinitivo \exref{ex:13-xi-b} poseen un
introductor subordinante; sin embargo, la relación de subordinación se sostiene
por los requerimientos del verbo selector correspondiente. Ambos ejemplos, las
subordinadas de naturaleza enfática (exclamativas e interrogativas parciales;
véase la nota~\ref{fn:13-exclamativas}) y las subordinadas con verbo no finito
(infinitivo, gerundio y participio), son casos típicos de oraciones subordinadas
con subordinante nulo (o sin expresión léxica). En el caso de las subordinadas
enfáticas, los ítems exclamativos o interrogativos iniciales (\forma{cuánta,
dónde, cómo}, etc.) no se ubican en el núcleo complementante, sino que
desempeñan su función sintáctica dentro de la oración subordinada de la que
forman parte. Por ejemplo, \forma{cuánta} es el especificador del SN \forma{presión
de Estados Unidos}, es decir que cumple la misma función sintáctica que le
correspondería en una oración exclamativa o interrogativa no subordinada (por
ejemplo, \forma{¿Cuánta presión de Estados Unidos soportará China?}), como se ve
en \exref{ex:13-diag3}.

\begin{exe}
  \ex\label{ex:13-diag3}
      \begin{forest} gramatica
        [{SC = OC}
          [SN\\Foco
            [D\\\textit{cuánta}]
            [N\prima [\textit{presión de Estados Unidos}, roof]]]
          [C\prima
            [C\\Ø]
            [SF [\textit{soportará China}, roof]]]]
      \end{forest}
\end{exe}

Por su lado, las oraciones subordinadas de formas verbales no finitas (verbos en
formas no personales) carecen de encabezadores léxicos que determinen la
naturaleza de la relación subordinativa \exref{ex:13-diag4}.

\begin{exe}
  \ex\label{ex:13-diag4}
      \begin{forest} gramatica
        [{SC = OC}
          [C\\Ø]
          [SF
            [SN\\PRO]
            [F\prima
              [F\\\textit{-r} (inf.)]
              [SV [\textit{ejercer presión sobre Estados Unidos}, roof]]]]]
      \end{forest}
\end{exe}

\pendiente{Diagramas 3 y 4: el original no conserva la imagen; se reconstruyeron}

\begin{ejercicio}\label{ej:13-paremias}
En muchas paremias (refranes, dichos y proverbios tradicionales) suelen aparecer
oraciones subordinadas. Identifique en los siguientes ejemplos la oración
subordinada (OS) dentro de la oración compuesta (Ocomp) y colóquela entre
corchetes, según se muestra en los ejemplos resueltos. Tenga en cuenta que puede
existir más de una oración subordinada en una matriz.

\begin{exe}
  \ex A la cama no te irás sin saber una cosa más.
  \ex A quien madruga, Dios lo ayuda.
  \ex Cuando el invierno primavea, la primavera invernea.
  \ex Cuando veas las barbas de tu vecino cortar, pon las tuyas a remojar.
  \ex Dando gracias por agravios negocian los hombres sabios.
  \ex Diciendo las verdades se pierden amistades.
  \ex Donde fueres, haz lo que vieres.
  \ex Donde hubo fuego, cenizas quedan.
  \ex El salir de la posada es la mejor jornada.
  \ex Es de bien nacidos ser agradecidos.
  \ex Hablando del rey de Roma, por la puerta asoma.
  \ex Lo que natura non da Salamanca non presta.
  \ex No por mucho madrugar amanece más temprano.
  \ex Nunca es tarde si la dicha es buena.
  \ex Pájaro que comió, voló.
  \ex Perro que ladra no muerde.
  \ex Preguntando se llega a Roma.
  \ex Si quieres ser bien servido, sírvete a ti mismo.
\end{exe}

\begin{resuelto}[Ejemplos resueltos]
1. [\etq{Ocomp} A la cama no te irás sin [\etq{OS} \textit{saber una cosa más}]].
2. [\etq{Ocomp} A [\etq{OS} \textit{quien madruga}], Dios lo ayuda].
\end{resuelto}
\end{ejercicio}

\section{Estructura interna de las oraciones subordinadas}\label{sec:13-estructura}

Dentro de este apartado atendemos a dos cuestiones estructurales: en primer
lugar, damos cuenta de la presencia o no de un introductor de la subordinada y
del tipo de encabezador que la introduce ---conjuntivo o relativo---; en segundo
lugar, nos ocupamos del tipo de flexión ---conjugada (finita o personal) o no
conjugada (no finita o no personal)--- que manifiesta el verbo de la subordinada.

\subsection{Presencia y tipo de encabezador}\label{sec:13-encabezador}

Las oraciones subordinadas de verbo finito (o conjugado en formas personales) son
introducidas por un ítem léxico, una conjunción subordinante (\forma{que, cuando,
si}, etcétera). Recordando lo visto en el \capref{5}, §\ref{sec:05-clases}, sobre
las clases de palabras funcionales o gramaticales, las conjunciones son unidades
léxicas que establecen relaciones entre distintos constituyentes de la oración
para conectarlos sintácticamente \parencite{giammatteo2009}. Esta función
gramatical de conexión involucra dos tipos diferentes de relaciones, según la
jerarquía sintáctica de los constituyentes involucrados: la coordinación, tratada
en el capítulo anterior, y la subordinación. Se reconocen así dos grandes
subclases de ítems conjuntivos: los coordinantes y los subordinantes. Las
conjunciones coordinantes, según lo ya visto, establecen relaciones entre
constituyentes de igual jerarquía sintáctica. Los elementos conjuntivos de
naturaleza subordinante, en contraste, marcan relaciones entre constituyentes de
distinta jerarquía sintáctica y tienen como función incluir una estructura
predicativa, la subordinada, en otra, que es la oración principal o matriz.

Asimismo, dentro de los subordinantes, se pueden reconocer dos tipos o subclases:
las conjunciones y los relativos. Si las conjunciones subordinantes (tales como
\forma{que, cuando, si}) se pueden caracterizar a partir de su función
subordinativa, los relativos (tales como \forma{que, el que, quien, cual, cuyo,
donde, cuando, como, cuanto}), que son de naturaleza pronominal, además de
incluir una oración subordinada en la matriz, cumplen una función sintáctica
determinada dentro de ella \parencite{giammatteo2009}.

Las conjunciones subordinantes, además de hacer explícita la relación de
subordinación, pueden aportar valores relacionados con la modalidad de la oración
subordinada. En los ejemplos de \exref{ex:13-modalidad}, las conjunciones
subordinantes marcan la modalidad declarativa \exref{ex:13-modalidad-a},
interrogativa \exref{ex:13-modalidad-b} o exhortativa \exref{ex:13-modalidad-c}
de las oraciones subordinadas que encabezan, en consonancia con los
requerimientos del verbo matriz correspondiente:

\begin{exe}
  \ex\label{ex:13-modalidad}
    \begin{xlist}
      \ex\label{ex:13-modalidad-a} Xi Jinping afirma [\etq{OC} \textit{que} la
          presión de Estados Unidos sobre China no tendrá éxito].
      \ex\label{ex:13-modalidad-b} Xi Jinping pregunta [\etq{OC} \textit{si} la
          presión de Estados Unidos sobre China tendrá éxito].
      \ex\label{ex:13-modalidad-c} Xi Jinping ordenó [\etq{OC} \textit{que} sus
          funcionarios no cedan ante la presión de Estados Unidos].
    \end{xlist}
\end{exe}

Los relativos, al igual que las conjunciones subordinantes, hacen explícita la
relación de subordinación y, además, suman una segunda función, interna a la
subordinada que introducen. En los ejemplos de \exref{ex:13-relativos}, los
relativos \forma{que} \exref{ex:13-relativos-a} y \forma{donde}
\exref{ex:13-relativos-b} cumplen funciones sintácticas en la oración
subordinada. En \exref{ex:13-relativos-a}, el relativo, cuyo antecedente es
\forma{la presión}, es el complemento objeto directo. Por su parte, en
\exref{ex:13-relativos-b}, el relativo, cuyo antecedente es \forma{las áreas}, es
el complemento de la preposición \forma{de}, con la que forma un complemento
circunstancial de lugar:

\begin{exe}
  \ex\label{ex:13-relativos}
    \begin{xlist}
      \ex\label{ex:13-relativos-a} Xi Jinping condena la presión [\etq{OR}
          [\textit{que}]\etq{COD} ejerce Estados Unidos sobre China].\\
          \textit{que} = la \textit{presión}.
      \ex\label{ex:13-relativos-b} Xi Jinping aumenta las áreas [\etq{OR} [de
          \textit{donde}]\etq{CCL} China extrae hierro y carbón].\\
          (de) \textit{donde} = (de) las \textit{áreas}.
    \end{xlist}
\end{exe}

Como ya adelantamos, típicamente carecen de introductor léxico explícito las
oraciones subordinadas articuladas alrededor de ítems interrogativos y
exclamativos, y también otras subordinadas que admiten la inserción de la
conjunción subordinante \forma{que}, así como las subordinadas articuladas
alrededor de un verbo no finito (o no conjugado en formas personales). En
síntesis, las oraciones subordinadas se pueden articular de dos modos: con o sin
introductor léxico (tabla~\ref{tab:13-encabezadores}).

{\small
\begin{longtable}{@{}>{\raggedright\arraybackslash}p{40mm}
    >{\raggedright\arraybackslash}p{83mm}@{}}
  \caption{Subordinación con y sin encabezadores léxicos}\label{tab:13-encabezadores}\\
  \toprule
  \endfirsthead
  \toprule
  \endhead
  Subordinación con encabezadores léxicos explícitos (marcan relaciones entre
    constituyentes de distinta jerarquía sintáctica) &
    \textbf{Conjunciones subordinantes.} Semánticamente, indican modalidad (son
    índices de modalidad declarativa, interrogativa, exhortativa).
    Sintácticamente, encabezan el sintagma complementante (SC), sin cumplir otra
    función sintáctica adicional. Morfológicamente, son invariables. \\
  \addlinespace
  & \textbf{Relativos.} Semánticamente, toman la referencia de un antecedente o
    suman valores generales (cosa, persona, cantidad, lugar, modo, tiempo).
    Sintácticamente, introducen el sintagma complementante (SC).
    Morfológicamente, pueden ser invariables o variables. \\
  \midrule
  Subordinación sin encabezadores léxicos explícitos &
    Oraciones subordinadas interrogativas y exclamativas. Oraciones subordinadas
    con \textit{que} facultativo. Oraciones subordinadas con formas verbales no
    finitas. \\
  \bottomrule
\end{longtable}}

\begin{ejercicio}\label{ej:13-introductor}
Retome los resultados del ejercicio~\ref{ej:13-paremias} y, si corresponde,
señale el tipo de introductor de cada subordinada, según se muestra en los
ejemplos resueltos.

\begin{resuelto}[Ejemplos resueltos]
1. Sin subordinante. 2. Con relativo \forma{quien}.
\end{resuelto}
\end{ejercicio}

\begin{ejercicio}\label{ej:13-principito}
A partir de los siguientes pares de oraciones simples, adaptadas de \textit{El
principito}, de Antoine de Saint-Exupéry (1943), construya una oración compuesta
(matriz-subordinada), según se muestra en el ejemplo resuelto~1.

\begin{exe}
  \ex Lo esencial es invisible a los ojos. Lo dijo el Principito.
  \ex A veces tenemos que aguantar a las orugas. Queremos disfrutar de las
      mariposas.
  \ex Eres el dueño de tu vida y tus emociones. Nunca te olvides de esto.
  \ex Los hombres ya no tienen tiempo. Compran cosas hechas a los comerciantes.
  \ex Me pregunto algo. ¿Las estrellas se iluminan con algún fin?
  \ex Nunca estamos contentos en un lugar. Estamos en ese lugar.
  \ex Todas las estrellas ríen. En ese momento brilla la luna.
  \ex Todas las personas grandes han sido niños antes. Pocas lo recuerdan.
  \ex Un hombre es vanidoso. Todos los demás hombres son admiradores.
  \ex Uno termina de arreglarse. Entonces hay que hacer cuidadosamente la
      limpieza del planeta.
\end{exe}

\begin{resuelto}
1. [\etq{Ocomp} El Principito dijo [\etq{OC} que lo esencial era invisible a los
ojos]].
\end{resuelto}
\end{ejercicio}

\subsubsection{Los encabezadores de las subordinadas}\label{sec:13-encabezadores}

Según hemos dicho, hay dos tipos de encabezadores léxicos de oraciones
subordinadas: las conjunciones subordinantes y los relativos. A continuación,
dedicamos un apartado a cada uno de estos tipos de introductores.

\paragraph{Las conjunciones subordinantes.}\label{sec:13-conjunciones}
Las conjunciones subordinantes son los núcleos funcionales de la oración (es
decir, del sintagma complementante, la máxima proyección oracional) y determinan
la naturaleza de la subordinada. Carentes de valor referencial, son responsables
de la interpretación modal de la subordinada; así, según se mostró en los
ejemplos de \exref{ex:13-modalidad}, \forma{que} es índice de modalidad
declarativa / exhortativa y \forma{si} marca modalidad interrogativa.

En los análisis clásicos de la subordinación oracional, las conjunciones
subordinantes y las locuciones conjuntivas derivadas de estas se asociaban a
distintos valores lógico-relacionales (tiempo, causa, concesión, condición,
finalidad). En otras palabras, se contaba con una lista de formas conjuntivas
simples y complejas dotadas de valores específicos: \forma{porque, como, que, a
causa de que, en razón de que, en la medida en que, ya que, puesto que, dado
que} (causales); \forma{aunque, aun cuando, si bien, eso que, a pesar de que,
pese a que, aun si, incluso si} (concesivas); \forma{si, de, siempre que, siempre
y cuando, a no ser que, en caso de que, a condición de que, con tal de que}
(condicionales); \forma{para que, a fin de que, con el objeto de que, con el fin
de que, con vistas a que} (finales); \forma{antes que, después que, siempre que,
ni bien, no bien} (temporales).

En la enumeración anterior, además de ciertas repeticiones de formas, se advierte
fácilmente que las conjunciones y locuciones subordinantes, en su gran mayoría,
pueden resumirse como una construcción encabezada por un adverbio o preposición
seguidos de una oración subordinada encabezada por las conjunciones \forma{que}
(mayoritariamente) o \forma{si} (en menor medida): \forma{antes (que), después
(que), siempre (que); por(que), a causa de (que), en razón de (que), en la medida
en (que), ya (que), puesto (que), dado (que); aun(que), (si) bien, eso (que), a
pesar de (que), pese a (que), aun (si), incluso (si); siempre (que), a no ser
(que), en caso de (que), a condición de (que), con tal de (que); para (que), a fin
de (que), con el objeto de (que), con el fin de (que), con vistas a (que)}. Por
tanto, si aceptamos la distinción que acabamos de hacer, se pueden admitir dos
versiones en el análisis de las conjunciones y locuciones conjuntivas
subordinantes:

\begin{itemize}
  \item una \termino{versión sintética}, que es la tradicional, en la que las
        locuciones conjuntivas subordinantes funcionan como encabezadores
        unitarios dotados de valores específicos ---\forma{aunque} (concesivo),
        \forma{porque} (causal), \forma{siempre que} (temporal)---;
  \item una \termino{versión analítica}, que es la alternativa propuesta en
        \textcite{bosque2009} y en \textcite{brucart2009,brucart2016}, en la que
        las conjunciones subordinantes funcionan como encabezadores de oraciones
        subordinadas dependientes de un núcleo preposicional o adverbial, que es
        el que posee el valor específico.
\end{itemize}

En forma esquemática, es posible representar la forma \forma{porque} de las
siguientes maneras:

\begin{exe}
  \ex\label{ex:13-porque}
    \begin{xlist}
      \ex\label{ex:13-porque-a} Versión sintética:\\[2pt]
          \begin{forest} gramatica
            [{SC = OS}
              [C\\\textit{porque}]
              [SF [\ldots, roof]]]
          \end{forest}
      \ex\label{ex:13-porque-b} Versión analítica:\\[2pt]
          \begin{forest} gramatica
            [SP
              [P\\\textit{por}]
              [{SC = OC}
                [C\\\textit{que}]
                [SF [\ldots, roof]]]]
          \end{forest}
    \end{xlist}
\end{exe}

En la versión sintética \exref{ex:13-porque-a}, la locución conjuntiva, como
cualquier conjunción subordinante, encabeza el sintagma complementante (SC), la
máxima proyección oracional; es el análisis clásico más extendido: [\etq{O}
\ldots\ [\etq{OS} \textit{porque} \ldots]]. En la versión analítica
\exref{ex:13-porque-b}, la conjunción subordinante encabeza el sintagma
complementante (SC), y esta proyección oracional es el complemento de un núcleo
preposicional o adverbial; es el análisis alternativo más reciente: [\etq{O}
\ldots\ [\etq{SP} \textit{por} [\etq{OC} \textit{que} \ldots]]].

Pese a ser el más extendido, el enfoque sintético de las conjunciones y
locuciones conjuntivas subordinantes pierde de vista una interesante reflexión
respecto de estas estructuras: no parece tomar en cuenta la sistemática aparición
de los dos subordinantes básicos (\forma{que, si}) en todas ellas. Esta
regularidad queda mejor analizada en el enfoque analítico, más simple y
regularizador (aunque inicialmente pueda parecernos más extraño porque no estamos
acostumbrados a pensar así estas construcciones): las conjunciones subordinantes
(\forma{que, si}) encabezan oraciones subordinadas que funcionan como
complementos de adverbios o de preposiciones que, a su vez, son las responsables
del valor que toma la construcción (temporal, causal, concesivo, condicional,
final). Así analizadas, estas expresiones adverbiales y preposicionales son las
mismas que, mayoritariamente, pueden tomar complementos no oracionales para
articular esos mismos valores: \forma{antes (de) \textbf{que comamos} / antes de
\textbf{la comida}; después (de) \textbf{que me visitaste} / después de
\textbf{tu visita}; por\textbf{que me lo pediste} / por \textbf{tu pedido}; a
causa de \textbf{que nevó copiosamente} / a causa de \textbf{una copiosa
nevada}; en razón de \textbf{que intervino el rector} / en razón de \textbf{la
intervención del rector}; dado \textbf{que cumple con los requisitos} / dado
\textbf{el cumplimiento de los requisitos}; aun\textbf{que coma poco} / aun
\textbf{poco alimentado}; a pesar de \textbf{que insistas} / a pesar de
\textbf{tu insistencia}; pese a \textbf{que se marchó} / pese a \textbf{su
salida}; aun \textbf{si estaba enferma} / aun \textbf{enferma}; en caso de
\textbf{que suceda un incendio} / en caso de \textbf{incendio}; para \textbf{que
estudies} / para \textbf{tus estudios}; con vistas a \textbf{que se promuevan
nuevos talentos} / con vistas a \textbf{la promoción de nuevos talentos}}.

\begin{ejercicio}\label{ej:13-analitica}
A partir de las siguientes oraciones compuestas, describa analíticamente las
conjunciones subordinantes y, donde corresponda, los valores que manifiestan los
adverbios, preposiciones o locuciones que las preceden, según se muestra en el
ejemplo resuelto~1.

\begin{exe}
  \ex Actuó así a pesar de que le aconsejaron lo contrario.
  \ex Aunque llueva, vamos a salir.
  \ex Habló con el rector con el fin de que se aprobase su curso.
  \ex Lo explico de nuevo para que no queden dudas.
  \ex Me alegro porque lo hice con cariño.
  \ex No pude llegar antes porque había mucho tráfico.
  \ex No va a aceptar, aunque se lo pidamos amablemente.
  \ex Nos alegra tu visita porque te extrañábamos.
  \ex Te lo repito para que lo entiendas.
  \ex Trabajamos con el fin de que el proyecto sea un éxito.
\end{exe}

\begin{resuelto}
1. [\etq{Ocomp} Actuó así [\etq{SP} \textit{a pesar de} [\etq{OC}
\textbf{\textit{que}} le aconsejaron lo contrario]]], con valor concesivo en
\forma{a pesar de} (locución preposicional).
\end{resuelto}
\end{ejercicio}

\paragraph{Los relativos.}\label{sec:13-relativos}
Al igual que las conjunciones subordinantes, los relativos son los núcleos
funcionales de la oración (es decir, del sintagma complementante, la máxima
proyección oracional) y determinan la naturaleza de la oración subordinada como
relativa. El carácter pronominal de estas unidades es el factor clave que
establece una diferencia frente a las conjunciones subordinantes: dado que los
relativos constituyen expresiones referenciales, cumplen funciones sintácticas,
tanto argumentales como no argumentales, dentro de la oración que encabezan
(tabla~\ref{tab:13-relativos}).

{\small
\begin{longtable}{@{}>{\raggedright\arraybackslash}p{40mm}
    >{\raggedright\arraybackslash}p{83mm}@{}}
  \caption{Clases de relativos}\label{tab:13-relativos}\\
  \toprule
  \endfirsthead
  \toprule
  \endhead
  Con antecedente léxico explícito & Adverbios: \textit{como, cuando, donde,
    cuanto}. Determinantes: \textit{cuya / cuyo / cuyas / cuyos}. Pronombres
    (simples y complejos): \textit{que; quien / quienes; la / el cual, las / los
    cuales; el / la / los / las que}. \\
  \midrule
  Sin antecedente léxico explícito & Libres (antecedente implícito):
    \textit{como, cuando, cuanto, donde; quien / quienes; cuanta / cuanto /
    cuantas / cuantos}. Semilibres (antecedente tácito o elidido): \textit{el /
    la / los / las que, lo que}. \\
  \bottomrule
\end{longtable}}

Los relativos pueden tener un antecedente léxico explícito con el cual establecen
una relación de correferencia, como en \forma{El día que me quieras}, en el que
el pronombre relativo \forma{que} toma la referencia de su antecedente nominal
\forma{día}. Pero los relativos también pueden emplearse sin un antecedente
léxico explícito: en esta función se los denomina \termino{relativos libres}. En
su significado, estos relativos incorporan una palabra general como antecedente
implícito, como en \forma{Cuando tallan los recuerdos}, donde el relativo
\forma{cuando} se interpreta como \sig{en el momento en que}.

Los relativos simples (\forma{que, quien, cuanto, donde, cuando\ldots}) y los que
incorporan el artículo definido (\forma{el que, las que, lo que, los cuales, la
cual, lo cual\ldots}) presentan comportamientos morfológicos distintos: unos son
invariables (\forma{que, como, cuando, donde}), otros flexionan en número
(\forma{quien, quienes; cual, cuales}) y otros varían en género y número
(\forma{cuanto, cuanta, cuantos, cuantas; cuya, cuyo, cuyas, cuyos}).

\begin{recuadro}[Para ampliar]
La mayoría de los relativos son formas simples (una sola palabra), ya se trate de
\textit{relativos con antecedente léxico explícito}, como en
\exref{ex:13-ambulancia-a} y \exref{ex:13-ambulancia-b}, o de \textit{relativos
libres}, como en \exref{ex:13-ambulancia-c} y \exref{ex:13-ambulancia-d}:

\begin{exe}
  \ex Formas relativas simples (una sola palabra):
    \begin{xlist}
      \ex\label{ex:13-ambulancia-a} La ambulancia [\etq{OR} \textit{que}
          robaron los asaltantes] llevaba un paciente.
      \ex\label{ex:13-ambulancia-b} La ambulancia [\etq{OR} \textit{donde}
          estaba el paciente] fue robada por los asaltantes.
      \ex\label{ex:13-ambulancia-c} [\etq{OR} \textit{Quienes} robaron la
          ambulancia] habían realizado otros asaltos previamente.
      \ex\label{ex:13-ambulancia-d} Los asaltantes abandonaron la ambulancia
          [\etq{OR} por \textit{donde} la habían robado].
    \end{xlist}
\end{exe}

En \exref{ex:13-ambulancia-a} y \exref{ex:13-ambulancia-b}, \forma{que} y
\forma{donde} son relativos ligados anafóricamente a un antecedente léxico
(\forma{ambulancia}): \forma{que robaron\ldots} = \sig{los asaltantes robaron la
ambulancia}; \forma{donde estaba\ldots} = \sig{en la ambulancia estaba un
paciente}. En \exref{ex:13-ambulancia-c} y \exref{ex:13-ambulancia-d},
\forma{quienes} y \forma{donde} son relativos libres, sin un antecedente léxico
explícito: \forma{quienes robaron\ldots} = \sig{las personas que robaron la
ambulancia}; \forma{por donde la habían\ldots} = \sig{por el lugar donde la
habían robado} (antecedentes implícitos).

Adicionalmente, también hay formas relativas compuestas: los pronombres relativos
complejos, los pronombres relativos semilibres y los pronombres relativos con
antecedente elidido \parencite[§\,44.1d--g]{rae2025}. Los pronombres
\textit{relativos complejos} de \exref{ex:13-compuestas-a} y
\exref{ex:13-compuestas-b}, ligados anafóricamente a un antecedente léxico
explícito, forman unidades que no admiten material léxico intermedio. En
contraste, los pronombres \textit{relativos semilibres} de
\exref{ex:13-compuestas-c} y \exref{ex:13-compuestas-d}, con antecedente tácito
independiente del contexto, y los pronombres \textit{relativos con elisión del
antecedente} de \exref{ex:13-compuestas-e} y \exref{ex:13-compuestas-f}, con
antecedente dependiente del contexto, carecen de un antecedente explícito y
admiten material léxico intermedio, ya que se constituyen de dos unidades
diferentes:

\begin{exe}
  \ex Formas relativas compuestas (más de una palabra):
    \begin{xlist}
      \ex\label{ex:13-compuestas-a} La ambulancia [\etq{OR} en \textit{la que}
          huyeron los asaltantes] fue robada a los recolectores.
      \ex[*]{\label{ex:13-compuestas-b} La ambulancia [en \textit{la única que}
          huyeron los asaltantes] fue robada a los recolectores.}
      \ex\label{ex:13-compuestas-c} En realidad, la Ø [\etq{OR} \textit{que} se
          acordó de tu cumpleaños] fue mi hermana.
      \ex\label{ex:13-compuestas-d} En realidad, la única Ø [\etq{OR}
          \textit{que} se acordó de tu cumpleaños] fue mi hermana.
      \ex\label{ex:13-compuestas-e} Había dos ambulancias. Los asaltantes
          huyeron en la Ø [\etq{OR} \textit{que} llevaba un paciente].
      \ex\label{ex:13-compuestas-f} Había dos ambulancias. Los asaltantes
          huyeron en la única Ø [\etq{OR} \textit{que} llevaba un paciente].
    \end{xlist}
\end{exe}

En \exref{ex:13-compuestas-a}, \forma{la que} (= \forma{la cual}) es un pronombre
relativo complejo con antecedente léxico explícito: \forma{en la que
huyeron\ldots} = \sig{los asaltantes huyeron en la ambulancia}. En
\exref{ex:13-compuestas-c}, \forma{la que} es un pronombre relativo semilibre,
sin antecedente léxico explícito: \forma{la \{Ø\} que se acordó\ldots} = \sig{la
persona que se acordó} (antecedente tácito no contextual). En
\exref{ex:13-compuestas-e}, \forma{la que} es un pronombre relativo con elisión
de un antecedente léxico: \forma{la \{Ø\} que llevaba\ldots} = \sig{la
ambulancia que llevaba} (antecedente elidido contextual).
\end{recuadro}

\subsection{Tipo de flexión de las subordinadas}\label{sec:13-flexion}

Las oraciones subordinadas pueden articularse alrededor de formas verbales
finitas (verbos con flexión temporal-aspectual-modal y rasgos de concordancia
personal) o alrededor de formas verbales no finitas (verbos sin flexión
temporal-modal ni rasgos de concordancia personal), como el infinitivo, el
gerundio y el participio, diferenciados entre sí por su forma y sus valores
aspectuales (véase el \capref{11}).

Por un lado, las oraciones subordinadas con encabezadores léxicos (conjunciones
subordinantes o relativos) tienen asociadas flexiones finitas, cuyas
posibilidades témporo-modales dependen tanto de las correlaciones establecidas
con los verbos de la oración matriz como de los encabezadores que las introducen,
ya que las conjunciones subordinantes son más selectivas de los modos verbales
que los relativos.

\begin{exe}
  \ex\label{ex:13-modo1}
    \begin{xlist}
      \ex\label{ex:13-modo1-a} Xi Jinping se pregunta [\etq{OC} \textit{si} China
          tendrá nuevos éxitos económicos].
      \ex[*]{\label{ex:13-modo1-b} Xi Jinping se pregunta [\textit{si} China
          tenga nuevos éxitos económicos].}
    \end{xlist}
  \ex\label{ex:13-modo2}
    \begin{xlist}
      \ex\label{ex:13-modo2-a} Por eso, [\etq{SAdv} siempre [\etq{OC}
          \textit{que} podamos]], hagamos bien a todos.
      \ex[*]{\label{ex:13-modo2-b} Por eso, [siempre \textit{que} podemos],
          hagamos bien a todos.}
    \end{xlist}
  \ex\label{ex:13-modo3}
    \begin{xlist}
      \ex\label{ex:13-modo3-a} Busco a una persona [\etq{OR} \textit{que} sabe
          hablar castellano].
      \ex\label{ex:13-modo3-b} Busco una persona [\etq{OR} \textit{que} sepa
          hablar castellano].
    \end{xlist}
\end{exe}

El contraste de gramaticalidad entre \exref{ex:13-modo1-a} y
\exref{ex:13-modo1-b} revela que el modo indicativo depende de la conjunción
subordinante (\forma{si}), que marca la modalidad interrogativa de la oración
subordinada (dado que el modo subjuntivo es incompatible con la modalidad
interrogativa). La misma diferencia de aceptabilidad se advierte en el contraste
entre \exref{ex:13-modo2-a} y \exref{ex:13-modo2-b}, donde la lectura
condicional\footnote{La lectura temporal de \forma{siempre que}, como en
\forma{Siempre que podemos hacemos el bien a todos}, aparece asociada al modo
indicativo.} del sintagma encabezado por el adverbio \forma{siempre} se encuentra
asociada al modo subjuntivo de la subordinada encabezada por la conjunción
\forma{que}. En cambio, la buena formación de \exref{ex:13-modo3-a} y
\exref{ex:13-modo3-b} muestra que, en estos casos, el modo verbal indicativo o
subjuntivo no depende del pronombre relativo \forma{que}, sino del carácter
definido y existente (o no) de la entidad referida por el nombre.

Por otro lado, las oraciones subordinadas sin encabezadores léxicos, como las
interrogativas, se construyen con formas verbales finitas y no finitas, además de
ser compatibles con una conjunción subordinante
(\forma{que}).\footnote{Como ilustración histórica, ya en obras clásicas del Siglo
de Oro se evidencia que los ítems interrogativos constituyen focos
interrogativos y no núcleos de la máxima proyección oracional (SC), dado que
pueden coexistir con el núcleo \forma{que}: \forma{el soldado y yo le
preguntamos que quién era aquel caballero} (F. de Quevedo, \textit{Historia de
la vida del Buscón}). Iguales en su forma, los relativos, en cambio, no coexisten
con la conjunción subordinante. En la actualidad esta característica se puede
encontrar en la lengua coloquial.}

\begin{exe}
  \ex\label{ex:13-interr}
    \begin{xlist}
      \ex\label{ex:13-interr-a} Xi Jinping se pregunta [\etq{OC} qué nuevos
          éxitos económicos producirá China].
      \ex\label{ex:13-interr-b} Xi Jinping se pregunta [\etq{OC} cómo producir
          nuevos éxitos económicos en China].
    \end{xlist}
  \ex\label{ex:13-excl}
    \begin{xlist}
      \ex\label{ex:13-excl-a} No sabés [\etq{OC} cuánto me alegra ver que estás
          bien].
      \ex\label{ex:13-excl-b} Ya vas a ver [\etq{OC} qué impresionante es la
          casona que está en venta].
    \end{xlist}
\end{exe}

Los ejemplos \exref{ex:13-interr-a} y \exref{ex:13-interr-b} muestran la
compatibilidad de la flexión finita y no finita con los ítems interrogativos. En
cambio, las subordinadas que contienen ítems exclamativos, \exref{ex:13-excl-a} y
\exref{ex:13-excl-b}, considerados como focos, solo se construyen con una forma
verbal finita. El esquema de \exref{ex:13-diag7} muestra la interpretación de las
subordinadas con foco interrogativo y exclamativo.

\begin{exe}
  \ex\label{ex:13-diag7}
    \begin{xlist}
      \ex Foco interrogativo:\\[2pt]
          \begin{forest} gramatica
            [{SC = OC}
              [SN\\Foco [\textit{qué nuevos éxitos económicos}, roof]]
              [C\prima
                [C\\Ø]
                [SF [\textit{producirá China}, roof]]]]
          \end{forest}
      \ex Foco exclamativo:\\[2pt]
          \begin{forest} gramatica
            [{SC = OC}
              [SAdv\\Foco [\textit{cuánto}]]
              [C\prima
                [C\\Ø]
                [SF [\textit{me alegra ver que estás bien}, roof]]]]
          \end{forest}
    \end{xlist}
\end{exe}

\pendiente{Diagrama 7 sin imagen en el original; reconstruido}

En oraciones como \forma{Solicito se me inscriba fuera de término en condición de
alumno regular} o \forma{Espero me concedan una extensión del plazo
establecido}, normalmente asociadas a registros formales, se puede reponer la
conjunción subordinante \forma{que} sin alterar el significado oracional:
\forma{Solicito \textbf{que} se me inscriba fuera de término en condición de
alumno regular} o \forma{Espero \textbf{que} me concedan una extensión del plazo
establecido}. La tabla~\ref{tab:13-tipos} sintetiza los tipos de subordinadas
según su encabezador y su flexión.

{\small
\begin{longtable}{@{}>{\raggedright\arraybackslash}p{40mm}
    >{\raggedright\arraybackslash}p{83mm}@{}}
  \caption{Tipos de subordinadas según el encabezador y la flexión}\label{tab:13-tipos}\\
  \toprule
  \endfirsthead
  \toprule
  \endhead
  Subordinación con encabezadores léxicos explícitos &
    \textbf{Oraciones subordinadas completivas o conjuntivas (OC).} Conjunciones
    subordinantes de naturaleza conjuntiva encabezan el sintagma complementante
    (SC) sin cumplir una función sintáctica adicional. Formas verbales finitas
    (flexión temporal-aspectual-modal y rasgos de concordancia personal). \\
  \addlinespace
  & \textbf{Oraciones subordinadas relativas (OR).} Relativos de distinta
    naturaleza categorial (adverbios, determinantes o pronombres relativos)
    encabezan el sintagma complementante (SC) y, además, cumplen una función
    sintáctica nuclear. Formas verbales finitas. \\
  \midrule
  Subordinación sin encabezadores léxicos explícitos &
    \textbf{Oraciones subordinadas interrogativas y exclamativas indirectas
    (parciales).} Con foco interrogativo (adverbios, determinantes y pronombres
    interrogativos) y formas verbales finitas o no finitas. Con foco exclamativo
    (adverbios, determinantes y pronombres exclamativos) y formas verbales
    finitas. \\
  \addlinespace
  & \textbf{Oraciones subordinadas de \textit{que} facultativo.} Con potencial
    inserción de la conjunción y con formas verbales finitas (registros
    formales). \\
  \addlinespace
  & \textbf{Oraciones subordinadas de infinitivo, de gerundio y de participio.}
    Formas verbales no finitas. \\
  \bottomrule
\end{longtable}}

\begin{ejercicio}\label{ej:13-perifrasis}
En los siguientes ejemplos hay oraciones subordinadas con verbo flexionado y
otras con verbo no flexionado, además de predicados constituidos alrededor de una
perífrasis verbal (véase el \capref{11}). Según se muestra en los ejemplos
resueltos, distinga entre oraciones subordinadas (con o sin verbo en formas
personales) y perífrasis verbales y, donde corresponda, señale el elemento
introductor de la oración subordinada.

\begin{exe}
  \ex Como estoy cansado, no voy a la fiesta.
  \ex Estando cansado, no voy a la fiesta.
  \ex No voy a la fiesta porque estoy cansado.
  \ex Aun estando cansados, trabajamos duro.
  \ex Aunque estábamos cansados, trabajamos duro.
  \ex No estando invitado, no puedo ir al cumpleaños. Si no me invitan, no puedo
      ir al cumpleaños.
  \ex Es importante comprender el funcionamiento del motor.
  \ex Es importante que comprendas el funcionamiento del motor.
  \ex Logró que sus amigos supieran la verdad.
  \ex Logró saber la verdad.
  \ex Quiso que sus amigos supieran la verdad.
  \ex Quiso saber la verdad.
  \ex Se negó a que sus amigos supieran la verdad.
  \ex Se negó a saber la verdad.
  \ex De esta manera comenzó a saber la verdad.
  \ex De esta manera pudo saber la verdad.
  \ex De esta manera va a saber la verdad.
  \ex El tenor está cantando el aria.
  \ex El tenor ingresó a la escena mientras cantaba el aria.
  \ex El tenor ingresó cantando el aria.
\end{exe}

\begin{resuelto}[Ejemplos resueltos]
1. [\etq{Ocomp} [\etq{OS} \textbf{\textit{Como}} \textit{estoy cansado}], no voy a
la fiesta]: oración subordinada con verbo en forma personal introducida por la
conjunción subordinante \forma{como}. 2. [\etq{Ocomp} [\etq{OS} \textit{Estando
cansado}], no voy a la fiesta]: oración subordinada con verbo en forma no
personal, gerundio. 18. [\etq{O} El tenor \textit{está cantando} el aria]:
oración simple con perífrasis verbal de \forma{estar} + gerundio, de valor
progresivo, evento en curso.\pendiente{En 1 y 2 el original habla de
«perífrasis verbal», pero \textit{estoy cansado} y \textit{estando cansado} no
lo son}
\end{resuelto}
\end{ejercicio}

\section{Función sintáctica de las oraciones subordinadas}\label{sec:13-funcion}

\subsection{De la equivalencia funcional a una agrupación basada en la estructura interna}\label{sec:13-equivalencia}

Las oraciones subordinadas cumplen distintas funciones sintácticas en la oración
matriz de la que forman parte. Se trata de las mismas funciones sintácticas que
cumplirían un sintagma nominal, un sintagma adjetival, un sintagma adverbial o un
sintagma preposicional, aun cuando una oración subordinada es la máxima
proyección de una categoría funcional (sintagma complementante) y no de una
categoría léxica. Se comprende, entonces, que el criterio tradicional para
clasificar las oraciones subordinadas haya sido el de su \termino{equivalencia
funcional} con clases de palabras léxicas. Así, tanto en las gramáticas como en
los diseños curriculares, se suele distinguir, por su potencial funcional, entre
oraciones subordinadas sustantivas, adjetivas y adverbiales:

\begin{enumerate}[label=(\alph*)]
  \item Son \textit{oraciones subordinadas sustantivas} las que pueden cumplir
        las funciones sintácticas del sustantivo (sintagma nominal), por
        ejemplo, la de sujeto de la oración.
  \item Son \textit{oraciones subordinadas adjetivas} las que pueden cumplir las
        funciones sintácticas del adjetivo (sintagma adjetival), por ejemplo, la
        de modificador del nombre.
  \item Son \textit{oraciones subordinadas adverbiales} las que pueden cumplir
        las funciones sintácticas del adverbio (sintagma adverbial), por
        ejemplo, la de circunstancial en el predicado.
\end{enumerate}

Dejando de lado esta clasificación tradicional, en este capítulo, hasta ahora
hemos presentado las oraciones subordinadas en términos de su \termino{estructura
interna}. Nos hemos concentrado en la forma en que se articula la oración
subordinada (con encabezadores conjuntivos o relativos o sin encabezador) y en la
forma que adopta el verbo (flexionada o no finita). Teniendo en cuenta dichos
parámetros, podemos agrupar las oraciones subordinadas del siguiente modo:

\begin{enumerate}[label=(\alph*)]
  \item \termino{Oraciones subordinadas completivas} o \termino{conjuntivas}
        (OC), encabezadas por conjunciones subordinantes y con flexión verbal
        finita.
  \item \termino{Oraciones subordinadas relativas} (OR), encabezadas por
        pronombres relativos y con flexión verbal finita.
  \item \termino{Oraciones subordinadas sin encabezadores léxicos explícitos}
        (OS). Son las oraciones subordinadas interrogativas o exclamativas
        indirectas (con foco interrogativo o exclamativo), las oraciones
        subordinadas de \forma{que} facultativo (en las que se puede reponer la
        conjunción \forma{que}) y las oraciones subordinadas con formas verbales
        no finitas (infinitivo, gerundio, participio).
\end{enumerate}

\begin{recuadro}[Para ampliar]
Por fuera de las regularidades estructurales marcadas para las oraciones
subordinadas (con o sin encabezadores léxicos explícitos más flexión verbal), se
incluyen como subordinadas los siguientes tipos, que constituyen objetos
lingüísticos que no responden a estructuras oracionales como las que venimos
viendo hasta ahora:

\begin{enumerate}[label=(\alph*)]
  \item Las \textit{oraciones mencionadas}, que son las que normalmente se
        destacan en la escritura por medio de comillas (como en \forma{El
        ministro de Economía manifestó: «\textbf{El que apuesta al dólar
        pierde}»}) y pueden ser tanto oraciones completas como fragmentos
        oracionales (como en \forma{El ministro de Economía respondió: «En
        Argentina»}).
  \item Las \textit{fórmulas sintácticas}, que responden a esquemas fijos de
        distinto tipo (como en \forma{\textbf{Sea quien sea}, no contesto
        llamados}). Estas fórmulas pueden tener distintos formatos, alternativos
        y no predecibles (como en \forma{\textbf{Sea como sea}, ya terminaron las
        clases} y en \forma{\textbf{Como sea}, ya terminaron las clases}, o en
        \forma{\textbf{Con todo y estando enfermo}, asistió a clase} y en
        \forma{\textbf{Con todo}, asistió a clase}).
\end{enumerate}
\end{recuadro}

\subsection{Grado de integración de las subordinadas y funciones argumentales y no argumentales}\label{sec:13-integracion}

Las funciones sintácticas son dependencias estructurales que se establecen entre
un núcleo y un constituyente subordinado a él. En el marco de la oración simple,
se distingue entre funciones argumentales (exigidas por un núcleo léxico) y no
argumentales (compatibles con un núcleo, aunque no requeridas por él). Esta misma
caracterización se recoge en el marco de la oración compuesta, donde las
funciones sintácticas de las subordinadas están determinadas por el tipo de
relación con la oración matriz:

\begin{enumerate}[label=(\alph*)]
  \item La oración subordinada puede \textit{estar integrada} dentro de la
        oración matriz, donde cumple una función argumental o actúa como
        predicativo.
  \item La oración subordinada puede \textit{modificar} a la oración matriz,
        donde cumple una función no argumental.
\end{enumerate}

Importa diferenciar ambos tipos de relaciones (integración y modificación) para
comprender que, si bien la subordinación como vínculo interoracional es una, su
proyección en la estructura de la oración compuesta puede ser diversa. En unos
casos, la oración subordinada contribuye con la buena formación de la oración
matriz, que sin ella no estaría completa, mientras que en otros, su participación
dentro de la oración matriz es más periférica. Todas las oraciones subordinadas,
tanto las completivas como las relativas, pueden cumplir funciones sintácticas
argumentales y no argumentales.

En suma, las oraciones subordinadas cumplen funciones sintácticas que se pueden
agrupar por su aporte a la estructura oracional y por el nivel sintagmático en
el que se proyectan:

\begin{enumerate}[label=(\alph*)]
  \item \textit{Funciones sintácticas argumentales y predicativas}. Las oraciones
        subordinadas completivas y relativas sin antecedente léxico pueden
        cumplir funciones argumentales (sujeto de la oración, complemento del
        verbo, complemento o término de preposición) y también pueden ser
        predicativos (subjetivos u objetivos).
  \item \textit{Funciones sintácticas no argumentales}:
        \begin{enumerate}[label=(b.\arabic*)]
          \item \textit{Modificadoras del núcleo nominal o adverbial}. Las
                oraciones subordinadas relativas con antecedente léxico siempre
                cumplen funciones no argumentales de modificación del nombre o
                del adverbio (adjunto nominal o adverbial).
          \item \textit{Modificadoras del núcleo verbal}. Las oraciones
                subordinadas relativas sin antecedente léxico pueden cumplir
                funciones no argumentales respecto del verbo en el predicado
                (adjunto circunstancial).
          \item \textit{Modificadoras del núcleo oracional}. Las oraciones
                subordinadas completivas, complemento de preposición o de
                adverbio, pueden formar construcciones que modifican el
                contenido proposicional de la oración (adjunto modificador de
                núcleo oracional, MNO).
          \item \textit{Modificadoras de la modalidad oracional}. Las oraciones
                subordinadas completivas, complemento de preposición o de
                adverbio, pueden formar construcciones que modifican la
                modalidad (adjunto modificador de modalidad, MMO).
          \item \textit{Modificadoras del grado}. Las oraciones subordinadas
                completivas pueden cumplir funciones de modificación
                (complementación o adjunción, según la perspectiva que se adopte)
                correlativas de cuantificadores de intensificación en estructuras
                cuantificadas complejas (coda comparativa y coda consecutiva).
        \end{enumerate}
\end{enumerate}

\subsection{Funciones de las oraciones subordinadas completivas}\label{sec:13-completivas}

Las oraciones subordinadas completivas, que son introducidas por una conjunción
subordinante y contienen un verbo en forma finita, las oraciones subordinadas
interrogativas y exclamativas indirectas y las oraciones subordinadas de
infinitivo pueden cumplir funciones argumentales respecto del verbo (sujeto y
complemento objeto directo) o funciones argumentales respecto de la preposición
(complemento o término). Se trata de funciones características de los nombres o
sintagmas nominales, por lo que tradicionalmente se las denominó
\textit{sustantivas}.

\begin{exe}
  \ex\label{ex:13-suj} Sujeto de la oración (argumento verbal):
    \begin{xlist}
      \ex\label{ex:13-suj-a} Me disgusta [\etq{OC/S} que plantees problemas].
          (Me disgusta \textit{tu planteo}).
      \ex\label{ex:13-suj-b} Finalmente se conoce [\etq{OC/S} qué cifra alcanzó
          la inflación el mes pasado]. (Finalmente se conoce \textit{la cifra de
          inflación del mes pasado}).
      \ex\label{ex:13-suj-c} Es recomendable [\etq{OC/S} salir temprano para la
          costa]. (Es recomendable \textit{una salida temprana}).
    \end{xlist}
  \ex\label{ex:13-cod} Complemento objeto directo (argumento verbal):
    \begin{xlist}
      \ex\label{ex:13-cod-a} Estamos esperando [\etq{OC/COD} que plantees
          problemas]. (Estamos esperando \textit{tus planteos}).
      \ex\label{ex:13-cod-b} El Presidente desconoce [\etq{OC/COD} qué cifra
          alcanzó la inflación el mes pasado]. (El Presidente desconoce \textit{la
          cifra de inflación del mes pasado}).
      \ex\label{ex:13-cod-c} Ayer decidimos [\etq{OC/COD} salir temprano para la
          costa]. (Ayer decidimos \textit{una salida temprana}).
    \end{xlist}
  \ex\label{ex:13-ct} Complemento o término de preposición (argumento
      preposicional):
    \begin{xlist}
      \ex\label{ex:13-ct-a} Mi respuesta va a depender [de [\etq{OC/CT} si
          planteás problemas]]. (Mi respuesta va a depender de \textit{tus
          planteos}).
      \ex\label{ex:13-ct-b} Ahora dan explicaciones [sobre [\etq{OC/CT} cuál fue
          la inflación del mes pasado]]. (Ahora dan explicaciones sobre \textit{la
          inflación del mes pasado}).
      \ex\label{ex:13-ct-c} Al final optamos [por [\etq{OC/CT} salir temprano para
          la costa]]. (Al final optamos por \textit{una salida temprana}).
    \end{xlist}
\end{exe}

En las oraciones de (a), la subordinada introducida por una conjunción cumple las
funciones argumentales de sujeto \exref{ex:13-suj-a}, complemento objeto directo
\exref{ex:13-cod-a} y complemento o término de preposición \exref{ex:13-ct-a}.
Asimismo, en las oraciones de (b), la subordinada interrogativa indirecta cumple
las funciones argumentales de sujeto \exref{ex:13-suj-b}, complemento objeto
directo \exref{ex:13-cod-b} y complemento o término de preposición
\exref{ex:13-ct-b}. Finalmente, en las oraciones de (c), la subordinada de
infinitivo cumple las mismas funciones argumentales de sujeto
\exref{ex:13-suj-c}, complemento objeto directo \exref{ex:13-cod-c} y complemento
o término de preposición \exref{ex:13-ct-c}. En todos los casos se ilustra entre
paréntesis la equivalencia nominal de tales oraciones subordinadas.

Aunque las oraciones subordinadas, según se ve, pueden cumplir funciones
sintácticas idénticas a las de los nombres o sintagmas nominales, esto no
significa que se trate de sustantivos, ya que existen diferencias interpretativas
entre las oraciones subordinadas y los nombres. En este sentido, las
subordinadas completivas, por su valor modal, pueden ser enunciativas, como en
\exref{ex:13-enun-a}, \exref{ex:13-enun-b} y \exref{ex:13-enun-c}, exhortativas
\exref{ex:13-enun-d}, interrogativas \exref{ex:13-indir}, tanto
totales\footnote{Las interrogativas totales o disyuntivas se caracterizan por
presentar opciones entre las que se debe elegir: \forma{¿Querés ir al cine (o
no)?; ¿Querés ir al cine o al teatro?} Las interrogativas parciales o
pronominales van encabezadas por pronombres o adverbios interrogativos que son el
foco de la pregunta sobre el que se debe responder: \forma{¿Dónde querés ir esta
noche?; ¿Quién te lo dijo?}} \exref{ex:13-indir-a} como parciales,
\exref{ex:13-indir-b} y \exref{ex:13-indir-c}, o
exclamativas\footnote{\label{fn:13-exclamativas}Las oraciones exclamativas son
siempre parciales, ya que se inician con un ítem exclamativo: \forma{¡Qué
barbaridad!; ¡Cuánta injusticia!}} \exref{ex:13-indign}, según los verbos que las
seleccionan y los valores semánticos que aportan. Esto marca una diferencia con
los sintagmas nominales, ya que las subordinadas son seleccionadas en función de
sus valores modales; en cambio, los nombres, que carecen de estos valores, son
seleccionados en función de su semántica.

\begin{exe}
  \ex Oraciones subordinadas completivas enunciativas y exhortativas:
    \begin{xlist}
      \ex\label{ex:13-enun-a} El estudiante responde [\etq{OC/COD} que los
          ángulos internos de un triángulo suman 180°].
      \ex\label{ex:13-enun-b} El estudiante espera [\etq{OC/COD} que los ángulos
          internos de un triángulo sumen 180°].
      \ex\label{ex:13-enun-c} El estudiante desea [\etq{OC/COD} aprobar el
          examen de matemática].
      \ex\label{ex:13-enun-d} El estudiante solicita [\etq{OC/COD} que se
          prorroguen las clases].
    \end{xlist}
  \ex\label{ex:13-indir} Oraciones subordinadas completivas interrogativas
      indirectas (totales y parciales):
    \begin{xlist}
      \ex\label{ex:13-indir-a} El maestro pregunta [\etq{OC/COD} si los ángulos
          internos de un triángulo suman 180°].
      \ex\label{ex:13-indir-b} El maestro pregunta [\etq{OC/COD} cuánto suman los
          ángulos internos de un triángulo].
      \ex\label{ex:13-indir-c} El maestro desconoce [\etq{OC/COD} hasta cuándo se
          prorrogarán las clases].
    \end{xlist}
  \ex\label{ex:13-indign} Oración subordinada completiva exclamativa indirecta
      (parcial): Es indignante [\etq{OC/S} cuánto demora el trámite de título].
\end{exe}

Las oraciones subordinadas completivas también pueden cumplir la función de
predicativo subjetivo \exref{ex:13-ps} u objetivo \exref{ex:13-po}.

\begin{exe}
  \ex\label{ex:13-ps} Predicativo subjetivo (predicado semántico):
    \begin{xlist}
      \ex Mi esperanza es [\etq{OC/PS} que no surjan problemas].
      \ex La duda es [\etq{OC/PS} qué pasa con el feriado del mes próximo].
      \ex El objetivo de la política fiscal es [\etq{OC/PS} eliminar los
          impuestos distorsivos].
    \end{xlist}
  \ex\label{ex:13-po} Predicativo objetivo (predicado semántico):
    \begin{xlist}
      \ex Ayer vi a tu hermano [\etq{OC/PO} entrar a la oficina]. (Ayer vi
          \textit{que tu hermano entraba a la oficina}).
      \ex En la esquina hay una cuadrilla de operarios [\etq{OC/PO} rompiendo la
          vereda].
    \end{xlist}
\end{exe}

\subsection{Funciones de las oraciones subordinadas relativas}\label{sec:13-funcrelativas}

Las oraciones subordinadas relativas (encabezadas por un relativo) pueden
construirse como ligadas a un antecedente léxico explícito de naturaleza nominal
o bien sin antecedente, caso en que se denominan \textit{relativas libres}.

\subsubsection{Relativas con antecedente}\label{sec:13-conantecedente}

Las relativas ligadas a un antecedente pueden cumplir funciones no argumentales
como modificadoras del nombre en calidad de adjuntos, que pueden ser
\termino{especificativos} o \termino{explicativos}. Los adjuntos especificativos,
como los de \exref{ex:13-as}, restringen la semántica del nombre al que
modifican; en cambio, los adjuntos explicativos, como los de \exref{ex:13-ax},
amplían y detallan la semántica del sintagma nominal al que se adjuntan. En ambos
casos, las funciones que cumplen son características de los adjetivos o
sintagmas adjetivales.

\begin{exe}
  \ex\label{ex:13-as} Adjunto nominal especificativo:
    \begin{xlist}
      \ex\label{ex:13-as-a} La policía detuvo a dos delincuentes [\etq{OR/AS} que
          atracaron un banco]. (La policía detuvo a dos delincuentes. \textit{Dos
          delincuentes} atracaron un banco).
      \ex\label{ex:13-as-b} Las personas [\etq{OR/AS} a quienes se detuvo]
          negaron su responsabilidad en el hecho. (Las personas negaron su
          responsabilidad en el hecho. Se detuvo \textit{a esas personas}).
      \ex\label{ex:13-as-c} Dijeron no conocer el banco [\etq{OR/AS} donde se
          produjo el atraco]. (Dijeron no conocer el banco. \textit{En el banco}
          se produjo el atraco).
      \ex\label{ex:13-as-d} Añadieron que para la fecha [\etq{OR/AS} cuando se
          produjo el hecho] estaban fuera de la ciudad. (Añadieron que para la
          fecha estaban fuera de la ciudad. \textit{En la fecha} se produjo el
          hecho).
    \end{xlist}
  \ex\label{ex:13-ax} Adjunto nominal explicativo:
    \begin{xlist}
      \ex\label{ex:13-ax-a} La causa, [\etq{OR/AX} la cual quedó a cargo del
          juez], suma testimonios de varios testigos. (La causa suma testimonios
          de varios testigos. \textit{La causa} quedó a cargo del juez).
      \ex\label{ex:13-ax-b} Según un testigo, [\etq{OR/AX} cuyo testimonio se
          considera clave], los ladrones fueron vistos en el banco. (Según un
          testigo, los ladrones fueron vistos en el banco. \textit{Su}
          testimonio se considera clave).
    \end{xlist}
\end{exe}

En ambas series, el relativo que encabeza la oración subordinada puede tener
funciones nominales, como en \exref{ex:13-as-a}, \exref{ex:13-as-b} y
\exref{ex:13-ax-a}, adverbiales, como en \exref{ex:13-as-c} y
\exref{ex:13-as-d}, o de determinante, como en \exref{ex:13-ax-b}. En todos los
casos, el relativo liga su interpretación al antecedente nominal (es decir, se
carga de significado a partir del nombre o sintagma nominal que lo antecede).

\subsubsection{Relativas libres}\label{sec:13-libres}

Las oraciones subordinadas relativas libres, encabezadas por un relativo no
ligado a un antecedente léxico explícito, pueden cumplir funciones nominales o
adverbiales.

\paragraph{Relativas libres con funciones adverbiales.}\label{sec:13-libres-adv}
Las relativas libres encabezadas por los adverbios \forma{donde, cuanto, cuando}
y \forma{como} pueden desempeñarse, dentro del predicado oracional, como
circunstanciales, argumentales o no argumentales, respecto del verbo, como se
presenta en \exref{ex:13-circ}. Las funciones que cumplen son características de
los sintagmas adverbiales y preposicionales.

\begin{exe}
  \ex\label{ex:13-circ} Circunstancial (complemento o adjunto verbal):
    \begin{xlist}
      \ex\label{ex:13-circ-a} Corrió en la cinta [\etq{OR/AC} cuanto le dieron
          sus piernas]. (Cantidad).
      \ex\label{ex:13-circ-b} Terminó la primaria [\etq{OR/AC} donde la había
          empezado]. (Lugar).
      \ex\label{ex:13-circ-c} Habla inglés [\etq{OR/AC} como le enseñaron en la
          casa]. (Modo o manera).
      \ex\label{ex:13-circ-d} Salió de su casa [\etq{OR/AC} cuando lo llamaron del
          campo]. (Tiempo).
    \end{xlist}
\end{exe}

Las oraciones de \exref{ex:13-circ} se corresponden con las relativas libres que,
en la tradición gramatical, fueron denominadas \textit{subordinadas adverbiales
propias} porque podían ser reemplazadas por adverbios. Así, podemos ver que las
relativas libres funcionan como adverbios: en \exref{ex:13-circ-a}, como adjunto
circunstancial de cantidad, encabezada por el relativo \forma{cuanto}; en
\exref{ex:13-circ-b}, como adjunto circunstancial de lugar, encabezada por el
relativo \forma{donde}; en \exref{ex:13-circ-c}, como adjunto circunstancial de
modo, encabezada por el relativo \forma{como}; y en \exref{ex:13-circ-d}, como
adjunto circunstancial de tiempo, encabezada por el relativo
\forma{cuando}.\footnote{En este apartado solo consideramos las subordinadas
relativas introducidas por \forma{cuando}; el resto de las construcciones que
expresan \sig{tiempo} se ven en §\ref{sec:13-temporales}.}

Es de notar que los mismos adverbios relativos libres de los ejemplos de
\exref{ex:13-circ} pueden estar ligados a un antecedente léxico explícito
(\forma{todo, escuela, manera, momento}), como en \exref{ex:13-circant}. En estos
casos la construcción subordinada se convierte en una relativa con antecedente
que funciona como adjunto especificativo. Asimismo, se advierte una
correspondencia con una serie de adverbios no relativos (\forma{mucho, allí, así,
entonces}), funcionalmente equivalentes a las respectivas oraciones subordinadas
relativas libres.

\begin{exe}
  \ex\label{ex:13-circant} Relativas con antecedente en estructuras con función
      de circunstancial:
    \begin{xlist}
      \ex Corrió en la cinta todo \textit{cuanto le dieron sus piernas}. Corrió
          \textit{mucho} en la cinta.
      \ex Terminó la primaria en la escuela \textit{donde la había empezado}. La
          terminó \textit{allí}.
      \ex Habla inglés de la manera \textit{como le enseñaron en la casa}. Lo
          habla \textit{así}.
      \ex Salió de su casa en el momento \textit{cuando lo llamaron del campo}.
          Salió \textit{entonces}.
    \end{xlist}
\end{exe}

\paragraph{Relativas libres con funciones nominales.}\label{sec:13-libres-nom}
Las relativas libres (encabezadas por el pronombre relativo \forma{quien} y su
variante plural) y semilibres (encabezadas por el pronombre relativo \forma{el
que} y sus variantes de género y número) pueden cumplir funciones argumentales
características de los nombres o sintagmas nominales, relacionadas con el verbo
(sujeto y complemento objeto directo) y la preposición (complemento o término).

\begin{exe}
  \ex\label{ex:13-abarca-s} Sujeto de la oración (argumento verbal):
    \begin{xlist}
      \ex\label{ex:13-abarca-s-a} {}[\etq{OR/S} Quien mucho abarca], poco
          aprieta. (\textit{La persona que mucho abarca}, poco aprieta).
      \ex\label{ex:13-abarca-s-b} {}[\etq{OR/S} Los que mucho abarcan], poco
          aprietan. (\textit{Los hombres que mucho abarcan}, poco aprietan).
    \end{xlist}
  \ex\label{ex:13-abarca-cod} Complemento objeto directo (argumento verbal):
    \begin{xlist}
      \ex\label{ex:13-abarca-cod-a} A [\etq{OR/COD} quien mucho abarca] en nada
          se lo considera idóneo. (A \textit{la persona que mucho abarca} en nada
          se la considera idónea).
      \ex\label{ex:13-abarca-cod-b} A [\etq{OR/COD} las que mucho abarcan] en nada
          se las considera idóneas. (A \textit{las personas que mucho abarcan} en
          nada se las considera idóneas).
    \end{xlist}
  \ex\label{ex:13-abarca-ct} Complemento o término de preposición (argumento
      preposicional):
    \begin{xlist}
      \ex\label{ex:13-abarca-ct-a} Todos se alejan [de [\etq{OR/CT} quienes mucho
          abarcan]]. (Todos se alejan de \textit{las personas que mucho
          abarcan}).
      \ex\label{ex:13-abarca-ct-b} Todos se alejan [de [\etq{OR/CT} la que mucho
          abarca]]. (Todos se alejan de \textit{la persona que mucho abarca}).
    \end{xlist}
\end{exe}

En las oraciones (a), la relativa libre cumple las funciones argumentales de
sujeto \exref{ex:13-abarca-s-a}, complemento objeto directo
\exref{ex:13-abarca-cod-a} y complemento o término de preposición
\exref{ex:13-abarca-ct-a}. El mismo esquema de funciones sintácticas se replica
con la relativa semilibre en la serie (b): sujeto \exref{ex:13-abarca-s-b},
complemento objeto directo \exref{ex:13-abarca-cod-b} y complemento o término de
preposición \exref{ex:13-abarca-ct-b}. En todos estos casos, las relativas
muestran una equivalencia funcional con sintagmas nominales encabezados por un
nombre general con rasgo [+humano] (\forma{hombre, mujer, persona}) o con
pronombres personales (\forma{él, ella, ellas, ellos}), hecho que normalmente
deriva en su consideración como \termino{relativas sustantivadas}.

En relación con su valor sustantivado conviene señalar que, aun siendo todas
oraciones subordinadas, existen diferencias interpretativas y sintácticas,
además de diferencias estructurales, entre las completivas y las relativas
sustantivadas. Así, las oraciones subordinadas completivas (OC) denotan nociones
abstractas (contenidos proposicionales) y se asimilan a oraciones
\exref{ex:13-esto}; en cambio, las relativas (OR) libres, aunque también pueden
cumplir funciones argumentales, aluden a cualquier referente y se asimilan a
nombres o pronombres \exref{ex:13-este}:

\begin{exe}
  \ex\label{ex:13-esto} Oraciones subordinadas completivas (OC), asimilables a
      oraciones, conmutables por el demostrativo neutro:
    \begin{xlist}
      \ex Me disgusta [\etq{OC/S} que plantees problemas]. Me disgusta
          \textit{esto}. (Pero *Me disgusta \textit{este / él}).
      \ex Finalmente se conoce [\etq{OC/S} qué cifra alcanzó la inflación el mes
          pasado]. Finalmente se conoce \textit{esto}. (Pero *Finalmente se conoce
          \textit{esta / ella}).
      \ex Es recomendable [\etq{OC/S} salir temprano para la costa]. Es
          recomendable \textit{esto}. (Pero *Es recomendable \textit{este / él}).
    \end{xlist}
  \ex\label{ex:13-este} Oraciones subordinadas relativas (OR), asimilables a
      nombres o pronombres, conmutables por pronombres demostrativos o
      personales:
    \begin{xlist}
      \ex {}[\etq{OR/S} Quien mal anda], mal acaba. \textit{Este / Él} mal acaba.
          (Pero *\textit{Esto} mal acaba). \textit{La persona que mal anda} mal
          acaba.
      \ex {}[\etq{OR/S} Lo que mal empieza], mal acaba. \textit{Esto} mal acaba.
          (Pero *\textit{Este / Él} mal acaba). \textit{Aquello que mal empieza}
          mal acaba.
    \end{xlist}
\end{exe}

Las mismas oraciones subordinadas relativas libres y semilibres (pronombres
relativos \forma{quien} y \forma{el que}, con sus variantes) pueden cumplir la
función de predicativo subjetivo.

\begin{exe}
  \ex\label{ex:13-predrel} Predicativo subjetivo (predicado semántico):
    \begin{xlist}
      \ex\label{ex:13-predrel-a} ¿Usted es [\etq{OR/PS} quien siempre avisa de
          los cortes de servicio a EDESUR]? (¿Usted es \textit{la persona que
          siempre avisa de los cortes de servicio a EDESUR}?)
      \ex\label{ex:13-predrel-b} Yo soy [\etq{OR/PS} la que lleva esta casa
          adelante]. (Yo soy \textit{la mujer que lleva esta casa adelante}).
    \end{xlist}
\end{exe}

En las oraciones de \exref{ex:13-predrel}, la relativa libre
\exref{ex:13-predrel-a} y la semilibre \exref{ex:13-predrel-b} cumplen la
función sintáctica de predicativo subjetivo del verbo copulativo \forma{ser}. El
pronombre relativo replica la referencia presente en el contexto oracional
cuando el verbo de la relativa muestra una concordancia diferente de la tercera
persona, como en \forma{Yo soy el que soy}.

\begin{ejercicio}\label{ej:13-funciones}
Volviendo sobre las paremias del ejercicio~\ref{ej:13-paremias}, identifique las
funciones sintácticas que cumplen las oraciones subordinadas reconocidas en
ellas, según se muestra en los ejemplos resueltos.

\begin{resuelto}[Ejemplos resueltos]
1. [\etq{Ocomp} A la cama no te irás \textit{sin} [\etq{OC/CT} \textit{saber una
cosa más}]]: la OC de infinitivo (sin subordinante) es el complemento o término
(CT) de la preposición \forma{sin} (valor de polarización negativa: \sig{no
sabiendo una cosa más}). 2. [\etq{Ocomp} A [\etq{OR/COD} \textit{quien
madruga}], Dios lo ayuda]: la OR libre encabezada por el relativo \forma{quien}
(\sig{alguien madruga}) es el complemento objeto directo (COD) del verbo
transitivo \forma{ayudar}.
\end{resuelto}
\end{ejercicio}

\begin{ejercicio}\label{ej:13-relativos}
En las siguientes oraciones compuestas, según se muestra en los ejemplos
resueltos, identifique el relativo y la función sintáctica que cumple dentro de
la oración subordinada y, donde corresponda, señale el antecedente del relativo.

\begin{exe}
  \ex Ayudemos a quienes lo necesiten.
  \ex Compré una valija. De hecho, compré la que me recomendaste.
  \ex Conocí a la persona para quien trabajó mi padre.
  \ex Decime lo que pensás.
  \ex Dieron un premio a quienes participaron.
  \ex El coche que está estacionado allí es el mío.
  \ex El libro que me prestaste es muy interesante.
  \ex El que llegó tarde no entró.
  \ex Esa es la casa cuya puerta estaba pintada de verde.
  \ex Esta es la lapicera con la cual escribí la carta.
  \ex Hizo la preparación como le enseñaron.
  \ex La docente con quien estudié inglés era británica.
  \ex La película por la cual vinimos al festival fue excelente.
  \ex Lo voy a comprar cuando esté en oferta.
  \ex Los estudiantes que estudiaron aprobaron el examen.
  \ex Los estudiantes, que estudiaron, aprobaron el examen.
  \ex Necesito lo que preparaste para el relleno de la tarta.
  \ex No me importa lo que digan.
  \ex Quien persevera, vence.
  \ex Quienes estudian, aprenden.
\end{exe}

\begin{resuelto}[Ejemplos resueltos]
1. [\etq{Ocomp} Ayudemos a [\etq{OR/COD} \textit{quienes} lo necesiten]]:
\forma{quienes}, sujeto de la OR; sin antecedente. 3. [\etq{Ocomp} Conocí a la
persona [\etq{OR/AS} para \textit{quien} trabajó mi padre]]: \forma{quien},
complemento o término de la preposición \forma{para}; antecedente,
\forma{persona}.\pendiente{El original atribuye a \textit{para} «valor
sociativo»}
\end{resuelto}
\end{ejercicio}

\section{Las construcciones con valores adverbiales}\label{sec:13-adverbiales}

Como ya hemos adelantado, partiendo del supuesto de la equivalencia funcional con
el adverbio, la gramática tradicional agrupaba, bajo la etiqueta de
\textit{subordinadas adverbiales}, construcciones sintáctica y semánticamente muy
diversas. En este capítulo, en cambio, hemos privilegiado un enfoque que parte de
la conformación interna de las subordinadas. Por tanto, las consideradas
tradicionalmente \textit{adverbiales propias} (temporales de \forma{cuando},
locativas, modales y de cantidad), que podían ser sustituidas por un adverbio de
significado equivalente, han sido incluidas entre las relativas libres por estar
encabezadas por adverbios relativos (respectivamente \forma{cuando, donde, como}
y \forma{cuanto}; véase §\ref{sec:13-libres}).

En este apartado nos vamos a ocupar de las temporales no relativas y de las que
recibían la denominación de \textit{adverbiales impropias} (causales, finales,
condicionales, concesivas, comparativas y consecutivas), a partir de que se
reconocía que no existen adverbios por los que puedan ser reemplazadas. En todos
estos casos, según el enfoque analítico que propiciamos, se trata de estructuras
encabezadas por conjunciones subordinantes (\forma{que, si}), que introducen
oraciones subordinadas completivas, que son seleccionadas por adverbios (p.~ej.,
\forma{aun, antes} o \forma{siempre}), preposiciones (p.~ej., \forma{para, por,
de}) y locuciones (\forma{dado que, puesto que}) o bien resultan correlativas de
elementos intensificativos presentes en la oración principal (\forma{tanto, más,
menos}). En este sentido, vamos a referirnos a todas ellas como
\termino{construcciones con valores adverbiales}, denominación amplia que
utilizamos para el conjunto formado por el adverbio, la preposición o locución
más la subordinada, que en conjunto forman un sintagma adverbial o
preposicional, según los casos. En cuanto a las comparativas y consecutivas, que
«están entre las que presentan mayor complejidad sintáctica de toda la gramática
española» \parencite[§\,45.1a]{rae2025}, la construcción abarca la estructura
formada por el intensificador, el núcleo al que este especifica y la subordinada
completiva introducida por \forma{que}: \forma{Esa foto tuya es [más actual que
la que está en el álbum]}. El valor semántico (temporal, causal, condicional,
comparativo, etc.) que toma la construcción lo aporta el adverbio, la
preposición, una locución o bien la correlación con el elemento intensificador.
Así, a modo de ejemplo, podemos mencionar casos como \forma{antes (de), después
(de), siempre} (temporal); \forma{por, a causa de, debido a} (causal); \forma{aun,
a pesar de, pese a} (concesivo); \forma{en caso de} (condicional); \forma{para, a
fin de, con vistas a} (final); \forma{más\ldots\ que\ldots} (comparativo);
\forma{tanto\ldots\ que} (consecutivo).

Antes de ver más en detalle los distintos tipos de construcciones con valores
adverbiales, vamos a presentar un enfoque que, desde la perspectiva semántica,
nos permite introducir una propuesta integradora que supere la atomización de
valores con que tradicionalmente se presentan estas subordinadas, a partir del
reconocimiento de las operaciones básicas que subyacen a las conexiones que
establecen estas construcciones.

\subsection{Una perspectiva integradora}\label{sec:13-integradora}

Partimos de la idea de que la oración constituye la representación de un evento
o estado de cosas del mundo exterior (como en \forma{Llueve}) o interior (como en
\forma{Tengo frío}). Cuando combinamos oraciones, esas representaciones se
entrelazan en la mente del hablante, lo que lleva a que las exprese combinadas
secuencialmente (como en \forma{Llueve; tengo frío}), ya sea en estructuras
independientes, coordinadas (como en \forma{Llueve y tengo frío}) o integradas
mediante subordinación (como en \forma{Tengo frío porque llueve}).

Según plantea \textcite[4]{duque2025}, «todos los hablantes comprendemos el
mundo y los textos mediante comparaciones, generalizaciones u ordenando la
información en función de relaciones temporales o causales». Estas relaciones
resultan intuitivamente simples y, por lo tanto, sumamente accesibles porque se
fundamentan en operaciones básicas. En esta perspectiva,
\textcite{rudolph1988,rudolph1996} propuso un esquema explicativo que permite dar
cuenta de las relaciones semánticas fundamentales que manifestamos al conectar
nuestras ideas mediante la combinación de oraciones. Para esta autora
\parencite[20]{rudolph1996}, existen cuatro formas básicas de conexión:

\begin{description}
  \item[Adición {[A + B]}:] dos proposiciones A y B están estrechamente
        relacionadas y son válidas simultáneamente.
  \item[Contraste {[A $-$ B]}:] para el hablante, dos proposiciones A y B son
        válidas simultáneamente y la proposición B marca un contraste con la
        proposición A.
  \item[Tiempo {[A, B]}:] dos proposiciones A y B no son válidas simultáneamente,
        y la proposición B denota un evento que sigue en el tiempo a la
        proposición A.
  \item[Causalidad {[A > B]}:] para el hablante, dos proposiciones A y B pueden
        ser válidas simultáneamente o no, y están ligadas en una cadena causal o
        de necesidad lógica o empírica.
\end{description}

Estas cuatro formas básicas de conexión se emplean para vincular eventos o
estados de cosas. Por un lado, las relaciones de adición y tiempo permiten al
hablante organizar la presentación de los hechos. Por otro lado, las relaciones
de contraste y causalidad hacen posible que, además de establecer un orden entre
los sucesos, el hablante también exprese sus opiniones y puntos de vista sobre
los hechos, estableciendo entre ellos semejanzas y diferencias o relaciones de
causa-efecto. Así, bajo la perspectiva aquí descripta, las construcciones que
manifiestan relaciones lógicas, como las temporales, condicionales, concesivas,
causales y finales, y las que expresan valores cuantificacionales-intensificativos,
como las comparativas y consecutivas, entre las que las descripciones
gramaticales tradicionales no solían establecer conexiones, se resuelven en
relaciones más básicas que van de la simple conexión y ordenación de los eventos
(simultaneidad o secuencia) a la manifestación de contrastes (oposiciones y
comparaciones) y diferentes tipos de causalidades (causa real, final, virtual e
ineficiente) entre sucesos.

Según plantea \textcite{rudolph1988}, los valores semánticos que expresan estas
cuatro formas básicas de conexión interoracional se pueden proyectar mediante la
simple secuenciación de oraciones independientes, es decir, sin utilizar
conectores léxicos. Sin embargo, el empleo de conjunciones y locuciones
conjuntivas que hagan explícitas las relaciones lógicas, discursivas y
argumentativas que contraen las oraciones constituye una excelente orientación
para el oyente en cuanto a la interpretación de los textos. En este sentido, los
conectores y conjunciones de argumentación (aditivos, adversativos, concesivos,
consecutivos), de reformulación (reformuladores, explicativos, rectificativos,
recapitulativos) y de estructuración (organizativos, ejemplificativos,
digresivos) presentados en el \capref{12}, §\ref{sec:12-marcadores}, funcionan
como formas lingüísticas especializadas en expresar esos señalamientos dirigidos
al receptor.

Las relaciones subordinativas (y también las coordinativas, vistas en el
\capref{12}), que involucran la combinación de unidades de nivel oracional, son
responsables de dar unidad y sentido a los textos estableciendo conexiones entre
las ideas. Desde esta perspectiva de integración, a continuación vamos a
presentar dos subtipos de construcciones con valores adverbiales: a) las que
manifiestan valores lógico-relacionales (temporales, causales, finales,
condicionales y concesivas) y b) las que manifiestan valores
cuantificacionales-intensificativos (comparativas y consecutivas).

\subsection{Construcciones con valores lógico-relacionales}\label{sec:13-logicas}

Según vimos, las oraciones subordinadas con funciones argumentales forman parte
de una estructura compleja, en la que completan el valor requerido por un núcleo
(verbal o preposicional) del que dependen. En cuanto a las oraciones relativas
ligadas a un antecedente nominal, contribuyen semánticamente a la restricción
(relativas especificativas) o a la explicación (relativas explicativas) de la
referencia, mientras que las relativas libres y semilibres introducen referentes
discursivos de naturaleza diversa (cantidad, lugar, manera, persona, tiempo). Es
así que, como plantea \textcite[327]{ditullio2014}, las consideradas
\textit{adverbiales propias} (temporales relativas, locativas, modales y de
cantidad) «denotan entidades individuales específicas como los momentos y los
lugares, así como también las maneras en que puede realizarse una cierta acción o
el grado o la intensidad con que se lleva a cabo». Dentro de este panorama, las
temporales introducidas por adverbios o preposiciones no relativas y las
tradicionalmente consideradas \textit{adverbiales impropias} ---causales,
finales, condicionales y concesivas--- suman valores que no dependen de la
estructura argumental ni del establecimiento de la referencia nominal, sino que
derivan de las relaciones lógicas de temporalidad y causalidad establecidas entre
los mundos representados en la oración principal y la subordinada.

Las construcciones con valores lógico-relacionales (temporales, causales,
finales, condicionales, concesivos) funcionan como modificadores adjuntos, ya sea
en el nivel del predicado (adjuntos circunstanciales) o en el nivel oracional
(modificadores del núcleo oracional y modificadores de la modalidad oracional).
Asimismo, establecen correlaciones temporales y modales con la oración principal
para orientar, junto con conectores y elementos conjuntivos, la interpretación
global de la oración compuesta. En este sentido, por ejemplo, la constitución de
un acto de habla compromisorio (advertencia o amenaza) puede derivar de la
interpretación condicional que surge en cualquiera de los ejemplos de
\exref{ex:13-gritando}.

\begin{exe}
  \ex\label{ex:13-gritando} Correlaciones temporales y modales más conectores y
      elementos conjuntivos:
    \begin{xlist}
      \ex\label{ex:13-gritando-a} \textit{Si me seguís gritando en la calle}, voy
          y te hago una denuncia.
      \ex\label{ex:13-gritando-b} \textit{Como me sigas gritando en la calle},
          voy y te hago una denuncia.
      \ex\label{ex:13-gritando-c} \textit{Donde me sigas gritando en la calle},
          voy y te hago la denuncia.
      \ex\label{ex:13-gritando-d} \textit{Gritándome cosas en la calle}, solo vas
          a conseguir que te haga la denuncia.
    \end{xlist}
\end{exe}

Esto es posible gracias, por un lado, a la presencia de conjunciones, ya sea la
típica de las condicionales, \forma{si} \exref{ex:13-gritando-a}, o también
otras, como \forma{como} \exref{ex:13-gritando-b} e incluso \forma{donde}
\exref{ex:13-gritando-c}, que no resultan características de prótasis
condicionales, y, por el otro, al mantenimiento de las correlaciones temporales y
modales entre las oraciones intervinientes. Estos mismos valores
lógico-relacionales se pueden obtener incluso sin encabezadores léxicos, a través
de estructuras oracionales organizadas alrededor de formas no finitas del verbo,
como en \exref{ex:13-gritando-d}.

A continuación, según las funciones sintácticas que pueden cumplir, vamos a
considerar tres subgrupos de construcciones con valores lógico-relacionales: a)
construcciones temporales que funcionan en el nivel del predicado como
circunstanciales (en este apartado exceptuamos las de \forma{cuando}, que ya
fueron incluidas con las relativas); b) causales y finales, que pueden funcionar
en el nivel del predicado como circunstanciales y en el nivel oracional como
modificadoras de la modalidad; c) condicionales y concesivas, que solo funcionan
en el nivel oracional como modificadoras del núcleo oracional o de la modalidad.

\subsubsection{Construcciones temporales}\label{sec:13-temporales}

Las construcciones que veremos en este apartado establecen una relación temporal
entre el evento manifestado en la oración principal y el presentado en la
subordinada. Manifiestan diversos valores semánticos relacionados con la conexión
básica de \sig{tiempo} \parencite{rudolph1996}, los cuales dependen del
significado del adverbio, preposición o locución temporal que las precede o de la
conjunción que las introduce. Pueden expresar simultaneidad
(\forma{mientras}\footnote{La naturaleza relativa o conjuntiva de
\forma{mientras} se encuentra en debate; no obstante, dada la índole
introductoria de este libro, no entraremos en estas cuestiones.}), anterioridad
(\forma{antes (de) que}) / posterioridad (\forma{apenas, después (de) que}),
delimitación (\forma{desde que / hasta que}), reiteración (\forma{siempre que,
cada vez que}), progresión (\forma{a medida que}), entre otros valores posibles.

Las construcciones temporales pueden cumplir funciones circunstanciales dentro
del predicado y, según sean requeridas o no por el verbo, se consideran
argumentales o no, es decir que se desempeñan como complementos o adjuntos
circunstanciales, como en \exref{ex:13-temp}.

\begin{exe}
  \ex\label{ex:13-temp} Valores temporales:
    \begin{xlist}
      \ex Podemos tomar un café \textit{\textbf{mientras} discutimos la propuesta
          laboral}. (Simultaneidad).
      \ex Arreglemos el orden del día \textit{\textbf{antes} de que empiece la
          reunión de directorio}. (Anterioridad).
      \ex Partimos para Mar del Plata \textit{\textbf{después} de que amanezca /
          \textbf{apenas} salga el sol}. (Posterioridad).
      \ex Mi hermano trabaja todo el día, \textit{\textbf{desde} que amanece
          \textbf{hasta} que anochece}. (Delimitación).
      \ex Me pongo a bailar \textit{\textbf{siempre} que escucho este tema}.
          (Reiteración).
      \ex Entendí el funcionamiento del cerebro \textit{\textbf{a medida} que
          avanzaba en la carrera}. (Progresión).
    \end{xlist}
\end{exe}

\begin{recuadro}[Una construcción particular: \textit{hacer} referido a \sig{transcurso de tiempo}]
El verbo \forma{hacer} seguido de un complemento temporal, como en \forma{Se fue
hace dos semanas}, constituye, según Martínez García (1996,
41),\pendiente{Martínez García (1996) no figura en la bibliografía} «una
estructura especial dentro de las subordinadas temporales durativas». En
principio, se trata de la única construcción del español con verbo conjugado en
indicativo que se subordina al verbo de la oración principal sin nexo:

\begin{exe}
  \ex\label{ex:13-hace1} Salieron para el aeropuerto [Ø \textit{hace dos
      horas}].
\end{exe}

Otra característica peculiar de esta construcción, que señala \sig{transcurso de
tiempo}, es que realiza un señalamiento temporal inverso, desde el presente hacia
el pasado, lo que permite explicar que \forma{hacer} en presente pueda
combinarse con verbos en pasado como \forma{salieron} \exref{ex:13-hace1} o
\forma{vivía} \exref{ex:13-hace2}:

\begin{exe}
  \ex\label{ex:13-hace2} \textit{Hace mucho tiempo} vivía en un lejano país.
\end{exe}

Además, el verbo \forma{hacer} conforma un esquema impersonal en tercera persona
singular, que presenta dos construcciones (con y sin \forma{que}) en apariencia
semejantes, aunque entre ellas existen «diferencias tanto semánticas como
sintácticas» \parencite[334]{ditullio2014}:

\begin{exe}
  \ex\label{ex:13-hace-I} \textit{Hace dos horas} salieron para el aeropuerto.
      (Tipo I: \textit{hace}\etq{temp} sin \textit{que}).
  \ex\label{ex:13-hace-II} \textit{Hace dos horas} que salieron para el
      aeropuerto. (Tipo II: \textit{hace}\etq{temp} más \textit{que}).
\end{exe}

Estas diferencias, estudiadas en \textcite{giammatteo2002}, muestran que en el
Tipo~I los constituyentes presentan un comportamiento unitario: no aceptan
permutación interna (\juicio{*}{\forma{Dos horas hace salieron para el
aeropuerto}}), pueden desplazarse conjuntamente dentro de la principal
(\forma{Salieron, hace dos horas, para el aeropuerto; Salieron para el aeropuerto
hace dos horas}), pueden usarse independientemente en respuesta a una pregunta
(\forma{¿Cuándo salieron para el aeropuerto? Hace dos horas} /
\juicio{*}{\forma{Dos horas}}),\pendiente{El original dice «no pueden usarse
independientemente», pero el ejemplo muestra lo contrario} pueden funcionar como
complemento de las preposiciones \forma{de, desde} y \forma{hasta} (\forma{Hasta
hace un rato estaban en casa}) y pueden ser conmutados por un adverbio
(\forma{Recientemente salieron para el aeropuerto}); por otra parte,
\forma{hacer} está prácticamente inmovilizado en presente. Las características
señaladas muestran que esta construcción ha sufrido un proceso de
gramaticalización que implica una descategorización: en este caso, una estructura
originalmente oracional ya no introduce otro evento diferente del de la
principal, sino que pasa a cumplir una función adverbial y a desempeñarse como un
adjunto temporal que sitúa temporalmente el suceso.

En el Tipo~II, el comportamiento de los constituyentes es diferente: pueden
permutarse internamente (\forma{Dos horas hace que salieron para el aeropuerto}),
no pueden desplazarse dentro de la principal ni tampoco omitirse
(\juicio{*}{\forma{Que salieron para el aeropuerto hace dos horas}}), pueden ser
usados independientemente en respuesta a una pregunta (\forma{¿Cuánto hace que
salieron? Dos horas}) y no funcionan como complemento de preposición ni pueden
ser reemplazados por un adverbio (\juicio{*}{\forma{Recientemente que salieron
para el aeropuerto}}); por otra parte, \forma{hacer} admite mucha más variación
en el uso de tiempos verbales.

Sintácticamente, mientras en el Tipo~I, sin \forma{que}, la construcción de
\forma{hacer} queda incluida como un adjunto en la oración principal, el Tipo~II
ejemplifica la situación inversa: la construcción de \forma{hacer}\etq{temp}
constituye la oración principal, a la que el evento de \forma{salir para el
aeropuerto} se subordina mediante la conjunción \forma{que}. Esta conformación
pone de relieve el \sig{transcurso}, mientras que el evento, expresado en la
completiva de \forma{que}, indica el momento a partir del cual se mide el
transcurso. De este modo, en \exref{ex:13-hace-II} interpretamos \sig{Transcurrieron
dos horas desde que salieron para el aeropuerto}. Por su valor semántico, esta
construcción constituye una subordinada adverbial que, como sostenía
\textcite[§\,778]{bello2002}, lleva omitida la preposición inicial \forma{de} o
\forma{desde} con el significado de \sig{a partir de} (cf. \forma{de tres a
cinco; de Victoria a Retiro}). Por lo general, la preposición se omite, aunque
puede eventualmente reponerse (\forma{Hace dos horas desde que salieron para el
aeropuerto}) y resulta obligatoria tanto cuando se pregunta por el complemento
temporal (\forma{¿De qué hace dos horas?} / \juicio{*}{\forma{¿Qué hace dos
horas?}}) como en la construcción alternativa en que la subordinada se reemplaza
por una nominalización (\forma{Hace dos horas de su salida}). Por último, solo la
construcción del Tipo~II admite otra especificación temporal inicial que indica
el límite final del transcurso (\forma{Hoy hace un año que se murió}).
\end{recuadro}

\subsubsection{Construcciones causales y finales}\label{sec:13-causales}

Las construcciones causales y finales manifiestan, en términos aristotélicos,
respectivamente, la \sig{causa eficiente} y la \sig{causa final} (o propósito).
Las causales pueden indicar la \sig{causa real} de lo que se sostiene en la
oración principal \exref{ex:13-cf-a} o una \sig{causa ilocutiva hipotética}, que
le permite al hablante hacer una inferencia respecto de lo que se dice en la
principal \exref{ex:13-mmo-a}. Las finales indican la \sig{causa final} o el
objetivo por el que se realiza un evento. Tienen carácter prospectivo, es decir
que se orientan a un tiempo posterior al de la principal, y también pueden
indicar un \sig{propósito real} \exref{ex:13-cf-b} o un \sig{propósito
ilocutivo}, que motiva al hablante a decir lo que dice en la principal
\exref{ex:13-mmo-d}.

Las oraciones subordinadas causales y finales pueden formar parte de
construcciones que cumplen funciones argumentales y no argumentales dentro del
predicado, es decir que pueden funcionar como adjuntos o complementos
circunstanciales \exref{ex:13-cf}.

\begin{exe}
  \ex\label{ex:13-cf} Circunstancial (complemento o adjunto verbal):
    \begin{xlist}
      \ex\label{ex:13-cf-a} Causa (construcción causal interna al predicado):\\
          {}[\etq{O} Habla inglés [\etq{SP/AC} por [\etq{OC/CT} que se lo
          enseñaron en la casa]]]. (Habla inglés \textit{porque se lo enseñaron
          en la casa}).
      \ex\label{ex:13-cf-b} Fin (construcción final interna al predicado):\\
          {}[\etq{O} Habla despacio [\etq{SP/AC} para [\etq{OC/CT} que lo
          entiendan bien]]]. (Habla despacio \textit{para que lo entiendan
          bien}).
      \ex\label{ex:13-cf-c} {}[\etq{O} Esta tijera sirve [\etq{SP/CC} para
          [\etq{OC/CT} cortar el césped]]]. (Esta tijera sirve \textit{para
          cortar el césped}).
    \end{xlist}
\end{exe}

En \exref{ex:13-cf-a}, la construcción causal formada por el sintagma
preposicional con núcleo en \forma{por} más la subordinada encabezada por
\forma{que} como su complemento o término es un adjunto circunstancial de causa,
es decir, un modificador del predicado que expresa la \sig{causa real} del estado
de cosas manifestado en la oración principal. En \exref{ex:13-cf-b}, la
construcción final, formada por el sintagma preposicional nucleado en
\forma{para} más la subordinada encabezada por \forma{que} como su complemento o
término, es un adjunto circunstancial de fin, un modificador del predicado que
expresa el propósito de la acción realizada. En \exref{ex:13-cf-c}, la
construcción final, formada por el sintagma preposicional nucleado en
\forma{para} más la subordinada de infinitivo como su complemento o término, es
un complemento circunstancial exigido por la semántica del verbo \forma{servir}.
Recordemos que es habitual que en las gramáticas se considere que \forma{porque}
es una conjunción que encabeza una subordinada causal y que \forma{para que} es
una locución conjuntiva que encabeza una subordinada final; sin embargo, la
posición aquí adoptada es la del enfoque analítico, que presentamos en
§\ref{sec:13-encabezadores}, dado que lo consideramos más simple y
regularizador.

Por otro lado, las construcciones con valores causales y finales también pueden
cumplir funciones no argumentales como adjuntos de nivel oracional, es decir,
como modificadores de la modalidad oracional, que constituyen funciones
periféricas dentro de la oración \exref{ex:13-mmo}.

\begin{exe}
  \ex\label{ex:13-mmo} Modificador de la modalidad oracional (adjunto
      oracional):
    \begin{xlist}
      \ex\label{ex:13-mmo-a} Causa (construcción causal externa al predicado):\\
          {}[\etq{O} Los chicos se habrán intoxicado en la cena, [\etq{SP/MMO}
          por [\etq{OC/CT} que faltaron a la escuela]]].
      \ex\label{ex:13-mmo-b} {}[\etq{O} [\etq{OS/MMO} Como tenemos trabajo
          atrasado], mañana lleguen temprano a la oficina].
      \ex\label{ex:13-mmo-c} {}[\etq{O} Mañana lleguen temprano a la oficina,
          [\etq{OS/MMO} que tenemos trabajo atrasado]].
      \ex\label{ex:13-mmo-d} Fin (construcción final externa al predicado):\\
          {}[\etq{O} [\etq{SP/MMO} Para [\etq{OC/CT} que quede constancia]], los
          estudiantes se quejan del reglamento].
      \ex\label{ex:13-mmo-e} {}[\etq{O} Ponete algo de abrigo, [\etq{SP/MMO} para
          [\etq{OC/CT} que no te enfermes]]].
    \end{xlist}
\end{exe}

La serie \exref{ex:13-mmo} muestra los valores epistémicos e ilocutivos que
pueden asumir las construcciones causales y finales cuando se orientan a la
enunciación y cambian la modalidad de la oración principal. En
\exref{ex:13-mmo-a}, la construcción causal modifica a la oración principal
expresando la razón que tiene el hablante para hacer una inferencia o conjetura,
con valor epistémico (\sig{porque faltaron a la escuela, \textit{deduzco} que los
chicos se intoxicaron en la cena}). En \exref{ex:13-mmo-b} y
\exref{ex:13-mmo-c}, la causal explicativa, antepuesta (\forma{como}) o
pospuesta (\forma{que}) a la oración principal, ofrece una explicación o una
justificación: al manifestar una atenuación se gana en cortesía (\sig{porque
tenemos trabajo atrasado, les \textit{pido} que mañana lleguen temprano a la
oficina}). De manera paralela, en \exref{ex:13-mmo-d}, la construcción final
modifica a la oración principal expresando la finalidad que motiva al hablante,
con valor ilocutivo (\sig{para que quede constancia, \textit{afirmo} que los
estudiantes se quejan del reglamento}), y en \exref{ex:13-mmo-e}, la final
modifica a la oración principal al dar una justificación que atenúa en forma
cortés el acto de habla (\sig{para que no te enfermes, \textit{sugiero} que te
pongas algo de abrigo}).

\subsubsection{Construcciones condicionales y concesivas}\label{sec:13-condicionales}

Condicionales y concesivas expresan otras dos variantes de la causa. En el caso
de las condicionales, manifiestan que la realización de lo enunciado en la
principal depende de una condición o requisito cuyo cumplimiento constituye una
\sig{causa virtual} o \sig{potencial} para que ocurra lo enunciado en la oración
principal. Las concesivas, en cambio, expresan una \sig{objeción} o dificultad
para la realización del evento manifestado en la principal que, sin embargo, no
impide su realización, por lo cual constituyen una \sig{causa inoperante} o
\sig{ineficaz}, es decir, una causa que no logra su objetivo; de modo que las
concesivas aúnan dos significados: \sig{causa} y \sig{contraste}
\parencite{giammatteo2019}. Si se confrontan las concesivas con las
condicionales, también se las puede interpretar como una condición negada, es
decir, como una «contra-causa en cuanto que el efecto es contrario al esperado»
\parencite[344]{fernandezlagunilla1995}.

Condicionales y concesivas cumplen funciones periféricas en la oración como
modificadoras del núcleo oracional o modificadoras de la modalidad oracional.
Junto con la principal forman construcciones binarias: los \termino{períodos
hipotéticos} condicional y concesivo, según se ejemplifica en
\exref{ex:13-mno}, \exref{ex:13-mmocond} y \exref{ex:13-mmoconc}.

\begin{exe}
  \ex\label{ex:13-mno} Modificador del núcleo oracional (adjunto oracional):
    \begin{xlist}
      \ex\label{ex:13-mno-a} Condición (período hipotético condicional):\\
          {}[\etq{O} [\etq{OS/MNO} Si te ofrecen algo de beber], no vayas a
          aceptarlo].
      \ex\label{ex:13-mno-b} Concesión (período hipotético concesivo):\\
          {}[\etq{O} [\etq{SAdv/MNO} Aun [\etq{OC/C} que te ofrezcan algo de
          comer]], no lo aceptes]. (\textit{Aunque te ofrezcan algo de comer},
          no lo aceptes).
    \end{xlist}
\end{exe}

Las oraciones de \exref{ex:13-mno} presentan los períodos hipotéticos que se
constituyen entre las construcciones condicionales \exref{ex:13-mno-a} o
concesivas \exref{ex:13-mno-b} y la oración principal. En \exref{ex:13-mno-a},
la \termino{prótasis} condicional encabezada por \forma{si} manifiesta una causa
hipotética del evento contenido en la \termino{apódosis}, que condiciona su
realización. En \exref{ex:13-mno-b}, la prótasis concesiva formada por el
sintagma adverbial nucleado por el adverbio \forma{aun} más la subordinada
encabezada por \forma{que} como su complemento manifiesta una \sig{causa
ineficaz} del evento contenido en la apódosis, que resulta contrario a las
expectativas. De modo semejante a lo que sucede con \forma{porque}, en las
gramáticas se suele considerar que \forma{aunque} es una conjunción que encabeza
una subordinada concesiva; no obstante, en la versión analítica que aquí
presentamos, la subordinada introducida por \forma{que} es una completiva que
funciona como complemento del adverbio \forma{aun}.

\begin{exe}
  \ex\label{ex:13-mmocond} Modificador de la modalidad oracional (adjunto
      oracional). Condición (período hipotético condicional):
    \begin{xlist}
      \ex\label{ex:13-mmocond-a} {}[\etq{O} [\etq{OS/MMO} Si adornó toda la casa],
          hoy es su cumpleaños y recibe a sus amigas].
      \ex\label{ex:13-mmocond-b} {}[\etq{O} [\etq{OS/MMO} Si no me falla la
          memoria], hoy es el cumpleaños de tu hermano].
      \ex\label{ex:13-mmocond-c} {}[\etq{O} [\etq{OS/MMO} Si te parece bien],
          celebremos su cumpleaños aunque falte un día].
      \ex\label{ex:13-mmocond-d} {}[\etq{O} [\etq{OS/MMO} Si consideramos la
          fecha], hoy deberíamos celebrar su cumpleaños].
    \end{xlist}
\end{exe}

Las oraciones de \exref{ex:13-mmocond} muestran los valores epistémicos e
ilocutivos de los períodos condicionales que se orientan a la enunciación: la
prótasis condicional modifica la modalidad de la oración principal. En
\exref{ex:13-mmocond-a}, la prótasis manifiesta una condición para inferir una
conjetura o una conclusión (el evento contenido en la apódosis) y el período
condicional asume un valor epistémico inferencial (\sig{si adornó toda la casa,
\textit{deduzco} que hoy es su cumpleaños y recibe a sus amigas}). En
\exref{ex:13-mmocond-b}, al igual que en \exref{ex:13-mmocond-c} y
\exref{ex:13-mmocond-d}, la prótasis modifica la fuerza ilocutiva (el acto de
habla que forma la apódosis) y el período condicional asume valores ilocutivos:
manifiesta una atenuación (\sig{si no me falla la memoria, \textit{afirmo} que
hoy es el cumpleaños de tu hermano}), imprime cortesía (\sig{si te parece bien,
\textit{propongo} que celebremos su cumpleaños, aunque falte un día}) o bien
orienta sobre el acto lingüístico (\sig{si consideramos la fecha, \textit{digo}
que hoy deberíamos celebrar su cumpleaños}).

\begin{exe}
  \ex\label{ex:13-mmoconc} Modificador de la modalidad oracional (adjunto
      oracional). Concesión (período hipotético concesivo):
    \begin{xlist}
      \ex\label{ex:13-mmoconc-a} {}[\etq{O} [\etq{SAdv/MMO} Aun [\etq{OC/C} que
          haya pruebas en su contra]], legalmente debe ser considerado
          inocente].
      \ex\label{ex:13-mmoconc-b} {}[\etq{O} [\etq{SAdv/MMO} Aun [\etq{OC/C} que
          parezca mentira]], es la primera vez que lo cita la justicia].
      \ex\label{ex:13-mmoconc-c} {}[\etq{O} [\etq{SAdv/MMO} Aun [\etq{OC/C} que
          no entiendo cómo sucedió]], tu hermano no tuvo la culpa].
    \end{xlist}
\end{exe}

Las oraciones de \exref{ex:13-mmoconc} muestran los valores epistémicos e
ilocutivos de los períodos concesivos orientados a la enunciación: la prótasis
concesiva modifica la modalidad oracional de la oración principal. En
\exref{ex:13-mmoconc-a}, la prótasis manifiesta una condición ineficaz para
inferir determinada conjetura o conclusión respecto del evento contenido en la
apódosis y el período concesivo toma un valor epistémico (\sig{aunque haya
pruebas en su contra, \textit{supongo} que legalmente es inocente}). En
\exref{ex:13-mmoconc-b} y \exref{ex:13-mmoconc-c}, la prótasis modifica la fuerza
ilocutiva (el acto de habla de la apódosis) y el período concesivo toma un valor
ilocutivo: atenúa el acto lingüístico (\sig{aunque parezca mentira,
\textit{afirmo} que es la primera vez que lo cita la justicia}) y ofrece una
orientación (\sig{aunque no entiendo cómo sucedió, \textit{digo} que tu hermano
no tuvo la culpa}).

\begin{recuadro}[Las construcciones con valores adverbiales con formas verbales no finitas]
Según lo señalado en §\ref{sec:13-encabezador}, algunas oraciones subordinadas se
construyen con las formas no finitas del verbo (infinitivo, participio y
gerundio). Estas subordinadas nunca van introducidas por la conjunción
\forma{que}, que requiere un verbo conjugado en la subordinada. Las de participio
y gerundio no llevan encabezadores obligatorios, aunque pueden admitir alguno
opcional: \forma{\textbf{Apenas} terminada la reunión, nos dirigimos al
aeropuerto; \textbf{Aun} considerando el mejor escenario, el rinde de la cosecha
va a resultar pobre}. Las de infinitivo pueden construirse sin encabezador o
formar parte de una construcción preposicional, en la que la preposición no
puede omitirse porque explicita el valor semántico de la subordinada:
\forma{El cielo se tiñó de rojo \textbf{al} caer la tarde; Recibió media falta
\textbf{por} haber llegado tarde / llegar tarde; Nos acercamos a Secretaría
\textbf{para} retirar el diploma}.

Una característica de las subordinadas con verbo no flexionado es que el sujeto
puede coincidir con el sujeto de la principal (como en \forma{Teniendo tiempo,
voy a revisar el trabajo}) o no coincidir con este (como en \forma{Muerto el
perro, se acabó la rabia}). Cuando el sujeto de la subordinada difiere del de la
principal, se pospone al verbo y la subordinada se suele separar de la oración
principal por una pausa (representada por coma en la escritura). Debido a estas
características, que les otorgan mayor independencia respecto de la principal,
estas subordinadas se han denominado \termino{absolutas} (desligadas):
\forma{Agotados los temas del orden del día, se da por concluida la asamblea; Se
da por concluida la asamblea una vez agotados los temas del orden del día}.

Según muestran los ejemplos anteriores, estas subordinadas pueden manifestar
diversos valores adverbiales. Cuando forman parte de una construcción
preposicional (\forma{\textbf{por} haber llegado tarde, \textbf{para} retirar el
diploma}), es la preposición la que marca el valor semántico de la subordinada;
pero muchas veces, cuando se construyen como absolutas, como sostiene
\textcite[672]{morenocabrera1991}, las subordinadas con formas no finitas se
utilizan «cuando no se quiere especificar el tipo de relación que existe entre el
estado de hechos denotado por la subordinada y el denotado por la oración». Es
decir que, al no tener un elemento previo que fije la interpretación, es el
contexto el que determina en cada caso el valor semántico específico de estas
oraciones, que muchas veces puede quedar ambiguo: \forma{Muerto el perro, se
acabó la rabia} (tiempo, causa, condición).
\end{recuadro}

\subsection{Construcciones con valores cuantificacionales-intensificativos}\label{sec:13-cuantificativas}

Se trata de construcciones que funcionan como complemento de un elemento
cuantificativo (\forma{tanto, más, menos}) que modifica el grado de una categoría
léxica (adjetivo, adverbio, verbo o nombre). Forman estructuras correlativas del
tipo \forma{tanto\ldots\ como, más\ldots\ que, tan(to)\ldots\ que}, etc. Las
oraciones subordinadas comparativas y consecutivas modifican el grado de un
elemento cuantificador o intensificador correlativo.

\subsubsection{Construcciones comparativas}\label{sec:13-comparativas}

De las cuatro relaciones básicas señaladas por \textcite{rudolph1996}, las
construcciones comparativas están vinculadas con las de adición y con las de
contraste, que permiten establecer semejanzas y diferencias entre entidades,
propiedades o sucesos. Así, en una estructura comparativa se puede realizar una
comparación sobre el número de elementos de un conjunto (\forma{Tengo
\textbf{menos} exámenes en el escritorio \textbf{que alumnos en la clase}}),
sobre la cantidad de alguna sustancia (\forma{Vi \textbf{más} basura en la playa
\textbf{que el verano pasado}}) o sobre el grado de una propiedad (\forma{La
ópera es \textbf{tan} intensa \textbf{como esperaba}}). En estas construcciones
se presentan valores cuantificables o graduables sujetos a comparación y la
estructura comparativa puede ser tanto una oración subordinada como un fragmento
oracional, siempre con las conjunciones \forma{que} y \forma{como} en su
encabezado. Los ejemplos de \exref{ex:13-compar} presentan las construcciones
comparativas como estructuras cuantificativas complejas.

\begin{exe}
  \ex\label{ex:13-compar} Coda comparativa (adjunto del cuantificador):
    \begin{xlist}
      \ex\label{ex:13-compar-a} {}[\etq{O} Este año la empresa va a trabajar
          \textit{más} horas [\etq{SC/CCM} que el año pasado]].
      \ex\label{ex:13-compar-b} {}[\etq{O} Este año la empresa tiene \textit{mucho
          menos} trabajo [\etq{SC/CCM} que el año pasado]].
      \ex\label{ex:13-compar-c} {}[\etq{O} Este año la empresa desarrolló un
          trabajo \textit{más} intenso [\etq{SC/CCM} que el año pasado]].
      \ex\label{ex:13-compar-d} {}[\etq{O} Este año la empresa trabajó
          \textit{tan} intensamente [\etq{SC/CCM} como el año pasado]].
      \ex\label{ex:13-compar-e} {}[\etq{O} Este año la empresa produjo \textit{dos
          veces más} [\etq{OC/CCM} de lo que se proyectó]].
    \end{xlist}
\end{exe}

\pendiente{Siglas CCM (coda comparativa) y CCN (coda consecutiva): se solapan con
otras abreviaturas; consultar a los autores}

Las oraciones de \exref{ex:13-compar} muestran la comparación del número en los
nombres contables (\forma{horas} en \exref{ex:13-compar-a}) y la comparación de
la cantidad en los nombres incontables (\forma{trabajo} en
\exref{ex:13-compar-b}). Asimismo, estas construcciones también permiten la
comparación del grado en los atributos y en los predicados, señalando el grado
del adjetivo \forma{intenso} en \exref{ex:13-compar-c}, el grado del adverbio
\forma{intensamente} en \exref{ex:13-compar-d} y el grado del verbo
\forma{producir} en \exref{ex:13-compar-e}. Tomados en conjunto, estos ejemplos
muestran que el resultado de la comparación puede indicar una desigualdad
(contraste), si los valores sujetos a comparación no coinciden
(\forma{más / menos\ldots\ que / de}), o una igualdad (semejanza), si los valores
sujetos a comparación coinciden (\forma{tan / tanto\ldots\ como}).

\subsubsection{Construcciones consecutivas}\label{sec:13-consecutivas}

Las construcciones consecutivas están vinculadas con la noción de causalidad.
Según \textcite[343]{ditullio2014}, «desde un punto de vista que solo contemple
las relaciones semánticas de orden lógico, estas oraciones son las inversas de
las causales, ya que expresan el resultado o efecto, de lo expresado por la
principal, la causa». Así, Di Tullio compara dos oraciones, \forma{Estoy agotado
porque caminé mucho} y \forma{Caminé tanto que estoy agotado}, para concluir que
difieren por el carácter ponderativo de la segunda, ausente en la primera, que
deriva de la presencia del intensificador \forma{tanto}.

En una estructura consecutiva se realiza una intensificación cuantitativa
(\forma{Había \textbf{tanta} basura en la playa \textbf{que no se podía
caminar}}) o cualitativa (\forma{Teníamos \textbf{tal} gentío en la puerta
\textbf{que debimos entregar números}}) y la consecuencia de esta intensificación
es un evento, manifestado por una subordinada encabezada por la conjunción
subordinante \forma{que}. Los ejemplos de \exref{ex:13-consec}, por su parte,
muestran las construcciones consecutivas como estructuras cuantificativas
complejas.

\begin{exe}
  \ex\label{ex:13-consec} Coda consecutiva (adjunto del cuantificador):
    \begin{xlist}
      \ex\label{ex:13-consec-a} {}[\etq{O} Tuvimos \textit{tantos} despachos
          [\etq{OC/CCN} que contratamos más empleados]].
      \ex\label{ex:13-consec-b} {}[\etq{O} Tuvimos \textit{tanta} inactividad
          [\etq{OC/CCN} que casi cerramos el local]].
      \ex\label{ex:13-consec-c} {}[\etq{O} El trabajo estuvo \textit{tan} flojo
          [\etq{OC/CCN} que redujimos los turnos a la mitad]].
      \ex\label{ex:13-consec-d} {}[\etq{O} Trabajé \textit{tan} bien [\etq{OC/CCN}
          que todavía no lo puedo creer]].
      \ex\label{ex:13-consec-e} {}[\etq{O} La empresa facturó \textit{tanto}
          [\etq{OC/CCN} que en un mes se alcanzó la meta bimestral]].
      \ex\label{ex:13-consec-f} {}[\etq{O} Las ventas fueron de \textit{tal /
          tanta} magnitud [\etq{OC/CCN} que reprogramamos muchos envíos]].
    \end{xlist}
\end{exe}

La serie \exref{ex:13-consec} presenta la intensificación cuantitativa del número
en los nombres contables (\forma{despachos} en \exref{ex:13-consec-a}) y de la
cantidad en los nombres incontables (\forma{inactividad} en
\exref{ex:13-consec-b}). Asimismo, estas construcciones permiten la
intensificación del grado en los atributos, en los predicados y en nombres
abstractos: el grado del adjetivo \forma{flojo} en \exref{ex:13-consec-c}, el
grado del adverbio \forma{bien} en \exref{ex:13-consec-d}, el grado del verbo
\forma{facturar} en \exref{ex:13-consec-e} y el grado del nombre \forma{magnitud}
en \exref{ex:13-consec-f}. En todos los casos, los ejemplos presentan la
intensificación extrema de lo expresado por una palabra léxica (nombre,
adjetivo, adverbio, verbo) como causa de la consecuencia codificada en la oración
subordinada consecutiva (\forma{tal / tan / tanto\ldots\ que}).

En las construcciones comparativas se da la correlación entre un cuantificador y
la coda comparativa y, paralelamente, en el caso de las consecutivas, la
correlación se establece entre un intensificador y la coda consecutiva. En cuanto
a las estructuras internas de tales codas, solo la coda comparativa admite tanto
oraciones subordinadas como sintagmas no oracionales, mientras que la coda
consecutiva solo admite oraciones subordinadas.

\begin{ejercicio}\label{ej:13-paratextos}
En los siguientes textos breves, formados por secuencias paratextuales (volanta,
título, copete) de noticias periodísticas, identifique las oraciones
subordinadas y las funciones sintácticas que cumplen, según se muestra en el
ejemplo resuelto~1.

% Textos periodísticos protegidos: se incorporan a mano desde el original.
\textoaincorporar{los diez paratextos periodísticos (volanta, título y copete)
del original: 1. Guido Kaczka y «Los 8 escalones» (\textit{La Razón},
21/07/2015); 2. Cacho Castaña en el Colón (\textit{La Razón}, 05/12/2016); 3.
camiones para tapar baches (\textit{La Razón}, 01/06/2016); 4. carta de Ángel Di
María (\textit{Infobae}, 05/08/2019); 5. declaraciones de Diego Maradona
(\textit{La Razón}, 22/07/2015); 6. la Selección en Houston (\textit{La Razón},
22/06/2016); 7. Jazmín Stuart (\textit{La Razón}, 24/11/2016); 8. «Maradona sin
vueltas» (\textit{La Razón}, 22/07/2015); 9. feriado por Güemes (\textit{La
Razón}, 01/06/2016); 10. fiesta clandestina en Moreno (\textit{La Razón},
22/09/2016)}

\begin{resuelto}
1. La OR encabezada por el pronombre relativo \forma{que} (\forma{que no le
gustó nada a Guido Kaczka}) es un adjunto especificativo (AS) del antecedente
nominal \forma{broma}. La OC de infinitivo (sin subordinante, \forma{participar})
es el complemento o término (CT) de la preposición \forma{a} (valor final). La
OC encabezada por la conjunción \forma{que} (\forma{que se cansó y se fue}) es la
coda consecutiva (CCN) correlativa del intensificador \forma{tanto}, con el que
forma una estructura cuantificativa compleja.
\end{resuelto}
\end{ejercicio}

\begin{recuadro}[Reflexión final]
Entre las relaciones interoracionales, la subordinación (hipotaxis) permite crear
sentidos complejos al vincular mundos completos (reales, irreales, posibles,
pasados, eventuales) que, en principio, no presentan vínculos entre sí. Más
estructurales o más periféricas, las relaciones de subordinación caracterizan a
la oración compuesta de la misma manera en que las relaciones de concordancia
caracterizan a la oración simple. Miradas en conjunto, las oraciones subordinadas
siguen patrones estructurales notablemente regulares: encabezadores léxicos
conjuntivos o relativos con verbos finitos, por un lado; verbos no finitos sin
encabezadores léxicos conjuntivos o relativos, por el otro. Los valores
semántico-pragmáticos que estas oraciones adoptan, al mismo tiempo, se pueden
resumir en dos ejes: el de la temporalidad (adición / tiempo) y el del punto de
vista (contraste / causalidad). Con estos pocos recursos gramaticales
articulamos la subjetividad de nuestros pensamientos complejos y, si acaso, los
convertimos en conocimientos intersubjetivos cuando los comunicamos a personas
que, aun sin conocerlas, tenemos por iguales.
\end{recuadro}

% --- Secciones finales -------------------------------------------------
\lecturasrecomendadas

\solucionario

\begin{solucion}{\ref{ej:13-paremias}}
3. [\etq{Ocomp} [\etq{OS} \textit{Cuando el invierno primavea}], la primavera
invernea]. 4. [\etq{Ocomp} [\etq{OS} \textit{Cuando veas} [\etq{OS} \textit{las
barbas de tu vecino cortar}]], pon las tuyas a [\etq{OS} \textit{remojar}]]. 5.
[\etq{Ocomp} [\etq{OS} \textit{Dando gracias por agravios}] negocian los hombres
sabios]. 6. [\etq{Ocomp} [\etq{OS} \textit{Diciendo las verdades}] se pierden
amistades]. 7. [\etq{Ocomp} [\etq{OS} \textit{Donde fueres}], haz [\etq{OS}
\textit{lo que vieres}]]. 8. [\etq{Ocomp} [\etq{OS} \textit{Donde hubo fuego}],
cenizas quedan]. 9. [\etq{Ocomp} El [\etq{OS} \textit{salir de la posada}] es la
mejor jornada]. 10. [\etq{Ocomp} Es de bien nacidos [\etq{OS} \textit{ser
agradecidos}]]. 11. [\etq{Ocomp} [\etq{OS} \textit{Hablando del rey de Roma}], por
la puerta asoma]. 12. [\etq{Ocomp} [\etq{OS} \textit{Lo que natura non da}]
Salamanca non presta]. 13. [\etq{Ocomp} No por [\etq{OS} \textit{mucho
madrugar}] amanece más temprano]. 14. [\etq{Ocomp} Nunca es tarde [\etq{OS}
\textit{si la dicha es buena}]]. 15. [\etq{Ocomp} Pájaro [\etq{OS} \textit{que
comió}], voló]. 16. [\etq{Ocomp} Perro [\etq{OS} \textit{que ladra}] no muerde].
17. [\etq{Ocomp} [\etq{OS} \textit{Preguntando}] se llega a Roma]. 18.
[\etq{Ocomp} [\etq{OS} \textit{Si quieres} [\etq{OS} \textit{ser bien
servido}]], sírvete a ti mismo].\pendiente{El ejercicio 1 del original no
incluye el refrán 9 (\textit{El salir de la posada}), que sí aparece en los
ejercicios 2 y 6: se agregó y se renumeró}
\end{solucion}

\begin{solucion}{\ref{ej:13-introductor}}
3. Con relativo \forma{cuando}. 4. Con relativo \forma{cuando}, y sin
subordinante en [\etq{OS} \textit{las barbas de tu vecino cortar}] y [\etq{OS}
\textit{remojar}]. 5. Sin subordinante. 6. Sin subordinante. 7. Con relativos
\forma{donde} y \forma{lo que}. 8. Con relativo \forma{donde}. 9. Sin
subordinante. 10. Sin subordinante. 11. Sin subordinante. 12. Con relativo
\forma{lo que}. 13. Sin subordinante. 14. Con conjunción \forma{si}. 15. Con
relativo \forma{que}. 16. Con relativo \forma{que}. 17. Sin subordinante. 18.
Con conjunción \forma{si}, y sin subordinante en [\etq{OS} \textit{ser bien
servido}].
\end{solucion}

\begin{solucion}{\ref{ej:13-principito}}
Respuestas posibles: 2. [\etq{Ocomp} [\etq{OS} \textbf{Si} queremos disfrutar de
las mariposas], a veces tenemos que aguantar a las orugas] / [\etq{Ocomp}
[\etq{OR} \textbf{Cuando} queremos disfrutar de las mariposas], tenemos que
aguantar a las orugas]. 3. [\etq{Ocomp} Nunca te olvides de [\etq{OC}
\textbf{que} eres el dueño de tu vida y tus emociones]]. 4. [\etq{Ocomp}
[\etq{OS} \textbf{Como} los hombres ya no tienen tiempo], compran cosas hechas a
los comerciantes] / [\etq{Ocomp} Los hombres compran cosas hechas a los
comerciantes por [\etq{OC} \textbf{que} ya no tienen tiempo]]. 5. [\etq{Ocomp}
Me pregunto [\etq{OC} \textbf{si} las estrellas se iluminan con algún fin]]. 6.
[\etq{Ocomp} Nunca estamos contentos en el lugar [\etq{OR} \textbf{donde}
estamos]] / [\etq{Ocomp} Nunca estamos contentos [\etq{OR} \textbf{cuando}
estamos en un lugar]]. 7. [\etq{Ocomp} Todas las estrellas ríen [\etq{OR}
\textbf{cuando} brilla la luna]]. 8. [\etq{Ocomp} Aun [\etq{OC} \textbf{que}
pocas lo recuerdan], todas las personas grandes han sido niños antes] /
[\etq{Ocomp} Aun [\etq{OR} \textbf{cuando} pocas lo recuerdan], todas las
personas grandes han sido niños antes]. 9. [\etq{Ocomp} [\etq{OS} \textbf{Si} un
hombre es vanidoso], todos los demás hombres son admiradores] / [\etq{Ocomp}
Todos los demás hombres son admiradores [\etq{OR} \textbf{cuando} un hombre es
vanidoso]]. 10. [\etq{Ocomp} Hay que hacer cuidadosamente la limpieza del
planeta [\etq{OR} \textbf{cuando} uno termina de arreglarse]] / [\etq{Ocomp} Hay
que hacer cuidadosamente la limpieza del planeta al [\etq{OC} terminar de
arreglarse]].
\end{solucion}

\begin{solucion}{\ref{ej:13-analitica}}
2. [\etq{Ocomp} [\etq{SAdv} \textit{Aun} [\etq{OC} \textbf{que} llueva]], vamos a
salir], con valor concesivo en \forma{aun}. 3. [\etq{Ocomp} Habló con el rector
[\etq{SP} \textit{con el fin de} [\etq{OC} \textbf{que} se aprobase su curso]]],
con valor final en \forma{con el fin de} (locución preposicional). 4.
[\etq{Ocomp} Lo explico de nuevo [\etq{SP} \textit{para} [\etq{OC} \textbf{que}
no queden dudas]]], con valor final en \forma{para}. 5. [\etq{Ocomp} Me alegro
[\etq{SP} \textit{por} [\etq{OC} \textbf{que} lo hice con cariño]]], con valor
causal en \forma{por}. 6. [\etq{Ocomp} No pude llegar antes [\etq{SP}
\textit{por} [\etq{OC} \textbf{que} había mucho tráfico]]], con valor causal en
\forma{por}. 7. [\etq{Ocomp} No va a aceptar, [\etq{SAdv} \textit{aun} [\etq{OC}
\textbf{que} se lo pidamos amablemente]]], con valor concesivo en \forma{aun}. 8.
[\etq{Ocomp} Nos alegra tu visita [\etq{SP} \textit{por} [\etq{OC} \textbf{que}
te extrañábamos]]], con valor causal en \forma{por}. 9. [\etq{Ocomp} Te lo
repito [\etq{SP} \textit{para} [\etq{OC} \textbf{que} lo entiendas]]], con valor
final en \forma{para}. 10. [\etq{Ocomp} Trabajamos [\etq{SP} \textit{con el fin
de} [\etq{OC} \textbf{que} el proyecto sea un éxito]]], con valor final en
\forma{con el fin de} (locución preposicional).
\end{solucion}

\begin{solucion}{\ref{ej:13-perifrasis}}
3. [\etq{Ocomp} No voy a la fiesta [\etq{SP} por [\etq{OC} \textbf{que} estoy
cansado]]]: oración subordinada con verbo en forma personal introducida por la
conjunción subordinante \forma{que}. 4. [\etq{Ocomp} [\etq{SAdv} Aun [\etq{OS}
\textit{estando cansados}]], trabajamos duro]: oración subordinada con verbo en
forma no personal, gerundio. 5. [\etq{Ocomp} [\etq{SAdv} Aun [\etq{OC}
\textbf{que} estábamos cansados]], trabajamos duro]: oración subordinada con
verbo en forma personal introducida por la conjunción subordinante \forma{que}.
6. [\etq{Ocomp} [\etq{OS} \textit{No estando invitado}], no \textit{puedo ir} al
cumpleaños]: oración subordinada con verbo en forma no personal, gerundio; en la
principal, perífrasis verbal de \forma{poder} + infinitivo, de valor modal
deóntico. [\etq{Ocomp} [\etq{OS} \textbf{Si} no me invitan], no puedo ir al
cumpleaños]: oración subordinada con verbo en forma personal introducida por la
conjunción subordinante \forma{si}; en la principal, perífrasis verbal de
\forma{poder} + infinitivo, de valor modal deóntico. 7. [\etq{Ocomp} Es
importante [\etq{OC} \textit{comprender el funcionamiento del motor}]]: oración
subordinada con verbo en forma no personal, infinitivo. 8. [\etq{Ocomp} Es
importante [\etq{OC} \textbf{que} comprendas el funcionamiento del motor]]:
oración subordinada con verbo en forma personal introducida por la conjunción
subordinante \forma{que}. 9. [\etq{Ocomp} Logró [\etq{OC} \textbf{que} sus
amigos supieran la verdad]]: oración subordinada con verbo en forma personal
introducida por la conjunción subordinante \forma{que}. 10. [\etq{Ocomp} Logró
[\etq{OC} \textit{saber la verdad}]]: oración subordinada con verbo en forma no
personal, infinitivo. 11. [\etq{Ocomp} Quiso [\etq{OC} \textbf{que} sus amigos
supieran la verdad]]: oración subordinada con verbo en forma personal
introducida por la conjunción subordinante \forma{que}. 12. [\etq{Ocomp} Quiso
[\etq{OC} \textit{saber la verdad}]]: oración subordinada con verbo en forma no
personal, infinitivo. 13. [\etq{Ocomp} Se negó [\etq{SP} a [\etq{OC}
\textbf{que} sus amigos supieran la verdad]]]: oración subordinada con verbo en
forma personal introducida por la conjunción subordinante \forma{que}. 14.
[\etq{Ocomp} Se negó [\etq{SP} a [\etq{OC} \textit{saber la verdad}]]]: oración
subordinada con verbo en forma no personal, infinitivo. 15. [\etq{O} De esta
manera \textit{comenzó a saber} la verdad]: oración simple con perífrasis verbal
de \forma{comenzar} + \forma{a} + infinitivo, de valor aspectual incoativo. 16.
[\etq{O} De esta manera \textit{pudo saber} la verdad]: oración simple con
perífrasis verbal de \forma{poder} + infinitivo, de valor modal, capacidad. 17.
[\etq{O} De esta manera \textit{va a saber} la verdad]: oración simple con
perífrasis verbal de \forma{ir} + \forma{a} + infinitivo, de valor temporal
futuro. 19. [\etq{Ocomp} El tenor ingresó a la escena [\etq{OS} \textbf{mientras}
cantaba el aria]]: oración subordinada con verbo en forma personal introducida
por la conjunción subordinante \forma{mientras}. 20. [\etq{Ocomp} El tenor
ingresó [\etq{OS} \textit{cantando el aria}]]: oración subordinada con verbo en
forma no personal, gerundio.
\end{solucion}

\begin{solucion}{\ref{ej:13-funciones}}
3. [\etq{Ocomp} [\etq{OR/AC} \textit{Cuando el invierno primavea}], la primavera
invernea]: la OR libre encabezada por el adverbio relativo \forma{cuando} es un
adjunto circunstancial (AC) de tiempo; atendiendo al sentido generalizador del
refrán, admite otra interpretación, en la que la oración subordinada funciona
como un adjunto modificador de núcleo oracional: la prótasis de un período
hipotético condicional (\sig{si el invierno primavea, la primavera invernea}).

4. [\etq{Ocomp} [\etq{OR/AC} Cuando veas [\etq{OC/COD} \textit{las barbas de tu
vecino cortar}]], pon las tuyas a [\etq{OC/CT} \textit{remojar}]]: la OR libre
encabezada por el relativo \forma{cuando} es un adjunto circunstancial (AC) de
tiempo, y admite otra interpretación como adjunto modificador de núcleo
oracional con valor condicional; la primera OC de infinitivo (sin subordinante)
es el complemento objeto directo (COD) del verbo transitivo \forma{ver}
(\sig{cuando veas que cortan las barbas de tu vecino}) y la segunda OC de
infinitivo (sin subordinante) es el complemento o término (CT) de la preposición
\forma{a} (valor final).

5. [\etq{Ocomp} [\etq{OS/AC} \textit{Dando gracias por agravios}] negocian los
hombres sabios]: la OS de gerundio (sin subordinante) es un adjunto
circunstancial (AC) de modo o manera.

6. [\etq{Ocomp} [\etq{OS/AC} \textit{Diciendo las verdades}] se pierden
amistades]: la OS de gerundio (sin subordinante) es un adjunto circunstancial
(AC) de modo o manera.

7. [\etq{Ocomp} [\etq{OR/AC} \textit{Donde fueres}], haz [\etq{OR/COD}
\textit{lo que vieres}]]: la OR libre encabezada por el adverbio relativo
\forma{donde} es un adjunto circunstancial (AC) de lugar y la OR semilibre
encabezada por el relativo \forma{lo que} es el COD del verbo transitivo
\forma{hacer}; por su valor generalizador como refrán, admite otra
interpretación, en la que la OS encabezada por \forma{donde} con verbo en
subjuntivo (futuro) es un adjunto modificador de núcleo oracional con valor
condicional (\sig{si estuvieres afuera, haz lo que vieres}).

8. [\etq{Ocomp} [\etq{OR/AC} \textit{Donde hubo fuego}], cenizas quedan]: la OR
libre encabezada por el adverbio relativo \forma{donde} es un adjunto
circunstancial (AC) de lugar; por su valor generalizador como refrán, admite otra
interpretación, en la que la OS encabezada por \forma{donde} es un adjunto
modificador de núcleo oracional con valor condicional (\sig{si hubo fuego,
cenizas quedan}).

9. [\etq{Ocomp} El [\etq{OC/S} \textit{salir de la posada}] es la mejor
jornada]: la OC de infinitivo (sin subordinante) es el sujeto (S) de la
oración; el determinante definido aporta valor expresivo (\sig{el salir de la
posada es la mejor jornada} = \sig{salir de la posada es la mejor jornada}).

10. [\etq{Ocomp} Es de bien nacidos [\etq{OC/S} \textit{ser agradecidos}]]: la
OC de infinitivo (sin subordinante) es el sujeto (S) de la oración; el
determinante definido aportaría valor expresivo (\sig{es de bien nacidos ser
agradecidos} = \sig{es de bien nacidos el ser agradecidos}).

11. [\etq{Ocomp} [\etq{OS/MNO} \textit{Hablando del rey de Roma}], por la puerta
asoma]: la OS de gerundio (sin subordinante) es un adjunto de nivel oracional, un
modificador de núcleo oracional (MNO) que introduce un referente discursivo que
da las condiciones de verdad de toda la oración; la OS puede tener una
interpretación temporal (\sig{cuando se habla del rey de Roma, por la puerta
asoma}).

12. [\etq{Ocomp} [\etq{OR/COD} \textit{Lo que natura non da}] Salamanca non
presta]: la OR semilibre encabezada por el relativo \forma{lo que} es el
complemento objeto directo (COD) del verbo transitivo \forma{prestar}; la
interpretación global del refrán, generalizadora, es la de un período hipotético
condicional (\sig{si natura non lo da, Salamanca non lo presta}).

13. [\etq{Ocomp} No por [\etq{OC/CT} \textit{mucho madrugar}] amanece más
temprano]: la OC de infinitivo (sin subordinante) es el complemento o término
(CT) de la preposición \forma{por} (valor causal); la construcción tiene el
mismo valor contrastivo que un período hipotético concesivo (\sig{aunque mucho
se madrugue, no amanece más temprano}).

14. [\etq{Ocomp} Nunca es tarde [\etq{OS/MNO} \textit{si la dicha es buena}]]:
la OS encabezada por la conjunción subordinante \forma{si} es un adjunto de nivel
oracional, un modificador de núcleo oracional (MNO) que introduce un estado de
cosas en el mundo que da las condiciones de verdad para el período hipotético
condicional; existe una lectura temporal de fondo, sobre la que se monta la
interpretación condicional, como se aprecia en la versión alternativa del refrán
(\forma{Nunca es tarde cuando la dicha es buena}).

15. [\etq{Ocomp} Pájaro [\etq{OR/AS} \textit{que comió}], voló]: la OR
encabezada por el relativo \forma{que} es un adjunto especificativo (AS) del
antecedente nominal \forma{pájaro}, cuya referencia se restringe.

16. [\etq{Ocomp} Perro [\etq{OR/AS} \textit{que ladra}] no muerde]: la OR
encabezada por el relativo \forma{que} es un adjunto especificativo (AS) del
antecedente nominal \forma{perro}, cuya referencia se restringe.

17. [\etq{Ocomp} [\etq{OS/AC} \textit{Preguntando}] se llega a Roma]: la OS de
gerundio (sin subordinante) es un adjunto circunstancial (AC) de modo o manera;
la OS de gerundio, carente de una marca (conjunción, adverbio o preposición) que
oriente en la interpretación, puede dar lugar a lecturas diversas:
\textit{temporal} (\sig{cuando se pregunta}), \textit{condicional} (\sig{si se
pregunta}) e incluso \textit{causal} (\sig{porque se pregunta}).

18. [\etq{Ocomp} [\etq{OS/MNO} \textit{Si quieres} [\etq{OC/COD} \textit{ser bien
servido}]], sírvete a ti mismo]: la OS encabezada por la conjunción subordinante
\forma{si} es un adjunto de nivel oracional, un modificador de núcleo oracional
(MNO) que introduce un estado de cosas en el mundo que da las condiciones de
verdad para el período hipotético condicional; la OC de infinitivo (sin
subordinante) es el complemento objeto directo (COD) del verbo transitivo
\forma{querer}.
\end{solucion}

\begin{solucion}{\ref{ej:13-relativos}}
2. [\etq{O} Compré una valija]. [\etq{Ocomp} De hecho, compré [\etq{OR/COD}
\textit{la que} me recomendaste]]: \forma{la que}, complemento objeto directo del
verbo transitivo \forma{recomendar}; antecedente contextual, \forma{valija}. 4.
[\etq{Ocomp} Decime [\etq{OR/COD} \textit{lo que} pensás]]: \forma{lo que},
complemento objeto directo del verbo transitivo \forma{pensar}; sin antecedente.
5. [\etq{Ocomp} Dieron un premio a [\etq{OR/CT} \textit{quienes}
participaron]]: \forma{quienes}, sujeto de la OR; sin antecedente. 6.
[\etq{Ocomp} El coche [\etq{OR/AS} \textit{que} está estacionado allí] es el
mío]: \forma{que}, sujeto de la OR; antecedente, \forma{coche}. 7. [\etq{Ocomp}
El libro [\etq{OR/AS} \textit{que} me prestaste] es muy interesante]:
\forma{que}, complemento objeto directo del verbo transitivo \forma{prestar};
antecedente, \forma{libro}. 8. [\etq{Ocomp} [\etq{OR/S} \textit{El que} llegó
tarde] no entró]: \forma{el que}, sujeto de la OR; sin antecedente. 9.
[\etq{Ocomp} Esa es la casa [\etq{OR/AS} \textit{cuya} puerta estaba pintada de
verde]]: \forma{cuya}, especificador de \forma{puerta}; antecedente,
\forma{casa}. 10. [\etq{Ocomp} Esta es la lapicera [\etq{OR/AS} con \textit{la
cual} escribí la carta]]: \forma{la cual}, complemento de la preposición
\forma{con} (valor instrumental); antecedente, \forma{lapicera}. 11.
[\etq{Ocomp} Hizo la preparación [\etq{OR/AC} \textit{como} le enseñaron]]:
\forma{como}, adjunto circunstancial de modo o manera del verbo
\forma{enseñar}; sin antecedente. 12. [\etq{Ocomp} La docente [\etq{OR/AS} con
\textit{quien} estudié inglés] era británica]: \forma{quien}, complemento de la
preposición \forma{con} (valor sociativo); antecedente, \forma{docente}. 13.
[\etq{Ocomp} La película [\etq{OR/AS} por \textit{la cual} vinimos al festival]
fue excelente]: \forma{la cual}, complemento o término de la preposición
\forma{por} (valor causal); antecedente, \forma{película}. 14. [\etq{Ocomp} Lo
voy a comprar [\etq{OR/AC} \textit{cuando} esté en oferta]]: \forma{cuando},
adjunto circunstancial de tiempo del verbo \forma{estar}; sin antecedente. 15.
[\etq{Ocomp} Los estudiantes [\etq{OR/AS} \textit{que} estudiaron] aprobaron el
examen]: \forma{que}, sujeto de la OR; antecedente, \forma{estudiantes}. 16.
[\etq{Ocomp} Los estudiantes, [\etq{OR/AX} \textit{que} estudiaron], aprobaron el
examen]: \forma{que}, sujeto de la OR; antecedente, \forma{estudiantes}. 17.
[\etq{Ocomp} Necesito [\etq{OR/COD} \textit{lo que} preparaste para el relleno
de la tarta]]: \forma{lo que}, complemento objeto directo del verbo transitivo
\forma{preparar}; sin antecedente. 18. [\etq{Ocomp} No me importa [\etq{OR/S}
\textit{lo que} digan]]: \forma{lo que}, complemento objeto directo del verbo
transitivo \forma{decir}; sin antecedente. 19. [\etq{Ocomp} [\etq{OR/S}
\textit{Quien} persevera], vence]: \forma{quien}, sujeto de la OR; sin
antecedente. 20. [\etq{Ocomp} [\etq{OR/S} \textit{Quienes} estudian], aprenden]:
\forma{quienes}, sujeto de la OR; sin antecedente.
\end{solucion}

\begin{solucion}{\ref{ej:13-paratextos}}
2. La OC encabezada por la conjunción \forma{que} (\forma{que no iba a volver a
hablar}) es el complemento objeto directo (COD) del verbo transitivo
\forma{pensar}. La OC encabezada por la conjunción \forma{que} (\forma{que está
feliz\ldots}) es el COD del verbo transitivo \forma{contar}. La OR encabezada
por el pronombre relativo \forma{que} (\forma{que dará en el escenario más
importante del país}) es un adjunto especificativo (AS) del antecedente nominal
\forma{concierto}. La OC (sin subordinante) con el pronombre interrogativo
\forma{cómo} (\forma{cómo fueron los momentos más duros de su enfermedad\ldots})
es el COD del verbo transitivo \forma{recordar}. La OR encabezada por el
pronombre relativo \forma{que} (\forma{que lo tuvo en coma dos meses}) es un
adjunto explicativo (AX) del antecedente nominal \forma{enfermedad}. La OS
encabezada por la conjunción \forma{que} (\forma{que es más lindo}) es un adjunto
circunstancial (AC) de causa (\sig{mejor me quedo y vivo, porque es más lindo}).
La oración completa entrecomillada en simples es una mención en estilo directo,
COD del verbo transitivo \forma{pensar}.

3. La OC de infinitivo (sin subordinante, \forma{tapar baches en 20 minutos}) es
el complemento o término (CT) de la preposición \forma{para}. La OR encabezada
por el pronombre relativo \forma{que} (\forma{que podían hacer las viejas
cuadrillas}) es un adjunto especificativo (AS) del antecedente elidido
\forma{aperturas} (\sig{la mitad de aperturas}). La OC encabezada por la
conjunción \forma{que} (\forma{que íbamos a tapar cada bache en no más de 15
días}) es el complemento o término (CT) de la preposición
\forma{a}.\pendiente{El original atribuye a \textit{a} «valor final» en
\textit{comprometerse a}}

4. La OR encabezada por el pronombre relativo \forma{que} (\forma{que
revolucionó al mundo del deporte}) es un adjunto especificativo (AS) del
antecedente nominal \forma{carta}. La OR encabezada por el pronombre relativo
\forma{que} (\forma{que así ganó la Supercopa de Francia}) es un adjunto
explicativo (AX) del antecedente nominal \forma{PSG}. La OS de gerundio (sin
subordinante, \forma{posteando un texto contra\ldots}) es un adjunto
circunstancial (AC) de modo o manera del verbo \forma{sorprender}
(alternativamente, la construcción acepta la interpretación temporal, \sig{al
postear un texto contra\ldots}). La OR semilibre encabezada por el pronombre
relativo \forma{los que} (\forma{los que critican a los futbolistas}) es el
complemento o término (CT) de la preposición \forma{contra}. La OC encabezada
por la conjunción \forma{que} (\forma{que el fútbol no tiene nada que ver con la
vida}) es el COD del verbo transitivo \forma{decir}. (Interpretamos que la
locución verbal \forma{tener que ver} se encuentra cuantificada en su grado,
\forma{tener \{nada / algo / bastante / mucho\} que ver}, con mantenimiento de
su régimen preposicional). La OC (sin subordinante) con el pronombre
interrogativo \forma{cuánto} (\forma{cuánto saben de la vida}) es el COD del
verbo transitivo \forma{saber}.

5. La OC encabezada por la conjunción \forma{si} (\forma{si un tipo de mi edad
se deja influenciar por alguien\ldots}) es el sujeto (S) de la oración.
(Interpretamos como perífrasis verbal \forma{dejarse influenciar}, con valor
modal radical, para indicar permiso). La OR encabezada por el pronombre relativo
\forma{que} (\forma{que está defendiendo sus cosas}) es un adjunto
especificativo (AS) del antecedente pronominal indefinido \forma{alguien}. La OR
encabezada por el adverbio relativo \forma{como} (\forma{como no lo hice
nunca}), cuyo antecedente es toda la oración anterior, constituye un adjunto
explicativo (AX) de esta (\sig{no voy a mandar al frente a nadie, de la misma
manera en que no lo hice nunca}). La OS encabezada por la conjunción subordinante
\forma{si} (\forma{si no me devuelven mi plata}) es un adjunto de nivel
oracional, un modificador de núcleo oracional (MNO) que introduce un estado de
cosas en el mundo que da las condiciones de verdad para el período hipotético
condicional (en términos ilocutivos, toda la oración representa una amenaza).

6. La OC de infinitivo (sin subordinante, \forma{volver a jugar otra final}) es
el predicativo subjetivo obligatorio (PSO) del verbo copulativo \forma{ser}. La
OC encabezada por la conjunción subordinante \forma{que} (\forma{que lo
merecemos}) es el complemento o término (CT) de la preposición \forma{por} (valor
causal).

7. La OR encabezada por el pronombre relativo \forma{el que} (\forma{en el que
protagoniza escenas jugadas}) es un adjunto especificativo (AS) del antecedente
nominal \forma{rol}. La OC encabezada por la conjunción \forma{que} (\forma{que
ese rol la puede reposicionar, tras\ldots}) es el COD del verbo \forma{sentir}.
La OC de infinitivo (sin subordinante, \forma{perder espacio para\ldots}) es
complemento o término (CT) de la preposición \forma{tras} (valor temporal). La
OC de infinitivo (sin subordinante, \forma{resguardar su intimidad}) es
complemento o término (CT) de la preposición \forma{para} (valor final). La OR
encabezada por el relativo \forma{que} (\forma{que cotiza alto en este oficio})
es un adjunto especificativo (AS) del antecedente nominal \forma{moneda}.

8. La OC encabezada por la conjunción \forma{que} (\forma{que si le regresan el
dinero\ldots}) es el COD del verbo \forma{asegurar}. La OS encabezada por la
conjunción subordinante \forma{si} (\forma{si le regresan el dinero\ldots}) es
un adjunto de nivel oracional, un modificador de núcleo oracional (MNO) que
introduce un estado de cosas en el mundo que da las condiciones de verdad para
el período hipotético condicional (en términos ilocutivos, toda la oración
representa una promesa). La OR encabezada por el pronombre relativo \forma{que}
(\forma{que le falta}) es un adjunto especificativo (AS) del antecedente nominal
\forma{dinero}. La OR libre encabezada por el pronombre relativo \forma{quien}
(\forma{quien ayer asistió a la primera audiencia}) es una aposición explicativa
(APX) del nombre propio, correferente y con el que puede cambiar posiciones, lo
que no sería esperable de una OR con antecedente nominal, que debe precederla.

9. La OC de infinitivo (sin subordinante, \forma{sumar un feriado en junio}) es
el complemento o término (CT) de la preposición \forma{de}. La OC de infinitivo
(sin subordinante, \forma{recordar la muerte del General Güemes\ldots}) es el
complemento de la preposición \forma{para}. La OS encabezada por la conjunción
subordinante \forma{si} (\forma{si obtiene la sanción definitiva del Congreso})
es un adjunto de nivel oracional, un modificador de núcleo oracional (MNO) que
introduce un estado de cosas en el mundo que da las condiciones de verdad para
el período hipotético condicional.

10. La OR encabezada por el pronombre relativo \forma{que} (\forma{que halló la
Policía}) es un adjunto especificativo (AS) del antecedente nominal
\forma{elementos}.

\pendiente{El original reproduce además cada paratexto con todos los corchetes;
se omite junto con los textos}
\end{solucion}

% =====================================================================
%  capitulos/cap14.tex — Capítulo 14. La epistemicidad: modalidad y
%  evidencialidad
%  Autora: Ana Marcovecchio
%  Fuente: Capitulo_14_normalizado.docx
% =====================================================================
\capitulo{La epistemicidad: modalidad y evidencialidad}{Ana Marcovecchio}

\section{Introducción}

En este capítulo se trata la modalidad, a partir de una definición gramatical
que opone el \textit{dictum} oracional, esto es, el contenido léxico de una
predicación, al \textit{modus}, o sea, la huella que imprime quien habla sobre
sus dichos, fundamentalmente, en relación con si los presenta, por ejemplo, como
aseveraciones, mandatos o deseos. La modalidad también comprende la polaridad,
ya que las predicaciones pueden ser afirmativas o negativas. Asimismo, la
manifestación de determinados posicionamientos modales está directamente
asociada al grado de conocimiento que quien enuncia tiene sobre los
acontecimientos a los que se refiere: esto se conecta con una noción que ha ido
ganando relevancia en la gramática, la de \termino{evidencialidad}. Si bien
\textit{modalidad epistémica} y \textit{evidencialidad} son nociones distintas,
interactúan en un dominio común que se conoce como el de la
\termino{epistemicidad}.

Al expresarse estos contenidos gramaticales mediante diferentes recursos, el
recorrido de este capítulo supone retomar varios de los temas ya vistos en este
libro: la morfología flexiva de los verbos; el potencial funcional de una clase
de palabras como el adverbio, que se caracteriza por la posibilidad de
comportarse como circunstancial de los verbos, pero también puede desempeñar
funciones extrapredicativas; las perífrasis verbales, sobre todo, las de
modalidad, y, en relación con esto, la organización oracional en tres capas (la
léxica, la flexiva y la informativa).

\section{Definición, delimitación y confluencia con otras categorías}\label{sec:14-definicion}

Indisociable de las propiedades de la estructura lingüística, aparece una
categoría cuyo deslinde acarrea controversia: la \termino{modalidad}. Esta
categoría lingüística codifica, mediante una notable variedad de recursos
gramaticales (véase §\ref{sec:14-recursos}), significados relativos a las
huellas que como hablantes imprimimos en nuestros enunciados, con respecto al
posicionamiento y la perspectiva que adoptamos sobre lo que decimos. Al
respecto, recordemos que, como se planteó en el \capref{6}, mientras que la
oración es una unidad de la gramática, el enunciado es su realización en
situaciones comunicativas concretas. En efecto, no existe la posibilidad de que
comuniquemos un determinado contenido sin que, en simultáneo, se articule con
una manera de enunciarlo: en esto consiste, en síntesis, la oposición entre el
\textit{dictum}, o contenido proposicional de la oración, y el \textit{modus}, o
modalidad, según la cual en todo enunciado (y en su forma abstracta, la oración)
se produce una combinación de un cierto contenido descriptivo con una
determinada intención de quien lo enuncia, que puede, por ejemplo, comunicar o
bien pedir información, pretender influir sobre su(s) interlocutor(es),
manifestar un deseo o formular una pregunta, entre otras
posibilidades.\pendiente{Comentario de la autora: no volver aquí sobre
«oración» frente a «enunciado»; señalar, en cambio, que la oración puede ganar
valores modales secundarios en el uso (p.~ej., interrogativas interpretadas como
pedidos)}

En su forma gramatical, las oraciones proyectan esta presencia de quien enuncia
en las tres capas estructurales que las conforman (véase el \capref{6},
§\ref{sec:06-capas}):

\begin{itemize}
  \item en la capa léxica, a través de la selección de las piezas que establecen
        la predicación sobre los acontecimientos, con su estructura argumental y
        la red de papeles temáticos que señalan a los participantes del evento;
  \item en la capa flexiva, mediante el anclaje temporal, que fundamentalmente
        indica simultaneidad, anterioridad o posterioridad del acontecimiento
        descripto por la predicación con respecto al \textit{hic et nunc} de la
        enunciación;
  \item en la capa complementante, de manera crucial, por ser esta la interfaz
        entre el contenido proposicional y su articulación con el discurso.
\end{itemize}

De ahí que podamos sostener que el \textit{dictum} oracional se estructura
también conforme al \textit{modus} que le imprime quien habla a través de
diferentes medios gramaticales. Pero antes de adentrarnos en ello, podemos
presentar las definiciones de modalidad que han propuesto
\textcite{kovacci1992} y \textcite{rodriguezramalle2005}, respectivamente.

\begin{cita}
La modalidad tiene carácter complejo: incluye al menos dos componentes: la
actitud adoptada con respecto al \textit{dictum} (declarativa, dubitativa, etc.)
y la indicación de la presencia del hablante como tal (como \sig{diciente}), que
implica su responsabilidad por la actitud \parencite[99]{kovacci1992}.
\end{cita}

\begin{cita}
La modalidad recoge las diferencias existentes entre oraciones como reflejo de
las distintas posturas que adopta el hablante, bien con respecto a la verdad del
contenido de la proposición que emite bien con respecto a la manera en que la
expresa. La modalidad se puede definir, por tanto, como la manifestación de las
actitudes u opiniones del hablante ante lo que dice y comprende
\parencite[470]{rodriguezramalle2005}.
\end{cita}

\subsection{Los recursos de la modalidad}\label{sec:14-recursos}

Los recursos que sirven para formalizar la modalidad oracional son múltiples y
variados:

\begin{itemize}
  \item La \textbf{entonación} o línea musical que, en la oralidad, se superpone
        al componente léxico-sintáctico organizado según la disposición de los
        diferentes ítems léxicos y gramaticales. En la escritura, la entonación
        queda esbozada en el empleo de los signos de entonación, y también de los
        de puntuación; así, la predicación de \exref{ex:14-juan1} se interpreta
        como una aserción; la de \exref{ex:14-juan2}, como una interrogación; la
        de \exref{ex:14-juan3}, como una reacción emotiva:
\end{itemize}

\begin{exe}
  \ex\label{ex:14-juan1} Viene Juan.
  \ex\label{ex:14-juan2} ¿Viene Juan?
  \ex\label{ex:14-juan3} ¡Viene Juan\ldots!
\end{exe}

\begin{itemize}
  \item La \textbf{selección del modo verbal}, expresada en la flexión del verbo,
        por ejemplo, en la oposición entre el modo indicativo y el modo
        subjuntivo, lo que conduce a interpretar aquello que comunica la
        profesora como una aserción en \exref{ex:14-profesora1} y como una orden
        en \exref{ex:14-profesora2}:
\end{itemize}

\begin{exe}
  \ex\label{ex:14-profesora1} La profesora dice que los estudiantes llegan
      temprano. (Aseveración).
  \ex\label{ex:14-profesora2} La profesora les dice a los estudiantes que lleguen
      temprano. (Pedido).
\end{exe}

\begin{itemize}
  \item El repertorio de \textbf{auxiliares de modalidad} presentes en las
        perífrasis verbales (véase el \capref{11}, §\ref{sec:11-modales}),
        unidades conformadas por dos verbos, uno ---en forma finita--- que
        aporta información de índole gramatical y otro ---en forma no finita---
        que manifiesta el contenido léxico, como expresiones de diferentes
        modalidades expresivas del hablante, \exref{ex:14-conducir} y
        \exref{ex:14-llover}:
\end{itemize}

\begin{exe}
  \ex\label{ex:14-conducir}
    \begin{xlist}
      \ex \textit{Puedo} conducir porque tengo licencia. (\sig{capacidad} o
          \sig{habilitación}).
      \ex \textit{Tengo que} conducir porque soy el conductor asignado.
          (\sig{necesidad}).
    \end{xlist}
  \ex\label{ex:14-llover} Mañana \textit{puede} llover. (\sig{posibilidad}).
\end{exe}

\begin{itemize}
  \item Ciertas \textbf{construcciones} también sugieren interpretaciones
        modales; así, las estructuras llamadas \textit{cuasirreflejas
        mediopasivas e impersonales} (véase el \capref{10},
        §§\ref{sec:10-mediopasivas} y \ref{sec:10-impersonales}), por ejemplo,
        se caracterizan por desencadenar lecturas deónticas
        \exref{ex:14-prenda} o epistémicas \exref{ex:14-incidentes}:
\end{itemize}

\begin{exe}
  \ex\label{ex:14-prenda} Esta prenda \textit{se lava a mano}. (\sig{debe ser
      lavada a mano}).
  \ex\label{ex:14-incidentes} A diferencia de los accidentes, los incidentes
      \textit{se evitan}. (\sig{son evitables} / \sig{pueden ser evitados}).
\end{exe}

\begin{itemize}
  \item En los márgenes de la predicación, la llamada \textbf{periferia}, que se
        configura en torno a la predicación primaria establecida a partir del
        verbo léxico, sus argumentos y sus circunstantes temporalizados, se
        caracteriza por alojar una serie de elementos que tienen ámbito
        habitualmente sobre todo el núcleo oracional y que expresan así diversos
        posicionamientos modales, materializados en sintagmas adverbiales,
        preposicionales o bien en oraciones subordinadas \exref{ex:14-nino}:
\end{itemize}

\begin{exe}
  \ex\label{ex:14-nino} \textit{Claramente} / \textit{por ahí} / \textit{según
      parece}, el niño descansó mal anoche: no deja de bostezar.
\end{exe}

Por caso, en \exref{ex:14-nino} se suscita la lectura de que los dichos sobre
cómo fue el descanso del niño se fundan, no en la comprobación directa de quien
habla, sino en una conclusión que obtiene a partir de ciertos indicios: los
continuos bostezos.

\begin{itemize}
  \item También se expresa la modalidad mediante \textbf{ítems léxicos
        específicos}, tales como nombres (\forma{posibilidad,
        probabilidad\ldots}), adjetivos (\forma{evitable, posible, probable,
        eventual\ldots}) y adverbios e índices de actitud (\forma{tal vez,
        acaso, posiblemente, capaz (que)\ldots}).
\end{itemize}

\section{La modalidad como expresión de la fuerza ilocutiva}\label{sec:14-ilocutiva}

De acuerdo con una mirada interlingüística, esto es, un enfoque en que se buscan
las propiedades comunes de las lenguas, es complejo definir la modalidad, dado
que las distintas lenguas no apelan exactamente a los mismos recursos ni
codifican el mismo tipo de contenidos modales. Sin embargo, lo que
\textcite{palmer1986} sugiere es que la distinción general entre contenido
proposicional y modalidad permite pensar en un paralelismo cercano al de
\termino{acto locutivo} y \termino{acto ilocutivo}. Esta última distinción,
propuesta por \textcite{austin1971}, se ha convertido en una herramienta
fundamental para estudiar la distancia entre lo que se dice literalmente y lo
que efectivamente se comunica, de modo que el acto locutivo supone la
articulación lingüística de un enunciado y la \termino{fuerza ilocutiva} refiere
a la intención o el propósito de ese acto de habla.

Por ejemplo, dos enunciados como \exref{ex:14-ordena} y \exref{ex:14-porque},
con distinta configuración gramatical (una oración exhortativa con verbo en modo
imperativo, en \exref{ex:14-ordena}; una interrogación negativa con \forma{por
qué} más verbo en modo indicativo, en \exref{ex:14-porque}), pueden servir a la
expresión de una misma intención: la de transmitir una orden, un pedido, una
sugerencia.

\begin{exe}
  \ex\label{ex:14-ordena} Orden\textit{á} tu cuarto.
  \ex\label{ex:14-porque} ¿\textit{Por qué no} orden\textit{ás} tu cuarto?
\end{exe}

\subsection{Modalidad epistémica y modalidad deóntica}\label{sec:14-epistemica}

Asimismo, \textcite{palmer1986}, luego de analizar una serie de nociones que
habitualmente se asocian a la de la modalidad, tales como \sig{subjetividad},
\sig{epistémico} / \sig{deóntico}, \sig{factualidad}, \sig{posibilidad} /
\sig{necesidad}, termina inclinándose por distinguir entre modalidad epistémica
y modalidad deóntica, como dos polos básicos que podemos suponer presentes en
todas las lenguas.

\begin{destacado}
La \termino{modalidad epistémica} no solo se aplica a los sistemas modales que
involucran las nociones de \sig{posibilidad} y \sig{necesidad} (como en la
lógica clásica), sino también a todo sistema modal que indique el grado de
compromiso de quien habla con lo que dice, de acuerdo con el nivel de
conocimiento que posee de los acontecimientos a los que se refiere. La
\termino{modalidad deóntica}, por otro lado, es aquella que se asocia a la
expresión de las capacidades, obligaciones, prohibiciones o permisos, por lo
cual incluye dos rasgos: la \sig{subjetividad} (las diversas marcas del «yo» en
aquello que dice) y la \sig{no factualidad}, esto es, la no asertividad: quien
habla no busca transmitir información ni aludir a acontecimientos de la
«realidad», sino que tiene otros propósitos vinculados con lo volitivo.
\end{destacado}

\subsection{Clases de modalidad y tipos de oraciones}\label{sec:14-clases}

De acuerdo con la definición clásica de modalidad como la «actitud subjetiva» del
hablante, se reconoce que las estructuras oracionales codifican, a través de
diferentes recursos formales ---suprasegmentales y morfosintácticos--- y de
ciertos índices léxicos, distintas posiciones enunciativas frente al
\textit{dictum}: el hablante «asevera, conjetura, exhorta, expresa un deseo»
\parencite[98]{kovacci1992}.

La tabla~\ref{tab:14-kovacci} muestra la propuesta de \textcite{kovacci1992},
quien a las dos modalidades distinguidas por \textcite{palmer1986} añade la
modalidad afectiva y se vale de estas tres modalidades para establecer las
diferentes configuraciones oracionales posibles y determinar sus propiedades
formales e índices léxicos característicos.

{\scriptsize\setlength{\tabcolsep}{2pt}
\begin{longtable}{@{}>{\raggedright\arraybackslash}p{23mm}
    >{\raggedright\arraybackslash}p{25mm} >{\raggedright\arraybackslash}p{45mm}
    >{\raggedright\arraybackslash}p{30mm}@{}}
  \caption{Modalidades y configuraciones oracionales \parencite[a partir
    de][]{kovacci1992}}\label{tab:14-kovacci}\\
  \toprule
  \textbf{Modalidades} & \textbf{Configuraciones oracionales} &
    \textbf{Propiedades formales} & \textbf{Índices léxicos} \\
  \midrule
  \endfirsthead
  \toprule
  \textbf{Modalidades} & \textbf{Configuraciones oracionales} &
    \textbf{Propiedades formales} & \textbf{Índices léxicos} \\
  \midrule
  \endhead
  \textbf{Intelectuales o epistémicas}: están relacionadas con el grado de
    conocimiento de los acontecimientos &
    \textbf{Aseverativas o declarativas}: exponen el \textit{dictum} como un
    hecho real &
    Figura tonal que se desarrolla sobre el rango del tono básico y juntura
    terminal descendente. Modo indicativo. &
    Pueden contar con refuerzos de la certeza y convicción del hablante con
    respecto a sus dichos: \textit{por supuesto, ciertamente}\ldots \\
  \addlinespace
  & \textbf{Dubitativas}: ubican el \textit{dictum} en el plano de lo posible,
    probable o hipotético &
    Figura tonal que no es distinta de la declarativa. Pueden incluir formas del
    modo indicativo con interpretación inferencial (futuro, condicional) y del
    modo subjuntivo, además de perífrasis verbales de modalidad (con
    \textit{poder, deber, tener que}). &
    Con índices de actitud: \textit{tal vez, a lo mejor, posiblemente, capaz
    que}\ldots \\
  \addlinespace
  & \textbf{Interrogativas}: indican la actitud de no certeza o de
    desconocimiento con respecto a todo el \textit{dictum} (preguntas totales) o
    a una parte (preguntas parciales) &
    Entonación interrogativa con diferentes configuraciones. Se opone a la
    entonación declarativa por una figura de movimientos alternativos del tono,
    que van desde los agudos a los graves o viceversa. Modo indicativo. &
    Con pronombres interrogativos (\textit{qué, quién}\ldots) y algunos otros
    índices: \textit{¿no?, ¿verdad?} \\
  \midrule
  \textbf{Volitivas}: están asociadas a lo deóntico &
    \textbf{Exhortativas} &
    Presentan variantes tonales e intervalos de altura e intensidad acentual
    según manifiesten un ruego o un mandato categórico. El modo propio de esta
    modalidad es el imperativo. El subjuntivo «completa» las formas vacantes de
    imperativo, aunque también el presente, el futuro del modo indicativo e
    incluso el infinitivo pueden manifestar actitud exhortativa. &
    Diversas interjecciones: \textit{¡vamos!, ¡por favor!, ¡dale!} \\
  \addlinespace
  & \textbf{Desiderativas} &
    Poseen una figura tonal semejante a la declarativa, aunque el hablante puede
    imprimir un particular grado de énfasis. En modo subjuntivo. Con el
    indicativo, solamente la perífrasis \textit{deber / tener que} + infinitivo
    en el pretérito imperfecto o el condicional puede expresar valor
    desiderativo. &
    Con índices desiderativos como \textit{que, ojalá (que)}, en posición
    inicial. \\
  \midrule
  \textbf{Afectivas}: expresan actitudes emocionales (alegría, disgusto,
    sorpresa, reproche\ldots) &
    \textbf{Exclamativas} &
    Las figuras tonales pueden adoptar diferentes configuraciones (ascendentes,
    descendentes, circunflejas) con un refuerzo de intensidad sobre el foco de
    la carga emotiva. Generalmente, se construyen en indicativo. &
    Con pronombres exclamativos, el relativo \textit{lo que}, intensificadores
    como \textit{tan, tanto, tal}\ldots\ y algunas interjecciones. \\
  \bottomrule
\end{longtable}}

Independientemente de la propuesta de clasificación oracional expuesta, es
necesario tener presente que es habitual cierta superposición de matices
modales. Por caso, ante una estructura como \forma{Que pase el que sigue}, es
posible interpretar que existe una intersección entre el deseo (\sig{quiero que
pase el que sigue}) y el mandato (\sig{pido que pase el que sigue}); asimismo,
frente a \forma{¿Cómo pudo haber hecho eso?}, se produce un cruce entre
interrogación y exclamación.

\begin{ejercicio}\label{ej:14-canciones}
Teniendo en cuenta los elementos destacados en cursiva en los fragmentos de
canciones de abajo, identifique la modalidad, el tipo de oración involucrada y
los recursos que la ponen de manifiesto. Dado que los diferentes medios se
integran en la estructuración lingüística, considere tanto los rasgos prosódicos
como la información morfosintáctica y léxica. Especifique cómo contribuye cada
uno de ellos en la interpretación modal predominante de la oración en que
aparecen.

% Letras de canciones protegidas: se incorporan a mano desde el original.
\textoaincorporar{los fragmentos de 1. \textit{Latinoamérica} (Calle 13), con las
formas destacadas a. \textit{se comparte}, b. \textit{que viva}, c. \textit{no
puedes comprar}; 2. \textit{El pianista del gueto de Varsovia} (J.~Drexler), con
a. \textit{tengo}, b. \textit{pude haber sido}, c. \textit{estoy, estabas}; 3.
\textit{Inoportuna} (J.~Drexler), con a. \textit{quien no lo sepa\ldots\
aprenderá}, b. \textit{para, espera, avisa}, c. \textit{tantos, tan}, d.
\textit{eran, decía, Cautela}, e. \textit{cuánta}, f. \textit{tan}}

\begin{resuelto}
1. a. Modalidad deóntica: la predicación \forma{Aquí se comparte}, solo posible
en 3.ª persona, se interpreta \sig{(las cosas) deben ser compartidas}; la
construcción mediopasiva \forma{se comparte} desencadena una lectura de
genericidad. b. Modalidad deóntica, con \forma{que} como índice desiderativo,
que exige el verbo flexionado en modo subjuntivo. c. Modalidad deóntica: la
perífrasis verbal con \forma{poder} en forma negativa indica
\sig{imposibilidad}, y el verbo en 2.ª persona del singular presenta
interpretación genérica (\sig{nadie puede comprar mi vida}).
\end{resuelto}
\end{ejercicio}

\subsection{Una relación compleja: modo y modalidad}\label{sec:14-modo}

Otro punto de confusión en este tema reside en la confluencia entre el
\termino{modo} y la \termino{modalidad}, algo que advierte
\textcite{dahl2015}, también desde un enfoque interlingüístico. El término
\textit{modalidad} alude al amplio dominio que relaciona el contenido
proposicional con la realidad e involucra nociones como \sig{necesidad},
\sig{posibilidad}, \sig{probabilidad}. El \textit{modo}, en cambio, se refiere a
distinciones gramaticales, o sea, a las inflexiones que se incluyen en las
categorías de tiempo, aspecto, modo y evidencialidad, que se expresan en la
sigla \textsc{tame} (véase §\ref{sec:14-evidencialidad}). Estas categorías, en
algunas lenguas como el español, se manifiestan por lo general en los verbos. El
modo, en ese sentido, no siempre resulta de una elección «libre» de quien habla,
sino que suele depender del tipo de oración en que el verbo se inserta; por
ejemplo, la oración desiderativa con \forma{ojalá que} exige, en el español, la
flexión verbal en subjuntivo \exref{ex:14-ojala}:

\begin{exe}
  \ex\label{ex:14-ojala} \textit{Ojalá} hay\textit{a} buen tiempo en el fin de
      semana.
\end{exe}

\textcite{dahl2015} argumenta que, típicamente, en el francés o en el alemán,
por caso, se usa el modo subjuntivo para marcar la \sig{no aserción} de ciertas
oraciones subordinadas: en el francés, aquellas que son complementos de los
predicados de \sig{deseo}; en el alemán, el subjuntivo se emplea, en cambio, en
las oraciones que expresan discurso indirecto.\footnote{Se denomina
\textit{discurso directo} (DD) «el que reproduce de forma literal palabras
pronunciadas, así como pensamientos transcritos en la misma forma en que se
originan» \parencite[§\,43.9a]{rae2025}, y \textit{discurso indirecto} (DI) al
que presenta la voz del hablante que narra lo que él mismo u otras personas han
dicho o pensado; estos contenidos aparecen en subordinadas sustantivas
introducidas por la conjunción subordinante \forma{que}, ausente del DD
\parencite[§\,43.9c]{rae2025}.} Aunque interlingüísticamente las formas
flexivas que no pertenecen al indicativo (considerado el modo «neutro» en el
conjunto de oposiciones que constituye con el subjuntivo y el imperativo en el
español) reflejan una considerable variación, se puede deducir que lo que sí
tienen en común es que se usan en general para expresar la \sig{no aserción}; por
ejemplo, para distinguir lo \sig{posible} de lo \sig{contrafáctico} (contrario a
los hechos efectivamente acaecidos), como en los períodos condicionales, que son
construcciones de dos constituyentes ---una prótasis o condicionante y una
apódosis o condicionado--- que ponen en un vínculo hipotético dos eventos,
\exref{ex:14-parque1} y \exref{ex:14-parque2} (véase el \capref{13},
§\ref{sec:13-condicionales}).

\begin{exe}
  \ex\label{ex:14-parque1} Si \textit{está} soleado, \textit{iré} al parque.
  \ex\label{ex:14-parque2} Si \textit{hubiera estado} soleado, \textit{habría ido}
      al parque.
\end{exe}

Mientras que en \exref{ex:14-parque1} ambos eventos, tanto el \forma{estar
soleado} como el \forma{ir al parque}, se presentan como hechos posibles, el
encadenamiento condicional de \exref{ex:14-parque2}, en cambio, se interpreta
como contrario a los hechos en el pasado: \forma{no estuvo soleado}, por lo
tanto, \forma{no fui al parque}.

\begin{recuadro}[Para ampliar]
\textcite{dahl2015} sostiene que en muchas lenguas emparentadas por razones
genéticas (o sea, lenguas que comparten un «antepasado común»), como el tagalo,
el sundanés o el javanés, de la familia austronesia ---habladas en la isla de
Madagascar, el archipiélago malayo y Oceanía---, se emplean formas de futuro
para comunicar la \sig{irrealidad}. De esto se puede desprender, entonces, la
existencia de un cierto patrón interlingüístico común que opone lo
\textit{realis} frente a lo \textit{irrealis}; y también que este último ámbito
de la manifestación de lo irreal revela una notable variación en lo que atañe a
las funciones específicas que recubre y al repertorio de medios disponibles
según la lengua que analicemos.\pendiente{El original dice «sudanés»}
\end{recuadro}

Ahora bien, ¿qué significa que se gramaticalice la modalidad en una lengua? Así
como algunos verbos plenos pueden ampliar su uso a un comportamiento auxiliar en
las perífrasis verbales, «cediendo» significado léxico y especializándose en
contenido gramatical, también los contenidos gramaticales de los paradigmas
flexivos codifican una información específica de índole muy abstracta. La
regularidad de las oposiciones manifestadas en estos conjuntos cerrados de formas
permite la contraposición interna dentro de un rango posible de valores, que se
expresan necesariamente en determinadas clases de palabras que establecen
conexiones morfosintácticas \parencite{lehmann2002,giammatteo2018}. Por caso, en
el español, la oposición de modo indicativo, subjuntivo e imperativo distingue
tres valores: \forma{cantás / cantes / cantá; comés / comas / comé; vivís / vivas
/ viví}, para la segunda persona voseante del singular de \forma{cantar, comer} y
\forma{vivir}, respectivamente. En esta oposición queda codificado el contraste
\sig{afirmación} / \sig{deseo} / \sig{mandato}, como las diferentes actitudes de
quien habla ante un mismo evento. Forzosamente, estas oposiciones de modo son
exhibidas por los verbos finitos y están vinculadas con las estructuras
sintácticas que integran; por ejemplo, el verbo \forma{comunicar} se combina con
el modo indicativo, en \exref{ex:14-comunica1}, dada la compatibilidad del acto
de «comunicar» con la introducción de una aseveración, que es lo codificado en
la flexión en modo indicativo de \forma{está}; pero rechaza, en este contexto,
el subjuntivo \exref{ex:14-comunica2}, puesto que es incompatible lo que
codifica \forma{esté} (básicamente, no aserción) con el predicado introductor.

\begin{exe}
  \ex\label{ex:14-comunica1} Comunica que est\textit{á} bien.
  \ex[*]{\label{ex:14-comunica2} Comunica que est\textit{é} bien.}
\end{exe}

\begin{ejercicio}\label{ej:14-modos}
A partir de las formas verbales presentes en los siguientes pares de oraciones,
explique qué significados modales se contraponen en ellas, según se muestra en
el ejemplo resuelto~1.

\begin{exe}
  \ex Los novios \textit{son} felices. / \textit{Sean} felices los novios.
  \ex El Señor los \textit{protege}. / El Señor los \textit{proteja}.
  \ex Esta semana \textit{escribiremos} el capítulo. / Esta semana
      \textit{escribiríamos} el capítulo.
  \ex ¡Ojalá \textit{venga} a la fiesta! / ¡Ojalá \textit{viniera} a la fiesta!
  \ex ¡Ojalá \textit{viniera} a la fiesta! / ¡Ojalá \textit{hubiera venido} a la
      fiesta!
  \ex Capaz \textit{juega} mañana. / Capaz \textit{juegue} mañana.
\end{exe}

\begin{resuelto}
1. La oposición de modo indicativo (\forma{son}) y modo subjuntivo
(\forma{sean}) contrapone modalidad epistémica aseverativa frente a modalidad
deóntica de deseo.
\end{resuelto}
\end{ejercicio}

\section{¿Modalidad o evidencialidad?}\label{sec:14-evidencialidad}

Según hemos planteado, la delimitación de la modalidad en una perspectiva
interlingüística no está exenta de controversias, entre otras razones porque la
difusión de la existencia de lenguas que exhiben un \termino{sistema gramatical
evidencial} ha provocado que se abriera el debate sobre si la modalidad y la
evidencialidad constituyen una misma categoría o si son dos categorías
diferentes. Por caso, \textcite{barnes1984} explica que en el tuyuca, una lengua
hablada en Brasil y Colombia, se gramaticalizan los siguientes significados: i)
\textit{evidencia visual}; ii) \textit{evidencia no visual} (proveniente de algún
otro sentido que no sea la vista, predominantemente el oído); iii)
\textit{apariencia} o \textit{deducciones} hechas a partir de la evidencia; iv)
evidencia de segunda mano, es decir, lo que se conoce como el \textit{reporte} o
la \textit{evidencia citativa}, y v) \textit{evidencia inferida}, que se expresa
cuando el hablante posee un conocimiento previo de un estado de cosas o de
ciertos patrones de conducta generales o habituales.

Esto implica que los hablantes del tuyuca, al decir algo, no pueden no marcar
con un afijo gramatical ---así como nosotros en español lo hacemos con el tiempo
o con el modo--- cuál ha sido la fuente de acceso a su conocimiento de los
acontecimientos referidos: si esa información proviene de algo que han visto
«con sus propios ojos», que han oído, que deducen de ciertos indicios, que les
han contado o que infieren a partir de otros datos.

Como vemos, entonces, hay lenguas de distintas familias (entre ellas, amerindias,
eslavas como el macedonio, túrquicas\ldots) que expresan algunos de estos
contenidos mediante categorías morfológicas específicas.\pendiente{El original
dice «amerindias, macedónicas, turcas»} Según \textcite{willett1988}, lo más
recurrente es que se distinga, en la morfología verbal, si la fuente para la
evidencia es directa o «atestiguada», indirecta (o citativa) o inferida (es
decir, la que se deriva de asignar causas, en función de una construcción
mental, a las situaciones observadas). Para estas lenguas, las marcas de
evidencia atestiguada proceden ---mediante fenómenos de fusión y reducción
sintáctica y fonológica--- de verbos de percepción sensorial; las de reporte, de
verbos de habla; y las de inferencia, de tiempos perfectos o de verbos de
observación visual. Este hecho muestra que las marcas evidenciales reflejan un
proceso que se conoce como \termino{gramaticalización}
\parencite{company2016}, por el que cierto material léxico (como los verbos de
\sig{percepción} y \sig{comunicación}) se vuelve, mediante un procedimiento de
reanálisis, un afijo gramatical. Una forma de gramaticalización que guarda
semejanzas con lo detectado por Willett se puede comprobar también en el
español: el verbo \forma{haber} en \forma{amar he > amaré, amar has >
amarás\ldots} se transformó, por un cambio que se produce a lo largo de la
historia de la lengua, en el sufijo de flexión del tiempo futuro; es decir,
\forma{haber}, que en un primer estadio «auxilia» al infinitivo en \forma{amar
he} (= \sig{he de amar}), termina integrándose en el verbo léxico, como parte de
la estructura morfológica de palabra.

En el caso de las marcas evidenciales, es importante destacar que, aunque sean
parte de la flexión de los verbos, su alcance es oracional: su ámbito de
operación recae sobre toda la estructura sintáctica, tal como ocurre, en el
español, con la flexión de tiempo, que se expresa en el verbo, pero «temporaliza»
toda la predicación. En efecto, la marcación evidencial especifica la fuente de
la información del hablante para la aserción y, por tanto, concierne a toda la
oración aseverada, aunque solamente se manifieste en un afijo gramatical.

Hemos visto, hasta aquí, que la evidencialidad, concebida como una categoría
lingüística (no meramente gramatical), involucra la fuente de información
(directa o indirecta ---inferencial y reportativa o citativa---), mientras que
la modalidad epistémica se refiere a los significados modales que manifiestan un
soporte, un respaldo completo, parcial o neutro (en aquellos casos en los que la
proposición no es objeto de una evaluación epistémica) de la aserción. Por
consiguiente, aunque se puede reconocer que modalidad y evidencialidad son
categorías conceptualmente diferentes, existe una interacción funcional entre
ambas: la epistemicidad \parencite{bybee2012}.

\subsection{Evidencialidad, deixis y perspectivización}\label{sec:14-deixis}

Por otra parte, en la codificación lingüística de la evidencialidad, no podemos
omitir la importancia de la \termino{deixis}. \textcite[319]{lyons1997} trata la
deixis como una operación de señalamiento referencial, es decir, una indicación
de «la relación que se establece entre las expresiones lingüísticas y lo que
representan en el mundo (o universo del discurso) en ocasiones particulares de
enunciado». En sintonía con esto, \textcite{bermudez2016} defiende que la
evidencialidad supone coordenadas que son propias de la deixis espacial
---\textit{puntos de referencia}, \textit{dirección} y \textit{distancia}---, que
se reinterpretan, metafóricamente, en el dominio evidencial:

\begin{itemize}
  \item los \textit{puntos de referencia} oscilan en un continuo entre el
        hablante, como fuente única, y una fuente general que condensa el
        conocimiento compartido por una comunidad, es decir, su
        \textit{folklore};
  \item la \textit{dirección} es el modo de acceso al conocimiento, que puede ser
        \sig{hacia} el hablante, a través de impresión sensorial, o \sig{desde}
        el hablante, mediante inferencia o razonamiento;
  \item la \textit{distancia} concierne, en un polo, a la información recibida de
        primera mano y, en el polo opuesto, a lo que proviene del rumor, o sea,
        de una fuente poco fidedigna.
\end{itemize}

\begin{destacado}
En función de esta confluencia entre la deixis y el sistema de perspectiva, la
evidencialidad puede ser entendida, según \textcite{hassler2010}, como un
\termino{fenómeno deíctico de perspectivización} desde el punto de vista del
hablante, que el oyente debe reinterpretar desde su propia posición discursiva.
\end{destacado}

Esta perspectivización implícita en la deixis permite distinguir, para la
autora, modalidad epistémica de evidencialidad: la primera aporta
monológicamente la actitud de quien habla; la otra requiere del oyente la
elaboración de la referencia a la fuente de información manifestada por el
hablante y la realización de todos los ajustes necesarios sobre la fiabilidad de
las fuentes.

En síntesis, el reconocimiento de estos sistemas evidenciales ha llevado a
\textcite{palmer1986} a admitir que es frecuente que no pueda establecerse una
distinción clara entre \textit{juicios} (asociados tradicionalmente a la noción
de modalidad en la lógica) y \textit{evidenciales} (marcas gramaticales del
acceso a la información), dado que los juicios de quien habla están naturalmente
atados al grado de fiabilidad de la evidencia que posee.

\section{Evidencialidad: paradigma frente a estrategias}\label{sec:14-estrategias}

En cuanto al español, podemos asumir, como plantea
\textcite{gonzalezvazquez2016}, que la discusión no queda zanjada, porque se
presentan las siguientes dificultades:

\begin{itemize}
  \item persiste la controversia en cuanto a qué relación guarda la modalidad con
        la evidencialidad;
  \item se admite que tanto en los sistemas gramaticalizados como en los no
        gramaticalizados los marcadores evidenciales pueden denotar
        polisémicamente varios significados;
  \item se usaron, a lo largo de la historia, diferentes criterios para
        clasificar los evidenciales, aunque la clasificación más extendida sea la
        de \textcite{willett1988}, que presentamos en
        §\ref{sec:14-evidencialidad}.
\end{itemize}

El recorrido que sigue expone la dificultad de estudiar la modalidad sin atender
a cómo inciden otras categorías asociadas a ella en la flexión verbal, como el
tiempo y el aspecto, y sin tener en cuenta la evidencialidad, categoría que, en
las lenguas románicas, no se organiza estrictamente en torno a oposiciones
flexivas y, en cambio, se expresa mediante una amplia variedad de mecanismos.
Como el español no es una lengua que cuente con un paradigma evidencial
sistematizado en las formas verbales, se considera que pertenece al grupo de
lenguas que desarrollan diversas \termino{estrategias evidenciales}, mediante el
empleo de ciertos ítems léxicos (verbos como \forma{oír, ver, parecer} y
adverbios como \forma{claramente, visiblemente, aparentemente}, entre otros) o
recursos gramaticales diversos, tales como ciertos usos del futuro, del
condicional y del pretérito imperfecto, o la configuración oracional de la
impersonalidad, entre todos los que ilustramos a continuación.

\subsection{El pretérito imperfecto}\label{sec:14-imperfecto}

\textcite[24]{hassler2017}, por ejemplo, analiza el pretérito imperfecto en
textos narrativos y periodísticos como marcador evidencial de que la información
procede de otra fuente, distinta de quien la reproduce, a partir de ejemplos
como \exref{ex:14-papa} y \exref{ex:14-wall}. De esta manera, el locutor se
distancia de sus enunciados y realiza simultáneamente una atenuación del
contenido proposicional:

\begin{exe}
  \ex\label{ex:14-papa} El Papa Francisco \textit{recibía} el lunes en el
      Palacio Apostólico a la última tanda de prelados de las provincias
      eclesiásticas del sur de España, incluidas las Islas Canarias, Ceuta y
      Melilla. (GlossaNet, \textit{ABC}, 13/03/2014).
  \ex\label{ex:14-wall} Wall Street abrió hoy sin tendencia definida y el Dow
      Jones de Industriales \textit{bajaba} un 0,07~\%, con los mercados
      cautelosos a la espera del inicio de la temporada de resultados
      trimestrales. (GlossaNet, \textit{El País}, 10/04/2014).
\end{exe}

Este uso del pretérito imperfecto ha sido relacionado con su función citativa
solo dentro del marco de sus valores pragmáticos y discursivos en el discurso
oral \parencite{reyes1990}; por ejemplo, cuando formulamos una pregunta como
\forma{¿A qué hora hab\textbf{ía} que salir?}, que busca confirmación de una
información ya dicha. Para \textcite{hassler2017}, a partir de una intersección
semántica entre la atenuación y la evidencialidad, el pretérito imperfecto ya no
expresa sus significados primarios o funciones prototípicas
\parencite[véase][]{giammatteo2018trombetta}, sino otros significados
secundarios que están relacionados con la actitud epistémica del hablante frente
a su enunciado y a la fuente de información, lo que se puede subsumir bajo la
noción de \termino{evidencialidad indirecta}.

Este significado secundario del pretérito imperfecto puede derivarse de su
aspecto imperfectivo, por medio del cual se contempla una acción en su
transcurso y desarrollo, lo que hace que la validez de la oración se presente
también como abierta o indeterminada: la aspectualidad imperfectiva conduce a la
modalidad debido a su indeterminación. Sin embargo, en otras lenguas, como el
alemán, donde no hay formas imperfectivas, esta función es asumida por verbos
modales que, además de su función deóntica, tienen usos evidenciales (por
ejemplo, el verbo \forma{wollen}).

Por otra parte, con una clara interpretación atenuativa de modalidad de
\sig{irrealidad}, el pretérito imperfecto puede emplearse en todas las lenguas
románicas como recurso de la cortesía lingüística \parencite[21]{hassler2017},
según los ejemplos reproducidos como \exref{ex:14-fr}, \exref{ex:14-port} y
\exref{ex:14-it}:

\begin{exe}
  \ex\label{ex:14-fr} fr. Je \textit{voulais} vous demander, Monsieur, de me
      permettre d'assister à votre cours.\\
      \sig{[Yo] quería pedirle, señor, me permita asistir a su clase}.
  \ex\label{ex:14-port} port. \textit{Podias} passar este texto a limpo?\\
      \sig{¿Podías pasar este texto en limpio?}
  \ex\label{ex:14-it} it. Scusi, \textit{volevo} un bicchier d'acqua.\\
      \sig{Disculpe, quería un vaso de agua}.
\end{exe}

Como muestran los datos de arriba y las glosas que explicitan el significado,
también en español es un recurso cortés el empleo del imperfecto como mitigador
de la fuerza del mandato. Las formas verbales \forma{quería} y \forma{podía}
indican cortesía en la expresión de lo pedido: \forma{asistir a su clase},
\forma{pasar este texto en limpio} y \forma{pedir un vaso de agua}, en
\exref{ex:14-fr}, \exref{ex:14-port} y \exref{ex:14-it}, respectivamente.

\subsection{El futuro y el condicional}\label{sec:14-futuro}

También \textcite[22]{hassler2017} muestra que el futuro simple puede usarse con
efectos de atenuación, sobre todo en contextos de cortesía, para suavizar una
petición, como en \exref{ex:14-dira}:

\begin{exe}
  \ex\label{ex:14-dira} Usted me di\textit{rá} lo que tengo que hacer. (En vez
      de \textit{Dígame lo que tengo que hacer}).
\end{exe}

En otros contextos, el uso del futuro simple puede mitigar el grado de certeza
con que se expresa un contenido proposicional, como en los empleos consabidos de
las formas de futuro con interpretación temporal de presente
\exref{ex:14-chico}, que sugieren que la aserción se asienta sobre inferencias:

\begin{exe}
  \ex\label{ex:14-chico} Ese chico ahora tend\textit{rá} unos diez años.
\end{exe}

\textcite[30]{hassler2017} se concentra en el \termino{futuro narrativo}
retrospectivo, al que pone en relación con la evidencialidad. Se trata de un uso
infrecuente en español, que aparece para referirse a situaciones pasadas con
respecto al momento de habla, pero que son presupuestas,\footnote{En el análisis
semántico, la presuposición alude al hecho de que hay expresiones lingüísticas
incrustadas dentro de una matriz que no se ve afectada por la negación: así,
tanto \forma{Me interesa que se estén jugando las Olimpiadas} como \forma{No me
interesa que se estén jugando las Olimpiadas} presentan como presupuesto el
hecho de que efectivamente se están realizando los juegos olímpicos. Dicho de
otro modo, las subordinadas introducidas por predicados que expresan reacciones
emotivas (\forma{gustar, interesar, alegrarse\ldots}) no quedan bajo el ámbito
de la negación, y su contenido, que se da por conocido y cuyo valor de verdad no
es cuestionado, es presentado «bajo el prisma de una evaluación, una emoción
[\ldots]» \parencite[§\,25.3z]{rae2025}.} que se exponen como dadas por el
locutor y que, como no son deícticas, son incompatibles con complementos
adverbiales referidos al momento presente \exref{ex:14-ragazzo}.\pendiente{Comentario
de F. Carranza: precisar la relación entre el futuro narrativo, la
evidencialidad y la presuposición}

\begin{exe}
  \ex\label{ex:14-ragazzo} it. Il ragazzo partì all'improvviso. \textit{Tornerà}
      solo dopo tre anni / *fa tre anni.\\
      \sig{El muchacho partió imprevistamente. Regresará solo después de tres
      años / *hace tres años}.
\end{exe}

Con este empleo del futuro narrativo retrospectivo, en general en tercera
persona, quien habla «simula» una dislocación retrospectiva con respecto al
momento de habla. En este sentido, Hassler interpreta que se puede atribuir al
futuro narrativo un valor evidencial: el hablante se presenta en tales
enunciados como alguien que conoce los hechos y comunica sus conclusiones, porque
conoce la historia hasta su final.

Como estamos viendo, la morfología verbal flexiva del español queda disponible
para sumar nuevas lecturas, no solamente las de modalidad, como los usos de
cortesía, sino también las que, más recientemente, tienden a ser analizadas como
empleos evidenciales \parencite{bermudez2016,garcianegroni2016,reyes1990}, como
el futuro inferencial \exref{ex:14-nadie}; el futuro citativo, que remite a algo
ya dicho y adopta valor concesivo \exref{ex:14-fortuna}, y el futuro citativo
refutativo, que se construye con la forma perifrástica
\exref{ex:14-conductor}:\pendiente{El original cita «García Negroni (2016)»;
la entrada es García Negroni y Marcovecchio (2016)}

\begin{exe}
  \ex\label{ex:14-nadie} No hab\textit{rá} nadie todavía, porque la puerta está
      cerrada. (= \sig{Supongo que no habrá\ldots}).
  \ex\label{ex:14-fortuna} Hab\textit{rá} trabajado mucho, pero esa fortuna no se
      hace con el trabajo de un año. (Réplica a una intervención previa, que el
      hablante concede parcialmente).
  \ex\label{ex:14-conductor} A: ---Es un excelente conductor.\\
      B: ---¿Qué \textit{va a ser} buen conductor? Es un irresponsable.
      (Réplica).
\end{exe}

De manera semejante al futuro, también admiten estos usos el condicional y el
imperfecto, cuando se emplean para hacer una suposición (valor inferencial)
\exref{ex:14-impacto} o retomar algo dicho (valor citativo)
\exref{ex:14-llamabas}:

\begin{exe}
  \ex\label{ex:14-impacto} El impacto se hab\textit{ría} producido por el estado
      de embriaguez del que conducía. (Como parte de la información de una
      crónica periodística).
  \ex\label{ex:14-llamabas} ¿Cómo te llama\textit{bas}? Me lo dijiste, ya sé,
      pero olvidé tu nombre.
\end{exe}

\begin{destacado}
De los datos anteriores se desprende la polifuncionalidad de muchas de las formas
verbales del español y, en lo relativo a la evidencialidad, la «ganancia» de
interpretaciones secundarias que sugieren, fundamentalmente, un acceso a la
información que se comunica mediante razonamiento deductivo ---por ejemplo, en
\exref{ex:14-nadie}, \forma{la puerta cerrada} funciona como un indicio para
desencadenar la inferencia \forma{no habrá nadie todavía}--- o a través de la
mención de dichos ajenos.
\end{destacado}

\begin{ejercicio}\label{ej:14-formas}
Identifique las formas verbales destacadas de las siguientes oraciones y
explique, en términos modales y evidenciales, las oposiciones que plantean,
según se muestra en el ejemplo resuelto~1.

\begin{exe}
  \ex \textit{Tenía} sesenta años cuando lo conocí. / \textit{Tendría} sesenta
      años cuando lo conocí.
  \ex Juan no está en la casa: ¿\textit{ha ido} a la facultad esta noche? /
      ¿\textit{habrá ido} a la facultad esta noche?
  \ex ---No encuentro mis llaves.\\
      ---Las \textit{dejaste} / \textit{habrás dejado} en la puerta.
  \ex El gobierno \textit{llamará} / \textit{llamaría} a elecciones
      extraordinarias.
  \ex Si lo \textit{hubiera sabido} / \textit{sabía}, no \textit{habría venido} /
      \textit{venía} a la reunión.
  \ex ¡Buen día! \textit{Quiero} / \textit{Quisiera} saber si hay alguna chance de
      probarla antes de comprar. Tengo una 6º y no me \textit{resulta} /
      \textit{estaría resultando}. (Consulta sobre una boquilla de un
      instrumento musical en un blog).\pendiente{El original dice «una 6*»; el
      asterisco se reemplazó por «º» para evitar que se lea como juicio de
      gramaticalidad: verificar}
\end{exe}

\begin{resuelto}
1. Las formas verbales \forma{tenía} (imperfecto) y \forma{tendría}
(condicional) oponen \sig{conocimiento efectivo} frente a \sig{conocimiento
inferido}.
\end{resuelto}
\end{ejercicio}

\subsection{Los auxiliares de perífrasis}\label{sec:14-auxiliares}

Algunos de los auxiliares de las perífrasis verbales poseen empleos modales y
evidenciales. Como se vio en el \capref{11}, \forma{poder, deber
(de)}\footnote{\forma{Deber} con \forma{de} o sin \forma{de} se revela, en el
uso del español general, como un fenómeno afectado por la variación. En lo
normativo, se recomienda \forma{deber} para el empleo que indica necesidad u
obligación, y \forma{deber de} para la expresión de probabilidad, pero el uso no
siempre registra esta diferencia.} y \forma{tener que} son habitualmente
reconocidos como índices de la modalidad deóntica, que manifiesta
\sig{habilidad}, \sig{permiso}, \sig{necesidad} u \sig{obligación}, y también de
la modalidad epistémica, que comprende desde lo \sig{posible} hasta lo
\sig{probable}. Pero, además, dichos auxiliares también pueden analizarse como
manifestación del grado de compromiso que el hablante asume según la fiabilidad
de los indicios que le permiten realizar ciertas deducciones
\parencite{cornillie2007}.

\begin{exe}
  \ex\label{ex:14-billetera} La billetera \textit{puede / debe de / tiene que}
      haber quedado en el auto.
\end{exe}

Los diferentes auxiliares señalan la menor o mayor fuerza de las presunciones de
quien habla: con \forma{poder} se sugiere que los indicios son débiles para una
aseveración, mientras que con \forma{deber} y \forma{tener que} se indica que
las «pistas» sobre dónde quedó la billetera permiten una deducción con mayor
grado de probabilidad. También con un auxiliar como \forma{soler} se produce una
imbricación entre el contenido aspectual de \sig{frecuencia} y el de modalidad
epistémica \parencite{giammatteo2011}.

\begin{exe}
  \ex\label{ex:14-suele} El profesor \textit{suele} venir por las tardes, pero hoy
      no concurrió al instituto.
\end{exe}

Asimismo, un predicado de \sig{comunicación} como \forma{prometer}, según vimos
en el \capref{11}, muestra un empleo auxiliarizado que sugiere un cierto grado de
adhesión de quien habla con respecto a la proyección de un acontecimiento
futuro:

\begin{exe}
  \ex\label{ex:14-promete} El encuentro \textit{promete} ser una fiesta.
\end{exe}

Nuevamente, así como planteamos con respecto a la morfología verbal, comprobamos
en algunos auxiliares una cierta hibridación entre contenidos modales y otros,
como el aspecto \exref{ex:14-suele}, o la descripción de la fuerza ilocutiva de un
acto de habla \exref{ex:14-senalo2}.\pendiente{El original remite aquí a (33) y
(34), que ilustran otra cosa} Ciertamente, la morfología verbal y las
perífrasis verbales constituyen dos extremos de la expresión de contenidos
gramaticales: en un caso, se recurre a la afijación; en otro, a la combinación de
dos palabras en una misma unidad verbal. Pero el comportamiento de ambos
mecanismos es complementario \parencite{marcovecchio2010}:

\begin{exe}
  \ex\label{ex:14-senalo1} El profesor le señaló que complet\textit{ara} el
      ejercicio.
  \ex\label{ex:14-senalo2} El profesor le señaló que \textit{debía} completar el
      ejercicio.
\end{exe}

Tanto en \exref{ex:14-senalo1} como en \exref{ex:14-senalo2}, la oración
dependiente de \forma{señalar} expresa la fuerza deóntica, pero en
\exref{ex:14-senalo1} lo hace mediante el uso del subjuntivo \forma{completara} y
en \exref{ex:14-senalo2} recurre al auxiliar \forma{debía} en modo indicativo, lo
que muestra la conexión supletoria entre la flexión en modo subjuntivo y el verbo
auxiliar modal.

\begin{ejercicio}\label{ej:14-parafrasis}
Establezca paráfrasis del significado de las siguientes oraciones utilizando
perífrasis verbales, según se muestra en el ejemplo resuelto~1.

\begin{exe}
  \ex El ACV \textit{es prevenible} y \textit{tratable}.
  \ex Eso \textit{es indiscutible}.
  \ex Las emociones \textit{no se medican}: \textit{se procesan}.
  \ex Esta prenda \textit{no es lavable} en máquina.
  \ex \textit{No se vive} del aire, sin trabajar.
  \ex En esta sala, \textit{no se grita}.
\end{exe}

\begin{resuelto}
1. El ACV \textit{puede ser prevenido y tratado} / \textit{se puede prevenir y
tratar}.
\end{resuelto}
\end{ejercicio}

\begin{ejercicio}\label{ej:14-gramaticalizacion}
Explique, en términos de gramaticalización, la distribución del modo subjuntivo
de la forma simple destacada o de la forma perifrástica en indicativo.

\begin{exe}
  \ex Dados los resultados del laboratorio, el médico le indicó al paciente que
      \textit{hiciera} (*\textit{hacía}) una dieta.
  \ex Dados los resultados del laboratorio, el médico le indicó al paciente que
      \textit{tenía que hacer} (*\textit{tuviera que hacer}) una
      dieta.\pendiente{¿*\textit{tuviera que hacer}? Parece aceptable:
      revisar el juicio}
\end{exe}
\end{ejercicio}

\subsection{Adverbios, locuciones y oraciones subordinadas}\label{sec:14-adverbios}

Ahora bien, el español, como las demás lenguas románicas, también se vale de
adverbios o locuciones equivalentes, así como de ciertas oraciones subordinadas
que, ubicados en la periferia oracional, se comportan como «estrategias» para
codificar contenidos evidenciales \parencite{marcovecchio2001}. En lo formal,
estos constituyentes quedan delimitados por un contorno entonativo propio,
indicio de su alcance oracional, extrapredicativo, y suelen ubicarse en el margen
preverbal, aunque justamente la delimitación melódica propia les brinda gran
movilidad posicional. En lo interpretativo, se trata de recursos que, por un
desplazamiento metafórico desde contenidos léxicos o menos gramaticalizados
hacia otros que manifiestan significados gramaticales o más gramaticalizados,
dejan de funcionar en el ámbito de la predicación y pasan a aludir a la esfera
del hablante.

Así, \forma{por ahí} es un SP que señala un contenido espacial en
\exref{ex:14-porahi1}, pero adquiere otra interpretación, vinculada a la
modalidad y a las inferencias que realiza quien enuncia, en
\exref{ex:14-porahi2}. Las huellas de la gramaticalización también se advierten
en la pronunciación, ya que en la interpretación modal ambos elementos pasan a
emitirse como una unidad y con desplazamiento acentual.

\begin{exe}
  \ex\label{ex:14-porahi1} Las llaves están \textit{por ahí}. (Pron.
      \fonol{poɾ aˈi}).
  \ex\label{ex:14-porahi2} \textit{Por ahí} (≡ \textit{capaz, a lo mejor}) viene
      esta tarde a visitarme. (Pron. \fonol{poɾ aˈi} \textasciitilde\
      \fonol{poɾˈai}).
\end{exe}

Justamente, acudiendo a la noción de evidencialidad, se advierte que la locución
\forma{por ahí} puede manifestar tanto el carácter citativo de un enunciado (a
cuyo contenido, habitualmente, se contrapone la opinión de quien habla)
\exref{ex:14-misa} como la lectura hipotético-inferencial que se le imprime a un
enunciado \exref{ex:14-calle} \parencite{marcovecchio2015}:\pendiente{El
original cita «Marcovecchio, 2015, p. 2020»}

\begin{exe}
  \ex\label{ex:14-misa} «\textit{Por ahí se dice} que no tendría que haber ido [a
      la misa]. Pero yo fui invitado y me parece que era una buena
      oportunidad» para homenajear a Juan Pablo II, uno de los principales
      responsables de que la Argentina y Chile no hayan entrado en guerra hace 30
      años. (lanacion.com.ar, 24/12/2008).
  \ex\label{ex:14-calle}
    \begin{xlist}
      \ex\label{ex:14-calle-a} Para finalizar, González advirtió: «No hay delito,
          ella está declarando porque desapareció. Si ella estuvo durmiendo en la
          calle, \textit{por ahí} no vio televisión [\ldots]».
          (elintransigente.com, 10/07/2018).
      \ex\label{ex:14-calle-b} (Estoy en condiciones de deducir que) Dado que
          durmió en la calle, es probable que no haya visto televisión.
    \end{xlist}
\end{exe}

Además, en \exref{ex:14-calle-a}, en que \forma{por ahí} manifiesta un vínculo
inferencial entre \forma{haber estado durmiendo en la calle} y \forma{no ver
televisión}, se puede llegar a establecer una correspondencia con las oraciones
subordinadas llamadas \textit{causales de la enunciación}, que tienen valor modal
en tanto se vuelven recursos epistémico-evidenciales \exref{ex:14-calle-b}
\parencite{marcovecchio2001}.

Se trata de un caso de evidencialidad inferida (el modo de acceso es el
razonamiento, no la experiencia sensible), y no es casual que en la locución
\forma{por ahí} sea un elemento deíctico el factor crucial para disparar la
lectura evidencial. El acercamiento a la esfera personal y subjetiva produce un
debilitamiento asertivo que es propio de procesos que convierten ciertos
elementos, en principio concernientes a la descripción de un estado de cosas en
el mundo, en recursos gramaticales asociados a la expresión de significados
relativos a la situación de habla y de los interlocutores: lo que se conoce como
\termino{pragmaticalización}.

Con respecto a las oraciones subordinadas (estudiadas en el \capref{13}), su
empleo como marcadores de la modalidad o de la evidencialidad no constituye una
excepción, sino más bien lo contrario: los períodos oracionales constituidos por
prótasis y apódosis ---todos los relacionados con el ámbito de la
\sig{causalidad}, que incluye subordinadas causales, consecutivas, finales,
condicionales y concesivas (véase el \capref{13},
§\ref{sec:13-logicas})--- son favorables a desencadenar lecturas que conciernen a
las inferencias que realiza el hablante \exref{ex:14-delpo} o a su propia
situación como «diciente» \exref{ex:14-murray}
\parencite{kovacci1992,marcovecchio2016}:

\begin{exe}
  \ex\label{ex:14-delpo}
    \begin{xlist}
      \ex Delpo se la jugó. Y jugó un partidazo. Le tocó sufrirlo a Ferrer, que
          también tenía sus inconvenientes ---un esguince en el tobillo
          izquierdo, por el que se infiltró para jugar---, pero no por eso era un
          adversario menos peligroso; a fin de cuentas, sus buenas razones tenía
          el alicantino \textit{para haber llegado hasta aquí también}.
          (\textit{La Nación}, Buenos Aires, 04/07/2013).
      \ex \sig{Si llegó hasta aquí también, (será porque / puedo deducir que)
          sus buenas razones tenía el alicantino}.
    \end{xlist}
  \ex\label{ex:14-murray}
    \begin{xlist}
      \ex «\textit{Para ser sincero}, ni siquiera yo lo sé», dijo Murray.
          (Canchallena, \textit{La Nación}, Buenos Aires, 21/01/2014).
      \ex \sig{Si hablo con sinceridad, (te digo que) ni siquiera yo lo sé}.
    \end{xlist}
\end{exe}

\begin{ejercicio}\label{ej:14-subjuntivo}
Explique qué significados modales manifestados en el modo verbal (e incluso en
las conjunciones) se contraponen en los siguientes pares de oraciones (por
ejemplo, \sig{conocimiento de primera mano} frente a \sig{conocimiento
inferido}, \sig{aserción} frente a \sig{deseo}, etc.), según se muestra en el
ejemplo resuelto~1.

\begin{exe}
  \ex Insiste en que los chicos siempre \textit{se portan} bien. / Insiste en que
      los chicos siempre \textit{se porten} bien.
  \ex Siento que \textit{cierra} esa puerta. / Siento que \textit{cierre} esa
      puerta.
  \ex Creo que no \textit{asistirá} a la clase. / No creo que \textit{asista} a
      la clase.
  \ex El juez no cree que \textit{existe} un delito. / El juez no cree que
      \textit{exista} un delito.
  \ex Depende de la historia que \textit{tengo} entre manos. / Depende de la
      historia que \textit{tenga} entre manos.
  \ex Literariamente, el fútbol es muy fértil y está bueno siempre que nos
      \textit{conduce} a otros lugares. / Literariamente, el fútbol es muy
      fértil y está bueno siempre que nos \textit{conduzca} a otros lugares.
  \ex \textit{Si} no \textit{presentás} el pasaporte, no salís del país. /
      \textit{Como} no \textit{presentes} el pasaporte, no salís del país.
\end{exe}

\begin{resuelto}
1. Las formas verbales en oposición \forma{portan} (indicativo) / \forma{porten}
(subjuntivo) contraponen modalidad epistémica aseverativa a modalidad deóntica
exhortativa.
\end{resuelto}
\end{ejercicio}

\subsection{Marcadores discursivos}\label{sec:14-marcadores}

Si bien el término \textit{marcadores discursivos} (concepto introducido en el
\capref{12}, §\ref{sec:12-marcadores}) se emplea habitualmente en el análisis
gramatical, no alude a una clase de palabras (como las categorías nombre o verbo,
por caso), sino a un conjunto categorialmente heterogéneo de ítems (adverbios,
conjunciones, locuciones prepositivas, incluso verbos, todos con otras funciones
vinculadas con la interacción) que comparten entre sí propiedades discursivas.
Estos atributos pueden sintetizarse en la «aislabilidad» respecto de la
sintaxis predicativa y la no dependencia del significado léxico que aporta la
predicación sobre la que inciden; el impacto interaccional sobre los intercambios
comunicativos, por el que se evoca a los participantes de la situación
dialógica; y su comportamiento ilocutivo, en tanto expresan el «estado mental»
de quienes hablan en lo relativo a sus emociones, el grado de conocimiento de los
acontecimientos referidos, el acuerdo o el desacuerdo. En síntesis, se
relacionan directamente con el acto de habla \parencite{sanchezlopez2020}.

Dentro del amplio grupo de marcadores discursivos, \forma{evidentemente, por lo
visto, al parecer} ilustran un subgrupo particular: los \termino{evidenciales
citativos e inferenciales}, que \textcite[8--9]{cornillie2015}, por ejemplo,
estudian en el español peninsular. Consisten en una serie de marcadores del
discurso referidos a la percepción visual, pero que se emplean para indicar una
inferencia o también una información de oídas. El análisis de estos autores
sigue un criterio discursivo interaccional y muestra que existe una asimetría
entre \forma{evidentemente}, por un lado, y \forma{por lo visto} y \forma{al
parecer}, por otro: mientras que el primero se asocia al señalamiento de
pensamientos compartidos \exref{ex:14-vidente}, los otros dos expresan una
evidencia indirecta, \exref{ex:14-coca} y \exref{ex:14-nazis}.

\begin{exe}
  \ex\label{ex:14-vidente} Entonces, yo lo que hice fue ponerme en contacto con
      la madre primero\ldots\ que ya la\ldots\ había habla<(d)>o con ella en otras
      ocasiones, entonces, \textit{evidentemente} no le conté lo de la muerte,
      pero yo pensé que si ella llamaba a la policía judicial la harían más caso
      que si llamaba yo, de parte de una vidente de las muchísimas que han
      llama<(d)>o, porque han llama<(d)>o\ldots\ mogollón [\ldots].
  \ex\label{ex:14-coca} <H2> Sabes que ayer pillaron a un\ldots\ un director de
      un banco con\ldots\\
      <H3> Sí.\\
      <H2> \ldots\ kilo y medio de coca, ¿no?\\
      <H1> Con quince kilos.\\
      <H4> Quince, quince [kilos de cocaína]. Sí. 15, 15.\\
      <H2> ¡Ah, quince! Ah. Se me\ldots\ se me había corrido la coma. <todos> ah,
      15. [\ldots]\\
      <H2> Está bien, está bien quince, ¿no?\\
      <H4> Está muy bien. Era un distribuidor de primera línea, \textit{por lo
      visto}.
  \ex\label{ex:14-nazis} <H3> De todas maneras, lo que sí que\ldots\ ha ido a más
      en Europa eh\ldots\ han sido los brotes de xenofobia\ldots\\
      <H1> Aquí no han llega<(d)>o aún, ¿no?\\
      <H4> Sí, pero Madrid se iba a convertir el domingo, o se va a convertir, ya
      veremos, en una\ldots\ van a venir todos los nazis de Europa, \textit{por
      lo visto}, o neonazis y esto va a ser\ldots\ un festejo nazi. Por eso,
      \textit{al parecer}, queda prohibido\ldots
\end{exe}

\pendiente{Falta la fuente de los ejemplos orales (corpus citado por Cornillie y
Gras Manzano)}

Los datos de \exref{ex:14-coca} y \exref{ex:14-nazis} sugieren que \forma{ser
distribuidor de primera línea}, por un lado, y \forma{venir todos los nazis de
Europa} y \forma{quedar prohibido el festejo}, por otro, bajo la incidencia de
\forma{por lo visto} y \forma{al parecer}, no comunican una percepción visual de
los hechos, sino una deducción que realiza quien habla a partir del conocimiento
de acontecimientos que son la base del razonamiento.\pendiente{El original dice
«por lo visto y evidentemente»} En \exref{ex:14-coca}, que la carga de cocaína
decomisada a un director de un banco sea de quince kilos (y no de un kilo y
medio) es la causa para suponer que se trataba de un gran distribuidor; en
\exref{ex:14-nazis}, que los brotes de xenofobia en Europa se hayan incrementado
desencadena la suposición de que un encuentro de manifestantes previsto en
Madrid se vuelva un festejo neonazi y que, por lo tanto, las autoridades lo
prohíban.

Como podemos ir advirtiendo, el uso de los marcadores discursivos está modelado
por razones dialectales y sociolectales (o sea, reflejan un potencial de
variación de acuerdo con la habitualidad de su empleo en diferentes variedades
del español según su distribución geográfica o según el grado de formalidad de la
situación comunicativa). Por esto, dado también que el contenido interaccional de
esos marcadores es abstracto, muchas de estas unidades no son conmutables en el
uso concreto de comunidades diversas de hablantes del español. Por ejemplo, si
comparamos el empleo de \forma{igual, lo mismo} y \forma{justo} en el español
peninsular y en el argentino
\parencite{garcianegroni2013,marcovecchio2022,marcovecchio2023}, encontramos que
los usos discursivos peninsulares pueden resultar opacos o no «naturales» en
nuestra variedad, más allá de otros valores compartidos. Por ejemplo,
\forma{justo}, en ambas variedades, puede explicitar una reformulación de
equivalencia que establece el hablante para los dos segmentos en que se
interpone \exref{ex:14-justo1}:

\begin{exe}
  \ex\label{ex:14-justo1} Simplemente, las pone en su sitio, \textit{justo} (≡
      \textit{precisamente, exactamente}) en su verdadera dimensión.
\end{exe}

Sin embargo, otros datos muestran que no siempre nos resultan «transparentes»
\forma{igual, lo mismo} y \forma{justo}: en \exref{ex:14-igual} y
\exref{ex:14-lomismo}, \forma{igual} y \forma{lo mismo} se comportan como
marcadores epistémicos asemejables al empleo de \forma{capaz (que)} en el Río de
la Plata; y en \exref{ex:14-justo2}, \forma{justo} expresa, en una interacción
dialógica, una confirmación de lo dicho previamente por otro, es decir, la
adhesión con el interlocutor:

\begin{exe}
  \ex\label{ex:14-igual} Lleva paraguas, que \textit{igual} llueve. (= \textit{Capaz
      que} llueve).
  \ex\label{ex:14-lomismo} ---¿Dónde está el profesor?\\
      ---No sé. \textit{Lo mismo} está en su despacho. (= \textit{Capaz que} está
      en su despacho).
  \ex\label{ex:14-justo2} ---¿Estás insinuando que el director va a dimitir?\\
      ---\textit{Justo}. (= \textit{Precisamente}).
\end{exe}

El ejemplo \exref{ex:14-justo2} exhibe un uso reformulador de \forma{justo}. De
hecho, la \termino{reformulación}, proceso que permite que el hablante vuelva
sobre un discurso anterior para presentarlo desde una nueva perspectiva o para
reinterpretarlo, es un campo fértil para considerar una extensión de funciones
interaccionales de ítems como \forma{es decir, o sea, a saber, mejor dicho} o
\forma{en / con otras palabras}, que claramente son expresiones que explicitan
esta operación metalingüística de establecimiento de lazos de relativa
equivalencia entre dos segmentos. No obstante, \textcite{rodriguezramalle2014},
que estudia los reformuladores \forma{es decir} y \forma{o sea} frente a
\forma{esto es, a saber, mejor dicho, con otras palabras}, etc., concluye que la
distribución de unos y otros en el uso no es semejante. Según la autora,
\forma{es decir} y \forma{o sea} se distinguen de los demás por poder ser
utilizados por un hablante distinto del que realiza la primera intervención y,
asimismo, pueden introducir conclusiones, como se ve en \exref{ex:14-ingeniero}
y \exref{ex:14-analogia}, fragmentos de intercambios dialógicos en los que los
reformuladores introducen la voz de un hablante distinto con respecto a los
dichos previos, que motivan una reinterpretación que se sugiere como resultado
del establecimiento de inferencias generalizadoras y conclusivas:

\begin{exe}
  \ex\label{ex:14-ingeniero} A: Cuando yo era niño, mi madre escuchaba música
      occidental [\ldots]. Luego mi padre compró un fonógrafo eléctrico y el
      sonido me pareció totalmente distinto\ldots\ Aquella diferencia hizo que me
      decidiese a hacerme ingeniero electrónico.\\
      B: ¿\textit{O sea que} se hizo ingeniero electrónico a causa de la música?
  \ex\label{ex:14-analogia} A: Que los niños son analogistas y sobrepasan siempre
      el grado de regularidad del sistema de la lengua adulta lo sabe todo el
      mundo.\\
      B: \textit{Es decir que} es la facultad de analogía sobre todo lo que
      consideras como innato. (Casado, 1991, pp.~107--108).\pendiente{Casado
      (1991) no figura en la bibliografía}
\end{exe}

También \textcite{rodriguezramalle2014} analiza como marcadores evidenciales las
locuciones conjuntivas \forma{conque, de manera que} y \forma{así que}, en
relación con los reformuladores \forma{es decir} y \forma{o sea}, dado que también
pueden introducir un discurso que remite a las fuentes de información (sin
especificaciones sobre si se trata de evidencialidad citativa o inferencial) y
que motiva la emisión de los dichos que siguen. De ahí las interpretaciones
expresadas en \exref{ex:14-conque-b} y \exref{ex:14-alvaro-b}:

\begin{exe}
  \ex\label{ex:14-conque}
    \begin{xlist}
      \ex ¡\textit{Conque} por fin te vas y nos dejas! $\rightarrow$ ¿\textit{Así
          que / o sea que} por fin te vas y nos dejas?
      \ex\label{ex:14-conque-b} \sig{Me he enterado, me han dicho que por fin te
          vas y nos dejas}.
    \end{xlist}
  \ex\label{ex:14-alvaro}
    \begin{xlist}
      \ex \textit{Es decir que} tú eres el famoso Álvaro, del que todo el mundo
          habla. $\rightarrow$ \textit{De manera que} tú eres el famoso Álvaro,
          del que todo el mundo habla.
      \ex\label{ex:14-alvaro-b} \sig{Me han dicho o he deducido por lo que me han
          contado que tú eres el famoso Álvaro}.
    \end{xlist}
\end{exe}

\begin{destacado}
Como vemos, es amplio el repertorio de recursos disponibles para explicitar o
sugerir significados modales y evidenciales. Además, el análisis de estos
contenidos gramaticales requiere atender a la variación de uso, de acuerdo con
diferencias geográficas, situacionales y de registro, así como a la extensión de
funciones y significados de ítems que, por procesos de pragmaticalización, suman
comportamientos discursivos.
\end{destacado}

\begin{ejercicio}\label{ej:14-evidenciales}
Indique qué tipo de evidencialidad y / o de modalidad codifican las formas
destacadas en las siguientes oraciones, según se muestra en el ejemplo
resuelto~1.

\begin{exe}
  \ex \textit{En principio}, no \textit{sería} nada grave [\ldots].
  \ex \textit{Debe ser} un fenómeno multicausal. La lógica \textit{indicaría} en
      principio que, al crecer la pobreza, \textit{deberían aumentar} este tipo
      de modalidades. \textit{Al parecer} lo único que está creciendo es el robo
      tipo motochorro.
  \ex ¿Vos te \textit{llamabas}\ldots?
  \ex ---¡\textit{Qué va a ser} hermoso! Es horrible.
  \ex El señor \textit{sabe venir} por aquí a la mañana.
  \ex Otro abogado explicó: «\textit{Puede que} los minoristas no \textit{acepten}
      ahora, pero \textit{si se concretara} el canje, algunos \textit{podrían
      aceptarlo}». (\textit{La Nación}, 16/04/2013).
\end{exe}

\begin{resuelto}
1. Tanto la locución \forma{en principio} como la flexión en condicional de
\forma{ser} codifican evidencialidad citativa y / o inferencial (la
interpretación es ambigua).
\end{resuelto}
\end{ejercicio}

\begin{destacado}
\textbf{Para tener en cuenta.} Estas diferentes actitudes subjetivas se combinan
con la \termino{polaridad} oracional: afirmativa o negativa. La polaridad
afirmativa no necesita una marca explícita; en cambio, la negativa, que siempre
expresa un distanciamiento, un cuestionamiento por parte de quien habla con
respecto a lo que dice, requiere marcas específicas (\forma{no, ni, nunca, nada,
nadie\ldots}), los llamados \termino{términos de polaridad negativa}.
\end{destacado}

\section{Modalidad del propio acto de enunciación}\label{sec:14-enunciacion}

Al estudiar la estructura oracional, podemos advertir que, típicamente, hay
constituyentes periféricos con respecto a la predicación que expresan contenidos
relativos a la situación de habla, lo que concierne tanto a los interlocutores
como a la misma estructuración del discurso. En el primer caso, se trata de
expresiones como \forma{francamente, sinceramente, en confianza, para ser
honesto\ldots} En el segundo, nos referimos a todos los elementos que establecen
indicios de la presencia del hablante como la persona que organiza los
contenidos en la secuencia discursiva: ordenándolos y jerarquizándolos
(\forma{primero / en primer lugar, segundo / en segundo lugar, por / para
empezar, en conclusión / para concluir, en síntesis\ldots}) o estableciendo
reformulaciones ---de equivalencia o no--- entre segmentos (\forma{en otras
palabras, dicho de otro modo, por poner un ejemplo, para ilustrar\ldots}).

\begin{ejercicio}\label{ej:14-tipos}
Reconozca las propiedades formales y los ítems léxicos que permiten incluir las
oraciones de abajo dentro de los tres tipos de modalidad diferenciados, según se
muestra en el ejemplo resuelto~1.

\begin{exe}
  \ex Modalidades intelectuales o epistémicas:
    \begin{xlist}
      \ex Por supuesto, guardé las llaves.
      \ex Por ahí, dejé las llaves sobre el escritorio.
    \end{xlist}
  \ex Modalidades volitivas (deónticas): exhortativa, desiderativa:
    \begin{xlist}
      \ex Que pase el que sigue.
      \ex Chicos, (les pido que) guarden los útiles en sus mochilas.
      \ex Chicos, pueden / deben guardar los útiles en sus mochilas.
      \ex Chicos, ¡a ver si guardan los útiles en sus mochilas (de una vez por
          todas)!
      \ex ¡Ojalá fuera rica! Si fuera rica, viajaría por el mundo.
    \end{xlist}
  \ex Modalidades afectivas: exclamativas:
    \begin{xlist}
      \ex ¡Qué lindo (que) es soñar!
      \ex Siempre dice qué (/ lo) lindo (que) es soñar.
    \end{xlist}
\end{exe}

\begin{resuelto}
1. a) Modo indicativo, con refuerzo de la aserción \forma{por supuesto}. b) Modo
indicativo, con marcador hipotético inferencial \forma{por ahí}.
\end{resuelto}
\end{ejercicio}

\begin{ejercicio}\label{ej:14-periferia}
Explique si los elementos destacados dependen del verbo de la predicación
oracional o se orientan hacia el acto de enunciación. Justifique su respuesta.

\begin{textolectura}
El hecho habría ocurrido a escasos kms. de San Guillermo. [\ldots]
\textit{Extraoficialmente}, se presume que el robo rondaría en una cifra cercana
a 300 mil pesos. [\ldots] \textit{Según fuentes consultadas por Radio Belgrano},
habría sido abordado por dos sujetos [\ldots]. Los delincuentes se habrían
llevado la recaudación.
\end{textolectura}
\end{ejercicio}

\begin{recuadro}[Reflexión final]
Sin lugar a duda, las sutilezas de la expresión de la modalidad requieren
continuar ahondando en los posibles significados que se explicitan o que apenas
se sugieren. No obstante, hasta aquí hemos confrontado datos provenientes del
uso real, que ilustran los recursos prosódicos, morfosintácticos, léxicos y de
la configuración oracional que inciden en la manifestación de la modalidad y de
la evidencialidad, como «estrategias» que permiten la reinterpretación de
ciertos morfemas de la flexión verbal, de algunos adverbios y conjunciones del
español en términos de empleos citativos o inferenciales. En este recorrido,
hemos consolidado temas tratados en capítulos anteriores (las categorías
morfológicas, las perífrasis verbales, la estructura oracional con funciones
nucleares ---argumentales o no, pero siempre dentro del ámbito del SV--- y
periféricas, las alternancias cuasirreflejas mediopasivas e impersonales) y los
hemos reconsiderado integrados en la combinatoria oracional. Asimismo, hemos
destacado la necesidad de atender al potencial de variación que existe en el
español: los empleos reales de la lengua muestran las particularidades de uso de
los recursos gramaticales asociados a la modalidad y la evidencialidad.
\end{recuadro}

% --- Secciones finales -------------------------------------------------
\lecturasrecomendadas

\solucionario

\begin{solucion}{\ref{ej:14-canciones}}
2. a. Modalidad epistémica aseverativa con verbos flexionados en presente del
modo indicativo. b. Modalidad epistémica dubitativa: la perífrasis verbal con
\forma{poder} encadenada a un infinitivo compuesto indica un hecho contrafáctico
(véase §\ref{sec:14-modo}) o lo que no sucedió (\sig{no fui el pianista del gueto
de Varsovia}). c. Modalidad epistémica aseverativa, con verbos en presente y
pretérito imperfecto del modo indicativo.

3. a. Modalidad deóntica exhortativa articulada en la oración principal sobre el
verbo \forma{aprender} en futuro, con valor de verdad general; la subordinada
introducida por el relativo \forma{quien} presenta el verbo en modo subjuntivo,
\forma{sepa}, que adopta valor de \sig{inespecificidad}, porque alude a
cualquier persona que no esté en conocimiento del contenido referido. b.
Modalidad epistémica aseverativa negativa con sentido general, con el verbo en
presente de indicativo. c. Modalidad afectiva, con intensificador resaltado por
la repetición (\forma{tantos, tan}). d. Modalidad epistémica declarativa con
verbos en indicativo (\forma{eran, decía}); \forma{Cautela} manifiesta modalidad
deóntica exhortativa con valor de advertencia. e. Modalidad afectiva, con el
intensificador \forma{cuánta}. f. Modalidad afectiva, con uso explícito y
enfático del pronombre de 2.ª persona y empleo del intensificador \forma{tan}
repetido.\pendiente{Las letras de los apartados (a, b, c, e, f, g en el
enunciado original; a--f en el solucionario) no coinciden}
\end{solucion}

\begin{solucion}{\ref{ej:14-modos}}
2. La oposición de modo indicativo (\forma{protege}) y modo subjuntivo
(\forma{proteja}) opone modalidad epistémica aseverativa a modalidad deóntica
desiderativa. 3. La oposición de futuro (\forma{escribiremos}) y condicional
(\forma{escribiríamos}), ambos dentro del modo indicativo, opone modalidad
epistémica aseverativa frente a modalidad epistémica dubitativa. 4. La oposición
de presente de subjuntivo (\forma{venga}) e imperfecto de subjuntivo
(\forma{viniera}) opone modalidad deóntica de deseo posible frente a modalidad
deóntica de deseo imposible (en el presente: \forma{¡Ojalá viniera hoy!}) o
improbable (en el futuro: \forma{¡Ojalá viniera uno de estos días!}). 5. La
oposición de imperfecto (\forma{viniera}) y pluscuamperfecto de subjuntivo
(\forma{hubiera venido}) contrapone modalidad deóntica de deseo imposible (en el
presente) o improbable (en el futuro) a modalidad de deseo irreal
contrafáctico. 6. La oposición de presente de indicativo (\forma{juega}) y
presente de subjuntivo (\forma{juegue}) contrapone modalidad epistémica de
incerteza débil frente a modalidad epistémica de incerteza fuerte.
\end{solucion}

\begin{solucion}{\ref{ej:14-formas}}
2. Las formas verbales \forma{ha ido} (perfecto compuesto) y \forma{habrá ido}
(futuro perfecto) oponen una pregunta que se formula desde el
\sig{desconocimiento} a otra que ofrece una inferencia. 3. Las formas verbales
\forma{dejaste} (perfecto simple) y \forma{habrás dejado} (futuro perfecto)
oponen \sig{conocimiento efectivo} a \sig{conocimiento inferido}. 4. Las formas
verbales \forma{llamará} (futuro) y \forma{llamaría} (condicional) oponen
\sig{conocimiento efectivo} a \sig{evidencialidad citativa} o bien
\sig{conocimiento inferido}. 5. Las formas verbales \forma{hubiera sabido / no
habría venido} (pluscuamperfecto de subjuntivo y condicional compuesto) y
\forma{sabía / venía} (ambas en imperfecto de indicativo) oponen registro formal
a registro coloquial informal; en ambos pares se mantiene la interpretación
contrafáctica. 6. Las formas \forma{quiero / resulta} (presente de indicativo) y
\forma{quisiera / estaría resultando} (imperfecto de subjuntivo y condicional)
oponen grados de cortesía: \forma{quisiera} y \forma{estaría resultando}, al
indicar \sig{irrealidad}, atenúan la fuerza deóntica del pedido o el comentario
negativo sobre la boquilla.
\end{solucion}

\begin{solucion}{\ref{ej:14-parafrasis}}
2. Eso \textit{no puede / no debe ser discutido} / \textit{no se puede
discutir}. 4. Esta prenda \textit{no puede / no debe ser lavada} / \textit{no se
puede / debe lavar} en máquina. 5. No \textit{se puede vivir} del aire, sin
trabajar. 6. En esta sala, \textit{no se puede / no se debe
gritar}.\pendiente{El original numera estas respuestas 2 a 5 y no resuelve
la oración 3 (\textit{Las emociones no se medican})}
\end{solucion}

\begin{solucion}{\ref{ej:14-gramaticalizacion}}
En las lenguas, se reconoce una especie de «complementariedad» entre recursos
gramaticales, aquí morfológicos (la flexión en modo subjuntivo de
\forma{hiciera}) o morfosintácticos (la perífrasis verbal \forma{tenía que
hacer}). El verbo \forma{indicar} introduce un COD que expresa una recomendación
mediante la modalidad deóntica; por lo tanto, el verbo de la subordinada o bien
manifiesta esa modalidad con el verbo flexionado en subjuntivo, o bien lo hace
con la perífrasis verbal de modalidad deóntica.
\end{solucion}

\begin{solucion}{\ref{ej:14-subjuntivo}}
2. En combinación con \forma{cierra} (indicativo), \forma{sentir} es un verbo de
\sig{percepción sensorial} que manifiesta modalidad epistémica aseverativa; con
\forma{cierre} (subjuntivo), \forma{sentir} se convierte en verbo psíquico y se
interpreta como \sig{lamentar}, lo que manifiesta modalidad afectiva. 3. Las
formas verbales \forma{asistirá} (indicativo) y \forma{asista} (subjuntivo)
contraponen modalidad epistémica aseverativa de certeza débil a modalidad
epistémica dubitativa. 4. Las formas verbales \forma{existe} (indicativo) y
\forma{exista} (subjuntivo) contraponen modalidad epistémica aseverativa con
posicionamiento fuerte del hablante (\sig{el juez no cree, pero el hablante
afirma que existe delito}) a un posicionamiento neutro de quien habla (\sig{el
juez no cree en la existencia del delito, y el hablante no asume tampoco un
compromiso epistémico}). 5. Las formas verbales \forma{tengo} (indicativo) y
\forma{tenga} (subjuntivo) contraponen modalidad epistémica aseverativa acerca de
la existencia de la historia mencionada al desconocimiento de su posible
existencia o la falta de especificidad al mencionarla. 6.\pendiente{El original
no resuelve la oración 6 (\textit{siempre que nos conduce / conduzca})} 7.
Ambas prótasis condicionales en oposición expresan un vínculo hipotético entre
\forma{la presentación del pasaporte} y \forma{la posibilidad de salir del país};
no obstante, con \forma{si} más predicado en modo indicativo, el posicionamiento
del hablante es neutro, mientras que con \forma{como} más predicado en modo
subjuntivo se sugiere un posicionamiento de advertencia o amenaza.
\end{solucion}

\begin{solucion}{\ref{ej:14-evidenciales}}
2. Las perífrasis con \forma{deber}, la flexión en condicional de
\forma{indicar} y \forma{deber} y el marcador discursivo \forma{al parecer}
expresan una intersección entre modalidad epistémica de probabilidad y
evidencialidad inferencial. 3. La forma en pretérito imperfecto del modo
indicativo \forma{llamabas} sugiere evidencialidad citativa. 4. La forma
perifrástica \forma{ir a} + infinitivo sugiere evidencialidad citativa, a la vez
que la no adhesión a los dichos previos. 5. La perífrasis verbal con
\forma{saber} expresa aspecto habitual y también involucra el matiz modal de
posibilidad. 6. \forma{Puede que} introduce una predicación modalizada
epistémicamente, con indicación de \sig{posibilidad}, que exige la forma
flexionada de \forma{aceptar} en modo subjuntivo. El período condicional que
sigue, con la prótasis en pretérito imperfecto del modo subjuntivo y la apódosis
en condicional, indica \sig{irrealidad}, específicamente, algo improbable en el
futuro.
\end{solucion}

\begin{solucion}{\ref{ej:14-tipos}}
2. Modalidades volitivas (deónticas): exhortativa, desiderativa. a) Modo
subjuntivo exigido por el índice de exhortación o deseo \forma{que}. b) Modo
subjuntivo exigido por el verbo \forma{pedir}, que explicita la fuerza
ilocutiva; inclusión del vocativo para nombrar a las personas a las que va
dirigida la exhortación. c) Perífrasis verbales que expresan modalidad deóntica
y empleo del vocativo para nombrar a las personas a las que va dirigida la
exhortación. d) Uso de estructuras que pueden expresar modalidad deóntica
(\forma{a ver si\ldots}), más un marcador del énfasis del pedido (\forma{de una
vez por todas}); empleo del vocativo para nombrar a las personas a las que va
dirigida la exhortación. e) Modo subjuntivo exigido por el índice desiderativo.

3. Modalidades afectivas: exclamativas. a) Uso de pronombres enfáticos que se
corresponden con la curva entonativa exclamativa. b) Uso de pronombres enfáticos
que indican la modalidad exclamativa.
\end{solucion}

\begin{solucion}{\ref{ej:14-periferia}}
\forma{Extraoficialmente} y \forma{según fuentes consultadas por Radio Belgrano}
no dependen de los verbos de las oraciones sobre las que inciden porque no
admiten focalización ni establecen alguna función sintáctica dentro de la
predicación. Desempeñan funciones periféricas y expresan que las declaraciones
hechas provienen de información obtenida de segunda mano.
\end{solucion}

% =====================================================================
%  capitulos/cap15.tex — Capítulo 15. Lengua, dialecto y descripción
%  gramatical: el caso del español rioplatense
%  Autores: Carlos Muñoz Pérez y Andrés Saab
%  Fuente: Capitulo_15_normalizado.docx
% =====================================================================
\capitulo{Lengua, dialecto y descripción gramatical: el caso del español rioplatense}{Carlos Muñoz Pérez y Andrés Saab}

\section{Introducción}

Este capítulo ofrece una introducción al \termino{español rioplatense}, la
variedad del español hablada en la región del Río de la Plata (Argentina y
Uruguay). El recorrido comienza con una discusión sobre las nociones de
\textit{lengua} y \textit{dialecto}; como veremos, si bien ambos términos son de
uso habitual, es difícil distinguirlos desde un punto de vista estrictamente
gramatical. Luego, se presentan los factores históricos y geográficos que
moldearon la variedad rioplatense y se describen sus rasgos más distintivos en
los planos léxico, fonético-fonológico, morfológico y sintáctico. Entre los
fenómenos analizados y discutidos se incluyen el lunfardo, el yeísmo rehilado,
el voseo y el doblado de clíticos, entre otros.

\section{Algunas precisiones terminológicas sobre lengua y dialecto}\label{sec:15-terminologia}

El español rioplatense es una de las variedades del español habladas en América
del Sur. Antes de describir sus rasgos particulares, quisiéramos problematizar
dos términos que utilizamos a lo largo de este capítulo, \textit{lengua} y
\textit{dialecto}, ya que, aunque parecen nociones sencillas, veremos que son más
complejas de lo que suele suponerse.

\subsection{Lengua-I y lengua-E}\label{sec:15-lenguai}

El término \textit{lengua} puede significar al menos dos cosas distintas. En un
primer sentido, una lengua, como se planteó en el \capref{1}
(§\ref{sec:01-generativa}), es un tipo de conocimiento interno: algo que cada ser
humano desarrolla durante sus primeros años de vida gracias a la combinación de
(i) una capacidad biológica innata para el lenguaje y (ii) los estímulos
lingüísticos que recibe de su entorno. Este conocimiento es mental e individual:
está en la cabeza de cada hablante. En la tradición generativa, siguiendo una
distinción propuesta por \textcite{chomsky1986}, este objeto se denomina
\termino{lengua-I}. En un segundo sentido, \textit{lengua} puede referirse a un
objeto externo: el conjunto de expresiones y enunciados que una comunidad
produce y que pueden ser observados, recopilados y analizados. Cuando hablamos de
«el español» o «el guaraní» como entidades que existen más allá de cualquier
hablante individual, estamos usando esta segunda acepción. A este objeto lo
llamamos \termino{lengua-E}. Esta segunda acepción es, a grandes rasgos, la
concepción de lengua que ha predominado en la mayor parte de la tradición
lingüística, desde la gramática tradicional hasta el estructuralismo.

¿Cuál es la relación entre ambas nociones? Podemos pensarlo así: la lengua-I es
el mecanismo que cada hablante tiene en su mente y que le permite producir y
comprender oraciones; la lengua-E es el producto observable de ese mecanismo, es
decir, el conjunto de enunciados que efectivamente se producen. Ahora bien, esta
relación no es perfecta: en el uso de la lengua intervienen factores externos
(interrupciones y errores por cansancio, distracción, etc.) que introducen
«ruido» en los datos. Por eso, parte del trabajo del lingüista consiste en
determinar qué fragmentos de la lengua-E reflejan genuinamente la lengua-I de los
hablantes y cuáles son producto de esas interferencias.

Para entender mejor la distinción, pensemos en un caso concreto. Imaginemos a una
niña nacida en la provincia de Buenos Aires en la década de 1940. Su padre es un
inmigrante italiano que habla un dialecto del norte de Italia; su madre es una
inmigrante gallega que habla español y gallego. Ambos padres quieren que su hija
aprenda español, la lengua oficial del país, que será imprescindible para su
escolarización. Además, tenemos que considerar que las lenguas de los inmigrantes
están fuertemente estigmatizadas. Por tanto, ¿qué datos lingüísticos recibe esta
niña? Podemos suponer que se trata de un conjunto bastante caótico: fragmentos de
español peninsular y de gallego (de su madre), fragmentos de un dialecto de
Italia (de su padre) y español rioplatense (de vecinos, familiares y, más tarde,
de la escuela). Esos estímulos lingüísticos no provienen de una única lengua-I,
ni constituyen un conjunto homogéneo, sino que son ejemplos de distintas
lenguas-E. Y, sin embargo, en un plazo relativamente breve, la niña comienza a
producir oraciones que responden a un sistema gramatical completo: su propia
lengua-I. Esta historia se repitió en innumerables hogares de Buenos Aires
durante las grandes olas migratorias: cada niño o niña desarrolló una lengua-I
estrictamente individual, pero, como los factores externos eran en gran medida
similares, sus producciones lingüísticas comparten numerosas propiedades
fonéticas, fonológicas, morfológicas y sintácticas. Ese conjunto de propiedades
compartidas es, justamente, lo que describimos cuando hablamos del español
rioplatense.

\begin{destacado}
\textbf{Para tener en cuenta.} Idealmente, una descripción completa de una lengua
o de un dialecto debería capturar ese conocimiento interno que cada persona
adquiere, es decir, debería ser una teoría acerca de la lengua-I. No obstante,
dada la naturaleza introductoria de este capítulo, lo que ofrecemos es una
recopilación de ciertos fragmentos de lengua-E. Para decirlo de otro modo,
presentamos un subconjunto de los datos lingüísticos que caracterizan de manera
diferencial al español hablado hoy en día en el área rioplatense.
\end{destacado}

\subsection{La escurridiza noción de dialecto}\label{sec:15-dialecto}

Pasemos entonces a una pregunta central: ¿existe alguna relación relevante entre
las nociones de \textit{lengua} y \textit{dialecto}? Lo primero que debemos hacer
es descartar una idea muy extendida, pero incorrecta: que un dialecto es una
«desviación» o una versión «inferior» de una lengua. En el uso popular, el
concepto de \textit{dialecto} suele referirse a formas no estándares de hablar, a
veces consideradas rústicas o inferiores, o a lenguas minoritarias sin estatus
oficial. Así, por ejemplo, en México es habitual oír que las comunidades
indígenas hablan «dialectos», en contraste con el español, que sí sería una
«lengua». Así presentada, esta distinción no tiene sustento lingüístico y la
rechazamos explícitamente.

Un criterio algo más sofisticado, pero igualmente problemático, es el de la
\termino{intercomprensión}: según este criterio, dos variedades son dialectos de
una misma lengua si sus hablantes pueden entenderse entre sí y, en cambio, son
lenguas distintas si no pueden hacerlo. El problema es que este criterio no
funciona de manera consistente. Hay lenguas reconocidas como tales que son mucho
más parecidas entre sí que ciertos dialectos de una misma lengua. El español y el
portugués, por ejemplo, son lenguas oficiales de estados distintos y se
consideran lenguas diferentes, pero un hablante del español rioplatense y un
hablante del portugués brasileño pueden entenderse con relativa facilidad en
muchas situaciones cotidianas. A su vez, hay diferencias específicas entre, por
ejemplo, el español caribeño hablado en República Dominicana y el español
rioplatense que pueden dificultar la mutua intercomprensión. Este criterio, por
tanto, no nos ofrece una frontera clara y sistemática.

¿Cómo definir entonces \textit{dialecto} de manera rigurosa?
\textcite{silvacorvalan2001} propone una solución práctica: si tomamos una
lengua-E como el español y la delimitamos a partir de factores históricos,
políticos, geográficos y culturales, podemos pensar que se organiza internamente
como una pirámide. En la cúspide de esta \termino{pirámide dialectal} se ubica la
lengua histórica o común (es decir, «el español» a nivel general); en un nivel
intermedio, los grandes dialectos regionales (es decir, el español peninsular y
el español de América); y en la base, los dialectos locales concretos (aquí es
donde encontramos el rioplatense, el español andaluz, etc.). Un punto importante
de esta concepción es que supone que la lengua general no es hablada por ninguna
comunidad en particular: lo que cada persona habla es siempre un dialecto
particular, situado en la base de la pirámide. En otras palabras, hablar una
lengua equivale siempre a hablar un dialecto de esa lengua. La
figura~\ref{fig:15-piramide} ilustra esta organización.

\begin{figure}[htbp]
  \centering
  \begin{tikzpicture}[x=1mm,y=1mm,font=\small]
    \draw[thick] (0,0) -- (90,0) -- (45,54) -- cycle;
    \draw (15,18) -- (75,18);
    \draw (30,36) -- (60,36);
    \node[align=center] at (45,42) {Lengua\\histórica};
    \node[align=center] at (45,26) {Dialectos regionales\\\scriptsize(español peninsular, español de América)};
    \node[align=center] at (45,8) {Dialectos locales\\\scriptsize(rioplatense, andaluz, caribeño\ldots)};
  \end{tikzpicture}
  \caption{Pirámide dialectal \parencite[adaptado de][14]{silvacorvalan2001}}\label{fig:15-piramide}
\end{figure}

\pendiente{Diagrama 1 sin imagen en el original; reconstruido a partir del texto}

Ahora bien, en un esquema como el de Silva-Corvalán, lo que ocupa la cúspide de
la pirámide, es decir, la «lengua histórica o común», no se define por criterios
puramente lingüísticos. Como ya se observó, son factores históricos,
geográficos, culturales y, de manera decisiva, políticos los que determinan qué
cuenta como «lengua» y qué cuenta como «dialecto». Es decir, lo que nos lleva a
clasificar el español rioplatense, el español caribeño o el español andaluz como
«dialectos del español» y no como lenguas independientes no es un criterio
gramatical, sino el hecho de que comparten una tradición histórica, geográfica y
cultural, un nombre y, sobre todo, un marco político-institucional que los
agrupa bajo una misma etiqueta.

\begin{ejercicio}\label{ej:15-weinreich}
Al lingüista Max Weinreich se le atribuye el aforismo «una lengua es un dialecto
con un ejército y una marina». Discuta el sentido de esta afirmación a la luz de
la discusión precedente. Específicamente: ¿a qué tipo de criterio para
distinguir \textit{lengua} de \textit{dialecto} refiere el aforismo?
\end{ejercicio}

\subsection{Una posible diferenciación: microparámetros frente a macroparámetros}\label{sec:15-parametros}

Cabe preguntarse, entonces, si existe alguna manera «lingüística» de establecer
una distinción entre lenguas y dialectos. Como veremos, la respuesta es «sí»,
pero bajo ciertos supuestos. Primero, debemos suponer que existen distintas
maneras en las que las lenguas pueden diferir entre sí. La idea es que existen
«opciones de configuración» del sistema lingüístico denominadas
\termino{parámetros} (véase el \capref{1}, §\ref{sec:01-generativa}), que se
fijan durante el proceso de adquisición de la lengua. Algunas de las propiedades
que determinan estos parámetros afectan a todo el sistema de la lengua y crean
una especie de «efecto dominó» con respecto al funcionamiento de varias áreas de
la gramática. Estos parámetros «de gran alcance» se denominan por lo general
\termino{macroparámetros}.\footnote{\textcite{roberts2019} propone una tipología
más fina, que distingue al menos cuatro niveles de variación paramétrica (macro-,
meso-, micro- y nanoparámetros), ordenados según el grado de especificidad del
cambio observado. Por simplicidad, en este capítulo trabajamos solo con el
contraste entre macro- y microparámetros.} Por ejemplo, muchas veces se asume
que existe un macroparámetro encargado del orden de las palabras dentro de los
sintagmas: en lenguas como el español, el núcleo aparece al inicio del sintagma;
en lenguas como el japonés, el núcleo aparece al final.

\begin{exe}
  \ex\label{ex:15-espanol} Núcleo inicial (español):
    \begin{xlist}
      \ex {}[\etq{SV} comprar un regalo]
      \ex {}[\etq{SP} en la ciudad]
    \end{xlist}
  \ex\label{ex:15-japones} Núcleo final (japonés):
    \begin{xlist}
      \ex
        \gll {}[\etq{SV} purezento-o kau]\\
             {} regalo-\textsc{acus} comprar\\
        \glt \sig{comprar un regalo}
      \ex
        \gll {}[\etq{SP} machi-de]\\
             {} ciudad-en\\
        \glt \sig{en la ciudad}
    \end{xlist}
\end{exe}

Por otra parte, la bibliografía también ha identificado parámetros de variación
que describen diferencias más «pequeñas» o de alcance más restringido. Estos se
conocen con el nombre de \termino{microparámetros}. Los microparámetros describen
cambios en las propiedades de los elementos funcionales de la lengua (p.~ej., en
pronombres, en morfemas flexivos, en artículos), básicamente en unidades de
clase cerrada dentro del sistema lingüístico. Así, por ejemplo, podemos concebir
el llamado \termino{leísmo} como un caso de variación microparamétrica en el
sistema de pronombres clíticos de ciertos dialectos del español. El leísmo
consiste en el uso del pronombre \forma{le} (originalmente dativo) en contextos
en los que otros dialectos emplean los pronombres acusativos \forma{la} o
\forma{lo}. Así, para referirse a una misma situación (por ejemplo, haber visto a
una persona el día anterior), un hablante de un dialecto no leísta y uno de un
dialecto leísta producirían oraciones distintas.

\begin{exe}
  \ex\label{ex:15-leismo}
    \begin{xlist}
      \ex La hemos visto ayer. \hfill \textit{variante no leísta}
      \ex Le hemos visto ayer. \hfill \textit{variante leísta}
    \end{xlist}
\end{exe}

A partir de la distinción entre macroparámetros y microparámetros es posible
ofrecer una definición explícita de lo que es una lengua y de lo que es un
dialecto. Así, si dos sistemas lingüísticos (es decir, dos lenguas-I en un
sentido técnico) difieren en su configuración macroparamétrica, podría decirse
que son lenguas diferentes (p.~ej., el español y el japonés son lenguas
diferentes porque, como muestran los ejemplos de \exref{ex:15-espanol} y
\exref{ex:15-japones}, difieren en al menos un macroparámetro). En cambio, si dos
sistemas lingüísticos solo difieren en términos de sus microparámetros, porque
sus macroparámetros están configurados de manera análoga, puede decirse que son
dialectos de la misma lengua (p.~ej., las variantes leístas y no leístas de
\exref{ex:15-leismo} representarían dialectos de la misma
lengua).\footnote{Esta es, en esencia, la propuesta de \textcite{baker2008}.}

Esta manera de entender qué es una lengua y qué es un dialecto descansa sobre la
distinción entre macroparámetro y microparámetro. Sin embargo, no todos los
lingüistas acuerdan en que ambas nociones son necesarias. Algunos autores
defienden la idea de que los macroparámetros no existen como tales y que las
grandes diferencias que se observan a nivel interlingüístico se deben
simplemente a la acumulación de diferencias microparamétricas. La idea suele
sintetizarse en la siguiente tesis de \textcite{kayne2005}:

\begin{exe}
  \ex\label{ex:15-kayne} Todo parámetro es un microparámetro.
\end{exe}

Bajo este supuesto, no existen distintos tipos de parámetro que justifiquen una
distinción cualitativa entre lengua y dialecto; es decir, dado que toda unidad
de cambio lingüístico es de la misma naturaleza, los términos \textit{lengua} y
\textit{dialecto} refieren al mismo tipo de variación. La única diferencia que
puede establecerse entre ambas nociones es, a lo sumo, de grado: el español y el
japonés serían lenguas distintas porque acumulan entre sí muchas pequeñas
diferencias; el español rioplatense y el español caribeño serían dialectos del
español porque acumulan entre sí muchas menos diferencias.

A pesar de nublar la distinción entre lengua y dialecto, esta segunda postura
tiene una consecuencia práctica interesante: convierte a la dialectología en una
herramienta de investigación poderosa para la teoría lingüística. Si la
variación entre dialectos es del mismo tipo que la variación entre lenguas,
entonces comparar dialectos cercanos nos permite aislar diferencias mínimas y
estudiar sus efectos en detalle, como en una suerte de experimento controlado.

\section{Cuestiones históricas y geográficas}\label{sec:15-historia}

Para contextualizar las propiedades del español rioplatense, es necesario conocer
algunos hechos históricos que incidieron directamente en la conformación de esta
variedad. No se trata de hacer un recorrido exhaustivo por la historia de la
región, sino de identificar los momentos clave en los que ciertos eventos
demográficos y políticos dejaron una huella en la lengua.

La fundación definitiva de Buenos Aires tuvo lugar en 1580, con la expedición de
Juan de Garay. De los sesenta y seis participantes de esa expedición, solo diez
provenían directamente de España; el resto eran mestizos procedentes de Asunción
\parencite{fontanella1987}. Esto significa que, desde sus orígenes, Buenos Aires
fue una ciudad donde convivían el español y el guaraní: una parte considerable de
sus primeros pobladores, en particular las mujeres, eran hablantes bilingües de
ambas lenguas. Esta situación inicial de contacto lingüístico se fue modificando
a lo largo de los siglos \textsc{xvii} y \textsc{xviii}, a medida que la llegada
de nuevos pobladores desde España impuso progresivamente el uso del español por
sobre el guaraní. Hacia 1776, con el establecimiento definitivo del Virreinato
del Río de la Plata, Buenos Aires se consolidó como centro administrativo y
cultural de la región, lo que reforzó el papel institucional del español como
lengua dominante.

Sin embargo, el factor histórico que más profundamente moldeó el perfil
lingüístico del español rioplatense fue la inmigración masiva de fines del siglo
\textsc{xix} y principios del \textsc{xx}. Como destaca
\textcite{ditullio2010a}, las políticas inmigratorias del Estado argentino,
consolidado hacia 1880, transformaron de manera drástica el panorama cultural y
lingüístico de la región. Para dimensionar la magnitud de este cambio, basta
considerar algunos datos: entre 1869 y 1914, la población del país se
cuadruplicó, pasando de aproximadamente dos millones a casi ocho millones de
habitantes. Hacia la época del Centenario (1910), uno de cada tres habitantes era
extranjero. De los cuatro millones de inmigrantes registrados en 1914, la mitad
eran italianos provenientes tanto del norte como del sur de Italia. Estos
hablantes, que traían consigo una enorme diversidad de dialectos italianos,
ejercieron una influencia considerable sobre el español de la zona,
especialmente en el plano de la entonación, la morfología y el
léxico.\footnote{Debe tenerse en cuenta, sin embargo, que, según
\textcite[94]{fontanella1987}, varias propiedades fonológicas y fonéticas de la
variedad, como el seseo, el yeísmo y la aspiración y elisión de /s/, ya habían
comenzado a consolidarse en el siglo \textsc{xviii}, antes de las grandes olas
migratorias (véase §\ref{sec:15-fonetica}).}

\subsection{El español rioplatense}\label{sec:15-rioplatense}

Ahora bien, es importante no confundir el español rioplatense con el español
argentino en general. La Argentina es un país extenso y, desde el punto de vista
dialectológico, presenta una diversidad interna que a menudo pasa inadvertida. En
un trabajo clásico, \textcite{vidal1964} distingue al menos cinco regiones
dialectales dentro del territorio argentino, muchas de las cuales se extienden
más allá de las fronteras nacionales.

\begin{exe}
  \ex\label{ex:15-regiones} Regiones dialectales de la Argentina
      \parencite{vidal1964}:
    \begin{xlist}
      \ex La \textit{región rioplatense}, que es el objeto de estudio de este
          capítulo.
      \ex La \textit{región guaranítica}, que comprende Corrientes, Misiones, el
          este de Formosa y del Chaco, el extremo noreste de Santa Fe y, de manera
          menos clara, zonas de Entre Ríos.
      \ex La \textit{región del noroeste}, que abarca Jujuy, Salta, Tucumán,
          Santiago del Estero, Catamarca, La Rioja, el norte de San Juan y San Luis
          y el noroeste de Córdoba.
      \ex La \textit{región de Cuyo}, que incluye Mendoza, San Juan y el norte de
          Neuquén.
      \ex La \textit{región central}, constituida principalmente por Córdoba y San
          Luis.
    \end{xlist}
\end{exe}

\figurapendiente{Regiones dialectales de la Argentina (adaptado de Vidal de Battini, 1964). Fuente: Wikimedia Commons, licencia CC BY-SA 3.0}{fig:15-mapa}

\begin{destacado}
\textbf{Para tener en cuenta.} La clasificación anterior permite apreciar que lo
que solemos llamar «español argentino» es en realidad un mosaico de variedades
con rasgos propios. Un hablante de Salta y un hablante de Buenos Aires comparten
la nacionalidad, pero sus maneras de hablar presentan diferencias notables en la
pronunciación, el vocabulario y ciertas construcciones gramaticales. El español
rioplatense es solo una de esas variedades y no debe tomarse como representativa
de todo el país.
\end{destacado}

¿Qué territorio abarca, entonces, el español rioplatense? Según
\textcite{vidal1964}, esta región incluye la provincia de Buenos Aires, la ciudad
de Buenos Aires, las provincias de Santa Fe, Entre Ríos y La Pampa, las zonas de
la Patagonia colonizadas por población proveniente de Buenos Aires y,
finalmente, amplias zonas del Uruguay. Se trata, en palabras de la autora, de
«la región lingüística más extensa y más europeizada de la República»
\parencite[76]{vidal1964}.

Un último punto merece atención. Como señalan \textcite[12]{ditullio2011}, el
español rioplatense no se corresponde con ninguna unidad política: no coincide
con las fronteras de un estado-nación, sino que se extiende a ambos lados del Río
de la Plata, abarcando territorio argentino y uruguayo. Esta observación, que
puede parecer obvia, tiene implicaciones importantes: nos recuerda que las
fronteras lingüísticas y las fronteras políticas rara vez coinciden. A pesar de
los intentos de homogeneización lingüística llevados adelante por el Estado
argentino entre fines del siglo \textsc{xix} y la primera mitad del \textsc{xx},
la diversidad dialectal del país sigue siendo rica y, en buena medida,
insuficientemente estudiada.

Veamos, entonces, cuáles son las particularidades que caracterizan al español de
la zona del Río de la Plata.

\section{Léxico}\label{sec:15-lexico}

Las lenguas no solo se diferencian por su pronunciación o su gramática: el
vocabulario es muchas veces el rasgo que los hablantes perciben con mayor
facilidad. En el caso del español rioplatense, el léxico presenta
particularidades que reflejan de manera directa la historia de contacto
lingüístico que describimos en la sección anterior.

En términos generales, las variedades del español del Cono Sur comparten un
conjunto de palabras tomadas de lenguas indígenas americanas, en particular del
quechua.

\begin{exe}
  \ex\label{ex:15-quechua} Préstamos del quechua en el español rioplatense:
    \begin{xlist}
      \ex \forma{cancha} (\sig{terreno delimitado para practicar deportes o
          realizar actividades})
      \ex \forma{choclo} (\sig{maíz tierno})
      \ex \forma{papa} (\sig{tubérculo comestible})
      \ex \forma{poroto} (\sig{tipo de legumbre})
    \end{xlist}
\end{exe}

En menor medida, el español rioplatense también conserva algunas voces de origen
guaraní, testimonio del contacto temprano entre ambas lenguas en la región del
Río de la Plata, según se señaló en §\ref{sec:15-historia}.

\begin{exe}
  \ex\label{ex:15-guarani} Préstamos del guaraní en el español rioplatense:
    \begin{xlist}
      \ex \forma{gurí} (\sig{niño})
      \ex \forma{pororó} (\sig{pochoclo} o \sig{palomitas de maíz})
      \ex \forma{yacaré} (\sig{tipo de caimán})
      \ex \forma{ananá} (\sig{piña})
    \end{xlist}
\end{exe}

Estas voces son de uso cotidiano en el español rioplatense, y muchos hablantes
las emplean sin saber que provienen de lenguas indígenas. La distribución de
estos préstamos no es homogénea entre las distintas variedades del español
americano: muchas voces de origen indígena que son comunes en una región resultan
desconocidas en otra. Estas diferencias se explican, en parte, por la diversidad
de \termino{lenguas de sustrato}\footnote{Se denominan \textit{lenguas de
sustrato} las lenguas que, en una situación de contacto lingüístico, son habladas
por la población originaria de un territorio cuando se impone una lengua
dominante (la lengua de superestrato). Aunque eventualmente la lengua de sustrato
puede dejar de hablarse, suele dejar huellas en la lengua dominante (típicamente
en el léxico, pero también en la fonética y la sintaxis). En el caso del español
americano, las lenguas indígenas de cada región funcionaron generalmente como
sustrato.} que han intervenido en la conformación léxica de cada variedad.

Sin embargo, el rasgo más distintivo del léxico rioplatense no son los
indigenismos, sino el llamado \termino{lunfardo}. ¿Qué es el lunfardo? En términos
simples, se trata de un conjunto de palabras y expresiones coloquiales,
originalmente asociadas al habla popular de Buenos Aires, que se han ido
incorporando al uso cotidiano de los hablantes rioplatenses a lo largo de más de
un siglo. Quien dice \forma{laburo} en lugar de \forma{trabajo}, \forma{mango} en
lugar de \forma{peso} (la moneda) o \forma{afanar} en lugar de \forma{robar}
está usando lunfardo, muchas veces sin ser consciente de ello. El lunfardo es un
repertorio dinámico y en constante transformación: cada generación incorpora
nuevas voces y deja de usar otras.\footnote{Varios de los fenómenos gramaticales
discutidos en este capítulo han sido estigmatizados por estar asociados al
lunfardo, que muchos veían como una jerga delictiva. Véanse
\textcite{gobello2001} y \textcite{conde2009} para una postura que reivindica el
lunfardo como una manifestación cultural legítima y un componente esencial de la
identidad rioplatense.}

La riqueza del lunfardo es el resultado de una interacción compleja entre varios
factores: la literatura y la música argentinas de comienzos del siglo
\textsc{xx} (en particular el tango, que difundió muchas de estas voces), el
proceso de consolidación de la variedad rioplatense y, de manera fundamental, la
contribución de los inmigrantes italianos y de sus dialectos.

\begin{recuadro}[¿Qué es el lunfardo?]
Según \textcite[14]{conde2004}, el lunfardo se conformó en su origen con términos
traídos por la inmigración, pero no es un vocabulario cerrado: hay un único
lunfardo, que en las décadas sucesivas se amplió con palabras provenientes de
ámbitos muy diversos, casi todas de creación local y sobre la base de la lengua
española. El lenguaje del fútbol y el del turf, las jergas de oficios y
profesiones, los ambientes de la droga, el mundo del rock y de las «tribus
urbanas», la jerga del psicoanálisis, la radio y la televisión, entre otros, han
aportado una cantidad innumerable de vocablos, extendidos ya a todo el espectro
social de buena parte del país.\pendiente{El original transcribe una cita
textual de Conde (2004, p.~14) de unas 120 palabras; se parafraseó: decidir si
se restituye}
\end{recuadro}

Numerosos términos del lunfardo son de origen italiano.

\begin{exe}
  \ex\label{ex:15-italianismos} Italianismos en el lunfardo rioplatense:
    \begin{xlist}
      \ex \forma{laburo} (\sig{trabajo}, del italiano \forma{lavoro})
      \ex \forma{birra} (\sig{cerveza}, del italiano \forma{birra})
      \ex \forma{bochar} (\sig{reprobar un examen}, del italiano
          \forma{bocciare})
      \ex \forma{gamba} (\sig{pierna}, del italiano \forma{gamba})
      \ex \forma{nona} (\sig{abuela}, del italiano \forma{nonna})
      \ex \forma{napia} (\sig{nariz}, del genovés \forma{nàppia}, \sig{nariz
          grande})
    \end{xlist}
\end{exe}

En menor medida, el lunfardo también incorporó voces de otras lenguas y
registros, como el francés, el portugués, el caló peninsular (la lengua de los
gitanos españoles), el guaraní y el quechua \parencite{conde2009}.

Un mecanismo particularmente interesante de creación léxica dentro del lunfardo
es el llamado \termino{vesre}. El término mismo ya nos da una pista de en qué
consiste: \forma{vesre} es la inversión silábica de la palabra \forma{revés}. El
procedimiento toma una palabra del léxico rioplatense y reordena sus sílabas, lo
que produce una forma con el mismo significado referencial, pero con un valor
expresivo diferente, generalmente más coloquial o lúdico.\footnote{Véase
\textcite{bohrn2015} para un análisis detallado de este fenómeno.}

\begin{exe}
  \ex\label{ex:15-vesre} Ejemplos de vesre en el español rioplatense:
    \begin{xlist}
      \ex \forma{café} $\rightarrow$ \forma{feca}
      \ex \forma{mujer} $\rightarrow$ \forma{jermu}
      \ex \forma{tango} $\rightarrow$ \forma{gotán}
      \ex \forma{pagar} $\rightarrow$ \forma{garpar}
    \end{xlist}
\end{exe}

Algunas formas de vesre se han convencionalizado de tal manera que su origen
resulta opaco, incluso para los propios hablantes rioplatenses. Un ejemplo
clásico es \forma{ortiba}, que significa \sig{delator} o \sig{aguafiestas}: pocos
hablantes reconocen que se trata de la inversión de \forma{batidor}, una palabra
del lunfardo que significa \sig{el que delata}.

El léxico del español rioplatense constituye, entonces, un terreno donde se
pueden rastrear con claridad las huellas del contacto lingüístico y de la
historia social de la región. En las secciones siguientes veremos que la
pronunciación, la morfología y la sintaxis de esta variedad presentan, cada una a
su manera, el mismo tipo de singularidad.

\begin{ejercicio}\label{ej:15-prestamos}
En la sección sobre léxico se presentaron préstamos del quechua, del guaraní y
del italiano en el español rioplatense. Elija tres palabras de uso cotidiano en
su variedad del español que sospeche que provienen de una lengua indígena o de
una lengua de inmigración. Investigue su etimología y discuta: ¿los hablantes son
generalmente conscientes del origen de estas palabras?
\end{ejercicio}

\section{Fonética y fonología}\label{sec:15-fonetica}

En esta sección describimos las propiedades más salientes de la pronunciación
del español rioplatense. Para representar los sonidos del habla, utilizamos el
Alfabeto Fonético Internacional (AFI), que se introdujo en el \capref{3}
(§\ref{sec:03-transcripcion}), un sistema de notación estándar en lingüística que
permite distinguir con precisión entre la forma escrita de una palabra y su
pronunciación efectiva. Como se vio en ese capítulo, las transcripciones
fonéticas aparecen entre corchetes ([~]), mientras que las representaciones
fonológicas aparecen entre barras (/~/).

Comenzamos por las propiedades \textit{segmentales}, es decir, las que conciernen
a los sonidos individuales (consonantes y vocales), para luego pasar a las
propiedades \textit{prosódicas}, que tienen que ver con la melodía y el ritmo del
habla. Recordemos que, como se dijo en el \capref{3}, al igual que la mayoría de
las variedades americanas del español, el rioplatense presenta \termino{seseo}:
la ausencia de la fricativa dental sorda /θ/, que en algunas variedades
peninsulares permite distinguir los grafemas <c> / <z> de la letra <s>. En
consecuencia, pares como \forma{casa} y \forma{caza} se pronuncian de manera
idéntica:

\begin{exe}
  \ex\label{ex:15-seseo} Seseo: \forma{casa}, \forma{caza} $\rightarrow$
      \fon{ˈka.sa}
\end{exe}

Este rasgo no es exclusivo del rioplatense, sino que lo comparte con
prácticamente todo el español americano. Más allá del seseo, el español
rioplatense presenta cuatro rasgos segmentales más o menos distintivos que
describimos a continuación.

En primer lugar, el rioplatense tiende a aspirar la /s/ cuando esta se encuentra
al final de una sílaba y antes de otra consonante.

\begin{exe}
  \ex\label{ex:15-aspiracion} Aspiración de /s/ en coda
      preconsonántica:\footnote{El concepto de coda silábica se introdujo en el
      \capref{3}, §\ref{sec:03-silaba}.}
    \begin{xlist}
      \ex \forma{asco} $\rightarrow$ \fon{ˈah.ko}
      \ex \forma{es lo mismo} $\rightarrow$ \fon{ˈeh.lo.ˈmih.mo}
    \end{xlist}
\end{exe}

\pendiente{El original transcribe \fon{ˈah.co}; se corrigió a \fon{ˈah.ko}}

También se observa la elisión completa de la /s/, especialmente al final de
palabra, aunque la aspiración es el fenómeno más frecuente. Este rasgo es
compartido con otras variedades del español argentino
\parencite{terrell1978,donni1984}.

En segundo lugar, el rioplatense exhibe patrones de diptongación de vocales
medias \parencite{luis1983}. En ciertos contextos, las vocales /e/ y /o/, cuando
están en contacto con otra vocal, tienden a pronunciarse como diptongos.

\begin{exe}
  \ex\label{ex:15-diptongacion} Diptongación de vocales medias:
    \begin{xlist}
      \ex \forma{línea} $\rightarrow$ \fon{ˈli.nja}
      \ex \forma{almohada} $\rightarrow$ \fon{al.ˈmwa.da}
    \end{xlist}
\end{exe}

En tercer lugar, se observa una pérdida progresiva de la distinción entre /ɲ/ y
la secuencia /nj/. La <ñ> del español (técnicamente, una consonante nasal que se
articula con la lengua contra el paladar) tiende a confundirse con la secuencia
formada por <n> + <i> seguida de vocal
\parencite{colantoni2013hualde,bongiovanni2015}. En la práctica, esto significa
que las secuencias ortográficas <ña>, <ñe>, <ño>, <ñu> no siempre se distinguen
de <nia>, <nie>, <nio>, <niu> en la pronunciación.

\begin{exe}
  \ex\label{ex:15-nj} Pérdida de la distinción /ɲ/ frente a /nj/:
      \forma{huraño}, \forma{uranio} $\rightarrow$ \fon{u.ˈɾa.njo}
\end{exe}

Por último, quizá el rasgo más emblemático del español rioplatense sea su forma
particular de \termino{yeísmo}, que se conoce como \termino{yeísmo rehilado}. El
término \textit{rehilamiento} alude a la fricción intensa o vibración acústica
característica de este tipo de sonidos sibilantes. En la mayoría de las
variedades yeístas del español, los grafemas <ll> e <y> se pronuncian de manera
idéntica, con una fricativa palatal sonora \fon{ʝ}. En el rioplatense, en cambio,
ambos grafemas se realizan mediante una consonante fricativa que se produce con
la lengua en una posición más adelantada, similar a la de <sh> en inglés
\parencite[entre muchos otros]{zamora1949,guitarte1955}. Además, se observa una
tendencia histórica al ensordecimiento de este sonido: mientras que las
generaciones mayores suelen producir una variante sonora \fon{ʒ} (similar a la
<j> del francés \forma{je}), los hablantes más jóvenes tienden a usar una
variante sorda \fon{ʃ} (similar a la <sh> del inglés \forma{she}).\footnote{Se
puede consultar \textcite{rohena2013} para una discusión detallada del proceso de
ensordecimiento.}

\begin{exe}
  \ex\label{ex:15-rehilado} Yeísmo rehilado:
    \begin{xlist}
      \ex \forma{pollo} $\rightarrow$ \fon{ˈpo.ʒo} (variante sonora,
          generalmente asociada a hablantes mayores)
      \ex \forma{pollo} $\rightarrow$ \fon{ˈpo.ʃo} (variante sorda,
          generalmente asociada a hablantes jóvenes)
    \end{xlist}
\end{exe}

Además de sus sonidos particulares, el español rioplatense se distingue por su
\termino{entonación}: la «melodía» con la que los hablantes producen las
oraciones. Se trata de un rasgo que los propios hablantes de otras variedades del
español perciben con facilidad y que, como veremos, tiene una conexión histórica
con la inmigración italiana.

Un primer rasgo prosódico distintivo es la llamada \termino{alineación temprana
del pico tonal}, que afecta a los acentos prenucleares, es decir, a los acentos
que preceden al acento más prominente del enunciado. En la mayoría de las
variedades del español, cuando un hablante pronuncia una oración declarativa
neutra, el punto más alto de la curva melódica ---el pico tonal--- tiende a
ubicarse después de la sílaba acentuada, en la sílaba siguiente. En el español
rioplatense, en cambio, ese pico se produce tempranamente dentro de la sílaba
acentuada, un patrón que se asemeja a lo que se observa en varias lenguas
habladas en Italia \parencite{colantoni2004}; la figura~\ref{fig:15-pico}
esquematiza este contraste.

\figurapendiente{Patrones de alineación del pico tonal en acentos de oraciones declarativas neutras (adaptado de Colantoni y Gurlekian, 2004)}{fig:15-pico}

El acento nuclear, es decir, el último y más prominente del enunciado, presenta
un comportamiento distinto, aunque igualmente característico. En el rioplatense,
la melodía desciende de manera abrupta dentro de la propia sílaba acentuada y
alcanza el punto más bajo de todo el enunciado en un lapso muy breve, en lo que
\textcite{kaisse2001} denomina una «larga caída». En otras variedades del
español, esta caída final es más gradual y se distribuye a lo largo de más
sílabas. El resultado auditivo en el rioplatense es una melodía que parece
«cortarse» abruptamente al final de la oración. Este patrón es también
característico de las lenguas de Italia habladas por una parte significativa de
los inmigrantes que se asentaron en Buenos Aires.

\textcite{colantoni2004} interpretan estos patrones como el resultado de un
proceso de convergencia: dos sistemas entonativos, el español y el italiano, que
ya eran tipológicamente similares, se volvieron aún más parecidos a consecuencia
del contacto. La influencia del italiano sobre el español rioplatense habría sido
tanto directa (contacto cotidiano con hablantes italianos) como indirecta, a
través del lunfardo y del \termino{cocoliche}.\footnote{\textit{Cocoliche} es el
nombre que recibe el español hablado por los inmigrantes, principalmente
italianos, llegados a Buenos Aires entre las últimas décadas del siglo
\textsc{xix} y las primeras del \textsc{xx}. Se trata del resultado de la
adquisición incompleta del español por parte de hablantes adultos, con fuerte
interferencia de su lengua de origen.}

Otro rasgo prosódico relevante es la posibilidad de acentuar los pronombres
clíticos cuando estos aparecen después del verbo. Estos son pronombres como
\forma{lo, la, me, se}, que normalmente no llevan acento propio y que se apoyan
prosódicamente en el verbo (véase el \capref{10}, §\ref{sec:10-tonicos}). En el
común de las variedades del español, cuando un clítico se añade al final de un
verbo (por ejemplo, \forma{aGÁrrenlo, DÁmelo}), el acento permanece en la sílaba
acentuada del verbo. En el rioplatense, sin embargo, es común que sea el clítico
el que recibe el acento.

\begin{exe}
  \ex\label{ex:15-encliticos} Acentuación de enclíticos:
    \begin{xlist}
      \ex agarrenLÓ
      \ex dameLÓ
      \ex disculpandoLÁ
    \end{xlist}
\end{exe}

Como observan \textcite{colantoni2013cuervo}, esta acentuación no tiene valor
enfático: no se trata de que el hablante quiera resaltar el pronombre, sino que
esta pronunciación parece ser la opción por defecto para muchos hablantes
rioplatenses. Las autoras sugieren, además, que este fenómeno podría vincularse
con la influencia italiana, ya que numerosos dialectos del sur de Italia también
presentan acentuación de clíticos en esta posición.\footnote{Para una discusión
más detallada de los aspectos entonativos de la variedad, véanse
\textcite{gabriel2010,gabriel2011} y \textcite{colantoni2014}.}

\begin{destacado}
En conjunto, las propiedades prosódicas del español rioplatense confirman una
observación que los estudiosos de esta variedad han hecho desde hace décadas: su
entonación tiene un carácter notablemente italianizante \parencite{vidal1964},
producto de la profunda huella que la inmigración italiana dejó en el habla de la
región.
\end{destacado}

\begin{ejercicio}\label{ej:15-transcripcion}
Considere las siguientes oraciones del español y compare su pronunciación en una
variedad peninsular estándar con su pronunciación en el español rioplatense.
Realice una transcripción fonética aproximada de cada una en ambas variedades,
prestando especial atención al yeísmo rehilado, la aspiración de /s/ en coda
preconsonántica, la realización de los grafemas <ll> e <y> y el seseo.

\begin{exe}
  \ex \textit{Ella mira el cielo.}
  \ex \textit{Es muy callado.}
  \ex \textit{Hasta mañana.}
\end{exe}
\end{ejercicio}

\section{Morfología}\label{sec:15-morfologia}

En lo que respecta a la morfología, nos concentramos en dos fenómenos: el voseo,
que pertenece al ámbito de la morfología flexiva, y el uso del prefijo
\forma{re-}, que corresponde a la morfología derivativa.

\subsection{El voseo}\label{sec:15-voseo}

Quizá el rasgo morfológico más conocido del español rioplatense sea el
\termino{voseo}. El voseo consiste en el uso de \forma{vos} en lugar de
\forma{tú} como pronombre de segunda persona singular. Este fenómeno tiene dos
dimensiones. Por un lado, está el \termino{voseo pronominal}: el hecho de que se
use el pronombre \forma{vos} en lugar de \forma{tú}, característico del español
general. Por otro lado, está el \termino{voseo verbal}: el hecho de que las
formas flexivas del verbo también cambien (p.~ej., se dice \forma{sabés} en lugar
de \forma{sabes}, o \forma{tenés} en lugar de \forma{tienes}). En el español de
Buenos Aires, ambas dimensiones van generalmente juntas: se usa \forma{vos} y se
conjugan los verbos con las flexiones voseantes correspondientes. Se trata de la
forma de tratamiento no marcada, es decir, la que se emplea por defecto en la
mayor parte de los contextos cotidianos. El uso de \forma{vos} alterna, sin
embargo, con el de \forma{usted}, que se reserva para situaciones de mayor
distancia o formalidad (trato con desconocidos, contextos profesionales o
jerárquicos, etc.).

\begin{exe}
  \ex\label{ex:15-voseo} Ejemplos de voseo en el español rioplatense:
    \begin{xlist}
      \ex ¿No tenés nariz, vos? ¿No olés como yo? ¡A podrido huele!
      \ex Lamponne, si me escuchás levantá el brazo derecho\ldots\ ¡El derecho,
          Lamponne!
      \ex ¡Sé vos nomás!
    \end{xlist}
\end{exe}

¿Cómo se conjuga un verbo en el sistema voseante? La
tabla~\ref{tab:15-voseo} presenta el paradigma básico para las tres
conjugaciones regulares (\forma{-ar, -er, -ir}) y para el verbo irregular
\forma{ser}. Aunque la tabla muestra las formas canónicas del voseo, el uso real
no siempre las sigue al pie de la letra. En particular, para el pretérito
perfecto simple, la mayoría de los hablantes rioplatenses prefiere las formas de
tuteo: por lo general se dice \forma{amaste} y no \forma{amastes},
\forma{comiste} y no \forma{comistes}. Algo similar ocurre con el presente del
subjuntivo, donde las formas voseantes (\forma{comás, vivás}) alternan con las
del tuteo (\forma{comas, vivas}).

\begin{table}[htbp]
  \centering\small
  \caption{Formas canónicas de voseo verbal en el español
    rioplatense}\label{tab:15-voseo}
  \begin{tabular}{@{}lllll@{}}
    \toprule
    & \textit{ser} & \textit{amar} & \textit{comer} & \textit{vivir} \\
    \midrule
    Presente del indicativo & \textit{sos} & \textit{amás} & \textit{comés} &
      \textit{vivís} \\
    Pretérito perfecto simple & \textit{fuistes} & \textit{amastes} &
      \textit{comistes} & \textit{vivistes} \\
    Presente del subjuntivo & \textit{seás} & \textit{amés} & \textit{comás} &
      \textit{vivás} \\
    Imperativo & \textit{sé} & \textit{amá} & \textit{comé} & \textit{viví} \\
    \bottomrule
  \end{tabular}
\end{table}

\pendiente{En la tabla, el imperativo de \textit{ser} estaba vacío; se completó
con \textit{sé} (cf. \exref{ex:15-voseo})}

El panorama del voseo no es uniforme dentro de la región rioplatense. Existe un
contraste importante entre Buenos Aires y Montevideo. En Buenos Aires, como
dijimos, lo habitual es la combinación de voseo pronominal y voseo verbal: se
dice \forma{vos cantás, vos debés}. Las formas del tuteo quedan restringidas a
contextos muy marcados, como ciertos registros mediáticos. En Montevideo, en
cambio, el sistema es más complejo: los hablantes alternan entre \forma{vos} y
\forma{tú}, y ambos pronombres pueden combinarse con la morfología verbal
voseante.

\begin{exe}
  \ex\label{ex:15-montevideo} Combinaciones de pronombre y morfología verbal en
      Montevideo:
    \begin{xlist}
      \ex \forma{vos tenés} (voseo pronominal + voseo verbal)
      \ex\label{ex:15-montevideo-b} \forma{tú tenés} (tuteo pronominal + voseo
          verbal)
      \ex \forma{tú tienes} (tuteo pronominal + tuteo verbal)
    \end{xlist}
\end{exe}

Estas combinaciones no se distribuyen de manera aleatoria, sino que están
asociadas a factores de registro, estilo y valoración social. En particular, el
patrón \exref{ex:15-montevideo-b}, que combina el pronombre \forma{tú} con la
morfología verbal voseante, funciona como una solución intermedia: permite
mantener el pronombre \forma{tú}, asociado al prestigio normativo y escolar, sin
abandonar las terminaciones verbales propias de la región
\parencite{bertolotti2011}. Otras variantes de voseo se documentan también dentro
del territorio argentino.\footnote{Se pueden consultar \textcite{ditullio2010b} y
\textcite{estomba2012} para una discusión relevante.}

\begin{ejercicio}\label{ej:15-voseo}
En el capítulo se describió el contraste entre el sistema de voseo de Buenos
Aires (donde se usa \forma{vos cantás}) y el de Montevideo (donde conviven
\forma{vos tenés}, \forma{tú tenés} y \forma{tú tienes}). Investigue qué sucede
con el voseo en al menos otras dos regiones de América Latina. ¿Se usa
\forma{vos}? ¿Con qué formas verbales? ¿Convive con \forma{tú} u otras formas de
segunda persona?
\end{ejercicio}

\subsection{El prefijo \textit{re-}}\label{sec:15-re}

Pasemos ahora a la morfología derivativa. Un fenómeno particularmente peculiar
del español rioplatense es el uso del prefijo \forma{re-}
\parencite{giammatteo2012albano,kornfeld2013}. A primera vista, \forma{re-}
funciona como un intensificador que se combina con distintas clases de palabras
(adjetivos, adverbios, nombres y verbos) de manera similar a \forma{muy}. Patrones
como los siguientes se documentan en diversas variedades del español.

\begin{exe}
  \ex\label{ex:15-re} \textit{re} lindo, \textit{re} alta, \textit{re}
      inteligente; \textit{re} bien, \textit{re} lejos; \textit{re} amigo,
      \textit{re} burro.\footnote{Según el uso consolidado en español, el prefijo
      \forma{re-}, en su valor de intensificador, se escribe separado de la base
      a la cual se aplica.}
\end{exe}

Con respecto a los verbos, \textcite{giammatteo2012albano} señalan que, a
diferencia de lo que sucede en el español general, en el español rioplatense son
usuales las combinaciones de \forma{re-} con verbos psicológicos (\forma{re
gustar, re querer, re coparse, re morir} (por algo o alguien)) con valor de
intensificación, y que con verbos télicos adopta sentido conclusivo
(\forma{recagar} (a alguien),\footnote{En la interpretación de \sig{superar o
vencer} a otro, propia del español de la Argentina
\parencite[80]{conde2004}.} \forma{requemarse}).

Asimismo, lo que hace peculiar al \forma{re-} rioplatense es que su uso va más
allá de las distintas clases de palabras. Como muestran
\textcite{kornfeld2013}, este prefijo puede combinarse incluso con oraciones
completas.

\begin{exe}
  \ex\label{ex:15-reext} Usos extendidos de \forma{re-} en el español
      rioplatense:
    \begin{xlist}
      \ex \forma{re-} + verbo finito (con interpretación similar a
          \forma{mucho}): \textit{re} estudiamos, \textit{re} trabajamos,
          \textit{re} jugamos
      \ex \forma{re-} + oración completa (con interpretación similar a
          \forma{seguro que\ldots}): \textit{re} que estudié, \textit{re} que
          trabajamos, \textit{re} que jugamos
    \end{xlist}
\end{exe}

Como puede verse, el nivel en el que se adjunta el prefijo produce cambios de
significado: cuando se combina con un verbo, cuantifica la acción (\forma{re
estudiamos} = \sig{estudiamos mucho}); cuando se combina con una oración
completa, adquiere un valor de certeza o énfasis (\forma{re que estudié} =
\sig{seguro que estudié}, \sig{de verdad que estudié}). Este comportamiento es
llamativo porque los prefijos normalmente se adhieren a palabras individuales,
no a oraciones enteras. Por eso, \textcite{kornfeld2013} analizan \forma{re-}
como un \termino{prefijo frasal}: un prefijo que puede operar sobre unidades más
grandes que la palabra. Evidencia adicional de esta independencia es que
\forma{re-} puede funcionar incluso como respuesta a una pregunta.

\begin{exe}
  \ex\label{ex:15-rerespuesta} Uso de \forma{re-} como respuesta:\\
      A: ¿Venís a mi casa?\\
      B: Re. (= \sig{Sí, seguro / por supuesto}).
\end{exe}

Los fenómenos presentados en esta sección, el voseo y el uso del prefijo
\forma{re-}, son apenas dos muestras de la riqueza morfológica del español
rioplatense, pero ilustran bien un punto general: las particularidades de un
dialecto no se limitan a la pronunciación o al vocabulario, sino que afectan
también la estructura interna de las palabras y sus combinaciones. En la sección
siguiente veremos que la sintaxis del rioplatense presenta, a su vez, rasgos
igualmente distintivos.

\section{Sintaxis}\label{sec:15-sintaxis}

En las secciones anteriores vimos que el español rioplatense presenta
particularidades en su léxico, su pronunciación y su morfología. En esta sección
nos ocupamos de su sintaxis, ya que también en este plano el rioplatense exhibe
rasgos distintivos, muchos de los cuales involucran elementos funcionales como
pronombres átonos, partículas negativas y operadores de foco, que se comportan de
manera diferente a como lo hacen en otras variedades del español. Organizamos la
presentación en torno a cuatro fenómenos: el doblado de clíticos, el doblado
verbal, el doblado negativo enfático y el uso de \forma{solo} como operador
focal adjetival.

\subsection{Doblado de clíticos}\label{sec:15-doblado}

Uno de los rasgos más prominentes de la sintaxis rioplatense es la posibilidad de
doblar clíticos acusativos con frases nominales plenas, un fenómeno que en otras
variedades del español está considerablemente más restringido. Para entender en
qué consiste el fenómeno, consideramos primero los casos básicos. Cuando un verbo
tiene un objeto, el español ofrece dos posibilidades: expresarlo mediante un
sintagma nominal completo o reemplazarlo por un pronombre átono (un clítico como
\forma{lo, la, le}, etc.).

\begin{exe}
  \ex\label{ex:15-od-sn} Objeto directo como sintagma nominal:
    \begin{xlist}
      \ex Santos miró a Rosa.
      \ex La vieja tomó al llorón de la mano.
    \end{xlist}
  \ex\label{ex:15-od-cl} Objeto directo como clítico:
    \begin{xlist}
      \ex Santos \textit{la} miró.
      \ex La vieja \textit{lo} tomó de la mano.
    \end{xlist}
\end{exe}

En muchas variedades del español, estas dos opciones son mutuamente excluyentes:
se usa una o la otra, pero no ambas a la vez. En el español rioplatense, en
cambio, es perfectamente natural que el clítico y el sintagma nominal completo
aparezcan juntos en la misma oración. Esto es lo que se denomina \termino{doblado
de clíticos}.

\begin{exe}
  \ex\label{ex:15-doblado} Doblado de objeto directo \parencite[adaptado
      de][148]{kany1969}:
    \begin{xlist}
      \ex Santos \textit{la} miró \textit{a Rosa}.
      \ex La vieja \textit{lo} tomó \textit{al llorón} de la mano.
    \end{xlist}
\end{exe}

En estos ejemplos, el clítico (\forma{la, lo}) y el sintagma nominal completo
(\forma{a Rosa, al llorón}) se refieren a la misma persona. El clítico «dobla» o
«duplica» al objeto directo. No se trata de una redundancia accidental, sino de un
patrón gramatical productivo y sistemático en la variedad. Conviene aclarar que el
doblado de clíticos como tal no es exclusivo del rioplatense: en todas las
variedades del español los pronombres tónicos deben ir doblados por clíticos
(p.~ej., \forma{me vio a mí}), y el doblado de objeto indirecto con \forma{le} es
frecuente en numerosos dialectos. Lo distintivo del rioplatense es, por un lado,
la extensión del doblado a objetos directos no pronominales y, como veremos más
abajo, la cuasiobligatoriedad del doblado de objeto indirecto.

¿Con qué tipos de objetos directos se produce el doblado? Los ejemplos anteriores
involucran sintagmas nominales definidos (nombres comunes precedidos por artículo
o nombres propios) y animados (se refieren a personas). Estos son los casos más
claros y aceptados. Sin embargo, el doblado también puede extenderse a objetos
directos indefinidos, siempre que se refieran a una entidad específica (es decir,
a alguien que el hablante tiene en mente, aunque el oyente no necesariamente sepa
quién es).\footnote{\textcite{suner1988} propone que el doblado de objeto directo
se restringe a sintagmas nominales que expresan especificidad. La cuestión de si
esta generalización se mantiene en todos los contextos ha sido objeto de debate,
para lo que se pueden consultar \textcite{zdrojewski2008} y
\textcite{ditullio2019}.}

\begin{exe}
  \ex\label{ex:15-indef} Doblado de objeto directo indefinido:
    \begin{xlist}
      \ex Dos personas encapuchadas \textit{la} asaltaron \textit{a una
          mujer}\ldots
      \ex Cuando iba al vestuario \textit{lo} escupió \textit{a un rival}.
      \ex \textit{La} fajaron \textit{a una piba de 12} en el colegio.
    \end{xlist}
\end{exe}

En lo que respecta al doblado de objeto indirecto, el español rioplatense va
considerablemente más lejos que otras variedades. En la mayoría de los dialectos
del español, el doblado del objeto indirecto con el clítico \forma{le} es
frecuente, pero no siempre obligatorio. En el rioplatense, en cambio, el doblado
resulta obligatorio en casi todos los contextos, incluso con pronombres
indefinidos y nombres propios, en casos donde algunas variedades peninsulares no
lo requieren.

\begin{exe}
  \ex\label{ex:15-oi} Doblado de objeto indirecto: contraste con algunas
      variedades peninsulares:
    \begin{xlist}
      \ex Eso no interesa a nadie. \hfill peninsular: sí / rioplatense: no
      \ex Eso no \textit{le} interesa \textit{a nadie}. \hfill peninsular: sí /
          rioplatense: sí
      \ex No daba dinero a Juana. \hfill peninsular: sí / rioplatense: no
      \ex No \textit{le} daba dinero \textit{a Juana}. \hfill peninsular: sí /
          rioplatense: sí
    \end{xlist}
\end{exe}

El contraste es claro: en algunas variedades peninsulares, las versiones sin
clítico son aceptables; en el rioplatense, resultan un tanto extrañas o
directamente inaceptables. Incluso en contextos donde otras variedades tienden a
rechazar el doblado, como con nombres desnudos sin determinante, el rioplatense
lo prefiere.

\begin{exe}
  \ex\label{ex:15-desnudos} Doblado de objeto indirecto con sustantivos sin
      determinante:
    \begin{xlist}
      \ex (*\textit{Les}) legó su fortuna \textit{a organizaciones de caridad}.
          \hfill otras variedades
      \ex Juan (\textit{les}) dio entradas \textit{a niños carenciados}. \hfill
          rioplatense
    \end{xlist}
\end{exe}

En síntesis, si bien en todas las variedades del español el doblado del objeto
indirecto presenta una extensión mucho mayor que el del objeto directo, en el
español rioplatense esta extensión es mucho mayor aún. La
tabla~\ref{tab:15-doblado} resume esta asimetría.

\begin{table}[htbp]
  \centering\small\setlength{\tabcolsep}{3pt}
  \caption{Distribución del doblado de clíticos según el tipo de
    objeto}\label{tab:15-doblado}
  \begin{tabular}{@{}lcccccc@{}}
    \toprule
    & \makecell{Pro-\\nombres} & \makecell{SN\\animados} &
      \makecell{SN\\indefinidos} & \makecell{SN no\\pronominales} &
      \makecell{N sin\\determinante} & \makecell{N\\genéricos} \\
    \midrule
    OD & Sí & Sí & \makecell{± (si son\\específicos)} & No & No & No \\
    OI & Sí & Sí & Sí & Sí & Sí & Sí \\
    \bottomrule
  \end{tabular}
\end{table}

\pendiente{La tabla da «No» para OD con «SN no pronominales», en contradicción
con el texto (el rioplatense extiende el doblado a OD no pronominales, como en
\exref{ex:15-doblado}). ¿Se quiso decir «SN inanimados»? Aun así, el doblado de
OD inanimados está atestiguado en el rioplatense: revisar con los autores}

\begin{ejercicio}\label{ej:15-juicios}
Como se señaló, el doblado de clíticos de objeto indirecto es prácticamente
obligatorio en el español rioplatense, mientras que en otras variedades es
opcional. Considere las siguientes oraciones y juzgue, según su variedad, si le
resultan naturales, extrañas o inaceptables.

\begin{exe}
  \ex
    \begin{xlist}
      \ex Esto no le gusta a nadie.
      \ex Esto no gusta a nadie.
    \end{xlist}
  \ex
    \begin{xlist}
      \ex Le regalé el libro a María.
      \ex Regalé el libro a María.
    \end{xlist}
\end{exe}

Compare sus juicios con los que se describen en el capítulo para el español
rioplatense y para el español peninsular. ¿Qué le dicen estas diferencias sobre
la variación sintáctica entre dialectos?
\end{ejercicio}

\subsection{Doblado verbal}\label{sec:15-dobladoverbal}

El segundo fenómeno sintáctico que presentamos es el llamado \termino{doblado
verbal}. En esta construcción, el verbo se repite al final del enunciado. Dado
que siempre debe haber un constituyente entre ambos verbos,
\textcite{kany1969} señala que en la Argentina este patrón se conoce
coloquialmente como \textit{hablar en sángüiche} (o también \textit{hablar en
capicúa}).

\begin{exe}
  \ex\label{ex:15-sanguche} Doblado verbal en el español rioplatense
      \parencite[adaptado de][315--316]{kany1969}:
    \begin{xlist}
      \ex ¿Son muchas las yeguas? ---No, señora. \textit{Son} ocho nomás,
          \textit{son}.
      \ex \textit{Tené} pasensia, \textit{tené}\ldots
      \ex \textit{Soy} un desgraciado, \textit{soy}.
      \ex \textit{Tengo} sentimientos, \textit{tengo}.
    \end{xlist}
\end{exe}

Kany describe este fenómeno como «propio de las personas incultas», un estigma
frecuentemente asociado a fenómenos relacionados con el lunfardo. Más allá de su
valoración social, el doblado verbal presenta propiedades gramaticales
sistemáticas que lo distinguen de una simple repetición de palabras. Un indicio
de esto es su comportamiento con los pronombres clíticos. En el español
rioplatense, como en la mayoría de las variedades del español, cuando un objeto
directo definido se omite, es obligatorio reemplazarlo por un pronombre clítico.

\begin{exe}
  \ex\label{ex:15-libro} Obligatoriedad del clítico con objetos definidos:\\
      A: ¿Compraste el libro?\\
      B: Sí, *(\textit{lo}) compré.
\end{exe}

Como nota \textcite{saab2008}, el doblado verbal presenta una restricción clara
con respecto a estos pronombres: ambas versiones del verbo deben coincidir en la
presencia o ausencia de clíticos. Si el primer verbo lleva un pronombre clítico,
el segundo debe llevarlo también; si el primero no lo lleva, el segundo tampoco
puede llevarlo.

\begin{exe}
  \ex\label{ex:15-simetria} Simetría de clíticos en el doblado verbal:
    \begin{xlist}
      \ex ¡Compró el libro, (*\textit{lo}) compró!
      \ex Lo compré ayer, *(\textit{lo}) compré.
    \end{xlist}
\end{exe}

Este comportamiento sugiere que el doblado verbal no consiste en dos oraciones
independientes: si este fuera el caso, el clítico \forma{lo} asociado a la
segunda ocurrencia de \forma{compró} debería ser obligatorio en todos los casos.
Por ello, \textcite{saab2008} sostiene que el doblado verbal involucra una única
estructura sintáctica en la que un mismo núcleo verbal aparece dos veces por
razones estrictamente gramaticales. Esto explica por qué los clíticos asociados al
verbo deben aparecer obligatoriamente en ambas ocurrencias.

\textcite{munozperez2021} propone que los patrones de doblado verbal del
rioplatense constituyen un tipo de construcción exclamativa. Evidencia para ello
es que ambos tipos de oración presentan restricciones similares. Por ejemplo,
tanto las exclamativas como el doblado verbal admiten adverbios de certeza, pero
rechazan adverbios de duda.

\begin{exe}
  \ex\label{ex:15-certeza}
    \begin{xlist}
      \ex ¡Estás equivocado seguro!
      \ex Voy seguro, voy.
      \ex[*]{¡Estás equivocado tal vez!}
      \ex[*]{Voy tal vez, voy.}
    \end{xlist}
\end{exe}

Esta restricción compartida sugiere que el doblado verbal y las exclamativas
pertenecen al mismo tipo de construcción gramatical, lo cual refuerza la idea de
que el fenómeno, lejos de ser un simple recurso retórico, tiene una estructura
sintáctica propia.

\subsection{Doblado negativo enfático}\label{sec:15-negativo}

El tercer tipo de construcción duplicativa que presentamos es el
\termino{doblado negativo enfático}, estudiado por \textcite{ditullio2008} (véase
también \cite{johnson2019}). Consideremos el siguiente intercambio:

\begin{exe}
  \ex\label{ex:15-nada1} Doblado negativo enfático:\\
      A: ¡Vení acá!\\
      B: ¡\textit{No} voy \textit{nada}!
\end{exe}

A primera vista, la respuesta de B podría parecer un uso normal de la palabra
\forma{nada}. Sin embargo, el verbo \forma{voy} en este ejemplo no selecciona un
complemento: \forma{nada} no funciona como un objeto directo, sino como una
partícula enfática que refuerza la negación. La oración \forma{¡No voy nada!}
significa algo así como \sig{de ninguna manera voy} o \sig{no voy en absoluto}.

Lo mismo puede ocurrir con verbos transitivos cuyo objeto ya está presente. En
\exref{ex:15-nada2}, el verbo \forma{comprar} ya tiene su objeto (\forma{lo}). El
elemento \forma{nada} no refiere a lo que se compra, sino que simplemente
enfatiza el rechazo.

\begin{exe}
  \ex\label{ex:15-nada2} A: ¡Comprá eso!\\
      B: ¡\textit{No} lo compro \textit{nada}!
\end{exe}

Otra diferencia entre el \forma{nada} enfático y el cuantificador \forma{nada}
tiene que ver con su posición en la oración. Como es sabido, los cuantificadores
negativos en español pueden aparecer antes o después del verbo, dependiendo de si
hay o no una negación explícita.

\begin{exe}
  \ex\label{ex:15-nadacuant} \forma{Nada} como cuantificador: alternancia
      posicional:
    \begin{xlist}
      \ex \textit{No} compró \textit{nada}.
      \ex \textit{Nada} compró.
    \end{xlist}
\end{exe}

El \forma{nada} enfático, en cambio, solo puede aparecer después del verbo; si
se lo mueve al inicio de la oración, el resultado es inaceptable.

\begin{exe}
  \ex\label{ex:15-nadafija} \forma{Nada} enfático: posición fija:\\
      A: ¡Comprá eso!\\
      B: *¡\textit{Nada} lo compro!
\end{exe}

Este comportamiento confirma que se trata de dos \forma{nada} distintos desde el
punto de vista gramatical, aunque suenen igual. De hecho, el patrón no se limita
a la palabra \forma{nada}: otras expresiones negativas enfáticas muestran un
comportamiento similar, lo que sugiere que estamos ante un fenómeno sintáctico
sistemático y no ante un caso aislado.

\begin{exe}
  \ex\label{ex:15-mierda} Otras expresiones con función de negación enfática:
    \begin{xlist}
      \ex No voy \textit{una mierda}.
      \ex No lo compro \textit{una mierda}.
    \end{xlist}
\end{exe}

\subsection{Operadores focales adjetivales}\label{sec:15-solo}

El último fenómeno sintáctico que presentamos involucra la palabra \forma{solo} y
un uso particular que se observa en el español rioplatense, aunque su
distribución dialectal precisa aún no ha sido establecida. Para entenderlo,
necesitamos distinguir primero dos usos diferentes de esta palabra que existen en
todas las variedades del español. Por un lado, \forma{solo} puede funcionar como
un adjetivo calificativo, con el significado de \sig{sin compañía}.

\begin{exe}
  \ex\label{ex:15-soloadj} \forma{Solo} como adjetivo calificativo:
    \begin{xlist}
      \ex Cosmo vino solo. $\approx$ \sig{Cosmo vino sin compañía}.
      \ex Los hombres solos son tristes. $\approx$ \sig{Los hombres que están sin
          compañía son tristes}.
    \end{xlist}
\end{exe}

Por otro lado, \forma{solo} puede funcionar como un \termino{adverbio de foco},
con un significado completamente diferente: indica que lo dicho se limita
exclusivamente a lo que se menciona, descartando otras posibilidades. En este
uso, \forma{solo} es equivalente a \forma{solamente} o \forma{únicamente}.

\begin{exe}
  \ex\label{ex:15-soloadv} \forma{Solo} como adverbio de foco:
    \begin{xlist}
      \ex Pedimos solo una docena de empanadas. $\approx$ \sig{Lo único que
          pedimos fue una docena de empanadas}.
      \ex Eliana trabaja solo los viernes. $\approx$ \sig{El único día que
          Eliana trabaja es el viernes}.
    \end{xlist}
\end{exe}

En los ejemplos anteriores, \forma{solo} introduce lo que podemos llamar una
\termino{lectura exhaustiva}: la referencia del elemento focalizado agota todas
las opciones relevantes: no se pidió nada más que empanadas; Eliana no trabaja
ningún otro día.

Lo interesante es que, en el español rioplatense, la forma adjetival de
\forma{solo}, es decir, la que concuerda en género y número con el sustantivo,
también puede producir esta lectura exhaustiva.

\begin{exe}
  \ex\label{ex:15-soloexh} \forma{Solo} adjetival con lectura exhaustiva:
    \begin{xlist}
      \ex Pedimos una docena sola de empanadas. $\approx$ \sig{Lo único que
          pedimos fue una docena de empanadas}.
      \ex María sola fue a la reunión. $\approx$ \sig{La única persona que fue a
          la reunión fue María}.
    \end{xlist}
\end{exe}

Esta lectura exhaustiva del adjetivo parece estar condicionada por su posición en
la oración. Específicamente, requiere que el adjetivo se combine directamente con
el nombre al que modifica. Esto se observa en el siguiente contraste.

\begin{exe}
  \ex\label{ex:15-eliana} Efecto de la posición del adjetivo:
    \begin{xlist}
      \ex\label{ex:15-eliana-a} Eliana llegó a la oficina sola.\\
          $\approx$ \sig{Eliana llegó sin compañía a la oficina}. \hfill
          \textit{lectura calificativa}
      \ex\label{ex:15-eliana-b} Eliana sola llegó a la oficina.\\
          $\approx$ \sig{Eliana llegó sin compañía a la oficina}. \hfill
          \textit{lectura calificativa}\\
          $\approx$ \sig{La única que llegó a la oficina fue Eliana}. \hfill
          \textit{lectura exhaustiva}
    \end{xlist}
\end{exe}

En \exref{ex:15-eliana-a}, \forma{sola} aparece al final de la oración, separada
del nombre, y solo admite la lectura calificativa, con interpretación de
\sig{sin compañía}. En \exref{ex:15-eliana-b}, en cambio, \forma{sola} aparece
inmediatamente después del nombre propio \forma{Eliana}, y esto habilita, además
de la lectura calificativa, la exhaustiva: la oración puede interpretarse como que
la única persona que llegó fue Eliana. A nuestro leal saber y entender, este
fenómeno no ha sido abordado en la literatura especializada. Si bien lo
presentamos aquí porque se atestigua de manera natural en el rioplatense, su
delimitación dialectal precisa requiere un estudio sistemático que confirme estas
observaciones preliminares y permita determinar en qué otras variedades del
español se manifiesta.

\begin{ejercicio}\label{ej:15-estigma}
En varias secciones del capítulo se mencionó la estigmatización de ciertos rasgos
lingüísticos: las lenguas de inmigración en la Buenos Aires del siglo
\textsc{xix}, el lunfardo como «jerga delictiva», el doblado verbal como rasgo de
«personas incultas». Elija un rasgo lingüístico que en su comunidad sea objeto de
valoración negativa (una pronunciación, una construcción, un uso léxico).
Descríbalo y discuta: ¿la valoración negativa tiene algún fundamento gramatical o
se basa enteramente en factores sociales? ¿Qué relación tiene esto con la
discusión sobre lengua y dialecto presentada en la introducción?
\end{ejercicio}

\begin{recuadro}[Reflexiones finales]
A lo largo de este capítulo hemos intentado ofrecer un panorama del español
rioplatense que, sin pretender ser exhaustivo, permita apreciar la riqueza y la
complejidad de esta variedad. Vale la pena detenerse un momento a mirar el camino
recorrido. Comenzamos por una cuestión que puede parecer puramente
terminológica, pero que tiene implicaciones profundas: ¿qué es una lengua y qué
es un dialecto? Vimos que la distinción habitual entre ambos términos descansa en
gran medida en factores extralingüísticos (históricos, políticos, culturales) y
que, desde un punto de vista estrictamente gramatical, resulta difícil
diferenciar ambas nociones, e incluso es posible ignorar esta distinción a nivel
teórico. Esta conclusión, lejos de ser un tecnicismo, tiene una consecuencia
metodológica concreta: estudiar un dialecto como el rioplatense no es estudiar
una versión «menor» o «derivada» del español, sino estudiar un sistema
lingüístico completo, tan legítimo y tan complejo como cualquier otro.

En cuanto a los factores históricos y geográficos, vimos que el perfil
lingüístico del rioplatense es el resultado de procesos de contacto y cambio que
se extienden a lo largo de varios siglos: desde la convivencia inicial con el
guaraní en los orígenes de Buenos Aires hasta la transformación radical que
produjo la inmigración europea masiva a fines del siglo \textsc{xix}. En el plano
del léxico, el lunfardo, con sus préstamos de diversas lenguas, se mostró como un
testimonio vivo de esa historia de contacto. En el plano fonético y fonológico,
describimos rasgos como la aspiración de /s/, el yeísmo rehilado y una
entonación de carácter «italianizante» que los hablantes de otras variedades
perciben de inmediato. En la morfología, analizamos el voseo y el uso extendido
del prefijo \forma{re-}, que ilustra cómo un elemento aparentemente menor puede
desarrollar comportamientos gramaticales complejos. Finalmente, en la sintaxis,
exploramos fenómenos como el doblado de clíticos, el doblado verbal, el doblado
negativo enfático y el uso focal del adjetivo \forma{solo}, cuyo estatus
dialectal preciso aún está por establecerse, todos los cuales revelan que las
particularidades del rioplatense no se limitan a qué palabras se usan o «cómo
suenan».

Naturalmente, mucho ha quedado fuera de este panorama. Cada uno de los temas
tratados aquí tiene una bibliografía extensa e involucra debates abiertos que no
hemos podido abordar. Nuestro propósito ha sido ofrecer un primer mapa de la
variedad, señalar sus rasgos más distintivos y, sobre todo, invitar a adentrarse
en el estudio del español rioplatense, un terreno que, a pesar de los avances
realizados, sigue ofreciendo enormes posibilidades de estudio.
\end{recuadro}

% --- Secciones finales -------------------------------------------------
\lecturasrecomendadas

\solucionario

\begin{solucion}{\ref{ej:15-weinreich}}
El aforismo de Weinreich señala con ironía que la distinción entre lengua y
dialecto no descansa en propiedades gramaticales, sino en factores políticos y de
poder estatal: lo que cuenta como «lengua» suele ser la variedad respaldada por
un Estado, mientras que las variedades sin ese respaldo institucional tienden a
ser clasificadas como «dialectos». Del mismo modo, un mismo sistema lingüístico
puede recibir diversas denominaciones y ser identificado como distintas lenguas
bajo criterios sociopolíticos. Por ejemplo, el serbocroata era la lengua que se
hablaba en Yugoslavia; si bien se reconocían variantes regionales, todas ellas
estaban sujetas a una misma norma. Tras la disolución de Yugoslavia en los años
noventa, lo que era una sola lengua se fragmentó políticamente en cuatro: serbio,
croata, bosnio y montenegrino. Las diferencias lingüísticas reales entre ellas
son comparables a las que existen entre dos variedades de español. Sin embargo,
cada nueva nación soberana necesitaba su propia lengua como símbolo de identidad
nacional, lo que llevó a la distinción actual.
\end{solucion}

\begin{solucion}{\ref{ej:15-prestamos}}
Tomando como referencia el español de México, pueden considerarse las siguientes
tres palabras: \forma{tomate, chocolate} y \forma{cenote}. \forma{Tomate}
proviene del náhuatl \forma{tomatl}; en México, designa específicamente el tomate
verde con cáscara, mientras que el tomate rojo recibe el nombre de
\forma{jitomate}. \forma{Chocolate} viene del náhuatl \forma{xocoatl}, compuesto
de \forma{xoco} (\sig{amargo}) y \forma{atl} (\sig{agua}). \forma{Cenote}, en
cambio, proviene del maya \forma{tz'ono'ot} y se traduce como \sig{caverna con
depósito de agua}; esta palabra designa los pozos naturales de agua subterránea
característicos de la península de Yucatán.
\end{solucion}

\begin{solucion}{\ref{ej:15-transcripcion}}
Para \forma{Ella mira el cielo}, una pronunciación posible en español peninsular
estándar es \fon{ˈe.ʝa.ˈmi.ɾa.el.ˈθje.lo}, mientras que en rioplatense podría ser
\fon{ˈe.ʃa.ˈmi.ɾa.el.ˈsje.lo}, en donde se observa yeísmo rehilado (en su
variante ensordecida) en \forma{ella} y seseo en \forma{cielo}. Para \forma{Es
muy callado}, la pronunciación peninsular estándar es \fon{es.muj.ka.ˈʝa.ðo} y
la rioplatense es \fon{eh.muj.ka.ˈʃa.ðo}, donde se observan dos fenómenos: la
aspiración de /s/ en coda preconsonántica (en \forma{es}) y el rehilamiento en
\forma{callado}. Para \forma{Hasta mañana}, la pronunciación peninsular estándar
sería \fon{ˈas.ta.ma.ˈɲa.na} y la rioplatense, \fon{ˈah.ta.ma.ˈɲa.na}; aquí se
observa la aspiración de /s/ en coda silábica.\pendiente{El original transcribe
\textit{cielo} con doble acento (ˈθjˈe.lo / ˈsjˈe.lo); se corrigió}
\end{solucion}

\begin{solucion}{\ref{ej:15-voseo}}
En el español de Chile predomina ampliamente el voseo verbal: si bien el
pronombre \forma{tú} es la forma pronominal habitual, las formas verbales
correspondientes al uso coloquial son muchas veces las del voseo (p.~ej.,
\forma{tú sabíh, tú veníh}); el uso del pronombre \forma{vos} se encuentra
estigmatizado y solo se atestigua en el habla subestándar. En el español de Costa
Rica, por otro lado, el voseo convive con el \textit{ustedeo}, es decir, el uso
informal y familiar del pronombre \forma{usted}; el pronombre \forma{tú}
prácticamente no se utiliza en la variedad. Como puede observarse, la variación
del voseo en diferentes dialectos del español involucra varios factores
simultáneos (forma del pronombre, morfología verbal, valor sociolingüístico de
cada combinación, etc.).
\end{solucion}

\begin{solucion}{\ref{ej:15-juicios}}
Para (2a) \forma{Le regalé el libro a María}, un hablante rioplatense la juzgará
natural; en cambio, (2b) \forma{Regalé el libro a María}, sin clítico, le
resultará claramente extraña. Para un hablante de algunas variedades
peninsulares, ambas son aceptables, aunque la versión con clítico también es muy
frecuente. La diferencia ilustra una variación microparamétrica entre dialectos:
el rioplatense ha extendido la obligatoriedad del doblado de clíticos dativos a
contextos donde algunas otras variedades admiten opcionalidad.
\end{solucion}

\begin{solucion}{\ref{ej:15-estigma}}
En el español rioplatense, el uso del determinante delante de nombres propios
(p.~ej., \forma{la Claudia le dijo al Sergio que ya no lo quería}) es objeto de
valoración negativa en ciertos contextos formales o normativos, donde se
considera «incorrecto» o propio de un registro descuidado. Sin embargo, la
construcción no tiene nada de gramaticalmente desviado: en otras lenguas y
dialectos esta misma estructura es la opción no marcada (p.~ej., el griego
clásico \textit{hē María} \sig{la María} o el portugués brasileño \forma{o Pelé}
\sig{el Pelé}). Lo que se está estigmatizando, entonces, no es una desviación
gramatical, sino una opción gramatical perfectamente sistemática que, por razones
extralingüísticas, se ha asociado con ciertos sectores sociales o registros menos
prestigiosos. Esto ilustra la idea, presentada en la introducción, de que la
distinción entre «lengua correcta» y «dialecto incorrecto» se sostiene en
factores históricos y sociales, no en propiedades gramaticales
intrínsecas.\pendiente{El original escribe el griego en alfabeto griego (ἡ
Μαρία); se transliteró}
\end{solucion}


\backmatter
% preliminares/siglario.tex
% Fuente: 02-siglario.docx (autor: Fernando Carranza)
\addchap{Siglario}
\addtocontents{toc}{\protect\autorentoc{Fernando Carranza}}
{\raggedright\normalfont\large\itshape Fernando Carranza\par}
\bigskip

A lo largo del libro, cuando se haga análisis sintáctico, se encerrarán los
constituyentes entre corchetes. Cada encorchetado estará etiquetado con una
etiqueta de la forma \textit{x/y}, de modo tal que la \textit{x} corresponderá a
la clase de palabra o tipo de sintagma (sintaxis categorial), mientras que la
\textit{y} corresponderá a la función sintáctica (sintaxis funcional). Por
ejemplo,

\begin{center}
  [\etq{SC = SF/O} [\etq{SN/Suj} Juan\etq{N/n}] [\etq{SV/Pred} vino\etq{V/n}]]
\end{center}

\noindent significa (o se lee) que \forma{Juan vino} es un sintagma
complementante / sintagma de flexión, que \forma{Juan} es un sintagma nominal
que funciona como sujeto, en el cual \forma{Juan} es un nombre que funciona como
núcleo, y que \forma{vino} es un sintagma verbal en el que \forma{vino} es un
verbo que funciona como núcleo. Para especificar tipo de categoría y función
sintáctica, se utilizan las abreviaturas que siguen.

\newcommand*{\grupositlas}[1]{\par\medskip{\sffamily\bfseries #1}\par\smallskip}
\newcommand*{\subgruposiglas}[1]{\par\smallskip{\itshape #1}\par}

\grupositlas{Etiquetas categoriales}

\subgruposiglas{Clases léxicas}
\begin{siglario}
  \item[N] Nombre / sustantivo (clase de palabra)
  \item[V] Verbo (clase de palabra)
  \item[A] Adjetivo (clase de palabra)
  \item[Adv] Adverbio (clase de palabra)
  \item[P] Preposición (clase de palabra)
\end{siglario}

\subgruposiglas{Clases funcionales}
\begin{siglario}
  \item[D] Determinante (clase de palabra, subtipo de determinativo)
  \item[Q] Cuantificador (clase de palabra, subtipo de determinativo)
  \item[F] Flexión (núcleo funcional, sintaxis categorial)
  \item[Conj] Conjunción coordinante (clase de palabra)
  \item[C] Complementante (núcleo funcional) y conjunciones subordinantes
        \textit{que} y \textit{si}
\end{siglario}

\subgruposiglas{Proyecciones sintagmáticas}
\begin{siglario}
  \item[X] Vale por cualquier categoría (clase de palabra o núcleo funcional)
  \item[X\prima] Proyección intermedia de cualquier categoría X (sintaxis
        categorial; no es necesario marcarla en el análisis sintáctico)
  \item[SX] Proyección máxima de una categoría X (sintaxis categorial)
\end{siglario}

\subgruposiglas{Categorías vacías}
\begin{siglario}
  \item[\textit{pro}] Pronombre silente valuado con caso (p.~ej., sujetos
        tácitos, sujetos expletivos silentes de verbos meteorológicos e
        impersonales, \textit{pro} arbitrario)
  \item[PRO] Pronombre silente no valuado con caso (p.~ej., sujeto de verbos en
        infinitivo)
  \item[h] Huella (la posición de base desde la cual se desplaza un objeto
        sintáctico)
\end{siglario}

\subgruposiglas{Otras etiquetas categoriales}
\begin{siglario}
  \item[Pr] Pronombre (clase de palabra)
  \item[PrEx] Pronombre exclamativo (clase de palabra)
  \item[PrInt] Pronombre interrogativo (clase de palabra)
  \item[R] Pronombre relativo (clase de palabra)
  \item[PRefl] Pronombre reflejo. Incluye los reflexivos puros y factitivos
  \item[PRecipr] Pronombre recíproco
  \item[IPCR] Índice pronominal cuasirreflejo. Incluye los pronombres
        cuasirreflejos que modifican la estructura sintáctica (alternantes e
        inherentes), los que modifican la estructura argumental (medio,
        mediopasivo, pasivo e impersonal) y los que modifican la estructura
        eventiva (delimitador e intensificador)
  \item[IPN] Índice de polaridad negativa
  \item[Ref] Reforzador (conector que refuerza a un coordinante, explícito o
        nulo; p.~ej., \textit{además, también})
  \item[Mat] Matizador (conector que matiza a un coordinante, explícito o nulo;
        p.~ej., \textit{sin embargo}). Matizadores y reforzadores son las mismas
        unidades, que funcionan como unos u otros según el coordinante al que se
        adjuntan: \textit{y además} (reforzador), \textit{pero además}
        (matizador)
  \item[PV] Perífrasis verbal (sintaxis categorial)
  \item[LV] Locución verbal
  \item[LAdv] Locución adverbial
  \item[LP] Locución preposicional
\end{siglario}

\grupositlas{Etiquetas funcionales}

\subgruposiglas{Funciones sintácticas básicas}
\begin{siglario}
  \item[n] Núcleo (esquema X con barra / sintaxis funcional)
  \item[Esp] Especificador (sintaxis funcional)
  \item[Comp] Complemento (esquema X con barra)
  \item[Adj] Adjunto
  \item[AS] Adjunto especificativo (sintaxis funcional)
  \item[AX] Adjunto explicativo (sintaxis funcional)
  \item[M] Modificador (sintaxis funcional)
  \item[ApR] Aposición restrictiva (sintaxis funcional)
  \item[ApE] Aposición explicativa (sintaxis funcional)
  \item[MF] Marca de función (sintaxis funcional). Es el caso de la \textit{a}
        personal; en estos casos, no corresponde marcar ninguna categoría
\end{siglario}

\subgruposiglas{Funciones sintácticas específicas del predicado verbal}
\begin{siglario}
  \item[COD] Complemento objeto directo (sintaxis funcional)
  \item[OI] Objeto indirecto no argumental. Incluye el dativo benefactivo, el
        dativo posesivo y el dativo ético (sintaxis funcional)
  \item[COI] Complemento objeto indirecto (sintaxis funcional)
  \item[CPR] Complemento preposicional de régimen (sintaxis funcional)
  \item[CAg] Complemento agente (sintaxis funcional)
  \item[CSim] Complemento simétrico (sintaxis funcional)
  \item[CC] Complemento circunstancial «a secas» (sintaxis funcional)
  \item[CCT] Complemento circunstancial de tiempo (sintaxis funcional)
  \item[CCL] Complemento circunstancial de lugar o complemento locativo
        (sintaxis funcional)
  \item[CCM] Complemento circunstancial de modo o complemento de manera
        (sintaxis funcional)
  \item[CCCant] Complemento circunstancial de cantidad (sintaxis funcional)
  \item[CCF] Complemento circunstancial de fin (sintaxis funcional)
  \item[AC] Adjunto circunstancial «a secas» (sintaxis funcional)
  \item[ACT] Adjunto circunstancial de tiempo (sintaxis funcional)
  \item[ACL] Adjunto circunstancial de lugar (sintaxis funcional)
  \item[ACM] Adjunto circunstancial de modo (sintaxis funcional)
  \item[ACC] Adjunto circunstancial de causa (sintaxis funcional)
  \item[ACCant] Adjunto circunstancial de cantidad (sintaxis funcional)
  \item[ACI] Adjunto circunstancial de instrumento (sintaxis funcional)
  \item[ACF] Adjunto circunstancial de fin (sintaxis funcional)
  \item[PSO] Predicativo subjetivo obligatorio (sintaxis funcional)
  \item[PSNO] Predicativo subjetivo no obligatorio (sintaxis funcional)
  \item[PO] Predicativo objetivo (sintaxis funcional)
  \item[MNO] Modificador del núcleo oracional
  \item[MMO] Modificador de modalidad oracional
  \item[MTP] Modificador temporal puntual (subtipo de circunstancial de tiempo)
  \item[MTD] Modificador temporal de duración (subtipo de circunstancial de
        tiempo)
  \item[MTPer] Modificador temporal de período (subtipo de circunstancial de
        tiempo)
  \item[MRep] Modificador temporal de repetición (subtipo de circunstancial de
        tiempo)
  \item[MFr] Modificador de frecuencia (subtipo de circunstancial de tiempo)
\end{siglario}

\subgruposiglas{Funciones sintácticas específicas del nivel de la oración}
\begin{siglario}
  \item[Suj] Sujeto (sintaxis funcional)
  \item[Pred] Predicado (sintaxis funcional)
  \item[O] Oración (sintaxis funcional)
  \item[OC] Oración subordinada completiva
  \item[OR] Oración subordinada de relativo
  \item[Ocomp] Oración compuesta (sintaxis funcional): la integrada por
        oraciones subordinadas o coordinadas
  \item[Mod.F] Modificador oracional (modificador al nivel de la capa flexiva,
        sintaxis funcional)
  \item[Mod.M] Modificador de modalidad (sintaxis funcional)
\end{siglario}

\subgruposiglas{Funciones informativas}
\begin{siglario}
  \item[Top] Tópico
  \item[Foc] Foco
\end{siglario}

\subgruposiglas{Otras etiquetas}
\begin{siglario}
  \item[ORed] Oración reducida (p.~ej., soporte para predicaciones secundarias)
  \item[\textit{i, j, k}] Índices de correferencialidad
  \item[\{\,/\,\}] Delimitación de opciones alternativas (p.~ej., \textit{tener
        \{nada / algo\} que ver})%
        \pendiente{Agregado a partir de la nota del original «AGREGAR LLAVES
        PARA DELIMITACIÓN DE OPCIONES»: revisar la redacción}
\end{siglario}

% =====================================================================
%  finales/glosario.tex — Glosario general (dos columnas)
%  Fuente: 18-glosario.docx (pendiente de conversión)
% =====================================================================
\addchap{Glosario}

\begin{glosariogeneral}
  \item[Aceptabilidad:] Condición de aceptación de las estructuras de una
        lengua. [Entrada de prueba: reemplazar con el glosario.]
  \item[Distribución]\subterminos
    \begin{glosario}
      \item[complementaria:] Entrada de prueba.
      \item[contrastiva:] Entrada de prueba.
    \end{glosario}
  \item[Gramaticalidad:] Condición de buena formación de las estructuras de
        una lengua. [Entrada de prueba.]
\end{glosariogeneral}

\bibliografiageneral

\end{document}
